\documentclass[10pt,a4paper]{article}

% ----------------- Paquetes -------------------%
\usepackage[T1]{fontenc}
\usepackage[utf8]{inputenc}
\usepackage[a4paper,margin=2.5cm]{geometry}
\usepackage{mathptmx}             
\usepackage{changepage} 
\usepackage{setspace}
\usepackage[spanish]{babel}
\usepackage[labelfont=bf,labelsep=space]{caption}
\usepackage{amssymb}
\usepackage{amsmath}
\usepackage{ebgaramond}
\usepackage{placeins}
\usepackage{etoolbox}
\usepackage{setspace}
\usepackage{parskip} 
\usepackage{tcolorbox}
\usepackage{needspace}
\usepackage{xparse}
\usepackage{ocgx2}
\usepackage{hyperref}
\usepackage{hypcap}
\usepackage{soul}\sethlcolor{yellow}% resaltado amarillo: \hl{texto} (Panel TCEE, ⌃H)


%%%%%%%%%%%%%%%%%%%%%%%%%%%%%%%%%%%%%%%%%%%%%%%%%%%%%%%%%%%%%%%%
% --- Fija primero tus valores globales "buenos" ---
\setstretch{1.2}                 % Interlineado global
\setlength{\parskip}{0.7\baselineskip}
\setlength{\parindent}{0pt}
\newlength{\GlobalParskip}
\setlength{\GlobalParskip}{\parskip}

\newcounter{eqnum}[subsection]
\renewcommand{\theeqnum}{\arabic{eqnum}}
\newcounter{alphaitem}
\newcounter{numitem}

%% ------------------- Recuadros----------------------%%
\newcommand{\puntosclave}[1]{%
  \begin{tcolorbox}[
    colframe=black,
    title={Puntos clave},
    before upper={\setlength{\parindent}{0.5cm}\setlength{\parskip}{0.5\baselineskip}}
  ]
  #1
  \end{tcolorbox}
}

\newcommand{\modificaciones}[1]{
  \begin{tcolorbox}[
    colback=white,
    colframe=red,
    title={Modificaciones a realizar},
    before upper={\setlength{\parindent}{0.5cm}\setlength{\parskip}{0.5\baselineskip}}
  ]
  #1
  \end{tcolorbox}
}

%% ------------------- Listas----------------------%%
    % --- Lista alfabética ---
    \newenvironment{la}
      {\begin{list}{\alph{alphaitem})}{%
          \usecounter{alphaitem}%
          \setlength{\leftmargin}{1cm}%
          \setlength{\parskip}{3pt}% 
          \setlength{\itemsep}{0pt}% ← espacio entre ítems
          \setlength{\parsep}{0pt}%  ← párrafos dentro de un ítem
      }}
      {\end{list}}
      
    % --- Lista numérica ---
    \newenvironment{lnum}
      {\begin{list}{\arabic{numitem}.}{%
          \usecounter{numitem}%
          \setlength{\leftmargin}{1cm}%
          \setlength{\parskip}{3pt}% 
          \setlength{\itemsep}{0pt}% ← espacio entre ítems
          \setlength{\parsep}{0pt}%  ← párrafos dentro de un ítem
      }}
      {\end{list}}
      
% ------------------- Imágenes----------------------%
    \graphicspath{{Img/}}% Rutas por defecto 
    % --- Imagen adaptativa ---
    % Registros de dimensión definidos una sola vez
    \newdimen\imgcap
    \newdimen\imgmin
    \newdimen\imgmax
    \newdimen\imgremain
    \newdimen\imgh
    \newdimen\imgneed
    
    \newcommand{\imagenfit}[3]{%
      \addvspace{10pt}% separación antes
      \par\begingroup
        % Parámetros de composición (asignaciones, no \newdimen)
        \imgcap=1\baselineskip            % alto estimado de título+fuente
        \imgmin=0.15\textheight
        \imgmax=0.15\textheight
        \imgremain=\pagegoal
        \ifdim\pagegoal=\maxdimen \imgremain=0pt\fi
        \advance\imgremain by -\pagetotal
        \advance\imgremain by -\imgcap
        % Decisión de altura destino
        \ifdim\imgremain<\imgmin
          \newpage
          \imgh=0.20\textheight
        \else
          \imgh=\imgremain
          \ifdim\imgh>\imgmax \imgh=\imgmax \fi
        \fi
        % Reservar espacio para evitar cortes del bloque completo
        \imgneed=\imgh \advance\imgneed by \imgcap
        \Needspace{\imgneed}
        % Cuerpo
        \noindent\begin{minipage}{\textwidth}
          \centering
          \captionof{figure}{\textemdash\ #2}\par\vspace{0.3em}% título arriba
          \includegraphics[width=\linewidth,height=\imgh,keepaspectratio]{\detokenize{#1}}\par
          \vspace{0.3em}\small\textit{Fuente: }#3% fuente abajo
        \end{minipage}%
      \endgroup
      \par\addvspace{10pt}%
    }

%------------ Bloque de RECUADRO ESPECIAL -------------%
    \newenvironment{specialblock}[1]{%
      \begin{quote} % crea sangría a izquierda y derecha
      \small % texto más pequeño
      \noindent\textbf{#1}\\[-0.3em] % título en negrita
      \rule{\linewidth}{0.4pt}\\[0.6em]}{
      \end{quote}}
      

%-------------------------- Author Font -----------------------%
    \newcommand{\authorfont}[1]{{\fontfamily{EBGaramond-TLF}\selectfont #1}}

%-------------------------- Anexos -----------------------%
    \makeatletter
    \newcounter{anexo}
    \renewcommand{\theanexo}{\Roman{anexo}}
    \newcommand{\anexo}[1]{%
      \clearpage
      \phantomsection
      \refstepcounter{anexo}%
      \noindent\textbf{Anexo \theanexo.- }#1\par\medskip
      \addcontentsline{stc}{section}{Anexo \theanexo.- #1}%
    }
    \makeatother

    %-------------------------- Lista de figuras (personal) -----------------------%
\makeatletter
\newcommand{\MyListOfFigures}{%
  \clearpage
  \phantomsection
  {\centering\bfseries\Large Índice de figuras\par}%
  \medskip
  % Entrada en nuestro índice de contenido (stc)
  \addcontentsline{stc}{section}{Índice de figuras}%
  % Recuperación del listado (sin cabecera propia)
  \begingroup
    \parindent=0pt\parskip=2pt
    \@starttoc{lof}%
  \endgroup
}
\makeatother

%-------------------------- Bibliografía -----------------------%
\makeatletter
% Contador para las referencias
\newcounter{myref}
% Cabecera de la bibliografía, entra en el índice 'stc'
\newcommand{\MyBibliography}{%
  \clearpage
  \phantomsection
  {\centering\bfseries\Large Bibliografía y referencias\par}%
  \medskip
  \addcontentsline{stc}{section}{Bibliografía y referencias}%
  \setcounter{myref}{0}%
}
\newcommand{\MyBibItem}[4]{%
  \refstepcounter{myref}%
  \noindent[\themyref]\ #1.\ \emph{#2}. #3. #4\par
}
\makeatother

%--------- Establecer cuatro niveles de índice --------%
    \setcounter{tocdepth}{4}   
    \setcounter{secnumdepth}{4}% numera hasta \paragraph

%%%%%%%%%%%%%%%%%%%%%%%%%%%%%%%%

\makeatletter

    %--------- Interlineado Global --------%
    \edef\GlobalBaseLineStretch{\baselinestretch}
    
    %--------- Espaciado de seccinones --------%
    \renewcommand\section{\@startsection{section}
      {1}{\z@}
      {2.5ex \@plus 1ex \@minus .2ex} % espacio antes
      {1.5ex \@plus .2ex}             % espacio después
      {\normalfont\Large\bfseries}    % formato del título
    }
    \renewcommand\subsection{\@startsection{subsection}{2}{\z@}
      {3ex \@plus 0.8ex \@minus .2ex}
      {1.5ex \@plus .2ex}
      {\normalfont\large\bfseries}
    }
    \renewcommand\subsubsection{\@startsection{subsubsection}{3}{\z@}
      {3ex \@plus 0.5ex \@minus .2ex}
      {1ex \@plus .2ex}
      {\normalfont\normalsize\bfseries}
    }
    \renewcommand\paragraph{\@startsection{paragraph}{4}{\z@}%
    {-2ex \@plus -0.5ex \@minus -.2ex}% espacio antes
    {0.8ex \@plus .2ex}% espacio después
    {\normalfont\normalsize\itshape}}% ← cursiva en lugar de negrita

    %--------- Ecuaciones --------%   
    % Variables globales para almacenar títulos y notas
    \gdef\leq@current@section{}%
    \gdef\leq@current@subsection{}%
    \gdef\leq@last@printed@section{}%
    \gdef\leq@last@printed@subsection{}%
    
    % Comando para añadir nota (invisible en el cuerpo)
    \newcommand{\eqnote}[1]{%
      \addcontentsline{leq}{leqnote}{#1}%
    }
    
    % Comando para ecuaciones
    \newcommand{\eqblock}[2]{%
      % Verificar si necesitamos imprimir el título de sección
      \ifx\leq@current@section\leq@last@printed@section\else
        \addcontentsline{leq}{leqsection}{\leq@current@section}%
        \global\let\leq@last@printed@section\leq@current@section
        \gdef\leq@last@printed@subsection{}% Reset subsección al cambiar sección
      \fi
      % Verificar si necesitamos imprimir el título de subsección
      \ifx\leq@current@subsection\leq@last@printed@subsection\else
        \ifx\leq@current@subsection\@empty\else
          \addcontentsline{leq}{leqsubsection}{\leq@current@subsection}%
          \global\let\leq@last@printed@subsection\leq@current@subsection
        \fi
      \fi
      % Ecuación
      \refstepcounter{eqnum}%
      \begin{equation}
        #1\tag{E.\theeqnum}
      \end{equation}%
      \addcontentsline{leq}{leqt}{[E.\theeqnum] -- #2}%
      \addcontentsline{leq}{leqm}{#1}%
    }
    
    %%% ----- Índice de ecuaciones ----------%%%
    % Tamaño para []
    \newcommand{\leqIndexFont}{\normalsize}
    \newcommand{\leqTitleFont}{\tiny}
    \newcommand{\leqNoteFont}{\tiny}
    % Cabecera de sección 
    \newcommand\l@leqsection[2]{%
      \addvspace{25pt}% Mayor separación antes de nueva sección
      \noindent\leqIndexFont
      \makebox[\linewidth]{\rule{.38\linewidth}{0.4pt}\ \ \textbf{  \MakeUppercase{#1}}\ \ \rule{.38\linewidth}{0.4pt}}%
      \par\vspace{-10pt}
    }
    % Cabecera de subsección 
    \newcommand\l@leqsubsection[2]{%
      \par\addvspace{15pt}%
      \noindent\leqIndexFont #1
      \par\vspace{-7pt}
    }
    % Título de ecuación 
    \newcommand\l@leqt[2]{%
    \par\addvspace{10pt}%
      \par\noindent\leqTitleFont #1\par
    }
    % Miniatura de ecuación
    \newcommand\l@leqm[2]{%
      \vspace{0.3\baselineskip}%
      \par\scriptsize\noindent\hspace{1.5em}\(#1\)\par%
    }
    % Nota de ecuación 
    \newcommand\l@leqnote[2]{%
      \vspace{0.2\baselineskip}%
      \par\noindent\leqNoteFont\hspace{1.5em}\textbf{Nota.--} #1\par\vspace{0.5\baselineskip}%
    }
    % Generación del índice
    \newcommand{\mylistofequations}{%
      \clearpage
      \phantomsection
      % Título visual de la lista (ajústalo a tu gusto):
      {\centering\bfseries\Large Índice de ecuaciones\par}
      % Entrada en nuestro índice 'stc':
      \addcontentsline{stc}{section}{Índice de ecuaciones}%
      % Llamada a tu lista real (usa el nombre exacto de tu macro):
      \@starttoc{leq}%
    }
    
    %--------- Captura de títulos de sección --------%
    \let\leq@original@section\section
    \let\leq@original@subsection\subsection
    % Flag para saber si estamos en una sección asterisco
    \newif\ifLeqStarSection
    \LeqStarSectionfalse
    
    % Redefinir \section CON CONTROL DE ASTERISCO
    \renewcommand{\section}{%
      \@ifstar{\leq@section@star}{\@dblarg\leq@section@nostar}%
    }
    \def\leq@section@star#1{%
      \global\LeqStarSectiontrue
      \gdef\leq@current@section{#1}%
      \gdef\leq@current@subsection{}%
      \leq@original@section{#1}%
      \global\LeqStarSectionfalse
    }
    \def\leq@section@nostar[#1]#2{%
      \global\LeqStarSectionfalse
      \gdef\leq@current@section{#2}%
      \gdef\leq@current@subsection{}%
      \leq@original@section[#1]{#2}%
    }
    % Redefinir \subsection
    \renewcommand{\subsection}{%
      \@ifstar{\leq@subsection@star}{\@dblarg\leq@subsection@nostar}%
    }
    \def\leq@subsection@star#1{%
      \global\LeqStarSectiontrue
      \gdef\leq@current@subsection{#1}%
      \leq@original@subsection{#1}%
      \global\LeqStarSectionfalse
    }
    \def\leq@subsection@nostar[#1]#2{%
      \global\LeqStarSectionfalse
      \gdef\leq@current@subsection{#2}%
      \leq@original@subsection[#1]{#2}%
    }

    % ----- Restauración de valores Globales de estilo ----------%
    % Flags
    \newif\ifInFootnote
    \InFootnotefalse
    \pretocmd{\@footnotetext}{\InFootnotetrue}{}{}
    \apptocmd{\@footnotetext}{\InFootnotefalse}{}{}
    % Profundidad de listas (soporta anidadas)
    \newcount\MyListDepth \MyListDepth=0
    \AtBeginEnvironment{la}{\advance\MyListDepth by 1}
    \AfterEndEnvironment{la}{\advance\MyListDepth by -1}
    \AtBeginEnvironment{lnum}{\advance\MyListDepth by 1}
    \AfterEndEnvironment{lnum}{\advance\MyListDepth by -1}
    \AtBeginEnvironment{itemize}{\advance\MyListDepth by 1}
    \AfterEndEnvironment{itemize}{\advance\MyListDepth by -1}
    \AtBeginEnvironment{enumerate}{\advance\MyListDepth by 1}
    \AfterEndEnvironment{enumerate}{\advance\MyListDepth by -1}
    % Restauración diferida: hazlo al INICIO del próximo párrafo real
    \newif\if@pendingRestore
    \@pendingRestorefalse
    \newcommand{\RequestRestore}{\global\@pendingRestoretrue}
    \everypar{%
      \if@pendingRestore
        \global\@pendingRestorefalse
        \ifInFootnote\else
          \ifnum\MyListDepth=0
            \setlength{\parskip}{\GlobalParskip}%
            \setstretch{\GlobalBaseLineStretch}%
          \fi
        \fi
      \fi
    }
    % Al final de bloques:
    \AfterEndEnvironment{equation}{\RequestRestore}
    \AfterEndEnvironment{equation}{\RequestRestore}
    \AfterEndEnvironment{la}{\RequestRestore}
    \AfterEndEnvironment{lnum}{\RequestRestore}
    % Al final de macros:
    \apptocmd{\eqblock}{\RequestRestore}{}{}
    \apptocmd{\imagenfit}{\RequestRestore}{}{}
    
\makeatother

% ===================== ToC propio con 4 niveles (único bloque) =====================
\makeatletter

% --- Escritores: añadimos una línea al .stc en cada cabecera numerada ---
% Guardamos las versiones actuales para llamar al comportamiento normal
\let\MyTOC@orig@section\section
\let\MyTOC@orig@subsection\subsection
\let\MyTOC@orig@subsubsection\subsubsection
\let\MyTOC@orig@paragraph\paragraph

% SECTION
\renewcommand{\section}{%
  \@ifstar{\MyTOC@section@star}{\@dblarg\MyTOC@section@nostar}%
}
\def\MyTOC@section@star#1{%
  \MyTOC@orig@section{#1}% sin entrada al .stc
}
\def\MyTOC@section@nostar[#1]#2{%
  \MyTOC@orig@section[{#1}]{#2}% comportamiento normal (escribe al .toc)
  % Escribimos también nuestra línea al .stc (texto con \numberline)
  \addcontentsline{stc}{section}{\protect\numberline{\thesection}#2}%
}

% SUBSECTION
\renewcommand{\subsection}{%
  \@ifstar{\MyTOC@subsection@star}{\@dblarg\MyTOC@subsection@nostar}%
}
\def\MyTOC@subsection@star#1{%
  \MyTOC@orig@subsection{#1}%
}
\def\MyTOC@subsection@nostar[#1]#2{%
  \MyTOC@orig@subsection[{#1}]{#2}%
  \addcontentsline{stc}{subsection}{\protect\numberline{\thesubsection}#2}%
}

% SUBSUBSECTION
\renewcommand{\subsubsection}{%
  \@ifstar{\MyTOC@subsubsection@star}{\@dblarg\MyTOC@subsubsection@nostar}%
}
\def\MyTOC@subsubsection@star#1{%
  \MyTOC@orig@subsubsection{#1}%
}
\def\MyTOC@subsubsection@nostar[#1]#2{%
  \MyTOC@orig@subsubsection[{#1}]{#2}%
  \addcontentsline{stc}{subsubsection}{\protect\numberline{\thesubsubsection}#2}%
}

% PARAGRAPH
\renewcommand{\paragraph}{%
  \@ifstar{\MyTOC@paragraph@star}{\@dblarg\MyTOC@paragraph@nostar}%
}
\def\MyTOC@paragraph@star#1{%
  \MyTOC@orig@paragraph{#1}%
}
\def\MyTOC@paragraph@nostar[#1]#2{%
  \MyTOC@orig@paragraph[{#1}]{#2}%
  % Si no numeras los \paragraph, cambia \theparagraph por texto plano sin \numberline
  \addcontentsline{stc}{paragraph}{\protect\numberline{\theparagraph}#2}%
}

% --- Impresor del índice propio (lee el .stc) ---
\newcommand{\MyTOC}{%
  \par\bigskip
  {\centering\bfseries\Large Índice de contenido\par}%
  \medskip
  \begingroup
    \parindent=0pt\parskip=2pt
    \@starttoc{stc}%
  \endgroup
}

\makeatother
% ===================== fin ToC propio =====================


%%%%%%%%%%%%%%


% --- Datos del tema ---
\title{Política Presupuestaria (IV): Sistema de Seguridad Social en España\\
\large Prestaciones y su financiación. La sostenibilidad del sistema público de pensiones}
\author{Marcelo Ortega}
\date{Septiembre 2026}

\begin{document}
\maketitle
\newpage

% --- Página de resumen y modificaciones ---
\puntosclave{ %% Escribir puntos clave Apuntes CECO
     La distribución aconsejada de tiempos es la siguiente: Introducción, un minuto; Desarrollo histórico de la Seguridad Social, dos minutos; Estructura de la Seguridad Social, ocho minutos; Sistema de pensiones, dieciséis minutos; el gasto en protección social en España en el contexto europeo: dos minutos y conclusiones: 1 minuto.
    }
\vspace{1cm}

\modificaciones{
    
    }

\newpage
\MyTOC
\newpage

\mylistofequations
\clearpage

\MyListOfFigures
\clearpage


\pagenumbering{arabic}
\setcounter{page}{1}


\section{Introducción}



\subsection{Contextualización}


\subsubsection{Relevancia}




\subsubsection{Problemática}



\subsection{Estructura de la exposición}





Este tema analiza la estructura, financiación y prestaciones de la Seguridad Social en España. Abordaremos sistemas asociados a la protección social, con especial énfasis en el Sistema de Pensiones.

En la introducción debe realizarse una breve descripción de qué se entiende por Seguridad Social y cómo se estructura en España, haciendo mención a la doble financiación a contribuciones sociales (parte contributiva) y vía impuestos (parte no contributiva). En la descripción de la evolución histórica de la Seguridad Social en España, deben describirse brevemente los periodo, sin olvidar la Ley de Bases de 1963 y los desarrollos legislativos durante la democracia (especialmente el Pacto de Toledo).

Sobre la estructura de la Seguridad Social, deben mencionarse: los principios rectores y los grandes números de los presupuestos. Debe exponerse claramente cuáles son las principales fuentes de financiación (cotizaciones y transferencias del Estado), así como las mayores prestaciones (pensiones, subsidios y otras prestaciones). Sobre el sistema de pensiones. después de unas breves notas sobre las características del sistema (que se pueden comparar con los países de nuestro entorno, haciendo mención a algunos de los datos que se presentan en el epígrafe 4.5), es importante describir las últimas reformas del mismo, es decir. los cambios paramétricos de la reforma de 2011. los dos mecanismos de ajuste introducidos en la reforma de 2013 y la última reforma llevada a cabo en 2021 y 2023.

A continuación, puede introducirse el debate de la sostenibilidad del sistema. Para ello, es importante explicar bien cuáles son los factores que explican el gasto en pensiones. perfilar su evolución reciente y describir cómo las últimas reformas afectan a la evolución futura. En este sentido. sería deseable ofrecer unas pinceladas del ahorro en términos de gasto en pensiones asociado a la reforma de 2011, el coste por a la derogación de la reforma de 2013 y las medidas adoptadas en 2021 y 2023 para contrarrestar dicha derogación (mecanismo de equidad intergeneracional, aumento de las transferencias del Estado a la Seguridad Social, aumento de la base máxima de cotización. recargo de solidaridad sobre los salarios por encima de la base máxima e incremento de los incentivos a posponer la edad de jubilación).

Tras el sistema de pensiones. se pueden describir brevemente el resto de sistemas asociados a la protección social. Una manera de exponer cada sistema es la siguiente: definición. organización, prestaciones. financiación y algunos datos sobre la estructura de gastos e ingresos. L<b conclusiones pueden resumir brevemente los puntos más importantes expuestos anteriormente y contextualizar el gasto en protección social en España con respecto a los países de nuestro entorno. utilizando los datos contenidos en la sección 6.




------------ Introducción

La Seguridad Social, de acuerdo con la Asociación Internacional de la Seguridad Social, se define como todo régimen o programa establecido por Ley, o por cualquier otra disposición obligatoria, que garantiza una protección, sea en metálico o en especie, en caso de accidentes de trabajo, enfermedades profesionales, desempleo, maternidad, enfermedad común, invalidez, vejez, jubilación, supervivencia o muerte, e incluye, entre otros, prestaciones por hijos y por otros miembros de la familia, prestaciones de salud, prevención, rehabilitación y cuidados de larga duración. Siendo esta una definición amplia de la Seguridad Social, en cada país tiene un significado distinto y las prestaciones asociadas a ella adoptan también formas y alcances muy diferentes.

Los sistemas de Seguridad Social tienen su origen en Alemania a finales del siglo XIX, cuando este país adoptó, por primera vez, un programa de seguro social de enfermedad y vejez para los trabajadores. Sin embargo, sería en Reino Unido, a partir del Informe BEVERIDGE de 1942, cuando se crearía el primer sistema unificado de Seguridad Social, abriendo paso a los sistemas de Seguridad Social modernos. El derecho a la misma aparece recogido en la Declaración Universal de los Derechos Humanos, en su art. $22$. 

En España, su origen se remonta también a finales del siglo XIX, aunque no es hasta la década de los sesenta del pasado siglo cuando se implanta un modelo integrado de protección social. Actualmente, cubre una serie de contingencias relacionadas con el envejecimiento, la supervivencia, la enfermedad, la dependencia y el desempleo. El sistema se organiza tanto desde un punto de vista contributivo o bismarckiano (donde las prestaciones guardan relación con las contribuciones), como de uno no contributivo o de tipo BEVERIDGE (basado en prestaciones de tipo asistencial que no dependen de las contribuciones realizadas). La parte contributiva del sistema se financia a través de las contribuciones sociales, mientras que la parte no contributiva se financia a través de la imposición general. Además, existe un nivel complementario de protección social que es libre, es decir, que puede ser de carácter privado.

En este capítulo se analiza el sistema público de Seguridad Social en España. Por un lado, de describe brevemente la evolución del sistema desde sus orígenes hasta su situación actual. Por otro lado, se analizan los pilares del sistema: el sistema de pensiones, el Sistema Nacional de Empleo, el Sistema Nacional de Salud y el Sistema para la Autonomía y Atención a la Dependencia. Por último, se analiza sucintamente el nivel de gasto en protección social en España en comparación con los países europeos.



\clearpage
\section{Seguridad Social: Marco jurídico-conceptual}
\puntosclave{
    El origen de las políticas de protección social en España data de finales del siglo XIX. La Ley General de la Seguridad Social de 1966 implantó efectivamente un Sistema de Seguridad Social en España. 
    
    Con la llegada de la democracia y la aprobación de la Constitución de 1978. se pone en marcha la configuración actual dd sistema de la Seguridad Social.
    
    Desde entonces, se han sucedido varias reformas de la Seguridad Social relacionadas con su gestión. su financiación y el alcance y el cálculo de las prestaciones que provee. Asimismo. se han producido reformas en el ámbito de la sanidad y el empico.

    El tamaño del gasto social en España se sitúa algo por encima de la media de europea. El gasto en salud es similar a la media de la Unión, mientras que el gasto en protección social se sitúa ligeramente por encima, debido principalmente al mayor gasto en pensiones de supervivientes y en prestaciones por desempleo.
    
    Pueden identificarse distintos grupos de países de acuerdo con sus niveles de gasto social. Los países nórdicos y algunos de Europa continental y del sur de Europa lideran los niveles de protección social. España se encontraría dentro del grupo de países con un gasto social medio, mientras que buena parte de los países de Europa del Este mostrarían niveles de gasto social bajo.

    Los principios rectores que sustentan el Sistema de la Segundad Social en España son la contributividad, la universalidad, la solidaridad intergeneracional, la equidad, la suficiencia y la unidad de caja.
    
    La gestión de la Seguridad Social se lleva a cabo a través de un entramado de actores, públicos y privados, con personalidad jurídica propia y dependientes, en último término, del Ministerio de Inclusión, Seguridad Social y Migraciones. 
    
    Desde el punto de vista presupuestario, la principal fuente de financiación de la Seguridad Social son las cotizaciones sociales y, a gran distancia, las transferencias del Estado, cuyo objeto es financiar las prestaciones que tienen carácter no contributivo.
    
    Por el lado de los gastos, la principal partida corresponde a las pensiones contributivas. Existen, además, otras prestaciones contributivas tales como la prestación por incapacidad temporal, los complementos a mínimos y el IMV.
}


\textcolor{blue}{Reunidas las piezas, es posible situar el sistema español en el mapa de los regímenes de bienestar con cierta precisión. España combina un núcleo bismarckiano en las pensiones y en la protección por desempleo, heredado del sistema de seguros sociales y consolidado por la Ley General de la Seguridad Social de 1966; una sanidad beveridgiana desde 1986; una tradición familiarista en los cuidados, corregida solo parcialmente con la creación del Sistema para la Autonomía y Atención a la Dependencia en 2006; una red asistencial tardía y fragmentada, construida por capas sucesivas (pensiones no contributivas en 1990, rentas mínimas autonómicas desde finales de los ochenta, IMV en 2020); y un dualismo laboral persistente, que se proyecta sobre la protección social. En términos de ESPING-ANDERSEN, es un régimen de raíz conservadora con los rasgos característicos del modelo meridional de FERRERA, en tránsito hacia un modelo más universalista. En términos de PALIER, ha recorrido con fidelidad las dos primeras fases de la reforma bismarckiana y ha retrocedido en la tercera. 

Esa trayectoria, con sus avances y sus retrocesos, es la que se reconstruye históricamente en el subapartado siguiente, y es también la que explica la composición singular de su gasto social: elevado en pensiones y desempleo, intermedio en sanidad y bajo en familia, vivienda y exclusión social.}

\subsection{Fundamentos: ¿Por qué es necesario establecer un Sistema de Seguridad Social?}
Antes de describir cómo está organizado el sistema español, hay que responder a una pregunta que los manuales suelen dar por resuelta: por qué el Estado obliga a los ciudadanos a asegurarse frente a la vejez, la invalidez, la muerte, la enfermedad o el desempleo, y por qué, además, organiza él mismo ese seguro en lugar de dejarlo al mercado. 

En una economía de mercado, la carga de la prueba la tiene la intervención. Nada impide a un trabajador ahorrar para su jubilación, contratar un seguro de vida o de invalidez o comprar al jubilarse una renta vitalicia a una aseguradora. Si el Estado le obliga a entregar en torno a un tercio de su salario en cotizaciones (España), sumando las partes de empresa y trabajador, y le promete a cambio unas prestaciones cuya cuantía fija la ley y no un contrato, esa obligación necesita una justificación. De hecho, es la justificación de la mayor intervención pública de la economía española por volumen de recursos.

La literatura ha dado tres grandes familias de respuestas. La primera es de \textbf{eficiencia}. Los mercados privados de seguros y de ahorro para la vejez fallan por razones informativas, por la naturaleza agregada de ciertos riesgos y por los sesgos de conducta de los individuos. Si fallan, la intervención puede mejorar el bienestar de todos, incluso sin propósito redistributivo alguno. La segunda es de \textbf{equidad}. Aunque esos mercados funcionaran perfectamente, una sociedad puede querer que la protección no dependa de la capacidad de pago de cada uno, bien porque considera la seguridad económica un derecho de ciudadanía, bien porque, puesta en la posición de quien todavía no sabe qué lugar ocupará en la distribución de la renta, elegiría asegurarse contra la mala fortuna. La tercera es \textbf{institucional y comparada}. Los argumentos anteriores no determinan un diseño único, y cada sociedad ha combinado seguro y redistribución de una forma históricamente contingente, dando lugar a regímenes de bienestar distintos cuya lógica ayuda a entender dónde se sitúa España y hacia dónde ha evolucionado.

Conviene, antes de entrar en esos argumentos, desactivar una confusión que contamina casi todo el debate público sobre pensiones. BARR y DIAMOND (2008) insisten en separar tres preguntas que suelen tratarse como una sola. La primera es si la protección debe ser obligatoria: la \textbf{pregunta del mandato}. La segunda es si debe gestionarla el sector público o puede confiarse a entidades privadas bajo regulación: la \textbf{pregunta de la provisión}. La tercera es si debe financiarse por reparto, con las cotizaciones de los activos pagando las pensiones de los jubilados, o por capitalización, con cada generación acumulando un fondo que financia sus propias prestaciones: la \textbf{pregunta de la financiación}. Las tres respuestas son lógicamente independientes, y la experiencia comparada ofrece casi todas las combinaciones posibles.\footnote{%
    Chile implantó en 1981 un sistema con mandato, provisión privada y capitalización. El Reino Unido, con la adhesión automática a planes de empresa desde 2012, ha construido algo parecido a un cuasimandato con provisión privada y capitalización. Suecia combina desde finales de los noventa un pilar de reparto con cuentas nocionales, gestionado públicamente, con un pilar obligatorio capitalizado en el que el individuo elige gestor. España tiene mandato, provisión pública y reparto.
} Y, como es natural, la consecuencia metodológica es importante para lo que sigue: casi todos los argumentos justifican el mandato; algunos justifican un papel público en la provisión o en la garantía; y muy pocos dicen algo concluyente sobre la forma de financiación, que es una cuestión distinta que analizaremos más adelante.

Además, BARR y DIAMOND distinguen en un sistema de pensiones dos tipos de objetivos: (i) \textbf{primarios}, de naturaleza individual, como alisar el consumo a lo largo del ciclo vital y asegurarse frente a riesgos como la longevidad, la invalidez o la muerte del sustentador; y, (ii) \textbf{secundarios}, de naturaleza colectiva, como el alivio de la pobreza en la vejez y la redistribución de la renta, tanto dentro de cada generación como entre generaciones. Esta distinción no es académica pues, como veremos en el último bloque, la tensión entre suficiencia y sostenibilidad, eje del debate español desde 2013, es en el fondo una tensión entre estos objetivos: cuanto más se exige al sistema que redistribuya y garantice mínimos, más se aleja de la lógica actuarial que protege su equilibrio financiero; y, cuanto más se acerca a esa lógica, más se debilita su capacidad para combatir la pobreza de quienes tuvieron carreras laborales cortas o mal pagadas.

\subsubsection{Razones de eficiencia: Fallos de mercado y de conducta}
El argumento de eficiencia parte de una idea sencilla: si los mercados privados de seguros y de ahorro para la vejez funcionaran como los mercados competitivos del manual, cada individuo podría elegir el nivel de protección que desea al precio que refleja su coste. La intervención pública solo reduciría el bienestar, porque impondría a todos una protección uniforme que para unos sería excesiva y para otros insuficiente. Pero la teoría económica de la información y la economía conductual han identificado razones poderosas por las que esos mercados no funcionan así. Cada una de ellas justifica la intervención de un modo distinto, con un alcance distinto y con límites distintos.

\paragraph{La selección adversa y el mercado imposible de las rentas vitalicias}
En el ámbito de las pensiones, el riesgo relevante es el de longevidad, es decir, el riesgo de vivir más de lo esperado y agotar los ahorros.\footnote{%
    El punto de partida es el célebre art. de AKERLOF (1970) sobre el mercado de coches usados, pero la formulación decisiva para el seguro es la de ROTHSCHILD y STIGLITZ (1976). Cuando el asegurado conoce su propio riesgo mejor que la aseguradora y esta no puede distinguir a los asegurados de alto y bajo riesgo, un mercado competitivo de seguros puede no tener equilibrio. Si lo tiene, es un equilibrio separador en el que los individuos de bajo riesgo solo pueden adquirir una cobertura parcial, inferior a la que desearían a un precio justo. La información asimétrica no produce solo una ineficiencia marginal: puede destruir el mercado o reducirlo a una fracción de su tamaño eficiente.
} El instrumento de mercado que lo cubre es la renta vitalicia: el individuo entrega un capital a una aseguradora y esta le paga una renta mientras viva. Pero la longevidad es exactamente el tipo de riesgo sobre el que el asegurado tiene información privada: conoce su estado de salud, sus hábitos, la longevidad de sus padres y abuelos, y su propia percepción de cómo se encuentra. Quien espera vivir mucho tiene un incentivo alto para comprar una renta vitalicia; quien sabe que su esperanza de vida es corta, un incentivo bajo o nulo.

Un ejemplo sencillo muestra el mecanismo en toda su crudeza. Supongamos una población de jubilados de sesenta y cinco años dividida a partes iguales en dos tipos: los que vivirán quince años más y los que vivirán veinticinco. Ignorando por simplicidad la tasa de descuento, una renta vitalicia de mil euros anuales cuesta a precio justo quince mil euros para el primer grupo y veinticinco mil para el segundo. Si la aseguradora no puede distinguir a unos de otros, cobrará el precio medio, veinte mil euros. Pero a ese precio los jubilados del primer grupo pagan veinte mil euros por algo que para ellos vale quince mil, y la mayoría se retira del mercado. Cuando lo hacen, la aseguradora descubre que su cartera está compuesta solo por longevos y eleva el precio a veinticinco mil euros. El mercado se ha desenredado hasta cubrir solo al grupo de alto riesgo, y los jubilados del primer grupo, que habrían deseado asegurarse a su precio justo, quedan sin cobertura. Un seguro obligatorio con prima única para todos rompe esta espiral por la vía más directa: al impedir la salida de los asegurados de bajo riesgo, mantiene a todos dentro del mismo grupo y hace viable un precio medio que, sin obligación, el mercado no sostendría.

FINKELSTEIN y POTERBA (2004), con datos de una gran aseguradora británica, mostraron que quienes eligen rentas vitalicias con más peso de la supervivencia, por ejemplo con pagos crecientes en el tiempo, viven efectivamente más que quienes eligen productos con más peso de los primeros años o con garantías de devolución en caso de fallecimiento temprano. MITCHELL, POTERBA, WARSHAWSKY y BROWN (1999) calcularon el llamado \textit{money's worth} de las rentas vitalicias estadounidenses, es decir, el valor actual de los pagos esperados por cada dólar de prima. Para un individuo representativo de la población general, ese valor era claramente inferior a la unidad. Para los compradores efectivos, más longevos que la media, se acercaba más a la unidad. La diferencia mide aproximadamente el coste de la selección adversa, al que se suman los gastos de gestión y los márgenes de las aseguradoras.

Esta evidencia conecta con uno de los enigmas clásicos de la economía de pensiones. YAARI (1965) demostró que un individuo racional, sin motivo herencia y con acceso a rentas vitalicias actuarialmente justas, debería convertir la totalidad de su riqueza en rentas vitalicias al jubilarse, porque la renta vitalicia domina a cualquier otra forma de ahorro: paga más mientras se vive y no deja nada que no se habría disfrutado. Sin embargo, en la práctica, casi nadie lo hace de forma voluntaria, lo que se conoce como el enigma de las rentas vitalicias (\textit{annuity puzzle}). La literatura ha ofrecido varias explicaciones complementarias, el más evidente sería el motivo herencia. Pero también la necesidad de liquidez ante gastos sanitarios imprevistos, la desconfianza hacia las aseguradoras, los sesgos de percepción que llevan a infravalorar la propia longevidad y, de forma destacada, el elevado precio que la selección adversa impone a estos productos.

Cuando el Reino Unido suprimió en 2015 la obligación de convertir en renta vitalicia la mayor parte del fondo de pensiones acumulado (las llamadas pension freedoms, anunciadas en el presupuesto de 2014), las ventas de rentas vitalicias se desplomaron, lo que confirma que muy pocos las compran si no se les obliga. La conclusión es que la pensión pública vitalicia e indexada es, en términos económicos, la mayor renta vitalicia del país, y es precisamente el producto que el mercado privado no consigue ofrecer a un precio razonable.

Ahora bien, el alcance de este argumento tiene un límite preciso que conviene no exagerar. La selección adversa justifica la obligatoriedad, pero no exige por sí sola la provisión pública ni el reparto. En teoría bastaría con obligar a todos los jubilados a comprar una renta vitalicia a aseguradoras privadas, prohibiendo la selección de riesgos y exigiendo primas uniformes, que es aproximadamente lo que hacen los sistemas de capitalización obligatoria en su fase de desacumulación. FRIEDMAN (1962) aceptaba explícitamente este razonamiento. Consideraba defendible obligar a asegurarse, pero rechazaba que el Estado gestionara el seguro. 

\paragraph{Los riesgos que el mercado no puede diversificar: inflación, longevidad agregada y riesgo entre generaciones}
El seguro privado funciona gracias a la ley de los grandes números. Si los riesgos individuales son independientes entre sí, la experiencia de un grupo suficientemente grande es muy predecible aunque la de cada individuo no lo sea: nadie sabe si él morirá este año, pero una aseguradora sabe con bastante precisión cuántos de cada cien mil asegurados de una edad determinada lo harán. El problema surge cuando el riesgo no es independiente, sino común a todos los asegurados a la vez. Un riesgo así no se diversifica agrupando individuos, porque todos sufren el mismo shock al mismo tiempo. BARR (2001) y BARR y DIAMOND (2008) identifican varios riesgos de este tipo en el ámbito de las pensiones, y son estos, más que la selección adversa, los que sostienen el argumento a favor de un papel público en la garantía de las prestaciones.

El primero es la inflación imprevista. Una renta vitalicia nominal contratada por un jubilado español a comienzos de los años setenta habría perdido la mayor parte de su poder adquisitivo en las dos décadas siguientes, con tasas de inflación que superaron el veinte por ciento a finales de esa década. Una aseguradora privada solo puede garantizar rentas indexadas a los precios si dispone de activos indexados de muy largo plazo con los que cubrirse, y esos activos han sido históricamente escasos o inexistentes en la mayoría de los mercados. El Estado, en cambio, puede garantizar la indexación porque su base de ingresos, los impuestos y las cotizaciones sobre salarios nominales, crece por sí misma con la inflación (progresividad en frio). Si los precios suben, suben los salarios nominales y con ellos las cotizaciones, lo que permite pagar pensiones nominalmente más altas sin necesidad de ningún activo de cobertura. Este es el fundamento económico, no solo político, de la garantía de revalorización con el IPC. Es también el fundamento del art. 50 CE cuando exige pensiones periódicamente actualizadas, el constituyente no pedía un privilegio, sino la cobertura de un riesgo que solo el sector público puede cubrir.

El segundo es la longevidad agregada. Una aseguradora puede diversificar el riesgo de que un asegurado concreto viva más de lo esperado, pero no el de que toda una cohorte viva más de lo que preveían las tablas de mortalidad con las que se fijaron las primas. Si la esperanza de vida a los sesenta y cinco años aumenta un año más de lo previsto, todas las rentas vitalicias de esa cohorte cuestan más a la vez, y ninguna agrupación de asegurados lo compensa. La historia demográfica del último medio siglo ha estado llena de estas sorpresas: las proyecciones de mortalidad han tendido a infravalorar sistemáticamente las ganancias de esperanza de vida en edades avanzadas. Un sistema público puede absorber este riesgo repartiéndolo entre generaciones, sea con más cotizaciones de los activos, con ajustes de la edad de jubilación o con mecanismos automáticos como los que veremos más adelante. Un asegurador privado solo puede trasladarlo al precio, encareciendo el producto para todos.

El tercero, y el más profundo desde el punto de vista teórico, es el riesgo entre generaciones. MERTON (1983) y, sobre todo, GORDON y VARIAN (1988) mostraron que los mercados privados no pueden repartir riesgos entre generaciones que todavía no han nacido, por la sencilla razón de que nadie puede firmar un contrato con quien aún no existe. Una generación que llega a la jubilación justo cuando estalla una crisis financiera, una guerra o una depresión sufre un shock que, en un sistema de capitalización individual, recae íntegramente sobre ella. Si esa misma generación hubiera podido firmar antes de nacer un contrato de seguro con las generaciones anteriores y posteriores, repartiendo entre todas el riesgo de nacer en mal momento, todas habrían estado mejor en términos esperados. SHILLER (1999) formuló esta idea como una de las justificaciones más potentes de los sistemas públicos de reparto: al obligar por ley a las generaciones futuras a participar en el sistema, estos sistemas completan un mercado que no puede existir, el de los seguros entre generaciones. España vivió implícitamente un ejemplo de este reparto entre 2008 y 2013. Los pensionistas no vieron caer sus pensiones al ritmo al que cayeron el empleo, la recaudación o el valor de los fondos de pensiones privados porque el shock se repartió entre generaciones mediante el déficit del sistema, el uso del Fondo de Reserva y, más tarde, la deuda con el Estado.

Pero este argumento tiene una contracara que conviene señalar desde ahora, porque recorre todo el tema.\footnote{%
    \textbf{OJO.-} La misma capacidad de un sistema público para repartir riesgos entre generaciones es la capacidad para trasladar cargas a generaciones que no votan porque todavía no existen o no tienen edad para hacerlo. La frontera entre un seguro intergeneracional eficiente y una transferencia intergeneracional políticamente motivada es borrosa, y es el núcleo de la crítica de la contabilidad generacional y de la economía política de la gerontocracia.
}

\paragraph{Miopía, sesgo presente y el argumento paternalista}
El tercer fallo no es de mercado, sino de conducta: afecta al comportamiento del propio individuo. DIAMOND (1977) situó la miopía en el centro de la justificación de la seguridad social. 

Una parte significativa de la población no ahorra lo suficiente para su jubilación, no porque haya decidido racionalmente consumir más hoy, sino por falta de información, por dificultad para calcular sus necesidades futuras o por una incapacidad genuina para planificar a tan largo plazo. Si eso es así, un mandato de ahorro no reduce el bienestar de estos individuos, sino que lo aumenta, porque corrige una decisión que ellos mismos, en un momento de reflexión, reconocerían como un error.

La economía conductual dio a este argumento un fundamento formal con el descuento cuasi-hiperbólico de LAIBSON (1997), que recupera una idea anterior de STROTZ (1955) y PHELPS y POLLAK (1968). En el modelo estándar, un individuo descuenta el futuro a una tasa constante, de modo que su valoración relativa de dos momentos futuros no cambia con el paso del tiempo. En el modelo de LAIBSON, en cambio, el individuo aplica un descuento adicional a todo lo que no es presente inmediato, lo que genera \textbf{inconsistencia temporal}: los planes que hace hoy para el futuro no son los que querrá cumplir cuando ese futuro llegue.\footnote{%
    Supongamos un individuo con un factor de sesgo presente de $0,7$ y sin ningún otro descuento. Visto desde hoy, un euro dentro de un año y un euro dentro de dos pesan lo mismo, $0,7$ cada uno, así que hoy está dispuesto a comprometerse a sacrificar cien euros del año próximo para obtener ciento cinco el siguiente, porque ciento cinco por $0,7$ supera a cien por $0,7$. Pero cuando llega el año próximo, esos cien euros se han convertido en presente y pesan uno, mientras que los ciento cinco del año siguiente pesan $0,7$, es decir, $73,5$. El plan de ahorro se abandona. Y como la misma lógica se repite cada año, el individuo pospone indefinidamente el ahorro que siempre tiene la intención de empezar el año que viene.
}

Un sistema obligatorio actúa en este contexto como un dispositivo de compromiso. La versión institucional de Ulises pidiendo a su tripulación que lo ate al mástil para no ceder al canto de las sirenas: el individuo no es obligado contra su voluntad reflexiva, sino protegido contra su voluntad impulsiva.

La evidencia empírica sobre la importancia de estos sesgos es abundante. El trabajo más citado es de MADRIAN y SHEA (2001), que estudiaron una empresa estadounidense que cambió la forma de adherir a sus empleados al plan de pensiones de empresa. Antes, el empleado debía solicitar expresamente su adhesión; después, quedaba adherido automáticamente salvo que solicitara expresamente la baja. La participación de los nuevos empleados pasó de algo más de un tercio a más del ochenta y cinco por ciento, pese a que las condiciones económicas del plan no habían cambiado en absoluto. Para la teoría estándar este resultado es inexplicable, porque la opción por defecto no altera ni los costes ni los beneficios de participar; para la economía conductual es la prueba de que la inercia y la procrastinación dominan las decisiones de ahorro de largo plazo. THALER y BENARTZI (2004) mostraron, con el programa \textit{Save More Tomorrow}, que comprometer a los trabajadores a destinar al ahorro una parte de sus futuras subidas salariales, y no de su salario actual, eleva drásticamente las tasas de ahorro, precisamente porque esquiva el sesgo presente: el sacrificio se sitúa en un futuro que el individuo valora con menos intensidad.

Como podemos suponer, estos hallazgos han abierto un debate de gran calado sobre la forma de intervenir. Si el problema es sobre todo la inercia y no una irracionalidad profunda, quizá no hace falta obligar: bastaría con diseñar bien las opciones por defecto y dejar al individuo la libertad de salirse. Es la propuesta del paternalismo libertario de THALER y SUNSTEIN (2008), y es la vía que siguieron Nueva Zelanda con KiwiSaver en 2007 y el Reino Unido con la adhesión automática a planes de pensiones de empresa, introducida por la Pensions Act de 2008 y aplicada gradualmente desde 2012. Sus defensores subrayan que preserva la libertad individual y respeta las preferencias heterogéneas: quien tiene buenas razones para no ahorrar, por ejemplo porque tiene deudas caras que amortizar, puede salirse. Sus críticos oponen dos objeciones de peso. La primera es que los individuos con mayor sesgo presente, los que más necesitan la protección, son también los que con más probabilidad usan la opción de salida en un momento de apuro, de modo que el empujón deja fuera precisamente a quienes pretendía proteger. La segunda es que la adhesión automática corrige la miopía, pero no la selección adversa ni los riesgos no diversificables examinados en los apartados anteriores. Un sistema de ahorro voluntario con opción por defecto sigue dejando a los jubilados expuestos a la inflación, a la longevidad agregada y a la mala suerte de jubilarse en plena crisis financiera.\footnote{%
    \textbf{NOTA.-} Este debate es directamente relevante para España, porque los planes de pensiones de empleo simplificados creados por la Ley 12/2022 intentan desarrollar el segundo pilar sin mandato y sin adhesión automática generalizada, es decir, confiando en la negociación colectiva. Su diseño y su escasa implantación se analizan en el bloque III (3.1.6).
}

El argumento paternalista tiene también críticos de principio, y su formulación más rigurosa es económica, no filosófica. FELDSTEIN (1985) planteó el dilema en términos formales. Cuanto más generosa es la pensión obligatoria, más protege a los miopes de su propia imprevisión; pero también más distorsiona las decisiones de los previsores, que habrían ahorrado por su cuenta de todas formas y que ahora ven sustituido su ahorro voluntario por uno obligatorio con otra rentabilidad, otro riesgo y otra liquidez. El nivel óptimo de la pensión pública depende entonces de la proporción de miopes en la población y de la magnitud de su miopía. Con una proporción moderada, FELDSTEIN concluía que el nivel óptimo era bastante inferior al de la mayoría de los sistemas europeos, y mucho más próximo a una pensión mínima que a una pensión proporcional al salario. La crítica de principio más influyente procede de HAYEK (1960), que aceptaba que el Estado obligara a los individuos a asegurarse, precisamente para evitar que el imprevisor se convirtiera en una carga para los demás, pero rechazaba el monopolio estatal del seguro por dos razones: (i) suprime la competencia entre gestores, que es la fuente de la innovación y de la eficiencia; y, (ii) convierte lo que debería ser un seguro en un instrumento de redistribución sometido a la lógica de la negociación política. Como podemos ver, este segundo argumento anticipaba la politización que se observa hoy en día de las prestaciones, y analizaremos en el último apartado.

\paragraph{El dilema del samaritano: cuando la previsión individual se convierte en una externalidad}
El cuarto argumento es quizá el más elegante, porque deriva la obligatoriedad no de un defecto del individuo, sino de una virtud de la sociedad. BUCHANAN (1975) lo formuló como el dilema del samaritano. 

Una sociedad altruista no dejará morir en la miseria a un anciano sin recursos, aunque su pobreza se deba a que no ahorró cuando pudo hacerlo. El individuo racional que conoce este altruismo tiene un incentivo a no ahorrar, porque sabe que, si no lo hace, la sociedad lo rescatará igualmente. El altruismo colectivo genera así un riesgo moral que reduce la previsión individual, y el samaritano queda explotado por quienes cuentan con su compasión. COATE (1995) formalizó el argumento y mostró que un mandato de ahorro o de aseguramiento es la respuesta óptima del samaritano para no ser explotado: obligando a todos a prever, la sociedad se asegura de que su altruismo no se convierta en una subvención a la imprevisión. En términos económicos, consiste en internalizar esa externalidad mediante un mandato legal.

Este argumento tiene una traducción muy directa en el diseño del sistema español, que conviene anticipar con un ejemplo. Supongamos que la pensión mínima garantizada para un jubilado es de novecientos euros mensuales y que un trabajador de salario bajo, que ha cotizado durante toda su vida, genera una pensión contributiva pura de ochocientos cincuenta euros. El sistema la completa con un complemento a mínimos de cincuenta euros hasta alcanzar los novecientos. Pero si ese mismo trabajador hubiera cotizado bastante menos, por ejemplo porque parte de su actividad transcurrió en la economía sumergida, y hubiera generado una pensión pura de setecientos euros, cobraría igualmente novecientos si cumple los requisitos de acceso al complemento. Para él, hasta el umbral del mínimo, el rendimiento marginal de cotizar más es prácticamente nulo, lo que equivale a un tipo impositivo implícito cercano al cien por cien sobre la cotización adicional. Esta versión contemporánea del dilema de BUCHANAN, explica por qué los sistemas con mínimos generosos necesitan combinar la obligatoriedad con un control estricto del cumplimiento, por qué el diseño de los mínimos afecta a la contributividad del sistema y por qué la frontera entre las prestaciones contributivas y las asistenciales no es solo contable, sino también de incentivos.

\begin{specialblock}{El riesgo moral: el límite interno de todos los argumentos anteriores}
    Los argumentos anteriores justifican la intervención, pero ninguno justifica una protección ilimitada. El seguro social está sujeto al mismo problema de información asimétrica que el privado, pero en sentido inverso: una vez asegurado, el individuo puede alterar su conducta de forma que aumente la probabilidad o la cuantía del siniestro. Es el riesgo moral, que ARROW (1963) identificó como un límite natural de todo seguro y que PAULY (1968) mostró que puede hacer ineficiente incluso un seguro completo desde el punto de vista del asegurado.

    En el sistema de pensiones, el riesgo moral se manifiesta de tres formas principales. La primera es la jubilación anticipada. Cuando las penalizaciones por adelantar la jubilación no son actuarialmente neutrales, es decir, cuando no reducen la pensión lo suficiente para compensar los años adicionales de cobro, el sistema subvenciona implícitamente la salida temprana del mercado de trabajo. Buena parte de la reducción de la edad efectiva de jubilación en la Europa de los años setenta y ochenta se explica así. La segunda es la incapacidad. La frontera entre \textit{no poder trabajar} y \textit{no querer trabajar} es, en muchas situaciones, difusa, y la generosidad de las pensiones de incapacidad afecta a su incidencia. La tercera, en el corto plazo, es la incapacidad temporal, cuya evolución reciente en España ilustra hasta qué punto la gestión y el control de una prestación importan para su coste.
    
    El marco moderno que integra ambas caras de la cuestión es el de CHETTY y FINKELSTEIN (2013). El nivel óptimo de un seguro social se alcanza cuando, en el margen, el valor del seguro para el asegurado, que depende de su aversión al riesgo y de la magnitud de la caída de consumo que sufriría sin protección, se iguala con el coste del riesgo moral, que depende de cuánto altera su conducta en respuesta a la protección. Esta fórmula convierte un debate ideológico sobre cuánta protección en una pregunta empírica sobre dos magnitudes que pueden estimarse y que analizaremos posteriormente bajo la regla de BAILY-CHETTY aplicada al seguro de desempleo.
\end{specialblock}


#### Los bienes preferentes: el argumento de valor

El último argumento no invoca ningún fallo de mercado ni de conducta, sino un juicio colectivo de valor. MUSGRAVE (1959) introdujo la categoría de los bienes preferentes (merit wants): bienes cuyo consumo la sociedad considera valioso con independencia de las preferencias individuales, de modo que su provisión se justifica aunque los individuos, dejados a su libre elección, consumirían menos de ellos. La seguridad económica en la vejez, la protección de la salud y la educación son los ejemplos canónicos. La justificación no descansa en que el individuo se equivoque, como en el argumento de la miopía, sino en que la sociedad considera que ciertos niveles de bienestar forman parte de lo que una vida digna exige, y los garantiza como tales. La crítica a este argumento es obvia y conocida: el concepto de bien preferente es tan elástico que puede justificar casi cualquier intervención, y abre la puerta a que la autoridad pública sustituya las preferencias de los ciudadanos por las suyas. Por eso buena parte de la literatura contemporánea ha preferido reformularlo en términos conductuales, de modo que lo que MUSGRAVE llamaba preferencia social se convierte en la corrección de sesgos individuales documentados empíricamente. 

El argumento de MUSGRAVE, aunque pueda presentar determinados sesgos a nivel teórico, es muy valioso para interpretar de forma finalista nuestro texto constitucional. Cuando el art. $50$ CE habla de suficiencia económica y de pensiones adecuadas, y cuando el art. $41$ habla de prestaciones sociales suficientes, el constituyente no está justificando el sistema por un fallo de mercado: está proclamando un valor, y ese valor es exactamente lo que MUSGRAVE llamaba un bien preferente.

#### Recapitulación: qué justifica cada argumento

Reunidos los argumentos, su alcance respectivo queda claro. La selección adversa, la miopía y el dilema del samaritano justifican con solidez la obligatoriedad de la protección, pero no exigen por sí mismos que la gestione el sector público. Los riesgos no diversificables, en cambio, justifican además un papel público en la garantía, porque solo un agente capaz de indexar sus ingresos a la inflación y de obligar a las generaciones futuras puede cubrir la inflación, la longevidad agregada y el riesgo entre generaciones. El riesgo moral fija el límite de la generosidad de la protección, y los bienes preferentes explican el lenguaje valorativo con que las constituciones consagran ese papel. Ninguno de estos argumentos, sin embargo, resuelve la cuestión de si el sistema debe financiarse por reparto o por capitalización. Esa es una pregunta sobre la rentabilidad relativa de ambos métodos, sobre sus efectos en el ahorro y el crecimiento y sobre el coste de pasar de uno a otro que se examinarán más adelante.

\subsubsection{Argumento de equidad: Seguro frente a redistribución}
El argumento de eficiencia tiene una limitación evidente: en un mundo donde todos los mercados funcionaran perfectamente, no justificaría ninguna intervención. Pero la mayoría de los ciudadanos y de los sistemas jurídicos no conciben la protección social solo como la corrección de un fallo de mercado. La conciben también, y a menudo sobre todo, como una forma de justicia. 

La sociedad garantiza a sus miembros un nivel de vida digno ante la vejez, la enfermedad o la pérdida del empleo, con independencia de cuánto hayan podido aportar. Examinamos ahora este segundo fundamento: cómo se relaciona con el primero, cuáles son sus versiones filosóficas, cuánto cuesta en términos de eficiencia y qué críticas ha recibido. La cuestión central que lo recorre es cuánto del sistema es seguro (cada uno recibe en proporción a lo que aporta) y cuánto es redistribución (unos reciben más de lo que aportan, a costa de otros). Esa proporción es la que define, en último término, la naturaleza de un sistema de protección social.

\paragraph{Hucha y Robin Hood: cuánto del Estado del bienestar es seguro}

BARR (2001) propuso distinguir dos funciones del Estado del bienestar que el debate político suele confundir. La primera es la de \textbf{hucha} (piggy bank). El Estado permite a los individuos trasladar recursos a lo largo de su propia vida, desde los años de actividad hacia los de vejez, enfermedad o desempleo, y asegurarse contra riesgos que el mercado no cubre bien, por las razones examinadas en el apartado anterior. La segunda es la de \textbf{Robin Hood}. El Estado traslada recursos de unos individuos a otros, de los más ricos a los más pobres, de los sanos a los enfermos, de los que tienen empleo estable a los que no lo tienen.

Su tesis central es que la primera función domina sobre la segunda en mucha mayor medida de lo que suele creerse, y que esta constatación cambia los términos del debate. Buena parte de lo que parece redistribución entre personas es, en realidad, redistribución de cada persona consigo misma a lo largo de su vida. El joven que paga impuestos y cotizaciones y no recibe prestaciones es, décadas después, el jubilado que recibe una pensión sin pagar apenas cotizaciones.\footnote{%
    El estudio empírico de referencia es el de FALKINGHAM y HILLS (1995), que, mediante microsimulación del ciclo vital completo de una población británica, estimaron que en torno a tres cuartas partes de las prestaciones que recibe un individuo a lo largo de su vida las financia él mismo con sus propios impuestos y cotizaciones, y solo la cuarta parte restante supone una redistribución neta entre personas.
}

La consecuencia es doble. Por un lado, el Estado del bienestar es menos redistributivo de lo que sugiere una fotografía anual, que compara a jóvenes que pagan con viejos que cobran. Por otro, su justificación descansa en buena medida en el argumento de eficiencia, y no solo en el de equidad, lo que lo hace más difícil de desmontar con argumentos puramente liberales.

Esta distinción ilumina directamente la arquitectura española. La pensión contributiva de jubilación es fundamentalmente hucha. Su cuantía depende de las bases por las que se ha cotizado y de los años de cotización, y su lógica es la del alisamiento del consumo a lo largo del ciclo vital. Los complementos a mínimos, las pensiones no contributivas y el IMV son fundamentalmente Robin Hood. Su cuantía no depende de lo aportado, sino de la situación de necesidad. El principio de separación de fuentes que el Pacto de Toledo consagró en 1995 (las prestaciones contributivas se financian con cotizaciones y las no contributivas con impuestos) es la traducción institucional de la distinción de BARR: cada función se financia con el instrumento que corresponde a su naturaleza. 

Sin embargo, la frontera real entre ambas funciones es mucho más borrosa de lo que sugiere el principio. La propia pensión contributiva contiene elementos redistributivos importantes: la base máxima que limita las cotizaciones, la pensión máxima que limita las prestaciones, la escala del porcentaje aplicable según los años cotizados y la integración de lagunas de cotización. La discusión sobre qué gastos son impropios de la Seguridad Social y deberían financiarse con impuestos es, en el fondo, una discusión sobre dónde está la frontera entre la hucha y Robin Hood.

\paragraph{Las dos lógicas fundacionales: Bismarck y Beveridge}

Una segunda vertiente teórica distingue entre un modelo contributivo o bismarckiano, en el que las prestaciones guardan relación con las contribuciones realizadas, y un modelo no contributivo o beveridgiano, basado en prestaciones de tipo asistencial que no dependen de lo aportado. La diferencia más clara es que el primero se financia con cotizaciones sociales y el segundo con imposición general, distinción correcta y clásica, pero que queda en la superficie.\footnote{%
    \textbf{NOTA.-} En realidad, el plan que BEVERIDGE propuso en 1942 no se financiaba con impuestos generales, sino con cotizaciones uniformes de igual cuantía para todos, a cambio de prestaciones también uniformes. La identificación entre modelo beveridgiano y financiación por impuestos, es una simplificación posterior, derivada sobre todo de la evolución del servicio nacional de salud británico y de los sistemas nórdicos. Lo verdaderamente beveridgiano no es la fuente de financiación, sino la uniformidad de la prestación y la universalidad del sujeto protegido.
} 

Detrás de los dos nombres no hay solo dos técnicas de financiación, sino dos teorías distintas sobre por qué existe la protección social y a quién debe proteger, que conviene exponer en su lógica interna:
\begin{la}
    \item \textbf{BISMARCK}. La lógica bismarckiana es la del seguro social profesional. El sujeto protegido no es el ciudadano, sino el trabajador asalariado y, a través de él, su familia. La prestación es la contrapartida de la cotización, y su cuantía es proporcional al salario por el que se ha cotizado, porque el objetivo no es garantizar un mínimo, sino mantener el nivel de vida alcanzado durante la vida activa. La financiación corresponde a las cotizaciones de empresarios y trabajadores, y la gestión, históricamente, a cajas o entidades con participación de los agentes sociales. El resultado es la reproducción del estatus: quien ganaba mucho recibe una pensión alta y quien ganaba poco, una pensión baja, de modo que la jerarquía social de la vida activa se prolonga en la vejez. El fundamento de legitimidad es el esfuerzo contributivo: se tiene derecho a la pensión porque se ha pagado por ella. Lógica que explica la fortaleza política de los sistemas bismarckianos, pues los cotizantes perciben la pensión como una propiedad y no como una dádiva. Pero también su principal debilidad, pues dejan desprotegidos a quienes no han tenido un empleo formal estable;
    \item \textbf{BEVERIDGE}. La lógica beveridgiana es la de la ciudadanía. El sujeto protegido es el ciudadano o el residente, con independencia de su situación laboral. La prestación es un derecho derivado de la pertenencia a la comunidad política, y su cuantía es uniforme, porque el objetivo no es mantener el nivel de vida de cada uno, sino garantizar a todos un mínimo nacional por debajo del cual nadie debe caer. El fundamento de legitimidad no es el esfuerzo contributivo, sino la necesidad y la ciudadanía compartida, y el resultado buscado no es la reproducción del estatus, sino la eliminación de la pobreza. Por encima de ese mínimo, cada uno queda libre de asegurarse privadamente el nivel de vida que desee.
\end{la}

Evidentemente, ningún sistema real responde hoy a ninguno de los dos modelos en estado puro. BONOLI (1997) propuso por eso clasificar los Estados del bienestar en dos dimensiones a la vez: el grado en que siguen la lógica bismarckiana o la beveridgiana (el cómo) y el tamaño de su gasto (el cuánto). Así, países igualmente bismarckianos pueden ser muy distintos según cuánto gastan, e igualmente beveridgianos según la generosidad de su mínimo.\footnote{%
    España es un híbrido característico. Sus pensiones y su protección por desempleo son de raíz bismarckiana; su sanidad es beveridgiana desde la Ley General de Sanidad de 1986, que convirtió la asistencia sanitaria en un derecho de ciudadanía; y ha ido construyendo una capa asistencial de tipo beveridgiano: las pensiones no contributivas de 1990, las rentas mínimas autonómicas desde finales de los ochenta y el IMV en 2020.
}

\paragraph{Ciudadanía social: Marshall, Titmuss y la crítica de la selectividad}
La lógica beveridgiana encontró su formulación filosófica más influyente en la sociología británica de posguerra. T. H. MARSHALL (1950), describió la ciudadanía moderna como el resultado de una acumulación histórica de derechos en tres fases sucesivas. Los derechos civiles, conquistados en el siglo XVIII, garantizan la libertad personal, la propiedad, el contrato y el acceso a la justicia. Los derechos políticos, del siglo XIX, garantizan la participación en el poder mediante el sufragio. Los derechos sociales, que el siglo XX estaba construyendo, abarcan, en su conocida fórmula, desde el derecho a un mínimo de bienestar económico y de seguridad hasta el derecho a participar plenamente en la herencia social y a vivir la vida de un ser civilizado conforme a los estándares predominantes en la sociedad. 

La tesis tiene una consecuencia normativa fuerte, que es la que interesa aquí. Si los derechos sociales forman parte de la ciudadanía en el mismo sentido que los civiles y los políticos, no pueden depender de la contribución previa, del mismo modo que el derecho de voto no depende de haber pagado impuestos. Un derecho que depende de lo aportado es un derecho de propietario o de asegurado, no de ciudadano. Es la raíz intelectual del principio de universalidad que consagra el art. $41$ de nuestra Constitución cuando habla de un régimen público de Seguridad Social para todos los ciudadanos. MARSHALL añadía una observación inquietante que conviene recordar: los derechos sociales conviven en tensión con el sistema de clases que produce el mercado, porque tienden a igualar lo que el mercado desiguala, y esa tensión no se resuelve nunca del todo.

R. TITMUSS, colega de MARSHALL en la London School of Economics, aportó dos ideas que siguen vivas en el debate contemporáneo. La primera es una tipología de tres modelos de política social, formulada en sus lecciones publicadas en 1974. El modelo residual sostiene que el Estado solo debe intervenir cuando fallan los dos canales naturales de satisfacción de necesidades, el mercado y la familia, y solo de forma temporal. El modelo de logro personal o de rendimiento industrial concibe las prestaciones como complemento de la economía y las hace depender del mérito, el trabajo y la productividad: es, en esencia, la lógica bismarckiana. El modelo institucional-redistributivo concibe la protección social como una institución integrada en la sociedad, que provee servicios universales fuera del mercado según el principio de necesidad. Esta tipología es el antecedente directo de la de ESPING-ANDERSEN.

La segunda idea de TITMUSS es su crítica de la selectividad. Frente a quienes defendían concentrar los recursos en los más necesitados mediante prestaciones condicionadas a la comprobación de recursos, TITMUSS sostuvo que esos programas producen tres efectos perversos. Generan estigma, porque solicitarlos equivale a declararse pobre. Tienen baja utilización, porque muchos de quienes tienen derecho no los solicitan, por vergüenza, por desconocimiento o por la complejidad del trámite. Y pierden apoyo político, porque la mayoría que los financia no se beneficia de ellos. Su conclusión, que a menudo se resume en la fórmula de que los servicios para los pobres acaban siendo servicios pobres, anticipa una de las tesis más influyentes de la sociología comparada, la \textbf{paradoja de la redistribución}.\footnote{%
    \textbf{NOTA.-} El problema de la baja utilización tiene un ejemplo muy próximo en el IMV: en diciembre de 2021, la prestación cubría a unos 284.000 hogares, en torno al cuarenta por ciento de sus posibles beneficiarios según la AIReF.
}

\paragraph{Velo de la ignorancia: la redistribución como seguro}

La filosofía política de la segunda mitad del siglo XX ofreció una forma de reconciliar los dos fundamentos, el seguro y la redistribución, que tiene consecuencias económicas de primer orden. 

HARSANYI (1953, 1955) y RAWLS (1971) propusieron evaluar la justicia de las instituciones desde una posición imaginaria en la que los individuos deciden las reglas de la sociedad sin saber qué lugar ocuparán en ella: si nacerán sanos o enfermos, con talento o sin él, en una familia rica o pobre, en un momento de prosperidad o de crisis. Es la célebre posición original tras un velo de ignorancia. HARSANYI, suponiendo que los individuos maximizan su utilidad esperada, dedujo que elegirían las reglas que maximizan la utilidad media: el utilitarismo. RAWLS, suponiendo una aversión extrema al riesgo ante una decisión tan trascendental, dedujo que elegirían las reglas que maximizan la situación de los peor situados, su célebre principio de diferencia, según el cual las desigualdades solo se justifican en la medida en que beneficien a los menos aventajados.

Lo decisivo no es qué versión es correcta, sino lo que ambas comparten. Si los individuos eligen las reglas sin saber qué suerte les tocará, entonces toda redistribución puede interpretarse como un seguro elegido ex ante contra el riesgo de nacer con malas cartas. La redistribución ex post, de los afortunados a los desafortunados, es la liquidación de un seguro que todos habrían contratado ex ante si hubieran podido. DWORKIN (1981) desarrolló una versión de esta idea con su mercado hipotético de seguros: la justicia distributiva exige compensar las desventajas que los individuos se habrían asegurado contra ellas si hubieran podido hacerlo en condiciones de igualdad, pero no las que resultan de sus propias elecciones. Esta distinción entre \textit{circunstancias} y \textit{elecciones} es la base filosófica de la diferencia entre la protección frente a la invalidez, que es una circunstancia, y el tratamiento de la jubilación anticipada voluntaria, que es una elección.

SINN (1995), extrajo de esta idea una consecuencia económica que cambia la forma de pensar la relación entre equidad y eficiencia. Si la redistribución es un seguro contra los riesgos de la vida, el Estado del bienestar no es solo un coste para la eficiencia, sino que puede aumentarla. Un individuo asegurado contra el fracaso está dispuesto a asumir riesgos productivos que no asumiría sin red: emprender un negocio, invertir en una formación incierta, cambiar de empleo o de ciudad en busca de una oportunidad mejor. SINN advertía, al mismo tiempo, que el mismo mecanismo que fomenta la asunción de riesgos productivos fomenta también el riesgo moral, de modo que existe un nivel óptimo de protección más allá del cual el efecto negativo domina, el mismo equilibrio que CHETTY y FINKELSTEIN formalizaron para cada programa concreto, trasladado ahora al conjunto del Estado del bienestar.

En una línea distinta pero complementaria, el enfoque de las capacidades de SEN (1985, 1999) ofrece un fundamento para la noción jurídica de suficiencia. Para SEN, lo que importa no es la renta en sí, sino las capacidades que permite: la libertad efectiva de las personas para realizar los funcionamientos que valoran, como estar bien alimentado, tener buena salud o participar en la vida de la comunidad. Desde esta perspectiva, una pensión suficiente no es la que alcanza una cifra determinada, sino la que permite al pensionista mantener esas capacidades. Es un criterio que se ajusta mejor que cualquier umbral monetario al mandato constitucional de garantizar la suficiencia económica en la tercera edad, y que explica la tendencia reciente a vincular las pensiones mínimas al umbral de la pobreza relativa.

\paragraph{Coste de redistribuir: del cubo agujereado a la paradoja de la redistribución}

Si la redistribución fuera gratuita, el debate sobre su alcance sería solo normativo. Pero OKUN (1975) formuló con una metáfora célebre por qué no lo es: redistribuir renta de ricos a pobres es como transportar agua en un cubo agujereado. Por el camino se pierde una parte, por los desincentivos que los impuestos generan en quienes los pagan, por los que las prestaciones generan en quienes las reciben y por los costes administrativos del propio traslado. La pregunta normativa que planteaba OKUN es cuánta agua estamos dispuestos a perder para que llegue la que llega. Una sociedad muy igualitaria aceptará perder mucha; una muy preocupada por la eficiencia, poca. La respuesta depende del tamaño de los agujeros, que es una cuestión empírica.

LINDERT (2004), planteó un contraargumento histórico que ha alimentado un debate intenso el \textbf{rompecabezas de la comida gratis} (free lunch puzzle). Si los agujeros del cubo fueran tan grandes como sugiere la teoría, los países con Estados del bienestar muy extensos, como los nórdicos, deberían haber crecido sistemáticamente menos que los países con Estados del bienestar reducidos. Pero la evidencia histórica no muestra esa relación. La explicación de LINDERT es que los Estados del bienestar grandes han aprendido a minimizar los costes de eficiencia: se financian con impuestos de base amplia y relativamente poco distorsionantes, en lugar de con impuestos muy progresivos sobre el capital, y diseñan sus prestaciones para que penalicen poco el trabajo, con políticas de activación, prestaciones en especie como la educación y la sanidad, y edades de jubilación relativamente altas.

Los críticos de LINDERT señalan que la comparación entre países no permite identificar causalidad, pues puede que los países ricos y bien gobernados se permitan Estados del bienestar grandes y no al revés, y que la muestra de países con Estados del bienestar muy grandes es pequeña y homogénea. La aportación más provocadora a este debate procede de la sociología comparada. KORPI y PALME (1998), en un art. que se ha convertido en un clásico, formularon la \textbf{paradoja de la redistribución}: cuanto más se focalizan las prestaciones en los pobres, menos se reduce la pobreza y la desigualdad. La paradoja se resuelve con un argumento de economía política que desarrolla las intuiciones de TITMUSS. Los programas universales o proporcionales al salario benefician también a las clases medias, que por eso los defienden políticamente y aceptan financiarlos con impuestos elevados. Los programas focalizados solo benefician a los pobres, que son pocos, tienen escasa influencia política y no pueden sostener por sí solos el apoyo presupuestario necesario, de modo que esos programas tienden a ser pequeños y a erosionarse. El resultado es que, aunque cada euro de un programa focalizado redistribuye más que cada euro de uno universal, los programas universales movilizan muchos más euros, y la redistribución total es mayor. En otros términos, el tamaño del pastel importa más que el reparto del pastel.

KENWORTHY (2011) y MARX, SALANAUSKAITE y VERBIST (2013) mostraron, con datos más recientes, que la relación negativa entre focalización y redistribución se ha debilitado desde los años noventa, y que algunos países, como Dinamarca o los Países Bajos, han logrado combinar universalidad en la mayor parte del sistema con una focalización eficaz en su parte baja. La lección no sería, por tanto, que la focalización es siempre contraproducente, sino que funciona cuando se inserta en un sistema con amplio apoyo social, y fracasa cuando sustituye a ese sistema.\footnote{%
    \textbf{NOTA.-} Esta discusión ilustra con precisión el problema español. El sistema contributivo de pensiones tiene una coalición defensora formidable: más de nueve millones de pensionistas contributivos, sus familias y los trabajadores próximos a la jubilación, lo que explica la reversión de la reforma de 2013. Las prestaciones focalizadas, como el IMV o las ayudas a la dependencia, tienen coaliciones mucho más débiles, lo que contribuye a explicar su menor desarrollo presupuestario, sus problemas de gestión y la persistencia de listas de espera.
}

\paragraph{Cara oculta de la redistribución: el gradiente de longevidad}
Hay un aspecto de la redistribución en los sistemas de pensiones que ha cobrado gran relevancia en la última década y que conviene tener presente, porque invierte algunas intuiciones comunes. CHETTY et al. (2016), con datos fiscales de toda la población estadounidense, estimaron que la diferencia de esperanza de vida a los cuarenta años entre el uno por ciento más rico y el uno por ciento más pobre de la distribución de la renta era de unos quince años en los hombres y de unos diez en las mujeres. En un sistema que paga pensiones vitalicias, quien vive más cobra durante más años. Por tanto, los grupos de renta alta, que viven más, recuperan más por cada euro cotizado que los de renta baja, que viven menos. De esta forma, parte de la redistribución progresiva que el sistema introduce a través de los topes de cotización y de prestación y de los complementos a mínimos, se deshace de forma inadvertida por el gradiente de longevidad.

Este hallazgo tiene dos implicaciones de gran alcance para el debate español. La primera afecta a los mecanismos que ligan la pensión o la edad de jubilación a la esperanza de vida media, como el factor de sostenibilidad de 2013. Estos mecanismos penalizan relativamente más a quienes tienen una esperanza de vida inferior a la media, que suelen ser los trabajadores de menor renta y los que desempeñaron trabajos más penosos, de modo que un instrumento aparentemente neutral tiene efectos distributivos regresivos. La segunda afecta a la valoración de las propuestas de introducir cuentas nocionales: al calcular la pensión dividiendo el capital nocional por la esperanza de vida media de la cohorte, trasladan implícitamente recursos de quienes viven menos a quienes viven más. Existen estudios para España que documentan un gradiente de longevidad por nivel educativo y de renta, de menor magnitud que el estadounidense.

\paragraph{Crítica de principio: Friedman, Hayek y Nozick}

El panorama normativo quedaría incompleto sin las críticas que niegan legitimidad o conveniencia a la redistribución a través del sistema de protección social. 

FRIEDMAN (1962) consideraba que la seguridad social pública estadounidense combinaba dos elementos separables y ambos objetables. El primero era la compra obligatoria de una renta vitalicia, que, a su juicio, restringía sin justificación suficiente la libertad de los individuos para decidir cómo prepararse para su vejez. El segundo era una redistribución encubierta, que no se sometía al control transparente de quién paga y quién recibe, porque se presentaba bajo la apariencia de un seguro. FRIEDMAN defendía, como alternativa, separar ambos elementos y sustituir la redistribución por un impuesto negativo sobre la renta, transparente y financiado con impuestos generales. NOZICK (1974), desde una posición filosófica más radical, negó la legitimidad de cualquier redistribución coactiva que fuera más allá de las funciones del Estado mínimo, por considerarla una violación de los derechos de propiedad de los individuos sobre los frutos de su trabajo. La crítica de HAYEK ya se ha expuesto en el apartado anterior.

Conviene tener presente que estas posiciones no son extravagancias marginales, sino que representan la frontera exterior del debate. Entre otras razones porque algunos de sus argumentos han sido asumidos por reformas perfectamente convencionales. La exigencia de transparencia y de separación entre seguro y redistribución que formulaba FRIEDMAN es exactamente la que inspira el sistema sueco de cuentas nocionales. En él, la parte actuarial, que relaciona estrechamente cotizaciones y prestaciones, se separa de la parte redistributiva, la pensión garantizada, que se financia con impuestos generales. La misma lógica inspira, en España, la separación de fuentes del Pacto de Toledo y el debate sobre los gastos impropios. En la dirección opuesta, las propuestas de renta básica universal, asociadas sobre todo a VAN PARIJS (1995), llevan al extremo la lógica beveridgiana de la ciudadanía y comparten con FRIEDMAN el rechazo a la condicionalidad y a la burocracia de las prestaciones categóricas. 

El debate no es solo más o menos Estado, sino también qué tipo de protección, y que los extremos del espectro ideológico pueden coincidir en sus críticas al modelo de seguro social condicionado.

\subsubsection{Regímenes de bienestar: cómo cada sociedad ha combinado seguro y redistribución}

Con todo lo expuesto, deberíamos tener mente todos los debates metafísicos que rodean al Sistema de Seguridad Social. Sin embargo, también queda claro que las ideas y razonamientos expuestos no determinan un diseño institucional único. Cada sociedad los ha combinado de forma distinta, en función de su historia, de la correlación de fuerzas entre sus clases sociales y sus partidos, de sus tradiciones religiosas y familiares y del momento en que construyó su sistema de protección. 

El resultado es una notable variedad de Estados del bienestar que no se explica por diferencias de renta o de gasto, sino por diferencias de lógica. Presentamos ahora la tipología que ha estructurado el análisis comparado durante más de tres décadas, sus críticas y ampliaciones más relevantes, la teoría sobre cómo se han reformado los sistemas de tipo continental y la cuestión de la dualización.

Con esas herramientas intenta situar a España, lo que permitirá leer su historia (1.2) y su gasto comparado (2.3.2) con criterio y no como una mera acumulación de datos.

\paragraph{Esping-Andersen y los tres mundos del capitalismo del bienestar}
ESPING-ANDERSEN (1990) transformó la forma de comparar los Estados del bienestar. Hasta entonces, la literatura comparada los había clasificado sobre todo por cuánto gastaban. ESPING-ANDERSEN sostuvo que el gasto es un indicador engañoso y que lo relevante es cómo se organiza la protección, para lo cual propuso tres dimensiones de análisis.

La primera, y la más conocida, es la \textbf{desmercantilización} (decommodification). El grado en que un individuo puede mantener un nivel de vida socialmente aceptable sin depender de la venta de su fuerza de trabajo en el mercado. El concepto es heredero directo de POLANYI, que ya vimos ([\authorfont{Ver Tema 4.B.12}]), y se mide a partir de la generosidad, la duración y las condiciones de acceso de las prestaciones de vejez, desempleo e incapacidad. Un sistema en el que las prestaciones son generosas, duraderas y accesibles con pocos requisitos desmercantiliza mucho, porque libera al trabajador de la necesidad de aceptar cualquier empleo en cualquier condición; uno con prestaciones escasas, breves y muy condicionadas desmercantiliza poco. 

La segunda dimensión es la \textbf{estratificación}. Si el sistema reproduce las diferencias de estatus existentes, con prestaciones distintas para colectivos profesionales distintos, las reduce, con prestaciones universales e iguales para todos, o crea un nuevo dualismo entre una minoría asistida por el Estado y una mayoría que se asegura en el mercado. 

La tercera es la \textbf{combinación} entre Estado, mercado y familia en la provisión del bienestar. ¿Quién asume, en cada sociedad, la responsabilidad principal de proteger a los individuos frente a los riesgos de la vida?

A partir de estas tres dimensiones, identificó tres regímenes. El régimen liberal, característico del mundo anglosajón (Estados Unidos, Reino Unido, Canadá o Australia), desmercantiliza poco. Sus prestaciones públicas son modestas, a menudo condicionadas a la comprobación de recursos, y se concentran en los más necesitados, mientras que la mayoría de la población se asegura en el mercado a través de planes de empresa o seguros privados. Su estratificación característica es un dualismo entre los asistidos, estigmatizados, y los asegurados privadamente, y el proveedor principal de bienestar es el mercado.

El régimen conservador o corporativo, característico de europa latina y continental (Alemania, Francia, Austria, Bélgica o Italia) desmercantiliza de forma intermedia. Sus prestaciones son generosas, pero dependen de la inserción laboral y son proporcionales al salario, y su lógica es la del seguro social bismarckiano. Su estratificación característica es la preservación del estatus, con regímenes distintos para colectivos profesionales distintos (funcionarios, asalariados, autónomos, agricultores), cada uno con sus propias reglas.\footnote{%
    Ha estado históricamente asociado a la influencia de la doctrina social de la Iglesia católica y a su principio de subsidiariedad, que otorga a la familia un papel central como proveedora de cuidados y hace del varón (individuo) sustentador el sujeto principal de la protección.
}

El régimen socialdemócrata, característico de los países nórdicos (Suecia, Noruega, Dinamarca y Finlandia), desmercantiliza mucho. Sus prestaciones son universales, generosas y, en buena medida, financiadas con impuestos. Su estratificación característica es la igualación, ya que las mismas prestaciones y los mismos servicios se ofrecen a todos, incluidas las clases medias, lo que, según la lógica de la paradoja de la redistribución, garantiza su apoyo político. El proveedor principal de bienestar es el Estado, que asume también muchas funciones de cuidado que en otros regímenes corresponden a la familia, facilitando así la plena incorporación de las mujeres al mercado de trabajo.

La lección metodológica de ESPING-ANDERSEN es la que interesa retener: dos países con el mismo gasto social pueden tener Estados del bienestar radicalmente distintos. Un país conservador puede gastar tanto como uno socialdemócrata y, sin embargo, redistribuir mucho menos, porque su gasto se concentra en pensiones proporcionales al salario que reproducen la desigualdad de la vida activa. Esta advertencia es la clave para leer correctamente las comparaciones que veremos después, donde España aparece con un gasto social en torno a la media europea, pero con una composición muy sesgada hacia las pensiones.

La tipología de ESPING-ANDERSEN ha sido, durante tres décadas, el punto de partida obligado de cualquier análisis comparado, pero también el blanco de críticas que la han enriquecido. La primera, y la más relevante para nosotros, se refiere a la existencia de un cuarto régimen propio de la Europa meridional. FERRERA (1996) sostuvo que Italia, España, Portugal y Grecia no encajaban en ninguno de los tres mundos y constituían un régimen propio: el régimen meriodional. Lo caracterizaban cuatro rasgos. El primero es un dualismo muy acusado en la protección de rentas: los trabajadores del núcleo del mercado laboral (asalariados estables del sector formal, funcionarios, trabajadores de grandes empresas) disfrutan de prestaciones generosas, a veces más que en los países continentales, mientras que los trabajadores de la periferia (temporales, informales, jóvenes, mujeres con carreras interrumpidas) tienen una protección débil o nula. El segundo es la coexistencia de ese sistema bismarckiano de rentas con una sanidad universal de tipo beveridgiano, creada en las décadas de los setenta y los ochenta. El tercero es la ausencia histórica de una red estatal de renta mínima, que dejaba la última protección en manos de la familia, la caridad o las administraciones locales. El cuarto es una cierta penetración del clientelismo y del particularismo en la distribución de las prestaciones, especialmente de las de invalidez.\footnote{%
    En España, MORENO (2000, 2001) matizó esta tesis con la noción de una vía media. Un modelo que comparte los rasgos mediterráneos, pero que ha seguido una trayectoria de convergencia gradual con los modelos continentales y, en algunos aspectos, como la sanidad universal, con los socialdemócratas. Otros autores, como GAL (2010), han discutido si el modelo meridional es un régimen propio o una variante tardía y menos desarrollada del conservador, y han propuesto hablar de una familia extendida de Estados del bienestar mediterráneos. CASTLES y MITCHELL (1993) defendieron, desde otra perspectiva, un cuarto mundo \textit{antípoda}, propio de Australia y Nueva Zelanda, que alcanzaba la igualdad por la vía de la regulación salarial más que por la de las transferencias. ARTS y GELISSEN (2002), en una revisión sistemática de la literatura, concluyeron que la tipología conserva valor heurístico, pero que ningún país real se corresponde limpiamente con un tipo ideal y que casi todos son híbridos.
}

La segunda gran crítica procede de la literatura feminista y ha tenido consecuencias teóricas profundas. LEWIS (1992) mostró que los Estados del bienestar se habían construido, en grados distintos, sobre el modelo del varón sustentador (male breadwinner): el hombre trabaja en el mercado y accede a los derechos sociales como trabajador asegurado, mientras la mujer se ocupa del hogar y de los cuidados y accede a la protección solo de forma derivada, como esposa o viuda. ORLOFF (1993) sostuvo que el concepto de desmercantilización era ciego a esta realidad: para muchas mujeres el problema no era liberarse del mercado, sino poder acceder a él; su trabajo de cuidados no remunerado quedaba fuera del análisis. ESPING-ANDERSEN (1999) aceptó la crítica e incorporó una cuarta dimensión, la desfamiliarización: el grado en que el Estado o el mercado liberan a los individuos de la dependencia de su familia para su bienestar, mediante servicios de cuidado infantil, atención a la dependencia o prestaciones individuales en lugar de derivadas. (OJO: la desfamiliarización es la clave conceptual para entender tres fenómenos del sistema español que se examinan más adelante. El primero es la llegada tardía de la protección pública de la dependencia, en 2006. El segundo es el hecho de que, en el Sistema para la Autonomía y Atención a la Dependencia, la prestación económica para cuidados en el entorno familiar, concebida como excepcional, se convirtiera durante años en la más frecuente. El tercero es la persistencia de una importante brecha de género en las pensiones, que el complemento creado en 2021 intenta corregir. También explica el peso singular de la pensión de viudedad en España, uno de los rasgos diferenciales de nuestro gasto social, y cuyo sentido se discute más adelante.

La tercera crítica procede de la economía política comparada y aporta una lógica económica que la tipología de ESPING-ANDERSEN, más sociológica, no explicita. ESTEVEZ-ABE, IVERSEN y SOSKICE (2001), en el marco de la teoría de las variedades de capitalismo de HALL y SOSKICE (2001), sostuvieron que la protección social no es solo un instrumento de redistribución o de desmercantilización, sino también un complemento institucional del sistema productivo. En las economías que compiten mediante trabajadores con cualificaciones específicas de una empresa o de un sector, como la alemana, esos trabajadores solo invertirán en adquirirlas si están protegidos frente al riesgo de perder su empleo y de tener que aceptar otro en el que esas cualificaciones no valgan nada. La protección del empleo, un seguro de desempleo generoso y proporcional al salario y unas pensiones que reproducen el estatus son, desde esta perspectiva, un seguro de la inversión en capital humano específico. En las economías que compiten con cualificaciones generales, como la estadounidense, esa protección es menos necesaria, porque las cualificaciones son transferibles. Esta vertiente ofrece una explicación económica muy razonable de por qué los sistemas bismarckianos han resistido tan bien: no son solo el producto de una coalición política, sino una pieza funcional de un modelo productivo, y su reforma afecta a los incentivos de trabajadores y empresas para invertir en formación.

#### Palier y el largo adiós a Bismarck

Si ESPING-ANDERSEN explica la diversidad de los Estados del bienestar, la literatura sobre su reforma explica cómo han cambiado desde la crisis de los años setenta. PALIER (2010) estudió las reformas de los sistemas continentales y propuso una secuencia típica en tres fases.\footnote{%
    España encaja en esta secuencia casi al milímetro. La primera fase tuvo en nuestro país una versión propia, con las jubilaciones anticipadas y las prejubilaciones asociadas a las reconversiones industriales de los ochenta. La segunda está representada por la Ley 26/1985 y, sobre todo, por el Pacto de Toledo de 1995 y la Ley 24/1997: la separación de fuentes que el Pacto consagró como su primera recomendación es literalmente la clarificación de responsabilidades de PALIER. La tercera fase tuvo un intento, el factor de sostenibilidad y el índice de revalorización de 2013, que introdujo mecanismos automáticos de ajuste. Pero su suspensión en 2018 y su derogación en 2021 muestran algo que la secuencia de PALIER no anticipaba: la tercera fase no es irreversible. Un sistema puede introducir mecanismos estructurales y retirarlos pocos años después si carecen del consenso político necesario.
}

En la primera fase, que abarca los años setenta y ochenta, los sistemas bismarckianos respondieron al aumento del desempleo con políticas de reducción de la oferta de trabajo: jubilaciones anticipadas, pensiones de invalidez utilizadas como vía de salida del mercado laboral y prolongación de la escolarización. Estas medidas se financiaron con aumentos de las cotizaciones, lo que encareció el trabajo y alimentó un círculo vicioso de más paro, más salidas anticipadas y más cotizaciones, que ESPING-ANDERSEN llamó \textit{welfare without work}. 

En la segunda fase, en los años noventa, la presión fiscal y las exigencias de la unión monetaria obligaron a contener el gasto y a clarificar las responsabilidades financieras. Se endureció la relación entre cotizaciones y prestaciones, alargando los periodos de cálculo y aumentando los años de cotización exigidos, y se separaron las prestaciones contributivas, que seguirían financiándose con cotizaciones, de las no contributivas o solidarias, que pasarían a financiarse con impuestos.\footnote{%
    La lógica de esta separación era doble: (i) reforzar la legitimidad del seguro, al hacer más visible la correspondencia entre lo que se paga y lo que se recibe; y, (ii) aliviar el coste laboral, al retirar de las cotizaciones la financiación de prestaciones que no correspondían a los cotizantes.
} 

En la tercera fase, desde finales de los noventa, se introdujeron cambios estructurales que alteraron la lógica misma del sistema: cuentas nocionales, pilares de capitalización complementarios, mecanismos de ajuste automático ligados a la demografía y políticas de activación que condicionan las prestaciones a la búsqueda de empleo.



La dualización: quién soporta el ajuste

La literatura más reciente ha puesto el acento no en cuánto se ha reformado, sino en quién ha soportado el coste de las reformas. RUEDA (2007) y, de forma más sistemática, EMMENEGGER, HÄUSERMANN, PALIER y SEELEIB-KAISER (2012), sostienen que las reformas de los sistemas continentales y meridionales han seguido una lógica de dualización. Los \textit{insiders} (trabajadores con contratos estables, carreras completas y fuerte representación sindical) han conservado en buena medida sus derechos, mientras que el ajuste se ha concentrado en los \textit{outsiders} (trabajadores temporales, a tiempo parcial involuntario, autónomos dependientes, jóvenes y trabajadores con carreras interrumpidas). 

La dualización no es solo un efecto del mercado de trabajo, sino también de la política. Los partidos y sindicatos que representan sobre todo a los insiders tienen incentivos para proteger sus derechos y aceptar recortes que recaen en otros colectivos, lo que RUEDA formuló como un conflicto de intereses dentro de la propia socialdemocracia. En este sentido, España es un caso paradigmático de dualización. Durante décadas, hasta la reforma laboral de 2021, la contratación temporal concentró el riesgo de desempleo en jóvenes, mujeres y trabajadores de baja cualificación, con tasas de temporalidad que figuraron entre las más altas de la Unión Europea. Y las consecuencias de esa dualidad no se agotan en el mercado de trabajo, sino que se proyectan décadas después sobre el sistema de pensiones. Las carreras de cotización intermitentes y con bases bajas se traducen en pensiones reducidas, en lagunas de cotización y en dependencia de los complementos a mínimos.\footnote{%
    \textbf{NOTA.-} Precisamente este es el vínculo entre el mercado de trabajo y la suficiencia de las pensiones que la descomposición del gasto en pensiones que veremos captura solo parcialmente a través del factor empleo, y que conviene tener presente al valorar las reformas que endurecen la relación entre cotizaciones y prestaciones. Cuanto más contributivo es un sistema, más traslada a la vejez las desigualdades de la vida activa.
}



\textcolor{blue}{Reunidas las piezas, es posible situar el sistema español en el mapa de los regímenes de bienestar con cierta precisión. España combina un núcleo bismarckiano en las pensiones y en la protección por desempleo, heredado del sistema de seguros sociales y consolidado por la Ley General de la Seguridad Social de 1966; una sanidad beveridgiana desde 1986; una tradición familiarista en los cuidados, corregida solo parcialmente con la creación del Sistema para la Autonomía y Atención a la Dependencia en 2006; una red asistencial tardía y fragmentada, construida por capas sucesivas (pensiones no contributivas en 1990, rentas mínimas autonómicas desde finales de los ochenta, IMV en 2020); y un dualismo laboral persistente, que se proyecta sobre la protección social. En términos de ESPING-ANDERSEN, es un régimen de raíz conservadora con los rasgos característicos del modelo meridional de FERRERA, en tránsito hacia un modelo más universalista. En términos de PALIER, ha recorrido con fidelidad las dos primeras fases de la reforma bismarckiana y ha retrocedido en la tercera. 

Esa trayectoria, con sus avances y sus retrocesos, es la que se reconstruye históricamente en el subapartado siguiente, y es también la que explica la composición singular de su gasto social: elevado en pensiones y desempleo, intermedio en sanidad y bajo en familia, vivienda y exclusión social.}



\subsection{Genealogía: de la caridad al seguro, y del seguro a la ciudadanía}

La historia de la protección social no avanza en línea recta desde la beneficencia hasta el Estado del bienestar maduro. Es una secuencia de soluciones parciales, cada una de las cuales resolvió un problema y generó otro. La asistencia pública a los pobres resolvió el problema de la indigencia extrema, pero lo hizo a costa de estigmatizar al asistido y de disuadir a quien podía trabajar. La previsión voluntaria mutualista evitó ese estigma, pero sucumbió a la selección adversa y a la insolvencia. El seguro social obligatorio bismarckiano resolvió la selección adversa, pero dejó fuera a quien no tenía un empleo asalariado estable. La universalización beveridgiana incorporó a todos los ciudadanos, pero se apoyó en el supuesto del pleno empleo y en unas tasas de natalidad que no iban a durar. 

Las reformas de las últimas cuatro décadas han intentado, con éxito desigual y reversiones frecuentes, adaptar ese edificio a un mundo de desempleo estructural, envejecimiento demográfico y restricciones fiscales. Leída así, la historia no es un catálogo de antecedentes, sino el mejor argumento para entender por qué el sistema actual es como es y por qué su reforma resulta tan difícil.

\subsubsection{El mundo: de la Ley de Pobres a la guerra de paradigmas}
La primera respuesta pública organizada a la pobreza en la vejez, la enfermedad o la falta de trabajo no fue el seguro, sino la asistencia. La Ley de Pobres isabelina de 1601 obligó a cada parroquia inglesa a sostener a sus pobres con un impuesto local, distinguiendo entre los pobres impotentes, como ancianos, enfermos o huérfanos, a los que se socorría, y los pobres capaces, a los que se obligaba a trabajar. 

El sistema combinaba ya las dos preocupaciones que acompañarían a toda la historia de la protección social: evitar la miseria extrema y evitar que la asistencia desincentivara el trabajo. Durante dos siglos, esa tensión se resolvió de forma descentralizada, parroquia por parroquia, con enormes diferencias locales.

A finales del siglo XVIII, en plena revolución industrial y con una fuerte subida del precio del grano, los magistrados de Speenhamland, en Berkshire, adoptaron en 1795 un sistema de prestaciones que completaba los salarios bajos percibidos por los trabajadores hasta un mínimo ligado al precio del pan y al tamaño de la familia; sistema que se extendió por buena parte del sur de Inglaterra.\footnote{%
    POLANYI (1944) convertiría este suceso en el eje de una de las interpretaciones más influyentes de la historia económica moderna. Según POLANYI, Speenhamland fue un intento de proteger a la sociedad frente a la conversión del trabajo en una mercancía sometida a las leyes de la oferta y la demanda. Su fracaso, que atribuía a que deprimió los salarios al permitir a los patronos pagar menos sabiendo que la parroquia completaría la diferencia, abrió el camino a la plena mercantilización del trabajo. La interpretación de POLANYI sobre los efectos de Speenhamland ha sido muy discutida por los historiadores económicos, que han encontrado efectos menos dramáticos de los que él suponía. Pero su tesis general, según la cual la economía de mercado tiende a mercantilizar el trabajo y la sociedad reacciona con un doble movimiento de autoprotección, sigue siendo el fundamento intelectual del concepto de desmercantilización, muy recurrente desde entonces.
} No obstante, la reacción en contra fue la Nueva Ley de Pobres de 1834, que introdujo el principio de \textit{less eligibility}: la situación del asistido debía ser siempre menos deseable que la del trabajador independiente más pobre. El instrumento fue el \textit{workhouse}, el asilo de trabajo, donde el pobre capaz solo recibía socorro a cambio de internarse en condiciones deliberadamente duras, separado de su familia y sometido a una disciplina estricta. Es la primera formulación institucional explícita del dilema entre protección y riesgo moral y el precio que se pagó fue el estigma, que TITMUSS denunciaría un siglo después. Las novelas de Dickens, con Oliver Twist (1838) como emblema, construyeron la imagen de este sistema que ha llegado hasta nosotros. Y, como en todo, el eco de la \textit{less eligibility} resuena todavía en los debates contemporáneos sobre la compatibilidad de las rentas mínimas con el empleo, incluido el diseño del IMV.

Junto a la asistencia pública, y en buena medida como reacción contra su carácter estigmatizante, se desarrolló la previsión voluntaria. Los gremios y cofradías medievales habían ofrecido ya socorros a sus miembros en caso de enfermedad o muerte. En el siglo XIX, con la industrialización, surgieron las \textit{friendly societies} británicas, las cajas de socorro alemanas (\textit{Hilfskassen}), las sociedades de socorros mutuos francesas y las mutualidades obreras en toda Europa. En España, la monarquía ilustrada del siglo XVIII había creado montepíos para militares y funcionarios (antecedente directo de las Clases Pasivas). Estas instituciones tenían un mérito evidente frente a la asistencia: sus miembros cobraban como asegurados y no como pobres, con dignidad y sin estigma. Pero padecían dos debilidades estructurales. Al ser voluntarias, adolecían de problemas de selección adversa: se afiliaban sobre todo quienes esperaban necesitar las prestaciones, y quedaban fuera los trabajadores más pobres, incapaces de pagar las cuotas. Al ser pequeñas y estar vinculadas a un oficio o a una localidad, eran vulnerables a los riesgos agregados (epidemias, crisis de un sector, inflación), que afectaban a todos sus miembros a la vez y las conducían a la insolvencia. 

La historia de la previsión voluntaria del siglo XIX es, en buena medida, una historia de quiebras, y ese fracaso preparó el terreno para la obligatoriedad.

El primer seguro social: la Alemania de Bismarck

El primer seguro social de la historia se sitúa en la Alemania de BISMARCK: un programa de seguro social de enfermedad y vejez para los trabajadores. Este programa se anunció en el mensaje imperial de 17 de noviembre de 1881 (Kaiserliche Botschaft), leído ante el Reichstag en nombre del emperador Guillermo I, y se materializó en tres leyes sucesivas: el seguro de enfermedad (1883), el seguro de accidentes de trabajo (1884) y el seguro de invalidez y vejez (1889). 

La motivación fue, ante todo, política. En 1878 BISMARCK había ilegalizado la socialdemocracia mediante las leyes antisocialistas (\textit{Sozialistengesetze}), aprovechando dos atentados contra el emperador. El seguro social fue la otra cara de esa represión: la zanahoria que acompañaba al palo. El objetivo explícito era integrar a la clase obrera industrial en el Estado imperial, apartarla del socialismo y ligar su bienestar a la supervivencia del régimen, convirtiendo al trabajador en un pequeño rentista del Estado con un interés material en su estabilidad. BISMARCK lo denominó \textit{Staatssozialismus}, socialismo de Estado, y en términos de la economía política moderna es un ejemplo temprano y explícito de compra de lealtad mediante transferencias.\footnote{%
    \textbf{NOTA.-} Curiosamente, a diferencia de los anarquistas, los marxistas admiten la lucha por las \textit{reformas}, es decir, por mejoras de la situación de los trabajadores que no lesionan el poder, dejándolo como estaba, en manos de la clase dominante. Pero, a la vez, los marxistas combaten con la mayor energía a los \textit{reformistas}, los cuales circunscriben directa o indirectamente los anhelos y la actividad de la clase obrera a las reformas. Para éstos, el reformismo es una manera que la burguesía tiene de engañar a los obreros, que seguirán siendo esclavos asalariados, pese a algunas mejoras aisladas, mientras subsista el dominio del capital. Vladímir LENIN consideraba que el reformismo era una herramienta de la burguesía para engañar a la clase trabajadora y apartarla de la lucha revolucionaria, palabra para la que acuñó el término \textit{engaño burgués}.
}

Tampoco debemos perder de vista el contexto intelectual contemporáneo. Los socialistas de cátedra (\textit{Kathedersozialisten}) agrupados desde 1872 en la \textit{Verein für Socialpolitik} (SCHMOLLER, WAGNER, BRENTANO) defendían la intervención del Estado para resolver la \textit{cuestión social} sin recurrir a la revolución. A WAGNER, en particular, formuló su célebre ley de expansión creciente de la actividad del Estado, que el opositor conoce de otros temas del temario y que tiene en el seguro social uno de sus ejemplos más claros.

El modelo bismarckiano presentaba tres rasgos que han perdurado hasta hoy en los sistemas continentales, incluido el español. El primero era la \textbf{obligatoriedad} para los asalariados por debajo de un umbral de renta, financiada con cotizaciones de empresarios y trabajadores. En el caso del seguro de vejez e invalidez, esas cotizaciones se complementaban con una subvención fija del Estado por pensión. El segundo eran las prestaciones \textbf{proporcionales} al salario, con el objetivo de mantener el estatus y no de garantizar un mínimo. El tercero era la \textbf{gestión descentralizada por cajas}, con participación de empresarios y trabajadores.\footnote{%
    Esta gestión corporativa tuvo un efecto que BISMARCK no había previsto: las cajas de enfermedad se convirtieron en uno de los principales espacios de formación de cuadros de la socialdemocracia alemana, lo que muestra que las instituciones de protección social tienen efectos políticos que a menudo escapan a las intenciones de sus creadores.
}

Conviene, además, tener en cuenta que la edad de jubilación fijada en 1889 era de setenta años, y no se rebajó a sesenta y cinco hasta 1916. Con la esperanza de vida de la época, muy pocos trabajadores industriales alcanzaban esa edad, y quienes lo hacían cobraban la pensión durante pocos años; por eso la mayor parte del gasto inicial del seguro correspondió a las pensiones de invalidez. La edad de sesenta y cinco años, que el siglo XX convirtió en una norma casi universal, se fijó cuando llegar a ella era la excepción. Hoy, en cambio, un español que cumple sesenta y cinco años puede esperar vivir en torno a otros veinte. Este contraste es el trasfondo histórico del debate sobre la conveniencia de ligar la edad de jubilación a la esperanza de vida: la edad legal no tiene un fundamento natural, sino que es una convención fijada en un contexto demográfico radicalmente distinto del actual.

\paragraph{Difusión y variantes: del seguro obrero a la seguridad social (1889-1939)}
El modelo alemán se difundió con rapidez, pero no siempre en su versión bismarckiana. La encíclica \textit{Rerum Novarum} de León XIII (1891) dio un impulso decisivo a la reforma social en los países católicos. Defendía la intervención del Estado para proteger a los obreros y la creación de asociaciones de previsión, al tiempo que rechazaba el socialismo. Su principio de subsidiariedad, formulado de forma más explícita en la Quadragesimo Anno de 1931, según el cual el Estado solo debe intervenir cuando las comunidades inferiores (familia, asociación, municipio) no pueden cumplir su función, se convirtió en uno de los rasgos definitorios del régimen conservador de ESPING-ANDERSEN. 

Austria adoptó seguros de accidentes y enfermedad en 1887 y 1888. Pero otros países optaron por una vía diferente, la asistencial. Dinamarca, en 1891, y Nueva Zelanda, en 1898, crearon pensiones de vejez no contributivas, financiadas con impuestos generales y sujetas a una comprobación de recursos. La razón era en parte económica y en parte social. En países con una población mayoritariamente agraria y pocos asalariados industriales, un seguro obrero de tipo alemán habría dejado fuera a la mayor parte de la población. Suecia creó en 1913 la primera pensión de vejez de carácter universal, que cubría a toda la población y no solo a los asalariados; el antecedente remoto del universalismo que caracterizaría después al régimen socialdemócrata.

El Reino Unido combinó ambas vías en dos pasos. La \textit{Old Age Pensions Act} de 1908, impulsada por el gobierno liberal y por su ministro de Hacienda, LLOYD GEORGE (al que asesoraba KEYNES en su etapa temprana), concedió una pensión no contributiva de cinco chelines semanales a los mayores de setenta años con recursos insuficientes y \textit{buen carácter}.\footnote{%
    La condición de buen carácter excluía, por ejemplo, a quienes hubieran estado en prisión recientemente o hubieran rehuido habitualmente el trabajo, y es un ejemplo de la condicionalidad moral que acompañó a la asistencia durante mucho tiempo.
} La \textit{National Insurance Act} de 1911 creó un seguro contributivo de enfermedad y el primer seguro obligatorio de desempleo del mundo, limitado inicialmente a algunos sectores especialmente expuestos al paro cíclico, como la construcción naval o la ingeniería. Francia aprobó en 1910 las \textit{retraites ouvrières et paysannes}, un seguro de vejez obligatorio de capitalización, pero su aplicación fue muy deficiente: los tribunales admitieron que la obligación no era exigible y la inflación posterior a la Primera Guerra Mundial destruyó el valor de los fondos. Finalmente, la Organización Internacional del Trabajo, convirtió el seguro social en un estándar internacional y aprobó durante el período de entreguerras una serie de convenios sobre seguros de enfermedad, vejez, invalidez y supervivencia.\footnote{%
    El lector atento habrá notado que, en esto, subyacía un argumento de economía política internacional que sigue vigente: \textbf{ningún país quería asumir en solitario los costes laborales de la protección social si sus competidores no lo hacían, de modo que la armonización internacional de los estándares mínimos protegía a los países más avanzados de la competencia a la baja.}
}

La Gran Depresión abrió la segunda gran oleada. En Estados Unidos, la \textit{Social Security Act} de 14 de agosto de 1935, aprobada en el marco del New Deal, creó un seguro federal de vejez, un sistema de seguro de desempleo gestionado por los Estados y varios programas asistenciales, y consagró jurídicamente la expresión \textit{social security}.\footnote{%
    La primera beneficiaria de una pensión mensual de la Seguridad Social de Estados Unidos, Ida May Fuller, una secretaria jubilada de Vermont, recibió su primer cheque en enero de 1940, por $22,54$ dólares. Había cotizado en total $24,75$ dólares durante los tres años en que estuvo sujeta al sistema. Vivió hasta los cien años y cobró en total unos $22.900$ dólares. Como vemos, este episodio ilustra perfectamente el concepto ganancia de lanzamiento (\textit{windfall gain}), ganancia que disfrutan las primeras generaciones de un sistema de reparto, que reciben pensiones completas habiendo cotizado solo una fracción de su vida laboral. Sin embargo, esa ganancia se financia con las cotizaciones de las generaciones siguientes, y su análisis formal, desde el modelo de SAMUELSON (1958), lo veremos más adelante.
} Nueva Zelanda aprobó en 1938 su propia \textit{Social Security Act}, que integró los distintos programas en un sistema unificado financiado con un impuesto específico y que suele considerarse el primer sistema de seguridad social en sentido moderno, anterior incluso a BEVERIDGE.

Beveridge y la edad de oro (1942-1973)

Con todo, vemos que el modelo de seguridad social predominante, casi único, durante el primer periodo de la provisión social fue de carácter eminentemente bismarckiano. Con alguna excepción, como las mencionadas, la \textbf{universalidad} y el derecho de acceso por \textbf{ciudadanía} no eran los elementos definitorios del sistema. Estas dos características llegaron durante el transcurso de la Segunda Guerra Mundial, a raíz de la publicación del más que célebre Informa BEVERIDGE (1942).\footnote{%
    Al Informe BEVERIDGE (1942) se le atribuye la creación en el Reino Unido del primer sistema unificado de Seguridad Social, que abrió paso a los sistemas modernos. El informe, titulado \textit{Social Insurance and Allied Services}, se publicó en diciembre de 1942, en plena guerra y pocas semanas después de la victoria de El Alamein, y fue un fenómeno editorial y político extraordinario: se vendieron más de seiscientos mil ejemplares y se difundió entre las tropas y en los territorios ocupados como símbolo de la sociedad por la que se luchaba. Su publicación tuvo un propósito explícito de movilización: ofrecer a una población exhausta la promesa de que la victoria no traería el regreso al desempleo masivo y la miseria de los años treinta.
} BEVERIDGE identificó cinco gigantes que la reconstrucción de posguerra debía derrotar: la necesidad (Want), la enfermedad (Disease), la ignorancia (Ignorance), la miseria (Squalor) y la ociosidad (Idleness). El seguro social se configuraba como instrumento contra el primero; los demás exigían un servicio de salud, un sistema educativo, una política de vivienda y una política de empleo. 

Su propuesta de seguro se basaba en seis principios: prestaciones uniformes de subsistencia, cotizaciones uniformes, unificación de la responsabilidad administrativa, suficiencia de las prestaciones, exhaustividad en la cobertura de riesgos y clasificación de la población según sus circunstancias. La expresión que resume su ambición, una protección \textit{de la cuna a la tumba}, procede del propio debate de la época.

El plan se apoyaba explícitamente en tres supuestos sin los cuales no funcionaría: la existencia de asignaciones familiares para los hijos (supuesto A), un servicio de salud integral y gratuito (supuesto B) y el mantenimiento del empleo, es decir, la prevención del desempleo masivo (supuesto C). El supuesto C es el más importante para entender la historia posterior. BEVERIDGE contaba con que la política económica keynesiana garantizaría el pleno empleo, y dedicó a ello un segundo informe, \textit{Full Employment in a Free Society} (1944). Con un paro bajo, el seguro de desempleo sería barato, las cotizaciones podrían ser moderadas y el sistema se financiaría holgadamente. Sin embargo, cuando el pleno empleo desapareció en los años setenta, el edificio perdió uno de sus pilares, y buena parte de la crisis del Estado del bienestar que se examina más adelante puede leerse como la consecuencia del incumplimiento del supuesto C.

El plan se materializó en la \textit{National Insurance Act} de 1946 y en la creación del National Health Service en 1948, bajo el gobierno laborista de ATTLEE. Como vemos, el NHS, financiado con impuestos y no con cotizaciones, se convirtió en el modelo de los sistemas sanitarios beveridgianos, y es el antecedente del Sistema Nacional de Salud español creado en 1986.

Como decíamos, la posguerra trajo también la consagración internacional de la seguridad social como derecho. La Declaración de Filadelfia de la OIT (1944) proclamó la extensión de la seguridad social para garantizar ingresos básicos a quienes los necesitaran. La DUDHH (1948) reconoce en su art. $22$ el derecho de toda persona, como miembro de la sociedad, a la seguridad social, y en su art. $25$ el derecho a un nivel de vida adecuado y a los seguros en caso de desempleo, enfermedad, invalidez, viudez y vejez. El Convenio $102$ de la OIT (1952), sobre la norma mínima de la seguridad social, definió nueve ramas de protección (asistencia médica, enfermedad, desempleo, vejez, accidentes de trabajo, familia, maternidad, invalidez y supervivencia) y fijó niveles mínimos de prestación para cada una. Su contenido y su valor para el ordenamiento español se examinan más adelante. El Pacto Internacional de Derechos Económicos, Sociales y Culturales (1966) reconoció en su art. $9$ el derecho a la seguridad social. En definitiva, la traducción jurídica de la ciudadanía social de MARSHALL, cuya conferencia de 1949 se publicó en 1950: el derecho a la protección se desliga de la condición de trabajador y se vincula a la condición de persona.

Las tres décadas siguientes, las \textit{trente glorieuses} en expresión de FOURASTIÉ, fueron la edad de oro de la expansión de los Estados del bienestar, favorecida por un crecimiento económico excepcional, el pleno empleo, una demografía favorable tras el baby boom y un amplio consenso político entre democristianos, socialdemócratas y liberales. Dos episodios ilustran bien el espíritu de la época y sus consecuencias a largo plazo.

El primero es la reforma alemana de 1957. El gobierno de Adenauer abandonó los restos de capitalización del sistema bismarckiano, cuyos fondos habían sido destruidos por dos hiperinflaciones y dos guerras, y adoptó un sistema de reparto con una pensión dinámica, indexada al crecimiento de los salarios. Las pensiones aumentaron de golpe en torno a un $60\%$, a pocos meses de unas elecciones que la CDU ganó con mayoría absoluta. Cuando el ministro de Economía, ERHARD, y otros críticos le advirtieron del riesgo demográfico que suponía hacer depender las pensiones de las cotizaciones de las generaciones futuras, Adenauer respondió con una frase que ha pasado a la historia: \textit{la gente siempre tiene hijos} (Kinder kriegen die Leute immer). Pocos años después, la natalidad alemana inició un descenso que la situaría entre las más bajas del mundo y el resto, como se suele decir, es historia. \textit{La lógica económica del reparto depende del crecimiento de la población y de los salarios}.

El segundo es la reforma sueca de 1959-1960, que creó una pensión complementaria obligatoria proporcional al salario, la ATP (allmän tilläggspension), gestionada públicamente y dotada de grandes fondos de reserva. Su aprobación exigió un referéndum consultivo, que ninguna de las tres opciones ganó por mayoría absoluta, povocando la ruptura de la coalición de gobierno entre socialdemócratas y agrarios, y unas elecciones anticipadas, y el proyecto se aprobó finalmente en 1959 por un solo voto de diferencia.\footnote{
El episodio ilustra hasta qué punto las pensiones son un asunto de alta política, capaz de hacer caer gobiernos; y, además, ostenta un significado teórico muy preciso: los socialdemócratas suecos incorporaron a las clases medias asalariadas a un sistema público con prestaciones proporcionales al salario. Estrategia que KORPI y PALME identificarían cuatro décadas después como la clave de la paradoja de la redistribución: ampliar la coalición que sostiene el Estado del bienestar incluyendo en él a las clases medias.
}

En los Estados Unidos de la Great Society, la creación de Medicare y Medicaid en 1965 completó el sistema con la protección sanitaria de mayores y pobres. 

En definitiva, en toda Europa, los años sesenta y comienzos de los setenta fueron de extensión de la cobertura a nuevos colectivos (autónomos, agricultores, trabajadores domésticos), de aumento de la generosidad de las prestaciones y de indexación generalizada de las pensiones a los precios o a los salarios. Fue, en esencia, una época de \textit{reclamación de créditos}: los gobiernos competían por atribuirse la mejora de las prestaciones, en un contexto en el que el crecimiento permitía financiarlas sin recortar otros gastos ni subir de forma visible los impuestos.

Crisis y retrenchment (1973-1990)

Las crisis del petróleo de 1973 y 1979, la estanflación y el retorno del desempleo masivo pusieron fin a la edad de oro. El supuesto C de BEVERIDGE se incumplió de forma persistente, y los Estados del bienestar se encontraron con un aumento simultáneo de sus gastos, por el paro y por las prestaciones asociadas, y con una reducción de sus ingresos, por la caída del empleo y del crecimiento. Además, y más notablemente, la crisis fue también intelectual: el consenso keynesiano se quebró, y ganaron influencia las críticas monetaristas y de la elección pública. Para estas, el Estado del bienestar generaba desincentivos al trabajo y al ahorro, una presión fiscal excesiva y una tendencia al crecimiento incontrolado del gasto, derivada de la lógica electoral y burocrática.

La respuesta, no obstante, no fue uniforme. Por un lado, los sistemas continentales adoptaron una vía de escape social (PALIER): reducir la oferta de trabajo mediante jubilaciones anticipadas y pensiones de invalidez utilizadas como vía de salida del mercado laboral y, de esta forma, el desempleo masivo.\footnote{%
    El caso más extremo fue el de los Países Bajos, donde el número de beneficiarios de pensiones de incapacidad laboral se aproximó al millón de personas a comienzos de los noventa, en un país de unos quince millones de habitantes. El primer ministro Ruud Lubbers llegó a declarar en 1990 que Holanda está enferma. Un ejemplo de cómo una prestación concebida como seguro contra una contingencia puede convertirse en un instrumento de gestión del desempleo, debido a los problemas implícitos de riesgo moral que trae aparejada.
} Por otro lado, el mundo anglosajón optó por la vía económica: la contención del gasto, o retrenchment.\footnote{%
    El gobierno de Thatcher ligó en 1980 la revalorización de la pensión básica británica a la evolución de los precios, y no de los salarios, lo que redujo progresivamente su valor en relación con el nivel de vida medio de la población. En 1986 recortó la pensión estatal proporcional al salario (SERPS) y ofreció incentivos para que los trabajadores la abandonaran a favor de planes de pensiones personales privados. Decisión que desembocaría en uno de los mayores escándalos financieros de la historia británica, el \textit{pension mis-selling scandal}: cientos de miles de trabajadores fueron inducidos a abandonar planes de empresa ventajosos por planes personales que les convenían menos, y las compensaciones posteriores alcanzaron miles de millones de libras. Un claro ejemplo de los problemas de información y de costes de gestión que presentan los sistemas privados, sobre los que se basa la capitalización. Paralelamente, en Estados Unidos, la Comisión GREENSPAN propuso en 1983 una reforma aprobada con apoyo de ambos partidos que elevaba gradualmente la edad de jubilación completa de sesenta y cinco a sesenta y siete años, entre 2000 y 2027, y adelantaba subidas de cotizaciones para acumular un gran fondo de reserva ante la jubilación del baby boom.
}

El estudio más influyente de estos años es el de PIERSON (1994, 1996), que analizó el retrenchment bajo Reagan y Thatcher y llegó a una conclusión contraintuitiva: pese a su retórica y a sus mayorías parlamentarias, ambos gobiernos consiguieron recortar el Estado del bienestar mucho menos de lo que se proponían. Su explicación, conocida como la tesis de la nueva política del Estado del bienestar, tiene tres elementos. El primero es que los programas de protección social crean sus propios grupos de apoyo, formados por los beneficiarios, los proveedores y las organizaciones que los representan, que no existían cuando los programas se crearon y que se movilizan para defenderlos. El segundo es que la política de expansión y la de recorte son asimétricas: expandir un programa produce beneficios concentrados y visibles, mientras que recortarlo impone pérdidas concentradas e inmediatas a cambio de ganancias difusas y futuras, de modo que el recorte es mucho más costoso electoralmente. El tercero es que, en consecuencia, los gobiernos que recortan recurren a estrategias de ocultación: cambios técnicos en las fórmulas de indexación, reformas con largos periodos transitorios que retrasan sus efectos hasta que los gobernantes que las aprobaron han dejado el poder, o recortes que afectan a los futuros beneficiarios y no a los actuales.

Esta lente es imprescindible para leer la historia española que sigue. Casi todas las reformas restrictivas de las pensiones en España, de la Ley 26/1985 a la Ley 27/2011, han utilizado periodos transitorios de diez o quince años. La reforma de 2013, en cambio, afectaba a todos los pensionistas actuales mediante la revalorización mínima, y su historia confirmará la predicción de PIERSON.

\subsubsection{Guerra de paradigmas: capitalización frente a reparto (1981-2008)}
A partir de los años ochenta, el debate dejó de ser solo cuánto gastar y pasó a ser cómo organizar la protección de la vejez. Durante un cuarto de siglo, la política de pensiones en el mundo estuvo dominada por una confrontación entre dos paradigmas: el reparto público heredado de la posguerra y la capitalización individual privada.\footnote{%
    El punto de partida fue Chile. El documento recoge en nota que el modelo más opuesto al español es el chileno, implantado bajo Pinochet por impulso del entonces ministro José Piñera. Sustituyó por completo el antiguo reparto por cuentas individuales de capitalización gestionadas por Administradoras de Fondos de Pensiones privadas (AFP), sirvió durante décadas de referencia internacional para los defensores de la capitalización y ha sido objeto en los últimos años de una contestación social muy intensa por la insuficiencia de las pensiones finalmente percibidas, hasta forzar reformas sustanciales. Se financiaba con una cotización del diez por ciento del salario, a cargo exclusivamente del trabajador y sin aportación del empleador. Los trabajadores que ya cotizaban al sistema antiguo podían elegir entre permanecer en él o trasladarse al nuevo, y recibían un bono de reconocimiento por las cotizaciones pasadas. Destaca el hecho que las Fuerzas Armadas y Carabineros quedaran excluidos de la reforma y conservaron su propio sistema de reparto. Es decir, los propios promotores del régimen no quisieron para sí el sistema que impusieron al resto. Un dato relevante a la luz de la tesis de PIERSON: una transición radical del reparto a la capitalización es extraordinariamente difícil en democracia.
}

El primer hito de gran influencia teórica fue el informe del Banco Mundial de 1994, \textit{Averting the Old Age Crisis}. El informe diagnosticaba que los sistemas públicos de reparto eran insostenibles ante el envejecimiento de la población, distorsionaban el mercado de trabajo con cotizaciones elevadas, reducían el ahorro nacional (con referencia expresa al trabajo de FELDSTEIN, 1974) y producían redistribuciones arbitrarias e incluso regresivas. El Informe proponía un sistema de tres pilares: (i) un pilar público, obligatorio, financiado con impuestos y de carácter redistributivo, con prestaciones modestas destinadas a prevenir la pobreza; (ii) un pilar obligatorio, gestionado privadamente y capitalizado, que proporcionaría la mayor parte de la pensión en función de lo ahorrado por cada trabajador; y, (iii) un pilar voluntario de ahorro complementario. Según el Informe, la capitalización obligatoria (el pilar más importante de todos), aumentaría el ahorro y la inversión, profundizaría los mercados financieros de los países en desarrollo, ofrecería mayores rentabilidades y protegería las pensiones de la manipulación política. El Informe tuve un triunfo notable: inspiró una oleada de reformas en América Latina durante los años noventa, en Perú en 1993, Colombia y Argentina en 1994 y México en 1997, entre otros, con sistemas totalmente privados o mixtos, y en Europa central y oriental, con Hungría en 1998 y Polonia en 1999, donde la transición al mercado ofrecía una oportunidad política excepcional para reformar los sistemas heredados del socialismo. 

Sin embargo, la réplica teórica más influyente llegó en 1999 y, paradójicamente, del propio Banco Mundial; cuando ORSZAG y STIGLITZ (economista jefe de la institución) publicaron \textit{Rethinking Pension Reform: Ten Myths About Social Security Systems}. Para los autores, todas las ventajas atribuidas a la capitalización individual eran mitos o dependían de supuestos que no se hacían explícitos. En primer lugar, la transición del reparto a la capitalización no aumenta el ahorro nacional si el coste de seguir pagando las pensiones de los jubilados actuales se financia con deuda pública, como ocurre en la práctica, porque el aumento del ahorro privado se compensa con la reducción del ahorro público. En segundo lugar, la mayor rentabilidad de los fondos privados frente al reparto no es una ganancia neta, sino la contrapartida de un mayor riesgo y, en buena medida, una ilusión contable que ignora el coste de la transición. En tercer lugar, los costes de gestión de los sistemas privados son muy superiores a los de los públicos y erosionan una parte significativa del ahorro de los trabajadores. Por último, la capitalización no resuelve el problema demográfico, porque, en último término, las pensiones de cualquier generación se pagan con la producción de los trabajadores activos en ese momento, sea mediante cotizaciones, en el reparto, o mediante la compra de los activos que los jubilados venden, en la capitalización.\footnote{%
    La propia posición del Banco Mundial, tras los experimentos de América Latina, evolucionó. HOLZMANN y HINZ (2005) ampliaron el esquema de 1994, a cinco pilares. Añadieron un pilar cero, no contributivo y destinado a combatir la pobreza en la vejez, y un pilar cuatro, que recogía el apoyo informal de la familia, el acceso a la vivienda y a la sanidad y otras fuentes de bienestar no monetarias. Mantenía el pilar público contributivo, el pilar obligatorio capitalizado y el pilar voluntario. El cambio fue más que terminológico: suponía la retirada implícita del dogma de 1994, porque el pilar capitalizado obligatorio dejaba de ser el centro del sistema y pasaba a ser una opción entre otras. Ya no había un modelo único, sino una caja de herramientas que cada país debía combinar según sus circunstancias. Un informe de evaluación independiente del propio Banco Mundial, publicado en 2006, fue aún más crítico con los resultados de las reformas que la institución había promovido, en particular por su escasa cobertura y sus elevados costes fiscales de transición.
}

En Europa occidental, en cambio, la influencia de ambos Informes fue mucho menor. La existencia de sistemas de reparto maduros, con enormes compromisos con los pensionistas actuales y futuros, hacía simplemente prohibitivo el coste de la transición.

No obstante, los acuciantes problemas financieros que acechaban el sistema provocaron provocaron el surgimiento de una tercera vía, casi como salvavidas en esta encrucijada de paradigmas, que ha resultado muy influyente en el debate regional (europeo): las \textbf{cuentas nocionales de contribución definida} (NDC, por sus siglas en inglés). Cada cotizante tiene una cuenta virtual en la que se anotan sus cotizaciones y que se revaloriza anualmente con un tipo de interés nocional, normalmente ligado al crecimiento de los salarios o de la masa salarial. En el momento de la jubilación, el saldo acumulado se divide por la esperanza de vida de su cohorte a esa edad, y el resultado es la pensión anual. Las cotizaciones no se invierten en ningún activo, sino que se destinan a pagar las pensiones corrientes, como en cualquier sistema de reparto; lo que cambia es la fórmula de cálculo, que establece una relación estrictamente actuarial entre lo aportado y lo recibido y ajusta automáticamente la pensión a la evolución de la esperanza de vida.\footnote{%
    Italia (Ley 335/1995) y Suecia (1994), adoptaron un sistema que sigue siendo de reparto pero que imita el funcionamiento de una cuenta de capitalización. No obstante, Suecia añadió en 2001 un mecanismo de equilibrio automático, conocido popularmente como bromsen, el freno, que reduce la revalorización de las cuentas y de las pensiones cuando los pasivos del sistema superan a sus activos. Se activó por primera vez en 2010, tras la crisis financiera, y provocó reducciones nominales de las pensiones que generaron una fuerte controversia política y obligaron a compensar a los pensionistas con rebajas fiscales. Letonia (1996) y Polonia (1999) adoptaron también el modelo, que el propio Banco Mundial acabó analizando con interés (HOLZMANN y PALMER, 2006). No es de extrañar que, hoy en día, las cuentas nocionales sean la alternativa más citada en el debate español.

Alemania, por su parte, siguió una vía intermedia entre la reforma paramétrica y el cambio de paradigma. La reforma Riester de 2001 introdujo un pilar de capitalización voluntario con subvenciones públicas, pensado para compensar una reducción gradual del nivel de la pensión pública. La reforma de 2004 incorporó a la fórmula de revalorización un factor de sostenibilidad (Nachhaltigkeitsfaktor), que reduce la revalorización cuando empeora la relación entre pensionistas y cotizantes. La reforma de 2007 elevó gradualmente la edad de jubilación de sesenta y cinco a sesenta y siete años, hasta 2029. El factor de sostenibilidad alemán es un antecedente directo del índice de revalorización español de 2013, pues ambos ligaban la revalorización al equilibrio financiero del sistema.
}

\paragraph{Reversiones, gran recesión e inversión social (2008-2019)}

La crisis financiera de 2008 golpeó a los sistemas de pensiones por dos vías. Por un lado, la caída de los mercados financieros redujo drásticamente el valor de los fondos de pensiones privados, lo que mostró de forma brutal el riesgo de mercado al que la capitalización expone a los trabajadores próximos a la jubilación. Por otro, la recesión y el aumento del desempleo redujeron los ingresos por cotizaciones de los sistemas de reparto, y la posterior crisis de deuda soberana en la zona del euro obligó a muchos gobiernos a recortar el gasto público, incluido el de pensiones.

Una de las consecuencias más notables fue la reversión de muchas privatizaciones. La OIT documentó en 2018 que, de treinta países que habían privatizado total o parcialmente sus sistemas de pensiones entre 1981 y 2014, dieciocho habían revertido total o parcialmente la privatización. Los casos más notorios fueron Argentina (2008), Hungría (2010) y Polonia (2014). Los motivos, no obstante, fueron diversos. En primer lugar, los elevados costes de gestión reducían las pensiones. La cobertura era baja entre los trabajadores informales, que en muchos de estos países eran mayoría. Y el coste fiscal de la transición agravaba el déficit público precisamente cuando la crisis y las reglas fiscales europeas exigían reducirlo.\footnote{%
    Este último motivo es especialmente significativo desde el punto de vista teórico. En los países de la Unión, las reglas de Maastricht computaban como déficit el coste de seguir pagando las pensiones de reparto, pero no registraban como activo el ahorro acumulado en los fondos privados obligatorios. Por lo que, en esencia, la capitalización era \textit{fiscalmente penalizadora}, y varios gobiernos la desmontaron para liberar recursos y reducir el déficit contable.
}

En Europa, la gran recesión trajo recortes directos de las pensiones en los países sometidos a programas de ajuste. Grecia sufrió desde 2010 sucesivas reducciones de las pensiones en vigor (Troika). Portugal congeló pensiones y aplicó recortes a las más elevadas. E Italia aceleró el aumento de la edad de jubilación y la ligó automáticamente a la esperanza de vida.\footnote{%
    La reforma italiana generó el problema de los esodati: trabajadores que habían aceptado salir anticipadamente de sus empresas contando con jubilarse pronto y que, tras el cambio de reglas, quedaron sin salario ni pensión durante años. En 2019, el gobierno italiano de coalición entre el Movimiento Cinco Estrellas y la Liga introdujo la Quota 100, que permitía jubilarse cuando la suma de la edad y los años cotizados alcanzaba cien, con un mínimo de sesenta y dos años de edad. Fue una reversión parcial de la reforma anterior que, como la derogación española de 2021, confirma que las reformas impuestas en momentos de crisis y sin consenso tienden a revertirse cuando cambian las condiciones políticas.
} En general, en Europa se aceleró la difusión de los mecanismos de ajuste automático, que ligan la edad de jubilación, la cuantía de la pensión inicial o la revalorización a la evolución de la esperanza de vida o del equilibrio financiero del sistema.\footnote{%
    Su atractivo es doble. Técnicamente, aseguran que el sistema se adapte a la demografía sin necesidad de reformas periódicas. Políticamente, permiten a los gobiernos evitar la culpa de los recortes, porque el ajuste se produce automáticamente y no por una decisión discrecional. Al mismo tiempo, la coordinación europea en materia de pensiones se intensificó. El Método Abierto de Coordinación en pensiones, iniciado tras el Consejo Europeo de Laeken (2001), ya había establecido objetivos comunes de suficiencia, sostenibilidad y modernización. En la década de 2010, el Semestre Europeo incorporó recomendaciones específicas sobre pensiones a cada Estado miembro, basadas en las proyecciones de los \textit{Ageing Reports} del Grupo de Trabajo sobre Envejecimiento y en los \textit{Pension Adequacy Reports}.
}\textsuperscript{,}\footnote{%
    En paralelo a esta agenda de contención, emergió en el debate académico un paradigma alternativo: la inversión social. Formulado por ESPING-ANDERSEN, GALLIE, HEMERIJCK y MYLES (2002), desarrollado por MOREL, PALIER y PALME (2012) y sistematizado por HEMERIJCK (2013, 2017), sostiene que el Estado del bienestar debe desplazar recursos desde la compensación de los riesgos, mediante pensiones y subsidios de desempleo, hacia la preparación de los individuos para afrontarlos, mediante educación infantil de calidad, formación a lo largo de la vida, políticas de conciliación y activación en el mercado de trabajo. El diagnóstico de fondo es que los Estados del bienestar de posguerra se diseñaron para proteger frente a los \textit{viejos riesgos sociales} de la sociedad industrial (vejez, enfermedad, desempleo cíclico del varón sustentador), pero no frente a los \textit{nuevos riesgos sociales} que definieron TAYLOR-GOOBY (2004) y BONOLI (2005): la dificultad para conciliar trabajo y familia, las familias monoparentales, la obsolescencia de las cualificaciones, el empleo precario y los cuidados de larga duración de una población envejecida. La crítica más citada a la inversión social es la de CANTILLON (2011), que identificó su riesgo de un \textit{efecto MATEO}, por el versículo según el cual \textit{al que tiene se le dará más}. Los servicios de inversión social, como las guarderías, la formación o las políticas de conciliación, tienden a ser aprovechados desproporcionadamente por las clases medias con mayor capital educativo y mejor información. Si su expansión se financia recortando las transferencias monetarias a los más pobres, el resultado puede ser un aumento de la pobreza. BONOLI, CANTILLON y VAN LANCKER (2017) documentaron empíricamente este efecto en el acceso a los servicios de cuidado infantil. El debate es importante para España, cuyo gasto social está fuertemente sesgado hacia la vejez y es comparativamente bajo en familia, infancia y vivienda. La inversión social sugiere reequilibrar esa composición. La crítica de CANTILLON advierte que el reequilibrio no puede hacerse a costa de la protección de rentas de los más vulnerables.
}


\paragraph{La década actual (2020-2026): pandemia, envejecimiento del baby boom y fragilidad política de las reformas}
La pandemia de 2020 revalorizó el seguro social como estabilizador automático y como instrumento de preservación del tejido productivo. 

Los esquemas de reducción temporal de jornada con apoyo público, cuyo modelo es el Kurzarbeit alemán y cuya versión española son los expedientes de regulación temporal de empleo (ERTE), se generalizaron en toda Europa y protegieron a decenas de millones de trabajadores. La Unión creó el instrumento SURE (Reglamento 2020/672), con una capacidad de hasta cien mil millones de euros en préstamos a los Estados miembros para financiar esos esquemas: un primer paso de solidaridad financiera europea en el ámbito de la protección social, tradicionalmente reservado a los Estados, y precedió al Mecanismo de Recuperación y Resiliencia. En España, la reforma de pensiones de 2021-2023 fue precisamente uno de los hitos comprometidos para recibir los fondos de ese Mecanismo.

Por lo que respecta a la vejez, la década está marcada por la llegada a la edad de jubilación de las generaciones del baby boom y por una creciente fragilidad política de las reformas. En Francia, la reforma de 2023, que elevaba la edad legal de jubilación de sesenta y dos a sesenta y cuatro años y el periodo de cotización exigido para la pensión completa a cuarenta y tres años, se aprobó sin votación en la Asamblea Nacional, tras meses de huelgas y manifestaciones masivas, y en 2025 se aprobó la suspensión de su vigencia hasta enero de 2028, como condición de los socialistas para no derribar al gobierno de Sébastien Lecornu; quedando así pendiente del resultado de las elecciones presidenciales de 2027. En Alemania, el Bundestag aprobó el 5 de diciembre de 2025 un paquete de pensiones que garantiza un nivel de pensión del $48\%$ hasta 2031, amplía desde 2027 el reconocimiento de periodos de crianza para las madres de hijos nacidos antes de 1992 y suprime la prohibición de volver a trabajar para el antiguo empleador después de la edad de jubilación. Se complementa con la llamada Aktivrente, que exime del impuesto sobre la renta una parte del salario de quienes sigan trabajando pasada la edad legal, desde 2026 (\textcolor{red}{su cuantía exacta está pendiente de verificar}). El gobierno deberá presentar antes de 2029 un informe sobre el sistema a partir de 2031, lo que en la práctica aplaza las decisiones difíciles.

Dinamarca, en el extremo opuesto, aprobó en mayo de 2025 elevar la edad de jubilación a setenta años en 2040, en aplicación de la regla que desde 2006 liga la edad legal a la evolución de la esperanza de vida. Será la edad de jubilación más alta de Europa, y su aprobación no estuvo exenta de controversia. (Ni siquiera un mecanismo automático aceptado durante casi dos décadas deja de generar tensiones cuando produce sus efectos...) Los Países Bajos, están transformando su enorme segundo pilar de empresa, uno de los más grandes del mundo en relación con el PIB, desde un modelo de prestación definida a otro de contribución definida colectiva, con una transición que debe completarse en 2028. El Reino Unido mantiene desde 2011 el llamado \textit{triple lock}, que garantiza la revalorización de la pensión estatal con el mayor de tres valores: la inflación, el crecimiento de los salarios o el $2,5\%$; un trinquete que solo puede subir y cuyo coste creciente es objeto de debate permanente que ningún gobierno se ha atrevido a suprimir, por la fuerza electoral de los pensionistas británicos.

La lección de la década, en definitiva, es que las reformas de pensiones han dejado de ser irreversibles. Francia suspende la suya a los dos años de aprobarla, Italia revirtió parcialmente la reforma Fornero, España derogó en 2021 los mecanismos de 2013 y Suecia compensó fiscalmente la activación de su freno automático. Hasta los años noventa, el problema central de la política de pensiones era técnico: diseñar reformas que garantizaran la sostenibilidad. Hoy es sobre todo político: diseñar reformas que sobrevivan a los cambios de gobierno, a la movilización de los afectados y a la fragmentación parlamentaria. 

\subsubsection{España: de la previsión social al sistema de Seguridad Social}
La historia de la protección social en España reproduce, con retraso y con rasgos propios, las etapas del apartado anterior. España llegó tarde al seguro social obligatorio y a la universalización, y ha construido su red asistencial cuando otros países ya la estaban reformando. Pero ha seguido con notable fidelidad la secuencia de reformas bismarckianas que describió PALIER y ha repetido, en sus reformas recientes, las dinámicas de reversión observadas en otros países. 

\paragraph{Primer período: la previsión social (1883-1963)}
La primera forma de protección pública de la vejez en España no se dirigió a los obreros, sino a los servidores de la Corona. La monarquía ilustrada del siglo XVIII creó una serie de montepíos para garantizar pensiones a las viudas y huérfanos de militares y funcionarios; el Montepío Militar se fundó en 1761, bajo Carlos III, y le siguieron otros para los ministerios civiles. Estos montepíos, financiados con descuentos sobre los sueldos y con aportaciones de la Hacienda, fueron absorbidos por el Estado a lo largo del siglo XIX, y sus obligaciones se integraron en el presupuesto como Clases Pasivas, con una regulación que se codificó en el Estatuto de Clases Pasivas de 1926. Por tanto, el régimen de Clases Pasivas es más antiguo que la propia Seguridad Social, y esa antigüedad explica que haya sobrevivido como régimen separado hasta hoy, financiado directamente con el presupuesto del Estado, aunque cerró su acceso a los nuevos funcionarios en 2011, que desde entonces se integran en el Régimen General.

La protección de los trabajadores industriales llegó con la \textit{cuestión social} del último cuarto del siglo XIX. El origen de las políticas de protección social en España se sitúa en la Comisión de Reformas Sociales (1883), encargada de estudiar las condiciones de vida y trabajo de la clase obrera y de proponer medidas para mejorarlas. Se creó por Real Decreto de 5 de diciembre de 1883, bajo un gobierno liberal y por impulso de Segismundo Moret, y su célebre Información oral y escrita, publicada entre 1889 y 1893, constituye todavía una fuente de primer orden sobre las condiciones de la clase trabajadora española de la época. En 1903 fue sustituida por el Instituto de Reformas Sociales, órgano técnico que preparó buena parte de la legislación laboral y social de las décadas siguientes. El reformismo español de este período, como el alemán, tenía un propósito explícito de integración social y de prevención de la conflictividad obrera, en un contexto de expansión del anarquismo y del socialismo. Combinaba, por tanto, influencias del krausismo institucionista, del catolicismo social de la Rerum Novarum y del conservadurismo reformista.

En 1900 se promulga la Ley de Accidentes de Trabajo, que declara por primera vez la responsabilidad directa y objetiva del empresario en los accidentes sufridos por sus trabajadores, la llamada Ley Dato (1900). Su importancia jurídica es considerable, pues sustituye el principio de responsabilidad por culpa del Código Civil, que obligaba al trabajador accidentado a probar la negligencia del patrono, por la teoría del riesgo profesional.\footnote{%
    Según esta, el empresario responde de los accidentes porque el riesgo es inherente a la actividad de la que obtiene beneficio, con independencia de que haya cometido falta alguna. La ley no obligaba a asegurarse, pero permitía al empresario sustituir su responsabilidad por un seguro, lo que impulsó el desarrollo de las mutuas patronales.
} De esta ley proceden dos rasgos que el sistema español conserva hoy. El primero es la cotización por accidentes de trabajo y enfermedades profesionales a cargo exclusivo de la empresa, con tarifas diferenciadas según el riesgo de cada actividad. El segundo es la existencia de las mutuas colaboradoras, herederas de aquellas mutuas patronales.

En 1908 nace el Instituto Nacional de Previsión (INP), con el fin de planificar y difundir la previsión social y de gestionar los seguros sociales que irían surgiendo. Se creó por Ley en 1908, bajo el gobierno de Antonio Maura, y su principal impulsor intelectual fue el actuario José MALUQUER Y SALVADOR. Su principio inspirador fue el de la libertad subsidiada: la previsión debía ser voluntaria, y el Estado se limitaba a incentivarla con bonificaciones a quienes se afiliaran. Por tanto, frente al modelo bismarckiano de obligatoriedad, el reformismo español de 1908 apostó por el voluntarismo, por razones tanto ideológicas (el liberalismo y el catolicismo social desconfiaban de la coacción estatal) como prácticas (un Estado como la España de la época, con escasa capacidad administrativa y fiscal, no podía hacer efectiva una obligación general). No obstante, el resultado constituye otra confirmación a las predicciones de la economía conductual: la afiliación fue muy baja, concentrada en los trabajadores con mayor capacidad de ahorro, mientras que los más pobres, que más necesitaban la protección, quedaron fuera. 

El fracaso del voluntarismo condujo, a partir de 1919, a los primeros seguros obligatorios. El primero de estos es el Retiro Obrero Obligatorio (1919). Funcionaba por capitalización, con aportaciones del Estado y de los empresarios, y garantizaba una pensión muy modesta a los sesenta y cinco años a los asalariados con ingresos inferiores a un umbral. No obstante, la inflación, la Guerra Civil y la desorganización posterior destruyeron el valor de sus fondos.\footnote{
Un ejemplo claro, en la propia historia española, de como el riesgo inflacionista no diversificable constituye un riesgo que la teoría acertadamente identifica como una de las debilidades inherentes a los sistemas de capitalización
} El segundo, el de desempleo. La Caja Nacional contra el Paro Forzoso (1931), creada durante los primeros meses de la Segunda República, se organizaba como un sistema de subvención pública a cajas voluntarias de sindicatos y mutualidades, con una cobertura muy reducida. Por tanto, aunque no era un seguro obligatorio, sí que actuaba de manera análoga a éste. España no tuvo un verdadero seguro obligatorio de desempleo hasta la Ley 62/1961, que implantó el Seguro Nacional de Desempleo en el contexto del Plan de Estabilización de 1959, cuyas consecuencias sobre el empleo exigían algún tipo de protección. Existió un último, que inicialmente se concibió como subsidio por maternidad (1923), y que acabaría convirtiéndose en un seguro obligatorio de maternidad en 1929.

La Constitución republicana de 1931, en su art. 46, proclamó el compromiso de la República con los seguros de enfermedad, accidente, paro forzoso, vejez, invalidez y muerte, en la estela de la Constitución de Weimar. Hubo proyectos para unificar los seguros sociales en un sistema integrado, pero la inestabilidad política y la Guerra Civil impidieron desarrollarlos.

El franquismo mantuvo la arquitectura de seguros sociales, pero la reorientó ideológicamente y la hizo crecer por yuxtaposición. El Fuero del Trabajo de 1938 proclamó la previsión como objetivo del Estado nacionalsindicalista. En 1939, el Subsidio de Vejez sustituyó al Retiro Obrero y abandonó la capitalización por el reparto, una decisión que respondía tanto a la destrucción de los fondos del sistema anterior como a la necesidad de pagar pensiones de inmediato. Paralelamente surgió el Seguro Obligatorio de Enfermedad (1942), gestionado por el INP, que es el antecedente directo del sistema sanitario público español y de su red de ambulatorios y residencias sanitarias. Más tarde, en 1947, el Subsidio de Vejez se transformó en el Seguro Obligatorio de Vejez e Invalidez (SOVI), que todavía hoy genera pensiones residuales para quienes cotizaron antes de 1967 y no tienen derecho a pensión en el sistema actual (curioso testimonio de la continuidad institucional del sistema).

A partir de 1946, bajo el ministerio de Trabajo de José Antonio Girón de Velasco, se establecieron formas complementarias de protección social organizadas por ramas profesionales a través de las Mutualidades Laborales. Estas ofrecían prestaciones complementarias, a menudo más generosas que las de los seguros sociales, a los trabajadores de cada rama de actividad. 

Como nos podemos imaginar, el resultado de todo este entramado institucional, económico y social, fue un mosaico de seguros sociales y mutualidades independientes entre sí, con cotizaciones, prestaciones y gestores distintos, muy desigual entre sectores y con graves problemas de coordinación y de solvencia; que la doctrina iuslaboralista describiría más tarde como un sistema de yuxtaposición. En términos de ESPING-ANDERSEN, es la estratificación corporativa en estado puro: cada colectivo profesional tenía su propia protección, en función de su capacidad económica y de su peso en la estructura del régimen.

\paragraph{Segundo período: la unificación (1963-1978)}

En 1963 se promulga la Ley de Bases de la Seguridad Social. Su objetivo principal era implantar un modelo unitario e integrado de protección social: unificó e integró los distintos seguros sociales e hizo emerger un nuevo Sistema de Seguridad Social con base financiera de reparto, gestión pública y participación del Estado en la financiación. De hecho, su preámbulo declaraba expresamente el tránsito desde los seguros sociales hacia la Seguridad Social, entendida como una técnica de protección integral de la población frente a las situaciones de necesidad, y no como una suma de seguros para riesgos específicos. Fruto de la influencia, tardía y adaptada a un régimen autoritario, de BEVERIDGE y del Convenio 102 de la OIT. Además, el contexto económico también importa para entender la transición. El desarrollismo de los años sesenta, con altas tasas de crecimiento, emigración masiva del campo a la ciudad y creciente asalarización de la población, hacía posible y necesaria una protección más amplia. El régimen buscaba, además, homologarse con los países europeos occidentales en un momento de apertura económica.

\textcolor{red}{La Ley General de la Seguridad Social de 1966, materialización de los objetivos de la Ley de Bases, entró en vigor el 1 de enero de 1967. Pese al objetivo integrador del nuevo sistema, pervivieron multitud de organismos de protección social heredados del antiguo mutualismo, que adolecían de desequilibrios financieros. La causa de este último problema eran las llamadas bases tarifadas: se cotizaba no por el salario efectivamente percibido, sino por unas bases fijas establecidas para cada categoría profesional, muy inferiores a los salarios reales. Tenía dos consecuencias perversas. Por el lado de los ingresos, la recaudación era insuficiente. Por el lado de las prestaciones, estas quedaban desligadas del salario real y del esfuerzo contributivo, lo que debilitaba la lógica bismarckiana que el sistema pretendía seguir.

El documento menciona a continuación la Ley de Financiación y Perfeccionamiento de la Acción Protectora del Régimen General, que intentó corregir esos problemas financieros definiendo una estructura de cotización vinculada a los salarios reales en el Régimen General y ajustada a la naturaleza de las distintas ocupaciones en los regímenes especiales. Añade que, en la práctica, aumentó la acción protectora sin generar los recursos necesarios para financiarla, lo que agravó los desequilibrios. Corrección de fecha: el documento la sitúa en 1971, pero se trata de la Ley 24/1972, de 21 de junio. Le siguió un nuevo texto refundido de la Ley General de la Seguridad Social, aprobado por el Decreto 2065/1974. La observación del documento sobre el aumento de la protección sin financiación suficiente merece subrayarse, porque tiene una lectura de economía política: en los últimos años del franquismo, con una conflictividad laboral creciente y un régimen que buscaba legitimidad, ampliar las prestaciones sin financiarlas adecuadamente fue una forma de comprar paz social. La democracia heredó esa hipoteca, que condicionó las reformas de los años ochenta.}

\paragraph{Tercer período: democracia y constitucionalización (1978-1995)}
La transición a la democracia coincidió con la crisis económica internacional de los años setenta, que en España se manifestó con especial intensidad y retraso, y con la necesidad de construir un sistema de protección social propio de un Estado democrático. Los Pactos de la Moncloa (1977), incluyeron entre sus compromisos la reforma de la Seguridad Social.

Con la llegada de la democracia y la aprobación de la Constitución de 1978 se pone en marcha la configuración actual del sistema de la Seguridad Social. El art. $41$ establece que los poderes públicos mantendrán un régimen público de Seguridad Social para todos los ciudadanos que garantice la asistencia y prestaciones sociales suficientes ante situaciones de necesidad, especialmente en caso de desempleo, y que la asistencia y las prestaciones complementarias serán libres. El art. $43$ atribuye a los poderes públicos la organización y tutela de la salud pública, y el art. $50$ les exige garantizar, mediante pensiones adecuadas y periódicamente actualizadas, la suficiencia económica de los ciudadanos durante la tercera edad. 

La primera gran reforma de este período se produce con el Real Decreto-ley 36/1978, que materializa lo acordado en los Pactos de la Moncloa. Crea un sistema de participación institucional de los agentes sociales y un nuevo sistema de gestión llevado a cabo por cinco organismos: el Instituto Nacional de la Seguridad Social, para la gestión de las prestaciones económicas; el Instituto Nacional de la Salud, para las prestaciones sanitarias; el Instituto Nacional de Servicios Sociales, para la gestión de los servicios sociales; el Instituto Social de la Marina, para la gestión de los trabajadores del mar; y la Tesorería General de la Seguridad Social, como caja única del sistema. El mismo instrumento creó también el Instituto Nacional de Empleo (INEM), antecesor del actual SEPE.

Su significado de fondo es el paso de una gestión por riesgos, en la que cada seguro tenía su propio gestor, a una gestión por tipos de prestación, en la que cada instituto gestiona una clase de prestaciones para todos los riesgos. Por tanto, supone el desmantelamiento definitivo de la yuxtaposición heredada del franquismo. 

En la década de los ochenta se adoptan una serie de reformas encaminadas a extender las prestaciones a colectivos no cubiertos y a dotar de mayor estabilidad económica al sistema. La Ley 51/1980, Básica de Empleo, generaliza la prestación por desempleo y la vincula a los periodos de ocupación cotizados. La Ley 31/1984, de protección por desempleo, estructuró la protección en los dos niveles que perviven hoy: el contributivo, con prestaciones proporcionales al salario y a los periodos cotizados, y el asistencial, con subsidios para quienes han agotado la prestación contributiva o no tienen derecho a ella y cumplen determinados requisitos de rentas. Reforma que respondía al enorme aumento del paro en los primeros años ochenta ($>20\%$) y a la necesidad de proteger a los desempleados de larga duración. Estos años fueron también los de la reconversión industrial, en la que las prejubilaciones y las jubilaciones anticipadas se utilizaron masivamente para reducir las plantillas de los sectores en crisis (siderurgia, naval, minería): la versión española de la primera fase de PALIER.

La Ley 26/1985, de medidas urgentes para la racionalización de la estructura y de la acción protectora de la Seguridad Social, adoptó medidas relevantes en materia de pensiones. La equiparación paulatina de las bases de cotización con los salarios reales, la revalorización de las pensiones en función del IPC, la ampliación de los periodos necesarios para acceder a las prestaciones y para calcular las pensiones, y la simplificación de la estructura del sistema. En concreto, el periodo mínimo de cotización para acceder a la pensión de jubilación pasó de $10\to15$ años, y el periodo de cálculo de la base reguladora, de los $2\to 8$ últimos años. 

Destaca porque es la primera reforma restrictiva de las pensiones de la democracia y, como es de esperar, provocó la primera huelga general de la democracia española, convocada por Comisiones Obreras el 20 de junio de 1985 contra el gobierno socialista de González. UGT no se sumó a la convocatoria, pero la reforma abrió una fractura entre el gobierno y su sindicato hermano que se agravaría en los años siguientes. El episodio ilustra con precisión la tesis de PIERSON, un coste político inmediato y concentrado a cambio de un beneficio financiero difuso y futuro, y también la técnica de evitación de la culpa, porque la reforma incluía la garantía de revalorización con el IPC, que funcionaba como compensación.

Por otro lado, la Ley 14/1986 General de Sanidad, universaliza la sanidad pública al crear un Sistema Nacional de Salud que descansa en las Comunidades Autónomas. Con esta, la sanidad española pasa de la lógica bismarckiana, en la que la asistencia sanitaria era un derecho del trabajador asegurado y sus beneficiarios, a la lógica beveridgiana, en la que es un derecho del ciudadano financiado con impuestos. Sin embargo, la transición no se completaría hasta finales del siglo, con la plena financiación por impuestos y la universalización de la cobertura.

La Ley 37/1988, de Presupuestos Generales del Estado para 1989, reforma la estructura financiera de la Seguridad Social, separando las fuentes de financiación de la asistencia sanitaria, mayoritariamente mediante una aportación finalista del Estado, y de las pensiones, a cargo principalmente de las cotizaciones sociales. El primer paso de la separación de fuentes que el Pacto de Toledo consagraría como su primera recomendación en 1995. La financiación de la sanidad exclusivamente con impuestos generales se completó a finales de los noventa, y desde la culminación de los traspasos sanitarios en 2002 se integra en el sistema de financiación de las Comunidades Autónomas.

A primeros de los años noventa se aprueba la Ley 26/1990, por la que se establecen prestaciones no contributivas en la Seguridad Social. Esta ley extiende la protección a las personas que, no habiendo cotizado el tiempo suficiente, no tienen derecho a prestaciones contributivas; universaliza las prestaciones del sistema, y establece que las no contributivas se financiarán con aportaciones del Estado. En términos de FERRERA, es el primer paso para cerrar la gran laguna del modelo mediterráneo: la ausencia de una red asistencial estatal. Introduce en el sistema una capa beveridgiana, fundada en la necesidad y en la ciudadanía y no en la contribución. En paralelo, desde finales de los ochenta, las Comunidades Autónomas empezaron a crear sus propias rentas mínimas de inserción, con el País Vasco como pionero en 1989, en ejercicio de su competencia en materia de asistencia social. Así nació una red asistencial fragmentada territorialmente que el IMV intentaría articular treinta años después.

En 1994 se aprobó, mediante el Real Decreto Legislativo 1/1994, de 20 de junio, un nuevo texto refundido de la Ley General de la Seguridad Social, que estuvo vigente hasta su sustitución por el actual texto refundido de 2015.

Por último, y en paralelo a este entramado de cambios legislativos, comienzan a surgir dificultades financieras, derivadas tanto de la crisis económica como de cambios demográficos y estructurales. La caída del número de cotizantes durante 1992-1993, la universalización de las prestaciones y el envejecimiento de la población pusieron de relieve el desequilibrio financiero del sistema. En esos años circularon en el debate público español, con apoyo de algunos servicios de estudios de entidades financieras, informes que pronosticaban la quiebra del sistema de reparto y proponían su sustitución, total o parcial, por un sistema de capitalización inspirado en el modelo chileno y en el paradigma del Banco Mundial (1994). El contexto intelectual del Pacto de Toledo es, por tanto, el de la guerra de paradigmas.

\paragraph{Cuarto período: el Pacto de Toledo y la edad del consenso (1995-2010)}
La respuesta al desequilibrio financiero fue la aprobación del Pacto de Toledo (1995), un documento que describía la evolución de la Seguridad Social en España y los factores futuros que podrían incidir en su financiación, y que proponía una serie de recomendaciones. Nació de una proposición no de ley presentada por CiU en 1994, y su elaboración se prolongó durante más de un año en una ponencia de la Comisión de Presupuestos. Su nombre procede del Parador de Toledo, donde se celebraron algunas de las reuniones decisivas. 

El Pleno del Congreso de los Diputados aprobó el informe el 6 de abril de 1995, con quince recomendaciones. Entre ellas destacan la separación de las fuentes de financiación, con las prestaciones contributivas a cargo de las cotizaciones sociales y las no contributivas a cargo de la imposición general; la constitución de reservas a partir de los excedentes en momentos de bonanza; la simplificación e integración de los regímenes especiales; la mejora de las bases, con mayor correspondencia con los salarios reales; el refuerzo del carácter contributivo del sistema; y el mantenimiento del poder adquisitivo de las pensiones. El seguimiento de estas recomendaciones se realiza en el Congreso a través de la Comisión permanente de Seguimiento y Evaluación de los Acuerdos del Pacto de Toledo.

Por tanto, frente a quienes proponían la capitalización, el Pacto optó por reforzar el reparto mediante reformas paramétricas y una clarificación de su financiación, en la línea de lo que Palier describiría como la segunda fase de la reforma bismarckiana.

Las quince recomendaciones del Pacto de Toledo constituyeron la base del acuerdo alcanzado en octubre de 1996 entre el Gobierno y los dos sindicatos mayoritarios, que se plasmó en la Ley 24/1997, de 15 de julio, de Consolidación y Racionalización del Sistema de la Seguridad Social. Esta ley articuló la separación financiera del sistema según la naturaleza de las prestaciones, estableció un único tope de cotización, amplió de $8\to 15$ años el periodo de cálculo de la base reguladora de la pensión de jubilación, consagró la revalorización automática de las pensiones en función de la variación de los precios y dispuso la constitución de reservas con cargo a los excedentes de cotizaciones sociales. El acuerdo lo firmó el primer gobierno de José María Aznar con CCOO y UGT, sin la patronal, y es un ejemplo temprano de concertación social en materia de pensiones. Su lógica negociadora, como se ve, es de carácter transaccional: la reforma restrictiva, con $15$ años de cálculo en lugar de ocho y una escala de porcentajes más exigente, se compensó con garantías, como la revalorización automática por ley, y con la separación de fuentes, que liberaba a las cotizaciones de financiar prestaciones no contributivas. 

El Acuerdo para la Mejora y el Desarrollo del Sistema de Protección Social de abril de 2001, firmado por el gobierno, CCOO y las organizaciones empresariales, pero no por UGT, desarrolló otras recomendaciones del Pacto. De él derivó la Ley 35/2002, de medidas para el establecimiento de un sistema de jubilación gradual y flexible, que reguló la jubilación parcial como respuesta a la décima recomendación del Pacto de Toledo. La idea de fondo, que reaparecerá en todas las reformas posteriores hasta la de 2024, es que la transición entre la actividad y la jubilación no tiene por qué ser un salto brusco, y que permitir combinar trabajo y pensión puede alargar la vida laboral efectiva. Posteriormente se aprobó la Ley 28/2003, de 29 de septiembre, reguladora del Fondo de Reserva de la Seguridad Social (la hucha de las pensiones), que estableció que los excedentes de ingresos de carácter contributivo resultantes de la liquidación de los presupuestos de la Seguridad Social de cada ejercicio se aplicarían prioritaria y mayoritariamente al Fondo de Reserva. Se desarrolló reglamentariamente mediante el Real Decreto 337/2004. 

En octubre de 2003 se procedió a la renovación del Pacto de Toledo: el Congreso aprobó el informe elaborado por la comisión no permanente de valoración de los resultados obtenidos. Ese informe valoró el grado de cumplimiento de las recomendaciones iniciales, incluyó algunas recomendaciones adicionales en materia de nuevas formas de trabajo, mujer y protección social, dependencia, discapacidad e inmigración, y contextualizó el sistema de pensiones en el marco de la Unión. La inclusión de la inmigración entre las nuevas recomendaciones no era casual: en esos años España se convirtió en uno de los principales destinos migratorios del mundo, y la afiliación de trabajadores extranjeros empezó a transformar las cuentas del sistema.

Siguiendo las prioridades marcadas por la renovación del Pacto, se aprobó la Ley 40/2007, de 4 de diciembre, de medidas en materia de Seguridad Social, con efectos sobre cuatro contingencias: incapacidad temporal, incapacidad permanente, jubilación y supervivencia. En materia de jubilación parcial, supeditó el acceso a $61$ años de edad, una antigüedad mínima de $6$ años en la empresa y un periodo de cotización de al menos $30$ años. En materia de viudedad, extendió la protección a las parejas de hecho. La ley desarrolló el acuerdo social de julio de 2006 entre el gobierno de José Luis Rodríguez Zapatero, los sindicatos y la patronal, en otro ejemplo de reforma pactada. En estos mismos años se dieron pasos en la integración de los regímenes especiales que el Pacto de Toledo recomendaba. Los trabajadores por cuenta propia agrarios se integraron en el régimen de autónomos en 2008. En 2012, los trabajadores por cuenta ajena agrarios y los empleados del hogar se incorporaron al Régimen General mediante sistemas especiales con reglas propias. 

En los otros ámbitos de la protección social, destacan tres normas de este período. La primera es la Ley 56/2003, de 16 de diciembre, de Empleo, que reformó la Ley de 1980 y estableció el marco de las políticas de empleo teniendo en cuenta la distribución constitucional de competencias entre el Estado y las Comunidades Autónomas, configurando el Sistema Nacional de Empleo. La segunda es la Ley 16/2003, de 28 de mayo, de cohesión y calidad del Sistema Nacional de Salud, que definió los ámbitos de actuación del Estado y de las Comunidades Autónomas en materia sanitaria tras completarse el proceso de transferencia de competencias en 2001.\footnote{%
    El primer traspaso sanitario fue a Cataluña en 1981, y los últimos, aprobados en diciembre de 2001 para las diez comunidades que aún no tenían la competencia, se hicieron efectivos el 1 de enero de 2002.
} La tercera es la Ley 39/2006, de 14 de diciembre, de Promoción de la Autonomía Personal y Atención a las personas en situación de dependencia, que creó el Sistema para la Autonomía y Atención a la Dependencia e incorporó la dependencia al sistema de protección social español.\footnote{%
    El primer gran paso de desfamiliarización de los cuidados en España. No obstante, la ley se aprobó en el cenit del ciclo expansivo, con superávit en las cuentas públicas, y se diseñó con un calendario de implantación gradual que la crisis de 2008 alteraría profundamente, dato que explica buena parte de su accidentado desarrollo posterior.
}

El balance financiero de este período fue excepcionalmente favorable. La llegada de varios millones de trabajadores inmigrantes, con un proceso extraordinario de regularización en 2005, y el crecimiento del empleo asociado al boom inmobiliario elevaron la afiliación a máximos históricos, mientras el sistema acumulaba superávits que nutrieron el Fondo de Reserva. Sin embargo, la importancia del Pacto de Toledo va más allá de su contenido. Es, ante todo, una innovación institucional: saca las pensiones de la competición partidista ordinaria y las somete a un consenso parlamentario amplio, con participación indirecta de los agentes sociales a través de los acuerdos que desarrollan sus recomendaciones. En términos de PIERSON y WEAVER, es un mecanismo de reparto de la culpa: si todos los grandes partidos firman las recomendaciones, ninguno puede utilizar electoralmente contra los demás las reformas que de ellas se derivan, y el coste político de las medidas impopulares se diluye. El Pacto ha funcionado razonablemente bien mientras el consenso se ha mantenido, y ha perdido eficacia cuando se ha roto, como en 2013. 


No obstante, la ilusión de un problema resuelto se desvaneció en muy poco tiempo.

\paragraph{Quinto período: crisis y ajuste (2010-2013)}
La crisis financiera de 2008 y la posterior crisis de deuda soberana transformaron radicalmente la situación. La destrucción de empleo, que llevó la tasa de paro por encima del $25\%$ en 2013, redujo drásticamente los ingresos por cotizaciones, mientras el gasto en pensiones seguía creciendo por la inercia demográfica y por el efecto sustitución. El sistema pasó del superávit al déficit en pocos años.

En mayo de 2010, en plena crisis de la deuda griega y bajo una presión extraordinaria de los mercados y de las instituciones europeas, el gobierno de Rodríguez Zapatero aprobó un plan de ajuste que incluía, mediante el Real Decreto-ley 8/2010, la congelación de las pensiones contributivas para 2011, con excepción de las mínimas y las no contributivas: la primera vez desde 1985 que se suspendía la revalorización de las pensiones, y la medida rompió la garantía que había sido la compensación de las reformas restrictivas anteriores.

El Acuerdo Social y Económico de 2 de febrero de 2011, firmado por el gobierno, CCOO, UGT, CEOE y CEPYME, y el Informe de Evaluación y Reforma del Pacto de Toledo aprobado por el Congreso en enero de 2011, sentaron las bases de la Ley 27/2011, sobre actualización, adecuación y modernización del sistema de Seguridad Social. Fue una reforma paramétrica pactada, aprobada por un gobierno socialista en plena crisis y bajo la presión de las instituciones europeas y de los mercados, con largos periodos transitorios que se extendían hasta 2027. Siguió, por tanto, la lógica de la evitación de la culpa y del consenso. Su contenido concreto se expone como regla vigente más adelante.

En el ámbito sanitario, el Real Decreto-ley 16/2012, de medidas urgentes para garantizar la sostenibilidad del Sistema Nacional de Salud y mejorar la calidad y seguridad de sus prestaciones, definió las prestaciones del sistema y su financiación e introdujo importantes novedades en el ámbito de las prestaciones farmacéuticas. Aprobado por el gobierno de Mariano Rajoy, también restringió el acceso a la asistencia sanitaria de los inmigrantes en situación irregular, lo que suponía un paso atrás en la universalización; restricción que fue revertida por el Real Decreto-ley 7/2018.

El Real Decreto-ley 28/2012, dejó sin efecto la actualización de las pensiones correspondiente a la desviación de la inflación de 2012 respecto de la prevista, que la legislación entonces vigente garantizaba; y que el Tribunal Constitucional (STC 49/2015) casó, con una doctrina de gran trascendencia: la actualización por la desviación del IPC no constituía un derecho adquirido de los pensionistas, sino una mera expectativa, que el legislador podía modificar antes de que se consolidara. Además, se endurecieron las condiciones de acceso a la jubilación anticipada y parcial, con el objetivo de elevar la edad efectiva de jubilación.

El Informe del Comité de Expertos sobre el factor de sostenibilidad del sistema público de pensiones, de 2013, sirvió de base a la Ley 23/2013, reguladora del Factor de Sostenibilidad y del Índice de Revalorización del Sistema de Pensiones de la Seguridad Social. Su informe proponía dos mecanismos de ajuste automático. El primero, un factor de equidad intergeneracional, ligaba la pensión inicial a la esperanza de vida. El segundo, un factor de revalorización anual, ligaba la revalorización de las pensiones al equilibrio financiero del sistema. La ley fue aprobada por el Partido Popular con su mayoría absoluta y sin el consenso del Pacto de Toledo, con el voto en contra del resto de grupos, en un contexto de recomendaciones europeas que instaban a España a introducir un factor de sostenibilidad y de un déficit creciente del sistema. 

En esencia, representó la ruptura del consenso y la clave para entender su posterior destino: una reforma impuesta por una mayoría absoluta carece de protección frente a un cambio de mayoría, y los partidos que votaron en contra quedaron libres para prometer su derogación; ejemplo de manual de inconsistencia temporal.

\paragraph{Sexto período: el ajuste silencioso y la reversión (2014-2023)}
Entre 2014 y 2017, el índice de revalorización introducido en 2013 se aplicó en su valor mínimo, el $0,25\%$ anual, porque el déficit del sistema impedía cualquier revalorización superior. Mientras la inflación se mantuvo próxima a cero o negativa, como ocurrió en 2014-2016, la pérdida de poder adquisitivo fue escasa y la medida pasó relativamente desapercibida (ajuste silencioso que la tesis de PIERSON predice como estrategia de recorte más viable). Al mismo tiempo, el déficit del sistema se financió primero con el Fondo de Reserva, que se agotó casi por completo, y desde 2017 con préstamos del Estado a la Tesorería General de la Seguridad Social, un mecanismo cuya evolución se examinará después.

El silencio se rompió cuando la inflación volvió. Desde enero de 2018, con una inflación que superaba ampliamente el $0,25\%$ de revalorización, los pensionistas se movilizaron de forma masiva y sostenida en toda España. Las concentraciones de los lunes en Bilbao se convirtieron en su emblema, y las manifestaciones de febrero y marzo de 2018 reunieron a cientos de miles de personas. Los PGE para 2018 (Ley 6/2018), aprobados gracias a un acuerdo del gobierno de Rajoy con el PNV, revalorizaron las pensiones un $1,6\%$, más en el caso de las mínimas, y aplazaron la entrada en vigor del factor de sostenibilidad hasta 2023 como máximo: la primera vez en España que un colectivo de beneficiarios consigue suspender un mecanismo automático de ajuste. 

Dos años después del cambio de Gobierno, en noviembre de 2020, el Congreso aprobó la renovación del Pacto de Toledo, con veinte recomendaciones, entre ellas la vuelta a la revalorización de las pensiones con el IPC, la financiación con impuestos de los llamados gastos impropios de la Seguridad Social y el fomento de la prolongación voluntaria de la vida laboral. Además, en plena pandemia, la protección social española dio un paso largamente pendiente. El Real Decreto-ley 20/2020, creó el IMV, tramitado después como Ley 19/2021. Con el IMV, España cerró la última gran laguna del modelo mediterráneo que FERRERA había identificado en 1996: la ausencia de una red de renta mínima estatal, establecida como derecho subjetivo a quienes se encuentran en situación de vulnerabilidad económica, y que complementa los programas de rentas mínimas de las Comunidades Autónomas. En paralelo, el Real Decreto-ley 3/2021 sustituyó el anterior complemento de maternidad por un complemento para la reducción de la brecha de género, tras la sentencia del TJUE, asunto C-450/18, que consideró discriminatorio para los hombres el complemento anterior.

Como última gran reforma, las dos disposiciones materializaron los cambios anteriores. La Ley 21/2021, de garantía del poder adquisitivo de las pensiones y de otras medidas de refuerzo de la sostenibilidad financiera y social del sistema público de pensiones. La segunda, el Real Decreto-ley 2/2023, de medidas urgentes para la ampliación de derechos de los pensionistas, la reducción de la brecha de género y el establecimiento de un nuevo marco de sostenibilidad del sistema público de pensiones. Sin embargo, las dos piezas de la reforma tuvieron un grado de consenso distinto, y esa diferencia es relevante. 

La Ley 21/2021 nació del acuerdo social entre el gobierno de Pedro Sánchez, con ESCRIVÁ al frente del Ministerio de Inclusión, Seguridad Social y Migraciones, los sindicatos y las organizaciones empresariales. El Real Decreto-ley 2/2023 se acordó solo con los sindicatos, porque la patronal rechazó una reforma que descansaba fundamentalmente en el aumento de las cotizaciones. Siendo, además, la primera gran reforma de las pensiones españolas que se apoya casi exclusivamente en el lado de los ingresos y no en el de los gastos, lo que la distingue de todas las anteriores y la sitúa a contracorriente de la mayoría de las reformas europeas del período. 

Completan este período otras tres normas. La Ley 12/2022, de 30 de junio, de regulación para el impulso de los planes de pensiones de empleo, que intentó desarrollar el segundo pilar, prácticamente inexistente en España. El Real Decreto-ley 13/2022, de 26 de julio, que implantó un nuevo sistema de cotización de los trabajadores autónomos por sus rendimientos netos reales, con una transición gradual hasta 2032. Y la Ley 3/2023, de 28 de febrero, de Empleo, que derogó el texto refundido de la Ley de Empleo de 2015 y previó la transformación del SEPE en la Agencia Española de Empleo, aún no materializada.

\paragraph{Séptimo período: la actualidad (2024-2026)}
Los últimos años se caracterizan por la continuación de las reformas en los márgenes del sistema y por una creciente fragilidad parlamentaria. 

El Real Decreto-ley 2/2024, reformó en profundidad el nivel asistencial de la protección por desempleo. Reorganizó los subsidios, amplió su cobertura a colectivos antes excluidos y creó el complemento de apoyo al empleo, que permite compatibilizar el subsidio con un salario durante un tiempo para evitar la trampa de la inactividad, un problema cuyo antecedente remoto es la \textit{less eligibility} de 1834. El Real Decreto-ley 11/2024, fruto de un acuerdo social alcanzado en el verano de 2024 con sindicatos y patronal, que rediseñó la jubilación activa, la demorada y la parcial, con el objetivo de acercar la edad efectiva de jubilación a la legal y de facilitar una transición gradual entre el trabajo y la jubilación, retomando la idea de la Ley 35/2002 y llevandola más lejos.

Más adelante, se promulgó el Real Decreto-ley 16/2025, que incluía en un único texto la revalorización de las pensiones para 2026, un $2,7\%$ con carácter general, junto con la prórroga de medidas del llamado escudo social y diversas medidas tributarias; por su heterogeneidad, un decreto ómnibus. Sin embargo, el Congreso de los Diputados no lo convalidó, de modo que la revalorización de las pensiones quedó momentáneamente sin base legal. Más adelante, el Gobierno separó materias, aprobando: (i) el Real Decreto-ley 3/2026, dedicado exclusivamente a la revalorización de las pensiones y a otras medidas urgentes de Seguridad Social, que el Congreso convalidó el 26 de febrero de 2026; y, (ii) el decreto que recogía las medidas del escudo social, que volvió a caer.

En septiembre de 2026, por último, el Congreso aprobó definitivamente la reforma de las leyes de dependencia y de discapacidad, que reconoce los sistemas públicos de servicios sociales como servicios esenciales, refuerza la asistencia personal y la atención centrada en la persona, y modifica los mecanismos de financiación estatal del Sistema para la Autonomía y Atención a la Dependencia. 

\paragraph{El hilo conductor: cuatro constantes de la historia española}

Leída con las herramientas teóricas y con el trasfondo internacional, la historia de la protección social española presenta cuatro constantes que conviene retener, porque estructuran el resto del tema.

La primera es la continuidad del reparto. Desde sus inicios, España, como la inmensa mayoría de los países de la Europa continental, ha financiado sus pensiones mediante un sistema de reparto, en el que las cotizaciones de los trabajadores en activo financian directamente las pensiones de los jubilados, sin acumulación intermedia. Salvo la excepción parcial del Fondo de Reserva, no existe una cuenta individual acumulada a nombre de cada cotizante. Esta continuidad no es un accidente. La destrucción de los fondos del Retiro Obrero en los años treinta y cuarenta vacunó al sistema contra la capitalización, y la guerra de paradigmas de los noventa se resolvió en España, con el Pacto de Toledo, a favor del reparto reformado. Las razones económicas de esta elección, y sus implicaciones, se examinan más adelante.

La segunda es el tránsito de la estratificación a la universalidad. El sistema ha pasado de un mosaico de seguros y mutualidades por colectivos profesionales a un sistema unificado, con la integración progresiva de los regímenes especiales, la universalización de la sanidad, la creación de las pensiones no contributivas y, por último, el IMV. Es un movimiento desde la lógica bismarckiana pura hacia un híbrido con una capa beveridgiana cada vez más amplia, aunque el núcleo del sistema de pensiones sigue siendo contributivo.

La tercera es la alternancia entre reformas pactadas y reformas impuestas. Las reformas aprobadas mediante el consenso del Pacto de Toledo y la concertación social (1997, 2002, 2007, 2011, 2021) han perdurado. La reforma aprobada por una mayoría absoluta sin consenso (2013) se suspendió a los cinco años y se derogó a los ocho. Por tanto, la calidad técnica de una reforma no garantiza su supervivencia; su legitimidad política, en cambio, sí contribuye a ella.

La cuarta es la creciente dependencia de los ingresos frente al ajuste de los gastos. Hasta 2013, todas las grandes reformas españolas actuaron sobre el gasto, con más años de cálculo, más años de cotización, más edad y mecanismos de ajuste automático. Desde 2021, la reforma se ha apoyado fundamentalmente en el aumento de las cotizaciones y de las transferencias del Estado. Si esta estrategia es suficiente para garantizar la sostenibilidad del sistema es la gran pregunta que deberemos responder. Si es la más adecuada desde el punto de vista de la eficiencia y de la equidad entre generaciones es una cuestión que deberemos abordar desde la teoría de la incidencia de las cotizaciones y desde la contabilidad generacional.

\subsection{Marco jurídico vigente: cómo se regula hoy el sistema y por qué}

Aunque parte del contenido se solape con el apartado histórico, aquí se tratará la regulación actual en sentido estricto. No interesa por qué se hicieron las reformas ni en qué contexto, que ya se ha expuesto, sino cómo se regula hoy el sistema, por qué y con qué objetivo. El apartado responde a esa exigencia en cuatro niveles, ordenados de lo más general a lo más concreto y del plano interno al supranacional. El primero es la \textbf{Constitución}, que fija los mandatos, los límites y el reparto de competencias. El segundo es la \textbf{legislación} ordinaria, que concreta esos mandatos en un sistema operativo. El tercero son los \textbf{principios} que dan coherencia al conjunto y orientan su interpretación. El cuarto es el \textbf{marco supranacional}, internacional y europeo, que condiciona cada vez más la libertad del legislador español.

\textbf{NOTA.-} La Seguridad Social es una institución en la que el derecho y la economía están entrelazados de forma inusualmente estrecha. Los principios jurídicos, como la suficiencia, la solidaridad o la contributividad, tienen una traducción económica inmediata en parámetros de gasto y de ingreso. A la inversa, las restricciones económicas, como la sostenibilidad financiera o la estabilidad presupuestaria, se han ido incorporando al propio ordenamiento como normas jurídicas. Buena parte de los conflictos que el tribunal espera ver analizados surgen precisamente en esa frontera: qué margen deja la Constitución al legislador para recortar las pensiones, si la sostenibilidad es ya un principio jurídico que compite con la suficiencia, o hasta dónde puede llegar la descentralización sin romper la caja única.

\subsubsection{La Constitución económica de la Seguridad Social}

\paragraph{El art. 41: un régimen público, para todos, con prestaciones suficientes}

El precepto central es el art. $41$ de la Constitución, cuyo tenor conviene tener presente literalmente: \textit{los poderes públicos mantendrán un régimen público de Seguridad Social para todos los ciudadanos, que garantice la asistencia y prestaciones sociales suficientes ante situaciones de necesidad, especialmente en caso de desempleo; la asistencia y prestaciones complementarias serán libres}. La doctrina iuslaboralista ha extraído de este enunciado, aparentemente sencillo, cinco mandatos distintos, cada uno con consecuencias propias:
\begin{la}
    \item \textbf{Público}. El primero es el carácter público del régimen. La Constitución no se limita a exigir que los ciudadanos estén protegidos, lo que podría lograrse con un seguro privado obligatorio, sino que exige que exista un régimen público de Seguridad Social. En los términos de la distinción de BARR y DIAMOND, el constituyente no solo impuso el mandato, sino también un papel público en la provisión o, como mínimo, en la organización y la garantía. La doctrina mayoritaria entiende que este mandato impide sustituir el sistema público por uno privado, aunque no impide que existan pilares privados complementarios, que el propio art. declara libres. Como podemos intuir, es un límite constitucional relevante para el debate sobre la capitalización: una transición completa al modelo de capitalización sería inconstitucional en España; una reforma mixta con un pilar capitalizado obligatorio junto al público sería, en cambio, discutible, pero no necesariamente contraria a la Constitución;
    \item \textbf{Universal}. El segundo mandato es la universalidad subjetiva: el régimen es para todos los ciudadanos. La expresión rompe con la lógica bismarckiana pura, en la que el sujeto protegido es el trabajador, y acoge la beveridgiana de la ciudadanía, en la estela de la ciudadanía social de MARSHALL. Pero la doctrina constitucional no ha interpretado este mandato como una exigencia de que todas las prestaciones sean universales, sino como la exigencia de que ningún ciudadano quede fuera de toda protección. Un sistema mixto, con un nivel contributivo para quienes trabajan y un nivel no contributivo para quienes no alcanzan la protección contributiva, cumple el mandato. Por eso la Ley 26/1990 de pensiones no contributivas y la Ley 19/2021 del IMV pueden leerse como el cumplimiento progresivo, y tardío, de la universalidad del art. $41$.
    \item \textbf{Suficiencia}. El tercero es la \textbf{suficiencia}. Las prestaciones deben ser suficientes, cuyo fundamento remoto es la categoría de bien preferente de MUSGRAVE. Evidentemente, es un concepto jurídico indeterminado, que la Constitución no define y que el legislador ha concretado de formas diversas a lo largo del tiempo. Además, como se verá, es también el principio que está en el centro del debate sobre la sostenibilidad. 
    \item \textbf{Necesidad}. El cuarto es la orientación a la situación de necesidad, con mención expresa del desempleo. La expresión situaciones de necesidad ha sido objeto de un debate doctrinal de largo alcance. Una lectura sostiene que el art. $41$ desplaza el fundamento del sistema desde la contingencia (el riesgo asegurado, que es la lógica del seguro) hacia la necesidad (la situación efectiva de carencia, que es la lógica de la asistencia). Según esta lectura, la Constitución habría asumido un modelo de Seguridad Social más asistencial que contributivo. La lectura opuesta, que ha prevalecido en la práctica legislativa y en la jurisprudencia constitucional, entiende que la necesidad puede presumirse cuando se produce la contingencia, de modo que un sistema contributivo que protege frente a la jubilación, la invalidez o el desempleo satisface el mandato aunque no compruebe la necesidad efectiva de cada beneficiario. Detrás de este debate técnico late la tensión entre la hucha y Robin Hood: si la Seguridad Social debe atender sobre todo a la necesidad, las pensiones contributivas elevadas de quienes tienen otros recursos serían difícilmente justificables; si debe atender a la contingencia, lo que resulta difícil de justificar es la comprobación de recursos.
    \item \textbf{Complementario}. El quinto es la libertad del nivel \textbf{complementario}. El documento señala que existe un nivel complementario de protección social que es libre, es decir, que puede ser de carácter privado. 
\end{la}

El art. $41$ cierra así un modelo de tres niveles: el contributivo y el no contributivo, ambos públicos, y el complementario, libre. Este último incluye las mejoras voluntarias de las prestaciones públicas, los planes y fondos de pensiones de empleo e individuales, y los seguros privados. Su ordenación corresponde a la legislación de planes y fondos de pensiones y a la de seguros, no a la de Seguridad Social, y su situación en España, marcada por un escasísimo desarrollo, se analiza más adelante.

Sin embargo, éste no agota la Constitución de la protección social. El art. $50$ obliga a los poderes públicos a garantizar, mediante pensiones adecuadas y periódicamente actualizadas, la suficiencia económica de los ciudadanos durante la tercera edad. También les obliga a promover su bienestar mediante un sistema de servicios sociales que atienda sus problemas específicos de salud, vivienda, cultura y ocio, con independencia de las obligaciones familiares. La cláusula periódicamente actualizadas ha sido una de las más debatidas del ordenamiento, porque de ella depende el alcance de la garantía de revalorización. El Tribunal Constitucional ha entendido que exige una actualización periódica, pero no una actualización necesariamente igual a la inflación, ni un mecanismo determinado. Es el legislador quien elige el método, y puede hacerlo de forma que las pensiones pierdan poder adquisitivo en determinadas circunstancias, siempre que no vacíe de contenido la garantía. Esta interpretación es la que permitió la constitucionalidad del índice de revalorización de 2013, con su suelo del $0,25\%$, y la que hace que la actual indexación al IPC sea una opción del legislador y no una exigencia constitucional. El legislador puede cambiarla por ley, sin reforma constitucional, como ya ocurrió en 2013 y en sentido inverso en 2021.

Por su parte, el art. $43$ reconoce el derecho a la protección de la salud y atribuye a los poderes públicos la organización y tutela de la salud pública a través de medidas preventivas y de las prestaciones y servicios necesarios. Su separación del art. $41$ es significativa: el constituyente concibió la protección de la salud como un derecho autónomo y no como una prestación de la Seguridad Social. Así preparó la transición del modelo bismarckiano de asistencia sanitaria al modelo beveridgiano del Sistema Nacional de Salud que la Ley General de Sanidad consumaría en 1986.

El art. $49$, reformado el 15 de febrero de 2024 en la primera reforma constitucional de contenido social desde 1978, sustituyó la expresión \textit{disminuidos físicos, sensoriales y psíquicos} por la de \textit{personas con discapacidad}. Reformuló además el precepto en clave de derechos, al imponer a los poderes públicos la promoción de su autonomía personal e inclusión social y la atención particular a las necesidades de las mujeres y los menores con discapacidad. Como sabemos, la reforma tiene relevancia directa para el Sistema para la Autonomía y Atención a la Dependencia y para las pensiones de incapacidad, y es el antecedente constitucional inmediato de la reforma legal de dependencia y discapacidad aprobada en septiembre de 2026.

El art. $39$ obliga a los poderes públicos a asegurar la protección social, económica y jurídica de la familia, fundamento de las prestaciones familiares y de nacimiento y cuidado de menor. Paralelamente, el art. $129.1$ dispone que la ley establecerá las formas de participación de los interesados en la Seguridad Social y en la actividad de los organismos públicos cuya función afecte directamente a la calidad de la vida o al bienestar general. Es el fundamento constitucional del principio de participación de los agentes sociales en la gestión del sistema, que examinaremos más adelante.

Además, los art.s enunciados están situados en el capítulo III del título I de la Constitución, bajo la rúbrica de principios rectores de la política social y económica, lo que tiene consecuencias jurídicas decisivas. 

En primer lugar, según el art. $53.3$, el reconocimiento, el respeto y la protección de estos principios informarán la legislación positiva, la práctica judicial y la actuación de los poderes públicos. Sin embargo, solo podrán ser alegados ante la jurisdicción ordinaria de acuerdo con lo que dispongan las leyes que los desarrollen. En otras palabras, los principios rectores no son derechos subjetivos directamente exigibles ante los tribunales: un ciudadano no puede reclamar una pensión amparándose solo en el art. $41$, sino en la Ley que lo desarrolla. El Tribunal Constitucional ha construido sobre esta base una doctrina que concede al legislador un amplio margen de configuración del sistema. En la STC 65/1987 afirmó que la Seguridad Social es una función del Estado y no un contrato de seguro entre el Estado y el asegurado. Por tanto, el derecho a las prestaciones nace de la Ley y no de las cotizaciones, y la ley puede modificarlo. En la STC 134/1987, sobre la fijación de un tope máximo a las pensiones, rechazó que la limitación de las pensiones más altas vulnerara la Constitución. Razonó que el legislador puede atender a las circunstancias económicas y a la necesidad de redistribuir recursos escasos, y que el art. $41$ no garantiza una correspondencia estricta entre cotizaciones y prestaciones. En la STC 37/1994 recurrió a la categoría de la garantía institucional: el art. $41$ garantiza la existencia de un régimen público de Seguridad Social reconocible como tal por la conciencia social del momento, y el legislador puede configurarlo libremente mientras no lo desnaturalice.\footnote{
La categoría de garantía institucional, procedente de la doctrina alemana de Weimar y en particular de SCHMITT, protege a una institución frente a su supresión o vaciamiento, no frente a su reforma. Aplicada a la Seguridad Social, significa que el legislador no puede suprimir el sistema público ni reducirlo hasta hacerlo irreconocible, pero puede alterar sus prestaciones, sus requisitos y su financiación con gran libertad. La STC 49/2015, de 5 de marzo, aplicó esta doctrina al Real Decreto-ley 28/2012, que había dejado sin efecto la actualización de las pensiones por la desviación del IPC de 2012. Concluyendo que esa actualización no era un derecho adquirido, sino una mera expectativa que solo se habría consolidado al final del ejercicio. Por tanto, su supresión antes de ese momento no constituía una expropiación de derechos ni una retroactividad prohibida por el art. $9.3$.

Una de las implicaciones claras de ésto, desde el punto de vista teórico, es que la pensión pública española no es un derecho de propiedad sobre las cotizaciones pagadas, sino una expectativa legal sometida a la voluntad del legislador. Este rasgo tiene dos caras. Por un lado, dota al sistema de la flexibilidad necesaria para repartir entre generaciones los shocks agregados, que es precisamente la ventaja que GORDON y VARIAN atribuían al reparto público. Por otro, expone a los cotizantes al riesgo político, es decir, al riesgo de que el legislador futuro cambie las reglas. Contrapartida exacta del riesgo de mercado que soportan los sistemas de capitalización individual, y argumento que utilizan los defensores de estos últimos. La literatura lo formula como una elección entre riesgo político y riesgo financiero. Ningún sistema elimina el riesgo; solo decide quién lo soporta y de qué tipo es.
}

Por otro lado, tiene gran relevancia práctica. El art. $86$ prohíbe que los decretos-leyes afecten a los derechos, deberes y libertades de los ciudadanos regulados en el Título I. El Tribunal Constitucional ha interpretado esta limitación de forma restrictiva, de modo que solo impide los decretos-leyes que regulen el régimen general o los elementos esenciales de esos derechos. Como el derecho a las prestaciones de Seguridad Social está en el capítulo III y es de configuración legal, el Gobierno ha podido utilizar con gran frecuencia el decreto-ley en materia de pensiones. Así se aprobaron la congelación de 2010, la supresión de la actualización de 2012, la segunda parte de la reforma de 2023, la reforma de la jubilación flexible de 2024 y, cada año, la revalorización de las pensiones. El episodio de enero y febrero de 2026, en el que el Congreso no convalidó el decreto-ley de revalorización por contener otras materias, pone de manifiesto la otra cara de esta práctica. Cuando no hay presupuestos aprobados, y en 2026 siguen prorrogados los de 2023, la revalorización depende de un decreto-ley que el Congreso puede no convalidar. Una garantía que el legislador ordinario ha establecido como automática queda así sometida en la práctica a una votación anual.

\paragraph{Cotización: una prestación patrimonial de carácter público}
La Constitución no menciona las cotizaciones sociales, pero su naturaleza jurídica tiene importancia tanto teórica como práctica. El art. $31.3$ establece que solo podrán establecerse prestaciones personales o patrimoniales de carácter público con arreglo a la ley. La STC 185/1995 definió estas prestaciones como aquellas impuestas coactivamente, es decir, establecidas unilateralmente por el poder público sin el concurso de la voluntad del obligado, y que tienen una finalidad de interés público. 

Las cotizaciones sociales cumplen ambos requisitos: son obligatorias y financian una función pública. Por eso la doctrina y la jurisprudencia las consideran prestaciones patrimoniales de carácter público, sujetas a la reserva de ley del art. $31.3$, aunque no sean tributos en sentido estricto ni se rijan por la Ley General Tributaria.

Ahora bien, esta calificación jurídica abre un debate económico que se desarrolla más adelante: si la cotización es una contraprestación por un seguro, y por tanto un precio, o un impuesto sobre el trabajo. La respuesta que demos tendrá consecuencias sobre su incidencia y sus efectos en el empleo.\footnote{%
    SUMMERS (1989) mostró que, en la medida en que los trabajadores perciban un vínculo entre lo que cotizan y lo que recibirán, la cotización no es un impuesto puro, sino en parte un precio. Su coste en términos de empleo es menor que el de un impuesto de igual cuantía, porque los trabajadores aceptan salarios netos menores a cambio de mayores prestaciones futuras. Las reformas que debilitan ese vínculo, como la subida de la base máxima sin aumento proporcional de la pensión máxima, la cuota de solidaridad o la cotización del MEI, que no genera derechos individuales, acercan la cotización a un impuesto. En consecuencia, aumentan su coste en términos de eficiencia.
} 

\paragraph{Reparto de competencias: el art. 149.1.17ª y la caja única}

El art. 149.1.17ª reserva al Estado la competencia exclusiva sobre la legislación básica y el régimen económico de la Seguridad Social, sin perjuicio de la ejecución de sus servicios por las Comunidades Autónomas. El precepto distingue, por tanto, tres ámbitos: la legislación básica, que es estatal; el régimen económico, que es íntegramente estatal, sin base ni desarrollo autonómico; y la ejecución de los servicios, que puede ser autonómica.

La expresión régimen económico es la clave del sistema. El Tribunal Constitucional (STC 124/1989) la ha interpretado como la garantía de la unidad patrimonial y financiera del sistema, es decir, de la caja única. Todos los recursos de la Seguridad Social, sean cotizaciones o transferencias del Estado, se integran en un patrimonio único, y todas las prestaciones se pagan con cargo a él, con independencia del territorio en que se recauden las cotizaciones o se paguen las prestaciones. Las Comunidades Autónomas pueden asumir la gestión de servicios pero no pueden tener su propia caja, ni fijar sus propias cotizaciones, ni pagar las pensiones contributivas con recursos propios del sistema.

¿Por qué esta excepción, en un Estado tan descentralizado como el español? La razón es económica y responde a la teoría del federalismo fiscal. Las pensiones de reparto son el instrumento de solidaridad interterritorial más potente, porque la estructura demográfica y laboral de las Comunidades Autónomas es muy desigual. Una comunidad envejecida, con muchos pensionistas y pocos cotizantes, como Asturias, Galicia o Castilla y León, genera un déficit territorial del sistema. Una comunidad joven y con alto empleo, como Madrid o Baleares, genera un superávit. Si cada comunidad tuviera su propia caja, las envejecidas tendrían que elevar mucho sus cotizaciones o reducir sus pensiones, y las jóvenes podrían ofrecer mejores condiciones con menores costes. Por tanto, la caja única es un mecanismo de mancomunación de riesgos entre territorios, que diversifica el riesgo demográfico regional del mismo modo que el seguro diversifica el riesgo individual.\footnote{%
    La teoría del federalismo fiscal (MUSGRAVE, OATES) asigna la función redistributiva al nivel central de gobierno, precisamente porque la movilidad de personas y empresas hace que la redistribución descentralizada sea inestable: los contribuyentes netos se trasladarían a los territorios con menor redistribución, y los beneficiarios netos, a los de mayor redistribución.
}

Ahora bien, la competencia de ejecución ha dado lugar a una jurisprudencia constitucional abundante sobre sus límites. Por ejemplo, la STC 239/2002,avaló que la Comunidad Autónoma de Andalucía concediera complementos a las pensiones no contributivas con cargo a su propio presupuesto. Entendió que esos complementos no forman parte del régimen económico de la Seguridad Social, sino que se amparan en la competencia autonómica exclusiva en materia de asistencia social del art. $148.1.20$ª. La sentencia abrió la puerta a que las Comunidades Autónomas complementen con sus propios recursos las prestaciones asistenciales del sistema, lo que ha contribuido a la heterogeneidad territorial de la protección, visible en las rentas mínimas autonómicas y en su articulación con el IMV.

Por otro lado, el debate sobre la caja única tiene hoy una actualidad particular. El Estatuto de Autonomía del País Vasco de 1979 prevé en su art. $18$ que la Comunidad Autónoma asuma la gestión del régimen económico de la Seguridad Social. Esa transferencia ha estado pendiente durante más de cuatro décadas y es una reivindicación histórica del nacionalismo vasco y el acuerdo de investidura incluyó el compromiso de negociar esa transferencia. En julio de 2026 el Gobierno presentó sus ofertas en la materia. Según la información disponible, la negociación versa sobre la gestión del régimen económico, excluye la gestión de las pensiones y preserva expresamente la caja única y la solidaridad interterritorial, pero no se ha cerrado. El debate enfrenta dos lecturas del art. $149.1.17$ª. La primera sostiene que la gestión de la recaudación y de los recursos es separable de su titularidad, y que puede transferirse sin afectar a la unidad patrimonial. La segunda advierte que la gestión autonómica de los recursos del sistema, en una comunidad cuyo saldo de Seguridad Social es deficitario, abre el camino a una ruptura de hecho de la caja única. Los socialistas vascos llegaron a cifrar en unos 4.300 millones de euros el déficit anual de la Seguridad Social en el País Vasco, lo que muestra hasta qué punto la caja única transfiere recursos hacia las comunidades más envejecidas, incluidas algunas de renta alta.

\paragraph{Estabilidad presupuestaria: el art. 135 y la sostenibilidad como norma constitucional}
La reforma del art. 135 de la Constitución, aprobada en septiembre de 2011 en plena crisis de la deuda soberana, añadió un elemento nuevo a la Constitución económica de la Seguridad Social. El precepto reformado establece que todas las Administraciones Públicas adecuarán sus actuaciones al principio de estabilidad presupuestaria y que el volumen de deuda pública no podrá superar el valor de referencia establecido en el TFUE. 

Su desarrollo es la Ley Orgánica 2/2012, de Estabilidad Presupuestaria y Sostenibilidad Financiera, que incluye expresamente a las Administraciones de Seguridad Social en su ámbito de aplicación como uno de los subsectores del sector público, junto con la Administración central, las comunidades autónomas y las corporaciones locales.

Por tanto, aunque el art. $135$ no menciona la Seguridad Social, la somete a la disciplina fiscal general. Así se abre un debate doctrinal sobre la relación entre los mandatos sociales de los art.s $41$ y $50$ y el mandato de estabilidad. Una lectura sostiene que el art. $135$ subordina los derechos sociales a la disponibilidad presupuestaria, de modo que la suficiencia de las pensiones queda condicionada a su sostenibilidad financiera. Una lectura opuesta, mayoritaria en la doctrina iuslaboralista, entiende que ambos mandatos tienen el mismo rango y deben ponderarse. Según ella, el art. $135$ no autoriza a vaciar de contenido la garantía institucional del art. $41$, pero sí legitima que el legislador tenga en cuenta la sostenibilidad al configurar las prestaciones. En la práctica, el TC ha atendido a la situación económica y a la necesidad de sostenibilidad para avalar medidas restrictivas, como en la STC 49/2015, sin necesidad de apoyarse expresamente en el art. $135$. 

\subsubsection{Legislación ordinaria: el texto refundido y su arquitectura}
La legislación actualmente en vigor se encuentra en el texto refundido de la Ley General de la Seguridad Social, aprobado por el Real Decreto Legislativo 8/2015, de 30 de octubre. Es el cuarto texto refundido de la historia del sistema, después de los de 1966, 1974 y 1994. Como todo texto refundido, no introdujo novedades sustantivas: se limitó a integrar en un único texto las numerosas modificaciones acumuladas desde 1994, incluidas las de las leyes 40/2007, 27/2011 y 23/2013. Desde su aprobación ha sido modificado en muchas ocasiones, y especialmente por la Ley 21/2021, el Real Decreto-ley 2/2023 y el Real Decreto-ley 11/2024, de modo que su versión consolidada actual difiere considerablemente de la original.

La arquitectura del texto refundido responde a una lógica interesante. Un título preliminar y un primer título de normas generales fijan el objeto, los principios, el campo de aplicación, la estructura del sistema, la acción protectora, las reglas comunes de afiliación y cotización, la gestión y el régimen económico. Siguen títulos específicos para el Régimen General, para la protección por desempleo, para el Régimen Especial de Trabajadores Autónomos y para las prestaciones no contributivas, además de otros sobre materias como la protección de los trabajadores a tiempo parcial.\footnote{%
    La ubicación de la protección por desempleo dentro del texto refundido, y no en una ley separada, tiene un significado jurídico preciso: el desempleo es formalmente una contingencia protegida por la Seguridad Social, aunque se gestione a través del Servicio Público de Empleo Estatal y se financie en parte con transferencias del Estado.
} Tres elementos del título general tienen especial importancia. 

El primero es el campo de aplicación: distingue entre la modalidad contributiva y la no contributiva. En la primera incluye a los españoles y a los extranjeros que residan o se encuentren legalmente en España y ejerzan su actividad en territorio nacional. En la segunda, a los españoles y extranjeros residentes en territorio nacional que cumplan los requisitos de cada prestación. El criterio de conexión es, por tanto, la actividad profesional en el nivel contributivo y la residencia en el no contributivo, lo que traduce jurídicamente las dos lógicas, bismarckiana y beveridgiana. 

El segundo es la estructura del sistema. El texto refundido organiza el sistema en un Régimen General, que es la regla, y unos regímenes especiales, que son la excepción, para colectivos que por la naturaleza de su actividad requieren reglas propias. Son el de trabajadores por cuenta propia o autónomos, el de trabajadores del mar, el de la minería del carbón, el de estudiantes y el de funcionarios públicos, civiles y militares, gestionado mediante el mutualismo administrativo. Dentro del Régimen General existen además sistemas especiales, para trabajadores agrarios por cuenta ajena y empleados del hogar, con particularidades de cotización y de acción protectora. Esta estructura es el resultado de la integración progresiva de los regímenes especiales que recomendaba el Pacto de Toledo y, como vemos, conserva todavía el rastro de la antigua estratificación corporativa.

El tercero es la acción protectora. El art. $42$ del texto refundido enumera las contingencias y situaciones protegidas. Incluye la asistencia sanitaria por maternidad, enfermedad común o profesional y accidente, sea o no de trabajo; la recuperación profesional; las prestaciones económicas por incapacidad temporal, nacimiento y cuidado de menor, riesgo durante el embarazo y la lactancia, incapacidad permanente, jubilación, desempleo, muerte y supervivencia; las prestaciones familiares; y el IMV, entre otras. Este precepto es decisivo para entender el perímetro del sistema, porque incluye formalmente la asistencia sanitaria y el desempleo. En términos jurídicos, ambos son acción protectora de la Seguridad Social. En términos financieros y de gestión, la asistencia sanitaria se financia hoy íntegramente con impuestos y se gestiona por los servicios de salud autonómicos, y el desempleo se gestiona por el Servicio Público de Empleo Estatal, fuera de las entidades gestoras de la Seguridad Social.


\paragraph{Las normas que rodean al texto refundido}

El texto refundido no es la única norma que regula el sistema, y un análisis económico no puede ignorar las demás. La Ley 47/2003, General Presupuestaria, integra el presupuesto de la Seguridad Social en los Presupuestos Generales del Estado y regula su elaboración, ejecución y control. La Ley Orgánica 2/2012, de Estabilidad Presupuestaria y Sostenibilidad Financiera, somete al subsector de Administraciones de Seguridad Social a los objetivos de estabilidad y de deuda. Las leyes anuales de Presupuestos Generales del Estado fijan cada año elementos esenciales del sistema: las bases máximas y mínimas de cotización, la pensión máxima, las cuantías de las pensiones mínimas y no contributivas y las transferencias del Estado. \footnote{%
    La prórroga presupuestaria tiene por eso consecuencias directas. En 2026 siguen prorrogados los Presupuestos de 2023, y la revalorización de las pensiones y la actualización de las bases han tenido que aprobarse por decreto-ley, con el episodio de no convalidación ya relatado. Cada año, una orden de cotización desarrolla las normas de cotización, y un amplio conjunto de reglamentos regula la afiliación, la cotización, la recaudación, la gestión de las prestaciones y la colaboración de las mutuas.
}

Además, fuera del texto refundido quedan dos regímenes que conviene mencionar. El primero es el de Clases Pasivas del Estado, regulado por su propio texto refundido de 1987 y financiado directamente con el presupuesto del Estado; aunque cerrado desde 2011 a los nuevos funcionarios, que se integran en el Régimen General. El segundo es el de planes y fondos de pensiones, regulado por su texto refundido de 2002 y modificado por la Ley 12/2022. Pertenece al nivel complementario y libre, no al sistema de Seguridad Social.

\paragraph{Límite de la reforma: qué puede hacer el legislador}
Reunida la Constitución y la legislación ordinaria, puede responderse a una pregunta con gran interés expositivo: ¿qué puede hacer el legislador con el sistema de pensiones sin reformar la Constitución?

Según la doctrina expuesta, puede modificar la edad de jubilación, el periodo de cálculo de la base reguladora, la escala de porcentajes, las bases y los tipos de cotización, los topes de las pensiones, el mecanismo de revalorización y las fuentes de financiación. Puede introducir mecanismos automáticos de ajuste y derogarlos. También puede crear pilares complementarios de capitalización, voluntarios o incluso obligatorios, junto al sistema público. 

Lo que no puede hacer, según la doctrina mayoritaria, es suprimir el régimen público ni reducirlo hasta hacerlo irreconocible; tampoco congelar indefinidamente las pensiones, porque el art. $50$ exige su actualización periódica, ni excluir de toda protección a los ciudadanos en situación de necesidad.

El margen es, por tanto, muy amplio. Esa amplitud explica tanto la flexibilidad del sistema para adaptarse como su vulnerabilidad al ciclo político.

\subsubsection{Principios del sistema}
Como cualquier ordenación básica de la acción del Estado, el sistema de Seguridad Social necesita principios por su propia naturaleza. Es un conjunto de normas extraordinariamente extenso, técnico y cambiante, sometido a reformas constantes, en el que conviven reglas procedentes de épocas distintas y de lógicas distintas. Sin unos principios que den coherencia al conjunto, el sistema se convertiría en una acumulación de reglas sin orientación.\footnote{%
    Los principios cumplen, en la teoría general del derecho, tres funciones. Tienen una función informadora, porque orientan al legislador cuando crea o reforma normas. Tienen una función interpretativa, porque permiten a los tribunales y a las entidades gestoras elegir, entre varias interpretaciones posibles de una norma, la más coherente con el sistema. Y tienen una función integradora, porque permiten colmar las lagunas del ordenamiento cuando no hay norma aplicable.
}

El art. $2.1$ del texto refundido dispone literalmente que el sistema de la Seguridad Social, configurado por la acción protectora en sus modalidades contributiva y no contributiva, se fundamenta en los principios de universalidad, unidad, solidaridad e igualdad. A estos cuatro principios legales se añade la suficiencia, que tiene rango constitucional (arts. $41$ y $50$). 

Examinados uno a uno, los principios parecen compatibles. Examinados en conjunto, revelan tensiones que constituyen el núcleo de todos los debates sobre el sistema, y que conviene hacer explícitas para cerrar este apartado. La contributividad tensiona con la solidaridad: cuanto más estrecha es la relación entre lo aportado y lo recibido, menos redistribuye el sistema. La suficiencia tensiona con la sostenibilidad: cuanto más generosas son las prestaciones, más difícil es financiarlas. La universalidad tensiona con la contributividad: cuanto más se protege a quienes no han cotizado, más se aleja el sistema de la lógica del seguro. La unidad tensiona con la autonomía territorial: cuanto más se centraliza, menos se adapta a las circunstancias de cada territorio. Ninguna de estas tensiones tiene una solución técnica. Todas exigen una elección de valores, que corresponde al legislador democrático dentro del amplio margen que le concede la Constitución. El Pacto de Toledo es, precisamente, el procedimiento que el sistema político español ha ideado para hacer esas elecciones de forma consensuada.


\textcolor{red}{En cambio, la contributividad, la equidad y la unidad de caja no son principios legales del sistema. La contributividad y la equidad son criterios de política legislativa que el Pacto de Toledo ha formulado en sus sucesivas recomendaciones y que la legislación ha ido aplicando. La unidad de caja es una manifestación del principio legal de unidad y de la competencia estatal sobre el régimen económico. La distinción no es formalista. Un principio legal vincula al intérprete y puede invocarse ante los tribunales en su función interpretativa; un criterio de política legislativa solo orienta al legislador, que puede apartarse de él. Por eso una reforma que reduce la contributividad, como la cuota de solidaridad, no es ilegal, mientras que una interpretación que vulnere la igualdad sí puede ser anulada.

En cuanto a los destinatarios, los principios se dirigen a tres sujetos. Al legislador, que debe respetarlos al configurar el sistema. A las entidades gestoras, que deben aplicarlos al gestionar las prestaciones. Y a los tribunales, que deben utilizarlos para interpretar las normas. Sin embargo, no se dirigen directamente a los ciudadanos, que no pueden reclamar prestaciones invocando solo un principio, sino las normas que lo concretan.

Con estas precisiones, se examinan a continuación los cuatro principios legales, el principio constitucional de suficiencia y los dos criterios del Pacto de Toledo que el documento incluye (contributividad y equidad), seguidos de los principios instrumentales y de la cuestión de si la sostenibilidad se ha convertido en un nuevo principio del sistema.}

\paragraph{Universalidad}
No hay que confundir la universalidad como la extensión máxima de la acción protectora. La doctrina distingue dos dimensiones del principio. La universalidad subjetiva se refiere a quiénes están protegidos y exige que ninguna persona quede fuera del sistema. La universalidad objetiva se refiere a qué está protegido y exige que el sistema cubra todas las situaciones de necesidad relevantes. La definición del documento solo recoge la segunda.

El principio está regulado en el art. $2.1$ del texto refundido y tiene su fundamento en el art. 41 de la Constitución, que exige un régimen público "para todos los ciudadanos". Su justificación teórica procede de tres fuentes que ya conocemos. La primera es la ciudadanía social de Marshall, según la cual los derechos sociales son un atributo de la ciudadanía y no del trabajo. La segunda es la crítica de Titmuss a la selectividad, según la cual los programas universales evitan el estigma y la baja utilización de los selectivos. La tercera es la paradoja de la redistribución de KORPI y PALME, según la cual los programas universales generan una coalición política más amplia y redistribuyen más.

Sus ventajas son las que se derivan de esas tres fuentes: legitimidad social, ausencia de estigma, apoyo político amplio y protección de todos frente a los riesgos no asegurables. Sus desventajas son el coste presupuestario y, desde la perspectiva de la eficiencia, la posibilidad de que se protejan con recursos públicos situaciones que no lo necesitan. Es el argumento clásico a favor de la focalización, que sostiene que una pensión mínima universal paga también a quien tiene otros recursos, mientras que una pensión focalizada concentra los recursos en quien los necesita.

El principio ha experimentado en España una expansión casi continua, aunque con retrocesos. La universalidad subjetiva avanzó con la Ley 26/1990 de pensiones no contributivas, con la universalización de la sanidad (Ley General de Sanidad de 1986, culminada con el Real Decreto-ley 7/2018 tras la restricción de 2012) y con el IMV en 2020. La universalidad objetiva avanzó con la creación del sistema de dependencia en 2006 y con la extensión de la protección a nuevas contingencias, como el cuidado de menores afectados por enfermedades graves o, para los autónomos, el cese de actividad. El retroceso más visible fue el Real Decreto-ley 16/2012, que restringió el acceso a la sanidad de los inmigrantes en situación irregular. Su reversión seis años después confirma la fuerza del principio en el ordenamiento.

Los debates abiertos se sitúan hoy en tres frentes. El primero es la baja utilización de las prestaciones universales de carácter asistencial, que muestra que la universalidad jurídica no garantiza la universalidad efectiva: una prestación a la que todos tienen derecho, pero que muchos no solicitan, no es de hecho universal. El segundo es la protección de los nuevos colectivos del mercado de trabajo, como los trabajadores de plataformas, los autónomos económicamente dependientes o los trabajadores a tiempo parcial con carreras intermitentes. Para ellos la universalidad formal convive con una protección efectiva débil, que es el problema de la dualización examinado en 1.1.3. El tercero es la propuesta, formulada desde posiciones diversas, de una renta básica universal, que llevaría el principio a su extremo lógico al desligar completamente la protección de la situación de necesidad y del trabajo.

\paragraph{Unidad y la caja única}
De la misma forma que el anterior, no hay que confundir el principio de unidad como unidad de caja, definida como que el Estado es el único titular de todos los recursos, obligaciones y prestaciones de la Seguridad Social. Este formulación, muy habitual, contiene un error que conviene corregir. 

El titular del patrimonio de la Seguridad Social no es el Estado, sino la propia Seguridad Social. El texto refundido declara que las cuotas, los bienes, los derechos, las acciones y los recursos de cualquier otro género de la Seguridad Social constituyen un patrimonio único afecto a sus fines, distinto del patrimonio del Estado, y atribuye su titularidad a la Tesorería General de la Seguridad Social. La distinción tiene consecuencias. El patrimonio de la Seguridad Social está afectado a una finalidad y no puede destinarse a otros fines del Estado; por eso los excedentes contributivos se destinaban al Fondo de Reserva y no al Tesoro. Y la relación financiera entre el Estado y la Seguridad Social es una relación entre dos sujetos distintos, lo que explica que el Estado preste dinero a la Seguridad Social, en lugar de simplemente pagar su déficit. Es la base jurídica de la deuda de la Seguridad Social con el Estado que se analiza en 2.2.3.

El principio de unidad, recogido en el art. $2.1$ del texto refundido, tiene un contenido más amplio que la caja única. Abarca la unidad normativa, porque todo el sistema se rige por una legislación común; la unidad de gestión, porque las prestaciones de cada tipo las gestiona una única entidad para todo el territorio; y la unidad financiera, que es la caja única. Su fundamento constitucional es la competencia estatal exclusiva sobre el régimen económico del art. $149.1.17$ª, que tratamos antes.

Las ventajas de la unidad son tres. La primera es la mancomunación del riesgo demográfico entre territorios, que protege a las comunidades envejecidas. La segunda son las economías de escala en la gestión: un sistema único tiene costes administrativos por prestación mucho menores que diecisiete sistemas autonómicos, y la gestión pública de pensiones en España tiene costes administrativos muy reducidos en relación con el volumen de prestaciones gestionadas. La tercera es la movilidad laboral sin fricciones, porque un trabajador que cambia de comunidad no pierde ni fragmenta sus derechos. Sus desventajas, señaladas por los defensores de la descentralización, son la menor capacidad de adaptación a las circunstancias territoriales y la ausencia de corresponsabilidad fiscal: las comunidades no soportan el coste de las decisiones que afectan al empleo y a la afiliación en su territorio. El debate sobre la transferencia de la gestión del régimen económico al País Vasco, expuesto en 1.3.1, es la manifestación actual de esta tensión.

\paragraph{Solidaridad: Intergeneracional, Intrageneracional e Interterritorial}
La solidaridad intergeneracional es solo una de las dimensiones de este principio: el hecho de que los trabajadores actuales sufragan las prestaciones actuales bajo la promesa de que, cuando sean población inactiva, los trabajadores de entonces financiarán sus prestaciones. No obstante, la doctrina distingue tres dimensiones del principio, recogido en el art. $2.1$ del texto refundido: la intergeneracional, la intrageneracional y la interterritorial.

La solidaridad intergeneracional es la más característica de un sistema de reparto. La definición del documento describe con precisión su mecánica: es el llamado contrato entre generaciones. La expresión evoca una célebre frase de Edmund Burke, que en sus Reflexiones sobre la Revolución en Francia (1790) definió la sociedad como una asociación no solo entre los vivos, sino entre los vivos, los muertos y los que están por nacer. Pero la metáfora contractual tiene un problema que la crítica ha señalado insistentemente: es un contrato que las generaciones futuras no han firmado, porque no existían o no tenían edad para votar cuando se estableció. La justificación teórica más sólida de esta solidaridad es la de GORDON y VARIAN (1988) y SHILLER (1999): el reparto completa un mercado de seguros entre generaciones que no puede existir de forma privada. La crítica más sólida es la de la contabilidad generacional de AUERBACH, GOKHALE y KOTLIKOFF (1991), que se examina más adelante. Esta muestra que los sistemas de reparto pueden transferir a las generaciones futuras cargas netas considerables, de modo que lo que se presenta como solidaridad es, en parte, una transferencia forzosa de los jóvenes y de los no nacidos hacia los mayores. La propia denominación del MEI de 2021 y 2023 es un reconocimiento legislativo de que esa solidaridad puede degenerar en inequidad si no se corrige.

La solidaridad intrageneracional es la redistribución entre miembros de la misma generación. Se realiza a través de varios elementos de diseño. La base máxima de cotización y la pensión máxima limitan, respectivamente, lo que aportan y lo que reciben los de mayores ingresos, y la pensión máxima tiende a crecer menos que la base máxima. Los complementos a mínimos elevan las pensiones más bajas. Las reglas de integración de lagunas de cotización protegen a quienes tuvieron periodos sin empleo. Y el complemento para la reducción de la brecha de género corrige parte de la desventaja de las carreras femeninas interrumpidas. Esta redistribución progresiva, como ya vimos, se ve parcialmente contrarrestada por el gradiente de longevidad, que hace que los más ricos cobren durante más años.

La solidaridad interterritorial es la que se realiza a través de la caja única, que transfiere recursos desde las comunidades con más cotizantes en relación con sus pensionistas hacia las más envejecidas. Es la dimensión del principio más directamente afectada por los debates sobre descentralización.

\paragraph{Igualdad}

El documento formula este principio como equidad e igualdad de derechos, con independencia del momento y lugar de residencia del asegurado. La formulación mezcla dos conceptos que conviene separar: la igualdad, que es un principio legal (art. $2.1$ del texto refundido) con fundamento en el art. $14$ de la Constitución, y la equidad, que es un criterio de política legislativa y se examina más adelante. La referencia al "lugar de residencia" alude a la dimensión territorial de la igualdad, cuyo fundamento constitucional es el art. $149.1.1$ª, que atribuye al Estado la regulación de las condiciones básicas que garanticen la igualdad de todos los españoles en el ejercicio de sus derechos.

El principio de igualdad exige que el sistema no establezca diferencias de trato injustificadas entre los asegurados. Su aplicación ha sido una de las fuentes más fecundas de reformas del sistema, por la vía de la jurisprudencia constitucional y europea. Tres casos lo ilustran bien. El Tribunal Constitucional (STC 103/1983) declaró inconstitucional la exigencia de requisitos adicionales a los viudos varones para acceder a la pensión de viudedad, que no se exigían a las viudas. Fue una de las primeras aplicaciones del principio de igualdad entre sexos en el ámbito de la Seguridad Social, y paradójicamente beneficiaba a los hombres. En la STC 61/2013, siguiendo la sentencia del TJUE en el asunto Elbal Moreno (C-385/12, de 22 de noviembre de 2012), el Tribunal Constitucional declaró inconstitucional la forma de computar los periodos de cotización de los trabajadores a tiempo parcial para acceder a las pensiones. Consideró que perjudicaba desproporcionadamente a las mujeres, que son la mayoría de esos trabajadores, y que constituía por tanto una discriminación indirecta por razón de sexo. En la STC 91/2019, tras la sentencia del Tribunal de Justicia en el asunto Villar Láiz (C-161/18), el Tribunal Constitucional anuló también el coeficiente de parcialidad aplicado al cálculo de la cuantía de la pensión, por la misma razón.

La igualdad entre regímenes es otro frente del principio. La integración progresiva de los regímenes especiales recomendada por el Pacto de Toledo ha tendido a homogeneizar la protección, pero persisten diferencias. Las más relevantes son las de los autónomos, que durante décadas cotizaron mayoritariamente por la base mínima y generaron pensiones bajas, y las de las Clases Pasivas, con reglas de cálculo propias. La reforma de la cotización de los autónomos por rendimientos reales, de 2022, puede leerse como una aplicación del principio de igualdad, porque aproxima su esfuerzo contributivo al de los asalariados con ingresos comparables. La brecha de género en las pensiones, de la que se trata más adelante, es hoy el debate más relevante en materia de igualdad. Sus causas están en las carreras laborales de las mujeres (más interrumpidas, más a tiempo parcial y con salarios más bajos), y su corrección enfrenta la igualdad de trato, que exige reglas neutras, con la igualdad material, que exige compensar desventajas.

\paragraph{Suficiencia}
La suficiencia se define como que las prestaciones sean suficientes para garantizar un nivel de bienestar adecuado. El principio tiene rango constitucional, porque el art. $41$ exige prestaciones suficientes y el art. $50$, pensiones adecuadas que garanticen la suficiencia económica en la tercera edad. Pero no figura entre los principios del art. $2.1$ del texto refundido. Su fundamento teórico es la categoría de bien preferente de MUSGRAVE y, en una formulación más moderna, el enfoque de las capacidades de SEN.

La principal dificultad del principio es su indeterminación: ¿cuándo es suficiente una pensión? La literatura y las instituciones internacionales han propuesto varias respuestas, que corresponden a distintas concepciones de la suficiencia. El \textit{Pension Adequacy Report} de la Comisión Europea mide la suficiencia en tres dimensiones: la protección frente a la pobreza, medida por el riesgo de pobreza de los mayores; el mantenimiento de la renta, medido por la tasa de reemplazo de la pensión respecto del salario previo; y la duración del periodo de jubilación. La primera concepción es la beveridgiana del mínimo, que se aproxima con el umbral de la pobreza relativa, fijado convencionalmente en el sesenta por ciento de la mediana de la renta equivalente. La segunda es la bismarckiana del mantenimiento del nivel de vida, que se mide con la tasa de reemplazo. El legislador español ha avanzado en los últimos años hacia una definición operativa de la suficiencia ligada a la primera concepción, al vincular las cuantías de las pensiones mínimas y de las no contributivas al umbral de la pobreza, con una senda de convergencia que se examinará más adelante.

El principio de suficiencia es el que entra en conflicto más directo con la sostenibilidad. El debate sobre la reforma de 2013 fue, en esencia, un conflicto entre ambos. Los mecanismos automáticos garantizaban la sostenibilidad a costa de reducir la tasa de sustitución, es decir, la pensión media en relación con el salario medio, y los críticos argumentaron que esa reducción vulneraba la suficiencia. La derogación de 2021 dio prioridad a la suficiencia, y el debate actual es si esa prioridad es compatible con la sostenibilidad (bloque IV).

\paragraph{Contributividad y la equidad: dos criterios del Pacto de Toledo}
El documento define la contributividad como la proporcionalidad entre lo percibido y lo aportado. Como se ha dicho, no es un principio legal del sistema, sino un criterio de política legislativa que el Pacto de Toledo ha formulado reiteradamente desde 1995, al recomendar el reforzamiento del carácter contributivo del sistema. Su fundamento teórico es la lógica bismarckiana del seguro y la función de hucha de BARR: la pensión contributiva es la devolución, a lo largo de la jubilación, de lo aportado durante la vida activa.

Sus ventajas se sitúan en tres planos. En el de la legitimidad, los cotizantes perciben la pensión como un derecho ganado y no como una dádiva, lo que refuerza el apoyo al sistema. En el de los incentivos, cuanto más estrecha es la relación entre cotizaciones y prestaciones, más se parece la cotización a un precio y menos a un impuesto, y menos desincentiva el empleo formal y la declaración de salarios reales, como mostró SUMMERS (1989). En el de la equidad, a igual esfuerzo, igual prestación. Sus desventajas son el reverso de la solidaridad. Un sistema estrictamente contributivo traslada a la vejez las desigualdades de la vida activa: quien tuvo una carrera corta, intermitente o mal pagada recibe una pensión baja y esto afecta especialmente a las mujeres y a los trabajadores precarios.

La contributividad ha experimentado en España una evolución ambivalente. Todas las reformas que alargaron el periodo de cálculo de la base reguladora, de dos a ocho años en 1985, a quince en 1997, a veinticinco en 2011 y a una opción de veintinueve años con exclusión de los peores en 2023, la reforzaron. También la reforzó la cotización de los autónomos por rendimientos reales en 2022. En la dirección opuesta, la debilitaron la subida de la base máxima por encima de la pensión máxima, la cuota de solidaridad y la cotización del MEI, que no generan derechos individuales. La reforma de 2021-2023 combinó, por tanto, elementos que refuerzan la contributividad en el cálculo de la pensión con otros que la debilitan en la financiación. Esta ambivalencia es el núcleo de la reforma silenciosa, que se analizará más adelante.

La equidad, que el documento une a la igualdad, es un concepto con varias acepciones que conviene separar. La equidad actuarial exige que el valor actual de las prestaciones que recibe cada individuo sea igual al de sus cotizaciones; es la contributividad llevada a su extremo, y es el principio que inspira las cuentas nocionales. La equidad horizontal exige que individuos con igual esfuerzo contributivo reciban igual trato. La equidad vertical exige que el sistema trate de forma diferente a individuos con circunstancias diferentes, lo que justifica la redistribución. La equidad intergeneracional exige que las distintas generaciones reciban un trato comparable, lo que el MEI pretende garantizar. La literatura actuarial estadounidense formuló ya en 1938, en un trabajo clásico de HOHAUS, la tensión fundamental de todo sistema de pensiones público entre la equidad individual (\textit{individual equity}), que exige que cada uno reciba en proporción a lo aportado, y la suficiencia social (\textit{social adequacy}), que exige que todos reciban lo necesario para vivir con dignidad. Todo el diseño del sistema español puede leerse como un compromiso entre ambas, y cada reforma como un desplazamiento del punto de equilibrio en una u otra dirección.

\paragraph{Principios instrumentales}
Junto a los principios que definen la naturaleza del sistema, existen otros de carácter instrumental que garantizan su funcionamiento y que el documento no menciona. El principio de participación, con fundamento en el art. 129.1 de la Constitución, exige la participación de los interesados en la gestión del sistema. Se concreta en la presencia de las organizaciones sindicales y empresariales en los órganos de control de las entidades gestoras, que se examina en 1.4. El principio de irrenunciabilidad declara nulo todo pacto por el que los trabajadores renuncien a los derechos que les confiere la ley, y protege al trabajador frente a presiones del empleador. El principio de protección de las prestaciones establece que no pueden ser objeto de cesión, embargo, retención, compensación o descuento, salvo excepciones legales. Protege así la función de sustitución de rentas que justifica su existencia. El principio de automaticidad de las prestaciones garantiza que el trabajador reciba sus prestaciones aunque el empresario haya incumplido sus obligaciones de afiliación o cotización. En determinados supuestos, la entidad gestora anticipa el pago y repite después contra el empresario responsable, de modo que el trabajador queda protegido frente al riesgo moral de su empleador.

\paragraph{¿Un nuevo principio? La sostenibilidad}
La última cuestión es si la sostenibilidad se ha convertido en un principio del sistema. Como hemos visto, no figura en el art. $2.1$ del texto refundido, pero ha ido adquiriendo en el ordenamiento un estatus creciente. Tiene fundamento constitucional indirecto en el art. 135 y en la Ley Orgánica 2/2012, que se refiere expresamente a la sostenibilidad financiera de las Administraciones de Seguridad Social. La Ley 23/2013 la convirtió en el eje de su regulación, con el propio nombre de factor de sostenibilidad. La Ley 21/2021, que derogó los mecanismos de 2013, mantuvo la referencia en su título, que habla de medidas de refuerzo de la sostenibilidad financiera y social del sistema. La expresión sostenibilidad social es significativa: intenta reconciliar la sostenibilidad con la suficiencia, al sugerir que un sistema que no garantiza pensiones suficientes tampoco es sostenible, porque pierde el apoyo social que lo sostiene. Es una idea que conecta directamente con la paradoja de la redistribución y con la reversión de la reforma de 2013.

El debate doctrinal sobre este punto es abierto. Para una parte de la doctrina, la sostenibilidad es ya un principio implícito del sistema, derivado del art. 135 y de la legislación de estabilidad, que debe ponderarse con los principios tradicionales. Para otra, es solo una restricción económica, no un principio jurídico, y convertirla en principio supondría subordinar los derechos sociales a la disponibilidad presupuestaria. En la práctica legislativa, la sostenibilidad funciona hoy como un principio procedimental. La cláusula de salvaguarda del MEI obliga a evaluar periódicamente la sostenibilidad del sistema y a adoptar medidas si se superan determinados umbrales, lo que la convierte en un parámetro jurídicamente vinculante, aunque no en un principio en sentido estricto.


\subsubsection{Marco supranacional}
El legislador español no configura el sistema en el vacío. Un conjunto creciente de normas internacionales y europeas condiciona su margen de maniobra, en tres planos de intensidad jurídica muy distinta. El primero es el derecho internacional de los derechos humanos y del trabajo, que fija estándares mínimos de protección. El segundo es el derecho de la Unión Europea, que coordina los sistemas nacionales y prohíbe la discriminación. El tercero es la gobernanza económica europea, que condiciona la sostenibilidad de los sistemas mediante la supervisión fiscal y, en los últimos años, mediante la financiación condicionada.

En conjunto, el marco supranacional produce un efecto de doble tenaza sobre el legislador español. Por un lado, los estándares internacionales de derechos humanos y las normas europeas de igualdad y de renta mínima presionan hacia la suficiencia y la ampliación de la protección. Por otro, la gobernanza económica europea presiona hacia la sostenibilidad y la contención del gasto. La tensión entre suficiencia y sostenibilidad que recorre todo el tema no es solo interna al ordenamiento español; se reproduce, amplificada, en el plano europeo. Allí el Pilar Europeo de Derechos Sociales y el Semestre Europeo expresan, respectivamente, las dos caras del mismo dilema.

\paragraph{Derecho internacional: derechos humanos y estándares de la OIT}
La Declaración Universal de Derechos Humanos (1948) reconoce en su art. $22$ el derecho de toda persona a la seguridad social, y en su art. $25$, el derecho a un nivel de vida adecuado y a los seguros en caso de desempleo, enfermedad, invalidez, viudez y vejez. El Pacto Internacional de Derechos Económicos, Sociales y Culturales (1966), ratificado por España en 1977, reconoce en su art. $9$ el derecho de toda persona a la seguridad social. En su desarrollo normativo se establece que el derecho a la seguridad social exige disponibilidad, cobertura de los riesgos sociales, suficiencia de las prestaciones, accesibilidad y no discriminación, y que las medidas regresivas deben justificarse plenamente y ser temporales, proporcionadas y no discriminatorias. Este principio de no regresividad ha sido invocado con frecuencia en el debate doctrinal sobre los recortes de la crisis, aunque su eficacia jurídica en el ordenamiento español es limitada.

El Convenio 102 de la OIT (1952), sobre la norma mínima de la seguridad social, es el instrumento internacional más importante en la materia. Define nueve ramas de protección y fija para cada una niveles mínimos de cobertura, de cuantía de las prestaciones y de condiciones de acceso. España lo ratificó en 1988, aceptando varias de sus ramas. En el ámbito del Consejo de Europa, el Código Europeo de Seguridad Social (1964) eleva esos estándares mínimos, y España lo ratificó en la década de los noventa. Estos instrumentos tienen, por el art. $96$ de la Constitución, fuerza de ley en el ordenamiento interno. Según el art. $10.2$, además, las normas relativas a los derechos fundamentales deben interpretarse de conformidad con ellos. Su eficacia práctica es limitada, porque sus estándares son mínimos y el sistema español los supera con holgura en la mayoría de las ramas.

Más relevante es la Carta Social Europea. Reconoce, entre otros, el derecho a la seguridad social (art. $12$), el derecho a la asistencia social y médica (art. $13$), el derecho de las personas mayores a la protección social (art. $23$) y el derecho a la protección contra la pobreza y la exclusión social (art. $30$). La ratificación del Protocolo de reclamaciones colectivas tiene una importancia práctica considerable. Permite a sindicatos, organizaciones empresariales y determinadas organizaciones no gubernamentales presentar reclamaciones ante el Comité Europeo de Derechos Sociales, que puede declarar que la legislación española no se ajusta a la Carta. Las decisiones del Comité, aunque no son formalmente vinculantes como una sentencia, han sido utilizadas por los tribunales españoles para inaplicar normas internas en materia laboral. Abren así una vía de control de la suficiencia de las prestaciones sociales que antes no existía.

Completa este plano una red de convenios bilaterales de Seguridad Social con países de emigración y de inmigración, que permiten sumar periodos de cotización y exportar prestaciones. A ella se añade el Convenio Multilateral Iberoamericano de Seguridad Social, firmado en 2007 y en vigor desde 2011, que coordina los sistemas de España, Portugal y un conjunto de países latinoamericanos.

\paragraph{Derecho de la Unión: coordinación, igualdad y derecho blando}
El derecho de la Unión no armoniza los sistemas de Seguridad Social, que siguen siendo competencia de los Estados miembros, pero los condiciona por varias vías. El art. $153$ del TFUE atribuye a la Unión una competencia de apoyo y complemento de la acción de los Estados en materia de seguridad social y protección social de los trabajadores. La Carta de los Derechos Fundamentales de la Unión Europea, en su art. $34$, reconoce y respeta el derecho de acceso a las prestaciones de seguridad social y a los servicios sociales de acuerdo con las modalidades establecidas por el derecho de la Unión y las legislaciones nacionales.

La vía más antigua e intensa es la \textbf{coordinación}. El art. $48$ del Tratado habilita a la Unión para adoptar las medidas necesarias en materia de seguridad social para garantizar la libre circulación de trabajadores. Su legislación de desarrollo no armoniza las prestaciones, pero establece cuatro reglas fundamentales.\footnote{%
    Las medidas se recogen hoy en el Reglamento (CE) 883/2004 y su reglamento de aplicación, el 987/2009.
} La primera es la aplicación de la legislación de un único Estado a cada trabajador. La segunda es la igualdad de trato entre nacionales y ciudadanos de otros Estados miembros. La tercera es la totalización de los periodos de seguro cumplidos en distintos Estados. La cuarta es la exportación de las prestaciones a otro Estado miembro. Su lógica económica es evitar que los sistemas nacionales de protección social se conviertan en barreras a la movilidad laboral, lo que sería contrario al mercado interior.

La segunda vía es la \textbf{igualdad} de trato entre hombres y mujeres. La Directiva 79/7/CEE prohíbe la discriminación por razón de sexo en los regímenes legales de seguridad social, tanto directa como indirecta. Su aplicación por el TJUE ha tenido un impacto considerable en el sistema español, como vimos en el principio de igualdad. Las sentencias en los asuntos Elbal Moreno (2012) y Villar Láiz (2019) obligaron a reformar el cómputo de los periodos de cotización a tiempo parcial. El asunto C-450/18, consideró discriminatorio para los hombres el complemento de maternidad de las pensiones, que se concedía solo a las mujeres con dos o más hijos; sentencia que obligó al legislador a sustituirlo por uno accesible también a los hombres que acrediten un perjuicio en su carrera por el cuidado de los hijos. 

La tercera vía es el \textbf{soft law}, que no obliga jurídicamente pero orienta las políticas nacionales. El Pilar Europeo de Derechos Sociales, proclamado en la cumbre de Gotemburgo de noviembre de 2017, recoge veinte principios. Entre ellos están el derecho a una protección social adecuada (p. 12), a prestaciones de renta mínima (p. 14), a pensiones de jubilación que garanticen una renta adecuada (p. 15) y a cuidados de larga duración asequibles y de calidad (p. 18). Su desarrollo ha dado lugar a varias recomendaciones del Consejo, en 2019 sobre el acceso a la protección social de los trabajadores por cuenta ajena y por cuenta propia y en 2023 sobre una renta mínima adecuada. Esta última insta a los Estados a garantizar que las rentas mínimas alcancen progresivamente un nivel adecuado, vinculado al umbral de riesgo de pobreza, y a reducir la falta de solicitud. En el ámbito del segundo y tercer pilar, la IORP II (Directiva 2016/2341) regula los fondos de pensiones de empleo, y el Reglamento (UE) 2019/1238 crea el producto paneuropeo de pensiones individuales (PEPP).

La asimetría entre estas tres vías ha sido objeto de una crítica influyente. SCHARPF (1999) distinguió entre la integración negativa, que elimina barreras al mercado y se impone mediante el derecho de la competencia y las libertades del mercado interior, y la integración positiva, que construye políticas comunes y exige acuerdos políticos difíciles. Según SCHARPF, la Unión ha avanzado mucho más en la primera que en la segunda. Los sistemas nacionales de protección social quedan así sometidos a las exigencias del mercado interior y de la disciplina fiscal sin que exista una política social europea capaz de compensarlas. FERRERA (2005) desarrolló una tesis complementaria: la integración europea ha erosionado las fronteras de los Estados del bienestar nacionales, que se construyeron sobre la base de una comunidad cerrada de contribuyentes y beneficiarios, sin construir unas fronteras europeas equivalentes.

\subsubsection{Gobernanza económica europea: supervisión fiscal y financiación condicionada}
El condicionamiento más intenso procede de la gobernanza económica de la Unión, dirigido a la sostenibilidad de las finanzas públicas. Las recomendaciones dirigidas, en el marco del Semestre Europeo, a España han incluido repetidamente la sostenibilidad del sistema de pensiones.\footnote{%
    Esas recomendaciones se apoyan en las proyecciones de largo plazo de los Ageing Reports del Grupo de Trabajo sobre Envejecimiento del Comité de Política Económica, que se elaboran cada tres años con una metodología común y se examinan más adelante.
}

El condicionamiento se hizo más intenso con el Mecanismo de Recuperación y Resiliencia (MRR), que vinculó el desembolso de los fondos NGEU al cumplimiento de hitos de reforma e inversión. La reforma de pensiones de 2021 y 2023 constituyó uno de los hitos pactados con la Comisión en el marco del Plan de Recuperación, Transformación y Resiliencia (PRTR). Es la primera vez que una reforma de pensiones española se vincula formalmente a la financiación europea y, en consecuencia, suscita un debate sobre la legitimidad democrática de las reformas: si la reforma respondió a una decisión soberana del legislador español o a una condición impuesta desde fuera.\footnote{%
    Un rasgo singular de la reforma de 2023 refuerza esta cuestión. La cláusula de salvaguarda del MEI utiliza como referencia las proyecciones de gasto del Ageing Report de la Comisión Europea. De ese modo, un documento técnico elaborado por un grupo de trabajo europeo se incorpora como parámetro de una norma jurídica española que puede activar automáticamente subidas de cotizaciones.
}

La reforma de las reglas fiscales europeas de 2024, estableció un nuevo marco basado en planes fiscales-estructurales nacionales a medio plazo. El plan de España para 2025-2028, prevé un periodo de ajuste de siete años, en lugar de los cuatro ordinarios. La ampliación está condicionada al cumplimiento de un conjunto de reformas e inversiones. Según la propia recomendación del Consejo, el plan tiene en cuenta las medidas de ingresos adoptadas en 2023 junto con la reforma de las pensiones, que mejoran la dinámica de la deuda a largo plazo mediante un aumento de las cotizaciones sociales. El nuevo marco refuerza, por tanto, la integración de la política de pensiones en la supervisión fiscal europea: la sostenibilidad del sistema de pensiones es ya un componente explícito de la evaluación de la sostenibilidad de la deuda pública española.



\subsection{Arquitectura institucional y perímetro}
Antes de describir las instituciones hay que resolver una cuestión previa: el perímetro, es decir, qué se entiende por Seguridad Social y qué queda dentro y fuera. 

Resuelta esa cuestión, se examinan la dirección política del sistema, la gestión directa por las entidades gestoras y los servicios comunes, la colaboración en la gestión de las mutuas y las empresas, las instituciones de protección social que están fuera del sistema en sentido estricto, y los órganos de control, evaluación y participación. 

La arquitectura institucional del sistema español presenta un balance ambivalente resultado de una racionalización incompleta: el Real Decreto-ley 36/1978 sustituyó la yuxtaposición franquista por una organización funcional coherente, pero la descentralización territorial, el reparto de competencias entre ministerios y la supervivencia de regímenes históricos han vuelto a fragmentar el sistema. 

Por un lado, el sistema dispone de una organización funcional coherente en su núcleo: entidades gestoras especializadas por tipo de prestación, una caja única gestionada por una Tesorería con amplia capacidad recaudatoria y tecnológica, y unos costes de administración reducidos en relación con el volumen de prestaciones. A ello se añaden mecanismos de control interno y externo consolidados, y una institución fiscal independiente que ha ganado protagonismo en la evaluación de su sostenibilidad. Por otro lado, esa organización está atravesada por varias líneas de fragmentación. La ministerial reparte entidades gestoras y prestaciones conexas entre varios departamentos. La territorial separa la gestión de la sanidad, la dependencia, las rentas mínimas y las políticas activas de empleo, por un lado, de la gestión de las pensiones y el desempleo, por otro. La histórica mantiene regímenes y entidades singulares como el Instituto Social de la Marina, las Clases Pasivas o el mutualismo administrativo. Y la público-privada confía a las mutuas la gestión de prestaciones públicas con un conflicto de intereses estructural. 

Los proyectos para corregir esta fragmentación, como la agencia única de la Seguridad Social o la Agencia Española de Empleo, llevan años sin culminarse, lo que confirma que la arquitectura institucional, como las propias prestaciones, está sometida a las inercias y a los equilibrios de la economía política que recorren todo el tema.

\subsubsection{Perímetro: de qué hablamos cuando hablamos de Seguridad Social}
La expresión Seguridad Social se utiliza en el debate público, en el derecho, en la estadística y en la contabilidad nacional con significados distintos, y buena parte de las confusiones sobre el déficit de la Seguridad Social o sobre el gasto en protección social proceden de mezclarlos. Conviene deshacer esa confusión distinguiendo tres conceptos.

El primero es el \textbf{concepto jurídico}, que es el del texto refundido de la Ley General de la Seguridad Social. Según ese concepto, el sistema está formado por el conjunto de regímenes (el General y los especiales) que protegen frente a las contingencias enumeradas en su art. $42$, y es gestionado por las entidades gestoras, los servicios comunes y las entidades colaboradoras que la ley designa. Como vimos, este concepto es amplio en su acción protectora, porque incluye formalmente la asistencia sanitaria y el desempleo, pero estricto en su organización, porque solo incluye las entidades que la ley califica como gestoras o colaboradoras del sistema.

El segundo es el \textbf{concepto funcional} o estadístico de protección social, que es el que utilizan las estadísticas internacionales comparadas. El Sistema Europeo de Estadísticas Integradas de Protección Social (ESSPROS) de Eurostat define la protección social como el conjunto de intervenciones de organismos públicos o privados destinadas a aliviar a los hogares y a los individuos de la carga de un conjunto definido de riesgos o necesidades, siempre que no exista una contrapartida simultánea y equivalente ni un acuerdo individual. Esos riesgos son la enfermedad y la asistencia sanitaria, la invalidez, la vejez, la supervivencia, la familia y los hijos, el desempleo, la vivienda y la exclusión social. La Clasificación COFOG ([\authorfont{Ver Tema 4.B.23}]), distingue por su parte la función de protección social de la de salud. Estos conceptos no dependen de qué institución gestiona el gasto, sino de su finalidad. Incluyen, por tanto, el gasto sanitario de las Comunidades Autónomas, las pensiones de Clases Pasivas pagadas por el Estado, la dependencia financiada por las Comunidades Autónomas y el Estado, y las rentas mínimas autonómicas, que el concepto jurídico deja fuera. Son los conceptos adecuados para las comparaciones internacionales que se examinan más adelante.

El tercero es el \textbf{concepto contable} de la contabilidad nacional y comprende las unidades públicas que gestionan cotizaciones y prestaciones sociales. Además del Sistema de la Seguridad Social, incluyen el Servicio Público de Empleo Estatal y el Fondo de Garantía Salarial. Es el subsector S.1314 del Sistema Europeo de Cuentas (SEC 2010), que la Intervención General de la Administración del Estado denomina Administraciones de Seguridad Social. Incluye las entidades gestoras y los servicios comunes del sistema, las mutuas colaboradoras (cuya actividad de gestión de prestaciones públicas se considera pública), el SEPE y el FOGASA. No incluye, en cambio, las pensiones de Clases Pasivas, que se registran en la Administración Central, ni las mutualidades de funcionarios, clasificadas también en ella, ni el gasto sanitario y de dependencia de las Comunidades Autónomas, que se registra en el subsector de las comunidades autónomas.

\paragraph{Por qué importa el perímetro: déficits distintos para una misma realidad}
La distinción entre estos tres conceptos no es un ejercicio de taxonomía, sino que tiene consecuencias económicas y políticas de primer orden. 

La primera es que \textbf{no existe un único déficit} de la Seguridad Social, sino varios, según el concepto que se utilice. Como sabemos, el déficit del presupuesto del Sistema de la Seguridad Social, medido en términos de liquidación presupuestaria, es la diferencia entre sus derechos y sus obligaciones reconocidos de carácter no financiero; incluyendo las transferencias del Estado, pero no los préstamos. El déficit del subsector de Administraciones de Seguridad Social en contabilidad nacional aplica los criterios del SEC 2010 e incluye el SEPE y el FOGASA, de modo que el superávit habitual del SEPE en épocas de bajo desempleo compensa parcialmente el déficit del sistema de pensiones. Y el déficit del sistema contributivo en sentido estricto, que algunos analistas calculan, compara las cotizaciones con el gasto en pensiones contributivas sin contar las transferencias del Estado; siendo mucho mayor que los anteriores.

La segunda consecuencia es todavía más importante: las transferencias entre el Estado y la Seguridad Social no alteran el déficit del conjunto de las Administraciones Públicas. Cuando el Estado transfiere recursos a la Seguridad Social para financiar un gasto, el déficit de la Seguridad Social se reduce en la misma cuantía en que aumenta el del Estado, y el déficit consolidado permanece inalterado.\footnote{%
    Por eso la decisión adoptada a partir de 2021 de aumentar las transferencias del Estado para financiar los llamados gastos impropios mejoró sustancialmente el saldo aparente de la Seguridad Social. Pero no mejoró en nada la situación financiera de las Administraciones Públicas en su conjunto: se limitó a trasladar el déficit de un subsector a otro.
} Lo mismo ocurre, en sentido inverso, con los préstamos del Estado a la Seguridad Social: son una operación financiera que no computa como déficit de ninguno de los dos subsectores, pero generan una deuda interna entre ambos que desaparece en la consolidación. Por tanto, cuando se oye que la Seguridad Social "ha reducido su déficit", la primera pregunta que debemos haceros es si lo ha hecho porque ha mejorado la relación entre cotizaciones y prestaciones o porque el Estado le ha transferido más recursos. Cuestión que, por su relevancia, examinaremos más tarde.

La tercera consecuencia afecta a las comparaciones internacionales. Un país que financia sus pensiones con impuestos generales a través del presupuesto del Estado y otro que las financia con cotizaciones a través de un fondo de seguridad social pueden tener el mismo gasto en pensiones y la misma presión fiscal. Sin embargo, el primero tendrá un subsector de Seguridad Social pequeño y el segundo, uno grande. Por eso las comparaciones deben hacerse siempre en términos funcionales (COFOG o ESSPROS) y nunca en términos institucionales.

Con esta clarificación, el resto del apartado examina las instituciones del sistema en sentido jurídico y las que, estando fuera del sistema, forman parte de la protección social en sentido funcional.

\subsubsection{Dirección política del sistema: Ministerio de Inclusión, Seguridad Social y Migraciones y la Secretaría de Estado}
Actualmente, el Ministerio de Inclusión, Seguridad Social y Migraciones (MISSM) es el departamento ministerial con competencias en materia de Seguridad Social y Clases Pasivas. Además, la gestión de la Seguridad Social se atribuye a diversos entes, públicos y privados, con personalidad jurídica propia y en general adscritos al MISSM, y que en la cúspide jerárquica se encuentra la Secretaría de Estado de la Seguridad Social y Pensiones, de la que dependen los entes con responsabilidad en la gestión.

El ministerio actual se creó en enero de 2020, en la reestructuración de departamentos ministeriales que acompañó la formación del primer gobierno de coalición de la democracia, y está dirigido desde noviembre de 2023 por Elma Saiz. Su denominación y su perímetro son el resultado de una larga serie de reorganizaciones. Durante casi toda la democracia, la Seguridad Social compartió ministerio con el trabajo y el empleo. La separación de 2020 entre un Ministerio de Trabajo y Economía Social y un Ministerio de Inclusión, Seguridad Social y Migraciones respondió en buena medida a la lógica del reparto de carteras en un gobierno de coalición. Pero tuvo consecuencias duraderas: las pensiones y el desempleo quedaron en ministerios distintos, pese a ser ambos acción protectora de la Seguridad Social según el art. $42$ del texto refundido.\footnote{%
    Esta fragmentación tiene explicaciones de economía política que la teoría de la burocracia ayuda a entender. NISKANEN (1971) y su modelo de burócrata maximizador de presupuesto, fue matizado por DUNLEAVY (1991) con su teoría del \textit{bureau-shaping}: los altos responsables no buscan tanto maximizar el presupuesto total como configurar sus organizaciones para concentrar en ellas las funciones de dirección y diseño de políticas, más prestigiosas, y desprenderse de las de gestión rutinaria. En un gobierno de coalición, a esta lógica se añade la del reparto de carteras entre socios, que tiende a multiplicar los departamentos y a fragmentar las competencias según criterios políticos más que funcionales.
}

El ministerio agrupa hoy tres grandes áreas. La primera es la Seguridad Social y las pensiones, que incluye desde 2020 la gestión de las Clases Pasivas, hasta entonces en el Ministerio de Hacienda. La segunda es la inclusión, cuyo instrumento principal es el IMV. La tercera son las migraciones. La agrupación tiene una lógica económica que conviene explicitar. Las tres áreas están ligadas por la sostenibilidad del sistema de pensiones: la inmigración es uno de los determinantes del número de cotizantes y el IMV es una prestación no contributiva del propio sistema. La integración de las Clases Pasivas bajo el mismo ministerio respondió a la voluntad de avanzar hacia una gestión unificada de todas las pensiones públicas, aunque las Clases Pasivas sigan siendo un régimen jurídico y financiero distinto.

La Secretaría de Estado de la Seguridad Social y Pensiones es el órgano superior que dirige y coordina la gestión del sistema. De ella dependen la Dirección General de Ordenación de la Seguridad Social, que prepara la normativa, tutela a las mutuas colaboradoras y elabora las proyecciones del sistema, y las entidades gestoras y servicios comunes adscritos al ministerio. Dos órganos de control y asesoramiento jurídico mantienen, en cambio, una dependencia funcional de otros departamentos, y su situación ilustra la tensión entre integración y control. La Intervención General de la Seguridad Social, que ejerce el control interno de la gestión económica del sistema, depende funcionalmente de la Intervención General de la Administración del Estado, en el Ministerio de Hacienda. El Servicio Jurídico de la Administración de la Seguridad Social asume la representación y defensa en juicio de las entidades del sistema.

La consecuencia más relevante de la organización política actual es la dispersión de las instituciones de protección social entre varios ministerios. Las cuatro entidades gestoras de la Seguridad Social dependen de tres ministerios distintos. El Instituto Nacional de la Seguridad Social y el Instituto Social de la Marina están adscritos al Ministerio de Inclusión. El Instituto de Mayores y Servicios Sociales, al Ministerio de Derechos Sociales. El Instituto Nacional de Gestión Sanitaria, al Ministerio de Sanidad. A ello se añade que la protección por desempleo depende del Ministerio de Trabajo, a través del SEPE, y que la Inspección de Trabajo y Seguridad Social, que controla el cumplimiento de las obligaciones de afiliación y cotización, es un organismo autónomo adscrito también a este último.

Por último, debemos tener en cuenta los costes de la fragmentación, que son costes de coordinación, y se manifiestan con especial claridad en las prestaciones que exigen la intervención de varias instituciones. El caso más ilustrativo es la incapacidad temporal. En ella, la baja médica la emiten los médicos de los servicios públicos de salud de las Comunidades Autónomas; el control lo ejercen los inspectores médicos de esos servicios, los del Instituto Nacional de la Seguridad Social y los de las mutuas; el pago lo realizan las empresas, el INSS o las mutuas según los casos; y la financiación corre a cargo de la Seguridad Social. Ninguna de estas instituciones soporta todos los costes de sus decisiones, lo que genera problemas de incentivos que contribuyen a explicar la evolución reciente del gasto en esta prestación. Otro ejemplo es la relación entre el IMV, gestionado por el INSS, y los subsidios por desempleo, gestionados por el SEPE, que atienden a poblaciones parcialmente coincidentes con reglas y sistemas de información distintos. La literatura sobre gestión pública ha denominado \textit{joined-up government} o \textit{whole-of-government} a los intentos de corregir estos problemas mediante la coordinación horizontal entre departamentos (CHRISTENSEN y LÆGREID, 2007), tras las décadas de fragmentación que produjo la Nueva Gestión Pública (HOOD, 1991).

\subsubsection{Gestión directa: entidades gestoras y servicios comunes}
El texto refundido de la Ley General de la Seguridad Social atribuye la gestión del sistema a las entidades gestoras, que son entidades de derecho público con personalidad jurídica propia, y a los servicios comunes, que prestan servicios a todo el sistema. Ambas categorías están sometidas a la tutela del ministerio al que están adscritas. 

La distinción entre ellas responde a una lógica funcional: las entidades gestoras se especializan por tipo de prestación, y los servicios comunes realizan funciones transversales que todas las prestaciones necesitan, como la recaudación, el pago, la gestión del patrimonio o los sistemas de información.

Esta arquitectura, responde a dos principios organizativos que conviene explicitar. El primero es la especialización funcional. Cada entidad gestora concentra la competencia técnica sobre un tipo de prestación y la ejerce para todos los regímenes, lo que genera economías de escala y de especialización frente al antiguo modelo de gestión por riesgos y colectivos. El segundo es la separación entre la recaudación y el pago de las prestaciones, por un lado, y el reconocimiento del derecho a ellas, por otro. La Tesorería General recauda y paga, pero no decide quién tiene derecho a una pensión; el Instituto Nacional de la Seguridad Social decide quién tiene derecho, pero no recauda ni gestiona los fondos. Es una aplicación del principio clásico de segregación de funciones del control interno, que reduce el riesgo de fraude y de errores al impedir que una misma unidad controle todo el ciclo de una operación. Es también una exigencia de la caja única: si cada entidad gestora recaudara sus propios recursos, la unidad financiera del sistema se rompería.

Cabe mencionar que hay algunas entidades particulares que no responden a la lógica completa del Sistema. Por ejemplo, el INGESA y el IMSERSO son jurídicamente entidades gestoras del sistema de Seguridad Social, aunque su actividad (la asistencia sanitaria residual y los servicios sociales) se financie hoy con impuestos y no con cotizaciones, y aunque estén adscritas a ministerios distintos del de Seguridad Social. El SEPE, en cambio, no es una entidad gestora del sistema, sino un organismo autónomo adscrito al Ministerio de Trabajo, aunque gestione una prestación, el desempleo, que el art. $42$ del texto refundido incluye en la acción protectora. Se trata, por tanto, de un caso en el que el concepto jurídico de acción protectora y el concepto orgánico de sistema no coinciden.

\paragraph{Instituto Nacional de la Seguridad Social (INSS): Gran gestor del Sistema}
El Instituto Nacional de la Seguridad Social (INSS) gestiona y administra las prestaciones económicas del sistema, con excepción de las atribuidas al Instituto de Mayores y Servicios Sociales o a los servicios competentes de las Comunidades Autónomas. Es la entidad gestora central del sistema. 

Reconoce y gestiona las pensiones contributivas de jubilación, incapacidad permanente, viudedad, orfandad y en favor de familiares; la incapacidad temporal cuando no la gestionan las mutuas; las prestaciones de nacimiento y cuidado de menor, de riesgo durante el embarazo y la lactancia y las prestaciones familiares; y, desde 2020, el IMV.\footnote{%
    La asignación del IMV al INSS, y no a las Comunidades Autónomas que gestionaban las rentas mínimas, fue una decisión de gran calado. Aprovechaba la capacidad administrativa y los sistemas de información del INSS, pero sometió a la entidad a una carga de gestión para la que no estaba preparada, con un tipo de prestación (condicionada a la comprobación de recursos del hogar) muy distinto de las prestaciones contributivas que gestionaba tradicionalmente. Las dificultades de gestión de los primeros años, con un volumen muy elevado de solicitudes denegadas o pendientes, y la escasez de personal del INSS, que produjo retrasos notables en la atención presencial, forman parte del debate sobre la baja utilización del IMV. La gestión del IMV en el País Vasco y en Navarra se transfirió a sus respectivas administraciones en virtud de sus regímenes forales, lo que constituye una excepción a la gestión estatal de esta prestación.
} Desde 2020 asume también funciones de gestión de las Clases Pasivas. Sus funciones excluyen expresamente las pensiones no contributivas, que gestionan las Comunidades Autónomas y, en Ceuta y Melilla, el IMSERSO, y la protección por desempleo, que gestiona el SEPE.
 

\paragraph{Tesorería General de la Seguridad Social}
La Tesorería General de la Seguridad Social (TGSS) gestiona los recursos económicos y la administración financiera del sistema. Es el servicio común más importante y la expresión institucional del principio de caja única. Sus funciones abarcan todo el ciclo financiero del sistema. En los ingresos, gestiona la inscripción de empresas y la afiliación, altas y bajas de los trabajadores, y la recaudación de las cotizaciones y demás recursos, tanto en periodo voluntario como en vía ejecutiva, con su propio procedimiento de apremio. 

En los pagos, realiza el pago de las prestaciones que reconocen las entidades gestoras. En la gestión patrimonial, es la titular del patrimonio único de la Seguridad Social, distinto del patrimonio del Estado, y administra el Fondo de Reserva bajo la supervisión de un comité de gestión. Es, además, la contraparte de los préstamos del Estado a la Seguridad Social, lo que la convierte en el deudor formal de la deuda del sistema con el Estado.

La TGSS tiene un papel central en la lucha contra el fraude y la economía sumergida, en colaboración con la Inspección de Trabajo y Seguridad Social, y ha protagonizado una intensa transformación digital en las dos últimas décadas. Desarrolló el Sistema RED para la transmisión electrónica de datos de las empresas, el sistema de liquidación directa de cotizaciones y la sede electrónica para trabajadores y empresas. La liquidación directa, que sustituyó la autoliquidación por parte de las empresas por un cálculo realizado por la propia Tesorería a partir de los datos comunicados, en un ejemplo de cómo la tecnología puede reducir los costes de cumplimiento y el fraude; y muy relevante para la Teoría de la Administración Tributaria, porque desplaza a la Administración la responsabilidad del cálculo de la obligación.

\paragraph{Gerencia de Informática y los demás servicios (GISS)}
La Gerencia de Informática de la Seguridad Social (GISS) tiene competencias relacionadas con el uso y aplicación de las nuevas tecnologías de la información y las comunicaciones en el ámbito de la Seguridad Social. 

Su importancia es mayor de lo que su discreción sugiere. La gestión de un sistema con más de $21$ millones de afiliados y más de $10$ millones de pensiones contributivas depende enteramente de sus sistemas de información. La capacidad del sistema para incorporar reformas complejas, como el cálculo dual de la base reguladora, la cotización de los autónomos por rendimientos reales o el IMV, depende de la capacidad de la GISS para adaptar sus aplicaciones. La cotización de los autónomos por rendimientos reales exige cruzar datos con la Agencia Tributaria para regularizar las cotizaciones provisionales, y su gestión tecnológica es uno de los desafíos más complejos de la reforma de 2022.

\begin{specialblock}{\textbf{EXTRA:} Otras Entidades Gestoras del Sistema de Seguridad Social}
    \textbf{Instituo Social de la Marina}. El Instituto Social de la Marina (ISM) se encarga del sector marítimo-pesquero y de gestionar el Régimen Especial de la Seguridad Social de los Trabajadores del Mar. Es la única entidad gestora organizada por colectivo y no por tipo de prestación, y constituye por eso un vestigio de la antigua gestión por riesgos y colectivos anterior a 1978.

    Su supervivencia se explica por la singularidad del colectivo protegido. Los trabajadores del mar tienen condiciones de trabajo muy específicas: largas estancias en el mar, riesgos laborales elevados y dispersión geográfica. Esas condiciones justifican reglas propias de cotización (con coeficientes reductores de la edad de jubilación y bases de cotización específicas para algunas actividades pesqueras) y servicios específicos. Entre ellos están la sanidad marítima, con la asistencia sanitaria a bordo y en el extranjero, la formación sanitaria de los tripulantes y los centros en el extranjero para marinos españoles. 
    
    El ISM ejerce así funciones que, para el resto de los trabajadores, corresponden a varias instituciones distintas. Ilustra tanto las ventajas de la gestión integrada, que permite ofrecer una atención completa a un colectivo con necesidades especiales, como sus inconvenientes, porque mantiene una estructura administrativa propia para un colectivo relativamente pequeño.

    \textbf{Instituto de Mayores y Servicios Sociales}. El Instituto de Mayores y Servicios Sociales (IMSERSO) tiene competencias sobre los planes, programas y servicios de ámbito estatal para personas mayores y en situación de dependencia, sin perjuicio de la gestión transferida a las Comunidades Autónomas. Es el heredero del Instituto Nacional de Servicios Sociales creado en 1978. Tras la transferencia de los servicios sociales a las Comunidades Autónomas, sus funciones de gestión directa se han reducido a Ceuta y Melilla, donde gestiona las pensiones no contributivas y los servicios sociales. 

    Conserva, en cambio, funciones de ámbito estatal de gran relevancia. La primera es la coordinación del Sistema para la Autonomía y Atención a la Dependencia, incluida la gestión de la financiación estatal del nivel mínimo de protección y el apoyo al Consejo Territorial. La segunda son los programas de turismo y termalismo social, conocidos popularmente como los \textit{viajes del IMSERSO}, que tienen una doble justificación: de política social para las personas mayores y de política económica para mantener la ocupación hotelera en temporada baja. 
    
    El IMSERSO está adscrito al Ministerio de Derechos Sociales, Consumo y Agenda 2030, no al de Seguridad Social, lo que ejemplifica la fragmentación descrita: una entidad gestora del sistema de Seguridad Social depende de un ministerio distinto del que dirige el sistema.

    \textbf{Instituto Nacional de Gestión Sanitaria} (INGS). El Instituto Nacional de Gestión Sanitaria (INGESA) tiene competencias de administración y gestión de los servicios sanitarios. Es el heredero del Instituto Nacional de la Salud (INSALUD), creado en 1978 para gestionar la asistencia sanitaria de la Seguridad Social. Está adscrito al Ministerio de Sanidad.

    Cuando se completaron los traspasos sanitarios a las Comunidades Autónomas en 2002, el INSALUD se transformó en el INGESA, con funciones residuales. Su principal cometido es la gestión de la asistencia sanitaria en Ceuta y Melilla, las únicas partes del territorio en las que la sanidad no está transferida, además de algunos servicios de ámbito estatal. 
    
    Su historia es la historia de la transformación del sistema sanitario español del modelo bismarckiano al beveridgiano. El INSALUD gestionaba la asistencia sanitaria como una prestación de la Seguridad Social financiada con cotizaciones. En cambio, el INGESA gestiona un residuo de esa asistencia en un sistema que ya es un servicio nacional de salud financiado con impuestos y gestionado por las Comunidades Autónomas. Que siga siendo formalmente una entidad gestora de la Seguridad Social es una reliquia institucional de esa transformación.
\end{specialblock}

\paragraph{Agencia Estatal de la Administración de la Seguridad Social: ¿por qué es un proyecto inacabado?}
La fragmentación de la gestión en varias entidades gestoras y servicios comunes con personalidad jurídica propia ha sido objeto de críticas desde hace décadas, y ha generado un proyecto de reforma que nunca se ha culminado. La disposición adicional séptima de la Ley 27/2011 habilitó al Gobierno para crear una Agencia Estatal de la Administración de la Seguridad Social que integrara las funciones de gestión de las entidades gestoras y los servicios comunes. Esa habilitación no se utilizó. El acuerdo social de julio de 2021 que dio lugar a la Ley 21/2021 incluyó el compromiso de presentar un proyecto para crear la agencia en un plazo de $6$ meses. En 2022 se anunció que el Ministerio aceleraba los trabajos, pero la agencia sigue sin crearse. El modelo de referencia es la Agencia Estatal de Administración Tributaria (AEAT), creada en 1991 como entidad de derecho público con autonomía de gestión, que integró las funciones de gestión, inspección y recaudación tributaria y aduanera. Se considera una de las administraciones tributarias más eficientes de Europa. 

Los argumentos a favor de una agencia única de la Seguridad Social son la eliminación de duplicidades entre entidades, la integración de los sistemas de información, la gestión unificada de los recursos humanos, que permitiría redistribuir personal entre funciones según las necesidades, y una mayor autonomía de gestión en materia de contratación y retribuciones. Esa autonomía permitiría afrontar la escasez de personal que ha afectado a la atención a los ciudadanos.\footnote{%
    Los argumentos en contra son de dos tipos. Unos son de economía política: la resistencia de las estructuras existentes, que es lo que predicen los modelos de NISKANEN y DUNLEAVY, y la de los sindicatos de empleados públicos ante los cambios en las condiciones de trabajo. Otros son de fondo: el temor a que la autonomía de gestión reduzca el control político y parlamentario, y la dificultad de integrar entidades adscritas a ministerios distintos.
}

El debate conecta con la literatura sobre la agencificación que siguió a la Nueva Gestión Pública (HOOD, 1991; POLLITT y BOUCKAERT, 2011). Sostiene que las agencias con autonomía de gestión, sometidas a contratos de objetivos y a evaluación de resultados, pueden ser más eficientes que las organizaciones burocráticas tradicionales. Una paradoja del caso español ilustra los vaivenes de esta doctrina. La LRJSP (40/2015) suprimió la figura de las agencias estatales creada en 2006, por considerar que no había producido los resultados esperados, y la Ley de Presupuestos para 2021 la recuperó. Si no se ha producido, no es por un obstáculo legal, sino por falta de voluntad o capacidad política.

\subsubsection{Colaboración en la gestión: mutuas y empresas}
La colaboración de entidades privadas en la gestión de un sistema público es uno de los rasgos más singulares de la arquitectura institucional española. Desde la Ley 35/2014, su denominación legal es mutuas colaboradoras con la Seguridad Social, lo que refleja la ampliación de sus funciones más allá de los accidentes de trabajo y las enfermedades profesionales.

\paragraph{Mutuas colaboradoras con la Seguridad Social: }
Las mutuas son asociaciones privadas de empresarios, sin ánimo de lucro, que se constituyen con autorización del Ministerio de Inclusión y bajo su tutela. Su objeto es colaborar en la gestión de determinadas prestaciones de la Seguridad Social, con responsabilidad mancomunada de los empresarios asociados. Su historia se remonta, como vimos, a la Ley de Accidentes de Trabajo de 1900, que permitió a los empresarios sustituir su responsabilidad objetiva por los accidentes de sus trabajadores mediante un seguro. Durante el franquismo y los primeros años de la democracia llegó a haber varios centenares de mutuas. No obstante, un intenso proceso de concentración, impulsado por la regulación, ha reducido su número a menos de una veintena, con economías de escala considerables.

Sus funciones actuales son cinco. En primer lugar, la gestión de las prestaciones por contingencias profesionales (accidentes de trabajo y enfermedades profesionales), incluida la asistencia sanitaria y la rehabilitación, respecto de los trabajadores de las empresas asociadas que opten por ellas. Segundo, la gestión de la prestación económica por incapacidad temporal derivada de contingencias comunes, si la empresa opta por ello. Tercero, la gestión de la prestación por cese de actividad de los trabajadores autónomos. En cuarto lugar, las prestaciones por riesgo durante el embarazo y la lactancia natural. Y, por último, la prestación por cuidado de menores afectados por cáncer u otra enfermedad grave. 

Para financiarlas, reciben de la Tesorería General las cotizaciones correspondientes a esas contingencias, que siguen siendo recursos públicos de la Seguridad Social y no ingresos propios. Deben aplicarlas exclusivamente a sus fines, y los excedentes que generan se destinan a reservas y, en parte, al Fondo de Contingencias Profesionales de la Seguridad Social. Las mutuas conservan, además, un patrimonio histórico privado, procedente de su actividad anterior a su integración en el sistema, cuya naturaleza ha sido objeto de controversia jurídica.

Ahora, la pregunta que nos puede estar rondando la cabeza es: ¿por qué un sistema público confía a asociaciones privadas de empresarios la gestión de prestaciones públicas? La respuesta tiene fundamento en la teoría de los incentivos. Los empresarios asociados a una mutua son los que soportan el coste de los accidentes de trabajo de sus trabajadores, a través de sus cotizaciones por contingencias profesionales, que son exclusivamente empresariales. Tienen, por tanto, un interés directo en la prevención de los riesgos laborales y en la recuperación rápida de los trabajadores accidentados. Las mutuas pueden canalizar ese interés mediante servicios de prevención, asistencia sanitaria especializada y rehabilitación.

La literatura económica sobre el seguro de accidentes de trabajo y el seguro de desempleo ha estudiado extensamente el papel de la tarificación por experiencia (experience rating), que hace depender la cotización de cada empresa de su siniestralidad o de sus despidos pasados. FELDSTEIN (1976) y TOPEL (1983) mostraron que, en el seguro de desempleo estadounidense, una tarificación por experiencia incompleta subvenciona a las empresas que despiden más y aumenta el desempleo temporal. La lógica es la misma en los accidentes de trabajo: si todas las empresas de un sector pagan la misma cotización con independencia de su siniestralidad, las empresas más seguras subvencionan a las menos seguras, y el incentivo a invertir en prevención se debilita. España aplica una forma limitada de tarificación por experiencia mediante un sistema de incentivos a las empresas que reducen su siniestralidad (conocido como \textit{bonus}), regulado desde 2017, que permite recuperar una parte de las cotizaciones por contingencias profesionales. Pero no existe un \textit{malus} equivalente que penalice a las empresas con mayor siniestralidad.

La colaboración de las mutuas plantea, sin embargo, un conflicto de intereses que es el eje de todos los debates sobre su papel. Como asociaciones de empresarios, las mutuas tienen incentivos para reducir el coste de las prestaciones que gestionan. Ese incentivo es socialmente deseable cuando se traduce en más prevención y en una recuperación más rápida, pero no lo es cuando se traduce en presionar para que los trabajadores vuelvan al trabajo antes de estar recuperados, o en calificar como comunes contingencias que son profesionales para trasladar su coste al sistema general. Es un problema de agencia en el sentido de la teoría económica: el principal (el sistema público y los trabajadores protegidos) delega en un agente (la mutua) cuyos intereses no coinciden plenamente con los suyos. La regulación intenta alinearlos mediante la tutela ministerial, el control de las altas médicas y la prohibición del ánimo de lucro.

El debate es hoy especialmente vivo en relación con la incapacidad temporal por contingencias comunes. El fuerte crecimiento de su gasto en los últimos años ha llevado a proponer que las mutuas tengan un papel mayor en el control y la gestión de las bajas, en particular en los procesos de origen traumatológico, que ellas podrían tratar con sus propios medios sanitarios. En mayo de 2026, la Comunidad Autónoma de Galicia, que registra una de las tasas de absentismo más elevadas de España, modificó su instrucción sobre gestión de la incapacidad temporal. Desde entonces, las propuestas de alta formuladas por las mutuas deben ser aceptadas por la inspección médica autonómica si están debidamente motivadas y completas, y el alta puede producirse si la inspección no responde. Los sindicatos denunciaron la medida como una privatización encubierta de la gestión sanitaria. Sus defensores la presentaron como una forma de reducir los tiempos de espera y el coste de las bajas prolongadas. Otras propuestas sindicales, como la de que las mutuas traten las lesiones físicas con cargo a la Seguridad Social, muestran que el debate no enfrenta simplemente a partidarios y adversarios de las mutuas, sino que versa sobre cómo repartir entre instituciones públicas y privadas la gestión de una prestación cuyo coste se ha disparado.

\paragraph{Empresas: ¿qué es la colaboración voluntaria?}
Junto a las mutuas, el sistema admite la colaboración voluntaria de las empresas en la gestión de determinadas prestaciones, principalmente la incapacidad temporal por contingencias profesionales y comunes, así como la colaboración obligatoria que consiste en el pago delegado de la prestación de incapacidad temporal. En virtud de esta, la empresa abona la prestación al trabajador junto con su salario y la deduce después de sus cotizaciones. El pago delegado es un mecanismo de eficiencia administrativa, porque aprovecha la relación de pago mensual entre la empresa y el trabajador para evitar un circuito de pago separado. Pero implica que las empresas anticipen recursos que después compensan. La colaboración voluntaria, que permitía a las grandes empresas gestionar directamente determinadas prestaciones, ha ido reduciéndose, y su régimen se ha ido restringiendo en los últimos años.

\subsubsection{Fuera del sistema, dentro de la protección}
Varias instituciones que forman parte de la protección social en sentido funcional no pertenecen al sistema de Seguridad Social en sentido orgánico. Su examen completa el mapa institucional.

\paragraph{Servicio Público de Empleo Estatal y el Fondo de Garantía Salarial}
El Servicio Público de Empleo Estatal (SEPE) tiene atribuida la gestión y control de las prestaciones por desempleo y de las políticas de formación y colocación en los territorios que no las tengan transferidas a las Comunidades Autónomas. Es un organismo autónomo adscrito al Ministerio de Trabajo y Economía Social, heredero del Instituto Nacional de Empleo creado en 1978. 

Su posición es singular en la arquitectura del sistema. Gestiona una prestación que forma parte de la acción protectora pero no es una entidad gestora del sistema, sino un organismo del Sistema Nacional de Empleo. Por eso se integra, junto con el Sistema de la Seguridad Social, en el subsector de Administraciones de Seguridad Social de la contabilidad nacional. Esta separación orgánica entre las prestaciones por desempleo y el resto de las prestaciones económicas responde a la vinculación del desempleo con las políticas activas de empleo, que las Comunidades Autónomas gestionan desde los traspasos de finales de los noventa y comienzos de la década de 2000. La Ley 3/2023 de Empleo previó su transformación en la Agencia Española de Empleo, que en 2026 sigue sin constituirse. 

El Fondo de Garantía Salarial (FOGASA), también organismo autónomo adscrito al Ministerio de Trabajo, garantiza a los trabajadores el cobro de los salarios y de las indemnizaciones por despido que las empresas no pueden pagar por insolvencia o concurso, dentro de ciertos límites. Se financia con una cotización a cargo exclusivo de los empresarios. Su lógica económica es la de un seguro colectivo frente al riesgo de insolvencia empresarial. Ese riesgo no es asegurable individualmente por el trabajador, porque su información sobre la solvencia de su empleador es limitada, y está correlacionado con el ciclo económico, de modo que las insolvencias se concentran en las recesiones. Es, por tanto, un riesgo sistemático que exige un mecanismo colectivo.

\paragraph{Mutualismo administrativo}
Las mutualidades de funcionarios (MUFACE, ISFAS y MUGEJU) responden a una lógica histórica singular dentro del sistema español de protección social. Preservan un régimen de mutualismo administrativo anterior a la generalización de la Seguridad Social contributiva ordinaria y permiten a cada mutualista optar cada año entre recibir asistencia sanitaria a través de la sanidad pública ordinaria o a través de una entidad de seguro privada concertada y financiada con cargo a la mutualidad. 

Es un sistema dual que ha sobrevivido, desde su origen preconstitucional, a la creación del sistema unificado de Seguridad Social. La Mutualidad General de Funcionarios Civiles del Estado (MUFACE), el Instituto Social de las Fuerzas Armadas (ISFAS) y la Mutualidad General Judicial (MUGEJU) son organismos autónomos que gestionan el régimen especial de Seguridad Social de los funcionarios civiles del Estado, de los militares y del personal de la Administración de Justicia, respectivamente. Su acción protectora cubre la asistencia sanitaria, la incapacidad temporal y algunas prestaciones sociales, pero no las pensiones, que para los funcionarios anteriores a 2011 corresponden al régimen de Clases Pasivas.

El rasgo más característico del mutualismo administrativo es el derecho de opción anual entre la sanidad pública y la privada concertada. Tradicionalmente, una amplia mayoría de los mutualistas, en torno a dos tercios, ha elegido la asistencia privada, lo que ha convertido al mutualismo en un modelo de competencia gestionada único en el sistema sanitario español. Es un experimento natural que los economistas de la salud han utilizado para comparar la eficiencia de la provisión pública y la privada. Sus defensores sostienen que ofrece a los funcionarios libertad de elección y un acceso más rápido a los especialistas, y que su coste por persona protegida es inferior al del Sistema Nacional de Salud. Sus críticos responden que esa comparación es engañosa. Los mutualistas que eligen la asistencia privada son, en promedio, más jóvenes y más sanos que los que eligen la pública; lo que lo convierte en un problema de selección de riesgos. Además, las aseguradoras privadas derivan a la sanidad pública los tratamientos más costosos, como los oncológicos o los de alta complejidad, y la duplicidad de sistemas genera un trato privilegiado a un colectivo profesional, contrario al principio de igualdad.

El debate tuvo un episodio de gran relevancia. En otoño de 2024, la licitación del concierto sanitario de MUFACE para 2025-2027 quedó desierta, porque ninguna aseguradora aceptó las condiciones económicas ofrecidas: alegaban que las primas de los años anteriores no cubrían el coste de atender a una población mutualista envejecida. Tras cinco meses de incertidumbre, dos licitaciones fallidas y dos prórrogas del concierto anterior, en marzo de 2025 el concierto se adjudicó a solo dos aseguradoras, Adeslas y Asisa. El incremento de las primas alcanzó el $41,2\%$ y el valor del contrato, unos $4.800$ millones de euros, para más de un millón y medio de mutualistas y beneficiarios. Otras aseguradoras que habían participado en conciertos anteriores declinaron presentarse. El episodio ilustra con precisión varios conceptos de este tema. El primero es la selección adversa en el mercado de seguros de salud: el envejecimiento de la población mutualista y la salida de los más sanos hacia la sanidad pública encarecieron el coste medio. El segundo es el poder de mercado de un pequeño número de aseguradoras frente a un comprador público. El tercero es la dificultad política de reformar un privilegio institucional con una coalición defensora poderosa: el Gobierno, que había anunciado estudios sobre la posible integración de los mutualistas en el Sistema Nacional de Salud, acabó renovando el modelo con un incremento sustancial de su coste.

\paragraph{Clases Pasivas}
El régimen de Clases Pasivas del Estado, protege a los funcionarios civiles y militares del Estado que ingresaron antes del 1 de enero de 2011, frente a la jubilación, la incapacidad y la muerte y supervivencia. Se financia directamente con cargo al presupuesto del Estado, en una sección presupuestaria propia, y no con cotizaciones de la Seguridad Social, aunque sus beneficiarios realizan aportaciones. Su gestión, tradicionalmente en manos del Ministerio de Hacienda, se trasladó en 2020 al Ministerio de Inclusión, Seguridad Social y Migraciones.

Es un régimen cerrado y en extinción: desde 2011 no se incorporan nuevos funcionarios, de modo que su número de pensionistas alcanzará un máximo en las próximas décadas y después decrecerá hasta desaparecer. Su integración progresiva en el Régimen General es un ejemplo de convergencia institucional que se ha hecho por la vía del cierre del régimen antiguo y no por la de su reforma, lo que evita el coste político de modificar derechos existentes.

\paragraph{Servicios de salud y los servicios sociales autonómicos}
Por último, la mayor parte del gasto en protección social en sentido funcional que no corresponde a pensiones y desempleo se gestiona por las Comunidades Autónomas. 

Sus servicios de salud gestionan la asistencia sanitaria, y sus servicios sociales gestionan la dependencia, las rentas mínimas y los servicios sociales en general. Su organización, su financiación a través del sistema de financiación autonómica y su relación con el Estado, a través del Consejo Interterritorial del Sistema Nacional de Salud y del Consejo Territorial de Servicios Sociales y del Sistema para la Autonomía y Atención a la Dependencia, se examinan más adelante. 

\begin{specialblock}{\textbf{AMPLIACIÓN:} Control, evaluación y participación}
    El control del sistema se ejerce en cuatro planos: el control interno de la legalidad y la gestión económica, el control externo de las cuentas, el control político parlamentario y la evaluación independiente de su sostenibilidad. A ellos se añade la participación de los agentes sociales, exigida por la Constitución.

    \textbf{control interno y el control externo}. La Intervención General de la Seguridad Social ejerce el control interno de la gestión económico-financiera de las entidades del sistema. Lo hace mediante la función interventora previa, el control financiero permanente y la auditoría pública. Elabora además la contabilidad del sistema y su cuenta general. Como se dijo, depende funcionalmente de la Intervención General de la Administración del Estado. El Tribunal de Cuentas, órgano constitucional de control externo, fiscaliza la cuenta general de la Seguridad Social y realiza fiscalizaciones específicas sobre aspectos de su gestión. En sus informes de los últimos años ha llamado la atención sobre el deterioro del patrimonio neto del sistema, que ha pasado a ser negativo como consecuencia de la acumulación de déficits financiados con préstamos del Estado. 

    \textbf{Control parlamentario: la Comisión del Pacto de Toledo}. El control político del sistema corresponde a las Cortes Generales. Se ejerce mediante la aprobación anual del presupuesto de la Seguridad Social dentro de los PGE, la convalidación de los decretos-leyes y la actividad de la Comisión de Seguimiento y Evaluación de los Acuerdos del Pacto de Toledo del Congreso de los Diputados. El seguimiento de las recomendaciones del Pacto se realiza en esa comisión. Su función no es solo de control, sino de elaboración de consensos. Es el foro en el que los grupos parlamentarios negocian las recomendaciones que después orientan las reformas. Como se vio, su capacidad para cumplir esa función depende de la existencia de un consenso político que se ha debilitado con la fragmentación parlamentaria. La reforma de 2023 le atribuyó además un papel en la cláusula de salvaguarda del MEI: si el gasto en pensiones supera el umbral de referencia, el Gobierno debe elevar a la Comisión una propuesta de medidas correctoras.

    \textbf{Evaluación independiente: la AiREF}. La Autoridad Independiente de Responsabilidad Fiscal (AIReF), es la institución fiscal independiente de España cuya función es velar por el cumplimiento de los principios de estabilidad presupuestaria y sostenibilidad financiera.\footnote{%
        La justificación teórica de las instituciones fiscales independientes procede de la literatura sobre la inconsistencia temporal de la política económica. En pensiones, el sesgo presente es especialmente intenso porque los beneficiarios votan y los que soportarán el coste, en buena medida, todavía no. CALMFORS y WREN-LEWIS (2011) sistematizaron las funciones de los consejos fiscales, que pueden contrarrestar el sesgo presente aportando información independiente, elaborando previsiones no sesgadas y evaluando públicamente el cumplimiento de las reglas, sin sustituir la decisión política, sino elevando su coste reputacional.
    } En materia de pensiones ha adquirido un papel central por tres vías. La primera son sus proyecciones de largo plazo del gasto en pensiones, con una metodología propia y alternativa a la del Ministerio y a la del Ageing Report. La segunda es la evaluación de la regla de gasto en pensiones que le atribuyó la reforma de 2023: cada tres años debe calcular el impacto de las medidas de ingresos adoptadas y compararlo con el gasto proyectado. La tercera es la evaluación del IMV, que la Ley 19/2021 le encomendó y que ha documentado sus problemas de cobertura y de baja utilización. A ellas se suman los ejercicios de revisión del gasto (spending reviews) que realizó entre 2018 y 2021 por encargo del Gobierno, algunos de ellos sobre prestaciones del sistema.
    
    El papel de la AIReF en las pensiones no está exento de tensiones. En 2025 se produjeron discrepancias públicas entre la AIReF y el Ministerio sobre las proyecciones de gasto y de ingresos del sistema. El Ministerio presentó sus propias estimaciones, elaboradas con un modelo propio, que resultaban más favorables. Esto ilustra una cuestión de fondo: la evaluación de la sostenibilidad del sistema depende de supuestos inciertos sobre la demografía, el empleo, la productividad y la inmigración, y la elección de esos supuestos no es políticamente neutral. Plantean también la cuestión de la legitimidad democrática de la delegación en instituciones técnicas. Una norma que hace depender la activación automática de subidas de cotizaciones de las proyecciones de la AIReF y del Ageing Report europeo traslada una decisión de gran relevancia distributiva a instancias no elegidas, un debate que la literatura sobre el Estado regulador y la despolitización ha planteado con carácter general.

    \textbf{Participación de los agentes sociales}. El art. $129.1$ de la Constitución, exige la participación de los interesados en la Seguridad Social.\footnote{%
        La participación de los agentes sociales tiene una justificación de legitimidad, porque son los representantes de quienes financian el sistema con sus cotizaciones, y una justificación de eficacia, porque los acuerdos negociados con ellos tienden a ser más estables. Pero plantea también la cuestión de la representación de los ausentes. Los sindicatos y las organizaciones empresariales representan principalmente a los trabajadores y a las empresas actuales. No representan, o representan de forma indirecta, a los pensionistas, a los desempleados de larga duración, a los trabajadores informales y, sobre todo, a las generaciones futuras, que soportarán el coste de las decisiones que hoy se adoptan. Es la versión institucional del problema de la dualización y del sesgo intergeneracional que justifica las instituciones fiscales independientes.
    } Se concreta en la presencia de las organizaciones sindicales y empresariales más representativas en los órganos de control de las entidades gestoras: los Consejos Generales y las Comisiones Ejecutivas, de composición tripartita con la Administración, del INSS y de las demás entidades. Estos órganos tienen funciones de control y seguimiento de la gestión y de elaboración de criterios de actuación, pero su influencia real en la gestión cotidiana es limitada. La participación de los agentes sociales en el sistema se ejerce mucho más a través del diálogo social, es decir, de la negociación de acuerdos entre el Gobierno, los sindicatos y la patronal que después se trasladan a la legislación. Como vimos, ha sido el cauce de la mayoría de las reformas del sistema desde 1996. El CES, órgano consultivo del Gobierno con participación de los agentes sociales, emite dictámenes sobre los anteproyectos de ley en materia socioeconómica, incluidas las reformas de la Seguridad Social. 
\end{specialblock}



\clearpage
\section{La Financiación del sistema}

Cada euro de pensión que el sistema promete es un euro que alguien, hoy o en el futuro, tendrá que aportar. La forma en que se organiza esa aportación determina quién soporta el coste, cuándo lo soporta, con qué riesgos y con qué efectos sobre el empleo, el ahorro y el crecimiento. Por eso la teoría de la hacienda pública trata la financiación de la Seguridad Social desde dos perspectivas complementarias, que corresponden a dos dimensiones del problema. La primera es temporal: en qué momento se acumulan los recursos que financian una pensión, si mientras el pensionista trabaja, como en la capitalización, o mientras cobra, como en el reparto. La segunda es instrumental: con qué figura se obtienen esos recursos, si con una cotización ligada al salario y a la prestación futura o con impuestos generales desligados de ella.

El bloque se organiza en tres apartados. El primero examina la teoría de la financiación, con esas dos dimensiones como subapartados: reparto frente a capitalización (2.1.1) y cotizaciones frente a impuestos (2.1.2). El segundo (2.2) describe las fuentes vigentes en 2026: las cotizaciones, las transferencias del Estado y los mecanismos de endeudamiento y de reserva, incluido el relato completo del Fondo de Reserva que el documento pedía reubicar. El tercero (2.3) analiza el presupuesto del sistema en términos consolidados por subsectores, como pedía tu nota, y lo sitúa en la comparación internacional del gasto social. 

La tesis del bloque, que se irá desarrollando, es que la financiación del sistema español ha evolucionado en los últimos quince años desde una lógica bismarckiana relativamente pura hacia un modelo híbrido. En él, la cotización se parece cada vez más a un impuesto y el Estado asume una parte creciente del coste, tanto por transferencias como por préstamos, sin que ese desplazamiento se haya planteado nunca abiertamente como una decisión de política económica.

\subsection{Teoría de la financiación: cuándo y con qué se pagan las pensiones}
Cualquier sistema de pensiones responde a un problema económico elemental que conviene formular con precisión antes de entrar en los debates. Los jubilados consumen bienes y servicios que no producen, porque ya no trabajan. Esos bienes y servicios tienen que ser producidos por los trabajadores en activo en ese mismo momento, porque la mayor parte del consumo no puede almacenarse durante décadas. BARR (2002) lo formuló en una frase que se ha convertido en punto de partida: \textbf{a los pensionistas no les interesa el dinero, sino el consumo}.

Lo que un sistema de pensiones hace, sea cual sea su forma, es organizar los derechos de los jubilados sobre la producción futura. El reparto lo hace mediante una promesa del Estado, que obliga por ley a los trabajadores futuros a ceder parte de su producción. La capitalización lo hace mediante la acumulación de activos financieros que los jubilados venden, o de los que cobran rendimientos, a los trabajadores futuros a cambio de esa misma producción. En ambos casos, el consumo de los jubilados se financia con la producción de los activos del momento. La diferencia está en el título jurídico que da derecho a ese consumo, en los riesgos que soporta ese título y en los efectos que cada forma de organización tiene sobre el ahorro, la inversión y el crecimiento, que pueden alterar el tamaño de la producción futura.

Esta observación, que a veces se denomina el principio de que \textit{la producción es lo central}, no zanja el debate, pero lo ordena. Descarta de entrada un argumento muy extendido en el debate público: que la capitalización resuelve el problema demográfico porque cada generación se paga su propia pensión. Si la población activa se reduce en relación con la jubilada, habrá menos producción por jubilado con cualquier sistema, y la única cuestión es cómo se reparte esa escasez. Pero deja abiertas dos cuestiones que sí son decisivas. La primera es si la forma de financiación afecta al tamaño de la producción futura, a través del ahorro y la acumulación de capital. La segunda es cómo reparte los riesgos entre generaciones y entre individuos. A estas cuestiones nos dedicamos primero. Después examinaremos la dimensión instrumental: la naturaleza, la incidencia y los efectos de la cotización frente al impuesto.

\subsubsection{Reparto frente a capitalización}
Conviene empezar por las definiciones. Desde sus inicios, España, como la inmensa mayoría de los países de la Europa continental, financia sus pensiones mediante un sistema de reparto (\textit{pay-as-you-go}): las cotizaciones de los trabajadores en activo financian hoy, directamente y sin acumulación intermedia, las pensiones de los jubilados de hoy, en un pacto de solidaridad intergeneracional continuo. No existe, salvo la excepción parcial del Fondo de Reserva, una cuenta individual acumulada a nombre de cada cotizante. El sistema nunca ha abandonado el reparto como columna vertebral, pero ha incorporado un elemento parcial de acumulación de reservas que, sin transformar su naturaleza, pretendía amortiguar sus fluctuaciones cíclicas más agudas. Por el contrario, el sistema alternativo de capitalización individual será aquel en el que cada trabajador acumula una cuenta de ahorro personal durante su vida activa, de la que dispone, con su propio rendimiento financiero, al jubilarse.

No obstante, estas dos definiciones mezclan dos dimensiones distintas que la literatura separa cuidadosamente. La primera es el \textbf{grado de capitalización}. Es decir, si las prestaciones se pagan con las aportaciones corrientes (reparto puro), con un fondo acumulado que cubre la totalidad de los derechos devengados (capitalización completa) o con una combinación de ambos (capitalización parcial). La segunda es la \textbf{fórmula de la prestación}. En la prestación definida, la pensión se calcula con una fórmula que depende del salario y de los años cotizados, y el riesgo de que las aportaciones no basten lo soporta el promotor del sistema. En la aportación definida, la pensión depende de lo acumulado y de su rentabilidad, y el riesgo lo soporta el individuo.

Las dos dimensiones pueden combinarse de cuatro formas. El sistema español es de reparto con prestación definida. Los planes de pensiones de empleo tradicionales de las grandes empresas británicas o estadounidenses eran de capitalización con prestación definida. El sistema chileno de 1981 es de capitalización con aportación definida. Y las cuentas nocionales, son de reparto con aportación definida: no acumulan activos, pero calculan la pensión como si lo hicieran. A ello se añaden los fondos de reserva o fondos colchón (buffer funds), que acumulan activos dentro de un sistema de reparto sin asignarlos a cuentas individuales. Ejemplos son los fondos AP suecos, el fondo de reserva del Canada Pension Plan, que desde la reforma de 1997 financia una parte significativa de las prestaciones con los rendimientos de su cartera, y el propio Fondo de Reserva español. El caso canadiense es especialmente ilustrativo. En 1997, ante las proyecciones de déficit, Canadá elevó rápidamente la cotización hasta un nivel estable, en torno al $9,9\%$ del salario, que se preveía suficiente para mantener el sistema a largo plazo gracias a los rendimientos de un fondo creciente gestionado de forma independiente. Es la llamada capitalización de estado estacionario (steady-state funding), una vía intermedia que España no ha explorado.

#### El modelo de Samuelson: el reparto como "contrato social"

La justificación teórica moderna del reparto procede de un art. de Samuelson (1958) que no trataba directamente de pensiones, sino del tipo de interés en una economía de generaciones solapadas. Samuelson imaginó una economía en la que cada individuo vive dos periodos, trabaja en el primero y está jubilado en el segundo, y en la que los bienes no pueden almacenarse. En esa economía no hay forma privada de ahorrar para la vejez. Los jóvenes no pueden prestar a los viejos, porque estos no vivirán para devolver el préstamo, y no hay capital en el que invertir. Sin intervención, cada generación consumiría todo lo que produce en su juventud y nada en su vejez. Samuelson mostró que un artificio social, una convención por la que cada generación de jóvenes cede parte de su producción a los viejos a cambio de la promesa de que la siguiente generación hará lo mismo por ellos, mejora el bienestar de todas las generaciones: la solidaridad intergeneracional que el documento describe.

La clave del modelo es la rentabilidad implícita del reparto. Si la población crece a una tasa n y la productividad a una tasa g, la masa salarial crece aproximadamente a la tasa $n + g$. Con un tipo de cotización constante, cada generación recibe en su vejez una pensión que crece con la masa salarial de la generación siguiente. Por tanto, la rentabilidad que obtiene de sus cotizaciones es aproximadamente $n + g$, lo que SAMUELSON denominó el tipo de interés biológico. AARON (1966) formuló la condición que lleva su nombre: el reparto es preferible a la capitalización, en términos de bienestar de todas las generaciones, si la tasa de crecimiento de la masa salarial $(n + g)$ es superior al tipo de interés real ($r$) que rendiría un sistema de capitalización.

Un ejemplo numérico ayuda a fijar la intuición. Supongamos que el tipo de cotización es el veinte por ciento del salario, que cada generación trabaja cuarenta años y que la masa salarial crece a una tasa anual del 2,5 por ciento, porque la población activa crece un 1 por ciento y la productividad un 1,5 por ciento. Una generación que cotiza durante toda su vida recibe una pensión cuya cuantía crece con la masa salarial de los activos que la financian: su rentabilidad implícita es del 2,5 por ciento anual. Si la alternativa fuera invertir esas mismas cotizaciones en un fondo con una rentabilidad real del 4 por ciento, al cabo de cuarenta años cada euro aportado al principio de la vida laboral valdría 1,04⁴⁰ ≈ 4,8 euros en capitalización y 1,025⁴⁰ ≈ 2,7 euros en reparto. La capitalización ofrecería, en esta comparación simplificada, casi el doble. En cambio, en la situación de los años cincuenta y sesenta, con un crecimiento de la población activa y de la productividad muy superior, la condición de AARON se cumplía holgadamente y el reparto era la mejor opción. En los términos del modelo, la historia de los sistemas de pensiones europeos es la historia de una condición de AARON que se cumplía en la edad de oro y que dejó de cumplirse con la caída de la natalidad y la desaceleración de la productividad desde los años setenta.

El propio SAMUELSON escribió en 1967, en su columna de la revista Newsweek, que la belleza de la seguridad social consistía en que era actuarialmente insostenible, porque cada jubilado recibía más de lo que había aportado, y que una nación en crecimiento era el mayor esquema piramidal (Ponzi game) jamás ideado. La frase, escrita con optimismo en plena edad de oro, ha sido citada después con ironía por los críticos del reparto, porque un esquema piramidal solo funciona mientras crece la base. Pero la analogía es engañosa: a diferencia de un esquema piramidal, el reparto no depende de que la base crezca indefinidamente, sino de que no se reduzca demasiado deprisa. Si la población y la productividad crecen a una tasa constante aunque sea baja, el sistema es sostenible con esa rentabilidad baja. El problema no es el reparto en sí, sino la transición demográfica: el paso de un crecimiento alto a uno bajo o negativo, con una generación numerosa que se jubila cuando las siguientes son mucho menos numerosas.

#### La ganancia de lanzamiento y la deuda implícita

El modelo de SAMUELSON contiene un elemento que es la clave de todo el debate sobre la transición entre sistemas. Cuando se crea un sistema de reparto, la primera generación de jubilados recibe pensiones sin haber cotizado, o habiéndolo hecho solo durante una parte de su vida. Es la ganancia de lanzamiento (windfall gain). Esa ganancia se financia con las cotizaciones de las generaciones siguientes y no se recupera nunca: mientras el sistema exista, cada generación paga la pensión de la anterior y recibe la suya de la siguiente, pero la primera recibió sin pagar. La consecuencia es que un sistema de reparto maduro lleva incorporada una deuda implícita, que es exactamente el valor actual de las pensiones que se deben a los jubilados y trabajadores actuales por las cotizaciones que ya han pagado y que no se han acumulado en ningún fondo.

BREYER (1989) y, de forma más completa, SINN (2000) extrajeron de esta observación una conclusión que desmonta el argumento más popular a favor de la capitalización. Aunque la rentabilidad de la capitalización ($r$) sea superior a la del reparto ($n + g$), la transición del reparto a la capitalización no puede mejorar a todas las generaciones. La generación que realiza la transición tiene que pagar dos veces: debe seguir pagando las pensiones de los jubilados actuales, a quienes no se puede dejar sin pensión, y a la vez acumular su propio fondo. Si esa doble carga se reparte entre las generaciones mediante deuda pública, las generaciones futuras deberán pagar los intereses de esa deuda. SINN demostró que el valor actual de esos intereses es exactamente igual a la ganancia de rentabilidad que obtienen de la capitalización. La diferencia entre $r$ y $n + g$ no es una ineficiencia del reparto que la capitalización pueda eliminar, sino un impuesto implícito que cada generación paga para servir la deuda implícita heredada de la primera generación. Cambiar de sistema no elimina ese impuesto, solo lo convierte en explícito.

GEANAKOPLOS, MITCHELL y ZELDES (1998) llegaron a la misma conclusión con datos estadounidenses. Su análisis mostró que la comparación entre la rentabilidad de un sistema de capitalización y la del reparto es engañosa si no se incluye el coste de la deuda implícita heredada, y que, una vez incluido, la ventaja de la capitalización desaparece. Es el segundo de los diez mitos de ORSZAG y STIGLITZ (1999) presentados: la mayor rentabilidad de las cuentas individuales es, en buena medida, una ilusión contable que ignora el coste de la transición.

El Sistema Europeo de Cuentas de 2010 obliga desde 2017 a los Estados miembros a publicar una tabla complementaria (tabla 29) con los derechos de pensión devengados. Es una estimación del valor actual de las pensiones que los sistemas públicos deben a los jubilados y trabajadores actuales por las cotizaciones ya pagadas. Para los sistemas de reparto maduros de Europa continental, esa deuda implícita representa varias veces el PIB, muy por encima de la deuda pública explícita. La cifra no significa que deba esa cantidad en el sentido en que debe su deuda pública, porque el legislador puede modificar las prestaciones futuras, pero mide la magnitud del compromiso que el sistema ha asumido y el coste que tendría una transición completa a la capitalización.

#### La eficiencia dinámica: cuándo el reparto reduce el bienestar

La condición de AARON tiene una contrapartida teórica que conviene conocer, porque determina si la comparación entre r y n + g favorece a uno u otro sistema. DIAMOND (1965) integró el modelo de Samuelson en un modelo de crecimiento con acumulación de capital y mostró que la economía puede estar en una situación de ineficiencia dinámica, con acumulación excesiva de capital, si el tipo de interés es inferior a la tasa de crecimiento ($r < n + g$). En ese caso, la economía ha ahorrado demasiado, y un sistema de reparto que reduzca el ahorro mejora el bienestar de todas las generaciones, porque libera recursos que la acumulación excesiva desperdiciaba. Es la condición de AARON en otro lenguaje. Si, por el contrario, $r > n + g$, la economía es dinámicamente eficiente. Un sistema de reparto que reduzca el ahorro reduce el capital, la producción y el bienestar de las generaciones futuras en el estado estacionario, aunque haya beneficiado a las primeras generaciones.

La evidencia empírica sobre esta cuestión era relativamente sólida. ABEL, MANKIW, SUMMERS y ZECKHAUSER (1989) mostraron, con un criterio que compara los flujos de beneficios del capital con la inversión, que las principales economías avanzadas eran dinámicamente eficientes. Sin embargo, el largo periodo de tipos de interés reales muy bajos, incluso negativos, entre la crisis financiera de 2008 y el repunte de la inflación de 2022 reabrió el debate. BLANCHARD (2019), en su conferencia presidencial ante la American Economic Association, argumentó que, si el tipo de interés seguro se mantiene por debajo de la tasa de crecimiento, el coste fiscal de la deuda pública es mucho menor de lo que se pensaba. Por la misma lógica, la ventaja relativa de la capitalización frente al reparto se reduce o incluso se invierte. La subida de los tipos de interés desde 2022 ha vuelto a situar la comparación en términos más favorables a la capitalización, lo que muestra que la condición de AARON no es un dato estructural fijo, sino una comparación entre dos tasas que fluctúan con el ciclo económico y financiero.

Efecto sobre el ahorro: Feldstein, Barro y un error informático

El argumento más influyente contra el reparto no se refiere a su rentabilidad, sino a su efecto sobre el ahorro nacional. FELDSTEIN (1974) sostuvo que un sistema de reparto reduce el ahorro privado por un efecto de sustitución de riqueza: los trabajadores que esperan recibir una pensión pública necesitan ahorrar menos por su cuenta para la vejez, y esa reducción del ahorro privado no se compensa con ningún ahorro público, porque las cotizaciones se gastan en pagar las pensiones corrientes. El resultado es una menor acumulación de capital, una menor productividad y una menor producción futura, precisamente lo que el principio de que la producción es lo central considera decisivo.\footnote{%
    FELDSTEIN añadía un efecto contrario, el de jubilación inducida: la pensión pública anima a jubilarse antes, lo que alarga el periodo de jubilación y aumenta la necesidad de ahorro. Pero estimó que el primer efecto dominaba. Con series temporales estadounidenses, concluyó que la seguridad social había reducido el ahorro personal en torno a la mitad.

No obstante, en 1982, LEIMER y LESNOY, dos economistas de la Administración de la Seguridad Social estadounidense, descubrieron que las estimaciones de FELDSTEIN contenían un error de programación informática en el cálculo de la variable de riqueza de la seguridad social. Corregido el error, el efecto negativo sobre el ahorro desaparecía o se reducía drásticamente. FELDSTEIN reconoció el error y presentó nuevas estimaciones que mantenían un efecto negativo, pero la controversia dañó la credibilidad del resultado y abrió una larga discusión econométrica sin conclusión definitiva.
}

La crítica teórica procedía de BARRO (1974), con su teorema de la equivalencia ricardiana. Si las generaciones están unidas por el altruismo, de modo que los padres se preocupan por el bienestar de sus hijos y les dejan herencias, un sistema de reparto que transfiere recursos de los jóvenes a los viejos no tiene efectos reales. Los viejos, que saben que la pensión que reciben la pagan sus hijos, aumentan en la misma cuantía las herencias que les dejan, y el ahorro nacional no cambia. La evidencia empírica no respalda la equivalencia ricardiana en su forma pura, porque las herencias no parecen responder de ese modo a las transferencias públicas. Pero el argumento de BARRO mostró que el efecto del reparto sobre el ahorro depende de supuestos sobre el comportamiento de las familias que no pueden darse por descontados.

El resultado de este debate es que el efecto del reparto sobre el ahorro nacional es teóricamente ambiguo y empíricamente incierto, aunque la mayoría de los estudios encuentra un efecto negativo moderado. Y aun si existiera, la transición a la capitalización solo aumentaría el ahorro nacional si se financiara con impuestos o reduciendo gasto, no con deuda: es el primer mito de ORSZAG y STIGLITZ. Una transición financiada con deuda aumenta el ahorro privado acumulado en los fondos, pero reduce en la misma cuantía el ahorro público, de modo que el ahorro nacional no cambia.

Capitalización ante la demografía: el \textit{desplome de los activos}

El segundo argumento popular a favor de la capitalización es que protege frente al envejecimiento. La respuesta de la literatura es que la capitalización no escapa a la demografía, sino que la sufre por un canal distinto: el precio de los activos. MANKIW y WEIL (1989) predijeron que la jubilación de la generación del baby boom provocaría una caída de los precios de la vivienda estadounidense, porque la demanda de vivienda depende de la estructura por edades de la población. La predicción no se cumplió en los términos previstos, pero abrió la línea de investigación sobre la hipótesis del desplome de los activos (asset meltdown). Según ella, cuando una generación numerosa se jubila y vende sus activos para financiar su consumo a una generación menos numerosa, el exceso de oferta de activos reduce su precio y su rentabilidad. POTERBA (2001) examinó empíricamente la hipótesis y encontró efectos de la estructura demográfica sobre la rentabilidad de los activos más moderados de lo que predecía la teoría. 

La conclusión, de nuevo, es que la producción es lo central: si hay menos trabajadores por jubilado, habrá menos producción por jubilado, y esa escasez se manifestará en el reparto como pensiones menores o cotizaciones mayores, y en la capitalización como menores rentabilidades o menores precios de los activos. La capitalización puede mitigar el problema si los fondos se invierten en países con demografías más jóvenes, porque así importa la producción de sus trabajadores. Pero esa posibilidad tiene sus propios riesgos, cambiarios y políticos, y su alcance es limitado.

Riesgos de cada sistema y el argumento de la diversificación

Con todo, si ninguno de los dos sistemas escapa a la demografía ni garantiza un mayor ahorro, la comparación relevante es la de los riesgos que cada uno impone y la de quién los soporta. 

Un sistema de reparto con prestación definida está expuesto a tres riesgos principales. El demográfico aparece si la relación entre cotizantes y pensionistas empeora. El económico aparece si el empleo o la productividad crecen menos de lo previsto. El político aparece si el legislador modifica las reglas. Un sistema de capitalización con aportación definida está expuesto a otros riesgos. El de mercado aparece si la rentabilidad de los activos es baja o negativa, especialmente cerca de la jubilación. El de inflación afecta si los activos no están indexados. El de longevidad aparece si el individuo no convierte su fondo en una renta vitalicia, o si la renta vitalicia es cara por la selección adversa. A ellos se añade el riesgo de gestión y los costes administrativos.

Los costes administrativos merecen una atención especial, porque su efecto se acumula de forma poco intuitiva. Una comisión de gestión del $1\%$ anual sobre el patrimonio acumulado parece pequeña, pero, a lo largo de una carrera de $40$ años, reduce el fondo final en torno a un $20\%$. Imaginemos, como ejemplo ilustrativo, un instrumento con una rentabilidad bruta del $4\%$. Un euro aportado al principio de la carrera crecería hasta unos $4,8$ euros sin comisiones y hasta unos $3,2$ euros con una comisión del $1\%$: una pérdida de un tercio para las aportaciones más tempranas y en torno a un quinto en promedio para toda la carrera. 

DIAMOND (1998) y WHITEHOUSE (2000) documentaron el elevado coste de los sistemas de cuentas individuales, especialmente en su fase inicial y en los países con mercados financieros poco desarrollados. La experiencia chilena y el escándalo británico de la venta abusiva de planes personales, confirmaron que la competencia entre gestores no garantiza costes bajos: los afiliados tienen poca información y poca sensibilidad a las comisiones, y los gestores compiten más en publicidad que en precio. Es el séptimo mito de ORSZAG y STIGLITZ. Los sistemas públicos de reparto, en cambio, tienen costes de gestión muy reducidos en relación con el volumen de prestaciones, por las economías de escala examinadas.

De la comparación de riesgos se deriva el argumento más sólido de la literatura contemporánea, que no es a favor de uno u otro sistema, sino de su combinación. MERTON, DUTTA, KAPUR y ORSZAG (2000) formularon el problema como una elección de cartera. La rentabilidad del reparto depende del crecimiento de la masa salarial; la de la capitalización, de la rentabilidad del capital. Como ambas no están perfectamente correlacionadas, un sistema que combine ambas fuentes de financiación diversifica el riesgo y ofrece, para cada nivel de rentabilidad esperada, un riesgo menor que cualquiera de los dos sistemas puros.\footnote{%
    Esta es, precisamente, la justificación teórica de los sistemas multipilar que recomendaron HOLZMANN y HINZ (2005) y que han adoptado Suecia, los Países Bajos, Suiza o Dinamarca, con un pilar público de reparto y un pilar de capitalización obligatorio o cuasiobligatorio de gran tamaño. En esos países, el segundo pilar representa varios puntos de PIB de gasto en pensiones. En España, el peso del segundo pilar es prácticamente nulo.
}

Desde la perspectiva de la diversificación, España tiene una cartera de pensiones concentrada casi exclusivamente en el riesgo de la masa salarial española, que es precisamente el riesgo más expuesto al envejecimiento demográfico del país.

Cuentas nocionales: el reparto que imita la capitalización

Las cuentas nocionales, merecen un análisis específico desde la teoría de la financiación, porque son una respuesta directa a los problemas del reparto de prestación definida que no requiere la transición a la capitalización. En un sistema de cuentas nocionales, la cotización de cada trabajador se anota en una cuenta virtual que se revaloriza cada año con un tanto nocional. Si ese tanto es igual a la tasa de crecimiento de la masa salarial, la rentabilidad que el sistema promete coincide exactamente con la que el reparto puede pagar de forma sostenible, que es la de SAMUELSON y AARON ($n + g$). SETTERGREN y MIKULA (2005) formalizaron esta relación y mostraron que un sistema de cuentas nocionales con tanto igual al crecimiento de la masa salarial y con un mecanismo de equilibrio automático es financieramente sostenible por construcción. Cualquier desequilibrio entre sus activos (el valor capitalizado de las cotizaciones futuras más el fondo de reserva) y sus pasivos (los derechos devengados) activa un ajuste automático de la revalorización.

Las ventajas de las cuentas nocionales son la transparencia, porque cada trabajador conoce el valor de sus derechos acumulados; la relación estrictamente actuarial entre cotizaciones y prestaciones, que reduce los desincentivos al trabajo; el ajuste automático a la esperanza de vida; y la ausencia de coste de transición, porque el sistema sigue siendo de reparto. Sus inconvenientes son que trasladan al individuo el riesgo demográfico y económico, que en el reparto de prestación definida soporta el compromiso socio-político (el Estado). También pueden reducir las pensiones de forma significativa si la esperanza de vida aumenta o la masa salarial crece poco, y su relación estrictamente actuarial elimina los elementos redistributivos del sistema, que deben financiarse aparte con impuestos. Traslada, además, a la vejez las desigualdades de la vida activa, en particular las de género, como vimos al examinar la contributividad.

\paragraph{Síntesis: la elección española y sus implicaciones}

Reunidos los argumentos, puede evaluarse la elección española por el reparto. Por un lado, la continuidad del reparto tiene justificaciones sólidas. La transición a la capitalización tendría un coste enorme, por la \textbf{deuda implícita del sistema}. Tampoco garantizaría un aumento del ahorro nacional si se financiara con deuda, ni escaparía al problema demográfico. Y la Constitución impide la sustitución completa del régimen público. Pero la teoría también identifica las debilidades de la elección española. Con una rentabilidad implícita del reparto ($n + g$) que la demografía está reduciendo, el sistema es cada vez menos atractivo en términos de rentabilidad para las generaciones jóvenes. Con un segundo pilar prácticamente inexistente, la cartera de pensiones de los españoles está concentrada en un único riesgo. 

Las reformas recientes han intentado mitigar la primera debilidad con una capitalización parcial, mediante la dotación del Fondo de Reserva con la cotización del MEI, y la segunda con los planes de pensiones de empleo de la Ley 12/2022. El alcance de ambas, muy limitado hasta ahora, se examina más adelante.

\subsubsection{Cotizaciones frente a impuestos}

#### La elección del instrumento: por qué cotizaciones

El sistema español se organiza desde un punto de vista contributivo o bismarckiano, en el que las prestaciones guardan relación con las contribuciones, y desde uno no contributivo o beveridgiano, en el que dependen de la situación de necesidad. La parte contributiva se financia a través de las cotizaciones sociales y la no contributiva, a través de la imposición general. La correspondencia entre la naturaleza de la prestación y la del instrumento que la financia es el principio de separación de fuentes que el Pacto de Toledo consagró en 1995 y la Ley 24/1997 desarrolló. 

Pero la elección del instrumento no es solo una cuestión de coherencia contable, sino una decisión con consecuencias económicas sobre la incidencia, la eficiencia y la legitimidad del sistema que conviene analizar con las herramientas de la teoría de la imposición.

La cotización social tiene tres rasgos que la distinguen del impuesto general. El primero es su \textbf{base}: recae sobre las rentas del trabajo, no sobre el conjunto de la renta o el consumo, y solo hasta un tope (la base máxima de cotización). El segundo es su \textbf{afectación}: sus recursos se destinan a financiar un gasto determinado, las prestaciones contributivas, y no se integran en el presupuesto general.\footnote{%
    La defensa teórica de la afectación (earmarking) procede de BUCHANAN (1963). Sostuvo que vincular un ingreso a un gasto concreto permite a los ciudadanos votar sobre el nivel de cada servicio público con conocimiento de su coste, en lugar de decidir sobre un presupuesto general en el que los costes y los beneficios se diluyen. La afectación mejoraría así la eficiencia de las decisiones colectivas y la disciplina del gasto. La posición contraria, tradicional en la hacienda pública y en los ministerios de hacienda, sostiene que la afectación produce rigidez presupuestaria: impide asignar los recursos a los usos más valiosos en cada momento y genera una ilusión de autonomía financiera que desaparece cuando el ingreso afectado no basta para cubrir el gasto. Exactamente lo que ha ocurrido con la Seguridad Social española desde 2012. Las cotizaciones afectadas no bastan para cubrir el gasto contributivo, y el Estado tiene que completarlas con transferencias y préstamos, de modo que la afectación se mantiene formalmente, pero se diluye en la práctica.
} El tercero es su \textbf{vínculo con la prestación}: lo que cada trabajador cotiza determina, en mayor o menor medida, lo que recibirá. Estos tres rasgos se justifican, en la teoría clásica de la hacienda pública, por el principio del beneficio: quien se beneficia de un servicio público debe pagarlo, y la cotización es el precio de un seguro que protege al propio cotizante.\footnote{%
    Este vínculo con la prestación encuentra también un argumento de economía política, en particular de la ilusión fiscal (PUVIANI, 1903), que justifica la cotización. Los ciudadanos perciben la cotización como una prima de seguro que les da derecho a una prestación futura, y no como un impuesto que desaparece en el presupuesto general. Por eso toleran niveles de cotización que no tolerarían en un impuesto de igual cuantía. Es una de las razones por las que los sistemas bismarckianos han podido financiarse con cotizaciones muy elevadas sin la resistencia que habría provocado una subida equivalente del impuesto sobre la renta. Es también la razón por la que su legitimidad depende de que se mantenga la percepción del vínculo entre cotización y prestación.
}



#### La incidencia: quién paga realmente la cotización

La primera cuestión económica que plantea la cotización es su incidencia: quién soporta realmente su coste. En España, como en la mayoría de los países europeos, la cotización se reparte legalmente entre la empresa y el trabajador. En 2026, de la cotización por contingencias comunes del $28,3\%$ del salario, un $23,6\%$ corresponde a la empresa y un $4,7\%$ al trabajador. En el conjunto de las cotizaciones (contingencias comunes, desempleo, formación profesional, FOGASA, contingencias profesionales y MEI), la parte empresarial supera ampliamente las tres cuartas partes del total. Sin embargo, sabemos por la Teoría Tributaria que la incidencia legal de cualquier figura poco, o nada, tiene que ver con su incidencia económica. Debemos, en cambio, estudiar la \textbf{traslación de la figura}.

La Teoría de la Imposición establece que, en un mercado de trabajo competitivo, la distribución de la carga económica de un impuesta, su incidencia, depende solo de las elasticidades relativas de la oferta y la demanda de trabajo. La parte del impuesto que soportan los trabajadores es igual a la elasticidad de la demanda de trabajo dividida por la suma de las elasticidades de la demanda y de la oferta. Como la oferta de trabajo es, en general, poco elástica, sobre todo en el margen (análisis marginal), la teoría predice que los trabajadores soportan la mayor parte de la cotización, sea cual sea su reparto legal, a través de salarios brutos menores de los que tendrían sin ella. Sencillo.

La evidencia empírica más citada procede de GRUBER (1997), que aprovechó la reforma chilena de 1981 como experimento natural. La reforma redujo drásticamente las cotizaciones a cargo de los empresarios, y GRUBER mostró que esa reducción se trasladó íntegramente a salarios más altos, sin efectos significativos sobre el empleo. Confirmación empírica de la predicción teórica. Sin embargo, estudios posteriores han matizado esta conclusión. SAEZ, MATSAGANIS y TSAKLOGLOU (2012), con una reforma griega que modificó las cotizaciones de forma distinta para distintas cohortes, encontraron que el reparto legal sí importaba a corto plazo: las cotizaciones empresariales no se trasladaban a los salarios y las de los trabajadores sí, lo que contradice la teoría estándar. SAEZ, SCHOEFER y SEIM (2019), con una reducción de las cotizaciones empresariales para los trabajadores jóvenes en Suecia, encontraron que la reducción no se trasladó a los salarios de esos trabajadores. Las empresas beneficiadas aumentaron el empleo y los salarios de todos sus trabajadores, no solo de los jóvenes, lo que sugiere que el mercado de trabajo no funciona como un mercado competitivo y que la negociación y el reparto de rentas dentro de la empresa median la incidencia. 

Nada fuera de lo que cabría esperar, por tanto. A largo plazo y en mercados competitivos, los trabajadores soportan la mayor parte de la cotización. Pero a corto plazo, y en mercados con negociación colectiva, rigideces salariales o poder de mercado, el reparto legal y las instituciones laborales pueden importar.

La relevancia de esta discusión para España es considerable. En primer lugar, las bonificaciones y reducciones de cotizaciones a la contratación, que han sido durante décadas uno de los principales instrumentos de la política de empleo española, solo crean empleo si la reducción de la cotización no se traslada íntegramente a los salarios. Las evaluaciones disponibles, incluida la revisión del gasto de la AIReF (2019), encontraron que buena parte de las bonificaciones tenía un peso muerto elevado: subvencionaban contrataciones que se habrían producido de todos modos. En segundo lugar, las subidas de cotizaciones de la reforma de 2023 se trasladarán, según la teoría, en buena medida a los salarios de los trabajadores a medio plazo. En tercer lugar, el hecho de que la mayor parte de la cotización sea formalmente empresarial tiene un efecto sobre la percepción. CHETTY, LOONEY y KROFT (2009) mostraron que los individuos reaccionan menos a los impuestos que no son visibles (poco salientes), y la cotización empresarial no aparece en la nómina del trabajador como una deducción, sino como un coste que solo ve la empresa. Los trabajadores españoles, por tanto, subestiman el coste total de su protección social, lo que refuerza la ilusión fiscal señalada antes.

#### La cotización como precio o como impuesto: el argumento de SUMMERS

La segunda cuestión es la eficiencia. ¿Cuánto distorsiona la cotización el funcionamiento del mercado de trabajo?

SUMMERS (1989), mostró que la respuesta depende de cómo valoren los trabajadores las prestaciones que la cotización financia.\footnote{
En su art. sobre \textit{la economía sencilla de las prestaciones obligatorias} (mandated benefits),
} Si un trabajador valora en un euro cada euro de cotización, porque espera recibir a cambio una prestación de igual valor actual, la cotización no es un impuesto, sino un precio, y no distorsiona su decisión de trabajar. Su oferta de trabajo responde al salario más el valor de la prestación, que no ha cambiado. Si, por el contrario, el trabajador no valora en absoluto la prestación, porque no espera recibirla o porque no depende de lo que cotice, la cotización es un impuesto puro y distorsiona su oferta de trabajo como cualquier impuesto sobre la renta del trabajo. En los casos intermedios, solo la parte de la cotización que no se valora como prestación actúa como impuesto.

Un ejemplo muestra la importancia práctica del argumento. Supongamos dos trabajadores con el mismo salario, cada uno de los cuales cotiza un euro adicional. El primero tiene una carrera larga y un salario por debajo de la base máxima: su euro adicional se traduce en una pensión futura algo mayor, cuyo valor actual es, por ejemplo, de setenta céntimos. Para él, la cotización adicional equivale a un impuesto de treinta céntimos. El segundo tiene un salario por encima de la base máxima: su euro adicional corresponde a una \textit{cuota de solidaridad} (o al MEI), que no generan ningún derecho de pensión adicional. Para él, la cotización adicional equivale a un impuesto de un euro. La misma cotización nominal tiene, por tanto, efectos distintos según el vínculo entre cotización y prestación que perciba cada trabajador. Es lo que convierte la contributividad en una variable de eficiencia y no solo de equidad.

DISNEY (2004), con datos de los países de la OCDE, contrastó empíricamente el argumento de SUMMERS. Encontró que las cotizaciones a los sistemas de pensiones tienen efectos sobre el empleo, sobre todo el femenino, más negativos en los países donde el vínculo entre cotizaciones y prestaciones es débil. Es decir, las cotizaciones actúan en mayor medida como un impuesto sobre el empleo en esos países. De ahí se deriva una conclusión relevante para el diseño del sistema: reforzar la contributividad no solo mejora la equidad horizontal, sino que reduce el coste en términos de empleo de la financiación, porque convierte una parte del impuesto en precio.

Cuña fiscal y sus efectos sobre el empleo

La suma de las cotizaciones de la empresa y del trabajador y del impuesto sobre la renta que grava el salario constituye la cuña fiscal sobre el trabajo. Es la diferencia entre lo que cuesta un trabajador a la empresa y lo que recibe en términos netos. Dicho de otra forma, es la diferencia entre el precio de equilibrio en el mercado de trabajo y el precio después de impuestos.

Según las estimaciones anuales de la OCDE en su publicación \textit{Taxing Wages}, la cuña fiscal para un trabajador soltero con salario medio en España se sitúa en torno al $40\%$ del coste laboral total, algo por encima de la media de la OCDE. La cotización social, sobre todo la empresarial, representa la mayor parte de esa cuña. De hecho, hay que tener en cuenta que el peso de las cotizaciones en la cuña fiscal española es superior al del impuesto sobre la renta, a diferencia de lo que ocurre en los países anglosajones y nórdicos, que financian una parte mayor de su protección social con impuestos generales.

La relación entre la cuña fiscal y el empleo ha sido objeto de una extensa literatura. DAVERI y TABELLINI (2000) encontraron que el aumento de la cuña fiscal sobre el trabajo explicaba una parte significativa del aumento del desempleo en los países europeos continentales desde los años setenta. Atribuyeron el efecto a la interacción entre los impuestos sobre el trabajo y la negociación colectiva centralizada, que traslada el impuesto a los costes laborales en lugar de a los salarios netos. Estudios posteriores han matizado la magnitud del efecto, que depende de las instituciones del mercado de trabajo de cada país. Pero la idea de que una cuña fiscal elevada perjudica al empleo, especialmente al de los trabajadores de baja cualificación, cuya demanda es más elástica, es ampliamente aceptada. Hasta el punto de que constituye el fundamento de las recomendaciones reiteradas de la OCDE, del FMI y de la Comisión a España: reducir cotizaciones sociales, sobre todo en los salarios bajos, y compensar la pérdida de recaudación con otros impuestos.

Devaluación fiscal: sustituir cotizaciones por IVA

Esta recomendación ha dado lugar a un debate específico que conviene conocer: el de la \textbf{devaluación fiscal}. FARHI, GOPINATH e ITSKHOKI (2014) demostraron formalmente que, en una unión monetaria donde un país no puede devaluar su tipo de cambio, una reducción de las cotizaciones sociales empresariales compensada con un aumento del IVA produce efectos reales equivalentes a una devaluación nominal. La reducción de las cotizaciones abarata los costes laborales y, por tanto, los precios de las exportaciones. El aumento del IVA, que grava las importaciones pero no las exportaciones, encarece los productos importados. DE MOOIJ y KEEN (2013) estimaron los efectos empíricos de estas operaciones y encontraron que eran significativos a corto plazo, pero transitorios, porque los salarios acaban ajustándose. La devaluación fiscal fue recomendada a España durante la crisis de deuda soberana como sustituto de la devaluación del tipo de cambio que la pertenencia al euro impedía. El gobierno aumentó el IVA en 2010 y 2012, pero no redujo sustancialmente las cotizaciones, de modo que la operación no se realizó en sus términos teóricos.

La devaluación fiscal plantea, sin embargo, un problema de coherencia con la lógica contributiva. Si las cotizaciones se sustituyen por IVA, el vínculo entre lo que cada trabajador aporta y lo que recibe se debilita, porque el IVA no se asocia a ningún derecho de prestación. El sistema se desplaza así hacia una financiación beveridgiana, aunque las prestaciones sigan siendo bismarckianas. La coherencia exigiría, entonces, que las prestaciones también se desligaran de las cotizaciones y se acercaran a un mínimo uniforme. Es el dilema que el siguiente apartado examina en la evolución reciente del sistema español.

Reforma silenciosa: de Bismarck a Beveridge por la puerta de atrás

Uno de los debates más relevantes de la financiación española, que continua lo mencionado en el apartado anterior, es la llamada \textbf{reforma silenciosa} de las pensione. El trabajo de referencia es el art. \textit{From Bismarck to Beveridge: the other pension reform in Spain} (CONDE-RUIZ y GONZÁLEZ, 2016) y su título resume la tesis principal: mientras el debate público se centraba en las reformas explícitas (la de 2011 y la de 2013), el sistema español estaba experimentando una transformación implícita de su naturaleza, de un modelo bismarckiano, con prestaciones proporcionales a las cotizaciones, a uno beveridgiano, con prestaciones cada vez más uniformes y desligadas de lo aportado. 

El mecanismo es sencillo. La base máxima de cotización, que limita lo que aporta cada trabajador, crece más deprisa que la pensión máxima, que limita lo que recibe. En consecuencia, una proporción creciente de trabajadores cotiza por encima del nivel que les daría derecho a la pensión máxima. Para ellos, una parte creciente de su cotización no genera ningún derecho adicional y funciona, en los términos de SUMMERS, como un impuesto puro.


Para probar esto empíricamente basta un rápido ejercicio de cálculo, muy sencillo. Si la base máxima crece más deprisa que la pensión máxima, la ratio entre la pensión máxima y la base máxima disminuye. En 2010, la base máxima mensual era de $3.198$ euros y la pensión máxima mensual, de $2.466,20$ euros, que se cobra en catorce pagas. En términos anuales, la pensión máxima, de unos $34.500$ euros, equivalía en torno al $90\%$ de la base máxima anual, de unos $38.400$ euros. En 2022, con una base máxima de $4.139,40$ euros y una pensión máxima de $2.819,18$ euros, la proporción anual había bajado al entorno del $79\%$. En 2026, con una base máxima de $5.101,20$ euros y una pensión máxima de $3.359,60$ euros mensuales ($47.034,40$ euros anuales), la proporción anual se sitúa en torno al $77$ por ciento de la base máxima anual de $61.214,40$ euros. 

La reforma de 2023 aceleró deliberadamente esta transformación. Estableció que la base máxima crecerá cada año con el IPC más $1,2$ puntos porcentuales hasta 2050, mientras que la pensión máxima crecerá con el IPC más un incremento adicional de solo $0,115$ puntos anuales entre 2025 y 2050. ¿Qué quiere decir? Si suponemos una inflación del $2\%$, el aumento real de la base máxima llegaría al $37,1\%$ en 2050, mientras que el crecimiento real acumulado de la pensión máxima entre 2025 y 2050 ascendería aproximadamente al $3\%$. A ello se añaden la cuota de solidaridad sobre los salarios que exceden la base máxima y la cotización del MEI, ninguna de las cuales genera derechos de pensión.

Las implicaciones de esta transformación son de tres tipos. En términos de equidad, aumenta la progresividad del sistema, porque los trabajadores de salarios altos aportan cada vez más en relación con lo que reciben. Sus defensores lo presentan como una forma de solidaridad intrageneracional y de financiación del sistema por quienes tienen mayor capacidad de pago. En términos de eficiencia, convierte una parte creciente de la cotización de los trabajadores más cualificados en un impuesto puro sobre el trabajo, lo que, según el argumento de SUMMERS y la evidencia de DISNEY, aumenta su coste en términos de empleo y de oferta de trabajo cualificado. Aumenta además los incentivos a la elusión, mediante la retribución en especie o la conversión de trabajadores en autónomos o en socios de sociedades profesionales. En términos de legitimidad, debilita la percepción del vínculo entre cotización y prestación que, como vimos, sostiene la tolerancia a la cotización, especialmente entre las clases medias altas. Conecta así con la paradoja de la redistribución de KORPI y PALME examinada: un sistema que redistribuye más a costa de las clases medias altas puede perder el apoyo político de estas.

No obstante, el problema de fondo, que CONDE-RUIZ y GONZÁLEZ señalan, es la incoherencia de la transformación. El sistema se desplaza hacia BEVERIDGE en su financiación sin reconocerlo abiertamente y sin adaptar su diseño a esa nueva naturaleza. Un sistema beveridgiano coherente financiaría con impuestos generales, y no con cotizaciones sobre el salario, una pensión uniforme de suficiencia. Dejaría el mantenimiento del nivel de vida por encima de ese mínimo a pilares complementarios de capitalización, como en el Reino Unido o los Países Bajos. El sistema español, en cambio, mantiene la retórica y la estructura bismarckiana de las prestaciones, y financia con cotizaciones cada vez menos contributivas una pensión máxima cada vez más alejada del salario de los trabajadores cualificados. Todo ello sin un segundo pilar que complemente sus pensiones. Es un ejemplo de lo que la literatura de PIERSON denominaría un cambio de política por deriva: la naturaleza de una institución cambia sin reforma explícita, por el efecto acumulado de decisiones parciales. Es una de las formas más eficaces de evitar la culpa, porque ningún gobierno anuncia que está transformando el sistema.

La financiación con impuestos de las prestaciones no contributivas y las cargas impropias

El reverso de la cotización es la financiación con impuestos de las prestaciones que no guardan relación con lo aportado, como la parte no contributiva del sistema español. El principio de separación de fuentes exige que las cotizaciones financien solo las prestaciones contributivas y que el Estado financie con impuestos las no contributivas, incluidos los complementos a mínimos de las pensiones contributivas, que se consideran de naturaleza no contributiva porque su cuantía no depende de lo cotizado. 

Pero la separación de fuentes plantea una cuestión más amplia, que el Pacto de Toledo de 2020 formuló como la de los gastos impropios. Gastos que el sistema contributivo financia con cotizaciones pero que, por su naturaleza, deberían financiarse con impuestos. Entre ellos las reducciones y bonificaciones de cotizaciones que el Estado concede como política de empleo, las jubilaciones anticipadas de sectores concretos por razones de política industrial, las prestaciones de nacimiento y cuidado de menor, que responden a una política de familia y no a un seguro de rentas, y los complementos para la reducción de la brecha de género. Desde 2021, las leyes de presupuestos han incrementado sustancialmente las transferencias del Estado a la Seguridad Social para financiar estos gastos.

Según una lectura, las transferencias por gastos impropios corrigen una anomalía histórica: el sistema contributivo venía financiando con cotizaciones políticas que correspondían al Estado, y su déficit aparente se debía en parte a esa anomalía. Según la lectura opuesta, las transferencias ocultan el déficit estructural del sistema contributivo, al trasladarlo al presupuesto del Estado, sin mejorar en nada la situación de las Administraciones Públicas en su conjunto. Ambas lecturas son compatibles con los mismos datos, porque la frontera entre lo propio y lo impropio del sistema no es técnica, sino normativa: depende de qué se considere la función del sistema contributivo, y es la misma frontera entre la hucha y Robin Hood (BARR). 

La conclusión que conviene retener es que el sistema español ya no se financia solo con cotizaciones. Se financia con una combinación de cotizaciones cada vez menos contributivas, transferencias crecientes del Estado y préstamos, y esa combinación lo ha convertido en un sistema mixto en su financiación, aunque siga siendo bismarckiano en la definición de sus prestaciones. Presentar hoy la financiación como la suma de cotizaciones y transferencias equivale a omitir más de $136.000$ millones de euros de deuda y un patrimonio neto negativo que el Tribunal de Cuentas acaba de señalar. 


\subsection{Fuentes de financiación vigentes}
Hemos establecido en términos teóricos cuándo y con qué pueden financiarse las pensiones. Ahora debemos describir cómo se financia hoy el sistema español y evalúa cada fuente con las herramientas teóricas ya expuestas. 

El texto refundido de la Ley General de la Seguridad Social enumera los recursos del sistema. Son las aportaciones progresivas del Estado, que se consignan con carácter permanente en sus presupuestos generales; las cotizaciones de las personas obligadas; las cantidades recaudadas por recargos, intereses y sanciones; los frutos, rentas e intereses de su patrimonio; y cualesquiera otros ingresos. 

Por el lado de los ingresos, la principal fuente de financiación de la Seguridad Social son las cotizaciones sociales y que, a gran distancia, siguen las transferencias del Estado, cuyo objeto es financiar las prestaciones de carácter no contributivo. Esa descripción, correcta en su jerarquía, es hoy incompleta en dos sentidos. Las transferencias del Estado ya no financian solo prestaciones no contributivas, sino también, desde 2021, una parte de los gastos del nivel contributivo calificados como impropios. Y a las dos fuentes tradicionales se ha añadido una tercera, el endeudamiento con el Estado, que ha dejado de ser un recurso extraordinario para convertirse en un componente estructural de la financiación del sistema.

Organizamos esta sección en tres partes, que corresponden a esas tres fuentes: las cotizaciones, las transferencias del Estado y el endeudamiento y las reservas, que incluyen el relato completo del Fondo de Reserva. La tesis que se desarrolla es la que anticipamos: el sistema español se financia hoy con una combinación inestable de cotizaciones cada vez menos contributivas, transferencias crecientes y préstamos que no se devuelven. 

Reunidas las tres fuentes, la financiación del sistema español puede describirse con precisión. Las cotizaciones siguen siendo la fuente principal y crecen a un ritmo elevado, impulsadas por el empleo y por las subidas de 2023. Pero una parte creciente de ellas no genera derechos de pensión y funciona, en términos económicos, como un impuesto sobre el trabajo. Las transferencias del Estado financian el nivel no contributivo y, desde 2021, una parte sustancial del contributivo calificada como impropia. Han desplazado hacia el presupuesto del Estado una parte del coste del sistema sin mejorar la situación de las Administraciones Públicas en su conjunto. El endeudamiento con el Estado, que en la práctica no se devuelve, cubre cada año el desfase restante, y ha llevado el patrimonio neto de la Seguridad Social a valores negativos superiores a los cien mil millones de euros. El Fondo de Reserva, alimentado ahora por una cotización específica, crece a un ritmo inferior al de esa deuda y está invertido en la misma deuda pública del Estado que financia los préstamos.

El resultado es un sistema cuya financiación es mixta, porque combina cotizaciones, impuestos y deuda; opaca, porque la frontera entre esas fuentes depende de convenciones contables discutibles; e inestable, porque su equilibrio formal depende de decisiones anuales del legislador sobre transferencias y préstamos que, en un contexto de presupuestos prorrogados y de fragmentación parlamentaria, no están garantizadas. Esta constatación no prejuzga todavía si el sistema es o no sostenible, que es la pregunta del bloque IV. Pero sí muestra que el saldo de la Seguridad Social es un mal indicador para responderla, y que la sostenibilidad debe evaluarse con el gasto total en pensiones, los recursos totales que lo financian y su evolución a largo plazo.

\subsubsection{Cotizaciones sociales}
La cotización del Régimen General, que agrupa a la gran mayoría de los asalariados, resulta de aplicar unos tipos a una base. 

La base de cotización está constituida por la remuneración total, cualquiera que sea su forma o denominación, que el trabajador tenga derecho a percibir mensualmente, con inclusión de la parte proporcional de las pagas extraordinarias y con algunas exclusiones tasadas, como determinadas dietas o indemnizaciones. La base está sometida a dos límites. El tope mínimo depende del grupo de cotización al que pertenezca la categoría profesional del trabajador; hay once grupos, y el más bajo está vinculado al SMI. El tope máximo es único para todos los grupos.\footnote{%
    \textcolor{red}{\textbf{NOTA.-} En 2026, la base máxima es de $5.101,20$ euros mensuales, un $3,9\%$ superior a la de 2025 en aplicación de la regla de la reforma de 2023, que la revaloriza con el IPC más $1,2$ puntos. El SMI para 2026, fijado por el Real Decreto 126/2026, de 18 de febrero, es de $1.221$ euros mensuales en catorce pagas. De ahí resulta una base mínima general, con prorrateo de pagas, de unos $1.424,50$ euros mensuales.}
}

El tipo general de cotización sobre los rendimientos brutos del trabajo se ha dividido, históricamente, en cinco tipos: contingencias comunes, desempleo, formación profesional, aportación FOGASA, y, además, un tipo por accidentes de trabajo y enfermedades profesionales que varía por actividad y recae íntegramente en el empleador.\footnote{%
    La cotización por accidentes de trabajo y enfermedades profesionales tiene una lógica distinta del resto, heredera de la Ley de Accidentes de Trabajo de 1900. Recae exclusivamente sobre la empresa, porque su fundamento es la responsabilidad objetiva del empresario por los riesgos inherentes a su actividad. Su tipo no es uniforme, sino que se fija en una tarifa de primas según la actividad económica de la empresa, de acuerdo con la clasificación nacional de actividades económicas (CNAE). Así refleja, aunque de forma grosera, el distinto riesgo de cada sector. Desde 2026, las empresas deben comunicar su actividad conforme a la nueva CNAE-2025; en caso contrario, se les aplica el tipo más alto de su actividad. La tarifa diferenciada por sectores es una forma elemental de tarificación por experiencia a nivel sectorial. Pero, como vimos, no discrimina entre empresas del mismo sector según su siniestralidad individual, salvo por el limitado sistema de incentivos a la reducción de la siniestralidad. Por eso las empresas seguras de un sector subvencionan a las inseguras, y el incentivo a la prevención se debilita. Los recursos de esta cotización se destinan en gran parte a las mutuas colaboradoras que gestionan las contingencias profesionales de las empresas asociadas, y sus excedentes nutren el Fondo de Contingencias Profesionales de la Seguridad Social.
} Cada uno de los tres primeros tipos presenta un tramo a cargo del empleador y otro del trabajador. No obstante, desde 2023, se añade la cotización del MEI.\footnote{%
    \textcolor{red}{\textbf{NOTA.-}} En 2026 los tipos del Régimen General son los siguientes. Por contingencias comunes, el $28,3\%$, del que el $23,6\%$ corresponde a la empresa y el $4,7\%$ al trabajador. Por desempleo, el $7,05\%$ en los contratos indefinidos ($5,5$ empresa y $1,55$ trabajador) y el $8,3\%$ en los temporales ($6,7$ empresa y $1,6$ trabajador). Para el FOGASA, el $0,2\%$, a cargo exclusivo de la empresa. Para formación profesional, el $0,7\%$ ($0,6$ empresa y $0,1$ trabajador). Y por el MEI, el $0,9\%$ ($0,75$ empresa y $0,15$ trabajador). A ellos se suma la cotización por contingencias profesionales, variable según la actividad y a cargo exclusivo de la empresa.
} El tipo conjunto (2026), sin contar las contingencias profesionales, asciende al $37,15\%$ para un contrato indefinido.

Un ejemplo numérico permite ver el efecto conjunto. Tomemos un trabajador con contrato indefinido y un salario bruto mensual de $2.000$ euros, con pagas prorrateadas y un tipo por contingencias profesionales del $1,5\%$ (próximo a la media de actividades de riesgo bajo o medio). La empresa cotiza $472$ euros por contingencias comunes, $110$ por desempleo, $12$ por formación profesional, $4$ para el FOGASA, $15$ por el MEI y unos $30$ por contingencias profesionales: en total, unos $643$ euros. El trabajador cotiza $94$ euros por contingencias comunes, $31$ por desempleo, $2$ por formación profesional y $3$ por el MEI: unos $130$ euros. El coste laboral total es de unos $2.643$ euros, y el trabajador percibe, antes del impuesto sobre la renta, unos $1.870$ euros. 

Las cotizaciones representan, por tanto, en torno al $29\%$ del coste laboral total, de las que más de las cuatro quintas partes son formalmente empresariales. Sin embargo, como vimos, la teoría de la incidencia sugiere que buena parte de esa cotización empresarial la soporta en realidad el trabajador, a través de un salario bruto menor del que tendría sin ella. Y la baja saliencia de la cotización empresarial, que no aparece como deducción en la nómina, hace que el trabajador perciba solo $130$ de esos $773$ euros.


\paragraph{RETA: de la base elegida a los rendimientos reales}
El Régimen Especial de Trabajadores Autónomos plantea un problema de financiación que no existe en el Régimen General: la medición de la base.

En el Régimen General, la base es el salario, que la empresa conoce y declara, y que la Administración puede verificar con información de terceros. En el caso de los autónomos, la base debería ser su renta, que solo ellos conocen y declaran. Durante décadas, la solución fue permitir que cada autónomo eligiera su base de cotización entre un mínimo y un máximo, con independencia de sus ingresos reales.

En torno a ocho de cada diez autónomos cotizaban por la base mínima o muy cerca de ella, lo que producía dos efectos. En el plano financiero, el régimen recaudaba poco en relación con la renta del colectivo y era estructuralmente deficitario. En el plano de la suficiencia, los autónomos generaban pensiones muy bajas, frecuentemente completadas con complementos a mínimos financiados por el Estado.

Como sabemos, el comportamiento de los autónomos era racional. Cotizar por la base mínima durante la mayor parte de la carrera y aumentarla en los años previos a la jubilación, cuando el periodo de cálculo de la base reguladora era corto, maximizaba la rentabilidad de las cotizaciones. Contar con el complemento a mínimos, además, protegía frente a una pensión demasiado baja, que es exactamente el dilema del samaritano de BUCHANAN. La prolongación del periodo de cálculo de la base reguladora y las restricciones al aumento de la base a partir de determinada edad redujeron esa estrategia, pero no eliminaron el problema de fondo.

La reforma llegó con el Real Decreto-ley 13/2022, de 26 de julio. Su funcionamiento combina tres elementos. El primero es una tabla de tramos de rendimientos netos, a cada uno de los cuales corresponde una base mínima y una máxima, que el autónomo elige provisionalmente según sus ingresos previstos. El segundo es la regularización posterior: la Tesorería General cruza la base provisional con los rendimientos efectivamente declarados a la AEAT y reclama o devuelve la diferencia. El tercero es un calendario de transición hasta 2032, en el que las bases de los tramos se irán ajustando progresivamente. En 2026, las cuotas de los tramos se han mantenido congeladas respecto de 2025, con la única excepción del aumento derivado del MEI, lo que refleja la dificultad política de avanzar en el calendario previsto. El tipo total del régimen se sitúa en torno al $31,5\%$.

La reforma ilustra dos problemas generales de la Teoría de la Imposición: la importancia de la percepción y la medición de las rentas autodeclaradas. Empezando por el segundo, KLEVEN, KNUDSEN, KREINER, PEDERSEN y SAEZ (2011), con un gran experimento de campo en Dinamarca, mostraron que la evasión es casi inexistente en las rentas sobre las que existe información de terceros, como los salarios declarados por las empresas. En cambio, es considerable en las rentas autodeclaradas, como los rendimientos de actividades por cuenta propia. Vincular la cotización de los autónomos a sus rendimientos fiscales traslada al sistema de Seguridad Social la calidad, y los defectos, de la medición fiscal de esos rendimientos. Si los autónomos infradeclaran sus rendimientos ante Hacienda, también cotizarán por debajo de sus ingresos reales. La reforma es, por tanto, un avance en equidad horizontal, pero su eficacia depende de la capacidad de la AEAT para medir las rentas de los autónomos. Creando, además, un incentivo adicional a infradeclarar, porque cada euro de rendimiento declarado ya no solo tributa por el impuesto sobre la renta, sino que eleva también la cotización. 

Pero, quizás, lo más interesante es la primera: ¿pagan mucho los autónomos? Si nos ajustamos a los datos, en realidad soportan menos carga fiscal. Sin embargo, la saliencia de sus cotización es infinitamente alta si lo comparamos con la saliencia de las cotizaciones empresariales (que es nula). De esta forma, el autónomo, que deba declarar y cumplir con sus obligaciones tributarias de forma individual y consciente, ve como una porción mucho mayor de su ingreso se ve mermado con respecto al de un trabajador que él perciba como \textit{análogo}. Sin embargo, la carga fiscal real, si comparamos los tipos totales de los regímenes, es de un $6\%$ inferior.

Nuevas figuras de 2023: base máxima, MEI y cuota de solidaridad

La reforma de 2021-2023 introdujo tres figuras de cotización que comparten un rasgo fundamental: no generan, o generan de forma muy limitada, derechos de pensión para quien las paga. En los términos de SUMMERS, se acercan a impuestos sobre el trabajo más que a precios de un seguro.

La primera es la \textbf{revalorización de la base máxima} por encima de la inflación, que analizamos en los apartados anteriores. Es el núcleo de la reforma silenciosa: los trabajadores con salarios por encima de la base máxima cotizarán cada vez más sin que su pensión máxima crezca en proporción.

La segunda es el \textbf{MEI}. Es un incremento finalista de las cotizaciones sociales durante el periodo 2023-2050, al objeto de dotar el Fondo de Reserva de la Seguridad Social.\footnote{%
    El aumento es progresivo: comenzó con $0,6$ puntos porcentuales y crece una décima al año hasta alcanzar $1,2$ puntos en 2029, manteniéndose aproximadamente la distribución entre empresa y trabajador de la cotización por contingencias comunes.
} En 2026 el tipo es del $0,9$ por ciento, del que $0,75$ corresponde a la empresa y $0,15$ al trabajador, y el Gobierno estima que recaudará unos $5.300$ millones de euros. Su naturaleza jurídica es la de una cotización afectada a un fin específico, la dotación del Fondo de Reserva, y no genera derechos de pensión individuales. Por eso funciona como un impuesto afectado sobre la masa salarial, cuya función en la estrategia de sostenibilidad del sistema se examina más adelante.

La tercera es la \textbf{cuota de solidaridad}. Ésta se introdujo al objeto de gravar la parte de los salarios que exceda de la base máxima de cotización, determinada por una escala progresiva de tres tramos.\footnote{%
    Los tipos varían entre el $0,92$ y el $1,17\%$ en 2025 y van subiendo anualmente hasta alcanzar un tipo mínimo del $5,5\%$ y uno máximo del $7\%$ en 2045. El tipo mínimo corresponde a los salarios entre la base máxima y un $10\%$ adicional; el máximo, a las retribuciones superiores al $50\%$ adicional de la base máxima. La distribución entre empresa y trabajador mantiene la proporción de la cotización por contingencias comunes, aproximadamente un $83/17\%$. En 2026 los tres tramos son: un $1,15\%$ sobre la parte del salario comprendida entre $5.101,21$ y $5.611,32$ euros mensuales, un $1,25\%$ sobre la parte entre $5.611,33$ y $7.651,80$ euros, y un $1,46\%$ sobre la parte que supera $7.651,80$ euros.
} Supongamos, para ilustrar este mecanismo, un directivo con un salario de $8.000$ euros mensuales. 

Éste pagaría en 2026 una cuota de solidaridad de unos $36$ euros al mes: $5,87$ euros por el primer tramo, $25,51$ por el segundo y $5,08$ por el tercero. Con los tipos previstos para 2045 ($5,5$, $6$ y $7$ por ciento) y la misma escala en términos reales, pagaría unos $175$ euros mensuales. La cuota de solidaridad convierte así la cotización en un instrumento de progresividad, como otras figuras de la reforma ya analizadas. Recordemos que la base máxima hacía que la cotización fuera regresiva en la parte alta de la distribución, porque los salarios por encima del tope pagaban un porcentaje decreciente de su salario. La cuota corrige parcialmente esa regresividad, pero a costa de desligar aún más la cotización de la prestación. Sus defensores la presentan como una forma de solidaridad intrageneracional y de ampliación de la base del sistema, en la línea de las propuestas de destope que se han debatido también en Estados Unidos. Sus críticos subrayan que es, en esencia, un impuesto adicional sobre las rentas del trabajo altas, que se suma a los tipos marginales máximos del impuesto sobre la renta. Aumenta así los incentivos a la elusión (retribución en especie, conversión en autónomos o en socios de sociedades) y a la localización en el extranjero del trabajo más cualificado.

\textcolor{blue}{\textbf{NOTA.-} Debemos destacar, aunque no entraremos, en que esta regresividad perenne del sistema puede estar detrás de las grandes divergencias salariales que existen en la distribución media-alta de ingresos por rentas del trabajo. Tengamos en cuenta, para ilustrar como esto afectaría sobre las rentas salariales, que cuando una empresa paga un salario que queda por encima de la base máxima de cotización, un euro adicional de salario no soporta carga fiscal en términos de cotizaciones, por lo que el mercado vuelve a competir en precios por los trabajadores altamente cualificados, capaces de ocupar estos cargos. De esta forma, aunque parezca paradógico teniendo en cuenta el principio de progresividad, la cuña fiscal puede tener una relación decreciente en salario a partir de la base máxima de cotización.} 

Bonificaciones y reducciones: la cotización como instrumento de política de empleo

La cotización no es solo un instrumento de financiación, sino que se ha utilizado durante décadas como instrumento de política de empleo, mediante bonificaciones y reducciones para la contratación de determinados colectivos o para el fomento del empleo autónomo. Las bonificaciones se financian con transferencias del Estado a la Seguridad Social, a través del presupuesto del SEPE; las reducciones, en cambio, las soporta el propio sistema, y son uno de los gastos impropios examinados más adelante. El Real Decreto-ley 1/2023, de 10 de enero, reformó el sistema de incentivos a la contratación con el objetivo de concentrarlos en los colectivos con mayores dificultades de inserción. La tarifa plana para los nuevos autónomos, una cuota reducida durante el primer año de actividad, es una de las medidas más conocidas y ha sido objeto de evaluaciones que cuestionan su eficacia para crear empleo sostenible.

Como vimos, la eficacia de estos incentivos depende de su incidencia y de su peso muerto, es decir, de cuántas contrataciones subvencionadas se habrían producido de todos modos. La revisión del gasto de la AIReF de 2019 sobre incentivos a la contratación encontró un peso muerto elevado en buena parte de ellos y recomendó reducir su número y focalizarlos. Desde la perspectiva de la financiación del sistema, las bonificaciones y reducciones plantean un problema adicional: convierten la Seguridad Social en un instrumento de políticas que no le son propias. Erosionan así la correspondencia entre cotizaciones y prestaciones y hacen más opaca la contabilidad del sistema, que es el fundamento del debate sobre los gastos impropios.

Recaudación y su dinámica

La recaudación del Sistema ha crecido con intensidad los últioms años. Entre enero y julio de 2026, los ingresos por cotizaciones alcanzaron $93.957$ millones de euros, un $7,4\%$ más que en el mismo periodo del año anterior, y el Gobierno estima que el conjunto de 2026 cerrará con una recaudación récord de unos $189.800$ millones de euros. La afiliación media se situó a finales de 2025 cerca de los $21,9$ millones de personas, tras cuatro años consecutivos de aumentos en torno al medio millón de afiliados. 

El crecimiento de la recaudación responde a cuatro factores que conviene distinguir, porque su sostenibilidad es muy diferente. El primero es el aumento del empleo, impulsado en gran medida por la inmigración. El segundo es el crecimiento de los salarios nominales, alimentado por la inflación de 2022-2023. El tercero es la reforma laboral de 2021, que al reducir la temporalidad aumentó la estabilidad de las cotizaciones. El cuarto son las subidas de tipos y de bases de la reforma de 2023. Los dos primeros son en buena medida cíclicos; los dos últimos, estructurales. La evaluación de la AIReF de 2026 sobre la regla de gasto en pensiones atribuyó una parte de la mejora de sus previsiones al mayor rendimiento de la nueva cotización de los autónomos, lo que muestra que el efecto recaudatorio de las reformas puede superar las estimaciones iniciales.

\subsubsection{Transferencias del Estado}
Entre los ingresos de la Seguridad Social destacan las transferencias corrientes procedentes del Estado, cuyo objeto es financiar las prestaciones no contributivas así como una serie de gastos impropios. Esta descripción general recoge las dos capas que las transferencias del Estado financian hoy, y que responden a dos lógicas distintas.

La primera capa es la financiación del nivel no contributivo, que responde al principio de separación de fuentes del Pacto de Toledo de 1995. Comprende las pensiones no contributivas de jubilación e invalidez, el IMV, las prestaciones familiares no contributivas y los complementos a mínimos de las pensiones contributivas.\footnote{%
    La financiación íntegra de los complementos a mínimos con transferencias del Estado se completó en 2013, tras un largo periodo transitorio iniciado con la Ley 24/1997 en el que una parte se financiaba con cotizaciones.
} Su justificación es clara: son prestaciones de Robin Hood, no de hucha, y deben financiarse con impuestos generales (BARR). La segunda capa es la financiación de los gastos impropios del nivel contributivo, que respondió a la primera recomendación del Pacto de Toledo de 2020. 

Desde los Presupuestos de 2021, el Estado transfiere a la Seguridad Social recursos adicionales para financiar gastos que, aun formando parte del nivel contributivo, se consideran ajenos a la lógica del seguro. Comprenden las reducciones de cotizaciones por políticas de empleo, las prestaciones de nacimiento y cuidado de menor, las jubilaciones anticipadas de determinados colectivos por razones de política sectorial, las ayudas a sectores en reconversión, los complementos para la reducción de la brecha de género y otros conceptos. 

La cuantía de estas transferencias ha sido considerable: en torno a $18.400$ millones de euros en 2022 y $19.900$ millones en 2023, equivalentes a un $1,4\%$ del PIB en este último año. Con los Presupuestos de 2023 prorrogados en 2024, 2025 y 2026, la cuantía se ha mantenido en niveles similares; la cifra exacta de cada año debe verificarse en la liquidación correspondiente.

Debate sobre los gastos impropios: corrección de una anomalía u ocultación del déficit

Como se anticipó, la transferencia por gastos impropios ha abierto un debate que no es contable, sino de interpretación económica. Los defensores de la medida, sostienen que corrige una anomalía histórica. Durante décadas, el sistema contributivo financió, con cotizaciones, políticas que correspondían al Estado. Por eso su déficit no reflejaba un desequilibrio entre cotizaciones y pensiones contributivas, sino el coste de esas políticas ajenas. Una vez depurado el sistema de esos gastos, su saldo refleja mejor su situación real, que es mucho menos grave de lo que sugerían las cifras anteriores.

Los críticos, con los análisis de Fedea como referencia más citada, sostienen que la transferencia oculta el déficit estructural del sistema contributivo sin resolverlo. Su argumento tiene una base empírica llamativa. Pese a las transferencias por gastos impropios de casi $40.000$ millones de euros entre 2022 y 2023, el déficit de la Seguridad Social se mantuvo prácticamente estable en torno al $0,5\%$ del PIB en ambos años. Es decir, el sistema necesita, además de esas transferencias, seguir recurriendo al endeudamiento. 

GARCÍA DÍAZ (2024), en un trabajo de Fedea sobre la situación financiera del componente contributivo de las pensiones públicas, que incluye tanto las de la Seguridad Social como las de Clases Pasivas, subraya que el aumento de las transferencias del Estado ha reducido el contenido informativo del déficit contable oficial. Estima un déficit del componente contributivo muy superior al oficial una vez descontadas todas las aportaciones del Estado. \textcolor{redç}{La cifra exacta de ese déficit está pendiente de verificar, pero se sitúa en varios puntos de PIB.}

Las dos lecturas son compatibles con los mismos datos, porque la frontera entre lo propio y lo impropio del sistema no es técnica, sino normativa. Pero hay un hecho que ninguna de las dos discute. Las transferencias del Estado, cualquiera que sea su calificación, no mejoran la situación financiera de las Administraciones Públicas en su conjunto: trasladan el déficit de la Seguridad Social al Estado. Por tanto, la sostenibilidad del sistema de pensiones no puede evaluarse con el saldo de la Seguridad Social, sino con el gasto total en pensiones en relación con el PIB y con el conjunto de los recursos que lo financian, enfoque que abordaremos en el último bloque, y es también el de la regla de gasto de la reforma de 2023, que no mide el saldo del sistema, sino el gasto bruto en pensiones y los ingresos adicionales aprobados.

\subsubsection{Endeudamiento y las reservas}

\paragraph{El Fondo de Reserva: la hucha que se vació en la crisis}

La segunda recomendación del Pacto de Toledo (1995) proponía constituir reservas con los excedentes de los años de bonanza, y la Ley 28/2003, desarrollada por el Real Decreto 337/2004, reguló su funcionamiento. Estableció que los excedentes de ingresos de carácter contributivo resultantes de la liquidación de cada ejercicio se aplicarían prioritaria y mayoritariamente al Fondo, y que sus recursos se invertirían en activos financieros públicos. Nació con un propósito muy concreto: acumular, durante los años de superávit del sistema contributivo, un colchón financiero capaz de amortiguar los años de déficit que el propio ciclo económico y la evolución demográfica harían inevitables. 

El Fondo alcanzó su máximo histórico el 31 de diciembre de 2011, con $66.814,99$ millones de euros, equivalentes en torno al $6,2\%$ del PIB de ese año. Resultado de una década de crecimiento sostenido financiado con los excedentes de gestión de las entidades gestoras y con el rendimiento de sus inversiones, mayoritariamente en deuda pública. A partir de 2012, el Fondo entró en una fase de disposición sistemática y acelerada: $7.003$ millones retirados en 2012, $11.648$ en 2013, $15.300$ en 2014 y $13.250$ en 2015. Los gobiernos recurrieron a él para cubrir los desajustes crecientes entre los ingresos por cotizaciones y el gasto en pensiones.

Para hacer posibles esas disposiciones se suspendió el límite legal que impedía retirar anualmente más del $3\%$ del gasto en pensiones contributivas. Las disposiciones continuaron en 2016 y 2017, con cuantías elevadas. El deterioro fue tan rápido que a finales de 2019 el Fondo apenas conservaba entre $1.500$ y $2.150$ millones de euros, menos del $3\%$ de su máximo de 2011, y hubo que complementar sus últimas disposiciones con préstamos crecientes del Estado a la Tesorería General de la Seguridad Social. 

Aparentemente, es la historia de un instrumento diseñado para blindar el sistema frente a los ciclos económicos adversos, que se agotó, precisamente, ante el ciclo más adverso de su corta historia. Sin embargo, la frase admite dos lecturas opuestas. Según la primera, el Fondo cumplió su función. Un fondo de reserva de un sistema de reparto se acumula en las expansiones para utilizarse en las recesiones, y eso es exactamente lo que ocurrió. Sin el Fondo, la crisis de 2008-2013 habría exigido recortes de pensiones o subidas de cotizaciones más intensas en el peor momento del ciclo, con efectos procíclicos. Según la segunda, el Fondo se desvió de su finalidad. Su propósito no era solo cíclico, sino demográfico: acumular recursos para la jubilación del baby boom, que empieza precisamente en la década de 2020. Al utilizarlo para cubrir el déficit de la crisis, se consumió el ahorro destinado a un problema futuro para resolver uno presente, sin reformas que corrigieran el desequilibrio de fondo.\footnote{%
    Hay, además, una tercera observación, menos conocida y más importante desde la teoría: la naturaleza de los activos del Fondo. Durante la mayor parte de su historia, su cartera estuvo compuesta sobre todo por deuda pública española. Esto tiene dos consecuencias que la literatura sobre los fondos de reserva ha señalado, sobre todo en el debate estadounidense sobre el Social Security Trust Fund. La primera es que un fondo invertido en deuda del propio Estado no constituye una capitalización real desde el punto de vista macroeconómico. Sus activos son créditos contra el Tesoro, de modo que, cuando la Seguridad Social los vende para pagar pensiones, el Tesoro tiene que refinanciarlos emitiendo nueva deuda o subiendo impuestos. Es un traslado de recursos dentro del sector público, no un ahorro acumulado frente a la economía. La segunda consecuencia es que, como la deuda pública se mide en términos consolidados según los criterios de Maastricht, la deuda del Estado en manos del Fondo se descuenta del total. Cuando el Fondo se acumulaba, reducía la deuda consolidada de las Administraciones Públicas. Cuando se vació, la deuda consolidada aumentó en la misma cuantía. En términos de las cuentas públicas en su conjunto, el Fondo de Reserva funcionó, por tanto, como una forma de reducción temporal de la deuda pública en la fase expansiva, cuyo efecto se deshizo en la recesiva. Un fondo invertido en activos extranjeros diversificados, como los fondos AP suecos o el fondo soberano noruego, habría proporcionado en cambio una capitalización real, con la diversificación de riesgos que la teoría de MERTON y otros (2000) recomienda.
}

\paragraph{MEI: la nueva hucha}

La lección aprendida de la historia del Fondo de Reserva: el MEI. En lugar de acumular reservas de forma oportunista con los excedentes que vayan surgiendo, el mecanismo las genera mediante una cotización específica y permanente, diseñada desde el principio para no depender de la coyuntura económica del momento, con el objetivo de amortiguar el impacto financiero de la jubilación de la generación del baby boom. La reforma de 2023 estableció, en efecto, un periodo de acumulación hasta 2032 y un periodo de disposición desde 2033, con un límite anual que se expresa en porcentaje del PIB y cuyo valor exacto debe verificarse en la norma.

Los resultados de la nueva estrategia son visibles. En 2025, el Fondo de Reserva recibió dotaciones por $4.307$ millones de euros, de los que $3.970$ procedían del MEI, $319$ de los rendimientos de sus inversiones y $18$ de las dotaciones de las mutuas colaboradoras. Desde el inicio del mecanismo en 2023, sus aportaciones acumuladas ascienden a $9.765$ millones de euros. El Fondo cerró 2025 por encima de los $14.000$ millones de euros, su nivel más alto desde diciembre de 2017, y a finales de marzo de 2026 superaba los $15.267$ millones.

La estrategia plantea, sin embargo, una paradoja que sus críticos han subrayado. Mientras el Fondo de Reserva se llena con la cotización del MEI, la Seguridad Social se endeuda con el Estado para pagar las pensiones corrientes. Entre diciembre de 2022 y marzo de 2026, el Fondo aumentó en unos $13.100$ millones de euros, mientras la deuda de la Seguridad Social con el Estado aumentó en unos $30.000$ millones entre diciembre de 2022 y diciembre de 2025. Es decir, el sistema está ahorrando en la hucha con una mano mientras pide prestado con la otra, y lo que pide prestado es más del doble de lo que ahorra. Como el Fondo sigue invertido casi íntegramente en deuda pública española, el resultado neto es un circuito en el que el Estado presta a la Seguridad Social y la Seguridad Social presta, a través del Fondo, al Estado. 

Los defensores del mecanismo responden con tres argumentos. El primero es de compromiso: afectar una cotización a un fondo que no puede utilizarse hasta 2033 es un dispositivo que protege el ahorro destinado al baby boom frente a la tentación de gastarlo antes. (\textit{Ulises atado al mástil}). El segundo es de visibilidad política: una hucha creciente genera confianza en el sistema, lo que, según la paradoja de la redistribución, contribuye a su legitimidad. El tercero es que la deuda con el Estado es una deuda dentro del sector público, que desaparece en la consolidación, mientras que el Fondo es un activo real del sistema que podrá movilizarse cuando se necesite.\footnote{%
    La crítica replica que, desde el punto de vista de las Administraciones Públicas en su conjunto, lo relevante es si el sector público ahorra más en la fase previa a la jubilación del baby boom. Llenar un fondo mientras se endeuda al sistema en mayor cuantía no aumenta ese ahorro.
}

\paragraph{Préstamos del Estado: una deuda que no se devuelve}

La tercera fuente de financiación del sistema son los préstamos del Estado. El Fondo de Reserva ha sido sustituido en la práctica, no por una acumulación renovada de reservas, sino por el mecanismo de préstamo estatal permanente que hoy sostiene el equilibrio presupuestario formal de la Seguridad Social. 

Desde 2017, las leyes de presupuestos y, en los años de prórroga presupuestaria, las normas de autorización correspondientes han concedido a la Tesorería General préstamos anuales del Estado para cubrir el desfase entre sus ingresos y sus gastos. Se conceden sin intereses y a largo plazo, y en la práctica no se devuelven: cada año se concede un nuevo préstamo, y el principal acumulado crece.

El Tribunal de Cuentas, en su declaración sobre la cuenta general del Estado publicada en junio de 2026, cuantificó la situación a finales de 2024. La deuda de la Seguridad Social con el Estado por préstamos ascendía a $126.171$ millones de euros. Estaba compuesta por unos $109.000$ millones de préstamos concedidos desde 2017 y por unos $17.169$ millones de préstamos concedidos entre 1992 y 1999, que siguen sin devolverse más de un cuarto de siglo después.\footnote{%
    Los préstamos de 1992-1999 muestran que el recurso al endeudamiento con el Estado no es nuevo. Se utilizó ya en la crisis de los noventa, y su principal nunca se devolvió, ni siquiera en los años de superávit de 2000-2008, en los que los excedentes se destinaron al Fondo de Reserva en lugar de a amortizar esa deuda.
} La deuda total, incluidas otras obligaciones históricas, alcanzaba los 135.253 millones. El patrimonio neto de la Seguridad Social era negativo en $110.625$ millones de euros. A finales de 2025, la deuda superaba ya los $136.000$ millones de euros, con un aumento de unos $30.000$ millones desde diciembre de 2022. 

Habida cuenta de esto, ha recomendado sustituir los préstamos por transferencias, con el argumento de que la financiación mediante préstamos carece de racionalidad económico-financiera. Según su razonamiento, los préstamos proporcionan liquidez, pero no ingresos. Aumentan el pasivo de la Seguridad Social sin compensar los déficits acumulados, y su contabilización como deuda no refleja la verdadera naturaleza económica de unas aportaciones que el Estado nunca espera recuperar. El Ministerio de Hacienda ha defendido la financiación mediante préstamos por ser jurídicamente correcta y porque impone una disciplina presupuestaria: convertirlos en transferencias podría fomentar la irresponsabilidad fiscal y debilitar el control del gasto.

El debate entre el Tribunal de Cuentas y Hacienda ilustra con precisión la cuestión del perímetro. Desde el punto de vista de las Administraciones en su conjunto, la distinción entre préstamo y transferencia es irrelevante. Ni el préstamo ni la transferencia alteran el déficit consolidado, y la deuda de la Seguridad Social con el Estado desaparece en la consolidación. Pero desde el punto de vista de la señal que transmiten las cuentas, la distinción es decisiva. Con préstamos, la Seguridad Social aparece como una entidad con un patrimonio neto negativo creciente, cuyo desequilibrio queda visible en su balance. Con transferencias, ese desequilibrio desaparecería de sus cuentas y se trasladaría al Estado, donde quedaría diluido en el déficit general. La posición de Hacienda puede interpretarse, como una defensa de la afectación y de la disciplina que impone: mantener visible el déficit del sistema obliga a afrontarlo. La del Tribunal de Cuentas, como una defensa de la imagen fiel de las cuentas: si el Estado nunca va a recuperar esos recursos, contabilizarlos como préstamos es una ficción. Ambas posiciones tienen fundamento, y la elección entre ellas es, en último término, una elección sobre cómo se quiere que el sistema político perciba el desequilibrio del sistema de pensiones.


\subsection{El presupuesto del sistema y el gasto comparado}
Los dos subapartados anteriores examinaron la financiación desde la teoría y desde sus fuentes concretas. Este cierra el bloque con dos preguntas empíricas complementarias. La primera es cuánto gasta el sistema de protección social español, en qué y quién lo gasta, considerando las Administraciones Públicas en su conjunto y no solo el presupuesto de la Seguridad Social. La segunda es cómo se sitúa ese gasto en comparación con otros países, y qué nos dice esa comparación sobre la naturaleza del modelo español.

Si en el apartado anterior se analizó el sistema desde una perspectiva orgánica e institucional, aquí se trata de su clasificación económica y funcional, pero considerando los subsectores de la Administración como un todo consolidado. El objetivo es saber, con independencia de quién gasta qué, cuánto se gasta y en qué, cómo se estructura el presupuesto. 

Entrar en los debates doctrinales y en la comparación académica, en los pros y los contras, en la situación de cada país, su gasto por funciones y sus razones, sus objetivos, sus fortalezas y debilidades. La comparación concreta del sistema de pensiones se reserva para el bloque IV.

\subsubsection{Presupuesto: del sistema a las Administraciones Públicas consolidadas}
Como se estableció al definir el perímetro, el gasto en protección social puede medirse en tres niveles distintos. El primer nivel es el presupuesto del Sistema de la Seguridad Social en sentido estricto, que responde a la pregunta de cuánto ingresa y gasta la institución que gestiona las pensiones y las demás prestaciones económicas del sistema. El segundo es el subsector de Administraciones de Seguridad Social en contabilidad nacional, que añade el SEPE y el FOGASA y aplica los criterios del SEC 2010. El tercero es el gasto funcional consolidado de todas las Administraciones Públicas, clasificado según la COFOG. 

\paragraph{Presupuesto del Sistema de la Seguridad Social}

\textcolor{red}{[DOC] El documento describe la liquidación del presupuesto del Sistema de la Seguridad Social de 2021, que abarca a los trabajadores por cuenta ajena y por cuenta propia, a los socios trabajadores de cooperativas de trabajo asociado, a los estudiantes y a algunos funcionarios públicos. [DOC, nota] Recuerda que ese presupuesto no incluye el Sistema Nacional de Empleo ni el sanitario, que se describen en otros epígrafes. [DOC] Según el documento, en 2021 el sistema cerró con un déficit del 0,9 por ciento del PIB, porque los ingresos, un 15 por ciento del PIB, fueron inferiores a los gastos, un 15,9 por ciento. Las cotizaciones sociales alcanzaron el 10,9 por ciento del PIB y las transferencias del Estado, el 3 por ciento. Por el lado de los gastos, el 98 por ciento correspondía a prestaciones, principalmente del nivel contributivo, que suponían el 85 por ciento de las transferencias corrientes. Dentro de ellas destacaban las pensiones contributivas, un 11,5 por ciento del PIB, que pueden ser de jubilación, invalidez, viudedad, orfandad y en favor de familiares, y los subsidios y otras prestaciones, un 1,8 por ciento, principalmente la incapacidad temporal y las prestaciones de nacimiento y cuidado de menor y de riesgo durante el embarazo y la lactancia. El nivel no contributivo representaba el 1,1 por ciento del PIB y comprendía las pensiones no contributivas de invalidez y jubilación, los complementos a mínimos, el IMV y las prestaciones de protección familiar. [DOC, nota] El documento aclara que sus cifras corresponden a derechos reconocidos netos y obligaciones reconocidas netas de la liquidación del presupuesto.}

El sistema cerró 2025 con un déficit de $7.387$ millones de euros, equivalente al $0,4\%$ del PIB. Es el nivel más bajo desde 2011, frente a los $9.834$ millones de 2024 y al máximo del $1,7\%$ del PIB alcanzado en 2016. Los ingresos por cotizaciones sociales alcanzaron los $176.918$ millones de euros, un $6,9\%$ más que el año anterior. Las transferencias del Estado ascendieron a $52.990$ millones, un $10,1\%$ más. Y el MEI recaudó $4.935$ millones, que se destinaron al Fondo de Reserva. Por el lado del gasto, las obligaciones no financieras alcanzaron los $240.169$ millones de euros, un $6,5\%$ más. De ellas, $206.360$ millones correspondieron a pensiones y prestaciones contributivas, $18.992$ millones a prestaciones no contributivas y $18.413$ millones a incapacidad temporal. Con un PIB de 2025 en torno a $1,69$ billones de euros, los ingresos no financieros del sistema se situarían en torno al $13,7\%$ del PIB y sus gastos, en torno al $14,2\%$. 

Tres observaciones sobre estos datos conectan con el análisis de los apartados anteriores. La primera es que la reducción del déficit de 2025 se explica por el crecimiento de las cotizaciones y, en no menor medida, por el de las transferencias del Estado, que crecieron más deprisa que las cotizaciones. Como se estableció, esa parte de la mejora no reduce el déficit de las Administraciones Públicas, sino que lo traslada al Estado. La segunda es el crecimiento del gasto en incapacidad temporal, un $11,8\%$ en un solo año, casi el doble que el del gasto total. La tercera es el crecimiento de las prestaciones no contributivas, un $11,5\%$, impulsado por el IMV y la revalorización de las pensiones no contributivas por encima de la inflación en aplicación de la senda hacia el umbral de la pobreza.

Desde la perspectiva de la clasificación económica, el presupuesto del sistema tiene una estructura muy singular en el conjunto del sector público. Casi la totalidad del gasto, en torno al $98\%$, son transferencias corrientes a las familias, es decir, prestaciones. Los gastos de funcionamiento del propio sistema representan, por tanto, una fracción muy pequeña, confirmación empírica sobre los bajos costes administrativos de la gestión pública de las pensiones, gracias a las economías de escala.

Desde la perspectiva de la clasificación funcional del propio presupuesto de la Seguridad Social, el gasto se organiza en grandes áreas: prestaciones económicas, asistencia sanitaria residual (INGESA e ISM), servicios sociales (IMSERSO) y funciones de tesorería, informática y gestión. Las prestaciones económicas concentran la práctica totalidad del gasto.

#### El subsector de Administraciones de Seguridad Social

El segundo nivel de medición es el del subsector de Administraciones de Seguridad Social de la contabilidad nacional, que además del sistema incluye el SEPE y el FOGASA. Su saldo difiere del saldo presupuestario del sistema por dos razones. La primera es de perímetro: el SEPE tiene habitualmente superávit en las fases de bajo desempleo, porque la cotización por desempleo supera el gasto en prestaciones, y ese superávit compensa parcialmente el déficit del sistema de pensiones. La segunda es de criterio contable: la contabilidad nacional aplica el principio de devengo y otros criterios del SEC 2010, que pueden diferir de los criterios presupuestarios de derechos y obligaciones reconocidos. Por eso el déficit de las Administraciones de Seguridad Social en contabilidad nacional ha sido en los últimos años algo inferior al déficit presupuestario del sistema. 

El dato de 2025 en contabilidad nacional debe verificarse en la publicación de la IGAE y en la notificación del procedimiento de déficit excesivo, pues es el saldo relevante para las reglas fiscales europeas y para el plan fiscal-estructural.

\paragraph{Gasto funcional consolidado de las Administraciones Públicas}
Pasemos ahora al análisis funcional del gasto consolidado de las Administraciones Públicas (COFOG).\footnote{%
    Según la Clasificación de las Funciones de las Administraciones Públicas (COFOG), elaborada por las Naciones Unidas e integrada en el SEC 2010, el gasto público se distribuye en diez divisiones: servicios públicos generales, defensa, orden público, asuntos económicos, protección del medio ambiente, vivienda, sanidad, ocio y cultura, educación y protección social.
} Según los datos de la Intervención General analizados por Fedea, el gasto total de las Administraciones Públicas españolas ascendió en 2024 a unos $725.000$ millones de euros, el $45,5\%$ del PIB. 

Las pensiones absorbieron unos $206.000$ millones, el $28,4\%$ del gasto público total, y la sanidad, el $14,2\%$. En términos de PIB, las pensiones representan en torno al $13\%$. Según el Ministerio de Sanidad, el gasto sanitario público alcanzó en 2024 los $101.739$ millones de euros, el $6,4\%$ del PIB. 

La evolución de largo plazo muestra un cambio profundo en la composición del gasto público español. Entre 1995 y 2024, el gasto público total en relación con el PIB aumentó solo $1,3$ puntos, pero la protección social ganó $4,3$ puntos, de los que $3,3$ corresponden a las pensiones, y la sanidad, $1,2$ puntos. Ese aumento se compensó con reducciones en transporte ($1,5$ puntos), vivienda ($0,6$), defensa ($0,5$) y educación ($0,2$). En tres décadas, por tanto, el margen fiscal generado por el crecimiento económico y la reducción de otros gastos se ha destinado casi íntegramente a pensiones y sanidad. Como advierte Fedea, la llegada a la jubilación del baby boom reducirá todavía más el margen para otras políticas públicas.

Esta evolución tiene una lectura teórica precisa, que conecta con dos leyes clásicas de la hacienda pública que el opositor conoce de otros temas. La primera es la Ley de WAGNER, según la cual el gasto público tiende a crecer con el desarrollo económico, en particular el destinado a funciones sociales. La segunda es la enfermedad de costes de BAUMOL (1967), los servicios intensivos en trabajo con escaso crecimiento de la productividad, como la sanidad, la educación o los cuidados de larga duración, ven encarecerse su coste relativo con el tiempo, porque sus salarios crecen al ritmo de la productividad del conjunto de la economía mientras su propia productividad crece menos.

Las pensiones no están sujetas a la enfermedad de BAUMOL en sentido estricto, porque son transferencias monetarias y no servicios. Pero su crecimiento se explica por la demografía y por el efecto sustitución: los nuevos pensionistas, con carreras más largas y salarios más altos, generan pensiones superiores a las de los pensionistas que fallecen.

\subsubsection{Gasto por subsectores: quién gasta en qué}
La distribución del gasto social entre subsectores responde a la arquitectura institucional descrita y al reparto constitucional de competencias, y conviene explicitarla.

Las pensiones contributivas y no contributivas del sistema de Seguridad Social, la incapacidad temporal, las prestaciones de nacimiento y el IMV se imputan al subsector de Administraciones de Seguridad Social, aunque una parte creciente se financie con transferencias del Estado. Las pensiones de Clases Pasivas se imputan a la Administración Central, porque se pagan directamente con el presupuesto del Estado. Las prestaciones por desempleo se imputan a las Administraciones de Seguridad Social a través del SEPE. Las políticas activas de empleo, en cambio, las ejecutan mayoritariamente las Comunidades Autónomas. El gasto sanitario público lo ejecutan en su inmensa mayoría las Comunidades Autónomas, con una participación menor de la Administración Central, las mutualidades de funcionarios y las corporaciones locales. Del gasto sanitario total, entre el $70\%$ y el $75\%$ corresponde a las Administraciones Públicas y el resto al sector privado. El gasto público en dependencia ascendió en 2021 al $0,8\%$ del PIB: un $79\%$ lo aportaron las Comunidades Autónomas y el $21\%$ restante la Administración General del Estado, a lo que se suman las aportaciones de los usuarios. 

Dicho esto, esta distribución tiene una consecuencia importante para el análisis de la sostenibilidad, pues el gasto sensible al envejecimiento se reparte entre subsectores que se financian de forma muy distinta. Las pensiones se financian con cotizaciones, transferencias del Estado y préstamos, dentro del subsector de Seguridad Social y de la Administración Central. La sanidad y la dependencia, en cambio, se financian dentro del sistema de financiación autonómica, con impuestos cedidos y transferencias del Estado. Por eso la presión del envejecimiento sobre la sanidad y la dependencia no aparece en las cuentas de la Seguridad Social, sino en las de las Comunidades Autónomas. 

Un análisis de sostenibilidad que se limite al sistema de pensiones infravalora el impacto fiscal total del envejecimiento. Es la razón por la que el Ageing Report y la AIReF proyectan conjuntamente el gasto en pensiones, sanidad, cuidados de larga duración y educación.

\subsubsection{Gasto social en perspectiva comparada}
Antes de comparar, hay que resolver un problema metodológico que invalida muchas comparaciones del debate público. 

El documento utiliza la COFOG y menciona, en nota, el Sistema Europeo de Estadísticas Integradas de Protección Social (ESSPROS), desarrollado por Eurostat y los Estados de la Unión para ofrecer un marco común sobre las prestaciones sociales a los hogares y su financiación. Las dos fuentes miden cosas distintas. La COFOG mide el gasto público por funciones, incluido el gasto de funcionamiento de las administraciones. ESSPROS mide las prestaciones sociales, públicas y privadas, que responden a riesgos sociales definidos (vejez, supervivencia, enfermedad, invalidez, familia, desempleo, vivienda y exclusión), con independencia de quién las pague, siempre que no exista una contrapartida simultánea y equivalente. Una tercera fuente, la base de datos de gasto social de la OCDE (SOCX), añade una dimensión decisiva para la comparación: la distinción entre gasto bruto y gasto neto.

ADEMA y LADAIQUE (2009), en los trabajos metodológicos de la OCDE sobre el gasto social neto, mostraron que las comparaciones del gasto social bruto exageran las diferencias entre países por tres razones. La primera es que los países con Estados del bienestar grandes recuperan una parte significativa de su gasto social mediante los impuestos directos e indirectos que gravan las prestaciones: una pensión sujeta al impuesto sobre la renta cuesta al Estado menos que su importe bruto, y el pensionista paga el IVA sobre lo que consume con ella. La segunda se deriva de los gastos fiscales, que algunos países contemplan como beneficios fiscales, que tienen una finalidad social (deducciones por hijos, por planes de pensiones o por seguros de salud) pero no aparecen como gasto, pero tienen el mismo efecto pues se dejan de ingresar. La tercera es que en algunos países el gasto social privado obligatorio o incentivado es muy importante (planes de pensiones de empresa, seguros de salud privados). 

Cuando se tienen en cuenta estos tres factores, las diferencias entre países se reducen drásticamente. Estados Unidos, cuyo gasto social público bruto es muy inferior al de los países europeos, se sitúa en términos de gasto social neto total a un nivel comparable al de muchos de ellos. La diferencia no está en cuánto gasta la sociedad en protección social, sino en quién lo gasta y cómo se distribuye.\footnote{%
    La idea tiene un antecedente clásico en TITMUSS (1958), que distinguió tres formas de bienestar. El bienestar social consiste en las prestaciones y servicios públicos directos. El bienestar fiscal consiste en los beneficios fiscales con finalidad social. El bienestar ocupacional consiste en las prestaciones que las empresas conceden a sus trabajadores. Las tres tienen efectos distributivos, pero solo la primera es visible en las estadísticas de gasto público. Las dos últimas suelen beneficiar desproporcionadamente a las clases medias y altas, que tienen mayor renta sujeta a impuestos y empleos más estables con mejores beneficios de empresa. HOWARD (1997) desarrolló esta idea para Estados Unidos con el concepto de Estado del bienestar oculto (\textit{hidden welfare state}), que se construye mediante el código tributario en lugar del presupuesto de gasto.
} 

La lección para el análisis comparado es que un Estado del bienestar pequeño en términos de gasto público puede ser grande en términos de bienestar fiscal y ocupacional, y que su efecto distributivo puede ser muy distinto del que sugiere su tamaño. 

\paragraph{España: un gasto en torno a la media, con una composición singular}
En el caso español, el bienestar fiscal y el ocupacional son relativamente pequeños, lo que hace que el gasto público bruto sea un indicador más fiel del esfuerzo total que en los países anglosajones. Así que, ¿cómo se sitúa España en el panorama internacional?

En 2024, el gasto público total de la Unión Europea ascendió al $49,1\%$ del PIB, del que la protección social representó el $40\%$ y la sanidad el $15\%$.\footnote{
De forma más desagregada encontramos los datos de 2023. El gasto público en protección social de la Unión se situó en el $19,2\%$ del PIB. Su distribución fue la siguiente: vejez, $10,4$ puntos; enfermedad e invalidez, $2,7$; familia e hijos, $1,9$; supervivencia, $1,4$; desempleo, $1,2$; exclusión social, $1,0$; y vivienda, $0,3$. El rango entre países iba del $8,1$ por ciento de Irlanda al $25,7$ por ciento de Finlandia.
} Por tanto, España, con un gasto público total del $45,5\%$ del PIB en 2024, se sitúa por debajo de la media de la Unión en gasto total. 

Su gasto en protección social, en torno al $18-19\%$ del PIB, y en sanidad, en torno al $6,4-6,6\%$, se sitúa cerca de la media en el primer caso y ligeramente por debajo en el segundo. \textcolor{red}{Las cifras exactas de la COFOG española de 2024 deben verificarse en la última publicación de la IGAE y de Eurostat.} Por tanto, a priori no parece que se lógica la magnitud mediática y política del gasto en protección social en España, ¿no?

Lo verdaderamente distintivo del caso español no es el nivel del gasto, sino su composición. Tiene tres rasgos que la literatura comparada ha destacado reiteradamente y que conectan con el régimen de bienestar descrito. El primero es un peso elevado del gasto en vejez y supervivencia. Las pensiones de jubilación y de viudedad absorben una proporción del gasto social superior a la media europea, que reflejan la herencia del modelo del varón sustentador y la tardía incorporación de la mujer al mercado de trabajo. El segundo es un gasto en desempleo elevado y muy procíclico, que refleja un desempleo estructural persistentemente superior a la media europea y una protección concentrada en el gasto pasivo más que en las políticas activas. El tercero es un gasto bajo en familia e infancia, en vivienda y en exclusión social, funciones que la literatura de la inversión social considera clave para afrontar los nuevos riesgos sociales.

Esta composición ha sido analizada por LYNCH (2006), que clasificó los Estados del bienestar según su orientación por edades. Distingue Estados orientados a los mayores (elderly-oriented), que concentran su gasto en pensiones y en sanidad de los mayores, y Estados orientados a los jóvenes (youth-oriented), que dedican más recursos a la educación, la familia y el empleo. Según su análisis, los países del sur de Europa, con Italia, Grecia y España como casos destacados, figuran entre los más orientados a los mayores. LYNCH atribuye esa orientación no a la demografía, sino a la estructura institucional de los Estados del bienestar: los sistemas bismarckianos, que ligan la protección al empleo y a la contribución, tienden a proteger mejor a quienes han completado una carrera laboral, es decir, a los mayores. La combinación de sistemas de protección fragmentados con la competencia política clientelar, solo hace que reforzar esa tendencia. VANHUYSSE (2013), con un indicador de sesgo hacia los mayores en el gasto social, que relaciona el gasto per cápita en mayores con el gasto per cápita en menores teniendo en cuenta la estructura demográfica, llegó a conclusiones similares, con los países mediterráneos entre los más sesgados.

\paragraph{Eficacia del gasto: redistribución y reducción de la pobreza}
Comparar cuánto se gasta es insuficiente sin preguntar qué se consigue con ese gasto. El indicador más utilizado es la reducción de la pobreza que producen las transferencias sociales: la diferencia entre la tasa de riesgo de pobreza antes y después de las transferencias. 

Según los datos de Eurostat, el sistema español reduce la pobreza de forma comparable a la media europea cuando se incluyen las pensiones, pero de forma notablemente inferior cuando se excluyen. Las transferencias sociales distintas de las pensiones (desempleo, familia, exclusión social, vivienda) tienen en España una capacidad de reducir la pobreza de las más bajas de la Unión. La consecuencia más visible es la pobreza infantil. La tasa de riesgo de pobreza o exclusión social de los menores en España se sitúa de forma persistente entre las más altas de la Unión, en contraste con una tasa de pobreza de los mayores de $65$ años que se ha reducido significativamente en las dos últimas décadas. Lo que se debe en buena medida a que las pensiones (mayores) han mantenido su poder adquisitivo mientras los salarios de los jóvenes se estancaban (ahora, con hijos).

Un diagnóstico que, de nuevo, enlaza con el debate teórico sobre la paradoja de la redistribución. El gasto social español es relativamente eficaz para los mayores, protegidos por un sistema contributivo con amplio apoyo político. Es poco eficaz para los niños y los hogares en edad de trabajar, que dependen de prestaciones focalizadas, fragmentadas y con baja utilización.

\paragraph{Por qué gastan distinto los países: las explicaciones de la literatura}
Ahora faltaría contestar a la última pregunta: ¿por qué los países difieren tanto en el nivel y la composición de su gasto social? La literatura ha ofrecido tres familias de explicaciones, no excluyentes.

La primera es institucional y de poder, y es la de ESPING-ANDERSEN ya examinada. Como breve recordatorio, el nivel y la composición del gasto social reflejan la movilización política de las clases trabajadoras y sus alianzas con las clases medias. Es la teoría de los recursos de poder (KORPI, 1983), según la cual los países con sindicatos fuertes y partidos socialdemócratas en el gobierno durante largos periodos desarrollaron Estados del bienestar más amplios y universales.

La segunda es electoral e institucional. IVERSEN y SOSKICE (2006) mostraron que los países con sistemas electorales proporcionales redistribuyen más que los de sistemas mayoritarios, porque la representación proporcional favorece coaliciones de centro-izquierda entre los partidos que representan a las clases medias y a las bajas. En los sistemas mayoritarios, en cambio, las clases medias tienden a aliarse con los partidos que representan a las clases altas, por temor a ser gravadas en beneficio de las bajas.

La tercera es cultural y de creencias. ALESINA y GLAESER (2004), atribuyeron la diferencia entre Estados Unidos y Europa a dos factores. El primero es la mayor heterogeneidad étnica estadounidense, que reduce la solidaridad entre grupos. El segundo son las distintas creencias sobre las causas de la pobreza. ALESINA y ANGELETOS (2005) formalizaron esta segunda idea. En las sociedades que creen que el éxito económico depende sobre todo del esfuerzo, la redistribución es baja, porque se ve como un castigo; en las que creen que depende sobre todo de la suerte, es alta, porque se considera que corrige una arbitrariedad. Además, debemos tener en cuenta que estas creencias se autorrefuerzan: en una sociedad con baja redistribución, el esfuerzo es más necesario y efectivamente determina más el éxito, lo que confirma la creencia inicial. Es un modelo de equilibrios múltiples que explica por qué sociedades con niveles de desarrollo similares pueden mantener Estados del bienestar muy distintos de forma estable.


\paragraph{Grupos de países y sus modelos}
Normalmente, se suelen identificar tres grupos de países según su gasto social en sanidad y protección social. Los de \textbf{gasto alto}, superior al $25\%$ del PIB, eran Francia, Finlandia, Dinamarca, Austria, Noruega, Italia, Alemania, Bélgica, Suecia y Grecia; los países nórdicos y algunos de Europa continental y del sur lideraban los niveles de protección social. Los de \textbf{gasto medio}, entre el $20\%$ y el $25\%$, eran España, Luxemburgo, Portugal, Países Bajos, Eslovenia, Polonia, Croacia, República Checa y Eslovaquia; España estaría en este grupo, junto con algunos países de Europa central y oriental. Los de \textbf{gasto bajo}, inferior al $20\%$, eran Estonia, Lituania, Islandia, Hungría, Rumanía, Letonia, Bulgaria, Chipre, Malta, Suiza e Irlanda. La clasificación es útil como primera aproximación, pero agrupa países con lógicas muy diferentes y oculta más de lo que muestra. Con las herramientas de este apartado puede reinterpretarse.

En el grupo de gasto alto conviven los países nórdicos, con gasto universal financiado mayoritariamente con impuestos, fuerte orientación a los servicios y a la inversión social, alta participación laboral femenina y baja pobreza infantil, y los países continentales y mediterráneos, con gasto contributivo proporcional al salario y fuerte orientación a las pensiones. Italia y Grecia están en ese grupo sobre todo por su gasto en pensiones, que figura entre los más altos del mundo, y no por la amplitud de su protección a otros riesgos. La fortaleza del modelo nórdico es su eficacia redistributiva y su compatibilidad con altas tasas de empleo; su debilidad, la elevada presión fiscal que exige y su dependencia de un alto nivel de confianza social. La fortaleza del modelo continental es la protección del nivel de vida de los trabajadores con carreras completas; su debilidad, el dualismo y el coste en términos de empleo de sus elevadas cotizaciones. El modelo mediterráneo combina la debilidad de ambos: un gasto elevado en pensiones con una protección débil de los jóvenes y de los hogares en edad de trabajar.

En el grupo de gasto medio, España comparte posición con países de lógicas también distintas. Los Países Bajos combinan un gasto público moderado con un enorme segundo pilar de pensiones de capitalización, lo que los situaría entre los de mayor gasto si se midiera el gasto neto total. Varios países de Europa central y oriental reformaron sus sistemas en los años noventa con pilares de capitalización, en parte revertidos después. La posición de España en ese grupo refleja, por tanto, algo distinto que la de sus compañeros: un gasto público moderado en el agregado, pero sin segundo pilar, y con una composición sesgada hacia las pensiones públicas, a costa de otro tipo de prestaciones dirigidas a personas en edad de trabajar.

En el grupo de gasto bajo, Suiza tiene un sistema de pensiones de tres pilares con un segundo pilar obligatorio muy desarrollado, que no aparece en el gasto público. Irlanda está distorsionada por su PIB.\footnote{%
    \textbf{NOTA.-} El gasto de Irlanda, el más bajo de la Unión en relación con el PIB, está fuertemente distorsionado a la baja porque su PIB está inflado por los beneficios contabilizados en el país por las grandes multinacionales tecnológicas y farmacéuticas, que no corresponden a la renta de los residentes. Por eso las autoridades irlandesas utilizan un indicador alternativo, la renta nacional bruta modificada (GNI). Medido con él, el gasto social irlandés se acerca mucho más a la media europea. Es un recordatorio de que el denominador de las ratios importa tanto como el numerador.
} Y los países bálticos y otros del este tienen Estados del bienestar genuinamente reducidos, de corte liberal, con baja presión fiscal y elevadas tasas de pobreza.

La conclusión que se deriva es la que anticipaba la introducción: el nivel agregado del gasto social español, en torno a la media europea, no es un buen indicador de la naturaleza del modelo. Lo que caracteriza a España es un Estado del bienestar de tamaño intermedio, centrado en las pensiones y en la protección de los trabajadores con carreras completas, con una capacidad limitada de redistribución fuera de la vejez y sin los pilares complementarios que en otros países de gasto público similar amplían la protección efectiva. Esa configuración es la que enfrenta el envejecimiento del baby boom. Como vimos, concentra el riesgo en la masa salarial española y en el presupuesto público, con un margen fiscal que se ha ido estrechando a medida que las pensiones y la sanidad absorbían el crecimiento del gasto.







\clearpage
\section{Bloque III. Las prestaciones del sistema}

La Seguridad Social es pieza fundamental del Estado del bienestar en España y debe garantizar a los ciudadanos su protección y la de sus familias ante situaciones de vejez, enfermedad, desempleo y otras carencias sociales que, en el transcurso de la vida, requieran de ayuda. 

Los bloques anteriores establecieron por qué existe el sistema, cómo ha evolucionado, cómo se regula y organiza y cómo se financia. Este bloque examina qué protege y cómo: las prestaciones concretas, sus requisitos, su cálculo y sus efectos económicos. 

La exposición sigue un criterio de importancia económica, no de orden legal. Empieza por las pensiones, que concentran la mayor parte del gasto y son el objeto directo del título del tema. Continúa con las prestaciones económicas de corto plazo, entre ellas la incapacidad temporal, cuyo crecimiento reciente es uno de los fenómenos más relevantes del sistema. Y termina con los tres sistemas de protección que, estando fuera del sistema de Seguridad Social en sentido orgánico, forman parte de la protección social en sentido funcional: el desempleo, la sanidad y la dependencia. En estos tres últimos, no se repite la plantilla jurídica del documento. 

Cada sistema se presenta desde su lógica económica propia: el seguro de desempleo desde la teoría del seguro óptimo, la sanidad desde la economía de la salud y la dependencia desde el problema de los cuidados de larga duración.

\subsection{Las pensiones del sistema}
El sistema de pensiones es la parte más importante del Sistema de la Seguridad Social. Cubre un conjunto de contingencias relacionadas con el envejecimiento (jubilación), el fallecimiento (viudedad, orfandad y en favor de familiares) y la enfermedad (incapacidad permanente). 

Tiene dos modalidades. La contributiva, obligatoria y financiada con cotizaciones, determina sus prestaciones según las contribuciones realizadas durante la vida laboral, en función de unas bases y unos tipos de cotización. La no contributiva o asistencial, financiada con impuestos, cubre a quienes, estando en situación de necesidad, no tienen rentas o patrimonio suficientes ni han cotizado el tiempo necesario para acceder a las prestaciones contributivas. 

\textcolor{blue}{A septiembre de 2026}, el sistema abonó $10,5$ millones de pensiones contributivas a más de $9,5$ millones de pensionistas, con una nómina mensual de unos $14.500$ millones de euros, un $6,3\%$ más que un año antes. La pensión media del sistema alcanzó los $1.374,60$ euros mensuales, un $4,6\%$ más que en septiembre de 2025, y la pensión media de jubilación, los $1.576,10$ euros.\footnote{
\textbf{NOTA.-} No todas las pensiones públicas están integradas en el Sistema de la Seguridad Social. Existe el Régimen de Clases Pasivas del Estado, que provee diversos tipos de pensiones, principalmente de jubilación, a funcionarios de carrera de las Administraciones del Estado y a militares de carrera.
}

Las diferencias entre regímenes son muy marcadas: la pensión media de jubilación del Régimen General era de $1.734,70$ euros, y la del régimen de autónomos, de $1.062,80$ euros. Esa brecha es el resultado acumulado de décadas de cotización por la base mínima ya examinadas.  Tres cuartas partes de la nómina, unos $10.650$ millones de euros mensuales, corresponden a las pensiones de jubilación; unos $2.290$ millones a las de supervivencia y unos $1.330$ millones a las de incapacidad permanente. La estructura del gasto no ha cambiado sustancialmente desde 2021, pero su volumen ha crecido con intensidad, impulsado por la revalorización con el IPC en años de alta inflación, por el aumento del número de pensiones y por el efecto sustitución: los nuevos jubilados, con carreras más largas y salarios más altos, generan pensiones muy superiores a las de los pensionistas que fallecen.


\subsubsection{Pensión de jubilación contributiva}
El carácter contributivo del sistema de pensiones significa que las prestaciones económicas se calculan teniendo en cuenta las aportaciones realizadas a lo largo de la vida laboral. La cuantía de la pensión de jubilación se determina aplicando a la base reguladora el porcentaje que corresponde según el número de años cotizados, aumentado o disminuido según el trabajador haya prolongado su vida laboral o adelantado su jubilación. 

La pensión de jubilación es, por tanto, el resultado de tres elementos: el \textbf{acceso} (cuándo y con qué requisitos se puede jubilar un trabajador), la \textbf{base reguladora} (qué salario de referencia se toma) y el \textbf{porcentaje} (qué proporción de ese salario se reconoce). A ellos se añaden los \textbf{límites} (pensiones máxima y mínima), las \textbf{modalidades de acceso} (anticipada, demorada, activa, parcial) y el complemento para la reducción de la brecha de género.

\begin{specialblock}{\textbf{Nivel complementario}: Planes y Fondos de previsión privados (Segundo Pilar)}
    El art. $41$ de la Constitución declara libres la asistencia y las prestaciones complementarias, de modo que existe un nivel complementario de protección social que puede ser de carácter privado. Ese nivel comprende las mejoras voluntarias de las prestaciones públicas pactadas en la negociación colectiva o concedidas por las empresas, los planes de pensiones de empleo, promovidos por las empresas para sus trabajadores, los planes individuales, contratados libremente por los particulares, y los planes de previsión asegurados y otros seguros de ahorro. El peso de este segundo pilar en España es muy pequeño, un $0,4\%$ del PIB en términos de gasto en prestaciones. Los activos de los fondos de pensiones equivalen a algo más del $10\%$ del PIB, muy por debajo de la media de la OCDE, que supera el $60\%$. Países como Suecia, los Países Bajos o el Reino Unido tienen sistemas privados de pensiones cuyo gasto supera el $3\%$ del PIB.

    El escaso desarrollo tiene explicaciones que conectan con la teoría. La primera es la elevada tasa de reemplazo del sistema público, que reduce la necesidad percibida de ahorro complementario. La segunda es la preferencia por la vivienda como forma de ahorro para la vejez: España tiene una de las tasas de vivienda en propiedad más altas de Europa, y la vivienda en propiedad funciona como un pilar informal de la jubilación, el pilar cuatro de HOLZMANN y HINZ. La tercera es la estructura empresarial, dominada por pymes con escasa capacidad para promover planes de empleo. La cuarta es la ausencia de mecanismos de adhesión automática, que la teoría conductual identifica como decisivos para la participación.
    
    La Ley 12/2022, de regulación para el impulso de los planes de pensiones de empleo, intentó corregir esta situación.\footnote{%
        Ley 12/2022 y límites de aportación. Antes de 2021, el límite general de aportación con derecho a reducción era de 8.000 euros anuales, sin distinción entre planes individuales y de empleo. Las leyes de presupuestos para 2021 y 2022 lo redujeron a 2.000 y a 1.500 euros para los planes individuales, y la Ley 12/2022 reforzó los planes de empleo. El objetivo era desplazar el ahorro complementario desde los planes individuales, regresivos, hacia los de empleo, de alcance más amplio.
    } Creó los fondos de pensiones de empleo de promoción pública, gestionados por entidades privadas seleccionadas por concurso y supervisadas por una comisión pública, que ofrecen planes de bajo coste a las empresas, en particular a las pymes, y a los autónomos. Creó también los planes de pensiones de empleo simplificados, que pueden promoverse en el marco de la negociación colectiva sectorial. Y reforzó los incentivos fiscales a las aportaciones empresariales. Al mismo tiempo, el legislador redujo drásticamente el límite de las aportaciones a los planes individuales con derecho a reducción en el impuesto sobre la renta: de $8.000$ euros anuales a $2.000$ en 2021 y a $1.500$ en 2022. Mantuvo un límite conjunto más elevado cuando se suman las contribuciones empresariales a planes de empleo.
    
    Esta combinación responde a un diagnóstico que conecta con la crítica de TITMUSS y de HOWARD al bienestar fiscal. Las reducciones fiscales por aportaciones a planes individuales benefician desproporcionadamente a los contribuyentes de renta alta, que tienen mayor capacidad de ahorro y un tipo marginal más elevado, de modo que constituyen un gasto fiscal regresivo. Sus defensores responden que el tratamiento fiscal de los planes de pensiones sigue el modelo de diferimiento (aportaciones deducibles, rendimientos exentos y prestaciones gravadas), que no es un beneficio fiscal en sentido estricto, sino una forma de gravar el consumo en lugar de la renta. El diferimiento puede generar un beneficio real si el tipo marginal en la jubilación es inferior al del periodo de aportación, lo que es habitual.
    
    Los resultados de la Ley 12/2022 han sido, hasta ahora, muy limitados. Los fondos de promoción pública han tenido una implantación lenta, y la reducción de los incentivos a los planes individuales ha provocado una caída de las aportaciones a estos, sin un aumento equivalente de las aportaciones a planes de empleo. La cifra exacta de partícipes y patrimonio en los fondos de promoción pública debe verificarse. Un contraste revelador procede del propio territorio español. En el País Vasco, las entidades de previsión social voluntaria (EPSV), con un régimen fiscal propio en virtud del concierto económico, han desarrollado un segundo pilar mucho más amplio que en el resto de España. Destaca Geroa, la entidad de previsión sectorial de Guipúzcoa creada por la negociación colectiva en los años noventa, que cubre a buena parte de los asalariados de ese territorio. 
    
    Desde la perspectiva de la teoría de la financiación, el escaso desarrollo del segundo pilar tiene una consecuencia que conviene recordar. Los españoles tienen una cartera de pensiones concentrada casi exclusivamente en el reparto público, es decir, en el riesgo de la masa salarial española, que es precisamente el riesgo más expuesto al envejecimiento demográfico del país. La cuestión es si puede desarrollarse sin mandato, confiando solo en la negociación colectiva y en los incentivos fiscales, o si exige algún mecanismo de adhesión automática como los del Reino Unido o Nueva Zelanda.
\end{specialblock}

\paragraph{Acceso: la edad y la carencia}
En 2026, la edad ordinaria de jubilación es de $66$ años y $10$ meses para quienes acrediten menos de $38$ años y $3$ meses de cotización, y de $65$ años para quienes acrediten $38$ años y $3$ meses o más.\footnote{%
    Edad ordinaria de jubilación (Ley 27/2011). Hasta 2012, la edad ordinaria de jubilación era de 65 años para todos los trabajadores. La Ley 27/2011 estableció su aumento gradual hasta los 67 años, que se completará en 2027, con la excepción de quienes acrediten carreras largas, que podrán seguir jubilándose a los 65 años. La reforma introdujo un aumento escalonado de la edad legal de jubilación desde los 65 hasta los 67 años hasta 2027. La reforma pretendía responder al aumento de la esperanza de vida y retrasar la edad efectiva de jubilación.
}

\textcolor{red}{Los trabajadores con carreras superiores a $38,5$ años podrán seguir jubilándose a los $65$ años incluso después de 2027. No obstante, el umbral de años de cotización que permite jubilarse a esta edad también aumenta gradualmente, y en 2027 será de $38$ años y $6$ meses.} 

El acceso exige, además, un periodo mínimo de cotización (carencia genérica) de $15$ años, de los que al menos $2$ deben estar comprendidos en los $15$ años inmediatamente anteriores al momento de causar el derecho (carencia específica), y estar en alta o en situación asimilada al alta. 

Como ya vimos, la edad de jubilación no tiene un fundamento natural, sino convencional: la edad de $65$ años se fijó cuando pocos trabajadores la alcanzaban. Hoy, la esperanza de vida a los $65$ años en España supera los $20$ años, y es de las más altas del mundo, especialmente en las mujeres. La doble edad de jubilación, que distingue según la longitud de la carrera, tiene una justificación de equidad. Quien empezó a trabajar joven, normalmente en ocupaciones manuales y con menor esperanza de vida, puede jubilarse antes que quien empezó tarde tras una larga formación. Es una forma rudimentaria de atender al gradiente de longevidad, que hace que los trabajadores de menor cualificación cobren la pensión durante menos años.

\paragraph{Base reguladora: qué salario de referencia}
En 2026, la base reguladora de la pensión de jubilación puede calcularse de dos formas, y se aplica la más favorable para el trabajador.

La primera, tradicional, es el cociente de dividir por $350$ la suma de las bases de cotización de los $300$ meses ($25$ años) inmediatamente anteriores al hecho causante. La segunda, introducida por la reforma de 2023 con un periodo transitorio, consiste en tomar, dentro de los $304$ meses anteriores, las $302$ bases más altas y dividir su suma por $352,33$.\footnote{%
    Periodo de cálculo de la base reguladora. El periodo de cálculo ha ido creciendo en sucesivas reformas: dos años hasta 1985, ocho años con la Ley 26/1985, quince años con la Ley 24/1997 y veinticinco años con la Ley 27/2011, con un periodo transitorio entre 2013 y 2022. La Ley 27/2011 introdujo un incremento progresivo, hasta 2022, del periodo considerado para calcular la base reguladora, desde los 15 hasta los 25 años. El Real Decreto-ley 2/2023 estableció la fórmula dual que culminará en 2044 con un periodo de 29 años con exclusión de los 24 peores meses. Entre 2023 y 2025 se mantuvo inalterada la fórmula de los últimos 25 años. Cada ampliación perseguía reforzar la contributividad y contener el gasto; la de 2023 añadía la protección de los trabajadores con finales de carrera interrumpidos.
} El calendario transitorio de esta segunda fórmula es simple. A partir de 2044, el periodo quedará fijado en los últimos $29$ años de cotización, excluyendo los $24$ meses más desfavorables. Entre 2026 y 2040 se toma la opción más favorable entre los últimos $25$ años y un periodo creciente que permite excluir un número creciente de meses: en 2028, los últimos $26$ años excluyendo los $6$ peores meses, y en 2034, los últimos $28$ años excluyendo los $18$ peores. En ambas fórmulas, las bases de los $24$ meses más recientes se toman por su valor nominal, y las anteriores se actualizan con la evolución del IPC hasta el mes anterior a ese periodo, para que la inflación no erosione el salario de referencia. El divisor ($350$ en la fórmula tradicional) es mayor que el número de meses ($300$) porque la pensión se paga en catorce pagas anuales: $300$ meses divididos por $12$ y multiplicados por $14$ dan $350$.





La longitud del periodo de cálculo es una de las variables más importantes del diseño del sistema, y su evolución histórica responde a una lógica precisa. Cuanto más largo es el periodo, más contributivo es el sistema: la pensión refleja mejor el esfuerzo de toda la carrera y no solo el de sus últimos años. Esto tiene tres efectos. El primero es reducir el comportamiento estratégico. Con un periodo corto, era rentable cotizar poco durante la mayor parte de la carrera y aumentar la cotización en los últimos años, como hacían muchos autónomos; con un periodo largo, esa estrategia pierde rentabilidad. El segundo es reducir la pensión media para la mayoría de los trabajadores, cuyos salarios crecen a lo largo de la carrera, porque se incorporan años de salarios más bajos. Es la razón financiera de fondo de todas las ampliaciones. El tercero es redistribuir entre perfiles de carrera: perjudica a quienes tienen carreras crecientes, típicamente los más cualificados, y favorece a quienes sufren caídas de ingresos al final de la carrera, como los desempleados de larga duración mayores de $50$ años, cuyas últimas bases son bajas.\footnote{%
    La opción dual de 2023, con la posibilidad de excluir los peores meses, responde precisamente a este último problema. La crisis de 2008-2013 dejó a muchos trabajadores mayores con largos periodos de desempleo o de cotización reducida al final de su carrera; un periodo de cálculo más largo, con exclusión de los peores meses, les permite diluir esos años malos. La reforma combina así dos objetivos que en principio se oponen: reforzar la contributividad, alargando el periodo, y proteger a los trabajadores con finales de carrera interrumpidos, permitiendo excluir los peores meses. Por eso es difícil calificarla como restrictiva o expansiva: su efecto depende del perfil de carrera de cada trabajador.
}


\begin{specialblock}{\textbf{Problema:} Integración de lagunas}
    Un problema técnico con grandes consecuencias es qué ocurre con los meses del periodo de cálculo en los que no existió obligación de cotizar, por ejemplo por desempleo sin prestación o por inactividad. Son las llamadas lagunas de cotización. Se definen como periodos en los que no existe la obligación de cotizar y que se rellenan, en el Régimen General, para el cálculo de la pensión. La regla general vigente integra las primeras $48$ mensualidades con lagunas con la base mínima de cada momento, y el resto, con el $50$ por ciento de esa base mínima. La reforma de 2023 mejoró la integración para las mujeres, y en determinadas condiciones para los hombres, mientras la brecha de género de las pensiones supere el $5\%$ por ciento: las mensualidades $49$ a $60$ se integran con el $100\%$ de la base mínima y las $61$ a $84$, con el $80\%$.\footnote{%
        Antes de 2023, las primeras 48 mensualidades se integraban con la base mínima y el resto con el 50 por ciento de la base mínima, sin distinción. La reforma estableció mejoras en la regla de integración de lagunas para las mujeres trabajadoras por cuenta ajena mientras la brecha de género de las pensiones supere el 5 por ciento. La reforma pretendía corregir el efecto de las carreras interrumpidas de las mujeres sobre su pensión.
    }
    
    La integración de lagunas es un elemento redistributivo dentro del nivel contributivo, porque concede derechos por periodos en los que no se ha cotizado. Su justificación es de suficiencia y de equidad: protege a quienes tuvieron carreras interrumpidas por causas ajenas a su voluntad. Su límite más discutido es que solo se aplica en el Régimen General y no en el de autónomos, cuyos periodos sin cotización cuentan como ceros.
\end{specialblock}





\paragraph{Porcentaje: cuánto de la base reguladora}
La pensión es el resultado de aplicar a la base reguladora un porcentaje que depende de los años cotizados. Con $15$ años cotizados, el porcentaje es del $50\%$. En 2026, a partir de los $15$ años, cada uno de los $49$ meses siguientes añade un $0,21\%$, y cada uno de los $209$ meses siguientes, un $0,19\%$, de modo que el $100\%$ se alcanza con $36$ años y $6$ meses cotizados. 

A partir de 2027, los primeros $248$ meses adicionales añadirán un $0,19\%$ cada uno y los $16$ siguientes un $0,18\%$, de modo que el $100\%$ exigirá 37 años.\footnote{%
    Antes de 2013, el 100 por ciento de la base reguladora se alcanzaba con 35 años cotizados, con una escala más favorable que la actual. La Ley 27/2011 aumentó de forma escalonada el número de años de cotización necesarios para recibir el 100 por ciento de la base reguladora, de 35 a 37 años en 2027. La reforma pretendía reforzar la relación entre años cotizados y pensión y contener el gasto.
}

Veamoslo con un ejemplo numérico. Supongamos que un trabajador que se jubila en 2026 a la edad ordinaria, con $30$ años cotizados y una base reguladora de $2.000$ euros mensuales, tiene derecho al $50$ por ciento por sus primeros $15$ años. Por los $180$ meses restantes suma $49 \cdot 0,21 = 10,29$ puntos y $131 \cdot 0,19 = 24,89$ puntos. Su porcentaje es del $85,18\%$, y su pensión, de unos $1.704$ euros mensuales en catorce pagas.

La forma de la escala tiene implicaciones económicas que conviene hacer explícitas. La primera es que es cóncava en su origen: los primeros $15$ años dan derecho al $50\%$, lo que equivale a más de $3,3$ puntos por año, mientras que cada año adicional solo añade en torno a $2,3-2,5$ puntos. Esa concavidad es un elemento redistributivo hacia las carreras cortas, que obtienen proporcionalmente más por cada año cotizado. La segunda es que la tasa de acumulación de derechos es muy elevada en comparación internacional: la tasa a la que se acumulan derechos de pensión (\textit{accrual rate}) se situa por encima de la media de la Unión, un $2,4\%$ frente al $1,4\%$. Es uno de los factores que explican la elevada tasa de reemplazo española, que veremos más adelante. La tercera es que la escala se detiene al alcanzar el $100\%$: a partir de los $36$ años y medio (o $37$ desde 2027), cotizar más años no aumenta la pensión, salvo por la vía de la jubilación demorada. GRUBER y WISE mostraron, en su amplio programa comparativo \textit{Social Security Programs and Retirement around the World}, que este tipo de diseño genera un impuesto implícito sobre la prolongación de la vida laboral. Un trabajador que ha alcanzado el $100\%$ y sigue trabajando paga cotizaciones sin obtener a cambio un aumento de su pensión, y pierde además un año de cobro de la pensión. Estos investigadores encontraron una fuerte correlación entre la magnitud de ese impuesto implícito y la tasa de salida temprana del mercado de trabajo en los países estudiados. BOLDRIN, JIMÉNEZ-MARTÍN y PERACCHI elaboraron el capítulo español de ese programa y mostraron que el sistema español de finales de los noventa contenía fuertes incentivos a la jubilación temprana.

\paragraph{Límites: pensión máxima y pensiones mínimas}

La pensión resultante está sometida a un límite máximo, que en 2026 es de $3.359,60$ euros mensuales en catorce pagas ($47.034,40$ euros anuales), fijado por el Real Decreto 39/2026 sobre limitación de la cuantía inicial de las pensiones y su revalorización. La comparación entre 2022 y 2026 ilustra, una vez más, la reforma silenciosa. La base máxima ha crecido un $23\%$, de $4.139\to 5.101$ euros, y la pensión máxima un $19\%$, de $2.819 \to 3.360$ euros. La diferencia se ampliará en las próximas décadas por la regla de la reforma de 2023.

Por debajo, el sistema garantiza unas pensiones mínimas, cuyas cuantías dependen de la edad y de la situación familiar. En 2026, para la jubilación con $65$ o más años, son de $17.592,40$ euros anuales ($1.256,60$ euros mensuales) con cónyuge a cargo, de $13.106,80$ euros ($936,20$ euros mensuales) sin cónyuge y de $12.441,80$ euros ($888,70$ euros mensuales) con cónyuge no a cargo. Cuando la pensión contributiva calculada es inferior a la mínima, el sistema concede un complemento a mínimos hasta alcanzarla, siempre que el pensionista resida en España y no supere un límite de rentas. La reforma de 2023 vinculó por primera vez la cuantía de las pensiones mínimas al umbral de la pobreza, con una senda de convergencia que debe completarse en 2027.\footnote{%
    Antes de 2023, las pensiones mínimas se revalorizaban de forma discrecional en cada ley de presupuestos, sin una referencia explícita. La reforma contempló una subida de las pensiones mínimas contributivas y asistenciales, que pasaron a estar ligadas al umbral de la pobreza. La senda prevé que la pensión mínima con cónyuge a cargo alcance en 2027 el umbral de la pobreza para un hogar de dos adultos, y que el resto de las mínimas y las no contributivas converjan con referencias vinculadas al umbral de la pobreza de un hogar unipersonal; los porcentajes exactos deben verificarse en la norma. La reforma pretendía concretar el principio constitucional de suficiencia.
} Es la concreción operativa del principio de suficiencia examinado.

Los complementos a mínimos ilustran de forma directa el dilema del samaritano de BUCHANAN. Para un trabajador cuya pensión contributiva calculada se sitúe por debajo del mínimo, el rendimiento marginal de cotizar más es prácticamente nulo, porque cualquier aumento de su pensión calculada se compensa con una reducción equivalente del complemento. Esto genera un impuesto implícito del $100\%$ sobre las cotizaciones de los trabajadores de bajos salarios y carreras cortas, y un incentivo a la economía sumergida. Es la razón por la que los sistemas con mínimos generosos combinan la garantía con la exigencia de un periodo mínimo de cotización y con controles de rentas.

\paragraph{Modalidades de acceso: anticipar, demorar, compatibilizar}
La edad ordinaria es la referencia del sistema, pero no la única forma de acceder a la jubilación. El legislador ha establecido modalidades que permiten anticiparla, demorarla o compatibilizarla con el trabajo, y su diseño es el principal instrumento para influir en la edad efectiva de jubilación, que es la variable decisiva para la sostenibilidad del sistema-

La jubilación anticipada tiene dos variantes. La involuntaria, derivada de un cese no imputable al trabajador (despido colectivo, despido objetivo por causas económicas y otros supuestos tasados), permite adelantar la jubilación hasta cuatro años respecto de la edad ordinaria, con un mínimo de $33$ años cotizados. La voluntaria permite adelantarla hasta dos años, con un mínimo de $35$ años cotizados, y exige que la pensión resultante sea superior a la pensión mínima que correspondería a los $65$ años. Haciendo cálculos, en 2026, la edad mínima de la anticipada voluntaria es de $64$ años y $10$ meses para quienes tengan menos de $38$ años y $3$ meses cotizados, y de $63$ años para quienes superen ese umbral. 

En ambas variantes, la pensión se reduce mediante coeficientes reductores que se aplican por cada mes de anticipación y que son menores cuanto más larga es la carrera. En la anticipada voluntaria, con menos de $38$ años y $6$ meses cotizados y la anticipación máxima de $24$ meses, el coeficiente es del $21\%$. En la involuntaria, con la anticipación máxima de $48$ meses y la carrera más corta, alcanza el $30\%$ por ciento.\footnote{%
    El Real Decreto-ley 5/2013 endureció las condiciones de la jubilación anticipada y vinculó su edad mínima al aumento de la edad ordinaria. Revisó las condiciones de acceso a la jubilación parcial y anticipada, y que la jubilación anticipada lleva asociados coeficientes de penalización sobre la pensión de entrada que se aplican durante toda la vida del jubilado. La Ley 21/2021 estableció un nuevo esquema: las penalizaciones pasaron a calcularse por meses de anticipación en lugar de por trimestres, y a aplicarse después de limitar la pensión al importe máximo, en lugar de antes. El objetivo era aumentar la edad efectiva de jubilación y eliminar la ventaja de quienes tenían pensiones calculadas por encima del tope. La aplicación gradual sobre la pensión calculada entre 2024 y 2033 matizó esta última medida.
} Además, tras la reforma de 2021, los coeficientes reductores se aplican progresivamente sobre la pensión calculada, y no sobre la pensión máxima. Con ello se elimina la ventaja de quienes, con pensiones calculadas muy superiores al tope, podían anticipar la jubilación sin apenas pérdida. Esta aplicación es gradual entre 2024 y 2033; en 2026, para una anticipación voluntaria de $24$ meses con menos de $38$ años y $6$ meses cotizados, la reducción aplicable a esas pensiones altas es del $9,10\%$.

Existe, además, una jubilación anticipada por razón de la actividad o de la discapacidad, que no se basa en coeficientes reductores de la pensión, sino en coeficientes reductores de la edad. Se aplica a colectivos que realizan actividades especialmente penosas, tóxicas, peligrosas o insalubres, con elevada morbilidad o mortalidad, como la minería, los trabajadores del mar, el personal de vuelo, los bomberos, los policías locales y las policías autonómicas, así como a personas con discapacidad elevada. Su justificación es la del gradiente de longevidad: quienes realizan trabajos penosos tienen menor esperanza de vida, y jubilarlos a la misma edad que el resto supondría un trato actuarialmente desfavorable. Su debate es de economía política: cada colectivo tiene incentivos para reclamar su inclusión, y la Seguridad Social tiene que valorar la penosidad sin reglas claras, lo que ha dado lugar a un procedimiento reglamentario de reconocimiento de nuevos colectivos y a reivindicaciones de muchos sectores. Sorprendentemente, el personal de Defensa no está adscrito a este régimen, lo que le perjudica en términos contributivos (los tipos de peligrosidad son menores) y en términos prestacionales (pueden acceder más tarde).

La jubilación demorada, en sentido opuesto, premia a quien prolonga su vida laboral más allá de la edad ordinaria. El trabajador puede elegir entre tres formas de compensación por cada año completo de demora: un porcentaje adicional del $4\%$ sobre su pensión, una cantidad a tanto alzado en un solo pago, cuya cuantía depende de su carrera y de su pensión, o una combinación de ambas. Desde abril de 2025, los periodos de demora de más de $6$ meses e inferiores a $1$ año, a partir del primer año, dan derecho a un incremento del $2\%$ por ciento.\footnote{%
    Antes de 2022, el porcentaje adicional por cada año de demora era del 2, 2,75 o 4 por ciento, según los años cotizados. La Ley 21/2021 aumentó el porcentaje adicional asociado a la demora e introdujo la posibilidad de recibir una cantidad a tanto alzado por cada año de trabajo entre la edad legal de jubilación y el acceso a la pensión. El Real Decreto-ley 11/2024 introdujo el cómputo por semestres a partir del primer año. El objetivo era aumentar la edad efectiva de jubilación.
}

La jubilación activa permite compatibilizar el cobro de una parte de la pensión con un trabajo por cuenta ajena o propia. Tras la reforma de 2024, en vigor desde abril de 2025, exige haber demorado al menos un año el acceso a la jubilación desde la edad ordinaria, pero ya no una carrera completa: basta con la carencia mínima de $15$ años. La parte compatible de la pensión es del $45\%$ por ciento si se accede tras un año de demora, crece con los años de demora hasta alcanzar el $100\%$, y aumenta además en cinco puntos por cada año que se mantenga la compatibilidad. Los autónomos que contraten a un trabajador indefinido pueden compatibilizar un porcentaje mayor.\footnote{%
    Antes de abril de 2025, la jubilación activa exigía haber alcanzado el 100 por ciento de la base reguladora y permitía compatibilizar el 50 por ciento de la pensión (o el 100 por ciento para los autónomos con al menos un trabajador). La reforma eliminó el requisito de carrera completa y estableció porcentajes crecientes con la demora, con el objetivo de ampliar su utilización.
}

La jubilación parcial permite, en sentido inverso, reducir la jornada y cobrar la parte proporcional de la pensión antes de la jubilación completa. Desde abril de 2025 puede anticiparse hasta $3$ años respecto de la edad ordinaria. Cuando la anticipación supera los $2$ años, la reducción de jornada durante el primer año debe estar entre el $20\%$ y el $33\%$. La jubilación parcial anticipada exige la celebración simultánea de un contrato de relevo con otro trabajador, que ahora debe ser indefinido y a tiempo completo. La jubilación parcial a partir de la edad ordinaria, sin contrato de relevo, admite reducciones de jornada de entre el 25 y el 75 por ciento.\footnote{%
    La jubilación parcial se reguló en 2002 y se endureció en 2007 y en 2013, por su utilización como vía de salida anticipada. Las reglas de la jubilación parcial anticipada son prolijas y que, en general, la edad más temprana de acceso pasa de 61 años en 2013 a 63 en 2027, con 33 años cotizados, y que la jubilación parcial con contrato de relevo no tiene penalización. El Real Decreto-ley 11/2024 amplió la anticipación a tres años, pero exigió que el contrato de relevo fuera indefinido y a tiempo completo, con el objetivo de que la jubilación parcial genere empleo estable.
}

La evaluación económica de estas modalidades gira en torno a un concepto central: la neutralidad actuarial. Una modalidad es actuarialmente neutral si el valor actual de la pensión que se recibe a lo largo de la vida es el mismo con independencia de la edad a la que se accede. Adelantar la jubilación supone cobrar durante más años una pensión más baja; demorarla, cobrar durante menos años una pensión más alta. Si las penalizaciones por anticipar son menores que las neutrales, el sistema subvenciona la jubilación temprana; si las bonificaciones por demorar son menores que las neutrales, la penaliza implícitamente.

Por ejemplo, a la edad ordinaria, con una esperanza de vida en torno a $20$ años, demorar un año la jubilación supone renunciar aproximadamente a $1/20$ parte del valor de la pensión vitalicia, en torno a un $5\%$. Un complemento del $4\%$ por año se sitúa, por tanto, algo por debajo de la neutralidad en términos brutos, aunque la valoración exacta depende del tipo de descuento y de la tabla de mortalidad utilizados. Los coeficientes de la anticipada voluntaria, en torno al $10,5\%$ por año para las carreras más cortas, se sitúan en cambio por encima de la neutralidad, de modo que la penalizan.

Los datos más recientes sugieren que el rediseño está teniendo efecto. En septiembre de 2026, el $11,6\%$ de las nuevas jubilaciones correspondía a demoras voluntarias, frente al $6,8\%$ de 2019, y la edad media de jubilación había aumentado hasta los $65,4$ años, desde los $64,4$ anteriores. La interpretación de estos datos exige cautela, porque la edad media de jubilación también aumenta mecánicamente por el aumento de la edad ordinaria establecido en 2011, que en 2026 alcanza los $66$ años y $10$ meses. Separar el efecto de los incentivos del efecto de la edad ordinaria es una cuestión empírica que se examina más adelante.

La jubilación flexible y gradual, que combina la parcial y la activa, responde a una idea que la literatura sobre la oferta de trabajo de los mayores ha defendido desde hace décadas: la transición entre la actividad y la jubilación no tiene por qué ser un salto brusco, y permitir combinar trabajo y pensión puede aumentar la participación laboral de los mayores. Su riesgo, que la experiencia española de la jubilación parcial con relevo puso de manifiesto, es que se utilice como vía de salida anticipada financiada por el sistema. En determinados sectores industriales se generalizó la jubilación parcial con reducciones de jornada máximas y concentración de la jornada restante, de modo que el trabajador dejaba de trabajar en la práctica años antes de la edad ordinaria. Por eso las reformas de 2013 y 2024 endurecieron los requisitos de reducción de jornada y de contrato de relevo.

\paragraph{Complemento para la reducción de la brecha de género}
Desde 2021, las pensiones contributivas de jubilación, incapacidad permanente y viudedad de las mujeres que hayan tenido uno o más hijos incluyen un complemento para la reducción de la brecha de género.\footnote{%
    Entre 2016 y 2021 existía un complemento por maternidad, un porcentaje adicional del 5 al 15 por ciento de la pensión para las mujeres con dos o más hijos. La sentencia del TJUE de 12 de diciembre de 2019 (asunto C-450/18) lo consideró discriminatorio para los hombres (1.3.4). El Real Decreto-ley 3/2021 lo sustituyó por el complemento actual, una cantidad fija por hijo desde el primero y accesible a los hombres cuya carrera se hubiera visto afectada. El Real Decreto-ley 2/2023 aumentó su cuantía. La reforma pretendía adaptar el complemento al derecho de la Unión y mejorar su focalización.
} Puede reconocerse a los hombres si acreditan haber interrumpido o visto afectada su carrera profesional con ocasión del nacimiento o la adopción de sus hijos. Si se reconoce al segundo progenitor, se extingue el complemento ya reconocido al primero. La vigencia del complemento está ligada a que la brecha de género de las pensiones, definida como el porcentaje que representa la diferencia entre el importe medio de las pensiones de jubilación contributivas causadas en un año por las mujeres y el de las causadas por los hombres, sea superior al $5\%$. En 2026, el complemento es de $36,90$ euros mensuales por hijo, con un máximo de cuatro hijos. En septiembre de 2026, unos $1,7$ millones de pensiones lo incluían, con una cuantía media de $76,70$ euros mensuales.

Sus defensores sostienen que compensa, aunque sea parcialmente, el perjuicio que la maternidad ha causado en las carreras laborales de las mujeres, en un sistema contributivo que traslada a la vejez las desigualdades de la vida activa. Sus críticos señalan tres problemas. El primero es de focalización: se concede por el número de hijos y no por el perjuicio efectivo en la carrera, de modo que lo reciben también mujeres con carreras completas y pensiones elevadas que no sufrieron ningún perjuicio. El segundo es de cuantía: un complemento de unos $37$ euros mensuales por hijo es muy inferior a la brecha que pretende corregir. El tercero es de diseño: corrige la brecha ex post, en el momento de la pensión, en lugar de hacerlo ex ante, mediante créditos de cotización por cuidado.\footnote{%
    Estos créditos existen en España, en forma de periodos de cotización reconocidos por el cuidado de hijos, pero son de alcance limitado, mientras que en países como Alemania o Suecia tienen un peso mucho mayor.
} 

La brecha de género de las pensiones de jubilación en España sigue siendo elevada, en torno a una tercera parte de la pensión media masculina; su cifra exacta en la última estadística debe verificarse.

\subsubsection{Revalorización de las pensiones}
Las pensiones contributivas, incluido el importe de las pensiones mínimas, se revalorizan al comienzo de cada año en un porcentaje igual a la media de las tasas de variación interanual del IPC de los doce meses previos a diciembre del año anterior. Para 2026, esa media fue del $2,7\%$, que es la revalorización general aplicada. Si la media fuera negativa, las pensiones se mantendrían sin cambios. 

Las pensiones mínimas y las no contributivas han tenido revalorizaciones superiores, en aplicación de la senda de convergencia con el umbral de la pobreza: un $11,4\%$ en 2026 para las no contributivas y el IMV.\footnote{%
    Desde 1985 hasta 2013, las pensiones se revalorizaban con el IPC previsto, con una paga compensatoria si la inflación real lo superaba. En 2011 se congelaron (Real Decreto-ley 8/2010) y en 2012 se suprimió la compensación por la desviación (Real Decreto-ley 28/2012). La Ley 23/2013 estableció un índice de revalorización ligado al equilibrio financiero del sistema, con un mínimo del 0,25 por ciento y un máximo del IPC más 0,5 puntos. Su aplicación entre 2014 y 2017 mantuvo la revalorización en el mínimo. Entre 2018 y 2021 la revalorización se fijó de forma discrecional, y la Ley 21/2021 restableció la indexación al IPC medio. El objetivo era garantizar el poder adquisitivo de las pensiones tras la contestación social al índice de 2013. La mecánica del índice de 2013 y la evaluación de ambos sistemas se examinan en 4.3.2 y 4.3.3.
}

La elección del índice de revalorización es una de las decisiones de diseño con mayor impacto a largo plazo, y la literatura distingue tres opciones con consecuencias muy distintas. La indexación a los precios, la española actual, mantiene el poder adquisitivo de la pensión a lo largo de la jubilación. Pero, si los salarios reales crecen, la pensión pierde valor relativo respecto del nivel de vida medio de la sociedad: un jubilado que cobra una pensión equivalente al $80\%$ del salario medio al jubilarse puede cobrar, veinte años después, una pensión equivalente a menos del $60\%$ del salario medio de ese momento, aunque su poder adquisitivo no haya cambiado. La indexación a los salarios, como en Alemania, mantiene el valor relativo, pero es más costosa y traslada a los pensionistas el riesgo de caída de los salarios reales. Los índices mixtos, como el suizo, que combina precios y salarios a partes iguales, buscan un punto intermedio. Y los índices de sostenibilidad, como el factor alemán de 2004 o el índice español de 2013, ligan la revalorización al equilibrio financiero del sistema.

La elección tiene una dimensión distributiva entre generaciones. En un periodo de alta inflación con salarios que crecen menos que los precios, como 2022-2023, la indexación a los precios protege a los pensionistas frente a los trabajadores. Los pensionistas mantuvieron su poder adquisitivo, con una revalorización del $8,5\%$ en 2023, mientras los salarios reales de los trabajadores caían. Es la función de seguro intergeneracional, pero en una dirección que ha suscitado críticas por su efecto sobre la equidad entre generaciones. En sentido inverso, en un periodo de crecimiento de los salarios reales, la indexación a los precios hace que los pensionistas pierdan terreno respecto de los trabajadores.

Como sabemos, además, la garantía de revalorización con el IPC es una opción del legislador ordinario y no una exigencia constitucional. La Constitución exige que las pensiones sean periódicamente actualizadas, pero no impone un índice concreto, y la STC 49/2015 confirmó que la actualización es una expectativa hasta que se consolida. 

\subsubsection{Incapacidad permanente y la protección por muerte y supervivencia}
La incapacidad permanente protege la pérdida de la capacidad de trabajo derivada de una enfermedad o un accidente, cuando, tras el tratamiento prescrito, el trabajador presenta reducciones anatómicas o funcionales graves, previsiblemente definitivas, que disminuyen o anulan su capacidad laboral. 

La ley distingue cuatro grados, según la intensidad de la pérdida de capacidad, cada uno con una prestación distinta. La incapacidad permanente parcial reduce el rendimiento en la profesión habitual sin impedir sus tareas fundamentales, y da derecho a una indemnización a tanto alzado de $24$ mensualidades de la base reguladora. La total inhabilita para la profesión habitual, pero no para otra distinta, y da derecho a una pensión del $55\%$ de la base reguladora, que se incrementa en un $20\%$ (la llamada total cualificada) para los mayores de $55$ años con dificultades para encontrar otro empleo. La absoluta inhabilita para toda profesión u oficio, y da derecho a una pensión del $100\%$ de la base reguladora. La gran invalidez requiere además la asistencia de otra persona para los actos más esenciales de la vida, y da derecho a la pensión de la incapacidad absoluta más un complemento destinado a remunerar a esa persona. Los requisitos de carencia y el cálculo de la base reguladora dependen de que la incapacidad derive de enfermedad común, accidente no laboral o contingencias profesionales. Para estas últimas no se exige periodo de carencia, en coherencia con la lógica de responsabilidad objetiva heredada de 1900. Al cumplir la edad ordinaria de jubilación, la pensión de incapacidad permanente pasa a denominarse pensión de jubilación, sin cambios en su cuantía.

La incapacidad permanente es la prestación en la que el riesgo moral se manifiesta con mayor claridad, porque la frontera entre la incapacidad real y la falta de empleo es difusa, especialmente para los trabajadores mayores de baja cualificación.\footnote{%
    La experiencia neerlandesa de los años ochenta, en la que las pensiones de incapacidad se convirtieron en una vía masiva de salida del mercado de trabajo, es el ejemplo de manual. En España, las tasas de incapacidad permanente muestran diferencias territoriales difíciles de explicar solo por diferencias de salud, lo que sugiere la influencia de factores económicos y de gestión.
}

Dos novedades recientes muestran que el régimen de la incapacidad permanente está en revisión. La primera es de origen europeo. La sentencia del TJUE (C-631/22), consideró contraria a la Directiva de igualdad de trato en el empleo la extinción automática del contrato de trabajo por la declaración de incapacidad permanente, sin que el empresario tuviera que valorar antes la posibilidad de realizar ajustes razonables. En respuesta, la Ley 2/2025, modificó el Estatuto de los Trabajadores para que la incapacidad permanente no extinga automáticamente el contrato. El empresario debe valorar ahora si puede adaptar el puesto o recolocar al trabajador. La segunda es de origen jurisprudencial interno. El TS modificó en 2024 su doctrina sobre la compatibilidad de la incapacidad permanente absoluta con el trabajo, y la limitó a actividades marginales. Esta doctrina ha abierto el debate sobre si el sistema debe favorecer la reincorporación laboral de los incapacitados, en línea con la lógica de la activación, o proteger su derecho a la pensión sin condiciones. 

\paragraph{Protección por muerte y supervivencia}
La protección por muerte y supervivencia protege a los familiares de un trabajador o pensionista fallecido. Su prestación principal es la pensión de viudedad, que consiste, con carácter general, en el 52 por ciento de la base reguladora del causante. El porcentaje se eleva al $60\%$ para los mayores de $65$ años sin otros ingresos significativos y al $70\%$ cuando el viudo tiene cargas familiares y la pensión constituye su principal fuente de ingresos. La pensión se extiende a las parejas de hecho que cumplan determinados requisitos de convivencia y dependencia económica, como estableció la Ley 40/2007. También se extiende, en proporción al tiempo de convivencia, a los cónyuges separados o divorciados que tuvieran derecho a una pensión compensatoria. Se completa con las pensiones de orfandad, de un $20\%$ de la base reguladora por hijo, con límites de edad que se extienden en caso de discapacidad o de ausencia de ingresos, y con las pensiones en favor de familiares, de carácter residual. 

La pensión de viudedad es una de las prestaciones más discutidas del sistema. En su origen, respondía a una lógica clara: en el modelo del varón sustentador, la muerte del marido privaba a la viuda, que no tenía ingresos propios, de su fuente de sustento. La pensión de viudedad era, por tanto, un seguro de rentas para la familia dependiente. Pero esa lógica encaja cada vez peor con la realidad social actual, por tres razones. La primera es que la incorporación de las mujeres al mercado de trabajo hace que muchas viudas tengan su propia pensión de jubilación. La viudedad es compatible sin límite con la pensión propia y con cualquier renta, lo que permite acumular dos pensiones que, en conjunto, pueden superar ampliamente el nivel de vida que tenía el hogar (mi abuela). La segunda es que la pensión se concede con independencia de la necesidad, lo que la aleja del principio de suficiencia y la acerca a un derecho patrimonial derivado del matrimonio. La tercera es que, en un sistema con presupuesto limitado, el gasto en viudedad compite con el gasto en otras prestaciones con mayor capacidad redistributiva, como las de familia o infancia.

Las propuestas de reforma, apuntan a reformular la pensión de viudedad para atender mejor a las situaciones de necesidad real y para diferenciar entre las viudedades de personas mayores, que dependían del causante, y las de personas jóvenes con capacidad laboral. Varios países europeos han adoptado pensiones de viudedad temporales para los supervivientes en edad de trabajar, con el objetivo de facilitar su transición al empleo, y han limitado su compatibilidad con otras rentas. Suecia fue pionera al suprimir en 1990 la pensión de viudedad vitalicia para los nuevos matrimonios. La reforma de la viudedad en España ha resultado políticamente muy difícil, por la misma lógica de PIERSON: afecta a un colectivo numeroso, mayoritariamente femenino y de edad avanzada, con gran capacidad de movilización y con una fuerte legitimidad social. Mi abuela...

\subsubsection{Nivel no contributivo: las pensiones no contributivas}
El nivel no contributivo cubre a quienes, estando en situación de necesidad, no tienen rentas o patrimonio suficientes y no han cotizado el tiempo necesario para acceder a las prestaciones contributivas. Fueron creadas por la Ley 26/1990 y tienen dos modalidades. La de jubilación, que exige tener $65$ o más años y haber residido en España durante al menos $10$ años entre los $16$ y la edad de solicitud, de los que dos deben ser inmediatamente anteriores. Y la de invalidez, que exige tener entre $18$ y $65$ años y un grado de discapacidad o de enfermedad crónica igual o superior al $65\%$. Ambas exigen carecer de rentas o ingresos suficientes, por debajo de un límite que se calcula por unidad económica de convivencia. 

En 2026, la cuantía íntegra de la pensión no contributiva es de $8.803,20$ euros anuales, $628,80$ euros mensuales en catorce pagas, un $11,4\%$ más que en 2025. A ella puede sumarse un complemento para quienes viven de alquiler. Su gestión corresponde a las Comunidades Autónomas, con la excepción de Ceuta y Melilla, donde la gestiona el IMSERSO, y su financiación, al Estado. Varias Comunidades Autónomas conceden complementos propios a las pensiones no contributivas con cargo a su presupuesto, amparadas en su competencia de asistencia social y en la doctrina de la STC 239/2002.\footnote{%
    Antes de 1990, quienes no tenían derecho a una pensión contributiva dependían de prestaciones asistenciales dispersas, como las del antiguo Fondo Nacional de Asistencia Social, de cuantía muy reducida. La Ley 26/1990 creó un derecho subjetivo a una pensión no contributiva, universalizando la protección del sistema.
}

Las pensiones no contributivas plantean los problemas característicos de las prestaciones sujetas a comprobación de recursos que ya examinamos. El primero es la trampa de la pobreza: como la pensión se reduce euro a euro cuando aumentan los demás ingresos del hogar, el beneficiario se enfrenta a un tipo marginal implícito cercano al $100\%$ sobre cualquier ingreso adicional. El segundo es el estigma y la baja utilización que TITMUSS asociaba a los programas selectivos. El tercero es la heterogeneidad territorial que introducen los complementos autonómicos. La convergencia de la cuantía con el umbral de la pobreza, impulsada por la reforma de 2023, ha reforzado su función de suficiencia, pero no ha resuelto estos problemas de diseño.

\subsubsection{Ingreso Mínimo Vital}
El IMV, establecido en 2020, se presente como un derecho subjetivo a recibir una prestación económica que garantiza un nivel mínimo de renta a quienes se encuentran en situación de vulnerabilidad económica. La prestación cubre la diferencia entre la suma de los ingresos de un hogar en el año anterior, o en ocasiones en el año en curso, y una renta mínima garantizada por ley, cuyo importe depende del número de miembros de la unidad de convivencia. Lo creó el Real Decreto-ley 20/2020, de 29 de mayo, luego tramitado como Ley 19/2021, de 20 de diciembre, y lo gestiona el Instituto Nacional de la Seguridad Social, salvo en el País Vasco y Navarra (1.4.3). Es una prestación no contributiva de la Seguridad Social, financiada con transferencias del Estado.

La renta garantizada de un beneficiario individual equivalía al $100\%$ de la pensión no contributiva anual prorrateada en doce mensualidades. Con la cuantía de la pensión no contributiva de 2026, la renta garantizada individual se situaría en torno a $733,60$ euros mensuales. Para las unidades de convivencia, la cuantía se incrementa en un $30\%$ por cada miembro adicional a partir del segundo, con un máximo del $220\%$. Existe un complemento para las unidades monoparentales o con algún miembro con discapacidad. Además, se abona un complemento de ayuda para la infancia por cada menor de la unidad de convivencia. Este último complemento puede percibirse aunque el hogar no tenga derecho al IMV completo, siempre que sus ingresos no superen un umbral más elevado.

\paragraph{Problemas: ¿qué alcance tiene y qué efectos genera?}
La cuarta opinión de la AIReF sobre el IMV, publicada en julio de 2025, actualizó el diagnóstico de cobertura. 

En 2024, la prestación llegaba al $42\%$ de los hogares con derecho a ella, y el complemento de ayuda para la infancia, solo al $23\%$. Dicho de otro modo, la tasa de no solicitud (non take-up) se situaba en el $55\%$ para el IMV y en el $72\%$ para el complemento de ayuda para la infancia, cifras que no habían mejorado respecto de años anteriores. 

La baja utilización es el problema que TITMUSS asociaba a las prestaciones selectivas. La literatura contemporánea lo ha analizado bajo el concepto de carga administrativa, formulado por HERD y MOYNIHAN (2018): los costes de aprendizaje (conocer la prestación y sus requisitos), de cumplimiento (reunir la documentación y completar el trámite) y psicológicos (el estigma, la frustración ante denegaciones) que la Administración impone a los potenciales beneficiarios. Según esta literatura, la carga administrativa no es un efecto secundario inevitable, sino una variable de diseño que actúa como mecanismo implícito de racionamiento: una prestación con mucha carga administrativa llega a menos beneficiarios y cuesta menos, sin necesidad de reducir su cuantía ni sus requisitos formales.\footnote{%
    CURRIE (2006) revisó la evidencia internacional sobre la utilización de las prestaciones sociales y concluyó que la baja utilización es un fenómeno generalizado en los programas condicionados a la comprobación de recursos, y que los factores administrativos pesan más que el estigma.
}

Por otro lado, la AIReF encontró un efecto negativo sobre el empleo: percibir el IMV reduce la probabilidad de trabajar en unos $3$ puntos porcentuales, un $12\%$, y el número de días trabajados al mes en un $11\%$. El efecto persiste incluso entre los beneficiarios que empezaron a percibirlo después de la introducción del incentivo al empleo en 2023 y supera el $20\%$ en algunos grupos, como las familias monoparentales y los menores de $30$ años. Y, además, constató una elevada persistencia: en torno al $60\%$ de los beneficiarios mantenía la prestación durante más de tres años.

Como sabemos, el efecto negativo sobre el empleo es la manifestación del riesgo moral y de la trampa de la pobreza. Si la prestación se reduce en la misma cuantía en que aumentan los ingresos laborales, el beneficiario no gana nada trabajando.\footnote{%
    El incentivo al empleo introducido en 2023, que permite conservar una parte de la prestación cuando se obtienen ingresos laborales, intentaba corregir este problema, pero la AIReF concluye que su diseño es demasiado complejo y su efecto demasiado diferido para que los beneficiarios lo perciban.
} La tensión entre la suficiencia de la prestación y el incentivo al trabajo es el dilema clásico de toda renta mínima. Su solución teórica, desde el impuesto negativo sobre la renta de FRIEDMAN, consiste en reducir la prestación en menos de un euro por cada euro de ingreso adicional. Pero, evidentemente, esa solución aumenta el coste de la prestación y el número de beneficiarios.

La coexistencia del IMV con las rentas mínimas autonómicas es el último problema de diseño. Las Comunidades Autónomas tenían sus propios programas desde finales de los ochenta, con enormes diferencias de cuantía, cobertura y requisitos entre ellas, y su competencia en materia de asistencia social les permite mantenerlos. Con la creación del IMV, las comunidades adoptaron estrategias diversas: algunas transformaron sus rentas en complementos del IMV, otras las mantuvieron como prestaciones independientes y otras las redujeron. El resultado es un mapa de protección de última instancia todavía muy heterogéneo territorialmente. La Recomendación del Consejo de enero de 2023 sobre una renta mínima adecuada, insta a los Estados a garantizar niveles adecuados, reducir la falta de solicitud y articular la renta mínima con la activación laboral, que son exactamente los problemas que la AIReF ha identificado en el caso español.

Con todo, la AIReF recomendó una reformulación completa del incentivo al empleo, para hacerlo más transparente e inmediato, y la automatización del reconocimiento de la prestación, en particular del complemento de ayuda para la infancia desde el nacimiento. El Ministerio respondió que la evaluación del incentivo al empleo era temprana y que convenía esperar a datos más consolidados.

\subsection{Prestaciones económicas de corto plazo}
Estas prestaciones comparten una lógica que las distingue de las pensiones y que conviene explicitar. Las pensiones sustituyen las rentas del trabajo de forma permanente, cuando la capacidad de trabajar se ha perdido o se considera socialmente legítimo dejar de trabajar. Las prestaciones de corto plazo sustituyen las rentas de forma temporal, durante un periodo en el que el trabajador conserva su vínculo con el empleo, pero no puede o no debe trabajar: porque está enfermo, porque acaba de tener un hijo o porque su trabajo supone un riesgo para su embarazo. Esta diferencia tiene consecuencias económicas decisivas. 

En las prestaciones de corto plazo, la decisión de acceder y la duración dependen en gran medida de la conducta del beneficiario y de quienes certifican su situación. Por eso el riesgo moral examinado en 1.1.1, y el diseño de los incentivos y del control, ocupan un lugar central. Además, el coste de estas prestaciones no lo soporta solo el sistema, sino también las empresas, a través de los días de baja que pagan y de la pérdida de producción, lo que convierte su diseño en un problema de reparto de costes entre instituciones.

Orahizaremos el subapartado en cuatro partes: la incapacidad temporal (3.2.1), a la que se dedica la mayor extensión por su peso y su actualidad; el nacimiento y cuidado de menor (3.2.2); las prestaciones por riesgo durante el embarazo y la lactancia y por cuidado de menores con enfermedad grave (3.2.3); y las prestaciones familiares (3.2.4).

\subsubsection{Incapacidad temporal}
La incapacidad temporal protege la pérdida de rentas del trabajador que, debido a una enfermedad o un accidente, está impedido temporalmente para trabajar y recibe asistencia sanitaria. Protege también los periodos de observación por enfermedad profesional. La prestación exige, cuando deriva de enfermedad común, un periodo de cotización de 180 días en los cinco años anteriores. No exige periodo de carencia cuando deriva de accidente, sea o no de trabajo, o de enfermedad profesional, en coherencia con la lógica de la responsabilidad objetiva heredada de 1900 (1.2.2).\footnote{%
    La Ley Orgánica 1/2023, que modificó la ley de salud sexual y reproductiva, creó tres situaciones especiales de incapacidad temporal: la menstruación incapacitante secundaria, la interrupción del embarazo y la gestación desde el primer día de la semana 39.Estas situaciones se pagan desde el primer día a cargo de la Seguridad Social, sin que el empresario soporte los días cuarto a decimoquinto, y sin periodo de carencia. La justificación es evitar que su coste recaiga en la empresa y genere discriminación en la contratación de mujeres.

Antes de 2023, la menstruación incapacitante, la interrupción del embarazo y la última fase de la gestación se protegían, en su caso, como incapacidad temporal ordinaria por enfermedad común, con los días de carencia y de pago empresarial correspondientes. La Ley Orgánica 1/2023 las configuró como situaciones especiales pagadas por la Seguridad Social desde el primer día y sin carencia, con el objetivo de proteger la salud de las mujeres y evitar la discriminación en la contratación.
}

La cuantía es un porcentaje de la base reguladora, que se calcula sobre la base de cotización del mes anterior a la baja. En las contingencias comunes es del 60 por ciento entre el cuarto y el vigésimo día de baja, y del 75 por ciento a partir del vigésimo primero. Los tres primeros días no dan derecho a prestación pública, aunque muchos convenios colectivos los retribuyen. En las contingencias profesionales, la prestación es del 75 por ciento desde el día siguiente a la baja. El pago se reparte entre instituciones en función de la duración de la baja. En las contingencias comunes, entre el cuarto y el decimoquinto día la prestación corre a cargo del empresario, y desde el decimosexto, a cargo del Instituto Nacional de la Seguridad Social o de la mutua colaboradora, según la opción de la empresa. En la práctica, la empresa adelanta el pago al trabajador en su nómina mediante el pago delegado y lo deduce después de sus cotizaciones (1.4.4). La duración máxima es de 365 días, prorrogables otros 180 cuando se prevé la curación. Agotado ese plazo, el trabajador es evaluado para determinar si procede declarar una incapacidad permanente (3.1.3).\footnote{%
    Un elemento que el debate público suele ignorar es que la prestación pública no es lo que realmente cobra la mayoría de los trabajadores de baja. Una proporción muy elevada de los convenios colectivos establece mejoras voluntarias a cargo de la empresa, que complementan la prestación hasta el 100 por ciento del salario, a veces desde el primer día. En consecuencia, la tasa de reemplazo efectiva de la incapacidad temporal es, para muchos trabajadores, completa. Esto tiene consecuencias sobre los incentivos que se examinan a continuación.
}

El control de la prestación es donde la arquitectura institucional revela su fragmentación. La baja y el alta las emiten los médicos de atención primaria de los servicios públicos de salud de las Comunidades Autónomas, que no pagan la prestación. El INSS, que sí la paga, dispone de inspectores médicos que pueden emitir altas y controlar los procesos, pero con recursos limitados. Las mutuas colaboradoras, que pagan la prestación en las empresas asociadas, pueden formular propuestas de alta que la inspección médica autonómica debe resolver. Desde abril de 2023, el trabajador ya no tiene que entregar a la empresa los partes de baja y confirmación: el servicio de salud los transmite telemáticamente al INSS, que los comunica a la empresa.\footnote{%
    Hasta abril de 2023, el trabajador debía entregar a su empresa una copia de los partes de baja y de confirmación en un plazo de tres días. El Real Decreto 1060/2022, que modificó el Real Decreto 625/2014 sobre la gestión de la incapacidad temporal en sus primeros 365 días, suprimió esa obligación y estableció la transmisión telemática de los partes del servicio de salud al INSS y de este a la empresa. El objetivo era reducir la carga administrativa para el trabajador enfermo.
}

La AIReF publicó en febrero de 2026 un estudio de revisión del gasto sobre la incapacidad temporal, dentro de su ciclo de evaluaciones del gasto público 2022-2026. Sus conclusiones son contundentes. Entre 2017 y 2024, la incidencia de las bajas aumentó un $60\%$ y su duración media, un $15\%$. El gasto alcanzó unos $16.500$ millones de euros anuales, la segunda partida del sistema después de las pensiones. El número de procesos por contingencias comunes casi se duplicó, de 4,7 millones a 8,6 millones. Las bajas por salud mental fueron las de crecimiento más intenso, con un aumento de su duración media de 67 a 98,5 días. El fenómeno está muy concentrado: el 25 por ciento de las personas acumula el 55 por ciento de los procesos. Los datos del Ministerio para 2025 confirman la tendencia: el gasto en incapacidad temporal alcanzó los 18.413 millones de euros, un 11,8 por ciento más que el año anterior. Es casi el doble del crecimiento del gasto total del sistema, del 6,5 por ciento (2.3.1).

La conclusión para España del creciente gasto en Incapacidad Temporal, como veremos, es que es, sobre todo, un problema de diseño institucional. Quien decide no paga, quien paga no decide, y quien soporta una parte del coste (la empresa) no tiene capacidad de control. Ninguna reforma que no aborde esa desconexión, sea mediante la corresponsabilidad financiera de los servicios de salud, sea mediante un mayor papel del INSS y de las mutuas en el control, sea mediante la reincorporación progresiva, es probable que contenga de forma duradera el crecimiento del gasto.

\begin{specialblock}{\textbf{Evidencia empírica:} ¿Qué nos dice la investigación económica sobre la Incapacidad Temporal?}
    La incapacidad temporal es un seguro de rentas de corto plazo frente al riesgo de enfermedad: un trabajador enfermo perdería su renta justo cuando más la necesita, y la presión económica le llevaría a trabajar enfermo, con riesgo para su salud y la de sus compañeros. Pero es también la prestación en la que el riesgo moral es más intenso, porque la frontera entre estar y no estar en condiciones de trabajar depende en gran medida de la valoración subjetiva del propio trabajador y de su médico, y porque el coste de una baja para el trabajador es bajo cuando la tasa de reemplazo es elevada.
    
    La literatura empírica sobre las bajas por enfermedad es abundante, y sus conclusiones son relevantes para el debate español. JOHANSSON y PALME (1996, 2005), con datos suecos, mostraron que las reformas que redujeron la tasa de reemplazo de la prestación por enfermedad redujeron significativamente el absentismo. ZIEBARTH y KARLSSON (2010) encontraron, con una reforma alemana que redujo la cobertura de la baja por enfermedad del 100 al 80 por ciento del salario, una reducción de las ausencias. La conclusión general es que la elasticidad del absentismo respecto de la generosidad de la prestación es positiva y significativa: cuanto mayor es la tasa de reemplazo, más bajas y más largas.\footnote{%
        Un segundo hallazgo se refiere al papel de los médicos como controladores del acceso. MARKUSSEN, RØED, RØGEBERG y GAURE (2011), con datos noruegos, mostraron que el absentismo depende no solo de las características de los trabajadores, sino también de los médicos que certifican las bajas: unos son sistemáticamente más proclives que otros a concederlas. MARKUSSEN, RØED y RØGEBERG (2013) mostraron que, cuando un trabajador cambia de médico de familia por razones ajenas a su salud, su probabilidad de baja cambia según la propensión del nuevo médico. Estos resultados cuestionan la idea de que la baja es una decisión puramente clínica, y señalan que el médico que certifica tiene un margen de discrecionalidad que responde a su propio criterio, a su carga de trabajo y a la presión del paciente, y no a los costes que la baja genera para el sistema.
    } También se ha estudiado la relación entre el absentismo y el ciclo económico. ASKILDSEN, BRATBERG y NILSEN (2005) mostraron que las bajas por enfermedad son procíclicas: aumentan en las expansiones y disminuyen en las recesiones. La literatura ofrece dos explicaciones complementarias. La primera es un efecto disciplina: cuando el desempleo es alto, los trabajadores temen perder su empleo y evitan las bajas. La segunda es un efecto composición: en las expansiones se incorporan al empleo trabajadores con peor salud o menor vinculación con el trabajo, que tienen más bajas. Más recientemente, se ha estudiado la adecuidad de las bajas parciales. Markussen, Mykletun y Røed (2012), mostraron que en Noruega la posibilidad de una baja parcial (trabajar a tiempo reducido mientras se cobra una prestación parcial) aumentaba la probabilidad de reincorporación completa al trabajo y reducía la duración total de las bajas. La razón es que mantener el vínculo con el trabajo durante la recuperación evita el deterioro de las capacidades y de la motivación que produce una ausencia prolongada. Fundamento teórico de las propuestas de "alta progresiva" que se examinan más adelante.
\end{specialblock}

\paragraph{España: ¿cuál es el diagnóstico empírico?}}
La AIReF identificó varios factores explicativos, que conectan con la literatura antes expuesta. El primero y principal es institucional: existe una desconexión entre quien concede la baja (el médico de atención primaria del servicio de salud autonómico) y quien la paga (el INSS o la mutua). Es exactamente el problema de coordinación examinado en 1.4.2, que en términos económicos es una externalidad fiscal vertical. El servicio de salud autonómico no soporta el coste de las bajas que concede, y puede incluso tener incentivos para prolongarlas cuando sus listas de espera retrasan el tratamiento que permitiría la curación. Es la versión institucional del hallazgo de Markussen y otros sobre el papel de los médicos: en España, quien certifica no solo tiene discrecionalidad, sino que no tiene ningún incentivo a considerar el coste de su decisión.
    
El segundo factor son varios cambios normativos y laborales que aumentaron la generosidad efectiva de la protección. La AIReF encontró que el restablecimiento del complemento al 100 por ciento del salario para los empleados públicos aumentó la probabilidad de iniciar una baja en torno a un 40 por ciento. Ese complemento se había reducido durante la crisis y se recuperó a partir de 2018.\footnote{%
    Durante la crisis, el Real Decreto-ley 20/2012 limitó el complemento que las administraciones públicas abonaban a sus empleados durante la incapacidad temporal, para reducir el absentismo y el gasto. A partir de la Ley de Presupuestos Generales del Estado para 2018, las administraciones recuperaron la posibilidad de complementar la prestación hasta el 100 por ciento de las retribuciones desde el primer día. La AIReF ha utilizado ese cambio como experimento natural para estimar el efecto de la tasa de reemplazo sobre el absentismo.
}
Es un experimento natural que confirma la elasticidad del absentismo respecto de la tasa de reemplazo documentada por Johansson y Palme. La AIReF encontró también que el paso de un contrato temporal a uno indefinido aumentaba la probabilidad de baja en torno a un 30 por ciento. Es la manifestación del efecto disciplina de Askildsen y otros: un trabajador temporal teme que una baja le cueste la renovación de su contrato, mientras que un trabajador indefinido tiene más protección. La reforma laboral de 2021, al reducir drásticamente la temporalidad, tuvo así un efecto colateral sobre el absentismo que sus diseñadores probablemente no previeron. El tercer factor son las listas de espera del sistema sanitario, que prolongan las bajas mientras el trabajador espera una prueba diagnóstica o una intervención. El cuarto es el ciclo económico expansivo, con el efecto composición antes descrito, y el envejecimiento de la población ocupada, porque los trabajadores mayores tienen más bajas y más largas.
    
El crecimiento de las bajas por salud mental merece una mención específica, porque plantea un dilema que no se resuelve con más control. Parte del aumento refleja una mayor visibilidad de trastornos que antes no se diagnosticaban ni se trataban, y una menor estigmatización de la salud mental, lo que es un avance. Parte refleja un deterioro real de la salud mental, especialmente entre los jóvenes y tras la pandemia. Y parte puede reflejar la dificultad de verificar objetivamente estos trastornos, lo que aumenta el margen de discrecionalidad. Las soluciones de control que funcionan para las bajas traumatológicas, como la evaluación funcional por las mutuas, son mucho menos aplicables a la salud mental, donde la recuperación depende del acceso a tratamiento psicológico, escaso en el sistema público.

\paragraph{Debate sobre la reforma}
El crecimiento del gasto ha abierto un debate de reforma en el que se enfrentan varias estrategias, cada una con fundamento teórico y con costes políticos distintos. 

La primera es la reincorporación progresiva. En junio de 2025, el Ministerio de Inclusión propuso a los agentes sociales un sistema de "altas progresivas": tras bajas largas, el trabajador podría reincorporarse durante un periodo de unos treinta días a media jornada, compatibilizando salario y prestación. Sindicatos y patronal recibieron la propuesta con frialdad por considerarla incompleta. Es la traducción española de la evidencia noruega de Markussen y otros sobre las bajas parciales, y su fundamento es sólido. Sus críticos, desde el lado sindical, temen que se convierta en una forma de presionar a los trabajadores para volver antes de estar recuperados. Desde el lado empresarial, temen los costes de organización que supone.

La segunda estrategia es institucional. En febrero de 2026, el Gobierno puso en marcha un Observatorio de la Incapacidad Temporal, con participación de los ministerios de Inclusión y de Sanidad, de las organizaciones empresariales y sindicales y de las Comunidades Autónomas. Su objetivo es analizar sistemáticamente la prestación, mejorar su gestión y apoyar el diseño de políticas basadas en la evidencia, con especial atención a la reincorporación progresiva en patologías graves como el cáncer. Es un intento de abordar el problema de coordinación identificado por la AIReF mediante el diálogo entre las instituciones implicadas, sin alterar el reparto de competencias.

La tercera estrategia es reforzar el papel de las mutuas en el control y la gestión, como se examinó en 1.4.4, con la experiencia gallega de mayo de 2026 como ejemplo más avanzado y controvertido. La cuarta, que defienden las organizaciones empresariales, que en el verano de 2026 intensificaron su presión sobre el Gobierno, es reforzar el control por el INSS y las mutuas y revisar el diseño de la prestación.

Menos presente en el debate español, pero muy relevante desde la teoría, sería optar por trasladar a las empresas una parte mayor del coste de las bajas, para que tengan incentivos a la prevención y a la reincorporación. Es la lógica de la tarificación por experiencia. El ejemplo internacional más citado es el de los Países Bajos, que tras la crisis de la incapacidad de los años ochenta (1.2.1) privatizaron prácticamente la prestación por enfermedad. Desde 2004 las empresas deben pagar el salario del trabajador enfermo durante dos años, y están obligadas a elaborar planes de reincorporación. El resultado fue una fuerte reducción del absentismo y de las entradas en la incapacidad permanente. El coste fue una mayor selección de riesgos en la contratación, porque las empresas tienen incentivos para evitar contratar a trabajadores con peor salud. 

En España, las empresas ya soportan los días cuarto a decimoquinto y las mejoras voluntarias de los convenios, pero no tienen capacidad para controlar las bajas, que dependen del servicio de salud. Combinación que alimenta la presión empresarial para reformar el sistema.

\subsubsection{Nacimiento y cuidado de menor}
La prestación por nacimiento y cuidado de menor protege la pérdida de rentas durante los periodos de suspensión del contrato por el nacimiento, la adopción, la guarda con fines de adopción o el acogimiento de un menor. 

Desde el Real Decreto-ley 9/2025, cada progenitor tiene derecho a 19 semanas. Las 6 primeras son obligatorias, deben disfrutarse a jornada completa inmediatamente después del parto y son simultáneas para ambos progenitores. Otras 11 pueden disfrutarse hasta que el menor cumpla doce meses. Las 2 restantes pueden disfrutarse hasta que el menor cumpla ocho años, para su cuidado. En las familias monoparentales, el permiso es de 32 semanas. Los derechos son individuales, iguales e intransferibles entre progenitores.\footnote{%
    Hasta 2019, la madre tenía un permiso de maternidad de 16 semanas, parcialmente transferible al padre, y el padre, un permiso de paternidad mucho más corto, que se introdujo en 2007 con dos semanas y se amplió progresivamente. El Real Decreto-ley 6/2019 equiparó gradualmente ambos permisos hasta alcanzar en 2021 las 16 semanas iguales e intransferibles para cada progenitor. El Real Decreto-ley 9/2025, de 29 de julio, convalidado por el Congreso en septiembre de 2025, amplió el permiso a 19 semanas por progenitor y a 32 en las familias monoparentales. Las dos semanas adicionales utilizables hasta los ocho años se aplican a los nacimientos desde el 2 de agosto de 2024 y pueden solicitarse desde el 1 de enero de 2026. La ampliación retribuye una parte del permiso parental previsto por la Directiva (UE) 2019/1158, relativa a la conciliación de la vida familiar y la vida profesional. El objetivo de ambas reformas era la corresponsabilidad y la igualdad de género en el mercado de trabajo.
}

La prestación económica es del 100 por ciento de la base reguladora, sin límite de duración distinto del propio permiso. Los requisitos de cotización son nulos para los menores de 21 años y crecientes con la edad para los demás. Existe también un subsidio no contributivo para las madres que no reúnen el periodo de cotización exigido. 

Como se estableció, esta prestación se considera un gasto "impropio" del sistema contributivo, porque responde a una política de familia y no a un seguro de rentas, y desde 2021 se financia con transferencias del Estado.



#### La economía de los permisos: penalización por maternidad, corresponsabilidad y natalidad

La transformación de la antigua prestación por maternidad en un permiso igual e intransferible para ambos progenitores responde a un diagnóstico económico preciso: la penalización por maternidad (child penalty). Kleven, Landais y Søgaard (2019), con datos administrativos daneses, mostraron que el nacimiento del primer hijo provoca una caída permanente de los ingresos de las madres, en torno al 20 por ciento a largo plazo, mientras que los ingresos de los padres apenas se alteran. Kleven, Landais, Posch, Steinhauer y Zweimüller (2019) extendieron el análisis a varios países y encontraron penalizaciones mucho mayores en los países de lengua alemana y en el Reino Unido que en los nórdicos. Para España, De Quinto, Hospido y Sanz (2021) estimaron una penalización considerable en los ingresos de las madres tras el primer hijo, sin efectos equivalentes en los padres. La penalización por maternidad es, según esta literatura, la principal causa de la brecha salarial de género persistente en los países desarrollados, y a través de ella, de la brecha en las pensiones.

Los permisos de paternidad intransferibles se justifican como una forma de reducir esa penalización, por dos vías. La primera es que, al hacer que ambos progenitores se ausenten del trabajo en torno al nacimiento, reducen la discriminación estadística de los empleadores contra las mujeres en edad fértil, porque el coste esperado de contratar a una mujer y a un hombre se iguala. La segunda es que, al implicar a los padres en el cuidado desde el principio, modifican el reparto de las tareas domésticas a largo plazo.

No obstante, hay un argumento adicional, directamente ligado al objeto del tema, que suele omitirse. En un sistema de pensiones de reparto, los hijos son un bien público para el sistema, porque las cotizaciones de los hijos de cada generación financian las pensiones de todos los jubilados, tengan o no hijos. Esta idea tiene una literatura teórica consolidada, con trabajos como el de Van Groezen, Leers y Meijdam (2003), que mostraron que en un sistema de reparto la decisión de tener hijos genera una externalidad positiva no internalizada: las familias soportan el coste de criar a los hijos, pero el beneficio de sus futuras cotizaciones se reparte entre todos los pensionistas. En consecuencia, sin intervención, la natalidad es inferior a la socialmente óptima, y las prestaciones familiares y los permisos de cuidado pueden justificarse como una corrección de esa externalidad. Los autores describieron las pensiones y las prestaciones por hijo como "gemelos siameses". Sinn (2000), examinado en 2.1.1, llevó el argumento más lejos y propuso que la pensión de cada individuo dependiera en parte de su número de hijos, o que quienes no los tienen ahorraran más en un pilar de capitalización. (OJO: la relación entre natalidad y sostenibilidad del sistema de reparto es el núcleo del factor demográfico que se examina en 4.2.1. España tiene una de las tasas de fecundidad más bajas del mundo, en torno a 1,1 hijos por mujer; la cifra exacta de 2025 debe verificarse con el INE.)

\begin{specialblock}{\textbf{Evidencia empírica:} ¿qué nos dice la investigación económica sobre estos permisos?}
    La evidencia empírica sobre los efectos de los permisos de paternidad es rica. Patnaik (2019), con la introducción en Quebec de una cuota de permiso reservada al padre, encontró un aumento persistente de la participación de los padres en las tareas domésticas y del empleo de las madres. Cools, Fiva y Kirkebøen (2015), con la cuota paterna noruega, encontraron efectos más modestos. Dahl, Løken y Mogstad (2014) mostraron que la decisión de un padre de tomar el permiso se contagia a sus compañeros de trabajo y a sus hermanos: un efecto de pares que amplifica el impacto de la política. Para España, Farré y González (2019) analizaron la introducción en 2007 de un permiso de paternidad de dos semanas y encontraron un resultado inesperado: el permiso redujo la fecundidad posterior de las familias afectadas. Lo atribuyeron a que los padres, al implicarse más en el cuidado, tomaron mayor conciencia de su coste, y a que las madres retrasaron su segundo hijo para afianzar su carrera. El resultado ilustra que los permisos persiguen objetivos que pueden entrar en tensión: la igualdad de género y la natalidad.

    La literatura sobre la duración de los permisos matiza sus beneficios. Ruhm (1998), con datos europeos, mostró que los permisos cortos aumentan el empleo de las madres, porque les garantizan la vuelta al trabajo, pero que los permisos muy largos reducen sus salarios, porque deprecian su capital humano y las alejan de la carrera profesional. Lalive y Zweimüller (2009), con una reforma austriaca que alargó el permiso parental, encontraron un aumento de la fecundidad a corto plazo y una reducción del empleo de las madres. La ampliación española a 19 semanas por progenitor, y en particular las dos semanas utilizables hasta los ocho años, se sitúa en un rango de duración que la literatura considera compatible con el empleo femenino, y su carácter igual e intransferible reduce el riesgo de que alargue la penalización por maternidad.
\end{specialblock}

\subsubsection{Riesgo durante el embarazo y la lactancia y el cuidado de menores con enfermedad grave}
Las prestaciones por riesgo durante el embarazo y por riesgo durante la lactancia natural protegen a la trabajadora cuyo puesto de trabajo supone un riesgo para su salud o la de su hijo cuando la empresa no puede adaptar su puesto ni trasladarla a otro compatible. En ese caso, su contrato se suspende y percibe una prestación del 100 por ciento de la base reguladora. 

Su singularidad es que se tratan como contingencias profesionales: derivan de las condiciones del puesto de trabajo, no de una enfermedad. Por eso no exigen periodo de carencia y su gestión corresponde, en las empresas asociadas, a las mutuas colaboradoras (1.4.4). Su justificación es la protección de la salud y la no discriminación: sin ella, el riesgo laboral del embarazo se convertiría en un coste para la empresa o en una causa de pérdida del empleo de la trabajadora.

La prestación por cuidado de menores afectados por cáncer u otra enfermedad grave protege a los progenitores que reducen su jornada, al menos en un 50 por ciento, para cuidar de forma directa, continua y permanente a un hijo hospitalizado o en tratamiento continuado por una enfermedad grave. Compensa la pérdida de salario derivada de la reducción con el 100 por ciento de la base reguladora en proporción a la reducción. Su ámbito de edad se ha ido ampliando en sucesivas reformas, y su gestión corresponde también a las mutuas en las empresas asociadas. 

Es un ejemplo de lo que la literatura de los nuevos riesgos sociales examinada en 1.2.1 identifica como protección frente a los riesgos de conciliación que el Estado del bienestar de posguerra no contemplaba.

\subsubsection{Prestaciones familiares}
Las prestaciones familiares de la Seguridad Social son hoy de alcance muy limitado y se concentran en situaciones específicas. La principal es la asignación económica por hijo o menor a cargo con discapacidad, que se concede por cada hijo con una discapacidad igual o superior al 33 por ciento, o al 65 por ciento si es mayor de edad. Se completa con una prestación económica de pago único por nacimiento o adopción en familias numerosas, monoparentales o con madres con discapacidad, y con una prestación por parto o adopción múltiples.\footnote{
La antigua asignación por hijo a cargo sin discapacidad, una prestación de cuantía muy reducida sujeta a un límite de ingresos, dejó de reconocerse a nuevos beneficiarios en 2021. Fue sustituida por el complemento de ayuda para la infancia del Ingreso Mínimo Vital.

Hasta 2021, la Seguridad Social concedía una asignación económica por cada hijo menor de 18 años a cargo, sin discapacidad, a las familias con ingresos inferiores a un límite, de cuantía muy reducida. La normativa del Ingreso Mínimo Vital impidió nuevos reconocimientos y la sustituyó por el complemento de ayuda para la infancia, de mayor cuantía y graduado por edad. El objetivo era concentrar el apoyo en la lucha contra la pobreza infantil.
}

El escaso desarrollo de las prestaciones familiares es uno de los rasgos más característicos del Estado del bienestar español, y conecta con varios hilos tratados. El gasto en la función de familia e infancia de España es claramente inferior a la media de la Unión Europea. Además, la pobreza infantil se sitúa de forma persistente entre las más altas de la Unión, en contraste con la reducción de la pobreza de los mayores. La explicación, en términos de 1.1.3, es la tradición familiarista del régimen de bienestar mediterráneo: la familia asumía el coste de la crianza sin apoyo público, y el Estado concentraba su esfuerzo en la protección de los trabajadores y de los mayores. En términos de la paradoja de la redistribución de Korpi y Palme, examinada en 1.1.2, las prestaciones familiares españolas han sido focalizadas y de escasa cuantía, con una coalición de apoyo débil, lo que ha limitado su desarrollo.

Una parte relevante del apoyo público a las familias españolas se canaliza, además, por la vía del bienestar fiscal de Titmuss examinado en 2.3.2. Destacan la deducción por maternidad del impuesto sobre la renta, que tras la reforma de 2023 se extendió a todas las madres con hijos menores de tres años con independencia de su situación laboral, y los mínimos por descendientes. Estas medidas tienen un alcance limitado para las familias con rentas bajas que no pagan impuesto sobre la renta, salvo en la medida en que la deducción sea reembolsable.

El debate sobre una prestación universal por crianza, una cantidad mensual por cada hijo menor sin condición de renta, al estilo de las prestaciones por hijo alemanas (Kindergeld) o de las de otros países europeos, ha estado presente en la política española de los últimos años sin concretarse. Sus defensores invocan la paradoja de la redistribución, la reducción de la pobreza infantil, la eliminación de la carga administrativa y de la baja utilización que afectan al complemento de ayuda para la infancia (72 por ciento de no solicitud según la AIReF, 3.1.5), y el argumento de la externalidad positiva de los hijos en un sistema de reparto. Sus críticos subrayan su coste, muy elevado por su carácter universal, y su escasa focalización en las familias con más necesidades. El paradigma de la inversión social examinado en 1.2.1 apunta en una dirección complementaria. Además de las transferencias monetarias, el apoyo a las familias exige servicios, como la educación infantil de 0 a 3 años, que favorecen el empleo de las madres y el desarrollo de los niños. Pero la crítica del efecto Mateo de Cantillon advierte que esos servicios tienden a beneficiar más a las familias de clase media.


\subsection{Protección por desempleo y el Sistema Nacional de Empleo}
España dedica una parte sustancialmente menor de su gasto total en empleo a las políticas activas que los países nórdicos que inspiran el modelo de flexiseguridad (flexicurity), con Dinamarca como referencia, donde el gasto en políticas activas supera de forma muy notable al español en proporción del PIB. [DOC] Añade que la crítica habitual señala que un sistema volcado desproporcionadamente en el gasto pasivo, sin un esfuerzo comparable en la reinserción efectiva, corre el riesgo de cronificar el desempleo estructural en lugar de resolverlo, y que este es el diagnóstico que está detrás de la última reforma del servicio de empleo.

El debate es relevante, pero es solo uno de los que plantea la protección por desempleo. Y no puede entenderse sin la teoría del seguro de desempleo, que es la pieza que falta en el documento. La protección por desempleo tiene una doble función: asegurar las rentas de quienes pierden su empleo (la función pasiva) y facilitar su vuelta al trabajo (la función activa). Ambas funciones están en tensión, porque cuanto más generosa es la protección de las rentas, menor es el incentivo a buscar y aceptar empleo. Resolver esa tensión es el problema central del diseño. España lo afronta en condiciones particularmente exigentes. En el segundo trimestre de 2026, según la Encuesta de Población Activa, la tasa de paro era del 9,87 por ciento, con unos 2,5 millones de desempleados y 22,8 millones de ocupados. Es la tasa más baja en casi dos décadas, pero todavía muy superior a la media europea, y con enormes diferencias territoriales: desde el 6,11 por ciento de las Islas Baleares hasta el 14,56 por ciento de Andalucía. [DOC] El documento recuerda que durante la crisis la tasa de paro creció del 10 al 25 por ciento, y el gasto en desempleo alcanzó el 3 por ciento del PIB.

El subapartado se organiza en cuatro partes: la teoría del seguro de desempleo (3.3.1), el marco jurídico e institucional, con la respuesta a tus notas sobre los sujetos y la desagregación vertical (3.3.2), las prestaciones vigentes (3.3.3) y la financiación, el gasto y el debate sobre las políticas activas (3.3.4).

El marco jurídico e institucional

#### La regulación vigente

[DOC, corregido] El documento define el Sistema Nacional de Empleo conforme al texto refundido de la Ley de Empleo aprobado por el Real Decreto Legislativo 3/2015: el conjunto de estructuras, medidas y acciones necesarias para promover y desarrollar la política de empleo. Sus objetivos serían, entre otros, mantener un sistema de protección frente al desempleo que comprende las políticas activas y las prestaciones, fomentar el empleo y la creación de puestos de trabajo, favorecer la integración y el progreso en el mercado laboral, garantizar la igualdad de oportunidades y mantener la unidad del mercado de trabajo en el territorio nacional. Ese texto refundido está derogado por la Ley 3/2023, de 28 de febrero, de Empleo, que es la norma vigente.\footnote{%
    El texto refundido de la Ley de Empleo de 2015 integró la Ley 56/2003 y sus modificaciones. La Ley 3/2023 lo sustituyó con el objetivo de reforzar los servicios de empleo, crear la Agencia Española de Empleo y cumplir uno de los compromisos del Plan de Recuperación, Transformación y Resiliencia.
} La nueva ley mantiene en lo esencial la definición y los objetivos del sistema, pero introduce cuatro novedades. Crea la Agencia Española de Empleo, que debe sustituir al SEPE. Define un catálogo de servicios garantizados a los demandantes de empleo. Identifica colectivos de atención prioritaria. Y establece la obligación de evaluar las políticas activas. La regulación de las prestaciones por desempleo se encuentra, en cambio, en el título III del texto refundido de la Ley General de la Seguridad Social. Recuérdese que el desempleo es formalmente una contingencia protegida por la Seguridad Social, según su art. 42 (1.3.2).








#### ¿Tiene sentido hablar de sujetos?

[DOC] Tu nota a este epígrafe planteaba precisamente esta pregunta. La respuesta es afirmativa, pero con una particularidad que distingue el desempleo de las demás prestaciones y que conviene explicar. En la sanidad, el sujeto protegido se define por la ciudadanía o la residencia: la protección es universal. En las pensiones contributivas, por la carrera de cotización. En el desempleo, el sujeto protegido se define por una situación jurídica: la situación legal de desempleo. Solo está protegido quien, pudiendo y queriendo trabajar, pierde su empleo de forma involuntaria: por despido, extinción del contrato temporal, finalización del periodo de prueba a instancias del empresario u otras causas tasadas. No lo está quien abandona voluntariamente su empleo. Este requisito es la traducción jurídica del riesgo moral examinado en 3.3.1: el seguro protege frente a un riesgo, la pérdida involuntaria del empleo, y no frente a una decisión.

El ámbito subjetivo se ha ampliado progresivamente. La protección contributiva comprende a los trabajadores por cuenta ajena del Régimen General y de los sistemas especiales, y desde octubre de 2022 incluye a las empleadas del hogar, que hasta entonces estaban excluidas.\footnote{%
    Hasta octubre de 2022, las trabajadoras del sistema especial de empleados del hogar no tenían derecho a la protección por desempleo, lo que el Tribunal de Justicia de la Unión Europea consideró una discriminación indirecta por razón de sexo en su sentencia de 24 de febrero de 2022 (asunto C-389/20). El Real Decreto-ley 16/2022 les reconoció la protección y estableció la correspondiente cotización.
} Los trabajadores autónomos no tienen protección por desempleo en sentido estricto, sino una prestación distinta, la prestación por cese de actividad, creada en 2010, gestionada por las mutuas colaboradoras (1.4.4) y financiada con una cotización específica. Su lógica es más compleja que la del desempleo asalariado, porque es difícil verificar si el cese de un autónomo es involuntario, y durante sus primeros años tuvo una utilización muy baja por la dureza de sus requisitos. Se flexibilizó durante la pandemia y en reformas posteriores.

Existe, por último, un tercer tipo de sujeto, que la Ley 3/2023 sitúa en el centro del sistema: el demandante de empleo, sea o no perceptor de prestaciones, como destinatario de las políticas activas. La ley identifica colectivos de atención prioritaria, como los jóvenes, los desempleados de larga duración, los mayores de 45 años, las personas con discapacidad y las víctimas de violencia de género, entre otros. Son los que requieren una atención específica de los servicios de empleo.

#### La desagregación vertical: por qué las prestaciones son estatales y las políticas activas autonómicas

[DOC] Tu nota pedía responder a cómo se distribuyen las competencias. [DOC] El documento señala que el Sistema Nacional de Empleo está integrado por el Servicio Público de Empleo Estatal y los servicios públicos de empleo de las Comunidades Autónomas. El SEPE es un organismo autónomo adscrito al Ministerio de Trabajo y Economía Social, que se encarga principalmente de la gestión y control de las prestaciones por desempleo, del análisis del mercado de trabajo y del registro público de ofertas, demandas y contratos, así como de las políticas de formación y colocación en los territorios que no las tienen transferidas.

El reparto responde al marco constitucional examinado en 1.3.1. El Estado tiene la competencia exclusiva sobre la legislación laboral, sin perjuicio de su ejecución por las Comunidades Autónomas (art. 149.1.7ª), y sobre el régimen económico de la Seguridad Social (art. 149.1.17ª). En consecuencia, las prestaciones por desempleo, que forman parte de la acción protectora de la Seguridad Social y se financian con cotizaciones, son de gestión estatal a través del SEPE. Las políticas activas de empleo, que son ejecución de la legislación laboral, se transfirieron a las Comunidades Autónomas a lo largo de las décadas de los noventa y dos mil.

Este reparto tiene una justificación económica que la teoría del federalismo fiscal permite formular. La protección de rentas debe centralizarse por las mismas razones que las pensiones (1.3.1): el desempleo tiene una fuerte dimensión territorial, y la caja única permite repartir entre regiones el riesgo de que una de ellas sufra una crisis. Persson y Tabellini (1996) analizaron este problema en su trabajo sobre las constituciones fiscales federales. Mostraron que la centralización del seguro frente a shocks regionales mejora el reparto de riesgos, pero genera un riesgo moral territorial: las regiones tienen menos incentivos a adoptar políticas que reduzcan su desempleo si el coste de las prestaciones lo soporta el conjunto. Las políticas activas, en cambio, deben descentralizarse por el argumento de Oates examinado en 1.3.1: requieren un conocimiento detallado del mercado de trabajo local, de sus empresas y de sus necesidades de formación, que las administraciones autonómicas tienen en mayor medida que el Estado.

La combinación de ambas lógicas genera, sin embargo, una externalidad fiscal vertical análoga a la examinada en la incapacidad temporal (3.2.1). Las Comunidades Autónomas gestionan las políticas activas que deberían reducir la duración del desempleo, pero el ahorro en prestaciones que resulta de su éxito lo obtiene el Estado. Una comunidad que invierte en reinsertar rápidamente a sus desempleados no se beneficia directamente del ahorro en prestaciones que genera. El sistema intenta corregir este problema distribuyendo una parte de los fondos estatales para políticas activas entre las comunidades en función del cumplimiento de objetivos, mediante la Conferencia Sectorial de Empleo y Asuntos Laborales, pero el peso de este criterio ha sido limitado. [DOC] Tus notas pedían completar este análisis con el gasto por subsector: como se estableció en 2.3.1, las prestaciones por desempleo se imputan al subsector de Administraciones de Seguridad Social a través del SEPE, y las políticas activas, en su mayor parte, al de las Comunidades Autónomas.

#### El SEPE y la Agencia Española de Empleo

[DOC, corregido] El documento señala que la Ley 3/2023 de Empleo ha planteado la transformación del SEPE en una Agencia Estatal, y afirma que es una figura que "se eliminó en 2015 con la LRJSP". Añade que esta transformación sugiere la necesidad de dotar a la gestión de las políticas de empleo de mayor autonomía y flexibilidad operativa de la que permite el régimen ordinario de organismo autónomo. Como se precisó en 1.4.3, la figura de la agencia estatal, suprimida por la Ley 40/2015, se recuperó en 2021, de modo que la transformación es jurídicamente posible. Sin embargo, más de tres años después de la aprobación de la Ley 3/2023, la Agencia Española de Empleo sigue sin constituirse, y el SEPE continúa funcionando como organismo autónomo. El retraso ilustra la dificultad de las reformas institucionales examinadas en 1.4: la creación de una agencia exige un estatuto, la adaptación de la plantilla y de los sistemas de información, y la negociación con los sindicatos de empleados públicos, y en ausencia de un impulso político sostenido, esas dificultades paralizan el proceso.

\subsubsection{Prestaciones por desempleo}

[DOC] El documento señala que existen en España más de diez tipos de prestaciones y ayudas por desempleo, de las que la principal es la prestación contributiva. La protección se estructura en dos niveles, cuyo origen es la Ley 31/1984 (1.2.2): un nivel contributivo, que protege a quienes han cotizado suficientemente, y un nivel asistencial, que protege a quienes no tienen derecho a la prestación contributiva o la han agotado, siempre que carezcan de rentas suficientes. A ellos se añade un mecanismo de protección del empleo en las crisis, los ERTE y el Mecanismo RED, que se examina al final.

\paragraph{Nivel contributivo}

[DOC] El documento describe la prestación contributiva como una prestación económica dirigida a los trabajadores que han perdido involuntariamente su empleo y tienen acumuladas suficientes cotizaciones por desempleo. Exige haber cotizado por desempleo al menos 360 días en los seis años anteriores, encontrarse en situación legal de desempleo e inscribirse como demandante de empleo. Exige además suscribir un acuerdo de actividad, por el que el beneficiario se compromete a buscar activamente empleo y a aceptar una colocación adecuada. [DOC] La cuantía es un porcentaje de la base reguladora, que se calcula con las bases de cotización de los últimos 180 días cotizados. Es el 70 por ciento durante los primeros seis meses y el 50 por ciento a partir del séptimo. La duración oscila entre un mínimo de cuatro meses y un máximo de dos años, según los días cotizados. Con 360 días cotizados se tiene derecho a 120 días de prestación, y cada 180 días adicionales de cotización añaden 60 días de prestación, hasta un máximo de 720 días con 2.160 días cotizados.\footnote{%
    La escala de duración, que vincula la prestación a los días cotizados en los seis años anteriores, se estableció en la Ley 31/1984 y se ha mantenido en lo esencial desde entonces. Su lógica es contributiva: cuanto más se ha cotizado, más larga es la protección.
}


La prestación está sometida a límites que se expresan en función del indicador público de renta de efectos múltiples (IPREM), que en 2026 sigue siendo de 600 euros mensuales, congelado por la prórroga presupuestaria. En 2026, la prestación mínima es de 560 euros sin hijos y de 749 euros con hijos. La máxima es de 1.225 euros sin hijos, de 1.400 euros con un hijo y de 1.575 euros con dos o más hijos. [DOC] El documento cifraba la cuantía media de la prestación contributiva en 2021 en 864 euros mensuales; la cifra actual está pendiente de verificar. Durante la percepción de la prestación, el SEPE ingresa en la Seguridad Social las cotizaciones del beneficiario, que mantiene así sus derechos a futuras prestaciones y a la pensión. [DOC] Son las "cotizaciones de desempleados" que el documento cifraba en el 8 por ciento de las cotizaciones del sistema, y que, como se precisó en 2.2.1, son un flujo interno entre dos entidades del mismo subsector.

Los topes de la prestación introducen un elemento redistributivo en un seguro que, en principio, es contributivo. El mínimo protege a los trabajadores de salarios bajos. El máximo, relativamente reducido, limita la protección de los trabajadores de salarios altos, cuya tasa de reemplazo efectiva es mucho menor que el 70 por ciento nominal. [Suposición: cálculo ilustrativo] Un trabajador sin hijos con una base reguladora de 3.000 euros mensuales tendría derecho, según el porcentaje, a 2.100 euros durante los primeros seis meses, pero cobra el máximo de 1.225 euros, lo que equivale a una tasa de reemplazo efectiva del 41 por ciento. En los términos de la regla de Baily-Chetty, el tope reduce el seguro precisamente para los trabajadores que, por tener más ahorros, menos lo necesitan, y cuya elasticidad de la duración del desempleo puede ser menor. Es una forma de focalizar el seguro en quienes sufren un mayor efecto liquidez.

\paragraph{Nivel asistencial: la reforma de los subsidios}

[DOC] El documento describe los subsidios por desempleo como prestaciones de tipo no contributivo dirigidas a los desempleados que han agotado la prestación contributiva o no reúnen las condiciones para recibirla. Menciona los de insuficiencia de cotización, los de mayores de 45 años que han agotado la contributiva y los de mayores de 55 años. Se refiere también a varios tipos de ayudas extraordinarias posteriores a la prestación contributiva y a los subsidios, dirigidas a colectivos con problemas especialmente graves de inserción, como los desempleados de larga duración, las personas con discapacidad o las personas en riesgo de exclusión. Esa descripción corresponde al sistema anterior a la reforma de 2024 y está en buena parte superada.

El Real Decreto-ley 2/2024, en vigor desde el 1 de noviembre de 2024, reformó en profundidad el nivel asistencial.\footnote{%
    Un primer intento de reforma, el Real Decreto-ley 7/2023, no fue convalidado por el Congreso en enero de 2024, por el voto en contra de uno de los socios parlamentarios del Gobierno, que se oponía a los cambios en la cotización para la jubilación de los mayores de 52 años. El Gobierno aprobó un nuevo texto, el Real Decreto-ley 2/2024, con modificaciones en ese punto, que sí fue convalidado. El episodio anticipa la fragilidad parlamentaria examinada en 1.2.2 con el decreto de revalorización de 2026. Antes de la reforma, los subsidios tenían una cuantía uniforme del 80 por ciento del IPREM, un mes de espera en algunos casos y una incompatibilidad casi total con el empleo. La reforma pretendía ampliar la cobertura, simplificar el sistema y eliminar la trampa de la inactividad. También respondía a un compromiso del Plan de Recuperación.
}
Simplificó los múltiples subsidios en dos modalidades principales: el subsidio por agotamiento de la prestación contributiva y el subsidio por cotización insuficiente, con un mínimo de tres meses cotizados para este último. Amplió el ámbito subjetivo a colectivos antes excluidos, entre ellos los menores de 45 años sin responsabilidades familiares que hayan agotado una prestación contributiva de al menos 360 días, los trabajadores eventuales agrarios y las víctimas de violencia de género o sexual. Estableció una cuantía decreciente en función del IPREM: el 95 por ciento durante los primeros seis meses (570 euros en 2026), el 90 por ciento entre el séptimo y el duodécimo mes (540 euros) y el 80 por ciento a partir del decimotercero (480 euros). La duración oscila entre 6 y 30 meses según la edad, las responsabilidades familiares y la duración de la prestación contributiva agotada. El acceso exige carecer de rentas superiores al 75 por ciento del salario mínimo interprofesional, 915,75 euros mensuales en 2026. La reforma suprimió además el mes de espera que existía antes de acceder a algunos subsidios, e integró en el nuevo sistema la renta activa de inserción, la antigua ayuda extraordinaria para colectivos de difícil inserción, que dejó de reconocerse a nuevos beneficiarios.

La reforma respondió a dos problemas que la teoría de 3.3.1 permite identificar. El primero era la fragmentación: la existencia de numerosos subsidios con requisitos distintos generaba lagunas de protección y una elevada carga administrativa, en los términos de Herd y Moynihan examinados en 3.1.5. El segundo era la trampa de la inactividad: los subsidios eran incompatibles con el empleo, de modo que aceptar un trabajo, sobre todo si era a tiempo parcial o temporal, suponía perder el subsidio. Para resolver el segundo problema, la reforma creó el complemento de apoyo al empleo (CAE), que permite compatibilizar el subsidio con un trabajo por cuenta ajena, a tiempo completo o parcial, durante un máximo de 180 días, con una cuantía decreciente. Es la aplicación a los subsidios de la lógica del impuesto negativo sobre la renta y del incentivo al empleo del Ingreso Mínimo Vital (3.1.5): al reducir el subsidio en menos de un euro por cada euro de salario, el trabajador gana dinero aceptando un empleo.

El subsidio para mayores de 52 años, que la reforma mantuvo con cambios, merece una atención específica, porque conecta el desempleo con las pensiones. Su cuantía es del 80 por ciento del IPREM durante toda su duración (480 euros en 2026), y se mantiene hasta que el beneficiario alcanza la edad de jubilación. Su rasgo distintivo es que el SEPE cotiza por el beneficiario a efectos de jubilación sobre una base superior a la mínima, lo que protege su futura pensión. Es en la práctica un puente hacia la jubilación para los desempleados de mayor edad con escasas posibilidades de reinserción. Su diseño plantea la tensión entre protección y activación examinada en 3.3.1: protege eficazmente a un colectivo vulnerable, pero puede desincentivar la búsqueda de empleo de quienes aún podrían reinsertarse.

[DOC] El documento menciona también el subsidio agrario y la renta agraria, dos prestaciones específicas para los trabajadores eventuales agrarios de Andalucía y Extremadura. Son herederas del antiguo empleo comunitario de los años setenta y ochenta y exigen un número de jornadas reales cotizadas muy inferior al de los demás subsidios. Su justificación original era la protección de una población rural con un empleo estacional y escaso. Su crítica, recurrente, es que se han convertido en un instrumento de sostenimiento de rentas en zonas rurales con escaso incentivo a la movilidad laboral y con riesgos de fraude en la acreditación de las jornadas. Es un ejemplo de la persistencia de programas diseñados para un contexto que ya no existe, en la línea de la tesis de Pierson examinada en 1.2.1.

#### Los ERTE y el Mecanismo RED: proteger el empleo en las crisis

[DOC] El documento señala que en 2020 el gasto en desempleo repuntó hasta el 3,7 por ciento del PIB por el gasto derivado de la COVID-19, en particular de los expedientes de regulación temporal de empleo (ERTE). Los ERTE permiten a las empresas suspender los contratos o reducir la jornada de sus trabajadores de forma temporal, por causas económicas, técnicas, organizativas o de producción, o por fuerza mayor. Los trabajadores afectados perciben una prestación por desempleo. Durante la pandemia, el Estado los reforzó con exoneraciones de cotizaciones para las empresas y con prestaciones que no consumían el derecho futuro de los trabajadores. En el momento álgido llegaron a proteger a más de tres millones y medio de trabajadores, como en el Kurzarbeit alemán, examinado en 1.2.1. La evaluación de los ERTE de la pandemia ha sido, en general, favorable. Evitaron una destrucción masiva de empleo y permitieron una recuperación rápida de la actividad, al mantener el vínculo entre trabajadores y empresas y preservar el capital humano específico. Es la lógica de Estevez-Abe, Iversen y Soskice examinada en 1.1.3.

La reforma laboral de 2021 institucionalizó esta experiencia mediante el Mecanismo RED de Flexibilidad y Estabilización del Empleo.\footnote{%
    Antes de la reforma laboral de 2021, los ERTE estaban regulados en el Estatuto de los Trabajadores, pero con escasa utilización fuera de la pandemia. Durante esta, sucesivos reales decretos-leyes establecieron regímenes extraordinarios. El Real Decreto-ley 32/2021 los incorporó de forma permanente al Estatuto de los Trabajadores, creó el Mecanismo RED y vinculó las exoneraciones de cotizaciones a la formación de los trabajadores, con el objetivo de convertir la reducción de jornada en el instrumento preferente de ajuste frente al despido.
} Tiene dos modalidades. La cíclica se activa por el Consejo de Ministros ante una coyuntura macroeconómica general adversa. La sectorial se activa ante cambios permanentes en un sector que exijan procesos de recualificación y de transición profesional de sus trabajadores. En ambos casos, las empresas pueden aplicar medidas de reducción de jornada o de suspensión de contratos con exoneraciones de cotizaciones condicionadas a la formación de los trabajadores afectados, y estos perciben una prestación específica. El Mecanismo RED es la aplicación institucional de la teoría de Landais, Michaillat y Saez examinada en 3.3.1: una protección más generosa en las recesiones, cuando el coste de riesgo moral es menor, y ligada a la preservación del empleo y a la recualificación.

\subsubsection{Financiación, el gasto y el debate sobre las políticas activas}

#### La financiación y el gasto

[DOC] El documento señala que los ingresos del SEPE provienen mayoritariamente de las cotizaciones sociales (80 por ciento) y de transferencias corrientes de otras Administraciones Públicas (20 por ciento). La cotización por desempleo, del 7,05 por ciento para los contratos indefinidos y del 8,3 por ciento para los temporales en 2026 (2.2.1), financia la prestación contributiva, una parte de las políticas activas y las bonificaciones a la contratación. Las transferencias del Estado financian los subsidios del nivel asistencial. El tipo más elevado para los contratos temporales es una forma elemental de tarificación por experiencia en el sentido de Blanchard y Tirole: los contratos con mayor probabilidad de terminar en desempleo pagan más por el seguro.

[DOC] El documento recoge que el gasto en desempleo alcanzó el 3 por ciento del PIB durante la crisis financiera, cuando la tasa de paro creció del 10 al 25 por ciento, y el 3,7 por ciento en 2020 por los ERTE. [DOC] Según el informe anual del SEPE, en 2021 el 70 por ciento del gasto en prestaciones correspondió al nivel contributivo y el 23 por ciento al asistencial; el resto, a la renta activa de inserción, el subsidio agrario y la renta agraria. Desde entonces, con la reducción del paro por debajo del 10 por ciento, el gasto en prestaciones por desempleo se ha reducido sustancialmente en relación con el PIB. Pero sigue siendo superior a la media europea, que en 2023 se situaba en el 1,2 por ciento del PIB según la COFOG (2.3.2). Las cifras exactas del gasto, del número de beneficiarios y de la tasa de cobertura en 2026 deben verificarse en la estadística del SEPE. El gasto en desempleo es el componente más cíclico del gasto social español, lo que refleja su función de estabilizador automático y la mayor volatilidad del empleo en España. [DOC] El documento lo ilustra con un gráfico que muestra cómo el gasto en desempleo siguió la evolución de la tasa de paro durante la crisis.

\paragraph{Debate sobre las políticas activas}

[DOC] El documento plantea, como se ha visto, el debate sobre el reparto entre políticas activas y pasivas. Para evaluarlo con rigor hay que responder a una pregunta previa: si las políticas activas funcionan. La referencia principal es el metaanálisis de Card, Kluve y Weber, publicado en su primera versión en 2010 y ampliado en 2018. Revisa más de doscientas evaluaciones de programas de políticas activas en todo el mundo. Sus conclusiones son matizadas y conviene conocerlas con precisión. Los programas de ayuda a la búsqueda de empleo y de orientación tienen efectos positivos a corto plazo, pero pequeños. Los programas de formación tienen efectos pequeños o nulos a corto plazo, porque retiran a los participantes de la búsqueda de empleo mientras se forman, pero positivos y crecientes a medio y largo plazo. Los programas de empleo público directo tienen efectos pequeños o incluso negativos, porque no mejoran la empleabilidad de los participantes en el mercado ordinario. Los incentivos a la contratación en el sector privado tienen efectos variables. Los efectos son mayores para los desempleados de larga duración y en las recesiones, y el diseño concreto de cada programa importa más que su tipo.

A la luz de esta evidencia, el diagnóstico del documento debe matizarse. El problema de las políticas activas españolas no es solo su escaso volumen en comparación con los países nórdicos, sino también su composición y su eficacia. Históricamente, una parte importante del gasto en políticas activas en España se ha destinado a bonificaciones a la contratación, cuyo elevado peso muerto documentó la AIReF (2.2.1), y a programas de formación para el empleo con escasa evaluación y frecuentes problemas de gestión. En cambio, la orientación e intermediación, que la evidencia muestra más eficaces a corto plazo, ha recibido recursos limitados. La AIReF, en su revisión del gasto en políticas activas, identificó como problemas principales la escasa cultura de evaluación, la débil coordinación entre el Estado y las Comunidades Autónomas y la baja capacidad de los servicios públicos de empleo para intermediar en el mercado de trabajo. Un gancho que suele citarse: la proporción de contrataciones que se producen a través de los servicios públicos de empleo en España ha sido tradicionalmente muy reducida, del orden de un pequeño porcentaje del total. Los trabajadores encuentran empleo sobre todo a través de contactos personales, de las propias empresas o de agencias privadas. La cifra exacta debe verificarse. La Ley 3/2023, con su catálogo de servicios garantizados, su atención a los colectivos prioritarios y su obligación de evaluación, intenta responder a estas debilidades, pero su aplicación depende de los servicios públicos de empleo autonómicos y de una Agencia Española de Empleo que todavía no existe.

El debate sobre la protección por desempleo se completa con una propuesta que se discutió en España en 2020 y 2021 y que conecta con la teoría de Blanchard y Tirole: la "mochila austriaca". Austria sustituyó en 2003 las indemnizaciones por despido por un sistema de cuentas individuales financiadas con una contribución empresarial mensual. El trabajador conserva su cuenta al cambiar de empleo y puede utilizarla al ser despedido o acumularla para la jubilación. Sus defensores sostienen que reduce el coste del despido para las empresas y elimina la penalización a la movilidad voluntaria de los trabajadores, que pierden su antigüedad al cambiar de empleo en el sistema de indemnizaciones tradicional. Sus críticos advierten que reduce la protección efectiva de los trabajadores despedidos y que su coste de transición es elevado, porque las empresas tendrían que financiar a la vez las nuevas cuentas y las indemnizaciones acumuladas por los trabajadores actuales. La propuesta se abandonó en España, y la reforma laboral de 2021 optó por combatir la dualidad restringiendo la contratación temporal, en lugar de reformar las indemnizaciones.

En conjunto, la protección por desempleo española ha avanzado en los últimos años hacia un diseño más coherente con la teoría. Dispone de una prestación contributiva decreciente, de subsidios simplificados y compatibles con el empleo y de un mecanismo institucionalizado de protección del empleo en las crisis. Sus debilidades persistentes se sitúan en el lado activo: unos servicios de empleo con escasa capacidad de intermediación, una coordinación territorial deficiente con una externalidad fiscal sin corregir y una cultura de evaluación todavía incipiente. Esas debilidades explican, en parte, que España siga teniendo, incluso en la fase expansiva más larga de su historia reciente, una tasa de paro que duplica la media europea.

\subsubsection{La economía del seguro de desempleo}
El desempleo es el ejemplo más claro de riesgo que el mercado privado no puede asegurar, por las razones examinadas en 1.1.1. La primera es la selección adversa: los trabajadores conocen mejor que cualquier asegurador su probabilidad de perder el empleo, y un seguro voluntario atraería sobre todo a los de mayor riesgo. La segunda es el riesgo moral: el asegurado influye en la probabilidad de perder el empleo y, sobre todo, en la duración del desempleo, con su esfuerzo de búsqueda y su disposición a aceptar ofertas. La tercera, y la más decisiva, es que el desempleo es un riesgo agregado, fuertemente correlacionado entre trabajadores. En una recesión, muchos trabajadores pierden su empleo a la vez, y un asegurador privado que hubiera vendido seguros de desempleo sufriría pérdidas simultáneas en toda su cartera, que no podría diversificar. Por eso no existe prácticamente en ningún país un mercado privado de seguros de desempleo, y la protección es pública y obligatoria en todas las economías avanzadas.

\paragraph{El seguro óptimo: la regla de Baily-Chetty}

El marco teórico de referencia es la regla de Baily-Chetty, formulada por Baily (1978) y generalizada por Chetty (2006). Es la aplicación más conocida del marco general de Chetty y Finkelstein examinado en 1.1.1. La regla establece que el nivel óptimo de la prestación iguala, en el margen, dos magnitudes. La primera es el valor del seguro, que depende de cuánto cae el consumo del trabajador al perder su empleo y de su aversión al riesgo: cuanto mayor es la caída del consumo, más valioso es el seguro. La segunda es el coste del riesgo moral, que depende de la elasticidad de la duración del desempleo respecto de la prestación: cuanto más se alarga el desempleo al aumentar la prestación, más costosa es esta.

Chetty (2008) introdujo una distinción que cambió la interpretación de la evidencia. Parte del aumento de la duración del desempleo cuando la prestación es más generosa no se debe al riesgo moral, sino a un efecto liquidez. Los trabajadores sin ahorros suficientes se ven obligados a aceptar la primera oferta disponible, aunque sea mala, porque no pueden permitirse seguir buscando; la prestación les permite buscar un empleo mejor ajustado a sus capacidades. Chetty estimó que una parte sustancial del efecto de la prestación sobre la duración del desempleo en Estados Unidos correspondía a este efecto liquidez, que no es una distorsión, sino una ganancia de bienestar. Nekoei y Weber (2017), con datos austriacos, confirmaron que prestaciones más largas pueden mejorar la calidad del emparejamiento entre trabajadores y empleos, con salarios más altos en el nuevo empleo.

\paragraph{Perfil temporal de la prestación}

La teoría ha examinado también cómo debe evolucionar la prestación a lo largo del periodo de desempleo. Shavell y Weiss (1979) y, sobre todo, Hopenhayn y Nicolini (1997) mostraron que, en presencia de riesgo moral, la prestación óptima debe ser decreciente en el tiempo. Una prestación que se reduce con la duración del desempleo mantiene el incentivo a buscar empleo pronto, sin dejar al trabajador desprotegido al principio, cuando la búsqueda es más eficaz. El diseño español, con una prestación contributiva del 70 por ciento de la base reguladora los primeros seis meses y del 50 por ciento después, y con un subsidio posterior también decreciente, responde a esta lógica.

La evidencia empírica confirma que los trabajadores responden a los incentivos del perfil temporal. Katz y Meyer (1990), con datos estadounidenses, documentaron un fuerte aumento de la tasa de salida del desempleo en las semanas inmediatamente anteriores al agotamiento de la prestación. Card, Chetty y Weber (2007), con datos austriacos, encontraron el mismo fenómeno, aunque de menor magnitud, y mostraron que parte de él se debía a que los trabajadores retrasaban la fecha de inicio del nuevo empleo hasta agotar la prestación. Para España, Rebollo-Sanz y Rodríguez-Planas (2020) aprovecharon la reforma de 2012 que redujo la prestación contributiva del 60 al 50 por ciento de la base reguladora a partir del séptimo mes. Encontraron que la reducción aceleró la salida del desempleo, sin deteriorar de forma significativa la calidad de los nuevos empleos. La referencia exacta debe verificarse.\footnote{%
    Hasta julio de 2012, la prestación contributiva era del 70 por ciento de la base reguladora durante los primeros seis meses y del 60 por ciento a partir del séptimo. El Real Decreto-ley 20/2012, en el marco del ajuste fiscal de la crisis, redujo el porcentaje al 50 por ciento a partir del séptimo mes para las nuevas prestaciones. El objetivo era reducir el gasto y aumentar el incentivo a la búsqueda de empleo.
}




\paragraph{Seguro de desempleo en la macroeconomía}
Además de su función de seguro individual, la protección por desempleo es uno de los estabilizadores automáticos más importantes de la economía. En una recesión, el gasto en prestaciones aumenta automáticamente y sostiene la demanda agregada, y en una expansión se reduce. Landais, Michaillat y Saez (2018) mostraron que, cuando el desempleo se debe a una insuficiencia de la demanda y no a la falta de esfuerzo de búsqueda, el coste de riesgo moral de la prestación es menor, porque un trabajador que busca más no crea un empleo nuevo, sino que le quita el puesto a otro. Por tanto, la prestación óptima debe ser anticíclica: más generosa en las recesiones y menos en las expansiones. Es la justificación teórica de las ampliaciones extraordinarias de la protección en las crisis, como los programas de reducción de jornada de la pandemia examinados en 1.2.1 y 3.3.3.

\paragraph{Seguro de desempleo, protección del empleo y dualidad}

La protección por desempleo no puede analizarse separada de la protección del empleo, es decir, de las indemnizaciones y restricciones que la legislación impone al despido. Ambas son formas alternativas de proteger a los trabajadores frente al riesgo de perder su empleo: la primera les compensa después de perderlo, y la segunda dificulta que lo pierdan. Blanchard y Tirole (2008) analizaron su diseño conjunto y mostraron que, si el seguro de desempleo es público, las empresas que despiden imponen un coste al sistema que no internalizan. Por eso la solución óptima combina un seguro de desempleo con un impuesto sobre el despido, cuyos ingresos lo financian, en la línea de la tarificación por experiencia examinada en 1.4.4. En este marco, las indemnizaciones por despido pueden interpretarse como un sustituto imperfecto de ese impuesto.

[DOC] El modelo de flexiseguridad que el documento menciona es la aplicación práctica más citada de esta idea. El "triángulo dorado" danés, descrito por Madsen (2004), combina tres elementos: una baja protección del empleo, que facilita el despido y la contratación; un seguro de desempleo generoso, que protege las rentas de quienes pierden el empleo; y unas políticas activas intensas, que facilitan su rápida reinserción. El trabajador danés no está protegido en su puesto de trabajo, sino en el mercado de trabajo. Su seguridad no depende de conservar un empleo concreto, sino de la facilidad para encontrar otro y de la protección de sus rentas mientras tanto.

España ha tenido durante décadas la combinación opuesta, que Bentolila, Cahuc, Dolado y Le Barbanchon (2012) analizaron en su comparación entre Francia y España durante la gran recesión: un mercado de trabajo dual. Los trabajadores indefinidos tenían una alta protección del empleo, con indemnizaciones por despido elevadas. Los temporales tenían una protección mínima y concentraban la destrucción de empleo en las recesiones. Es la dualización examinada en 1.1.3. Las consecuencias sobre la protección por desempleo eran perversas. Los temporales, con carreras de cotización cortas e intermitentes, accedían a prestaciones contributivas breves o solo a subsidios, mientras que los indefinidos, que rara vez perdían el empleo, acumulaban derechos elevados. La reforma laboral de 2021 redujo drásticamente la temporalidad y alteró este equilibrio, con efectos colaterales sobre la incapacidad temporal examinados en 3.2.1.

\subsection{Sistema Nacional de Salud}
El Sistema Nacional de Salud se define como el conjunto de servicios de salud que prestan el Estado y las Comunidades Autónomas, que integra todas las funciones y prestaciones sanitarias que son responsabilidad de los poderes públicos. Recuerda que la Constitución, en su art. $43$, reconoce el derecho a la protección de la salud y a la atención sanitaria de todos los ciudadanos.

La inclusión de la sanidad en un tema sobre la Seguridad Social exige una justificación, porque hoy la sanidad no forma parte del sistema de Seguridad Social en sentido financiero ni orgánico, como se estableció en 1.4.1. Formalmente, el art. 42 del texto refundido de la Ley General de la Seguridad Social sigue incluyendo la asistencia sanitaria en la acción protectora. Pero desde finales de los noventa se financia íntegramente con impuestos generales, y desde 2002 la gestionan los servicios de salud de las Comunidades Autónomas. Su historia es, precisamente, la de su salida del sistema de Seguridad Social: el paso del modelo bismarckiano, en el que la asistencia sanitaria era una prestación del seguro obligatorio de enfermedad para los trabajadores y sus familias, al modelo beveridgiano de un servicio nacional de salud universal financiado con impuestos, en la estela del National Health Service británico (1.2.1 y 1.2.2). Entender ese tránsito, su fundamento económico y sus consecuencias es lo que justifica el tratamiento de la sanidad en este tema, y es lo que el subapartado intenta hacer.

El subapartado se organiza en cuatro partes: la economía de la sanidad pública, que explica por qué el mercado no provee adecuadamente la asistencia sanitaria y qué determina su gasto (3.4.1); el marco jurídico, los sujetos y la desagregación vertical, en respuesta a tus notas (3.4.2); las prestaciones, con la cartera de servicios y la prestación farmacéutica (3.4.3); y el gasto, su sostenibilidad y los debates abiertos (3.4.4).

La regulación básica del Sistema Nacional de Salud empieza con la Ley 14/1986, de 25 de abril, General de Sanidad. La Ley 16/2003, de cohesión y calidad del Sistema Nacional de Salud. La Ley 29/2006, de garantías y uso racional de los medicamentos y productos sanitarios, hoy refundida en el texto refundido de 2015; la Ley 33/2011, General de Salud Pública; y el Real Decreto-ley 16/2012, de medidas urgentes para garantizar la sostenibilidad del Sistema Nacional de Salud. A ellas deben añadirse el Real Decreto-ley 7/2018, sobre el acceso universal al Sistema Nacional de Salud, y, en 2026, el decreto-ley de reforma del copago farmacéutico y el proyecto de Ley de Gestión Pública e Integridad del Sistema Nacional de Salud, que se examinan en 3.4.3 y 3.4.4.

#### Los sujetos: la universalidad y sus vaivenes

[DOC] El documento señala que, en general, son titulares del derecho a la protección de la salud y a la atención sanitaria todas las personas con nacionalidad española y las personas extranjeras que tengan establecida su residencia en territorio español. Para acceder a la asistencia con cargo a fondos públicos, los titulares deben encontrarse en alguno de estos supuestos: tener nacionalidad española y residencia habitual en España; tener reconocido el derecho a la asistencia sanitaria en España por cualquier otro título jurídico, aun sin residir en España, siempre que no exista un tercero obligado al pago; o ser persona extranjera con residencia en España sin obligación de acreditar la cobertura sanitaria por otra vía. [DOC] Añade que las personas extranjeras no registradas ni autorizadas como residentes en España tienen derecho a la protección de la salud y a la atención sanitaria en las mismas condiciones que las personas con nacionalidad española.

Esta descripción corresponde a la regulación vigente tras el Real Decreto-ley 7/2018, que culminó la universalización del sistema.\footnote{%
    Tras la Ley General de Sanidad de 1986, la universalización avanzó gradualmente, y la Ley 33/2011 de Salud Pública la extendió a todos los españoles residentes. El Real Decreto-ley 16/2012 introdujo la figura del "asegurado", ligó de nuevo el derecho a la asistencia a la afiliación o a otras situaciones tasadas y excluyó a los extranjeros en situación irregular, salvo en urgencias, embarazo y para los menores. El objetivo era contener el gasto en plena crisis. El Real Decreto-ley 7/2018 recuperó la universalidad basada en la ciudadanía y la residencia, con el objetivo de restablecer el carácter universal del sistema.
} El criterio de acceso es la ciudadanía o la residencia, y no la afiliación a la Seguridad Social, lo que confirma el carácter beveridgiano del sistema. La inclusión de los extranjeros en situación irregular, con algunas condiciones de acreditación, es la aplicación más extensa del principio de universalidad examinado en 1.3.3. Es también la más debatida políticamente. Sus defensores invocan razones de salud pública (una enfermedad contagiosa no tratada afecta a toda la población) y de eficiencia (la atención primaria es más barata que las urgencias hospitalarias a las que acudirían los excluidos). Sus críticos invocan el coste y el posible efecto llamada.

#### La desagregación vertical: por qué la sanidad sí se descentralizó

[DOC] El documento señala que se trata de un sistema descentralizado, cuyas competencias se reparten entre el Estado, las Comunidades Autónomas y las entidades locales. El proceso de descentralización se inició en 1981 con el traspaso a Cataluña y culminó en 2001. [DOC] Corresponden al Estado el establecimiento de las normas que fijan las condiciones y requisitos mínimos del sistema, la coordinación general, la sanidad exterior y la legislación y autorización de medicamentos y productos sanitarios. Corresponden a las Comunidades Autónomas la planificación sanitaria, la salud pública y la gestión y prestación de los servicios de salud, de modo que cada una cuenta con su propio servicio de salud. Las corporaciones locales tienen competencias en materia de salubridad. El Consejo Interterritorial del Sistema Nacional de Salud es el órgano de coordinación entre el Estado y las Comunidades Autónomas. Como se precisó en 1.2.2, los últimos traspasos, aprobados en diciembre de 2001 para las diez comunidades que aún no tenían la competencia, se hicieron efectivos el 1 de enero de 2002.

[DOC] Tu nota pedía responder a la pregunta de cómo se distribuyen las competencias y completarla con el gasto por subsector. El fundamento constitucional es el art. 149.1.16ª, que reserva al Estado las bases y la coordinación general de la sanidad y la legislación sobre productos farmacéuticos, y deja a las Comunidades Autónomas el desarrollo y la gestión. La comparación con las pensiones es reveladora y responde a la pregunta implícita en tus notas: ¿por qué la sanidad se descentralizó y las pensiones no? La respuesta está en la distinta naturaleza de ambas prestaciones. Las pensiones son transferencias monetarias de carácter redistributivo y asegurador, cuya centralización permite mancomunar el riesgo demográfico entre territorios, como se examinó en 1.3.1. La sanidad es un servicio en especie, cuya provisión exige un conocimiento detallado de las necesidades locales y una gestión de proximidad. En los términos de Oates, examinado en 1.3.1, la sanidad es el tipo de servicio en el que la descentralización genera ganancias de eficiencia, porque permite adaptar la provisión a las preferencias y circunstancias de cada territorio.

La literatura del federalismo fiscal identifica otras dos ventajas de la descentralización sanitaria. La primera es la competencia por comparación (yardstick competition): los ciudadanos evalúan la gestión de su gobierno autonómico comparándola con la de las comunidades vecinas, lo que incentiva la eficiencia. Besley y Case (1995) mostraron empíricamente este mecanismo en las políticas fiscales de los Estados norteamericanos. La segunda es la experimentación: las comunidades funcionan como "laboratorios" de políticas públicas, en la expresión clásica del juez Brandeis, y las innovaciones que tienen éxito en una pueden extenderse a las demás.

La descentralización tiene también costes, que la experiencia española ha puesto de manifiesto. El primero es la desigualdad territorial en el acceso y la calidad de la asistencia, visible en las listas de espera (3.4.4). El segundo es la restricción presupuestaria blanda. Kornai (1986) formuló el concepto, y Rodden (2006) lo aplicó al federalismo fiscal. Cuando los gobiernos subnacionales esperan que el gobierno central los rescate si incurren en déficits, tienen incentivos a gastar más de lo que pueden financiar. La experiencia española del Fondo de Liquidez Autonómico, creado en 2012 para financiar a las comunidades que no podían acceder a los mercados, se interpreta a menudo en esta clave, porque una parte sustancial de la deuda autonómica se generó por el gasto sanitario. El tercero es la dificultad de coordinación, que la pandemia de 2020 puso de manifiesto con respuestas dispares de las comunidades y tensiones entre estas y el Estado.

La sanidad se financia a través del sistema de financiación autonómica, sin recursos afectados.\footnote{%
    Hasta 2001, la sanidad transferida se financiaba mediante un sistema específico, separado del resto de la financiación autonómica, con recursos afectados. La Ley 21/2001 integró la financiación sanitaria en el sistema general de financiación autonómica, sin afectación, aunque con una garantía temporal de mantenimiento del gasto sanitario. La Ley 22/2009 mantuvo ese enfoque. El nuevo modelo aprobado por el Consejo de Política Fiscal y Financiera en septiembre de 2026 introduce instrumentos específicos para asegurar el reparto de los recursos sanitarios, y está pendiente de tramitación parlamentaria.
} Las comunidades reciben ingresos por los impuestos cedidos y por las transferencias del Estado, y deciden cuánto destinar a la sanidad. El sistema incluye mecanismos de nivelación, basados en la "población ajustada", que pondera la población por su edad para reflejar las distintas necesidades de gasto sanitario. En septiembre de 2026, el Consejo de Política Fiscal y Financiera aprobó, con el apoyo del Gobierno, Cataluña y Canarias y la oposición de la mayoría de las comunidades gobernadas por el Partido Popular, un nuevo modelo de financiación autonómica. Supone unos 21.000 millones de euros adicionales para las comunidades y eleva la cesión del impuesto sobre la renta y del IVA. Incorpora dos instrumentos para garantizar el reparto de los recursos sanitarios: un fondo de compensación interterritorial y un sistema de incentivos. El modelo debe tramitarse en las Cortes, donde su aprobación es incierta. Su relevancia para este tema es que la sostenibilidad del gasto sanitario, que crece con el envejecimiento y la tecnología, depende de la capacidad fiscal de las comunidades y de los mecanismos de nivelación del sistema de financiación, no del sistema de Seguridad Social.

[DOC] En cuanto al gasto por subsector, el documento señala que, del gasto sanitario total, entre el 70 y el 75 por ciento corresponde a las Administraciones Públicas y el resto al sector privado. Como se estableció en 2.3.1, la inmensa mayoría del gasto sanitario público lo ejecutan las Comunidades Autónomas, con una participación menor de la Administración Central (el Ministerio de Sanidad, el INGESA y las mutualidades de funcionarios) y de las corporaciones locales. La proporción de gasto sanitario privado, en torno a una cuarta parte del total, es relativamente elevada para un sistema beveridgiano. Refleja el peso del seguro privado duplicado y de los pagos directos de los hogares, en particular en odontología, óptica y medicamentos, que se examinan en 3.4.4.

\subsubsection{Prestaciones}
El Real Decreto-ley 16/2012 diferencia tres modalidades de prestaciones. La cartera común básica de servicios asistenciales se cubre completamente con financiación pública. La cartera común suplementaria incluye las prestaciones de dispensación ambulatoria (como la prestación farmacéutica, la ortoprotésica o los productos dietéticos) y está sujeta a aportación del usuario. La cartera común de servicios accesorios incluye actividades no esenciales o de apoyo a patologías crónicas, sujetas también a aportación. [DOC] Las Comunidades Autónomas pueden aprobar sus propias carteras de servicios, que deben incluir las carteras comunes básica y suplementaria. La cartera común de servicios se detalla en el Real Decreto 1030/2006 y sus sucesivas actualizaciones. Su contenido se decide en el Consejo Interterritorial con apoyo de la Red Española de Agencias de Evaluación de Tecnologías Sanitarias, que valora la eficacia y el coste de las nuevas prestaciones antes de su inclusión.

La existencia de una cartera común garantiza un mínimo de prestaciones igual en todo el territorio. Es la concreción sanitaria de la competencia estatal sobre las condiciones básicas de igualdad del art. 149.1.1ª. Las carteras complementarias autonómicas permiten a las comunidades ampliar ese mínimo con cargo a sus propios recursos. Es una aplicación del modelo de federalismo examinado en 3.4.2: un suelo común garantizado por el Estado y un margen de diferenciación autonómica. Pero genera desigualdades en prestaciones como la atención dental, la salud mental o las ayudas a determinados tratamientos.

\paragraph{Prestación farmacéutica: financiación selectiva y precios de referencia}

[DOC] El documento explica que la financiación de la prestación farmacéutica está sometida al sistema de precios de referencia. Se definen conjuntos homogéneos de medicamentos con el mismo principio activo e idéntica vía de administración, entre los que debe figurar un genérico o un biosimilar, salvo que el medicamento o su principio activo lleven autorizados menos de diez años en la Unión. Para cada conjunto, el Estado fija un precio de referencia, que actúa como límite máximo de financiación pública, de modo que en cada conjunto existe un genérico de precio igual o inferior. El farmacéutico debe sustituir el medicamento prescrito por otro que no supere el precio de referencia si aquel lo supera, aunque el usuario puede optar por el medicamento prescrito pagando la diferencia.

El sistema de precios de referencia es un instrumento de contención del gasto que utiliza la competencia de los genéricos para presionar a la baja los precios de los medicamentos cuya patente ha expirado. La literatura de economía de la salud ha mostrado que es eficaz para reducir los precios, aunque su efecto depende del grado de competencia en cada conjunto y de la respuesta de los fabricantes de marca, que a menudo reducen sus precios hasta el de referencia para evitar la sustitución. Complementan el sistema la financiación selectiva, por la que el Ministerio decide qué medicamentos se financian con fondos públicos, y la negociación de precios de los medicamentos innovadores, cada vez más apoyada en los informes de posicionamiento terapéutico y en acuerdos de riesgo compartido con los fabricantes.

\paragraph{Copago farmacéutico: de 2012 a la reforma de 2026}

[DOC] El documento describe el sistema de aportación del usuario en la prestación farmacéutica ambulatoria que estableció el Real Decreto-ley 16/2012. Los medicamentos dispensados en el hospital no tienen copago. La aportación del usuario se fija según tres criterios (renta, edad y gravedad de la enfermedad) y puede ser del 0 por ciento (por ejemplo, para los parados que han perdido el derecho al subsidio y los perceptores de pensiones no contributivas), del 10 por ciento (para los pacientes con enfermedades graves y los pensionistas) o de entre el 40 y el 60 por ciento para los activos, en función de la renta, con límites máximos mensuales para los pensionistas. [DOC, corregido] El documento resume esta última horquilla como "entre el 0 y el 60 por ciento en función de la renta", lo que no es preciso: para los activos, los porcentajes eran del 40, el 50 y el 60 por ciento según los tramos de renta.

Ese sistema ha sido reformado en 2026 mediante un real decreto-ley que el Congreso convalidó el 28 de mayo de 2026, con 164 votos a favor, 33 en contra y 164 abstenciones.\footnote{%
    Antes de 2012, los pensionistas no pagaban por los medicamentos y los activos pagaban un 40 por ciento con independencia de su renta. El Real Decreto-ley 16/2012 introdujo el copago de los pensionistas, del 10 por ciento con topes mensuales, y graduó el de los activos según la renta (40, 50 o 60 por ciento, sin topes), con el objetivo de contener el gasto y aumentar la progresividad. La reforma de 2026, convalidada el 28 de mayo de ese año, añadió tramos y estableció topes para los activos de rentas bajas y medias, con el objetivo de corregir la inequidad entre activos y pensionistas.
} La reforma aumenta el número de tramos de renta y establece por primera vez topes mensuales para los trabajadores activos de rentas bajas y medias. Para los activos, fija seis tramos. Las rentas inferiores a 9.000 euros anuales pagan un 40 por ciento con un tope de 8,23 euros mensuales. Las rentas de 9.000 a 17.999 euros pagan un 40 por ciento con un tope de 18,52 euros. Las rentas de 18.000 a 34.999 euros pagan un 45 por ciento con un tope de 61,75 euros. Por encima, el porcentaje es del 45, el 50 o el 60 por ciento, sin tope. Para los pensionistas, se mantiene con carácter general el 10 por ciento, con topes mensuales graduados según la renta, salvo para las rentas superiores a 100.000 euros, que pagan un 60 por ciento con un tope de 61,75 euros. El Gobierno estima un ahorro de unos 500 euros anuales para los pacientes crónicos de renta baja.

El copago plantea el dilema clásico que el experimento de la RAND permitió cuantificar. Su justificación es doble: reducir el riesgo moral, es decir, el consumo innecesario de medicamentos que se produce cuando son gratuitos, y recaudar recursos para financiar el sistema. Su riesgo es que reduzca también el consumo necesario, sobre todo entre los pacientes con menos renta y más enfermedades crónicas, con consecuencias para su salud y, a largo plazo, para el propio gasto sanitario, si la falta de adherencia al tratamiento provoca hospitalizaciones evitables. La evidencia sobre la reforma de 2012 en España, analizada entre otros por Puig-Junoy y sus coautores, mostró que el copago redujo el consumo de medicamentos, especialmente entre los pensionistas, que pasaron de la gratuidad al 10 por ciento. Mostró también indicios de reducción de la adherencia en algunos tratamientos crónicos. La referencia exacta debe verificarse. La reforma de 2026, con sus topes mensuales para los trabajadores de rentas bajas y medias, pretende corregir la inequidad vertical del sistema anterior. En él, un trabajador activo con renta baja y una enfermedad crónica podía pagar en medicamentos mucho más que un pensionista de renta alta, porque los activos no tenían topes. Es una aplicación directa del principio de igualdad en su dimensión material examinado en 1.3.3.

\paragraph{Gasto farmacéutico}

[DOC] El documento señala que, según los datos del Ministerio de Hacienda, el gasto farmacéutico alcanzó en 2021 el 2,3 por ciento del PIB. Se distribuía en gasto farmacéutico hospitalario (31 por ciento), gasto en productos farmacéuticos y sanitarios por receta médica (46 por ciento) y gasto sin receta (23 por ciento). [DOC] Añade que las medidas de contención del gasto por receta se aplicaron desde 2009, de modo que en 2019 ese gasto se había reducido en tres décimas de PIB respecto de su máximo, pero que en 2020 y 2021 recuperó parte de la caída hasta situarse en torno al 1 por ciento del PIB.

El gasto farmacéutico público a través de receta médica del Sistema Nacional de Salud alcanzó en 2024 los 13.865 millones de euros, un 4,9 por ciento más que en 2023, lo que supone unos 285 euros por habitante. En 2025 siguió creciendo en torno al 4 por ciento, con un nuevo máximo histórico cercano a los 14.000 millones de euros. Las cifras de ambos años proceden de fuentes con métodos de cómputo algo distintos y deben homogeneizarse. El crecimiento del gasto se concentra cada vez más en el ámbito hospitalario, impulsado por los medicamentos oncológicos y de terapias avanzadas, de precios muy elevados. En el ámbito de la receta ha contribuido la expansión de los agonistas del receptor GLP-1, como la semaglutida, utilizados en la diabetes y cada vez más en la obesidad. Es la manifestación actual del efecto expansión del progreso tecnológico que el documento describe: tratamientos nuevos, eficaces y muy caros, que amplían el número de condiciones tratables y de pacientes tratados.

\subsubsection{Gasto, su sostenibilidad y los debates abiertos}

#### La evolución del gasto sanitario

[DOC] El documento describe la evolución del gasto sanitario público según el Sistema de Cuentas de Salud. Creció de forma continuada entre 2003 y 2009, hasta alcanzar un máximo del 7 por ciento del PIB. Las restricciones presupuestarias lo redujeron hasta el 6,6 por ciento en 2019. En 2020, por la pandemia, repuntó hasta el 8 por ciento. Según el Ministerio de Sanidad, el gasto sanitario público alcanzó en 2024 los 101.739 millones de euros, el 6,4 por ciento del PIB. El gasto ha recuperado el nivel previo a la pandemia en términos del PIB, aunque su volumen absoluto es muy superior. Como se estableció en 2.3.2, se sitúa algo por debajo de la media de la Unión Europea en porcentaje del PIB. La reducción del gasto entre 2009 y 2019 fue un caso de ajuste presupuestario en un servicio con demanda creciente, y sus consecuencias se manifestaron en el deterioro de las listas de espera, la reducción de plantillas y la congelación salarial de los profesionales.

[DOC] En cuanto a las proyecciones, el documento cita el Ageing Report de 2011, que preveía un crecimiento del gasto público en cuidados sanitarios en España del 5,7 por ciento del PIB en 2013 al 7,1 por ciento en 2060. Esa proyección está superada. El Ageing Report de 2024 proyecta para España un aumento moderado del gasto sanitario en porcentaje del PIB hasta 2070, en el escenario de referencia, con escenarios alternativos que incorporan distintos supuestos sobre la salud en las edades avanzadas y sobre el progreso tecnológico. La cifra exacta debe verificarse en el informe. Las proyecciones son muy sensibles a esos supuestos, lo que refleja la incertidumbre sobre la compresión o expansión de la morbilidad y sobre el ritmo de la innovación examinada en 3.4.1.

[DOC] El documento señala que el incremento esperado del gasto sanitario, común a los países desarrollados y asociado al envejecimiento, ha situado su control en un lugar prioritario del diseño de las políticas públicas. Enumera como medidas de contención el copago de servicios sanitarios, la delimitación de los servicios cubiertos, la mejora de la gestión, el control de la calidad, la optimización del gasto farmacéutico y el aumento de la eficiencia en la provisión. A la luz de 3.4.1, ese diagnóstico debe matizarse: el envejecimiento es solo uno de los factores del crecimiento del gasto, y no el principal. Las medidas de contención eficaces a largo plazo no son las que reducen el acceso (copagos, delimitación de la cartera), que tienen un efecto limitado y un coste en equidad. Son las que mejoran la eficiencia: la evaluación económica de las tecnologías antes de su incorporación, la gestión de las enfermedades crónicas en atención primaria y la reducción de la variabilidad injustificada de la práctica clínica.

#### Las listas de espera y el seguro privado

Las listas de espera son el principal indicador de las tensiones del sistema y el problema que más preocupa a los ciudadanos. A 31 de diciembre de 2025, 853.509 pacientes esperaban una intervención quirúrgica en el Sistema Nacional de Salud, con un tiempo medio de espera de 121 días, algo inferior a los 126 días de un año antes. Más de cuatro millones de pacientes esperaban una primera consulta con el especialista, con un tiempo medio de 102 días. Las diferencias territoriales son muy marcadas: Andalucía registraba el tiempo medio de espera quirúrgica más elevado, 173 días, y Cantabria, Cataluña y La Rioja, las tasas más altas de pacientes en espera por habitante. Las listas de espera son un mecanismo de racionamiento no monetario: cuando el precio de un servicio es nulo y la oferta es limitada, la demanda se ajusta mediante el tiempo de espera. Es una alternativa al copago que no discrimina por renta, pero que impone costes elevados en salud y en bienestar, y que, como vimos en 3.2.1, prolonga las bajas laborales.

Las listas de espera han impulsado el crecimiento del seguro privado de salud en España, que cubre ya a una proporción muy significativa de la población, con más de diez millones de asegurados. La cifra exacta debe verificarse con los datos del sector. En la mayoría de los casos se trata de un seguro duplicado: los asegurados mantienen su derecho a la sanidad pública y pagan además un seguro privado para acceder más rápidamente a especialistas y pruebas diagnósticas. Este fenómeno plantea un problema de equidad que la literatura conoce bien. Tudor Hart (1971) formuló la "ley de cuidados inversos" (inverse care law): la disponibilidad de una buena asistencia sanitaria tiende a variar inversamente con la necesidad de la población atendida. Cuando las clases medias y altas resuelven sus necesidades en el sector privado, disminuye su interés en la calidad del sistema público, lo que conecta con la paradoja de la redistribución de Korpi y Palme (1.1.2): un sistema público que pierde a las clases medias pierde también su apoyo político. El mutualismo administrativo y su crisis de 2024-2025, examinados en 1.4.5, son la manifestación institucional de este mismo fenómeno.

#### La gestión: del modelo Alzira a la Ley de Gestión Pública

La forma de gestión de los servicios sanitarios es el último gran debate del sistema. La Ley 15/1997, sobre habilitación de nuevas formas de gestión del Sistema Nacional de Salud, permitió que las comunidades gestionaran los servicios sanitarios mediante fórmulas diversas, incluidas las concesiones administrativas y los conciertos con entidades privadas. El ejemplo más conocido fue el modelo Alzira de la Comunidad Valenciana: una empresa privada construía y gestionaba un hospital y la atención sanitaria de un área de salud a cambio de un pago capitativo por habitante. La Generalitat Valenciana revirtió su gestión a la administración pública en 2018. La evidencia sobre la eficiencia de estas fórmulas es controvertida. Sus defensores señalan menores costes por habitante y mayor productividad. Sus críticos, la selección de pacientes, la derivación de los casos más costosos a los hospitales públicos y la falta de transparencia.

El 12 de mayo de 2026, el Consejo de Ministros aprobó el proyecto de Ley de Gestión Pública e Integridad del Sistema Nacional de Salud, que deroga la Ley 15/1997.\footnote{%
    La Ley 15/1997, de art. único, habilitó nuevas formas de gestión, incluidas las de carácter privado, con el objetivo de introducir flexibilidad y eficiencia. Sobre esa base se desarrollaron las concesiones administrativas, como la del modelo Alzira, y otras fórmulas de colaboración público-privada. El proyecto de Ley de Gestión Pública e Integridad del Sistema Nacional de Salud, aprobado por el Consejo de Ministros el 12 de mayo de 2026, pretende derogarla y establecer la preferencia por la gestión pública directa.
} El proyecto establece la gestión pública directa como pilar del sistema y convierte la gestión indirecta en estrictamente excepcional: las administraciones deberán demostrar que la gestión directa es imposible y que la fórmula elegida es sostenible y eficiente. Exige una evaluación previa obligatoria, da preferencia a las entidades sin ánimo de lucro en caso de empate y prohíbe los contratos en los que una misma empresa construya la infraestructura y preste la asistencia sanitaria. Los contratos existentes se mantienen hasta su expiración. El proyecto llegó al Parlamento sin apoyos parlamentarios suficientes asegurados, y su aprobación es incierta. El debate enfrenta dos concepciones del sistema. Una sostiene que la titularidad y la financiación públicas son lo esencial, y que la forma de gestión debe decidirse por criterios de eficiencia. Otra sostiene que la gestión pública es inseparable del carácter público del sistema, porque la gestión privada introduce incentivos contrarios al interés de los pacientes. Es, en el ámbito sanitario, la misma tensión entre provisión pública y privada que se examinó en 1.1 con la distinción de Barr y Diamond entre mandato, provisión y financiación.

El sistema sanitario español presenta, en conjunto, un balance que conviene subrayar para cerrar el subapartado. Con un gasto por habitante inferior a la media de los países de su entorno, consigue uno de los mejores resultados de salud del mundo. España tiene una de las esperanzas de vida más altas del planeta, y algunos índices internacionales la han situado entre los sistemas más eficientes. Su fortaleza es la universalidad y la atención primaria, que durante décadas fue un modelo de referencia. Sus debilidades son las listas de espera, la presión sobre la atención primaria por la falta de profesionales y el envejecimiento de las plantillas, la desigualdad territorial y la creciente dualización entre quienes pueden pagar un seguro privado y quienes dependen exclusivamente del sistema público. Su sostenibilidad no depende del sistema de Seguridad Social, sino de la financiación autonómica y de la capacidad de incorporar la innovación tecnológica con criterios de eficiencia.

Las fuentes oficiales sobre la dependencia discrepan hasta el punto de que la magnitud del problema depende de a quién se pregunte. En el segundo trimestre de 2026, el Ministerio de Derechos Sociales cifraba la lista de espera en unas 143.000 personas. Unos meses antes, el Observatorio Estatal de la Dependencia la había cifrado en más de 258.000 para 2025, porque incluye también a quienes aún no han sido valorados, y estimaba en más de 32.000 las personas fallecidas ese año mientras esperaban. El documento describe el sistema con datos de 2021 y en una sola página, y no permite entender esa discrepancia ni el debate de fondo. Ese debate es por qué España llegó a la protección de la dependencia treinta años más tarde que sus vecinos, y por qué su sistema se ha desarrollado como un modelo "de bajo coste" basado en pagos a las familias en lugar de servicios. Lo que procede del documento va marcado [DOC], con las correcciones en el propio texto, y las reformas, en nota al pie.

\subsubsection{Economía de la sanidad pública}

#### Arrow y la singularidad de la asistencia sanitaria

El punto de partida de la economía de la salud moderna es el art. de Arrow (1963), "Uncertainty and the Welfare Economics of Medical Care", que se citó en 1.1.1 al examinar el riesgo moral. Arrow sostuvo que la asistencia sanitaria se aparta de los supuestos del mercado competitivo en tantos aspectos que las instituciones que la organizan no pueden entenderse con los instrumentos ordinarios del análisis económico. Identificó cuatro rasgos singulares. El primero es la incertidumbre sobre la incidencia de la enfermedad: nadie sabe cuándo enfermará ni cuánto le costará curarse, y los costes pueden ser catastróficos. El segundo es la incertidumbre sobre la eficacia del tratamiento, que el paciente no puede evaluar. El tercero es la asimetría de información entre el médico y el paciente: el médico sabe mucho más sobre la enfermedad y su tratamiento, de modo que el paciente no puede actuar como un consumidor informado. El cuarto es la relación de agencia: el paciente delega en el médico decisiones que no puede tomar por sí mismo, y confía en que actúe en su interés y no en el propio.

De estos rasgos se derivan varios fallos de mercado que conectan con la teoría de 1.1.1. La incertidumbre sobre la incidencia justifica el seguro, pero el mercado de seguros de salud sufre selección adversa, porque quienes esperan enfermar más buscan más cobertura, y selección de riesgos, porque las aseguradoras tienen incentivos para evitar a los enfermos crónicos y a los mayores. Es el mecanismo examinado en 1.1.1 con Rothschild y Stiglitz y, en el caso español, con la crisis de MUFACE de 2024-2025 (1.4.5). El seguro genera a su vez riesgo moral, porque el asegurado consume más asistencia cuando no paga su coste. Pauly (1968), en su respuesta a Arrow, formalizó este problema. La asimetría de información y la relación de agencia generan además el riesgo de demanda inducida por la oferta: el médico puede recomendar más tratamientos de los necesarios cuando su remuneración depende del número de actos que realiza.

La evidencia empírica más sólida sobre el riesgo moral en sanidad procede del experimento de seguros de salud de la RAND, realizado en Estados Unidos entre 1974 y 1982, en el que se asignaron aleatoriamente familias a seguros con distintos niveles de copago. Manning y otros (1987) estimaron una elasticidad-precio de la demanda de asistencia sanitaria de alrededor de -0,2. Es decir, un aumento del 10 por ciento del precio que paga el paciente reduce su consumo de asistencia en torno a un 2 por ciento. El experimento encontró además que los copagos reducían por igual el consumo de asistencia necesaria e innecesaria, y que no tenían efectos significativos sobre la salud de la población media, pero sí sobre la de los pobres y los enfermos crónicos. Es la base empírica del debate sobre los copagos que se examina en 3.4.3.

Existe, además, un argumento de equidad que no depende de ningún fallo de mercado. Tobin (1970) formuló el igualitarismo específico: aunque la sociedad acepte la desigualdad en la distribución de la mayoría de los bienes, puede considerar que ciertos bienes básicos, como la asistencia sanitaria, deben distribuirse según la necesidad y no según la capacidad de pago. Es una versión de la categoría de bien preferente de Musgrave (1.1.1) y del enfoque de las capacidades de Sen (1.1.2), aplicada a la salud, y es el fundamento del art. 43 de la Constitución.

#### Los modelos de sistema sanitario

A partir de estos fundamentos, los países han construido tres grandes modelos de sistema sanitario, que reproducen la tipología de 1.1.2. El modelo beveridgiano o de servicio nacional de salud, cuyo arquetipo es el NHS británico, se caracteriza por la cobertura universal por ciudadanía o residencia, la financiación con impuestos generales y la provisión mayoritariamente pública. El modelo bismarckiano o de seguro social de salud, cuyo arquetipo es el alemán, se caracteriza por la cobertura vinculada al empleo, la financiación con cotizaciones gestionadas por cajas de enfermedad y una provisión mixta, pública y privada. El modelo de mercado, cuyo ejemplo más cercano es Estados Unidos, combina el seguro privado vinculado al empleo con programas públicos para los mayores y los pobres. Enthoven (1978, 1993) propuso una vía intermedia, la competencia gestionada (managed competition), en la que aseguradoras públicas o privadas compiten por los asegurados bajo reglas que impiden la selección de riesgos. Inspiró las reformas de los Países Bajos y el funcionamiento del mutualismo administrativo español (1.4.5).

España pasó del modelo bismarckiano al beveridgiano en un proceso gradual que se relató en 1.2.2. El Seguro Obligatorio de Enfermedad de 1942 era bismarckiano: cubría a los trabajadores y sus familias y se financiaba con cotizaciones. La Ley General de Sanidad de 1986 inició la transición, al crear un servicio nacional de salud universal. La Ley de Presupuestos para 1989 empezó a sustituir las cotizaciones por impuestos. La financiación exclusivamente con impuestos se completó a finales de los noventa. La universalización plena de la cobertura se produjo en varias etapas y culminó con el Real Decreto-ley 7/2018 (3.4.2). El resultado es un sistema de modelo beveridgiano en su cobertura y financiación, con una provisión mayoritariamente pública pero con una participación relevante de la provisión privada concertada. Conserva además un rasgo de competencia gestionada en el mutualismo administrativo de los funcionarios.

#### Los determinantes del gasto sanitario

[DOC] El documento examina los determinantes de la evolución del gasto sanitario desde el lado de la demanda y de la oferta, y conviene reproducir su análisis antes de ampliarlo con la literatura. [DOC] Por el lado de la demanda, el número de personas que necesitan servicios sanitarios depende del envejecimiento de la población, de su estructura por sexo y de su estado de salud. El envejecimiento presiona al alza el gasto por dos mecanismos: la mayor tasa de dependencia, que reduce la proporción de contribuyentes respecto de los demandantes de servicios, y el aumento de la longevidad, que exige más gasto sanitario si no va acompañado de una mejora de la salud en edades avanzadas. [DOC] Por el lado de la oferta, el gasto depende de las normas que regulan el acceso a los servicios, de su disponibilidad y del progreso tecnológico. Este tiene un efecto ambiguo. Si reduce el coste unitario de tratamientos ya existentes, reduce el gasto (efecto sustitución). Si permite tratar a más personas o condiciones clínicas que antes no se trataban, lo aumenta (efecto expansión).

La literatura ha precisado la importancia relativa de cada factor, con conclusiones que corrigen la intuición común. El primer hallazgo se refiere al envejecimiento. Zweifel, Felder y Meiers (1999), en un art. titulado provocativamente "Ageing of Population and Health Care Expenditure: A Red Herring?", mostraron que el gasto sanitario no depende tanto de la edad como de la proximidad a la muerte: la mayor parte del gasto sanitario de una persona se concentra en sus últimos años de vida, sea cual sea la edad a la que muera. Si el aumento de la esperanza de vida retrasa la muerte, retrasa también ese gasto concentrado, y el efecto del envejecimiento sobre el gasto total es mucho menor de lo que sugieren las proyecciones basadas en perfiles de gasto por edad. Es la tesis de la compresión de la morbilidad formulada por Fries (1980): si la mejora de la salud acompaña al aumento de la longevidad, los años adicionales de vida son años sanos, y el gasto no aumenta en proporción. La tesis contraria, la de la expansión de la morbilidad, sostiene que los años adicionales son años de enfermedad crónica y dependencia, y aumentan el gasto. La evidencia es mixta y varía según las enfermedades. Por eso el Ageing Report proyecta el gasto sanitario con escenarios alternativos que incorporan distintos supuestos sobre la evolución de la salud.

El segundo hallazgo se refiere a la tecnología. Newhouse (1992), en un art. muy influyente, mostró que el envejecimiento, el aumento de la renta y la extensión del seguro solo explicaban una parte del crecimiento del gasto sanitario estadounidense desde 1940. La mayor parte del residuo correspondía al cambio tecnológico, es decir, a la introducción de nuevos tratamientos, en la línea del efecto expansión que el documento describe. Cutler y McClellan (2001) se preguntaron si ese gasto adicional merecía la pena y concluyeron que, en la mayoría de las enfermedades estudiadas, el valor de las mejoras de salud superaba ampliamente su coste. La tecnología encarece la sanidad, pero produce a cambio más salud. La cuestión no es si el gasto crece, sino si su crecimiento produce un valor suficiente, que es la pregunta de la evaluación económica de las tecnologías sanitarias.

El tercer hallazgo se refiere a la renta. Newhouse (1977) mostró que el gasto sanitario de los países crece más que proporcionalmente con su renta, lo que sugiere que la asistencia sanitaria es un bien de lujo a nivel agregado. Getzen (2000) precisó que la elasticidad-renta es inferior a uno a nivel individual, porque la asistencia se financia colectivamente, pero superior a uno a nivel de los países. A medida que una sociedad se enriquece, decide destinar una proporción creciente de su renta a la salud. El cuarto factor es la enfermedad de costes de Baumol examinada en 2.3.1: la sanidad es intensiva en trabajo cualificado con escaso crecimiento de la productividad, de modo que su coste relativo aumenta con el tiempo.

La conclusión de esta literatura es relevante para el análisis de la sostenibilidad del bloque IV. El envejecimiento, que el debate público presenta como la principal amenaza para el gasto sanitario, es en realidad un factor de segundo orden frente a la tecnología, la renta y el efecto Baumol. El crecimiento del gasto sanitario no es principalmente un problema demográfico, sino el resultado de decisiones colectivas sobre qué tecnologías incorporar y cuánto de la renta destinar a la salud.



\subsection{El Sistema para la Autonomía y Atención a la Dependencia}

[DOC] El documento define el Sistema para la Autonomía y Atención a la Dependencia (SAAD) como el conjunto de servicios y prestaciones económicas destinados a la promoción de la autonomía personal y a la atención y protección de las personas en situación de dependencia. Se estableció como una acción coordinada entre el Estado y las Comunidades Autónomas, con la participación de las entidades locales y de centros privados acreditados. [DOC] Añade que, para ello, se creó el Consejo Territorial de Servicios Sociales y del Sistema para la Autonomía y Atención a la Dependencia, para acordar el marco de cooperación interadministrativa, así como diversos órganos consultivos. [DOC, corregido] El documento adscribe ese Consejo al "Ministerio de Sanidad, Servicios Sociales e Igualdad", una denominación ministerial superada. Desde 2020, la competencia estatal sobre la dependencia corresponde al Ministerio de Derechos Sociales, Consumo y Agenda 2030, al que está adscrito también el IMSERSO (1.4.3).

La dependencia es el último gran riesgo social que el Estado del bienestar español ha asumido, y el que mejor ilustra dos rasgos del régimen mediterráneo examinados en 1.1.3: el familiarismo y la baja desfamiliarización. Hasta 2006, el cuidado de las personas mayores o con discapacidad que no podían valerse por sí mismas era responsabilidad casi exclusiva de las familias, y en particular de las mujeres, con una red pública de servicios sociales muy limitada y fragmentada territorialmente. La Ley 39/2006 convirtió por primera vez la atención a la dependencia en un derecho subjetivo universal. Su historia desde entonces, marcada por una crisis que llegó apenas dos años después de su aprobación, por los recortes de 2012, por la pandemia, que golpeó con especial dureza a las residencias, y por la reforma de 2026, es la de un derecho proclamado con ambición y desarrollado con recursos insuficientes.

El subapartado se organiza en cuatro partes: la economía de los cuidados de larga duración (3.5.1), el marco jurídico, los sujetos y la desagregación vertical (3.5.2), las prestaciones (3.5.3) y la financiación, el desarrollo del sistema y la reforma de 2026 (3.5.4).


La Ley 39/2006, de 14 de diciembre, de Promoción de la Autonomía Personal y Atención a las personas en situación de dependencia, creó el SAAD e incorporó la dependencia al sistema público de protección social en España. [DOC] El art. 2 de la ley define la dependencia como el estado de carácter permanente en que se encuentran las personas que, por razones derivadas de la edad, la enfermedad o la discapacidad, y ligadas a la falta o pérdida de autonomía física, mental, intelectual o sensorial, precisan de la atención de otra u otras personas o de ayudas importantes para realizar las actividades básicas de la vida diaria. En el caso de las personas con discapacidad intelectual o enfermedad mental, también la precisan quienes necesitan otros apoyos para su autonomía personal.

La ley clasifica la dependencia en tres grados, según la intensidad de la ayuda necesaria. El grado I o dependencia moderada corresponde a quien necesita ayuda al menos una vez al día. El grado II o dependencia severa, a quien la necesita dos o tres veces al día, sin necesidad de apoyo permanente. El grado III o gran dependencia, a quien la necesita varias veces al día y precisa el apoyo indispensable y continuo de otra persona. El grado se determina mediante un baremo de valoración aplicado por los órganos de las Comunidades Autónomas. La valoración es el primer paso de un procedimiento largo, que incluye después la elaboración de un programa individual de atención, en el que se determinan los servicios o prestaciones concretos. Como se verá en 3.5.4, la duración de ese procedimiento es uno de los principales problemas del sistema.

#### Los sujetos

[DOC, corregido] El documento enumera como requisitos generales para ser beneficiario "ser español", tener cualquier edad (con peculiaridades para los menores de tres años), encontrarse en situación de dependencia en alguno de los grados establecidos y residir en territorio español y haberlo hecho durante los cinco años anteriores a la solicitud. [DOC] Añade que existen otros requisitos para las personas que carecen de la nacionalidad española y para los emigrantes españoles retornados. La formulación "ser español" como requisito general es imprecisa. La ley reconoce el derecho a los españoles y, en los términos de la legislación de extranjería y de los tratados internacionales, a los extranjeros residentes, de modo que la nacionalidad no es, en rigor, un requisito excluyente. El requisito de residencia exige haber residido en España durante cinco años, de los que dos deben ser inmediatamente anteriores a la solicitud, con reglas especiales para los menores. Es un requisito que expresa la lógica beveridgiana del sistema, en el sentido de 1.1.2 y 1.3.2: el derecho se vincula a la residencia, no a la cotización, porque el SAAD se financia con impuestos y no forma parte del sistema contributivo de la Seguridad Social. El periodo mínimo de residencia es, a la vez, una forma de limitar el acceso de los recién llegados y de evitar el llamado "turismo de prestaciones".

#### La desagregación vertical: un sistema compartido

[DOC] Tu nota a este epígrafe pedía responder a cómo se distribuyen las competencias y completarlo con el gasto por subsector. El reparto competencial de la dependencia es el más complejo de los sistemas de protección examinados, porque combina tres títulos competenciales. Las Comunidades Autónomas tienen competencia exclusiva en materia de asistencia social (art. 148.1.20ª de la Constitución), que incluye los servicios sociales y, en principio, la atención a la dependencia. El Estado no tiene sobre la dependencia la competencia del art. 149.1.17ª, porque el SAAD no forma parte del sistema de Seguridad Social. Por eso la Ley 39/2006 se apoyó en el art. 149.1.1ª, que atribuye al Estado la regulación de las condiciones básicas que garanticen la igualdad de todos los españoles en el ejercicio de sus derechos. Con ese título, el Estado define el derecho, los grados, las prestaciones y un nivel mínimo de protección, mientras que las Comunidades Autónomas gestionan el sistema y pueden mejorar la protección. Esta construcción jurídica explica la naturaleza compartida del sistema y su principal fuente de tensión: el Estado define un derecho cuya gestión y buena parte de cuya financiación corresponden a las comunidades.

La ley estableció tres niveles de protección. El nivel mínimo lo financia el Estado con una cantidad fija por beneficiario según su grado de dependencia. El nivel acordado lo financian conjuntamente el Estado y cada comunidad mediante convenios bilaterales, y se distribuye según criterios pactados en el Consejo Territorial. El nivel adicional lo pueden establecer libremente las comunidades con sus propios recursos. La ley preveía, en su memoria económica, que el Estado y las comunidades financiaran el sistema a partes iguales. Esa previsión nunca se cumplió. [DOC] El documento señala que en 2021 el gasto público en dependencia ascendió al 0,8 por ciento del PIB, del que las Comunidades Autónomas aportaron el 79 por ciento y la Administración General del Estado el 21 por ciento restante. El nivel acordado se suspendió entre 2012 y 2020 en el marco de los recortes de la crisis, y se recuperó con el plan de choque de 2021.\footnote{%
    Entre 2012 y 2020, la financiación estatal se limitó al nivel mínimo, y el nivel acordado quedó suspendido por las sucesivas leyes de presupuestos. El plan de choque aprobado por el Consejo Territorial en 2021, en el contexto del Plan de Recuperación, Transformación y Resiliencia, recuperó el nivel acordado, aumentó las cuantías del nivel mínimo y comprometió incrementos anuales de la financiación estatal, con el objetivo de reducir las listas de espera y mejorar la calidad de la atención.
} Las comunidades soportaron así la mayor parte del coste de un derecho definido por el Estado. Es un ejemplo de lo que la literatura del federalismo fiscal denomina mandato no financiado (unfunded mandate), y es una de las principales causas del desarrollo desigual del sistema entre comunidades.

La tensión competencial se resolvió en parte con la reforma de 2026, que obliga por ley al Estado a asumir exactamente la mitad del coste del SAAD (3.5.4). [DOC] Como se estableció en 2.3.1, el gasto en dependencia se imputa en su mayor parte al subsector de las Comunidades Autónomas, con una participación menor de la Administración Central a través del IMSERSO. [DOC] El documento señala que el desarrollo del sistema ha generado un desarrollo desigual entre Comunidades Autónomas, y que el grado de cobertura varía significativamente entre ellas. Es la manifestación de los costes de la descentralización examinados en 3.4.2: la desigualdad territorial en el acceso a un derecho que la ley define como universal.

\subsubsection{Prestaciones}

[DOC] El documento señala que las prestaciones del SAAD son de dos tipos: servicios y prestaciones económicas. Los servicios se prestan a través de la oferta pública de la red de servicios sociales de las Comunidades Autónomas, mediante centros públicos o privados concertados, y comprenden, entre otros, la teleasistencia, la ayuda a domicilio, los centros de día y de noche y la atención residencial. A ellos se añaden los servicios de prevención de la dependencia y de promoción de la autonomía personal. [DOC] Las prestaciones económicas son de tres tipos: la vinculada al servicio, cuando no es posible el acceso a un servicio público o concertado por el grado de dependencia o la capacidad económica; la de cuidados en el entorno familiar y apoyo a cuidadores no profesionales, cuando el beneficiario es atendido por su cónyuge o un pariente; y la de asistencia personal, dirigida a la contratación de un asistente personal.

La ley estableció una jerarquía entre estas prestaciones: los servicios son la prestación prioritaria; la prestación vinculada al servicio, un sustituto cuando no hay servicio público disponible; y la prestación por cuidados en el entorno familiar, una excepción, reservada a los casos en que el beneficiario ya estuviera siendo atendido en su domicilio y se dieran condiciones adecuadas de convivencia y habitabilidad. La jerarquía respondía a un objetivo de política social explícito en la exposición de motivos de la ley: la desfamiliarización de los cuidados y la creación de empleo profesional en el sector.

En la práctica, la excepción se convirtió en la regla. Durante los primeros años del sistema, la prestación por cuidados en el entorno familiar fue la más frecuente, con mucha diferencia. Según el Observatorio Estatal de la Dependencia, todavía en 2025 representaba el 44,6 por ciento de las prestaciones, con una cuantía media de solo 264,55 euros mensuales. Las razones fueron varias. Era la opción más barata para las Comunidades Autónomas, que soportaban la mayor parte del coste. Era la preferida por muchas familias, que así recibían un ingreso y mantenían al dependiente en su domicilio. Y la oferta de servicios públicos era insuficiente para atender la demanda. El resultado fue lo que el Observatorio califica como un sistema de bajo coste (low cost): una protección basada en pagos modestos a las familias, que en términos de 3.5.1 refamiliarizó el cuidado en lugar de desfamiliarizarlo. Contradijo así el objetivo original de la ley.

Los datos más recientes del Ministerio muestran un cambio de tendencia. En el segundo trimestre de 2026, el 57 por ciento de los servicios profesionales se prestaba en entornos comunitarios, la ayuda a domicilio crecía un 13,6 por ciento anual y la teleasistencia, un 21,8 por ciento. El número medio de servicios por persona había aumentado de 1,27 en 2020 a 1,45. El cambio responde a la recuperación de la financiación estatal desde 2021 y a una reorientación deliberada del sistema hacia los servicios de proximidad, en línea con la estrategia de desinstitucionalización que se examina a continuación.

La prestación de asistencia personal, pensada originalmente para personas con discapacidad que quieren llevar una vida independiente, ha tenido un desarrollo muy escaso y concentrado en unas pocas comunidades. La reforma de 2026 la configura como un derecho garantizado y reconoce la figura profesional del asistente personal. Es un avance en la línea del modelo de vida independiente defendido por el movimiento de las personas con discapacidad y por la Convención de las Naciones Unidas sobre los Derechos de las Personas con Discapacidad (2006), que España ratificó en 2007.

La atención residencial merece una mención específica por lo ocurrido durante la pandemia. En la primavera de 2020, las residencias de mayores fueron uno de los principales focos de mortalidad por COVID-19 en España y en otros países europeos, con decenas de miles de fallecidos en ellas. La cifra exacta es objeto de controversia entre distintas fuentes. La crisis puso de manifiesto las debilidades del modelo residencial: centros de gran tamaño, escasa integración con el sistema sanitario, plantillas insuficientes y precarias, y un peso creciente de la gestión privada con ánimo de lucro. En respuesta, el Consejo Territorial aprobó en 2022 nuevos criterios comunes de acreditación y calidad de los centros, que priorizan las unidades de convivencia de menor tamaño y la atención centrada en la persona. El Gobierno impulsó además una estrategia de desinstitucionalización, que favorece la permanencia de las personas en su entorno con apoyos domiciliarios y comunitarios.

\subsubsection{La financiación, el desarrollo del sistema y la reforma de 2026}
[DOC] El documento señala que el SAAD recibe recursos de tres fuentes: los créditos que fijan anualmente los Presupuestos Generales del Estado, los presupuestos de las Comunidades Autónomas y las aportaciones de los usuarios, que dependen de su capacidad económica y del tipo de servicio. Las aportaciones de los usuarios, o copagos, son significativas en los servicios residenciales y tienen en cuenta tanto la renta como, en algunos casos, el patrimonio. Su regulación varía entre comunidades, lo que añade una fuente de desigualdad territorial. Según el Observatorio Estatal de la Dependencia, el gasto total en el sistema alcanzó en 2025 unos 13.500 millones de euros, el 0,8 por ciento del PIB, lo que el Observatorio estima en torno a la mitad de la media de la Unión Europea en cuidados de larga duración. La comparación exacta debe verificarse con el Ageing Report.

El Gobierno aprobó en 2026 un aumento de 6.200 millones de euros adicionales de financiación estatal para 2026 y 2027. Con él, la aportación estatal alcanzará en 2027 los 7.239 millones de euros, el doble que en 2025, y el Estado pasará a financiar el 50 por ciento del coste total del sistema. Es la primera vez desde la creación del SAAD que se cumpliría el reparto a partes iguales previsto en la memoria económica de la ley de 2006.

#### El desarrollo del sistema: de la crisis al plan de choque

[DOC] El documento señala que el desarrollo del SAAD desde su implantación ha estado fuertemente condicionado por la crisis económica y el posterior proceso de consolidación fiscal, con referencia a González Ortega (2013) y Cabrero y otros (2022). Ello generó numerosos cambios normativos y un desarrollo desigual entre comunidades. [DOC] En 2021 había algo más de 1.200.000 personas beneficiarias con prestación, frente a las poco más de 220.000 de finales de 2008, y algo más de 200.000 personas seguían pendientes de recibir prestación pese a tener reconocido un grado de dependencia.

La historia del SAAD tiene tres etapas, que ilustran la economía política de 1.2.1. La primera, entre 2007 y 2011, fue de implantación, con un calendario gradual que incorporó primero a los grandes dependientes y después a los grados inferiores. El sistema creció rápidamente, pero con predominio de las prestaciones por cuidados familiares. La segunda, entre 2012 y 2020, fue de recorte y estancamiento. El Real Decreto-ley 20/2012 suspendió el nivel acordado, redujo las cuantías del nivel mínimo y de la prestación por cuidados familiares, suspendió la cotización a la Seguridad Social de los cuidadores no profesionales y retrasó la incorporación de los dependientes moderados.\footnote{%
    El Real Decreto-ley 20/2012 redujo las cuantías del nivel mínimo, rebajó en torno a un 15 por ciento la prestación por cuidados en el entorno familiar, suprimió la cotización a la Seguridad Social a cargo del Estado de los cuidadores no profesionales, que pasó a ser voluntaria y a cargo del propio cuidador, e introdujo un plazo suspensivo para el cobro de las prestaciones económicas. El objetivo era contener el gasto en el marco del ajuste fiscal. La cotización de los cuidadores no profesionales a cargo del Estado se recuperó en 2019, mediante el Real Decreto-ley 6/2019.
} Las listas de espera se dispararon, y el sistema fue el ámbito de la protección social española más afectado por la austeridad. Es un ejemplo de la tesis de la paradoja de la redistribución de Korpi y Palme (1.1.2): un derecho reciente, con una coalición de apoyo débil, fue mucho más vulnerable a los recortes que las pensiones, con su poderosa coalición de pensionistas. La tercera etapa, desde 2021, es de recuperación. El plan de choque para la dependencia de 2021, acordado con las comunidades, recuperó el nivel acordado, aumentó la financiación estatal y fijó objetivos de reducción de las listas de espera y de mejora de los servicios.

Los resultados de esta tercera etapa son visibles en los datos del Ministerio. En el segundo trimestre de 2026, el número de personas atendidas alcanzó el máximo histórico de 1.707.328, un 11 por ciento más que un año antes y un 57,5 por ciento más que en 2021. La lista de espera se redujo a 142.887 personas, un 21 por ciento menos que un año antes y un 54 por ciento menos que en 2021. El tiempo medio de tramitación, desde la solicitud hasta la resolución, bajó a 314 días, frente al máximo de 421 días de 2021.

El Observatorio Estatal de la Dependencia, que elabora la Asociación Estatal de Directoras y Gerentes en Servicios Sociales, ofrece en su dictamen de marzo de 2026 una lectura más crítica de los datos de 2025. Cifra en unas 258.000 las personas en espera, al incluir también a quienes aún no han sido valorados, y en 341 días el tiempo medio de resolución. Estima que unas 32.700 personas fallecieron en 2025 mientras esperaban ser valoradas o recibir la prestación a la que tenían derecho. Solo seis comunidades superan el aprobado en su valoración, con Castilla-La Mancha y Castilla y León en cabeza, y Extremadura y Ceuta y Melilla en último lugar. La diferencia entre las cifras del Ministerio y las del Observatorio no se debe a un error de ninguna de las dos fuentes, sino a una diferencia de definición. El Ministerio cuenta como lista de espera a las personas con grado reconocido que esperan su prestación. El Observatorio incluye también a las que esperan ser valoradas. Es el llamado "limbo de la dependencia". Es un recordatorio de que, como la medición de la pobreza o del déficit (1.4.1), la medición de las listas de espera no es técnicamente neutral: depende de dónde se sitúe el punto de partida del derecho. Y la elección de la definición tiene consecuencias políticas.

#### La reforma de 2026

En septiembre de 2026, el Congreso aprobó definitivamente la reforma de las leyes de dependencia y de discapacidad, que se había aprobado en su primera lectura en julio con la abstención del Partido Popular.\footnote{%
    La reforma modifica la Ley 39/2006 y el texto refundido de la Ley General de derechos de las personas con discapacidad y de su inclusión social. Su tramitación siguió a la reforma del art. 49 de la Constitución de febrero de 2024 (1.3.1). Su número de ley y su fecha de publicación en el Boletín Oficial del Estado están pendientes de verificar.
} Sus principales contenidos son cinco. El primero es la obligación legal de que el Estado financie exactamente la mitad del coste del SAAD, que hasta ahora era solo una previsión nunca cumplida. Resuelve así, al menos formalmente, el problema del mandato no financiado examinado en 3.5.2. El segundo es la eliminación de las incompatibilidades entre prestaciones, que impedían combinar, por ejemplo, una prestación económica con un servicio residencial o de teleasistencia. El tercero es el reconocimiento automático de la discapacidad a quienes tienen reconocida la dependencia, lo que elimina una doble valoración y reduce la carga administrativa en el sentido de Herd y Moynihan (3.1.5). El cuarto es el reconocimiento de la asistencia personal como derecho y de la figura profesional del asistente personal. El quinto es la consideración de los servicios sociales como servicios esenciales, con garantía de atención en situaciones de emergencia. La reforma se completa con el aumento de la financiación estatal de 6.200 millones de euros ya examinado.

La reforma responde a los principales problemas identificados en este subapartado: la infrafinanciación estatal, la rigidez del catálogo de prestaciones, la carga administrativa y el escaso desarrollo de la asistencia personal. Deja sin resolver, sin embargo, las cuestiones de fondo que la teoría de 3.5.1 plantea. La primera es la de la sostenibilidad a largo plazo de un sistema que tendrá que atender a la generación del baby boom a partir de la década de 2040, con una financiación que depende de las decisiones presupuestarias anuales del Estado y de las comunidades, sin una fuente de ingresos afectada como el seguro social alemán o japonés. La segunda es la de la escasez de cuidadores profesionales, en un sector con salarios bajos y condiciones de trabajo duras, que depende en gran medida de la inmigración. La tercera es la del equilibrio entre servicios y prestaciones monetarias, que determina si el sistema contribuye a la desfamiliarización de los cuidados o la refuerza. Algunos expertos han propuesto crear un seguro social de dependencia financiado con una cotización específica, como el alemán, para dotar al sistema de una fuente de financiación estable y visible. Sus críticos advierten que una nueva cotización aumentaría la cuña fiscal sobre el trabajo, en un país donde ya es elevada (2.1.2), y preferirían una financiación con impuestos generales. El debate conecta directamente con la elección entre cotizaciones e impuestos examinada en el bloque II.

\subsubsection{Economía de los cuidados de larga duración}

#### El rompecabezas del seguro de dependencia

La necesidad de cuidados de larga duración es un riesgo con las características que, según la teoría de 1.1.1, justifican un seguro. Es un riesgo incierto, porque nadie sabe si llegará a ser dependiente. Tiene un coste potencialmente catastrófico, porque una estancia prolongada en una residencia puede agotar el patrimonio de toda una vida. Y está concentrado en la vejez, cuando la capacidad de generar ingresos ha desaparecido. Un mercado de seguros eficiente debería ofrecer seguros de dependencia, y los individuos racionales deberían contratarlos. Sin embargo, en casi todos los países el mercado privado de seguros de dependencia es muy pequeño. Es el llamado rompecabezas del seguro de cuidados de larga duración (long-term care insurance puzzle), análogo al de las rentas vitalicias examinado en 1.1.1.

La literatura ha ofrecido varias explicaciones. Pauly (1990), en "The Rational Nonpurchase of Long-Term-Care Insurance", sostuvo que no contratar un seguro de dependencia puede ser racional por dos razones. La primera es que muchas personas prefieren ser cuidadas por su familia, y no por profesionales, lo que reduce el valor de un seguro que paga cuidados profesionales. La segunda es que los hijos, que heredarán el patrimonio de sus padres, tienen un incentivo para cuidarlos, y ese incentivo se debilitaría si los padres contrataran un seguro. Bernheim, Shleifer y Summers (1985) habían formulado la teoría de la herencia estratégica: los padres utilizan la promesa de herencia para obtener atención de sus hijos, y un seguro de dependencia debilitaría ese mecanismo. Brown y Finkelstein (2008), en su estudio sobre el mercado estadounidense, mostraron que buena parte de la explicación está en la interacción con el seguro público. El programa Medicaid paga los cuidados de larga duración de quienes han agotado su patrimonio, y por eso actúa como un impuesto implícito muy elevado sobre el seguro privado: para la mayoría de las personas de renta media y baja, un seguro privado pagaría cuidados que Medicaid habría pagado igualmente. Es el dilema del samaritano de Buchanan (1.1.1) aplicado a la dependencia: la existencia de una red pública de último recurso desincentiva la previsión privada.

La conclusión de esta literatura, recogida por Barr (2010) en su análisis del seguro de dependencia, es que los cuidados de larga duración son un caso claro para el seguro social obligatorio. Lo son por las mismas razones que las pensiones: la selección adversa, la miopía, el dilema del samaritano y la existencia de riesgos agregados, como el aumento de la longevidad y de la prevalencia de la demencia, que el mercado privado no puede diversificar.

#### Los modelos de protección de la dependencia

Los países que han asumido la protección de la dependencia lo han hecho siguiendo modelos que reproducen la tipología de 1.1.2 y 1.1.3. Los países nórdicos la integraron desde los años sesenta y setenta en sus servicios sociales universales, financiados con impuestos y gestionados por los municipios, con una fuerte orientación a los servicios profesionales. Alemania creó en 1995 un seguro social de dependencia (Pflegeversicherung) de tipo bismarckiano, financiado con cotizaciones y con prestaciones de cuantía fija según el grado de dependencia. Permite elegir entre servicios profesionales y una prestación monetaria de menor cuantía para cuidados familiares. Japón creó en 2000 un seguro social de dependencia financiado a partes iguales con cotizaciones e impuestos, que se ha convertido en una referencia internacional por su éxito en la expansión de los servicios domiciliarios. Los países mediterráneos, con Italia como caso más característico, desarrollaron tardíamente prestaciones monetarias de escasa cuantía, como la italiana indennità di accompagnamento. Dejaron así la mayor parte del cuidado a las familias y, cada vez más, a trabajadoras domésticas inmigrantes.

Este último fenómeno ha sido analizado por Bettio, Simonazzi y Villa (2006) con el concepto de "fuga de cuidados" (care drain). En los países mediterráneos, la combinación de prestaciones monetarias sin control sobre su uso, de mercados de trabajo doméstico informales y de flujos migratorios procedentes de América Latina y Europa del Este produjo un modelo de "inmigrante en la familia". Las familias utilizan las prestaciones públicas para contratar, a menudo de forma irregular, a mujeres inmigrantes que cuidan a sus mayores en el propio domicilio. El modelo resolvió el problema de los cuidados a bajo coste para las familias y para el Estado, pero a costa de la precariedad laboral de las cuidadoras y de un "drenaje de cuidados" en sus países de origen, donde sus propias familias quedaban sin atención. España reproduce en buena medida este modelo, y la inclusión de las empleadas del hogar en la protección por desempleo en 2022 (3.3.2) es una de las respuestas a su precariedad.

#### Prestaciones monetarias frente a servicios

La elección entre prestaciones monetarias (cash for care) y servicios es el debate central de la política de dependencia, y Da Roit y Le Bihan (2010) lo analizaron comparando seis países europeos. Las prestaciones monetarias tienen tres ventajas: respetan la libertad de elección de las familias, permiten que el cuidado lo presten personas de confianza del dependiente y tienen un coste presupuestario mucho menor que los servicios profesionales. Sus inconvenientes son también tres. Tienden a refamiliarizar el cuidado, reforzando la dependencia de las mujeres del hogar y reduciendo su participación laboral, lo que contradice el objetivo de desfamiliarización de Esping-Andersen (1.1.3). No garantizan la calidad de los cuidados, porque la Administración no controla su uso. Y no generan empleo profesional de calidad, a diferencia de los servicios. Los servicios profesionales tienen las ventajas e inconvenientes opuestos: garantizan la calidad y generan empleo formal, pero son mucho más caros y limitan la elección. Los cuidados de larga duración están además sujetos a la enfermedad de costes de Baumol examinada en 2.3.1, en su versión más pura, porque son intensivos en trabajo y apenas admiten ganancias de productividad. Su coste relativo aumentará con el tiempo, lo que presiona hacia las soluciones monetarias de menor coste.

#### La dimensión demográfica

La presión sobre los sistemas de dependencia crecerá intensamente en las próximas décadas por el envejecimiento de la población, y en particular por el aumento del número de personas de más de 80 años, el grupo con mayor prevalencia de dependencia. En España, la llegada de las generaciones del baby boom a esas edades a partir de la década de 2040 multiplicará la demanda. A la vez, la reducción del tamaño de las familias, la mayor participación laboral de las mujeres y la dispersión geográfica de los hijos reducirán la disponibilidad de cuidadores informales, que hoy prestan la mayor parte de los cuidados. Es la manifestación en la dependencia del mismo cambio demográfico que afecta a las pensiones, pero con un agravante: mientras las pensiones son transferencias monetarias, los cuidados exigen trabajo que alguien tiene que prestar, y la escasez de cuidadores no se resuelve solo con más recursos financieros. Las proyecciones del Ageing Report sobre el gasto en cuidados de larga duración se examinan en el bloque IV.




\clearpage
\section{Sistema Público de Pensiones: ¿es sostenible?}
\puntosclave{
    
}


Los tres bloques anteriores han descrito un sistema que protege frente a los riesgos de la vida con una lógica mixta de seguro y redistribución (bloque I). Lo financia con una combinación cada vez más opaca de cotizaciones, impuestos y deuda (bloque II). Y ofrece un conjunto de prestaciones cuyo núcleo, las pensiones, absorbe más de la cuarta parte del gasto público total (bloque III). La pregunta que el título del tema reserva para el final es si ese sistema puede mantenerse en el tiempo. Puede formularse así: si las promesas que hoy hace a los trabajadores podrán cumplirse cuando se jubilen, y a qué coste para quienes tendrán que financiarlas.

[DOC] El documento resume en sus puntos clave la respuesta que daba a esa pregunta en su versión original. El aumento de la longevidad y el progresivo envejecimiento de la población imponen una presión muy importante sobre el gasto en pensiones durante las próximas décadas. Las reformas de 2011 y, sobre todo, de 2013 constituyeron un avance sustancial en favor de la sostenibilidad del sistema. Sin embargo, la reforma de 2013 implicaba una caída muy pronunciada de la pensión media en relación con el salario medio, lo que minoraba la suficiencia de las prestaciones y llevó a su posterior derogación. [DOC] Añade que aún era pronto para determinar el alcance de las reformas de 2021 y 2023, y que la AIReF había proyectado que supondrían un incremento del déficit de 1,1 puntos de PIB en 2050 y de 1 punto en 2070 respecto de un escenario en el que el índice de revalorización y el factor de sostenibilidad de 2013 hubieran permanecido vigentes. Ese resumen sigue siendo válido en su estructura, pero está superado en sus datos. Desde entonces se han publicado el Ageing Report de 2024, dos evaluaciones de la regla de gasto de la AIReF, la de 2025 y la de 2026, y nuevas proyecciones demográficas y del Ministerio.

El bloque se organiza en cinco apartados. El primero (4.1) construye el concepto de sostenibilidad y examina cómo se mide. El segundo (4.2) realiza el diagnóstico, mediante la descomposición del gasto en pensiones en sus factores determinantes. El tercero (4.3) examina las reformas como respuestas a ese diagnóstico y evalúa sus efectos. El cuarto (4.4) analiza la economía política de las pensiones y las alternativas de reforma que se discuten. El quinto (4.5) sitúa el sistema español en la comparación internacional. La tesis que recorre el bloque es la siguiente. La sostenibilidad del sistema español no es un problema técnico con una solución única, sino una elección colectiva entre objetivos que están en tensión: la suficiencia de las pensiones, el coste para los trabajadores y la equidad entre generaciones. Desde 2021, esa elección se ha resuelto a favor de la suficiencia y de la financiación mediante ingresos. Si esa opción es sostenible depende de supuestos demográficos y económicos inciertos, y sobre todo de que el consenso político que la sostiene se mantenga.



\subsection{Sostenibilidad: ¿qué es, por qué importa y cómo se mide?}

[DOC] El documento parte de una observación que justifica la importancia de la cuestión: como las pensiones son la principal fuente de ingresos de la población mayor de 65 años, garantizar su sostenibilidad y su suficiencia es crucial. [DOC] Define la sostenibilidad del sistema de pensiones como el equilibrio entre los ingresos y las prestaciones del sistema a lo largo del ciclo económico y en el largo plazo. Define la suficiencia como la capacidad del sistema de proporcionar unos ingresos adecuados para prevenir la pobreza y permitir una vida digna durante la vejez.

Ambas definiciones son correctas como punto de partida, pero dejan abiertas tres cuestiones que este apartado aborda. La primera es qué tipo de equilibrio se exige: contable, fiscal, económico o político (4.1.1). La segunda es cómo se relaciona la sostenibilidad con la suficiencia y con la equidad entre generaciones, que forman con ella un dilema del que no se puede escapar (4.1.2). La tercera es cómo se mide, porque los distintos indicadores ofrecen diagnósticos muy distintos (4.1.3). Una última parte examina por qué importa la sostenibilidad y quién debe garantizarla (4.1.4).

\subsubsection{Sostenibilidad: ¿qué hay que tener en cuenta?}

\paragraph{Sostenibilidad financiera: el equilibrio del sistema}

El concepto más inmediato, y el que utiliza el documento, es la sostenibilidad financiera del sistema de pensiones considerado como una entidad separada. Un sistema es financieramente sostenible si sus ingresos propios bastan para cubrir sus gastos, sin necesidad de recursos externos, a lo largo del ciclo y a largo plazo. En un sistema de reparto puro, esto exige que las cotizaciones de los trabajadores en activo cubran las pensiones de los jubilados en cada momento, o que los déficits de unos años se compensen con los superávits de otros, por ejemplo a través de un fondo de reserva.

Este concepto es útil, pero tiene un problema grave que se examinó en el bloque II: depende de una frontera contable entre el sistema y el Estado que el legislador puede mover a voluntad. Como se estableció en 1.4.1 y 2.2.2, cuando el Estado aumenta sus transferencias a la Seguridad Social para financiar gastos que califica de "impropios", el déficit del sistema se reduce en la misma cuantía en que aumenta el del Estado, sin que cambie nada en la economía. Cuando el Estado presta a la Seguridad Social en lugar de transferirle recursos, el sistema acumula una deuda que el Tribunal de Cuentas considera una ficción, porque nadie espera que se devuelva (2.2.3). La sostenibilidad financiera del sistema en sentido estricto es, por tanto, un concepto convencional: su medida depende de decisiones contables, y un sistema puede aparecer como sostenible o insostenible según cómo se definan sus ingresos propios. Por eso el diagnóstico debe completarse con conceptos más amplios.

\paragraph{Sostenibilidad fiscal: las pensiones en la restricción presupuestaria del Estado}

El segundo concepto, que es el que utilizan la Comisión Europea, el FMI y la AIReF, es la sostenibilidad fiscal. En lugar de preguntar si el sistema de pensiones se autofinancia, pregunta si el conjunto de las Administraciones Públicas puede hacer frente a sus compromisos futuros, incluido el gasto en pensiones, sin que su deuda crezca de forma explosiva. El fundamento teórico es la restricción presupuestaria intertemporal del sector público: el valor actual de los gastos futuros no puede superar el valor actual de los ingresos futuros más los activos netos actuales. En otras palabras, la deuda actual debe poder pagarse con superávits primarios futuros. Blanchard (1990), en un trabajo para la OCDE, propuso medir la sostenibilidad mediante indicadores de brecha fiscal (sustainability gap): el ajuste permanente del saldo primario necesario para cumplir la restricción intertemporal, dados los gastos proyectados.

La Comisión Europea utiliza dos indicadores de este tipo en sus informes de sostenibilidad fiscal. El indicador S1 mide el ajuste necesario para reducir la deuda al 60 por ciento del PIB en un horizonte determinado. El indicador S2 mide el ajuste necesario para estabilizar la deuda en un horizonte infinito. Una parte significativa del S2 de España procede del coste del envejecimiento, es decir, del aumento proyectado del gasto en pensiones, sanidad y cuidados de larga duración. En las últimas evaluaciones de la Comisión, España ha figurado entre los países con riesgo alto de sostenibilidad fiscal a medio plazo, por la combinación de una deuda pública elevada y del aumento proyectado del gasto ligado al envejecimiento. La clasificación exacta en el último Debt Sustainability Monitor debe verificarse. Como se vio en 1.3.4, el nuevo marco fiscal europeo integra esta evaluación en el plan fiscal-estructural de medio plazo, y la Recomendación del Consejo que aprobó el plan español en enero de 2025 tuvo en cuenta expresamente las medidas de ingresos de la reforma de pensiones.

El concepto de sostenibilidad fiscal resuelve el problema del perímetro: es indiferente a si las pensiones se financian con cotizaciones, transferencias o préstamos, porque lo que importa es el saldo del conjunto de las Administraciones Públicas. Por eso es el concepto que utiliza la regla de gasto de la reforma de 2023, que no mide el saldo del sistema, sino el gasto bruto en pensiones en relación con el PIB y el efecto de las medidas de ingresos (4.3.3). Tiene, sin embargo, su propia limitación: trata las pensiones como un gasto más del Estado. Así ignora su naturaleza de derechos adquiridos mediante cotizaciones, que les da una rigidez política mucho mayor que la de otros gastos.

\paragraph{Sostenibilidad económica: la parte de la producción que se destina a los jubilados}

El tercer concepto es el más fundamental, y deriva del principio de que "la producción es lo central" examinado en 2.1. Con independencia de cómo se financien las pensiones, su sostenibilidad económica depende de la parte de la producción de cada momento que la sociedad destina a los jubilados, y de si esa parte es compatible con el bienestar de los trabajadores y con el crecimiento económico. Como mostró Barr (2002), ningún mecanismo financiero puede alterar el hecho de que los bienes y servicios que consumen los jubilados los producen los trabajadores del mismo momento. Si la proporción de jubilados sobre trabajadores aumenta, la parte de la producción destinada a los jubilados tiene que aumentar, o la producción por jubilado tiene que disminuir.

Desde esta perspectiva, la sostenibilidad económica depende de tres magnitudes: la relación entre jubilados y trabajadores, la productividad de los trabajadores y la generosidad relativa de las pensiones respecto de los salarios. Son exactamente los factores de la descomposición del gasto en pensiones que el documento presenta y que se examina en 4.2. Un sistema es económicamente sostenible si la parte de la producción que destina a los jubilados no crece hasta un nivel que reduzca excesivamente el consumo de los trabajadores, desincentive su oferta de trabajo o reduzca la inversión y el crecimiento. Ese nivel no es un umbral técnico, sino una decisión colectiva. Un país puede destinar el 15 o el 18 por ciento de su PIB a pensiones si su sociedad lo acepta y lo financia sin efectos negativos excesivos sobre el crecimiento, como ocurre en algunos países europeos. La cuestión no es si hay un límite absoluto, sino cuál es el coste de acercarse a él.

\paragraph{Sostenibilidad social y política: la legitimidad del sistema}
El cuarto concepto, cada vez más presente en el debate, es la sostenibilidad social o política. Un sistema de pensiones es sostenible solo si mantiene el apoyo de quienes lo financian y de quienes se benefician de él. Si las pensiones pierden suficiencia, los pensionistas presionan para mejorarlas. Si las cotizaciones son excesivas o los jóvenes perciben que recibirán menos de lo que aportan, los trabajadores pierden la confianza en el sistema. En ambos casos el sistema puede quedar deslegitimado. La reforma de 2013 ilustra el primer riesgo. Era financieramente sostenible, porque ajustaba el gasto automáticamente a los ingresos, pero socialmente insostenible, porque la pérdida de poder adquisitivo que imponía a los pensionistas provocó una movilización que forzó su suspensión en 2018 y su derogación en 2021 (1.2.2). La expresión "sostenibilidad financiera y social" que figura en el título de la Ley 21/2021 recoge precisamente esta idea: un sistema que no garantiza pensiones suficientes tampoco es sostenible, porque pierde el apoyo que lo mantiene.

Este concepto conecta con la paradoja de la redistribución de Korpi y Palme (1.1.2) y con la economía política de Pierson (1.2.1). La sostenibilidad de un sistema de pensiones no depende solo de sus cuentas, sino también de la coalición política que lo sostiene. Un sistema que concentra los costes en una generación, o que reduce drásticamente las prestaciones de otra, puede ser técnicamente sostenible y políticamente inviable. Su riesgo, examinado en 4.4.1, es el inverso. Un sistema que prioriza la sostenibilidad política a corto plazo, protegiendo a los pensionistas actuales, puede trasladar costes crecientes a las generaciones futuras, que no votan, y comprometer así su propia legitimidad a largo plazo.

\subsubsection{El dilema: sostenibilidad, suficiencia y equidad entre generaciones}

#### El triángulo de hierro

La relación entre la sostenibilidad y la suficiencia puede formularse con una identidad contable sencilla, cuya fuerza explicativa es notable. En un sistema de reparto en equilibrio, los ingresos por cotizaciones deben igualar el gasto en pensiones. Los ingresos son el producto del tipo de cotización por la masa salarial, que es el número de trabajadores por el salario medio. El gasto es el producto del número de pensionistas por la pensión media. Igualando ambas expresiones, el tipo de cotización necesario para el equilibrio es igual a la relación entre pensionistas y trabajadores (la tasa de dependencia del sistema) multiplicada por la relación entre la pensión media y el salario medio (la tasa de sustitución).

Esta identidad, que la literatura denomina a veces el "triángulo de hierro" de las pensiones, establece que, dada la evolución de la tasa de dependencia, que viene determinada en gran medida por la demografía, el equilibrio del sistema solo puede mantenerse ajustando alguna de estas tres variables. La primera es el tipo de cotización, o más en general los ingresos: cotizar más. La segunda es la tasa de sustitución: pensiones menores en relación con los salarios. La tercera es la propia tasa de dependencia, mediante variables que la política puede influir, como la edad de jubilación, el empleo o la inmigración: más trabajadores por cada pensionista. No hay una cuarta opción. Cualquier reforma que se presente como la solución del problema demográfico sin subir los ingresos, sin reducir la generosidad relativa de las pensiones y sin aumentar el número de trabajadores por pensionista, simplemente pospone el ajuste o lo traslada a la deuda pública.

Un ejemplo numérico muestra la magnitud del dilema en el caso español. El Ageing Report de 2024 proyecta que la tasa de dependencia demográfica de España, la relación entre la población de 65 y más años y la de 15 a 64, pasará de en torno a 0,33 en la actualidad a en torno a 0,64 en 2050: prácticamente se duplicará. [Suposición: cálculo ilustrativo con la identidad simplificada] Si la relación entre pensionistas y trabajadores se duplicara en la misma proporción, mantener constantes la tasa de sustitución y el equilibrio del sistema exigiría duplicar el tipo de cotización. Mantener constante el tipo de cotización exigiría, en cambio, reducir a la mitad la pensión media en relación con el salario medio. Ninguna de las dos opciones extremas es política ni socialmente viable, y por eso todas las reformas combinan los tres ajustes. Más empleo y una jubilación más tardía reducen el aumento de la tasa de dependencia del sistema. Una cierta reducción de la tasa de sustitución modera el gasto. Y un aumento de los ingresos, por cotizaciones o por impuestos, cubre el resto. El propio Ageing Report de 2024 proyecta, en efecto, una reducción de la tasa de sustitución de las pensiones españolas desde en torno al 64 por ciento hasta en torno al 56 por ciento en 2050, que compensa parcialmente el aumento de la tasa de dependencia (4.2.3).

#### La equidad entre generaciones

El triángulo de hierro tiene una dimensión temporal que añade un tercer vértice al dilema: la equidad entre generaciones. Cada una de las tres formas de ajuste reparte el coste del envejecimiento de forma distinta entre las generaciones. Subir las cotizaciones hace que el coste lo soporten los trabajadores actuales y futuros. Reducir la tasa de sustitución de las nuevas pensiones hace que lo soporten los futuros jubilados. Reducir la revalorización de las pensiones en curso hace que lo soporten los jubilados actuales. Endeudarse lo traslada a las generaciones futuras, que tendrán que pagar la deuda. La elección entre estas opciones no es técnica, sino distributiva: determina qué generaciones pagan el coste del envejecimiento.

La equidad entre generaciones puede definirse de varias formas, y la elección de la definición condiciona la valoración de las reformas. Una definición exige que todas las generaciones obtengan del sistema la misma rentabilidad, es decir, la misma relación entre lo que aportan y lo que reciben a lo largo de su vida. Otra exige que todas las generaciones soporten la misma carga en relación con su renta. Una tercera, inspirada en el principio de ahorro justo de Rawls (1971), exige que cada generación deje a la siguiente al menos las mismas oportunidades que recibió de la anterior. El factor de sostenibilidad de 2013 respondía a una idea de equidad basada en la igualdad de la pensión total a lo largo de la vida: si una generación vive más, su pensión anual debe ser menor para que el valor total de las pensiones que recibe sea similar al de las generaciones anteriores. [DOC] Como explica el documento en su descripción del factor, la equidad intergeneracional se alcanzaría en la medida en que las generaciones futuras disfrutarían de una pensión inicial menor, pero la recibirían durante más tiempo. El Mecanismo de Equidad Intergeneracional de 2021 y 2023 responde, pese a su nombre, a una idea distinta: que las generaciones actuales contribuyan con una cotización adicional a un fondo que aliviará la carga de las futuras cuando se jubile el baby boom. Son dos concepciones de la equidad entre generaciones que conducen a políticas opuestas, y su comparación se desarrolla en 4.3.

#### La suficiencia como restricción

La suficiencia, que el documento define como la capacidad de prevenir la pobreza y permitir una vida digna en la vejez, actúa como una restricción del dilema. Como se estableció en 1.3.3, tiene rango constitucional (art.s 41 y 50) y se ha concretado recientemente vinculando las pensiones mínimas al umbral de la pobreza (3.1.1). El Pension Adequacy Report de la Comisión Europea mide la suficiencia en tres dimensiones: la protección frente a la pobreza, el mantenimiento de la renta respecto de la vida activa y la duración del periodo de jubilación. La dimensión de la duración es la que más se relaciona con la sostenibilidad. Una generación que vive más tiempo recibe una pensión durante más años, lo que aumenta su suficiencia en términos de pensión total a lo largo de la vida, aunque su pensión anual sea menor. Es el argumento del factor de sostenibilidad de 2013. Sus críticos respondían que la suficiencia se mide en términos de la pensión mensual, que es la que determina el nivel de vida, y no de la pensión total a lo largo de la vida. La elección entre ambas perspectivas es, de nuevo, normativa. Y conecta con el gradiente de longevidad examinado en 1.1.2: los ajustes basados en la esperanza de vida media perjudican a quienes viven menos, que suelen ser los trabajadores de menor renta.

\subsubsection{Cómo se mide la sostenibilidad}

La sostenibilidad no es directamente observable, porque se refiere al futuro. Su medida exige proyecciones basadas en supuestos inciertos, y la elección de los indicadores y de los supuestos condiciona el diagnóstico. Se examinan aquí los cinco enfoques principales, de menor a mayor complejidad.

#### El saldo anual y sus limitaciones

El indicador más sencillo es el saldo anual del sistema: la diferencia entre sus ingresos y sus gastos en un año determinado. Como se vio en 2.3.1, el sistema cerró 2025 con un déficit de 7.387 millones de euros, el 0,4 por ciento del PIB, el nivel más bajo desde 2011. El saldo anual es un indicador muy limitado de la sostenibilidad por tres razones. La primera es que depende del perímetro: incluye las transferencias del Estado, cuya cuantía decide el legislador (2.2.2). La segunda es que depende del ciclo económico: el saldo mejora en las expansiones, cuando aumenta el empleo y la recaudación, y empeora en las recesiones. La tercera, y la más importante, es que no dice nada sobre el futuro. Un sistema puede tener un saldo equilibrado hoy y un desequilibrio creciente en las próximas décadas, porque la jubilación del baby boom todavía no ha alcanzado su intensidad máxima. Por eso el documento, siguiendo la práctica de la literatura, complementa el saldo con las proyecciones de largo plazo.

#### Las proyecciones de gasto: el Ageing Report, la AIReF y el Ministerio

[DOC] El documento explica que el Grupo de Trabajo sobre Envejecimiento de la Comisión Europea (Ageing Working Group) elabora regularmente proyecciones de largo plazo del gasto público en pensiones para todos los países de la Unión con una metodología común. Se publican cada tres años en los Ageing Reports, y el análisis de los últimos informes es muy útil para entender la evolución de la sostenibilidad del sistema español. El último Ageing Report, publicado en 2024, proyecta para España un gasto medio en pensiones del 15,1 por ciento del PIB en el periodo 2022-2050, frente al compromiso del 13,3 por ciento asumido en el Plan de Recuperación. El siguiente informe está previsto para 2027.

Junto al Ageing Report, existen en España otras dos fuentes principales de proyecciones. La AIReF, en su opinión sobre la sostenibilidad de las Administraciones Públicas a largo plazo y en sus evaluaciones de la regla de gasto de pensiones, utiliza un modelo propio. Según su evaluación de 2026, el gasto bruto en pensiones se situaría en un promedio del 14,6 por ciento del PIB hasta 2050. Descontado el efecto de las medidas de ingresos, de 1,6 puntos, el gasto neto sería del 13,0 por ciento. El Ministerio de Inclusión utiliza su propio modelo, conocido como INTegraSS, cuyas proyecciones han resultado en algunos casos más favorables que las de la AIReF. Esas diferencias han dado lugar a discrepancias públicas examinadas en 1.4.6.

Las proyecciones de las distintas instituciones difieren porque dependen de supuestos sobre variables muy inciertas a horizontes de varias décadas. La fecundidad determina el número de trabajadores futuros. La esperanza de vida determina el número de pensionistas y la duración de las pensiones. La inmigración afecta al número de trabajadores a corto plazo y al de pensionistas a largo plazo. La productividad determina el crecimiento de los salarios y del PIB. Y las tasas de empleo por edad y sexo determinan cuántas personas en edad de trabajar trabajan efectivamente. Pequeñas diferencias en estos supuestos producen grandes diferencias en las proyecciones a cuarenta o cincuenta años, y la elección de los supuestos no es políticamente neutral. Un supuesto optimista sobre la inmigración o la productividad reduce el gasto proyectado y, con él, la necesidad de reformas. Por eso la cláusula de salvaguarda del Mecanismo de Equidad Intergeneracional toma como referencia el Ageing Report, elaborado por una instancia técnica europea, y no las proyecciones del propio Gobierno. Es una forma de protegerse frente al sesgo optimista que la teoría de las instituciones fiscales independientes examinada en 1.4.6 atribuye a los gobiernos.

#### La deuda implícita: los derechos de pensión devengados

Un enfoque distinto consiste en medir la deuda implícita del sistema, es decir, el valor actual de las pensiones que el sistema ha prometido pagar. La literatura distingue dos conceptos. La deuda de grupo cerrado (closed-group liabilities) es el valor actual de las pensiones devengadas por los jubilados y trabajadores actuales por las cotizaciones ya pagadas, sin considerar las cotizaciones y pensiones futuras de nuevos trabajadores. Es la deuda que el sistema tendría que pagar si se cerrara hoy. La deuda de grupo abierto (open-group liabilities) incluye también las cotizaciones y pensiones de las generaciones futuras. Mide el desequilibrio del sistema en su funcionamiento continuado. Como se vio en 2.1.1, el Sistema Europeo de Cuentas obliga desde 2017 a publicar una tabla complementaria con los derechos de pensión devengados (la tabla 29), que corresponde aproximadamente al concepto de grupo cerrado.

La deuda implícita de grupo cerrado de los sistemas de reparto maduros es de varias veces el PIB, muy superior a la deuda pública explícita. Pero, como se estableció en 2.1.1, esta cifra no mide un desequilibrio, sino simplemente el volumen del compromiso del sistema, que en un sistema de reparto se financia con las cotizaciones futuras y no con activos acumulados. Es un indicador relevante para evaluar el coste de una transición a la capitalización, pero no para evaluar la sostenibilidad del reparto. Para esto último es más útil la deuda de grupo abierto, cuya estimación es mucho más sensible a los supuestos.

#### El balance actuarial

Suecia desarrolló, como complemento de su sistema de cuentas nocionales, un instrumento que permite evaluar la solvencia de un sistema de reparto de forma análoga al balance de una empresa: el balance actuarial. Su innovación consiste en valorar como activo del sistema no solo el fondo de reserva, sino también las cotizaciones futuras que financiarán las pensiones, calculadas mediante la llamada "duración de rotación" (turnover duration): el tiempo medio que transcurre entre el pago de una cotización y el cobro de la pensión que genera. Si el valor de los activos, cotizaciones más fondo, iguala o supera el valor de los pasivos, los derechos de pensión devengados, el sistema es solvente. Si no, se activa el mecanismo de equilibrio automático examinado en 1.2.1. Boado-Penas, Valdés-Prieto y Vidal-Meliá (2008) aplicaron esta metodología al sistema español de jubilación y encontraron un desequilibrio actuarial significativo: los pasivos superaban a los activos, lo que indicaba que el sistema, con sus reglas de entonces, no era solvente a largo plazo. La magnitud exacta y su actualización deben verificarse. El balance actuarial tiene la ventaja de ofrecer una medida sintética de la solvencia, comparable en el tiempo y entre países, y el inconveniente de depender, como todas las medidas prospectivas, de los supuestos sobre las cotizaciones futuras.

#### La contabilidad generacional y las tasas internas de rendimiento

El enfoque más ambicioso es la contabilidad generacional, propuesta por Auerbach, Gokhale y Kotlikoff (1991). En lugar de medir el saldo del sector público en un año, calcula para cada generación el valor actual de los impuestos que pagará menos las transferencias que recibirá a lo largo de su vida. Si la carga neta de las generaciones futuras es superior a la de las actuales, la política fiscal es intergeneracionalmente desequilibrada: traslada costes a quienes todavía no han nacido. La contabilidad generacional permite visualizar de forma directa la dimensión intergeneracional del dilema examinado en 4.1.2. En los países con sistemas de reparto generosos y poblaciones envejecidas, suele mostrar una carga neta muy superior para las generaciones futuras.

La contabilidad generacional ha sido objeto de críticas influyentes. Haveman (1994) sostuvo que ignora los beneficios del gasto público que no son transferencias, como la educación o las infraestructuras, que benefician a las generaciones futuras. También señaló que sus resultados dependen de supuestos arbitrarios sobre el tipo de descuento y sobre quién soporta la carga de los impuestos. Buiter (1997) criticó que presupone que la política fiscal actual se mantendrá indefinidamente, lo que es irreal. Una metodología relacionada, las cuentas nacionales de transferencias (National Transfer Accounts) desarrolladas por Lee y Mason (2011), mide cómo se reparten el consumo y la renta entre las edades a lo largo del ciclo vital, incluidas las transferencias privadas dentro de las familias. Ha permitido mostrar que en España, como en otros países mediterráneos, las transferencias familiares hacia los jóvenes compensan parcialmente las transferencias públicas hacia los mayores.

Un indicador más sencillo y muy utilizado en España es la tasa interna de rendimiento (TIR) que obtiene cada generación de sus cotizaciones: la rentabilidad implícita que igualaría el valor actual de lo que cotiza y de lo que recibe. Los trabajos de Devesa, Devesa, Meneu y sus coautores han estimado, para distintas cohortes y perfiles de carrera, que la TIR del sistema español de jubilación ha sido superior a la tasa de crecimiento de la masa salarial que, según la condición de Aaron examinada en 2.1.1, el reparto puede pagar de forma sostenible. Es decir, el sistema ha sido más generoso de lo que su propia lógica financiera permite. La diferencia debe financiarse con más ingresos, con menores pensiones futuras o con deuda. Las referencias exactas y sus estimaciones deben verificarse. La comparación entre la TIR de distintas cohortes permite además evaluar la equidad intergeneracional: si las generaciones que se jubilan hoy obtienen una TIR muy superior a la que obtendrán las que se jubilen dentro de treinta años, el sistema está trasladando recursos de las segundas a las primeras.

\subsubsection{Por qué importa la sostenibilidad y quién debe garantizarla}

#### Por qué importa

La sostenibilidad importa, en primer lugar, por el peso del sistema de pensiones en la economía y en las finanzas públicas. Como se vio en 2.3.1, las pensiones absorben en torno al 28 por ciento del gasto público total y en torno al 13 por ciento del PIB. Cualquier desequilibrio del sistema tiene, por tanto, consecuencias macroeconómicas y presupuestarias de primer orden, y reduce el margen para otras políticas públicas, como la educación, la inversión o la protección de los jóvenes (2.3.2). Importa, en segundo lugar, por la dependencia de los mayores respecto de las pensiones. [DOC] Como señala el documento, las pensiones son la principal fuente de ingresos de la población mayor de 65 años. Una crisis del sistema no afectaría solo a las cuentas públicas, sino a las condiciones de vida de más de nueve millones de pensionistas y de sus familias.

Importa, en tercer lugar, por la credibilidad. Como se examinó en 1.1.1, el valor del seguro que ofrece un sistema de pensiones depende de que sus promesas sean creíbles. Un trabajador que duda de que el sistema vaya a pagarle su pensión cuando se jubile valora menos sus cotizaciones y las percibe más como un impuesto que como un precio, con los efectos sobre el empleo y la economía sumergida examinados en 2.1.2. La teoría de la inconsistencia temporal de Kydland y Prescott (1977), examinada en 1.4.6, explica por qué la credibilidad es frágil: un gobierno tiene siempre la tentación de prometer pensiones generosas hoy y dejar su financiación a los gobiernos futuros. Un sistema cuya sostenibilidad está garantizada por reglas creíbles reduce esa tentación y refuerza la confianza. Uno cuya sostenibilidad depende de decisiones discrecionales anuales, como la revalorización aprobada por decreto-ley en 2026 (1.2.2), la debilita.

Importa, en cuarto lugar, por la justicia entre generaciones, examinada en 4.1.2. Un sistema insostenible traslada a las generaciones futuras la carga de pagar las promesas hechas a las actuales, sin que aquellas hayan podido participar en la decisión. Es la versión en materia de pensiones de un problema general de la política económica. Tiene paralelos en el debate sobre el cambio climático, en el que la elección de la tasa de descuento de los costes y beneficios futuros determina cuánto peso se da al bienestar de las generaciones venideras.

#### Quién debe garantizarla

La garantía de la sostenibilidad corresponde, en último término, al legislador, que es quien fija los parámetros del sistema dentro del amplio margen que le concede la Constitución (1.3.1). Pero la experiencia española y la internacional muestran que el legislador, sometido a la presión del ciclo electoral y de la coalición de los pensionistas, tiende a posponer los ajustes. Por eso los sistemas de pensiones han desarrollado instituciones y reglas que limitan su discrecionalidad. En España, esas instituciones y reglas son cuatro, y se han examinado en los bloques anteriores. La primera es el Pacto de Toledo, como mecanismo de consenso que reparte la culpa de las reformas impopulares (1.2.2 y 1.4.6). La segunda es la AIReF, como evaluador independiente de la sostenibilidad (1.4.6). La tercera es la regla de gasto del Mecanismo de Equidad Intergeneracional, que obliga a adoptar medidas si el gasto supera un umbral de referencia (4.3.3). La cuarta es la gobernanza económica europea, que integra la sostenibilidad de las pensiones en la supervisión fiscal (1.3.4). La cuestión que queda abierta, y que se examina en 4.4, es si estas instituciones son suficientes para garantizar la sostenibilidad frente a las presiones políticas, o si, como la historia de las reformas de 2013 sugiere, cualquier regla puede ser suspendida o derogada cuando su aplicación se vuelve impopular.

Con este apartado quedan establecidos los conceptos necesarios para el diagnóstico. La sostenibilidad del sistema español debe evaluarse con el gasto total en pensiones en relación con el PIB y con los recursos totales que lo financian, no con el saldo del sistema. Exige elegir entre ajustes que reparten el coste del envejecimiento de forma distinta entre las generaciones. Y su medida depende de proyecciones inciertas y de supuestos que no son políticamente neutrales. El apartado siguiente aplica estos conceptos a la descomposición del gasto en pensiones, que permite identificar qué factores explican su evolución pasada y proyectada.


\subsection{El diagnóstico: la descomposición del gasto en pensiones}
Antes del texto, un aviso que condiciona todo el apartado. La descomposición del gasto que el documento presenta como "manera común de analizar el sistema" es una identidad contable, no una teoría. Se cumple siempre, por construcción, y por eso no explica nada por sí sola. Si en el examen se usa para atribuir causas ("el 60\% del aumento se debe a la generosidad") sin advertir esa limitación, el tribunal puede leerlo como un uso ingenuo del instrumento. He redactado el apartado para que la identidad sirva de hilo conductor y, a la vez, para enseñar sus límites. Además, los datos 2008-2021 del documento tienen un sesgo cíclico importante en los años de referencia, que corrijo y actualizo con el Ageing Report 2024 y las proyecciones del INE de junio de 2026.

Hasta aquí hemos definido qué significa que un sistema de pensiones sea sostenible, por qué esa noción tiene al menos cuatro dimensiones y con qué instrumentos se mide. Falta la pregunta que interesa al decisor público: por qué se mueve el gasto. Es la pregunta previa a cualquier reforma. Solo sabiendo qué fuerzas empujan el gasto al alza puede decidirse sobre cuál actuar, a qué coste y con qué reparto de la carga entre generaciones, entre cotizantes y pensionistas, y entre trabajo y capital. [DOC] El documento resuelve esta cuestión con una herramienta muy extendida en la literatura aplicada y en los organismos que proyectan el gasto (la Comisión Europea, la AIReF, el Banco de España, la OCDE): descomponer el gasto en pensiones contributivas, medido en proporción del PIB, en el producto de cuatro factores.

La identidad se construye multiplicando y dividiendo por las magnitudes adecuadas. Si llamamos G al gasto en pensiones, Y al PIB, N al número de pensiones, p a la pensión media, P a la población de 16 a 64 años, L al número de ocupados y w al salario medio, el gasto es G = N·p, y puede escribirse:

G/Y = (N/P) · (P/L) · (p/w) · (wL/Y)

Basta simplificar para comprobarlo: N/P por P/L da N/L; al multiplicarlo por p/w queda Np/(Lw); y al multiplicarlo por wL/Y queda Np/Y, que es G/Y. (OJO: la nota 22 del documento está ilegible en la versión de trabajo. Esta es su formulación correcta, y conviene escribirla así en el examen, con la simplificación visible, porque demuestra que se entiende que es una identidad.)

Cada cociente tiene una lectura económica. El primero, pensiones sobre población en edad de trabajar, es el factor demográfico. Más exactamente, mezcla demografía y cobertura: no mide cuántos mayores hay por cada persona en edad de trabajar, sino cuántas pensiones hay, y eso depende también de cuándo se jubila la gente, de cuántas pensiones de viudedad, incapacidad u orfandad existen y de la concurrencia de pensiones. El segundo, población en edad de trabajar sobre ocupados, es la inversa de la tasa de empleo y constituye el factor empleo. Cuanto mayor sea la proporción de ocupados, mayor será el PIB con la misma población y menor el peso relativo de las pensiones. El tercero, pensión media sobre salario medio, es la tasa de sustitución en sentido agregado, o factor de generosidad. El cuarto, masa salarial sobre PIB, es la participación de los salarios en la renta. [DOC] Como recoge la nota 23 del documento, los dos últimos factores pueden fundirse en uno: (p/w)·(wL/Y) = p/(Y/L), es decir, la pensión media sobre la productividad media por ocupado. Dado el salario medio, cuanto menor sea la productividad, mayor será el gasto sobre el PIB.

Antes de recorrer los cuatro factores, conviene fijar tres advertencias metodológicas que la literatura considera imprescindibles y que el documento omite.

La primera: identidad no es causalidad. La descomposición reparte contablemente la variación del gasto entre los factores, pero no dice nada de las relaciones entre ellos, y esas relaciones son intensas. Un aumento de la inmigración reduce hoy el factor demográfico y mejora el de empleo, pero en un sistema contributivo genera derechos que elevarán el primer factor dentro de treinta o cuarenta años. Un aumento de la productividad reduce hoy la ratio pensión/productividad, pero si los salarios crecen, las pensiones iniciales, calculadas sobre esos salarios, también crecerán, y la generosidad acabará reflejándolo en parte. Un retraso de la edad de jubilación reduce el número de pensiones y aumenta el empleo, pero también eleva la cuantía de las pensiones que se causan más tarde, por la bonificación por demora y por las carreras más largas. La descomposición toma una fotografía de cada factor por separado; las políticas actúan sobre todos a la vez. Por eso los organismos que proyectan el gasto no se limitan a descomponerlo: construyen modelos de simulación o de equilibrio general (los modelos de generaciones solapadas que veremos en 4.3) que capturan esas interdependencias. La descomposición es el lenguaje con el que se presentan los resultados, no el motor que los genera.

La segunda: la variante oficial es distinta. La que utilizan el Grupo de Trabajo sobre Envejecimiento (Ageing Working Group, AWG) de la Comisión y los Ageing Reports no coincide con la del documento, aunque responde a la misma lógica, como advierte su nota 25. El AWG separa en dos lo que el documento funde en el primer factor. Por un lado, la tasa de dependencia, que es la población de 65 años o más sobre la de 20 a 64 (no la de 16 a 64). Por otro, la tasa de cobertura, que es el número de pensionistas sobre la población de 65 años o más. Divide además el efecto del mercado de trabajo en la tasa de empleo y la intensidad laboral (las horas trabajadas), y define la ratio de prestación (benefit ratio) como la pensión media sobre el salario medio de la economía. La separación no es un detalle técnico. La tasa de dependencia es casi pura demografía y apenas responde a la política de pensiones a medio plazo. La tasa de cobertura, en cambio, es el lugar donde actúan la edad efectiva de jubilación, las jubilaciones anticipadas y la incapacidad permanente, es decir, donde el legislador sí tiene palanca. Fundir ambas en un único "factor demográfico" atribuye a la demografía lo que en parte es regulación.

La tercera: el factor demográfico no es independiente del saldo. Esta advertencia es la más profunda y enlaza la identidad con la teoría del reparto del bloque II. Los ingresos por cotizaciones, en proporción del PIB, son el tipo efectivo de cotización τ multiplicado por la masa salarial sobre el PIB: C/Y = τ·(wL/Y). El saldo del sistema contributivo puro es entonces:

(C − G)/Y = (wL/Y) · [τ − (N/L) · (p/w)]

La participación salarial aparece multiplicando a ingresos y gastos a la vez. En un sistema financiado exclusivamente con cotizaciones sobre salarios, su evolución no altera el signo del saldo, solo su escala. El equilibrio exige que el tipo de cotización iguale el producto de dos cocientes: el número de pensiones por ocupado, o tasa de dependencia del sistema (N/L), y la tasa de sustitución (p/w). Es la condición de equilibrio del reparto, versión contable de lo que vimos con Samuelson (1958) y Aaron (1966): en un sistema de reparto maduro, el tipo que equilibra es igual a la ratio de dependencia efectiva por la tasa de sustitución. Una ilustración [Suposición: cifras redondeadas con fines didácticos]. Con unas 0,43 pensiones por afiliado (el inverso de las 2,35 personas afiliadas por pensión que registraba la Seguridad Social en 2026) y una tasa de sustitución agregada en torno al 0,6, el tipo de equilibrio rondaría el 26%. Si la ratio de dependencia del sistema subiera hacia 0,7 hacia 2050 y la tasa de sustitución se mantuviera, el tipo de equilibrio superaría el 40%. Este cálculo sencillo es lo más útil que ofrece la descomposición. Hace visible el triángulo de hierro de 4.1.2: si N/L sube por razones demográficas, el equilibrio solo se mantiene subiendo τ (más cotizaciones), reduciendo p/w (menos generosidad) o conteniendo N/L por otras vías (más empleo o jubilación más tardía). No hay cuarta vía dentro del sistema. La única salida es aportar recursos de fuera, mediante transferencias del Estado financiadas con otros impuestos, y eso desplaza la pregunta de la sostenibilidad del sistema a la sostenibilidad fiscal (4.1.1).

Con estas advertencias, recorremos los cuatro factores. Cada epígrafe explica qué mide el factor, cómo ha evolucionado y se proyecta en España, qué debates teóricos y empíricos suscita y qué margen de actuación ofrece. Cerramos con la evidencia agregada: la descomposición del período 2008-2021 que aporta el documento y la que proyecta el Ageing Report 2024.

\subsubsection{Factor demográfico: natalidad, longevidad, baby boom y migración}

[DOC] El documento parte de una afirmación correcta. Al ser un sistema de reparto, su sostenibilidad depende crucialmente de la ratio entre beneficiarios y contribuyentes: los cotizantes actuales sostienen las prestaciones actuales con la promesa de que la generación siguiente pagará las suyas. En España, la caída de la natalidad, el aumento de la esperanza de vida, una entrada proyectada de inmigrantes que el documento califica de débil y el peso de la generación del baby boom transformarán la pirámide de población hasta parecerse cada vez más a una pirámide invertida, según recoge Hernández de Cos (2021) a partir de las proyecciones del INE, la AIReF y Eurostat. El diagnóstico es correcto en lo esencial, pero agrupa bajo una misma etiqueta cuatro fenómenos con lógicas temporales, causas e implicaciones de política muy distintas. Conviene separarlos.

a) La longevidad: el componente permanente y, en sí mismo, una buena noticia.

El envejecimiento tiene dos motores: se vive más (envejecimiento por la cúspide) y nacen menos niños (envejecimiento por la base). La longevidad es el más persistente y el menos discutido. Según las proyecciones de población 2026-2076 que el INE publicó el 17 de junio de 2026 [Probable en las cifras exactas], la esperanza de vida al nacer alcanzaría hacia 2075 unos 87 años en los hombres y 90 en las mujeres. A los 65 años llegaría a unos 23,5 y 26,3 años adicionales, respectivamente, entre 2,7 y 3,7 años más que ahora. En términos de pensiones, lo relevante es la esperanza de vida a los 65, porque determina durante cuánto tiempo se cobra. Un sistema diseñado cuando un jubilado de 65 años vivía de media unos 13 o 14 años más, como ocurría en España en los años setenta y ochenta [Probable], paga hoy prestaciones durante veinte años o más sin que la edad de acceso se haya movido en la misma proporción.

El enfoque ortodoxo extrae de aquí una conclusión casi aritmética: si la vida se alarga y la edad de jubilación no, la duración de la pensión crece y, con ella, el valor actuarial de lo prometido, sin que las cotizaciones lo reflejen. De ahí proceden los mecanismos que vinculan algún parámetro a la esperanza de vida, que analizaremos en 4.3: el factor de sostenibilidad de 2013, la vinculación de la edad legal a la longevidad en Dinamarca, Países Bajos, Finlandia o Italia, o el divisor de las cuentas nocionales en Suecia (Settergren y Mikula, 2005). La defensa teórica viene de Diamond (2004) y Barr y Diamond (2008, 2009): si la longevidad aumenta, lo eficiente es repartir las ganancias de vida entre más años de trabajo y más años de jubilación, en lugar de destinarlas íntegramente a la jubilación.

Frente a ese razonamiento se han formulado tres críticas de peso.

La primera es distributiva. La esperanza de vida no es homogénea: presenta un gradiente pronunciado por renta, nivel educativo y ocupación, como vimos con Chetty et al. (2016) para Estados Unidos. En España los estudios del INE sobre esperanza de vida por nivel de estudios muestran diferencias de varios años a los 65 entre los extremos educativos. Vincular la edad de jubilación a una media que crece sobre todo por la parte alta de la distribución castiga desproporcionadamente a quienes tienen trabajos manuales, carreras más largas y menor longevidad. La literatura sobre la regresividad de los sistemas de pensiones ante la mortalidad diferencial (Ayuso, Bravo y Holzmann, 2017; National Academies of Sciences, 2015) subraya que este efecto puede neutralizar en parte la progresividad de la fórmula de cálculo.

La segunda es conceptual, y es un buen gancho para el examen. Sanderson y Scherbov (2005, Nature; 2010, Science) propusieron medir la edad de forma prospectiva: por los años que quedan por vivir, no por los vividos. Con esa métrica, "ser mayor" es haber llegado a la edad en la que la esperanza de vida restante es de, por ejemplo, quince años, y ese umbral se desplaza con la longevidad. Recalculada así, la tasa de dependencia crece mucho menos que la convencional. Spijker y MacInnes (2013, BMJ) llegaron a titular un art. "Population ageing: the timebomb that isn't?", y en España los demógrafos del CSIC Julio Pérez Díaz y MacInnes (2009, Sociological Review) describen el proceso como una revolución reproductiva: una forma más eficiente de reproducir poblaciones, en la que se necesitan menos nacimientos porque casi todos los nacidos llegan a la edad adulta. No es una patología que haya que corregir. Esta crítica tiene fuerza para relativizar el alarmismo, pero no disuelve el problema financiero: aunque se viva con más salud, mientras la edad de acceso a la pensión esté fijada en años de calendario, cada año adicional de vida es un año adicional de prestación. La edad prospectiva es un argumento para mover la edad de jubilación, no para ignorar la longevidad.

La tercera es de economía política. Vincular automáticamente un parámetro a la longevidad es precisamente un mecanismo de evitación de la culpa (Weaver, 1986): el recorte se produce sin que ningún gobierno tenga que votarlo. Sus defensores lo presentan como una virtud, porque despolitiza y da previsibilidad (Kydland y Prescott, 1977, en la lógica de la regla frente a la discrecionalidad que vimos en 4.1.4). Sus críticos, como una erosión de la rendición de cuentas. La reversión española de 2018-2021 demuestra que estas reglas son revocables cuando su efecto se hace visible, y lo veremos en 4.3.

b) La natalidad: la causa estructural con efectos diferidos.

España tiene una de las fecundidades más bajas del mundo. Según el INE, el indicador coyuntural de fecundidad fue de 1,10 hijos por mujer en 2024, muy por debajo del nivel de reemplazo de 2,1. Las proyecciones de 2026 suponen una leve recuperación hasta 1,16 en 2040 [Probable]. España fue, junto con Italia, el caso que llevó a Kohler, Billari y Ortega (2002, Population and Development Review) a acuñar el término fecundidad muy baja (lowest-low fertility), por debajo de 1,3. Sus causas enlazan con lo visto en el bloque II: la emancipación tardía, la precariedad laboral de los jóvenes, el coste de la vivienda, la penalización por maternidad que documentan Kleven et al. y De Quinto, Hospido y Sanz, y la "revolución incompleta" de Esping-Andersen (2009), en la que la incorporación de la mujer al empleo no ha ido acompañada de un reparto de los cuidados ni de servicios públicos que la hagan compatible con la crianza. Bongaarts y Feeney (1998) añadieron un matiz técnico: parte de la caída del indicador es un efecto calendario (tempo effect). Las mujeres retrasan la maternidad, lo que deprime temporalmente el indicador coyuntural aunque la descendencia final caiga menos.

La verdad incómoda de este epígrafe es la siguiente. Una recuperación de la natalidad no resuelve el problema de las pensiones en el horizonte que importa, y a corto plazo lo empeora. Un niño nacido en 2027 no entrará en el mercado de trabajo, en el mejor de los casos, hasta mediados de los años cuarenta. Hasta entonces, eleva la tasa de dependencia total (la de los menores) y el gasto en educación, sanidad infantil y prestaciones familiares. Como la tensión máxima del sistema se concentra en torno a 2050, según veremos, las políticas de natalidad llegan tarde para aliviar el pico. Pueden estar justificadas por otras razones: la igualdad de género, la libertad reproductiva (la brecha entre los hijos deseados y los tenidos es muy grande en España [Probable, según la Encuesta de Fecundidad del INE de 2018]) o la sostenibilidad demográfica a muy largo plazo. Pero no deben presentarse como política de pensiones. (OJO: en el examen conviene formularlo así, porque es un error frecuente en el debate público y el tribunal lo valorará.)

c) El baby boom: ¿joroba transitoria o meseta permanente?

El baby boom español fue tardío respecto al europeo y al estadounidense, que se produjeron en la posguerra inmediata. Se extendió aproximadamente de finales de los cincuenta a finales de los setenta, con un máximo en torno a 1976, cuando nacieron cerca de 680.000 niños, más del doble de los alrededor de 320.000 de 2024. Estas cohortes, las más numerosas de la historia de España, han empezado a jubilarse y lo harán masivamente entre mediados de los veinte y mediados de los cuarenta. Su paso de cotizantes a pensionistas produce un doble efecto: resta del denominador y suma al numerador de la ratio de dependencia.

La lectura habitual presenta este fenómeno como una joroba (hump): un bulto transitorio de gasto que desaparecerá cuando estas cohortes fallezcan, a partir de mediados de siglo. La interpretación tiene consecuencias de política importantes. Si el problema es transitorio, lo racional es financiarlo con un colchón, es decir, un fondo de reserva que se acumula antes y se agota durante la joroba, o con deuda, sin ajustar permanentemente los parámetros. Esta es, en esencia, la lógica del Mecanismo de Equidad Intergeneracional (MEI) de 2021-2023, que analizaremos en 4.3. Es también la tesis de los defensores de las reformas recientes: el sistema debe "pasar la joroba" sin recortar derechos.

Sin embargo, las proyecciones del Ageing Report 2024 matizan mucho esta lectura. [Probable en las cifras] La tasa de dependencia de la vejez (65 años o más sobre 20-64) casi se duplicaría, del 33,3% en 2022 al 63,9% en 2050, y se estabilizaría después, en torno al 64,5% en 2070, sin volver a bajar. El gasto bruto en pensiones pasaría del 13,1% del PIB en 2022 a un máximo del 17,3% en torno a 2050, para descender solo ligeramente hasta el 16,7% en 2070. Las proyecciones del INE de 2026 van en la misma dirección: la proporción de mayores de 65 pasaría del 21,1% actual al 30,9% en 2076, su máximo al final del horizonte, y la tasa de dependencia alcanzaría su máximo en 2076. No hay joroba: hay una escalera con un pequeño rellano. Cuando las cohortes del baby boom desaparezcan, su lugar en la pirámide lo ocuparán cohortes menos numerosas pero mucho más longevas, que se sostendrán sobre una base de fecundidad muy baja. El envejecimiento es estructural, no coyuntural.

La implicación para el debate es directa, y es una verdad incómoda para los defensores de un enfoque puramente transitorio. Un fondo de reserva puede suavizar la transición, pero no puede sustituir a un ajuste permanente de ingresos, gastos o edad, porque el nuevo equilibrio demográfico es permanente. Es un argumento que, como veremos en 4.3, han formulado la AIReF, el Banco de España y la propia Comisión Europea al valorar el MEI.

d) La migración: la variable más volátil y la más politizada.

La inmigración es el único componente demográfico que puede alterar la estructura por edades de forma rápida, porque los inmigrantes llegan mayoritariamente en edad de trabajar. España lo comprobó entre 2000 y 2008, cuando recibió varios millones de inmigrantes, en uno de los mayores flujos relativos de la OCDE en ese período. En aquellos años el número de afiliados creció con fuerza y el Fondo de Reserva se acumuló. Ha vuelto a comprobarlo desde 2022, con saldos migratorios muy elevados. Las proyecciones del INE de 2026 suponen un saldo migratorio de unas 615.000 personas en 2026, decreciente después, y estiman que la población nacida en España pasaría del 79,8% en 2026 al 59,6% en 2076. El Ageing Report 2024 suponía flujos netos medios de unas 220.000 a 230.000 personas anuales entre 2030 y 2050. [DOC] La calificación de "débil" que el documento da a la inmigración proyectada, heredada de proyecciones anteriores, está desactualizada: tanto el INE como el AWG han revisado al alza sus hipótesis migratorias.

La literatura ofrece aquí tres advertencias.

La primera es la de la migración de reemplazo. En 2000, la División de Población de Naciones Unidas publicó un informe célebre, Replacement Migration: Is It a Solution to Declining and Ageing Populations?, que calculaba cuántos inmigrantes necesitaría cada país para mantener constante su ratio de apoyo potencial (población de 15 a 64 años por cada persona de 65 o más). La conclusión fue que las cifras eran tan elevadas que resultaban inverosímiles: en varios casos habrían multiplicado la población total. En España, el Banco de España estimó en su Informe Anual de 2018 un orden de magnitud de decenas de millones de inmigrantes en edad de trabajar hasta mediados de siglo para mantener constante la ratio de dependencia. La inmigración puede mitigar, no neutralizar, el envejecimiento.

La segunda es la del esquema Ponzi demográfico. En un sistema contributivo, el inmigrante que cotiza hoy adquiere derechos que cobrará mañana. Si la solución al envejecimiento de la población actual es la inmigración, el envejecimiento de esos inmigrantes exigirá más inmigración en el futuro. Storesletten (2000, Journal of Political Economy) formalizó que el efecto fiscal neto de la inmigración depende de manera crucial de la edad de llegada y de la cualificación: los inmigrantes cualificados que llegan jóvenes aportan un valor actual neto positivo, mientras que los que llegan con baja cualificación o en edades avanzadas pueden aportar uno negativo. En España, la concentración de la inmigración en empleos de baja productividad y salarios bajos limita su contribución: cotizan por bases bajas y, si llegan a la pensión con carreras incompletas, pueden necesitar complementos a mínimos.

La tercera es de economía política. Las hipótesis migratorias de las proyecciones son las más inciertas y las más sensibles a la política migratoria, al ciclo económico y a la opinión pública. Una proyección de sostenibilidad que dependa sobre todo de flujos migratorios elevados y sostenidos durante cincuenta años es frágil. (OJO: en el contexto de 2025-2026 conviene mencionar la regularización extraordinaria impulsada por la iniciativa legislativa popular admitida a trámite en 2024 y el nuevo Reglamento de Extranjería aprobado en noviembre de 2024 y en vigor desde mayo de 2025, que facilita el arraigo. Su efecto sobre la afiliación es positivo a corto plazo, pero su efecto neto sobre la sostenibilidad depende de la productividad y de las carreras de cotización de los regularizados. Pendiente de verificar el estado de tramitación de la regularización a la fecha del examen.)

Síntesis del factor demográfico. La demografía es el factor que más empuja el gasto al alza y el que menos responde a la política de pensiones en el horizonte relevante. La longevidad es permanente, la natalidad actúa con décadas de retraso y la inmigración es volátil y solo mitiga el problema. Por eso las reformas se han concentrado en los otros tres factores, y por eso el debate se ha desplazado de "cómo evitar el envejecimiento" a "cómo repartir su coste".

---

\subsubsection{Factor empleo: tasa de ocupación, edad efectiva de salida y cobertura}

[DOC] El segundo factor de la identidad es la inversa de la tasa de empleo: la población en edad de trabajar sobre los ocupados. Cuanto mayor es la proporción de ocupados, mayor es el PIB por persona en edad de trabajar y menor el peso de las pensiones sobre él. El documento atribuye a este factor en torno al 10% del aumento del gasto entre 2008 y 2021, y señala que la crisis y la pandemia hundieron el número de cotizantes. El factor merece más atención de la que recibe, por una razón que el documento no explicita. Es el factor en el que España tiene mayor margen respecto a sus socios europeos, y a la vez aquel cuyo efecto sobre la sostenibilidad es más ambiguo en un sistema contributivo.

a) La anomalía española: un mercado de trabajo que desperdicia recursos.

El mercado de trabajo español se ha caracterizado durante décadas por tasas de paro que duplican la media europea, una elevada temporalidad, una fuerte sensibilidad del empleo al ciclo y tasas de actividad bajas en los extremos de la vida laboral: jóvenes que tardan en incorporarse y trabajadores mayores que salen pronto. Según la EPA del segundo trimestre de 2026, publicada por el INE, la tasa de paro fue del 9,87%, la primera vez que baja del 10% desde 2008, y la ocupación alcanzó unos 22,8 millones de personas, máximo histórico. La tasa de empleo de la población de 20 a 64 años se sitúa, aun así, varios puntos por debajo de la media de la UE y lejos del objetivo europeo del 78% para 2030.

El ejercicio mental más útil es el contrafactual. [Suposición: orden de magnitud ilustrativo] Si España tuviera la tasa de empleo de los países nórdicos o de los Países Bajos, con el mismo número de pensiones y la misma pensión media, su gasto en pensiones sobre el PIB sería varios puntos menor, simplemente porque habría más producción sobre la que medirlo. Este es el núcleo del argumento de quienes sostienen, desde posiciones heterodoxas o socialdemócratas, que el problema de las pensiones es un problema de empleo: Navarro, Zubiri o los economistas agrupados en Economistas Frente a la Crisis, y también buena parte del discurso sindical. Tiene un fundamento sólido en la identidad, y enlaza con la tesis de Barr (2001) que vimos en 4.1.1: "la producción es lo central", porque lo que financia las pensiones es la producción de cada momento, no los mecanismos financieros.

b) El contraargumento ortodoxo: más empleo hoy son más pensiones mañana.

La réplica viene de la propia naturaleza del sistema. En un sistema de reparto contributivo, el aumento del empleo genera más cotizaciones hoy, pero también más derechos de pensión en el futuro. Un trabajador adicional que cotiza treinta y cinco años cobrará después una pensión proporcional a lo cotizado. El efecto neto sobre la sostenibilidad a largo plazo depende de la rentabilidad implícita que el sistema ofrece a esas cotizaciones. Si la tasa interna de rendimiento del sistema supera la tasa de crecimiento sostenible de la masa salarial (la condición de Aaron vista en 2.2), cada cotizante nuevo es, en valor actual, un pasivo neto, no un activo. Los trabajos de Devesa, Devesa, Domínguez, Encinas, Meneu y Nagore sobre la TIR del sistema español, que vimos en 4.1.3, apuntan a que esa rentabilidad implícita ha sido elevada en términos comparados [Probable]. Desde esta perspectiva, que representan De la Fuente, Conde-Ruiz o Doménech, el empleo alivia la transición pero no corrige el desequilibrio estructural si no va acompañado de ajustes en la relación entre cotizaciones y prestaciones.

Las dos posiciones son compatibles si se distingue el horizonte temporal. A corto y medio plazo, el empleo mejora el saldo sin ambigüedad. A largo plazo, su efecto depende del diseño del sistema. El aumento del empleo es, sin embargo, especialmente valioso durante la joroba, o más bien durante la subida de la escalera: los cotizantes adicionales de 2025-2050 pagan pensiones del baby boom, y sus propios derechos vencerán cuando la estructura por edades se haya estabilizado. (OJO: esta distinción temporal es la clave para superar el falso dilema "empleo o recortes" en la exposición.)

c) La edad efectiva de salida y la tasa de cobertura: el factor empleo en el tramo final de la vida laboral.

La palanca con más efecto sobre este factor, y sobre el demográfico en la variante del AWG, es la edad efectiva de salida del mercado de trabajo. Retrasarla actúa por partida doble: reduce el número de pensiones (el numerador del primer factor o la tasa de cobertura) y aumenta el número de ocupados (el denominador del segundo). La literatura de Gruber y Wise (1999, 2004), que ya vimos en 3.2, demostró que los incentivos implícitos de los sistemas de pensiones (el "impuesto implícito" sobre seguir trabajando) explican una parte sustancial de la caída de la participación laboral de los mayores de 55 años en los países de la OCDE entre los años setenta y noventa. Para España, Boldrin, Jiménez-Martín y Peracchi documentaron el peso de las prejubilaciones y de las jubilaciones anticipadas, a menudo financiadas indirectamente con la prestación por desempleo o con el subsidio para mayores de 52 años, como mecanismo de ajuste de las plantillas en las reconversiones industriales y bancarias.

El Ageing Report 2024 proyecta que la edad efectiva media de salida pasaría de 64,0 años en 2022 a 66,4 en 2070, un aumento de 2,4 años que atribuye en buena medida a los incentivos introducidos en 2021: el endurecimiento de los coeficientes reductores de la jubilación anticipada, que pasan a ser mensuales y dependen de la carrera, y la mejora de la bonificación por demora, del 4% anual o el pago único. Proyecta también que la tasa de participación de 55 a 64 años subiría del 65,4% al 77,5%, y la tasa de empleo de 20 a 64 años del 69,6% al 76,4%. Son hipótesis optimistas respecto a la experiencia histórica española y dependen de que el comportamiento responda a los incentivos con la elasticidad estimada. Si no ocurre, el efecto amortiguador del factor empleo será menor. (OJO: la evaluación de si estas medidas funcionan corresponde a 4.3; aquí basta dejar claro que el factor empleo proyectado es condicional a los incentivos.)

d) La participación femenina: un factor que se agota.

La incorporación masiva de la mujer al mercado de trabajo desde los años ochenta ha sido una de las principales fuentes de crecimiento del empleo en España y, por tanto, de alivio del factor empleo. Pero es un efecto de nivel, no de tasa: una vez que la participación femenina converge con la masculina, deja de aportar. Además, tiene la contrapartida señalada en el apartado anterior: las mujeres que se incorporaron en los ochenta y noventa se jubilan ahora con carreras más largas y bases más altas que sus predecesoras, que a menudo solo accedían a pensiones de viudedad o a pensiones mínimas. Esto eleva la tasa de sustitución media, es decir, el tercer factor. El Ageing Report 2024 proyecta aumentos de la carrera media de cotización mucho mayores para las mujeres (casi ocho años) que para los hombres (menos de tres). El efecto es positivo para la suficiencia y la equidad de género (reduce la brecha de género en pensiones, que vimos en 3.2), pero negativo, en términos contables, para el gasto.

e) La ratio de afiliados por pensión: un indicador engañoso.

En 2026 la Seguridad Social informaba de una ratio de unas 2,35 personas afiliadas por pensión, la más alta en casi veinte años, y el dato se ha utilizado en el debate público como prueba de que el sistema "mejora". El indicador es engañoso por tres razones. Primera: las afiliaciones no son personas, porque una misma persona puede tener varias, ni las pensiones son pensionistas, porque hay concurrencia. Segunda: la ratio mide solo el primer y el segundo factor de la identidad, e ignora el tercero. El gasto puede crecer más que las cotizaciones aunque la ratio mejore, si la pensión media crece más que el salario medio, y eso es exactamente lo que ha ocurrido en varios años recientes, como veremos. Tercera: es un indicador del presente, sensible al ciclo, que no dice nada de la tendencia demográfica que llegará en 2030-2050. (OJO: es un buen ejemplo para mostrar en el examen el uso crítico de los datos. La mejora de la afiliación es real y positiva, pero no es un indicador de sostenibilidad a largo plazo.)

\subsubsection{Factor de generosidad: la tasa de sustitución y el efecto sustitución}

[DOC] El tercer factor es la pensión media sobre el salario medio, que el documento llama "factor económico o de cuantía" y describe como la tasa de sustitución agregada: cuanto mayor es, más generoso es el sistema y mayor el gasto. Según su descomposición, más del 60% del aumento del gasto entre 2008 y 2021 procede de este factor. Es la cifra más llamativa del diagnóstico, y también la que más matices necesita.

a) Precisión conceptual: tasa de sustitución frente a tasa de reemplazo.

La literatura y los organismos internacionales distinguen dos conceptos que en el debate español se confunden a menudo.

La tasa de sustitución agregada, o ratio de prestación (benefit ratio, en la terminología del AWG), es una ratio de stocks: la pensión media de todos los pensionistas, dividida por el salario medio de todos los trabajadores en un momento dado. Mide la posición relativa de los pensionistas en conjunto frente a los ocupados en conjunto, y es la que aparece en la identidad.

La tasa de reemplazo, en cambio, es una ratio de flujos o individual: la primera pensión de un trabajador dividida por su último salario, o por su salario medio de la carrera, según la definición. Mide cuánto renta pierde una persona al jubilarse. La OCDE, en Pensions at a Glance, calcula tasas de reemplazo teóricas para trabajadores tipo con carreras completas. El AWG calcula tasas de reemplazo "en el momento de la jubilación" (replacement rate at retirement) con datos reales.

Las dos se mueven de forma distinta. Un sistema que indexa las pensiones en curso a los precios, y no a los salarios, puede tener una tasa de reemplazo alta en el momento de la jubilación y una tasa de sustitución agregada decreciente, porque la pensión de cada jubilado pierde terreno frente a los salarios año tras año. El Ageing Report 2024 proyecta para España que la ratio de prestación total caería del 64% en 2022 al 51% en 2070, y la tasa de reemplazo en el momento de la jubilación del 77% al 64%. Ambas caen, pero la brecha entre ellas refleja el deterioro relativo de las pensiones a lo largo de la jubilación. (OJO: la comparativa internacional de las tasas de reemplazo de la OCDE, en la que España figura entre las más altas, corresponde a 4.5.)

b) El efecto sustitución: por qué la pensión media sube más que la revalorización.

[DOC] La nota 15 del documento define el efecto sustitución: el gasto en pensiones aumenta porque los trabajadores que se jubilan lo hacen con pensiones más altas que las de los pensionistas que fallecen. Es un concepto fundamental y poco conocido fuera del ámbito especializado. La pensión media del sistema crece cada año por dos vías. La primera es la revalorización: la actualización de las pensiones en curso, hoy según el IPC medio (Ley 21/2021). La segunda es el efecto composición, o efecto sustitución: las bajas del sistema (fallecimientos, sobre todo de cohortes con carreras cortas, bases bajas y, en el caso de las mujeres, a menudo solo con pensión de viudedad) son sustituidas por altas de cohortes con carreras más largas, salarios más altos y mayor presencia femenina en pensiones propias. En España, el efecto sustitución ha añadido de manera recurrente del orden de un punto porcentual o algo más al crecimiento anual de la pensión media, por encima de la revalorización. La pensión media de las nuevas altas de jubilación del Régimen General supera con claridad la pensión media del conjunto del sistema.

El efecto sustitución tiene dos lecturas.

La lectura ortodoxa lo considera una presión estructural sobre el gasto que ningún índice de revalorización neutraliza, porque no depende de la actualización de las pensiones existentes, sino de las nuevas. Por eso el índice de revalorización de 2013 lo incorporaba explícitamente en su fórmula como un término que restaba de la revalorización. Como el efecto sustitución era siempre positivo, empujaba la revalorización hacia el suelo del 0,25%: la fórmula obligaba a los pensionistas existentes a "pagar" la mayor generosidad de los nuevos. Lo veremos en 4.3.

La lectura crítica subraya que el efecto sustitución no es una "generosidad" del sistema, sino la maduración de un sistema contributivo en una sociedad cuyas carreras laborales se han alargado y cuyos salarios reales han crecido a largo plazo. Llamarlo generosidad sugiere que el sistema paga más de lo debido, cuando en realidad paga lo que corresponde a lo cotizado según la fórmula vigente. Es la posición de buena parte de la literatura heterodoxa y del discurso sindical. Ambas lecturas son compatibles: el efecto sustitución es a la vez la maduración legítima del sistema y una presión real sobre el gasto. La cuestión normativa es si la fórmula vigente, que determina la pensión inicial, es actuarialmente equilibrada. Ahí remitimos al debate sobre la contributividad y la TIR del sistema (4.1.3).

c) Por qué subió la tasa de sustitución entre 2008 y 2021: el problema del denominador.

El documento atribuye a la mayor tasa de sustitución más del 60% del aumento del gasto entre 2008 y 2021. La lectura inmediata, "el sistema se volvió más generoso", es incompleta. Una ratio puede subir porque crece el numerador o porque cae el denominador. En este período ocurrieron ambas cosas.

Por el lado del numerador, operaron el efecto sustitución, las subidas de las pensiones mínimas por encima de la revalorización general y, en 2008-2011, revalorizaciones con inflación baja o negativa, que mejoraban el poder adquisitivo de las pensiones.

Por el lado del denominador, el salario medio se estancó en términos reales durante la crisis y la devaluación interna de 2010-2014. Hubo, además, un efecto composición del empleo en sentido contrario al que se esperaría: la destrucción de empleo se concentró en la construcción y en los trabajos de baja cualificación, lo que elevó el salario medio en 2008-2009, pero la recuperación posterior se produjo con empleo de salarios más bajos y más parcialidad.

En 2022-2023 se produjo un episodio muy ilustrativo. Con la inflación de la crisis energética, las pensiones se revalorizaron un 8,5% en 2023, según el IPC medio de diciembre de 2021 a noviembre de 2022, mientras los salarios pactados en convenio crecieron bastante menos. La tasa de sustitución subió de golpe, no porque el sistema cambiara, sino porque la indexación al IPC protegió a los pensionistas de una pérdida de poder adquisitivo que sí sufrieron los asalariados.

La conclusión es relevante para la evaluación. Con salarios reales estancados, la indexación al IPC no reduce la tasa de sustitución: la mantiene o la eleva. La caída de la tasa de sustitución que proyecta el Ageing Report 2024, del 64% al 51%, depende de que los salarios reales crezcan de forma sostenida, es decir, de la productividad, el cuarto factor. Si la productividad no crece, la tasa de sustitución no cae y el gasto sube más de lo proyectado.

d) El efecto tope y la "reforma silenciosa".

El Ageing Report 2024 identifica otro mecanismo de caída de la tasa de sustitución: el efecto tope (cap effect). Una proporción creciente de pensiones de jubilación alcanzará la pensión máxima, porque las bases máximas de cotización crecen más deprisa que la pensión máxima. Enlaza con la reforma silenciosa de Conde-Ruiz y González (2016, SERIEs) que ya vimos en 3.2. Si la base máxima de cotización se actualiza más deprisa que la pensión máxima, el sistema cobra cada vez más a los salarios altos sin aumentar proporcionalmente sus prestaciones, y se desliza de un modelo bismarckiano, de pensión proporcional al salario, hacia uno más beveridgeano, de pensión que tiende a una cantidad uniforme en la parte alta. El RDL 2/2023 lo formaliza al fijar una subida de la base máxima del IPC más 1,2 puntos anuales hasta 2050, más una cuota de solidaridad sobre los salarios por encima de la base máxima, mientras la pensión máxima crece con el IPC más un incremento adicional mucho menor (0,115 puntos acumulativos anuales hasta 2050, según la disposición correspondiente). (Pendiente de verificar la cifra exacta del incremento adicional de la pensión máxima.) El resultado es una reducción de la contributividad en la parte alta de la distribución.

El debate es de fondo. Los defensores de la progresividad lo consideran una mejora de la equidad vertical: quien más gana contribuye más a la solidaridad del sistema. Sus críticos (Conde-Ruiz, Devesa, De la Fuente) advierten tres riesgos: que se rompa la legitimidad contributiva entre los trabajadores de salarios altos, que tienen más capacidad para salir del sistema o reducir su base declarada; que el cambio de modelo se produzca sin debate explícito; y que se incentive la elusión, por ejemplo con retribuciones en forma de dividendos o mediante sociedades interpuestas. Es exactamente el mecanismo de Korpi y Palme (1998): si las clases medias-altas dejan de percibir el sistema como propio, su apoyo político se erosiona y, con él, la suficiencia de todo el sistema.

e) Generosidad y suficiencia: la tensión de fondo.

Toda reducción de la tasa de sustitución es un ajuste en el eje de la suficiencia. El triángulo de hierro de 4.1.2 lo hacía evidente: si la tasa de sustitución es la variable de ajuste, se mantiene la sostenibilidad a costa de la posición relativa de los pensionistas. La cuestión normativa es si una ratio de prestación del 51% en 2070 es compatible con el mandato constitucional de "pensiones adecuadas y periódicamente actualizadas" (art. 50 CE). Hay dos elementos que la hacen más aceptable de lo que parece. En primer lugar, la indexación al IPC garantiza que la pérdida sea relativa y no absoluta. En segundo lugar, la tasa de reemplazo en el momento de la jubilación, del 64%, sigue siendo elevada en términos comparados. Hay dos que la hacen más problemática: la caída se concentra en los pensionistas de mayor edad, cuya pensión lleva más años indexada a los precios y que son a la vez los más vulnerables (sobre todo las viudas), y la protección real depende de que la productividad crezca. Cómo evaluar esta tensión es el objeto de 4.3 y 4.4.

\subsubsection{Participación salarial y la productividad}

[DOC] El cuarto factor es la masa salarial sobre el PIB, es decir, la participación de los salarios en la renta. El documento señala que, dado el salario medio, cuanto menor sea la productividad, mayor será el gasto, y que la caída de la participación salarial entre 2008 y 2021 compensó en parte el aumento del gasto. Es el factor menos tratado y, desde el punto de vista teórico, el más interesante, porque en él se cruzan tres debates: el de la distribución funcional de la renta, el de la productividad española y el de la base de financiación del sistema.

a) El papel ambiguo de la participación salarial.

En la identidad del gasto, una menor participación salarial reduce el gasto sobre el PIB, porque, con la misma masa de pensiones, el PIB es mayor en relación con los salarios. Esto explica que el documento diga que la caída de la participación salarial "compensó" el aumento. Pero esta lectura es engañosa si se aísla del lado de los ingresos. Como vimos en la introducción, en un sistema financiado con cotizaciones sobre salarios, la participación salarial multiplica a la vez los ingresos y los gastos. Una caída de la participación salarial reduce el gasto sobre el PIB, pero reduce en la misma proporción las cotizaciones sobre el PIB, y no mejora el saldo. En realidad lo empeora si el sistema tiene gastos que no dependen de los salarios o si la revalorización de las pensiones sigue a los precios y no a los salarios. La "compensación" que señala el documento es un artefacto de medir el gasto sobre el PIB en lugar de sobre la base de cotización.

En la práctica, los ejercicios de proyección del AWG suponen una participación salarial constante a largo plazo, de modo que este factor desaparece de las proyecciones por construcción. Es una hipótesis cómoda pero discutible, a la vista de la literatura reciente.

b) El declive de la participación salarial: una tendencia global con implicaciones fiscales.

Desde los años ochenta, la participación del trabajo en la renta ha caído en la mayoría de las economías avanzadas. Karabarbounis y Neiman (2014, Quarterly Journal of Economics) la atribuyeron al abaratamiento relativo de los bienes de inversión, que incentiva sustituir trabajo por capital. Autor, Dorn, Katz, Patterson y Van Reenen (2020, Quarterly Journal of Economics) la vincularon a la aparición de empresas superestrella con márgenes elevados y participaciones salariales bajas. Piketty (2014) la enmarcó en una dinámica de largo plazo en la que el rendimiento del capital supera el crecimiento. Acemoglu y Restrepo (2019, 2020) documentaron el efecto desplazamiento de la automatización. En España la participación salarial cayó durante la crisis de 2008-2014 y la devaluación interna, y se ha recuperado solo parcialmente.

La implicación para las pensiones es directa, y enlaza con lo visto en 2.2.3 sobre la estructura de financiación. Si la participación salarial cae de forma persistente, un sistema financiado exclusivamente con cotizaciones sobre salarios ve reducirse su base de financiación en relación con la renta nacional, aunque el PIB crezca. De ahí proceden varias propuestas de ampliar la base de financiación a rentas no salariales: la financiación con imposición general (el modelo beveridgeano), los impuestos sobre el valor añadido afectados, al modo de la Contribution Sociale Généralisée francesa, que grava todas las rentas, o, en su versión más mediática, los impuestos sobre los robots, propuestos por Bill Gates en 2017 y analizados formalmente por Guerreiro, Rebelo y Teles (2022, Review of Economic Studies), que concluyen que un impuesto moderado sobre la automatización puede ser óptimo de forma transitoria, pero no a largo plazo. Abbott y Bogenschneider (2018) sostienen, desde el derecho tributario, que los sistemas fiscales favorecen la automatización al gravar más el trabajo que el capital. Frente a estas propuestas, la ortodoxia recuerda que la incidencia económica de las cotizaciones recae en gran medida sobre los trabajadores (Gruber, 1997; lo vimos en 2.2), de modo que el cambio de base altera la distribución, pero no crea recursos nuevos, y que gravar el capital en una economía abierta tiene costes de eficiencia y de fuga de base. El debate queda abierto y será uno de los ejes de 4.4.

c) La productividad: el factor más poderoso y el más dependiente del régimen de indexación.

[DOC] La nota 23 del documento permite reescribir los dos últimos factores como pensión media sobre productividad por ocupado. Esta formulación revela que la productividad es, en principio, el factor más poderoso de todos: si la producción por ocupado crece deprisa, la carga de las pensiones sobre el PIB se reduce aunque la demografía empeore. Es la versión contable de la tesis de Barr: con suficiente crecimiento de la producción, cualquier ratio de dependencia es soportable.

Pero la potencia de la productividad depende crucialmente del régimen de indexación de las pensiones, y esta es la idea más importante del epígrafe. Si las pensiones iniciales se calculan sobre los salarios de la carrera y las pensiones en curso se indexan a los salarios, un aumento de la productividad eleva por igual salarios, pensiones iniciales y pensiones en curso, y la ratio pensión/productividad no cambia. La productividad no alivia el sistema: lo hace más rico, pero no más sostenible. Si, en cambio, las pensiones en curso se indexan a los precios, como en España desde la Ley 21/2021, un aumento de la productividad eleva los salarios y el PIB, mientras las pensiones en curso solo mantienen su poder adquisitivo. La ratio cae y el sistema sí gana sostenibilidad, pero a costa de que los pensionistas no participen de las ganancias de productividad. En ese régimen, la productividad alivia el gasto precisamente porque reduce la tasa de sustitución. Es el mismo ajuste por la vía de la suficiencia relativa, solo que menos visible.

Así se entiende que las proyecciones del Ageing Report 2024 dependan tanto de la hipótesis de productividad. La caída de la ratio de prestación, que contribuye con unos −4,1 puntos del PIB al cambio del gasto entre 2022 y 2070, es, en gran parte, el resultado de suponer que los salarios reales crecen con la productividad (en torno al 1,5% anual a largo plazo, según la metodología común del AWG) mientras las pensiones en curso se indexan a los precios.

d) El problema español de la productividad.

Aquí aparece otra verdad incómoda. La productividad total de los factores española se ha estancado, o ha crecido muy poco, desde mediados de los noventa, según los trabajos de la Fundación BBVA y el Ivie (Mas, Pérez y sus colaboradores), del Banco de España y de la Comisión Europea. Las causas son bien conocidas: el pequeño tamaño medio de las empresas, la especialización en sectores de baja productividad (construcción, hostelería, comercio minorista), la mala asignación del capital hacia sectores protegidos durante la burbuja (Gopinath, Kalemli-Özcan, Karabarbounis y Villegas-Sánchez, 2017, Quarterly Journal of Economics), el bajo esfuerzo en I+D+i, la dualidad del mercado de trabajo y la escasa formación en el empleo. Si el crecimiento de la productividad a largo plazo se quedara por debajo de la hipótesis del AWG, la caída proyectada de la tasa de sustitución no se produciría y el gasto se acercaría al del escenario de riesgo que el propio Ageing Report publica como análisis de sensibilidad. La sostenibilidad del sistema español descansa, en buena medida, en una hipótesis de productividad que la economía española no ha cumplido en los últimos treinta años.


\begin{specialblock}{\textbf{Evidencia empírica:} Período 2008-2021 y las proyecciones del Ageing Report 2024}
    Con los cuatro factores definidos, podemos leer la evidencia agregada. Primero la del documento para el período 2008-2021, después las proyecciones más recientes.

    \textbf{Descomposición 2008-2021 del documento y su sesgo cíclico}. El gasto en pensiones contributivas, excluidos los complementos a mínimos, pasó del 7,1% del PIB en 2008 al 11,5% en 2021, un aumento de 4,4 puntos. Tomando logaritmos en la identidad, de modo que la variación del logaritmo del gasto se reparte aditivamente entre las variaciones de los logaritmos de cada factor, algo más del 35% del aumento procede del factor demográfico, en torno al 10% del empeoramiento del empleo y más del 60% de la mayor tasa de sustitución, compensado en parte por la menor participación salarial.
    
    Tres observaciones críticas son imprescindibles.
    
    La primera es el sesgo de los años de referencia. 2008 fue el último año de la expansión, con un PIB en máximos cíclicos y un empleo en niveles récord. 2021 fue el año posterior a la caída del PIB de 2020, la mayor en tiempos de paz (cerca del 11%), con un PIB aún por debajo de su nivel previo a la pandemia. La ratio del 11,5% está inflada por un denominador deprimido. Tomando como año final 2024 o 2025, con el PIB ya recuperado y la afiliación en máximos, el aumento del gasto contributivo sobre el PIB sería menor, y el peso relativo de los factores cambiaría: menos empleo, más demografía y generosidad. (OJO: en el examen, citar el dato del documento advirtiendo el sesgo cíclico y con cifras más recientes.)
    
    La segunda es la del método. La descomposición logarítmica atribuye a cada factor una parte del cambio suponiendo que los demás están fijos. Como vimos en la introducción, los factores no son independientes, y el reparto no puede leerse como causal. Que "el 60% se deba a la generosidad" significa, en rigor, que la ratio pensión media/salario medio aumentó en una magnitud que, a igualdad de lo demás, explicaría ese porcentaje, no que un cambio normativo de generosidad lo haya provocado. Buena parte de ese aumento fue, como vimos en 4.2.3, efecto sustitución y estancamiento salarial.
    
    La tercera es la del perímetro. Las cifras del documento (gasto contributivo sin complementos a mínimos) no son comparables con las del Ageing Report, que mide el gasto bruto en pensiones públicas: incluye las Clases Pasivas, las pensiones no contributivas, los complementos a mínimos y otras prestaciones asimiladas. Por eso el Ageing Report 2024 sitúa el gasto de 2022 en el 13,1% del PIB, bastante por encima del 11,5% del documento para 2021. (OJO: es un error frecuente mezclar perímetros. Hay que decir siempre de qué gasto se habla.)
    
    b) Los saldos y el Fondo de Reserva.
    
    [DOC] En el mismo período, las cotizaciones sociales aumentaron solo del 9,7% al 10,9% del PIB, y el saldo no financiero del sistema pasó de un superávit del 1,4% del PIB en 2008 a un déficit del 0,9% en 2021. El Fondo de Reserva, que alcanzó su máximo en 2011, se redujo hasta el 0,2% del PIB en 2021, el equivalente a apenas el 1,5% del gasto anual en pensiones contributivas. Conviene precisar, como hicimos en 2.2.3, que el máximo fue de unos 66.800 millones de euros, en torno al 6,2% del PIB, no el 6,3% que indica el documento, y que su agotamiento se debió a las disposiciones de 2012-2017 para pagar las pagas extraordinarias. Desde 2023, con las dotaciones del MEI, el Fondo ha empezado a recuperarse: según el Ministerio de Inclusión, Seguridad Social y Migraciones, cerró 2025 por encima de 14.000 millones y superó los 15.200 millones en abril de 2026, máximo de la última década, aunque todavía por debajo del 1% del PIB y muy lejos de su máximo histórico. (Remisión: la estructura de financiación y las transferencias del Estado que "cubren" parte del déficit, a 2.2.3; el diseño y la evaluación del MEI, a 4.3.)
    
    c) Las proyecciones del Ageing Report 2024: la descomposición prospectiva.
    
    [Probable en todas las cifras, a partir de la ficha de España del Ageing Report 2024] El gasto bruto en pensiones pasaría del 13,1% del PIB en 2022 al 14,3% en 2030, al 16,2% en 2040 y a un máximo del 17,3% en 2050, para descender al 16,7% en 2070: un aumento de 3,6 puntos en todo el horizonte. La descomposición de ese aumento es la síntesis de todo el apartado:
    
    - La tasa de dependencia aportaría en torno a +10,3 puntos. Es la presión demográfica bruta: si nada más cambiara, el gasto casi se duplicaría.
    - La ratio de prestación restaría unos −4,1 puntos. Es el ajuste por la generosidad, con la indexación a los precios, el crecimiento de los salarios reales y el efecto tope.
    - El mercado de trabajo (tasa de empleo e intensidad laboral) restaría en torno a −1,9 puntos.
    - La tasa de cobertura restaría alrededor de −0,4 puntos, por el retraso de la edad efectiva de salida.
    
    Por el lado de los ingresos, las cotizaciones públicas para pensiones, en la definición del AWG, pasarían del 12,9% del PIB al 14,6% en 2050 y al 14,0% en 2070. El AWG proyecta, a política constante, un déficit del sistema de pensiones cercano al 2,7% del PIB en 2070. (Pendiente de verificar: la definición exacta de ingresos del AWG incluye transferencias. La cifra de déficit debe citarse con esa advertencia.)
    
    Lectura del conjunto. La descomposición prospectiva confirma lo que hemos ido desgranando en cada epígrafe. La demografía empuja con una fuerza que ningún otro factor compensa por sí solo. El empleo y la edad de jubilación amortiguan, pero menos de lo que suele creerse. El grueso del ajuste proyectado recae sobre la tasa de sustitución, es decir, sobre la posición relativa de los pensionistas, y depende de una hipótesis de productividad exigente. Dicho de otro modo: aun después de las reformas de 2021-2023, que eliminaron el factor de sostenibilidad y el índice de revalorización, el sistema español sigue confiando buena parte de su sostenibilidad a un deterioro relativo de las pensiones respecto a los salarios. Ya no procede de una regla explícita, sino de la combinación de la indexación al IPC, el crecimiento de la productividad y el efecto tope. La diferencia entre las reformas de 2013 y de 2021-2023 no está tanto en si cae la tasa de sustitución, sino en cuánto, a través de qué mecanismo y con qué grado de visibilidad. Esta es la pregunta con la que entramos en 4.3.
\end{specialblock}




\subsection{Reformas como respuestas: de 2011 a la cláusula de salvaguarda}
Una verdad incómoda antes de empezar, porque cambia la forma de presentar este apartado. La cláusula de salvaguarda que el documento describe como garantía de cierre del sistema ya se aplicó en marzo de 2025. Se "cumplió" por un margen de una décima de PIB, y solo después de que el Gobierno aprobara, un mes antes, un real decreto de metodología que incluía en el cálculo transferencias del Estado y una revisión al alza del PIB. La propia AIReF concluyó que la regla se cumplía formalmente pero no garantizaba la sostenibilidad, y lo ratificó en mayo de 2026. Si en el examen se presenta la cláusula como algo por estrenar, el tema queda desactualizado en su punto más sensible. He redactado 4.3 incorporando ese episodio. También he corregido una errata del documento que afecta a la cronología: el "2011 Ageing Report" no existe, es el de 2012.

El apartado anterior dejó planteada una pregunta: si la demografía empuja el gasto hacia arriba con una fuerza que ningún otro factor compensa por sí solo, sobre quién recae el ajuste. Toda reforma de pensiones es, en el fondo, una respuesta a esa pregunta. La identidad de 4.2 ofrece los cuatro lugares donde puede buscarse. El factor demográfico, que en la variante del AWG es la tasa de cobertura, se mueve con la edad de jubilación y el acceso a la pensión. El factor empleo, con los incentivos a trabajar más tiempo. La generosidad, con la fórmula de cálculo, la revalorización o los mecanismos automáticos ligados a la longevidad. Y los ingresos, con los tipos, las bases o las transferencias del Estado. Cada opción traslada la carga a un grupo distinto: a los futuros jubilados, a los pensionistas actuales, a los cotizantes, a los salarios altos o a los contribuyentes en general.

En los últimos quince años, España ha recorrido esas opciones en secuencia, y en direcciones a veces opuestas. La reforma de 2011 actuó sobre la edad y la fórmula de cálculo: sobre los futuros jubilados. La de 2013 introdujo dos mecanismos automáticos que hacían recaer el ajuste en la tasa de sustitución: sobre los pensionistas, presentes y futuros. La de 2021-2023 desmontó esos mecanismos y desplazó el ajuste hacia los ingresos: sobre los cotizantes, los salarios altos y el Estado. Además, añadió una cláusula de cierre cuyo mecanismo por defecto vuelve a ser una subida de cotizaciones.

Tres precisiones ordenan el apartado.

La primera es de reparto con el bloque III. La regulación vigente de cada elemento (edad, base reguladora, coeficientes, revalorización, MEI, base máxima y cuota de solidaridad) ya se ha expuesto allí, con notas sobre su evolución. Aquí no se repite la norma: se evalúa su efecto.

La segunda es el hilo conductor. Utilizamos los sucesivos Ageing Reports de la Comisión Europea (2009, 2012, 2015, 2018, 2021 y 2024) como hilo, porque cada uno es, en esencia, la evaluación ex ante de la reforma que acababa de aprobarse. [DOC] La tabla 5 del documento recoge esa secuencia. (OJO: el documento habla de un "2011 Ageing Report" que proyectaría un aumento de 3,6 puntos. Los informes son trienales y no hubo informe en 2011: es el 2012 Ageing Report el que recoge por primera vez los efectos de la reforma de 2011, como reconoce el propio documento unas líneas antes. Conviene citarlo bien.)

La tercera es un marco analítico de la economía política de las reformas que ya presentamos en el bloque II y en 4.1: la evitación de la culpa de Weaver (1986), la dependencia de la trayectoria y la "nueva política del Estado de bienestar" de Pierson (1994, 1996, 2001), y el dilema entre reglas y discrecionalidad de Kydland y Prescott (1977). La tesis que se sostiene a lo largo del apartado es la siguiente. En España, la duración de una reforma de pensiones ha dependido más de su grado de consenso que de su calidad técnica. La de 2011, pactada, sigue en vigor. La de 2013, técnicamente la más ambiciosa, se aprobó sin consenso y fue revocada en ocho años. La de 2021-2023, aprobada con apoyo de los agentes sociales, pero sin el del principal partido de la oposición en su segunda fase, está sometida a una prueba de credibilidad cuyo primer episodio fue la aplicación de la cláusula en 2025.

\subsubsection{Punto de partida y la reforma de 2011: el ajuste por la edad y la fórmula}

a) El diagnóstico del 2009 Ageing Report.

[DOC] Según el Ageing Report de 2009, el sistema español era difícilmente sostenible antes de la reforma de 2011: a largo plazo proyectaba un aumento del gasto en pensiones de unos 6,7 puntos del PIB respecto al año base. Era uno de los mayores aumentos proyectados en la UE y colocaba a España en el grupo de países de "alto riesgo" en la evaluación de sostenibilidad a largo plazo de la Comisión. El informe se elaboró antes de que la Gran Recesión se hiciera sentir en el mercado de trabajo español, y con proyecciones demográficas que suponían la continuidad de la inmigración de los años 2000. La situación de partida se agravó rápidamente cuando el empleo se desplomó y el superávit del sistema, del 1,4% del PIB en 2008, desapareció en tres años. (Remisión: el contexto político de 2010-2011, la crisis de deuda soberana, la carta del BCE y el giro de mayo de 2010, en 1.2.2.)

b) Qué hizo la reforma y sobre qué factor actuó.

La Ley 27/2011, de 1 de agosto, sobre actualización, adecuación y modernización del sistema de Seguridad Social, fue la traducción legal del Acuerdo Social y Económico de 2 de febrero de 2011, firmado por el Gobierno, CCOO, UGT, CEOE y CEPYME. Previamente, la Comisión del Pacto de Toledo había aprobado su informe de evaluación y reforma de enero de 2011. Es la última gran reforma de pensiones española que combinó el consenso parlamentario en la Comisión del Pacto de Toledo y el pacto con los agentes sociales. Sus tres piezas centrales fueron:

- el aumento gradual de la edad legal de jubilación de 65 a 67 años hasta 2027, con la excepción de quienes acrediten 38 años y 6 meses de cotización, que pueden seguir jubilándose a los 65;
- la ampliación del período de cálculo de la base reguladora de 15 a 25 años hasta 2022;
- el aumento de los años necesarios para alcanzar el 100% de la base reguladora, de 35 a 37 en 2027.

(OJO: el documento dice "35 a 37 años en 2027", que es correcto. La cifra errónea de "36 años (37 a partir de 2017)" que aparecía en otra parte del documento ya se corrigió en 3.2.) Complementariamente, el Real Decreto-ley 5/2013 endureció las condiciones de la jubilación anticipada y parcial, alineándolas con la nueva edad legal. (Regulación vigente y detalle de la transición, en 3.2.)

En el lenguaje de 4.2, la reforma actuó sobre dos factores. El aumento de la edad reduce la tasa de cobertura, porque hay menos pensiones por persona mayor de 65, y aumenta la tasa de empleo, porque hay más ocupados de 65 y 66 años. La ampliación del período de cálculo y los años exigidos para el 100% reducen la pensión inicial y, por tanto, la tasa de sustitución de las nuevas cohortes. Ampliar el período de cálculo reduce la pensión porque, en carreras con salarios reales crecientes, los años más lejanos tienen bases más bajas, incluso actualizadas por el IPC, y la media baja. El efecto es mayor para quienes tienen perfiles salariales crecientes (los trabajadores cualificados) y menor, o incluso favorable, para quienes sufren una caída salarial o un desempleo al final de la carrera. Esta ambivalencia distributiva estuvo en el centro del debate de 2011 y ha vuelto a estarlo en 2023 con el sistema dual de cómputo.

c) La evaluación: un ahorro de entre un tercio y dos quintos.

[DOC] Las estimaciones disponibles coinciden en que la reforma de 2011 contrarrestó en parte las presiones al alza, pero no garantizaba la sostenibilidad si se confirmaban las proyecciones demográficas. El ahorro estimado se situaba entre el 30% y el 40% del aumento de gasto previsto a largo plazo sin la reforma:

- el Gobierno, en la Actualización del Programa de Estabilidad 2012-2015, en torno al 40% del aumento proyectado entre 2010 y 2050;
- De la Fuente y Doménech (2013), un ahorro del 33% en 2050;
- Conde-Ruiz y González (2013), un 29%;
- el Banco de España (2011), un 43%.

En el 2012 Ageing Report, el aumento proyectado del gasto se redujo a unos 3,6 puntos del PIB sobre el año base, frente a los 6,7 del informe de 2009.

Tres observaciones críticas matizan esa valoración.

La primera: la reforma era lenta. Su calendario de aplicación, de 2013 a 2027, hacía que el grueso del ahorro se produjera en las décadas de 2030 y 2040, cuando la presión ya sería máxima. Era una elección deliberada, que respetaba las expectativas de quienes estaban cerca de la jubilación, y es habitual en las reformas paramétricas pactadas. Pero dejaba sin resolver el déficit que la crisis había abierto a corto plazo.

La segunda: no tocaba ni la revalorización ni la longevidad. La ley introdujo el principio de un factor de sostenibilidad ligado a la esperanza de vida (art. 8 de la Ley 27/2011), pero lo aplazó a 2027 y lo remitió a un desarrollo posterior. Esa puerta es la que abrió la reforma de 2013, adelantando su aplicación.

La tercera: las excepciones erosionaban su alcance. La posibilidad de jubilarse a los 65 con 38,5 años cotizados, y los regímenes de jubilación anticipada por actividades penosas o peligrosas, reducían el efecto sobre la edad efectiva.

Conde-Ruiz y González (2013, Hacienda Pública Española) cuantificaron que la reforma dejaba sin cubrir la mayor parte del desequilibrio a largo plazo y propusieron, como otros autores de la época (Boldrin, Jiménez-Martín, la propuesta de "Reformas estructurales" de Fedea de 2010), avanzar hacia un sistema de cuentas nocionales al estilo sueco, en el que la sostenibilidad se garantiza por construcción. La propuesta no prosperó, pero reaparecerá en 4.4.

Desde la perspectiva de la economía política, la reforma de 2011 ofrece la primera lección. Es la única reforma de las tres que no ha sido revertida, ni siquiera parcialmente, en su núcleo: la edad de 67 años y el período de cálculo de 25 años siguen vigentes, y la reforma de 2023 amplió este último. Su legitimidad descansaba en el consenso con los agentes sociales, que asumieron el coste de aceptar la jubilación a los 67 a cambio de las excepciones por carrera larga y de mantener la revalorización. La prueba empírica de la tesis de Pierson sobre la "política de la austeridad permanente" es clara: en un sistema maduro, con millones de beneficiarios organizados, las reformas que perduran son las que reparten la culpa entre todos los actores con poder de veto.

\subsubsection{Reforma de 2013: los mecanismos automáticos y el ajuste por la generosidad}

a) De la edad a la regla: el cambio de lógica.

La Ley 23/2013, de 23 de diciembre, reguladora del factor de sostenibilidad y del índice de revalorización del sistema de pensiones de la Seguridad Social, supuso un cambio cualitativo. Las reformas anteriores eran paramétricas y discrecionales: el legislador fijaba nuevos valores de los parámetros y la sostenibilidad dependía de que la realidad se ajustara a las proyecciones. La de 2013 introdujo mecanismos de ajuste automático (automatic balancing mechanisms), que corrigen los parámetros de forma continua ante las desviaciones, sin necesidad de nuevas leyes. Seguía una tendencia internacional bien documentada: la OCDE (Pensions at a Glance, 2021, cap. 2) contaba en torno a dos tercios de sus miembros con algún mecanismo automático, ya fuera un vínculo entre la edad legal y la esperanza de vida, un factor de sostenibilidad en la pensión inicial o un mecanismo de equilibrio financiero. Los modelos de referencia eran el sistema sueco de cuentas nocionales con su "freno" (bromsen), el factor de sostenibilidad alemán de 2004 (Nachhaltigkeitsfaktor) y el factor de esperanza de vida finlandés. (Remisión: comparativa detallada en 4.5.)

La reforma tuvo un origen técnico explícito. El Gobierno encargó en 2013 un Comité de Expertos sobre el factor de sostenibilidad, presidido por el sociólogo Víctor Pérez-Díaz y formado por doce miembros. Su informe de junio de 2013 propuso tanto el factor de sostenibilidad intergeneracional como un factor de revalorización anual, con una única opinión discrepante (la de Santos Ruesga). El Gobierno adoptó el esquema básico con modificaciones, y fue aprobado por la mayoría absoluta del Partido Popular sin el acuerdo del Pacto de Toledo ni de los agentes sociales. Es un dato decisivo para entender lo que vino después.

b) El índice de revalorización: una regla de equilibrio presupuestario aplicada a los pensionistas.

[DOC] El índice de revalorización de las pensiones (IRP) determinaba la actualización anual de todas las pensiones contributivas, incluidas las mínimas. Lo hacía con una fórmula que recogía la evolución reciente y prevista de los ingresos y gastos del sistema y de los principales determinantes del gasto: el número de pensiones y el efecto sustitución. Por ley no podía ser inferior al 0,25% ni superior al IPC más el 0,50%.

La fórmula puede escribirse, de forma simplificada, así:

IRP = ḡ(I) − ḡ(P) − ḡ(s) + α · (I − G) / G

Aquí ḡ(I) es el crecimiento medio de los ingresos del sistema, ḡ(P) el del número de pensiones y ḡ(s) el del efecto sustitución, todos ellos en medias móviles de once años centradas en el año de referencia (cinco pasados, el corriente y cinco proyectados). El último término corrige una fracción α, entre 0,25 y 0,33, del desequilibrio medio entre ingresos (I) y gastos (G).

La lógica es transparente si se lee con la identidad de 4.2. Si los ingresos crecen al mismo ritmo que el número de pensiones y el efecto sustitución, y el sistema está en equilibrio, el IRP es cero: la pensión media crece solo por el efecto sustitución y el gasto crece al ritmo de los ingresos. Si el número de pensiones crece más deprisa que los ingresos (el factor demográfico), o si las nuevas pensiones son más altas que las que se extinguen (el efecto sustitución), o si hay déficit acumulado, el IRP cae. Es una regla de equilibrio presupuestario del sistema contributivo, cuya única variable de ajuste es la revalorización de las pensiones en curso. Los factores demográfico, de empleo y de sustitución se traducen en una menor actualización.

[DOC] En la práctica, dado el déficit del sistema y las proyecciones demográficas, el IRP implicaba que, sin ingresos adicionales, la revalorización permanecería en el suelo del 0,25% durante muchos años, lo que supondría una notable pérdida de poder adquisitivo de los pensionistas. La nota 26 del documento cuantifica el efecto: con una inflación del 2%, la pérdida anual de poder adquisitivo sería del 1,75%, y un pensionista con veinte años de jubilación vería reducido el poder adquisitivo de su pensión en torno a un 30% respecto al inicial. La cifra se comprueba fácilmente: 0,9825 elevado a 20 es aproximadamente 0,70. El documento añade una observación que conviene destacar, porque disuelve una aparente contradicción. La pensión media del sistema seguiría creciendo en términos nominales, e incluso reales durante algunos años, porque el efecto sustitución (las altas con pensiones más altas) compensaría con creces la pérdida de las pensiones en curso. El ajuste no se veía en la media: se veía en cada pensionista individual a lo largo de su jubilación.

c) El factor de sostenibilidad: la longevidad en la pensión inicial.

[DOC] El factor de sostenibilidad, cuya entrada en vigor se fijó inicialmente en 2019, ajustaba la pensión inicial de jubilación a la evolución de la esperanza de vida. Como esta tiende a aumentar, las generaciones futuras recibirían una pensión inicial menor, para una misma carrera, pero durante más tiempo, de modo que el importe total a lo largo de la vida sería similar. Con las proyecciones del INE 2016-2066, una persona que se jubilara en 2030 vería reducida su pensión inicial en torno a un 7% respecto a otra a la que no se aplicase el factor.

Técnicamente, el factor multiplicaba la pensión inicial por el cociente entre la esperanza de vida a los 67 años en un período de referencia y la de un período posterior, con revisiones quinquenales y elevado a 1/5 para suavizar el ajuste. Es un mecanismo de neutralidad actuarial parcial ante la longevidad: garantiza que el valor esperado de la pensión a lo largo de la vida no crezca por el mero hecho de que se viva más. Su justificación normativa es la equidad intergeneracional en el sentido estricto de 4.1.2: dos cohortes con la misma carrera reciben el mismo valor actual esperado de prestaciones. Su debilidad distributiva es la que vimos en 4.2.1. Se aplicaba con la esperanza de vida media, lo que castiga en términos relativos a los grupos de menor longevidad, que ven reducida su pensión anual sin disfrutar de los años adicionales que la justifican.

d) La evaluación ex ante: el gasto se congela, la suficiencia se desploma.

[DOC] Los Ageing Reports de 2015 y 2018, que ya incorporaban la reforma de 2013, proyectaban un gasto en pensiones aproximadamente constante en proporción del PIB durante todo el horizonte, pese al fuerte aumento de la tasa de dependencia. El mecanismo de ajuste era, según la tabla 5 del documento, una caída muy grande de la tasa de sustitución. La pensión media caería frente al salario medio porque este crecería con la productividad, mientras las pensiones en curso quedarían limitadas por la revalorización del 0,25%. [DOC] El 2018 Ageing Report llegaba a proyectar una revalorización del 0,25% durante prácticamente todo el horizonte. En esos informes, la ratio de prestación de España llegaba a caer por debajo del 40% hacia 2060-2070, una de las caídas más intensas de la UE.

[DOC] Los modelos de equilibrio general de generaciones solapadas calibrados para España llegaron a la misma conclusión. Es el caso de De la Fuente et al. (2017), Díaz-Giménez y Díaz-Saavedra (2017) y Sánchez (2014), del Banco de España. El factor de sostenibilidad y el IRP habrían avanzado sustancialmente hacia el equilibrio financiero a largo plazo, pero como el ajuste se producía principalmente a través de la pensión media, la suficiencia quedaba muy comprometida. El trabajo de Sánchez (2014) mostró además que, con el IRP, el ajuste era muy sensible a la inflación. Con inflación baja o nula, el suelo del 0,25% apenas suponía pérdida real. Con una inflación del 2%, la pérdida acumulada era muy severa. Esta observación es crucial para entender lo que ocurrió después.

e) La economía política del fracaso: por qué se revirtió.

La reforma de 2013 funcionó durante sus primeros años en la mejor coyuntura posible para su supervivencia política. Entre 2014 y 2016 la inflación fue nula o negativa, por lo que la revalorización del 0,25% supuso ganancias de poder adquisitivo. Los pensionistas no notaron el ajuste. Cuando la inflación repuntó en 2017, en torno al 2% en media anual, la pérdida de poder adquisitivo se hizo visible, y en 2018 surgió un movimiento de pensionistas sin precedentes, con manifestaciones semanales en Bilbao y concentraciones masivas en toda España.

La respuesta política fue inmediata. El acuerdo presupuestario de 2018 entre el Gobierno del PP y el PNV fijó una revalorización del 1,6% para 2018 y 2019 (con subidas mayores para las mínimas) y aplazó la entrada en vigor del factor de sostenibilidad hasta 2023 como máximo. [DOC] El documento lo resume: como resultado de la contestación social, en 2018 se suspendió la aplicación del IRP y se pospuso el factor de sostenibilidad. La Ley 21/2021 derogó ambos mecanismos.

La lectura teórica de este fracaso es doble.

Desde la evitación de la culpa (Weaver, 1986), el IRP era un diseño ideal: el ajuste no requería votaciones anuales y se atribuía a una fórmula técnica. Pero la estrategia falla cuando el efecto se hace visible y se concentra en un colectivo numeroso, organizado y con alta participación electoral. Los pensionistas son, en España, más de nueve millones de votantes con tasas de participación muy superiores a las de los jóvenes. La literatura sobre la gerontocracia o la política de la edad (Lynch, 2006; Vanhuysse, 2013; Sinn y Uebelmesser, 2002, que estimaron que Alemania se convertiría en una "gerontocracia" en la que la mayoría votante tendría más de 50 años hacia 2016) predice exactamente esta dinámica: una regla automática es tan creíble como el poder político del grupo que la sufre le permita ser.

Desde el dilema entre reglas y discrecionalidad (Kydland y Prescott, 1977), la reforma de 2013 ilustra el problema de la inconsistencia temporal a la inversa. Una regla solo es creíble si es costoso abandonarla, y una regla aprobada por una sola fuerza política es fácil de abandonar cuando esa fuerza pierde la mayoría o cambia de estrategia. Los defensores de la reforma, entre ellos buena parte de la profesión académica (De la Fuente, Doménech, Conde-Ruiz, los informes del Banco de España y la Comisión), sostuvieron que el problema no era el diseño, sino la falta de ingresos adicionales que atenuaran el ajuste. Con más ingresos, el IRP habría permitido revalorizaciones mayores. Sus críticos (sindicatos, economistas heterodoxos como Zubiri o Navarro, y la mayor parte de la izquierda parlamentaria) respondieron que el sistema arrastraba un déficit artificial: durante años había financiado con cotizaciones gastos que no le correspondían (reducciones de cuotas, prestaciones familiares, tarifas planas), y la regla lo hacía pagar a los pensionistas. Este argumento, conocido como el de los gastos impropios o la separación de fuentes, tendría consecuencias legales en 2021.

\textcolor{red}{La reforma de 2013 es el ejemplo perfecto de por qué la sostenibilidad política es una dimensión autónoma, como anticipamos en 4.1.1. Un sistema financieramente sostenible sobre el papel dejó de serlo cuando su coste social se hizo visible.}

\subsubsection{Reforma de 2021-2023: el ajuste por los ingresos}

a) La lógica de la reforma: cambiar de vértice del triángulo.

La reforma de 2021-2023 se articuló en dos normas. La Ley 21/2021, de 28 de diciembre, de garantía del poder adquisitivo de las pensiones y otras medidas de refuerzo de la sostenibilidad financiera y social del sistema público de pensiones, surgió del acuerdo de julio de 2021 con los agentes sociales. El Real Decreto-ley 2/2023, de 16 de marzo, de medidas urgentes para la ampliación de derechos de los pensionistas, la reducción de la brecha de género y el establecimiento de un nuevo marco de sostenibilidad del sistema público de pensiones, se pactó con los sindicatos, pero no con la patronal, y se convalidó sin el apoyo del PP. [DOC] La reforma fue uno de los hitos comprometidos con la Comisión Europea en el Plan de Recuperación, Transformación y Resiliencia (PRTR). Es decir, su aprobación condicionaba el desembolso de fondos Next Generation EU. Esta condicionalidad europea es un elemento de economía política nuevo respecto a 2011 y 2013. La Comisión validó el diseño como suficiente para garantizar la sostenibilidad, y ese aval ha sido utilizado por el Gobierno como argumento de legitimidad técnica frente a las críticas.

Leída con el triángulo de hierro, la reforma supuso un cambio de vértice. Donde la de 2013 hacía recaer el ajuste casi íntegramente en la tasa de sustitución, la de 2021-2023:

- garantizó el poder adquisitivo de las pensiones en curso (revalorización con el IPC medio);
- eliminó el ajuste por longevidad en la pensión inicial (derogación del factor de sostenibilidad);
- elevó la generosidad en la parte baja, al vincular las pensiones mínimas al umbral de la pobreza, reforzar el complemento de brecha de género y mejorar la integración de lagunas;
- trasladó el grueso del ajuste a los ingresos: el MEI; la subida de la base máxima por encima del IPC; la cuota de solidaridad; el nuevo sistema de cotización de los autónomos por ingresos reales; y el aumento de las transferencias del Estado para financiar los "gastos impropios", que reconoció legalmente la tesis de la separación de fuentes;
- introdujo incentivos al retraso de la jubilación (coeficientes reductores mensuales y aplicados sobre la pensión máxima, bonificación por demora del 4% anual o pago único);
- amplió el período de cálculo, en su versión dual, a 29 años con exclusión de 24 meses, en la línea de 2011.

(Regulación detallada, en 3.2 y 2.2.3.)

La reforma respondía a un diagnóstico explícito del Ministerio de Inclusión, entonces dirigido por José Luis Escrivá, anterior presidente de la AIReF: el problema del sistema era transitorio, la joroba del baby boom, y debía resolverse con ingresos y un colchón financiero, no con recortes permanentes de prestaciones. Como vimos en 4.2.1, las proyecciones del Ageing Report 2024 cuestionan esa premisa: la tasa de dependencia no vuelve a bajar después de 2050.

b) El coste de la reversión: qué suponía desmontar 2013.

[DOC] Antes incluso de que la reforma se aprobara, la ficha de España del Ageing Report 2021 incluyó estimaciones del coste de revalorizar con el IPC y de no aplicar el factor de sostenibilidad. Ambas medidas elevarían el gasto en pensiones sobre el PIB en 2060 en unos 2,5 y 1 puntos, respectivamente. Las estimaciones de la AIReF (2020) eran similares. De la Fuente et al. (2019) elevaban el coste de la revalorización con el IPC a unos 4 puntos del PIB en 2050. La AIReF (2023) hizo la estimación más completa del efecto neto de las reformas de 2021 y 2023, frente a un escenario con el IRP y el factor de sostenibilidad vigentes: un aumento del déficit de 1,1 puntos del PIB en 2050 y de 1 punto en 2070, que incluye un aumento del gasto de 2,7 puntos en 2050 y 2,5 en 2070 solo por la revalorización con el IPC (tabla 6 del documento).

La lectura es importante. Las medidas de ingresos compensan una parte sustancial del mayor gasto, pero no todo: según la AIReF, el coste de las medidas de gasto supera en torno a un punto del PIB el rendimiento de las medidas de ingresos a largo plazo. Las estimaciones del rendimiento de los ingresos divergen mucho. El Ministerio de Inclusión lo cifraba en torno al 1,7% del PIB de media en 2022-2050; la AIReF, en su evaluación de la primavera de 2023, en torno al 1,0%; el Ageing Report 2024, en torno al 1,2%. La diferencia no es menor: 0,7 puntos del PIB durante casi tres décadas es mucho dinero, y, como veremos, la cláusula de salvaguarda depende directamente de qué cifra se tome.

c) El Ageing Report 2024: la primera evaluación europea completa.

[Probable, a partir de la ficha de España] El Ageing Report 2024 fue el primero en incorporar la reforma completa. Proyectaba un gasto bruto del 13,1% del PIB en 2022, un máximo del 17,3% en torno a 2050 y un 16,7% en 2070. La media de 2022-2050 se situaba en torno al 15,1%, ligeramente por encima del umbral de referencia del 15% de la cláusula de salvaguarda. Comparado con las proyecciones de los informes de 2015 y 2018, que mostraban un gasto prácticamente plano, el cambio es de varios puntos del PIB a mediados de siglo, y la diferencia la explica casi por completo la reversión de 2013. La Comisión señaló que la reforma aumentaba el gasto proyectado sobre todo por la indexación a la inflación y la eliminación del factor de sostenibilidad, y que su sostenibilidad dependía del rendimiento efectivo de las medidas de ingresos.

d) El MEI: ¿mecanismo de equidad o de inequidad intergeneracional?

[DOC] El Mecanismo de Equidad Intergeneracional (MEI) sustituyó al factor de sostenibilidad. En su diseño inicial de la Ley 21/2021, era una cotización finalista de 0,6 puntos entre 2023 y 2032, complementada con un mecanismo de ajuste que se activaría en 2032 si el gasto proyectado en 2050 superaba el del Ageing Report 2024, descontado el efecto del factor derogado. El RDL 2/2023 lo modificó sustancialmente. Ahora es una cotización finalista que empieza en 0,6 puntos en 2023, aumenta una décima al año hasta 1,2 en 2029 y se mantiene hasta 2050, destinada a dotar el Fondo de Reserva, cuyos recursos podrán utilizarse a partir de 2033 con límites anuales. En 2026, el tipo del MEI es del 0,9%, y el Fondo de Reserva superó los 15.200 millones de euros en abril de 2026.

El MEI concentra buena parte de la crítica académica a la reforma, articulada en cuatro argumentos.

El primero es el nombre. Conde-Ruiz, entre otros, ha señalado la paradoja de que un mecanismo llamado "de equidad intergeneracional" sea una cotización que pagan los trabajadores actuales, más jóvenes, para financiar las pensiones de los baby boomers, sin ningún ajuste en las prestaciones de estos. El factor de sostenibilidad al que sustituía repartía el coste de la longevidad entre cohortes de jubilados. El MEI lo traslada a los cotizantes. Desde la perspectiva del principio de ahorro justo de Rawls y de la contabilidad generacional (4.1.2 y 4.1.3), es un aumento de la carga de las generaciones jóvenes.

El segundo es la escala. El MEI recauda en torno a medio punto del PIB al año en su nivel máximo [Probable en la cifra], y el aumento del gasto proyectado supera los 4 puntos. Como vimos en 4.2.1, el problema no es una joroba transitoria, sino una meseta permanente, y un fondo que se agota en pocos años a partir de 2033 no puede sustituir a un ajuste estructural.

El tercero es la circularidad financiera. El Fondo de Reserva invierte mayoritariamente en deuda pública española. Desde la óptica del conjunto de las Administraciones Públicas, una cotización adicional que se destina a comprar deuda del Estado no aumenta el ahorro nacional si el Estado, a la vez, incurre en déficit. Es un traspaso contable entre subsectores. El balance consolidado solo mejora si esa recaudación se traduce en un menor déficit del conjunto. Es la vieja advertencia de Barr (2001) sobre los fondos públicos: lo que importa es si cambia la producción futura, no el nombre del activo. El agotamiento del Fondo original entre 2012 y 2019 (2.2.3) mostró además que un fondo sin reglas estrictas se utiliza en la primera crisis.

El cuarto es la cuña fiscal y el empleo. España tiene ya una de las cuñas fiscales sobre el trabajo más altas de la OCDE para los salarios medios, según Taxing Wages. La evidencia sobre la incidencia de las cotizaciones (Gruber, 1997; lo vimos en 2.2) sugiere que, a largo plazo, recaen en gran parte sobre los salarios, y a corto plazo, con rigideces salariales, pueden reducir el empleo. BBVA Research, Fedea y la CEOE han advertido de este riesgo. El Gobierno ha respondido que el aumento es moderado y gradual, y que el empleo ha seguido creciendo desde 2023 a ritmos históricamente altos. Ninguno de los dos argumentos es concluyente todavía: el efecto de una subida de un punto sobre el empleo es difícil de identificar en un contexto de fuerte crecimiento y de inmigración elevada.

e) La base máxima, la cuota de solidaridad y la contributividad.

[DOC] El RDL 2/2023 estableció que la base máxima de cotización aumentará con el IPC más 1,2 puntos al año entre 2024 y 2050. Con una inflación del 2%, eso supone un aumento real acumulado del 37,1% en 2050. La pensión máxima, en cambio, crecerá solo 0,115 puntos anuales por encima del IPC entre 2025 y 2050, un aumento real acumulado de en torno al 3%. La cuota de solidaridad grava los salarios por encima de la base máxima con tipos que empiezan entre el 0,92% y el 1,17% en 2025 y llegan al 5,5% y el 7% en 2045, en tres tramos, repartidos entre empresa y trabajador en proporción 83/17. (Confirmada la cifra de 0,115 que dejé pendiente en 4.2.)

Es la formalización legal de la "reforma silenciosa" de Conde-Ruiz y González (2016), ahora explícita. El sistema cobra cada vez más a los salarios altos por encima de lo que les devuelve, y se desliza hacia un modelo más redistributivo en la parte alta. La valoración depende del criterio. Desde la equidad vertical, es una mejora: aumenta la progresividad de la financiación. Desde la contributividad y la lógica bismarckiana del sistema, es una erosión: rompe la correspondencia entre lo que se aporta y lo que se recibe, que es el fundamento de la legitimidad de un sistema contributivo (Korpi y Palme, 1998). Desde la eficiencia, genera un impuesto implícito sobre el trabajo cualificado que puede incentivar la elusión. En cualquier caso, es la medida de ingresos más cuantiosa de la reforma a largo plazo [Probable], y su rendimiento efectivo depende de la respuesta de las rentas altas.

f) Los primeros efectos: incentivos, edad efectiva y brecha de género.

Los primeros datos parecen confirmar la dirección esperada en la edad efectiva de jubilación. La edad media de acceso a la jubilación ha subido hasta cerca de 65,3 años a principios de 2026, frente a unos 64,5 antes de la reforma. Ha aumentado el número de jubilaciones demoradas y la proporción de jubilaciones anticipadas ha bajado ligeramente. El Real Decreto-ley 11/2024, en vigor desde abril de 2025, reforzó la compatibilidad entre pensión y trabajo (jubilación activa, parcial y flexible), en la línea de las recomendaciones de la OCDE y de la Comisión. Es pronto para atribuir estos cambios a la reforma y no a la composición de las cohortes o al buen momento del mercado de trabajo; la identificación rigurosa exigirá evaluaciones con microdatos de la Muestra Continua de Vidas Laborales. En cuanto a la brecha de género en pensiones, el complemento para su reducción ha contribuido a que la brecha en las nuevas altas baje, aunque sigue muy por encima del umbral del 5% que la ley fija para mantenerlo.

\subsubsection{Cláusula de salvaguarda: diseño, primera aplicación y debilidades}

a) El diseño.

[DOC] El RDL 2/2023 introdujo un mecanismo de cierre. Desde marzo de 2025, y cada tres años, la AIReF calcula el impacto medio anual en 2022-2050 de las medidas de ingresos aprobadas desde 2020. Si ese impacto alcanza el 1,7% del PIB, no hacen falta nuevas medidas mientras el gasto medio en pensiones de 2022-2050 no supere el 15% del PIB. El umbral sube o baja en la misma cuantía en que el impacto de los ingresos se desvíe del 1,7%. Si los ingresos rinden 0,5 puntos más, el umbral pasa al 15,5%. Si rinden solo el 1%, baja al 14,3%. El gasto de referencia es el del Ageing Report más reciente. Si se supera el umbral, el Gobierno debe presentar a la Comisión del Pacto de Toledo una propuesta de corrección, con medidas de ingresos, de gasto o de ambos tipos. Si no hay acuerdo parlamentario, la cotización del MEI sube automáticamente el 1 de enero siguiente en la cuantía necesaria para corregir el 20% de la desviación, y sigue subiendo cada año hasta corregirla o hasta que se adopten otras medidas.

La condición equivale a otra más fácil de interpretar: el gasto neto de las medidas de ingresos no debe superar el 13,3% del PIB (15% menos 1,7%), en promedio de 2022-2050. Es la forma en que la AIReF la ha aplicado. Su lógica es que el gasto que exceda el nivel previo a la reforma debe estar financiado por las nuevas medidas de ingresos.

b) La primera aplicación (marzo de 2025) y la controversia metodológica.

[Probable en las cifras, a partir de la nota de prensa de la AIReF y la prensa especializada] El 31 de marzo de 2025, la AIReF publicó su primer informe sobre la regla de gasto de pensiones:

- impacto de las medidas de ingresos: 1,4% del PIB de media en 2022-2050, por debajo del 1,7% previsto por el Gobierno;
- gasto proyectado: del 12,7% del PIB en 2022 al 16,1% en 2050, con una media en torno al 14,6%;
- gasto neto: en torno al 13,2%, por debajo del umbral del 13,3% por una décima.

Conclusión: la regla se cumplía y no era necesario activar la cláusula. Pero la AIReF añadió advertencias poco habituales en un informe de cumplimiento. La sostenibilidad del sistema no había mejorado respecto a sus proyecciones de 2023, sino que se había deteriorado en unas cuatro décimas. Es "perfectamente compatible" que la regla se cumpla y la sostenibilidad no mejore. La regla es "excesivamente sensible" al momento del cálculo y ofrece un análisis solo "parcial" de la sostenibilidad. Y, sin nuevas medidas, las transferencias del Estado a la Seguridad Social tendrían que aumentar en unos 2,4 puntos del PIB hasta 2050.

El margen de una décima se obtuvo tras una decisión normativa controvertida. Un mes antes, el Gobierno aprobó el Real Decreto 100/2025, que desarrollaba la metodología de cálculo. Permitía computar como medidas de ingresos determinadas transferencias del Estado a la Seguridad Social y aplicar una corrección al alza de la senda del PIB proyectada, basada en la revisión estadística de la contabilidad nacional y en el crecimiento observado de 2023, que reducía mecánicamente el gasto sobre el PIB. De la Fuente (2025, Fedea, Apuntes 2025/18) documentó el efecto de estas decisiones. Aplicando la regla sin esos ajustes, el gasto neto se habría situado en torno al 14,05% del PIB, 0,75 puntos por encima del umbral. Con la corrección del PIB del real decreto, bajaba al 13,18%. Las transferencias computadas, unas tres décimas, resultaron decisivas. De la Fuente habló del Gobierno como "juez y parte", al fijar las reglas de cálculo de una evaluación que le afectaba, y sostuvo que una revisión estadística del PIB crea una ilusión: mejora la ratio sin que aumenten los ingresos del sistema. El Ministerio de Inclusión defendió que la revisión del PIB era real y que las transferencias financian gastos que no deben pesar sobre las cotizaciones, en coherencia con la separación de fuentes.

En mayo de 2026, la AIReF actualizó su análisis con los datos de 2024 y 2025 y ratificó el resultado. Revisó al alza el impacto de las medidas de ingresos, hasta el 1,6% del PIB, sobre todo por los datos reales de cotización de los autónomos con el nuevo sistema. Estimó un gasto neto en torno al 13%, por debajo del umbral del 13,3%, y reiteró que el cumplimiento de la regla no garantiza la sostenibilidad del sistema. Estimó también que el efecto a largo plazo de la regularización de inmigrantes sobre la regla era prácticamente nulo. (Pendiente de verificar: la naturaleza exacta del informe de 2026, si es una actualización voluntaria o prevista en la norma, dado que la evaluación formal es trienal y la siguiente correspondería a 2028, con el Ageing Report 2027 como referencia.)

c) Las debilidades del diseño: una crítica estructurada.

La primera aplicación de la cláusula permite identificar cinco debilidades, que no son solo de los críticos de la reforma, sino que la propia AIReF ha señalado en parte.

La primera: mide el gasto, no el saldo. La condición compara el gasto neto de las nuevas medidas de ingresos con un umbral, pero ignora la evolución del resto de ingresos del sistema y el déficit de partida. Un sistema puede cumplir la regla con un déficit creciente financiado con transferencias del Estado. Es precisamente lo que proyecta la AIReF, con esas transferencias adicionales de 2,4 puntos del PIB. De la Fuente (2025) propuso reformular la condición en términos de déficit, que es la magnitud relevante para la sostenibilidad fiscal (4.1.1).

La segunda: es excesivamente sensible a las proyecciones. El resultado depende de cifras como el PIB proyectado, el rendimiento de los ingresos o las hipótesis del Ageing Report, cuya incertidumbre es muy superior al margen de una décima con el que se cumplió. Un cambio menor en cualquiera de ellas invierte el veredicto.

La tercera: la metodología la fija el Gobierno evaluado. El episodio del RD 100/2025 muestra la vulnerabilidad de una regla cuyas normas de cálculo pueden modificarse por reglamento. Enlaza con la literatura sobre las instituciones fiscales independientes (Calmfors y Wren-Lewis, 2011; Beetsma y Debrun, 2018): su eficacia depende no solo de su independencia al evaluar, sino de que las reglas que evalúan no puedan redefinirse a conveniencia.

La cuarta: el ajuste por defecto es asimétrico. Si no hay acuerdo, la cláusula corrige siempre por el lado de los ingresos, con más MEI, nunca por el de las prestaciones. Es coherente con la filosofía de la reforma, pero traslada a los cotizantes el coste de cualquier desviación futura, incluidas las debidas a una mayor longevidad. Es una diferencia de fondo con los mecanismos automáticos de otros países, como el freno sueco o el factor alemán, que reparten el ajuste entre cotizantes y pensionistas.

La quinta: la credibilidad es la de cualquier regla revocable. Como el IRP en 2018, la subida automática del MEI puede suspenderse por una ley posterior si su coste se hace visible. La lección de 2013 es que una regla automática es tan creíble como el consenso que la sostiene.

\subsubsection{Balance: quién ha soportado el ajuste}
Leídas en conjunto, las tres reformas muestran un patrón que la identidad de 4.2 hace visible. La de 2011 ajustó por la edad y la fórmula: la carga recayó sobre los futuros jubilados, en forma de más años de trabajo y pensiones iniciales algo menores, con una transición larga y pactada. La de 2013 ajustó por la generosidad: la carga recayó sobre los pensionistas actuales y futuros, a través de la revalorización y la pensión inicial, con un mecanismo automático, opaco para el ciudadano medio y sin consenso. La de 2021-2023 ajusta por los ingresos: la carga recae sobre los cotizantes actuales, sobre los salarios altos, en forma de base máxima y cuota de solidaridad, y sobre el conjunto de los contribuyentes, en forma de transferencias del Estado, con consenso parcial y una cláusula de cierre de eficacia discutida.

Tres conclusiones se desprenden de este recorrido.

Primera: el consenso ha determinado la durabilidad. La reforma pactada de 2011 sigue vigente. La impuesta de 2013 se revirtió. La de 2021-2023 descansa en un consenso parcial que la aplicación de la cláusula ha tensionado. La lección para el diseño institucional, que retomaremos en 4.4, es que el Pacto de Toledo no es un adorno, sino la condición de credibilidad de cualquier regla de largo plazo.

Segunda: ningún ajuste ha sido indoloro, pero ha cambiado quién lo paga y cómo de visible es. La reforma de 2021-2023 no ha eliminado la caída de la tasa de sustitución: según el Ageing Report 2024, sigue cayendo del 64% al 51% en 2070 (4.2.3). La ha reducido y la ha hecho depender de la productividad, no de una regla explícita. Y ha añadido una carga visible sobre los cotizantes y una menos visible sobre el presupuesto del Estado.

Tercera: la sostenibilidad del sistema se ha desplazado hacia la sostenibilidad fiscal. Con la separación de fuentes y el aumento de las transferencias, el saldo de la Seguridad Social ha dejado de ser el indicador relevante. El déficit se financia con impuestos generales o con deuda del Estado, y la pregunta pertinente ya no es "¿es sostenible la Seguridad Social?", sino "¿es sostenible una política fiscal que dedica una parte creciente de sus recursos a las pensiones, en un país con una deuda pública por encima del 100% del PIB y sujeto a las nuevas reglas fiscales europeas?". Es la pregunta de 4.1.1 sobre la sostenibilidad fiscal frente a la financiera. De ella se ocupa el apartado siguiente, que examina la economía política de ese reparto y las alternativas de diseño que plantea la literatura.

La verdad incómoda con la que abre este apartado es que la literatura no ofrece ninguna alternativa de diseño que elimine el coste del envejecimiento. Capitalización, cuentas nocionales, segundo pilar o ampliación de la base de financiación cambian quién paga, cuándo paga y cuánto se ve, pero no suprimen el hecho de que una población más envejecida debe destinar una parte mayor de su producción a las personas mayores. Es la tesis de Barr ("la producción es lo central") y el llamado principio de Mackenroth (1952). Si en el examen se presenta alguna alternativa como "la solución", el tribunal lo leerá como falta de rigor. He redactado 4.4 para que cada alternativa se evalúe con ese criterio. La he precedido de la economía política, porque sin ella no se entiende por qué España ha elegido las opciones que ha elegido.

\subsection{Economía política y alternativas de diseño}
El recorrido de 4.3 deja una paradoja. Existe un consenso técnico amplio sobre la magnitud del problema: la Comisión Europea, la AIReF, el Banco de España, la OCDE y la mayor parte de la academia coinciden en que el gasto en pensiones aumentará varios puntos del PIB y en que las medidas de ingresos aprobadas no lo compensan del todo. Y, sin embargo, las reformas han oscilado entre direcciones opuestas en apenas una década. La explicación no está en la técnica, sino en la política. Las pensiones son la mayor partida del gasto público español, afectan directamente a más de nueve millones de pensionistas y, de forma diferida, a más de veintidós millones de ocupados. Además, redistribuyen recursos entre grupos con intereses, capacidad de organización y peso electoral muy distintos. Cualquier alternativa de diseño debe evaluarse no solo por sus propiedades económicas, sino por su viabilidad política: su capacidad para aprobarse y, sobre todo, para mantenerse cuando su coste se hace visible.

El apartado se organiza en dos partes. La primera, en los epígrafes 4.4.1 y 4.4.2, examina la economía política del sistema: por qué los sistemas de reparto son tan resistentes al cambio, qué papel desempeñan los mayores como grupo de presión, cómo se construyen las coaliciones que hacen posibles las reformas y hasta qué punto el llamado "conflicto intergeneracional" es real o retórico. La segunda, en los epígrafes 4.4.3 a 4.4.5, evalúa las alternativas de diseño que plantea la literatura: las que operan dentro del reparto (la vinculación de la edad a la longevidad, las cuentas nocionales, la separación de fuentes, la reforma de la viudedad), las que operan fuera (la capitalización y el segundo pilar) y las que actúan sobre la base de financiación.

Una advertencia de solapamiento: la comparación teórica entre reparto y capitalización (Samuelson, Aaron, Feldstein, Barr y Orszag-Stiglitz, el Banco Mundial de 1994 y su revisión) ya se expuso en 2.2, y la incidencia de las cotizaciones y la devaluación fiscal, en 2.2.2 y 2.2.3. Aquí se aplican esos marcos al caso español actual, no se repiten.

\subsubsection{Por qué el reparto se resiste al cambio: votantes, coaliciones e instituciones}

a) El modelo del votante mediano: la tendencia a la sobreexpansión.

La economía política formal de las pensiones parte de Browning (1975, Journal of Political Economy), cuyo razonamiento es sencillo y potente. En un sistema de reparto, el valor de ampliar las pensiones es muy distinto según la edad del votante. Un jubilado recibe el aumento sin haber pagado las cotizaciones adicionales. Un trabajador de 55 años paga unos pocos años de cotizaciones más altas y recibe toda su jubilación mejorada. Un joven de 25 paga cuarenta años de cotizaciones más altas a cambio de una pensión lejana e incierta. Si las decisiones se toman por mayoría y el votante mediano tiene una edad avanzada, la mayoría preferirá un sistema más grande del que elegiría un planificador que ponderase por igual a todas las generaciones, incluidas las que aún no votan o no han nacido. El envejecimiento, al desplazar la edad del votante mediano hacia arriba, refuerza esa tendencia.

La literatura posterior refinó el modelo en tres direcciones.

La primera: si el votante mediano prefiere siempre más reparto, el sistema debería crecer sin límite, y no lo hace. Tabellini (2000, Scandinavian Journal of Economics) explicó los límites y la supervivencia del sistema por su redistribución intrageneracional: los trabajadores de rentas bajas apoyan el reparto porque, a través de la progresividad de las prestaciones, reciben más de lo que aportan, y forman con los mayores una coalición mayoritaria. Conde-Ruiz y Galasso (2005, Journal of Public Economics) desarrollaron esa "aritmética positiva del Estado de bienestar", y Conde-Ruiz y Galasso (2003, Review of Economic Dynamics; 2004) explicaron con ese marco la persistencia de la jubilación anticipada: los trabajadores mayores de baja productividad y los jubilados forman una coalición que la sostiene aunque sea ineficiente. Ambos son economistas españoles y su trabajo se aplica directamente a las prejubilaciones españolas.

La segunda: los modelos de votación de una sola vez no explican por qué los jóvenes no desmantelan el sistema cuando son mayoría. La respuesta, formalizada por Cooley y Soares (1999, Journal of Political Economy) y resumida por Galasso y Profeta (2002, European Journal of Political Economy), es la sostenibilidad como equilibrio. Cada generación de trabajadores mantiene el sistema porque, si lo abandonara, perdería sus propios derechos futuros, dado que las generaciones siguientes no tendrían ningún incentivo para pagarle. El reparto es un contrato implícito que se mantiene por las expectativas, y su ruptura es cara para quien la inicia.

La tercera, empírica. Breyer y Craig (1997, European Journal of Political Economy) encontraron, con datos de la OCDE, una relación positiva entre la edad del votante mediano y el gasto en pensiones. Sinn y Uebelmesser (2002, European Journal of Political Economy) proyectaron que Alemania se convertiría en una gerontocracia, con más de la mitad de los votantes por encima de la edad a partir de la cual se gana con el reparto, en torno a 2016. Pero Disney (2007, Economic Policy) planteó un contraargumento importante. En los países de la OCDE, el envejecimiento se asocia a un aumento del gasto total en pensiones, pero a una reducción de la generosidad por pensionista: los sistemas absorben el aumento del número de beneficiarios recortando la prestación relativa. El "poder gris" no protege la pensión media, sino, como mucho, el gasto agregado.

b) El caso español: un electorado envejecido y muy movilizado.

En España, los mayores de 65 años representan en torno a una cuarta parte del censo electoral, con tasas de participación muy superiores a las de los menores de 35. La edad mediana del electorado ronda los 50 años y seguirá subiendo. El movimiento de pensionistas de 2018, analizado en 4.3.2, es la ilustración más clara de su capacidad de movilización: forzó en pocos meses la suspensión de una reforma aprobada cinco años antes. El episodio de enero de 2025 es igual de revelador. El Congreso rechazó un real decreto-ley "ómnibus" que incluía la revalorización de las pensiones junto con otras medidas, y en pocos días el Gobierno aprobó un nuevo decreto que recogía la revalorización, convalidado con apoyos amplios. Ningún partido con aspiraciones de gobierno está dispuesto a aparecer como responsable de una pensión no revalorizada. La consecuencia es una asimetría política: las medidas que mejoran las prestaciones se aprueban con facilidad y amplias mayorías, y las que las contienen requieren consenso previo o mecanismos que diluyan la responsabilidad.

c) Cómo se reforma pese a todo: evitación de la culpa y construcción de coaliciones.

Si el modelo del votante mediano fuera la última palabra, ninguna reforma restrictiva sería posible. Y, sin embargo, casi todos los países europeos han aprobado reformas que reducen las prestaciones futuras. La ciencia política ha identificado dos mecanismos.

El primero es la evitación de la culpa (Weaver, 1986; Pierson, 1994). Las reformas restrictivas se diseñan para que su coste sea diferido (se aplica a generaciones futuras, con largas transiciones), difuso (se reparte entre muchos) u opaco (se esconde en fórmulas técnicas o mecanismos automáticos). La reforma de 2011, con transición hasta 2027, es un ejemplo de lo primero. El IRP de 2013, de lo tercero. La "reforma silenciosa" de la base máxima, de lo segundo y lo tercero. El fracaso del IRP muestra el límite de la opacidad: funciona hasta que el efecto se hace visible.

El segundo es la construcción de coaliciones mediante paquetes. Bonoli (2000, The Politics of Pension Reform) mostró, comparando Suiza, Reino Unido y Francia, que las reformas restrictivas que prosperan suelen combinarse con contrapartidas para grupos concretos: recortes generales a cambio de mejoras para las mujeres, los trabajadores precarios o las rentas bajas. Häusermann (2010, The Politics of Welfare State Reform in Continental Europe) generalizó la idea. En sociedades postindustriales, los conflictos sobre las pensiones no son unidimensionales (más o menos Estado), sino multidimensionales (quién gana, quién pierde y qué riesgos se cubren), y eso permite ingeniería de coaliciones: juntar a actores con preferencias distintas en un mismo paquete. La reforma española de 2021-2023 encaja en este marco. Combinó más cotizaciones (MEI, base máxima, cuota de solidaridad) con mejoras visibles para grupos amplios (revalorización con el IPC, pensiones mínimas ligadas al umbral de la pobreza, complemento de brecha de género, integración de lagunas), y así obtuvo el apoyo de los sindicatos y de la mayoría parlamentaria de izquierdas y nacionalista. El coste recayó en grupos con menor capacidad de veto: los salarios altos, las empresas y, de forma difusa, los contribuyentes futuros.

d) Las instituciones del consenso: el Pacto de Toledo y la Unión Europea.

El Pacto de Toledo es la institución española de gestión de este conflicto. Nacido en 1995 como un acuerdo parlamentario para sacar las pensiones de la confrontación electoral, se ha renovado con informes de evaluación y reforma en 2003, 2011 y 2020. El último, de octubre de 2020, aprobó una veintena de recomendaciones: separación de fuentes, revalorización con el IPC, derogación del factor de sostenibilidad, mejora de la contributividad, retraso voluntario de la jubilación, impulso de la previsión social complementaria de empleo y reforma de la viudedad. La literatura sobre la "política de la austeridad permanente" (Pierson, 2001) y los estudios comparados sobre el consenso en las reformas de pensiones (Myles y Pierson, 2001; Natali, 2008) coinciden en que los grandes pactos transversales, como el sueco de 1994-1998 (el pensionsgruppen, que agrupaba a cinco partidos) o el Pacto de Toledo, son la condición de durabilidad de las reformas. Su debilidad es la contraria: las recomendaciones consensuadas suelen ser genéricas, y el consenso se rompe precisamente cuando hay que concretar quién paga, como ocurrió en 2013 y parcialmente en 2023. Desde 2024-2025 la Comisión del Pacto de Toledo ha abierto un nuevo proceso de revisión de sus recomendaciones. Pendiente de verificar su estado a la fecha del examen.

La Unión Europea se ha convertido en el segundo actor institucional. Lo es a través de tres canales: las recomendaciones específicas por país del Semestre Europeo, que durante años pidieron a España reforzar la sostenibilidad del sistema; la condicionalidad del Mecanismo de Recuperación y Resiliencia, que convirtió la reforma de pensiones en hito para el desembolso de fondos (4.3.3); y el nuevo marco de gobernanza fiscal de 2024, que vincula las sendas de gasto a la sostenibilidad a largo plazo y tiene en cuenta el coste del envejecimiento. Este anclaje externo cumple una función conocida en la literatura de la evitación de la culpa: los gobiernos pueden atribuir a Bruselas las decisiones impopulares. Pero, como mostró la aplicación de la cláusula de salvaguarda, el aval europeo no sustituye a la credibilidad interna.

\subsubsection{Conflicto intergeneracional: entre la realidad y la retórica}

a) El argumento: una transferencia creciente de jóvenes a mayores.

El debate público español ha incorporado en la última década el marco del conflicto intergeneracional: la idea de que el sistema de pensiones, junto con otros rasgos del Estado de bienestar y del mercado de trabajo, transfiere recursos de forma creciente de generaciones jóvenes, precarias y con peores perspectivas a generaciones mayores, con patrimonio y rentas protegidas. El argumento tiene sustento empírico.

- Pobreza por edades. Según la Encuesta de Condiciones de Vida del INE, la tasa de riesgo de pobreza de los menores de 18 años supera el 25%, una de las más altas de la UE, mientras que la de los mayores de 65 se sitúa bastante por debajo, sobre todo si se tiene en cuenta el alquiler imputado de la vivienda en propiedad. Hace treinta años el patrón era el contrario.
- Riqueza. La Encuesta Financiera de las Familias del Banco de España muestra una concentración creciente de la riqueza neta, sobre todo inmobiliaria, en los hogares de más edad, y una caída de la tasa de propiedad de la vivienda entre los jóvenes.
- Rentas relativas. La pensión media de jubilación de las nuevas altas del Régimen General supera en muchos casos el salario mediano de los menores de 30 años.
- Contabilidad generacional. Los ejercicios que vimos en 4.1.3, de Auerbach, Gokhale y Kotlikoff aplicados a España por Bonin y Patxot, y las Cuentas Nacionales de Transferencias de Lee y Mason, con los trabajos de Patxot, Rentería y Souto, muestran que las generaciones futuras soportan una carga fiscal neta mayor que las actuales.

Vanhuysse (2013), con su índice de justicia intergeneracional para la OCDE, sitúa a los países del sur de Europa entre los que más sesgan su política social a favor de los mayores, y Lynch (2006) relacionó ese sesgo con el origen bismarckiano y particularista de sus Estados de bienestar. Los defensores de este diagnóstico, en España sobre todo en Fedea (Conde-Ruiz, De la Fuente), en parte de la prensa económica y en think tanks liberales, concluyen que las reformas de 2021-2023 agravan el desequilibrio al proteger a los pensionistas actuales y cargar el ajuste sobre los cotizantes jóvenes a través del MEI.

b) El contraargumento: familias, ciclo vital y clase.

La crítica a este marco tiene tres componentes.

El primero es familiar. En los regímenes de bienestar familistas del sur de Europa (Ferrera, 1996; Moreno, lo vimos en 1.4), las transferencias públicas a los mayores se redistribuyen en buena medida dentro de la familia hacia los jóvenes. Kohli (1999, European Societies) y Albertini y Kohli (2013, European Sociological Review) documentaron que los flujos privados netos van mayoritariamente de mayores a jóvenes (ayudas monetarias, vivienda, herencias, cuidado de los nietos), especialmente en el sur de Europa. En España, durante la crisis de 2008-2014, las pensiones actuaron como un colchón familiar decisivo ante el desempleo de hijos y nietos. Recortar las pensiones no transferiría esos recursos a los jóvenes de forma automática: a menudo les quitaría un sostén.

El segundo es de ciclo vital. Las generaciones no son grupos de interés permanentes: todos los jóvenes de hoy serán mayores mañana. Un sistema de pensiones generoso es, para un joven, también su propio seguro futuro. El conflicto relevante no es entre personas, sino entre cohortes. Se mide por las tasas internas de rendimiento o la contabilidad generacional, no por la comparación entre la situación actual de jóvenes y mayores.

El tercero es de clase. La literatura crítica (Bristow, 2015; Willetts, 2010, desde el otro lado del espectro, con The Pinch) advierte que el marco generacional puede ocultar desigualdades de clase dentro de cada generación. Hay mayores pobres, sobre todo mujeres viudas con pensiones mínimas, y jóvenes ricos, herederos de patrimonio familiar. Presentar el problema como "los mayores contra los jóvenes" puede servir para legitimar recortes que afectan sobre todo a los pensionistas de rentas bajas, sin tocar la desigualdad de patrimonio que se transmite por herencia.

c) La confianza y la narrativa del "no habrá pensiones".

Las encuestas del CIS y de otros institutos muestran de forma recurrente que una parte muy elevada de los jóvenes españoles cree que no cobrará pensión o que será muy inferior a la actual. La sociología y la ciencia política señalan dos riesgos de esa desconfianza. Erosiona la legitimidad del contrato implícito del reparto, que, como vimos con Cooley y Soares, se sostiene por las expectativas. Y puede reducir el esfuerzo contributivo, por ejemplo con la infracotización de los autónomos o la economía sumergida. Desde posiciones críticas (Navarro, Zubiri y parte del sindicalismo) se ha sostenido que el alarmismo sobre las pensiones responde en parte a los intereses de la industria financiera, que se beneficiaría de un mayor peso de los planes privados. La ortodoxia responde que las proyecciones proceden de instituciones públicas independientes (la Comisión, la AIReF, el Banco de España) y no de la industria. Ambas cosas pueden ser ciertas a la vez: las proyecciones son sólidas y hay intereses privados que aprovechan su comunicación. La respuesta institucional adecuada es la transparencia. Volveremos sobre ello a propósito de la información individual en 4.4.3.

\subsubsection{Alternativas dentro del reparto}

a) La vinculación de la edad de jubilación a la esperanza de vida.

Es la alternativa con más consenso académico y más extensión internacional. Dinamarca (desde 2006, con aplicación efectiva desde 2030 y la edad legal camino de 70 años a partir de 2040), Países Bajos (desde 2012, con una revisión en 2019 que suavizó el vínculo: dos tercios de las ganancias de esperanza de vida se trasladan a la edad), Finlandia (2017), Portugal, Italia, Grecia, Chipre, Estonia y Eslovaquia vinculan de algún modo la edad legal a la longevidad. Según la OCDE (Pensions at a Glance 2023), la edad normal de jubilación futura en los países que lo han hecho superará los 67 años, y en Dinamarca alcanzará 74 años hacia 2070 si se mantiene el vínculo.

El fundamento es el que expusimos en 4.2.1 y 4.3.2: si la vida se alarga, repartir las ganancias entre años de trabajo y años de jubilación mantiene constante la proporción de la vida adulta en que se cobra la pensión. La regla de los dos tercios, que propuso la Comisión Europea en su Libro Blanco de 2012 y aplican Finlandia y los Países Bajos, destina dos tercios de las ganancias de longevidad a trabajar más y un tercio a la jubilación, y estabiliza aproximadamente la relación entre años de trabajo y de jubilación.

Frente al factor de sostenibilidad español de 2013, que ajustaba la cuantía para una edad dada, la vinculación de la edad tiene dos ventajas. Mantiene la tasa de sustitución, porque el ajuste se hace en años, no en euros, lo que la hace más compatible con la suficiencia. Y actúa a la vez sobre la tasa de cobertura y la de empleo (4.2.2), mientras que el factor solo actuaba sobre la generosidad. Sus inconvenientes son los señalados en 4.2.1: el gradiente de longevidad por nivel socioeconómico la hace regresiva si no se acompaña de vías de salida para quienes tienen trabajos penosos o carreras largas, y la capacidad real de trabajar a edades avanzadas es desigual. El Banco de España y la OCDE en sus estudios económicos de España han recomendado explorar esta vía. El Gobierno y los sindicatos la han rechazado en favor de incentivos a la jubilación voluntaria más tardía.

b) Las cuentas nocionales de contribución definida (NDC).

Las cuentas nocionales, introducidas en Suecia (1994-1998), Italia (1995, reforma Dini), Letonia (1996), Polonia (1999) y Noruega (2011), mantienen la financiación de reparto pero calculan las prestaciones como en un sistema de capitalización de aportación definida. Cada trabajador tiene una cuenta virtual en la que se anotan sus cotizaciones, que se capitalizan a un tanto nocional, ligado al crecimiento de los salarios o de la masa salarial. En el momento de la jubilación, el saldo se convierte en una renta vitalicia dividiendo por un divisor basado en la esperanza de vida de la cohorte. La sostenibilidad descansa en tres elementos: la rentabilidad nocional no supera, por construcción, el crecimiento de la base de financiación (la condición de Aaron-Samuelson de 2.2); el divisor absorbe automáticamente la longevidad; y, en el caso sueco, un mecanismo de equilibrio automático (el freno, bromsen) reduce la revalorización si los pasivos superan a los activos. (Remisión: la teoría de las NDC y la contribución de Holzmann y Palmer, en 2.2.)

En España, las cuentas nocionales han sido la propuesta de reforma estructural preferida por una parte relevante de la academia. Destacan la propuesta de Fedea de 2010 para un sistema de cuentas nocionales, los trabajos de Boado-Penas, Vidal-Meliá y Valdés-Prieto sobre el balance actuarial (4.1.3), los de Devesa y colaboradores, y Conde-Ruiz y González, que la presentaron como alternativa a la reforma de 2013. Sus ventajas: transparencia, porque cada trabajador conoce el valor de su cuenta; neutralidad actuarial, porque se elimina el incentivo a la jubilación anticipada y se premia el retraso; resistencia a la manipulación política, porque las reglas son automáticas; y equidad entre cohortes.

Las críticas son también de peso. Barr (2006) y Barr y Diamond (2008) sostuvieron que las NDC no hacen nada que no pueda lograrse con reformas paramétricas bien diseñadas en un sistema de prestación definida, que su "sostenibilidad automática" depende de reglas que pueden suspenderse (como ocurrió con el freno sueco en 2010, cuando el Gobierno compensó con rebajas fiscales a los pensionistas el recorte automático) y que eliminan la redistribución intrageneracional, que debe reconstruirse con una pensión mínima garantizada financiada con impuestos. La experiencia italiana muestra además la lentitud de la transición: la reforma de 1995 aplicaba el sistema nocional íntegramente solo a quienes empezaron a cotizar después de 1996, de modo que sus efectos plenos no se producirán hasta mediados de este siglo. Para España, una transición a cuentas nocionales exigiría un consenso político que hoy no existe y generaría pérdidas visibles para los trabajadores de salarios crecientes y carreras largas, que hoy se benefician de una fórmula no actuarial.

c) La información individual: el "sobre naranja" que nunca llegó.

Una alternativa de bajo coste, que la literatura conductual respalda, es la información individual sobre la pensión futura. Suecia envía cada año desde 1999 el "sobre naranja" (orange kuvert) con la estimación de la pensión de cada trabajador en distintos escenarios. La Ley 27/2011 obligaba a la Seguridad Social española a informar a los trabajadores de su pensión estimada (disposición adicional vigésima novena), pero su aplicación se aplazó reiteradamente y se limitó al simulador en línea "Tu Seguridad Social", sin envío proactivo. La evidencia sobre el efecto de estos envíos es mixta: en Estados Unidos, el Social Security Statement mejora el conocimiento, pero con efectos pequeños sobre el comportamiento (Mastrobuoni, 2011). En Suecia y en otros países, su valor es sobre todo de transparencia y legitimidad. En el contexto español de desconfianza juvenil (4.4.2), la información individual tendría un valor específico: sustituir la narrativa de "no habrá pensiones" por una estimación verificable.

d) La separación de fuentes y la reforma de la viudedad.

La separación de fuentes, principio del Pacto de Toledo desde 1995, establece que las prestaciones contributivas se financian con cotizaciones y las no contributivas y universales con impuestos. La Ley 21/2021 dio un paso importante al reconocer como "gastos impropios" y trasladar al Estado conceptos como las reducciones de cuotas, las prestaciones de maternidad y paternidad o el complemento de brecha de género. Su efecto contable es mejorar el saldo de la Seguridad Social, pero no el de las Administraciones Públicas. Como vimos en 4.3.5, desplaza el problema del sistema al presupuesto del Estado. Su valor está en la transparencia: permite distinguir qué parte del déficit corresponde al núcleo contributivo y cuál a decisiones de política social.

La pensión de viudedad es la gran asignatura pendiente del Pacto de Toledo. Diseñada para un modelo de familia con un único sustentador masculino, hoy se concede con independencia de la renta del superviviente y es compatible con rentas del trabajo y con la pensión propia. Beneficia sobre todo a mujeres mayores, muchas de las cuales dependen de ella, pero también a viudos y viudas con rentas altas. Las recomendaciones del Pacto de Toledo de 2011 y 2020 plantearon reformularla para concentrarla en quienes carecen de otros ingresos, respetando los derechos adquiridos. Es un ejemplo de la dificultad política de las reformas: afecta sobre todo a un colectivo vulnerable y muy numeroso (más de dos millones de pensiones), y cualquier recorte se percibe como un ataque a las viudas.

\subsubsection{Alternativas fuera del reparto: capitalización y segundo pilar}

a) La transición a la capitalización: el doble pago.

La propuesta más radical, sustituir total o parcialmente el reparto por la capitalización individual, perdió buena parte de su atractivo académico tras las críticas de Orszag y Stiglitz (2001) y Barr (2002), y su atractivo político tras la experiencia latinoamericana y de Europa del Este. (Remisión: teoría en 2.2.) En el caso español, su principal obstáculo es el problema del doble pago. En la transición, una generación debe financiar a la vez las pensiones de los jubilados del sistema de reparto y su propio ahorro en el sistema de capitalización. Breyer (1989) demostró que, en una economía dinámicamente eficiente, no existe una transición Pareto-mejorante del reparto a la capitalización: alguna generación debe perder. El coste de la transición es la deuda implícita del sistema, que en España equivale a varias veces el PIB (4.1.3). Hacerla explícita o financiarla con impuestos no crea riqueza: redistribuye la carga entre generaciones.

La experiencia internacional ha reforzado el escepticismo. Ortiz et al. (2018, OIT) documentaron que, de treinta países que privatizaron total o parcialmente sus sistemas entre 1981 y 2014, dieciocho habían revertido la privatización total o parcialmente, entre ellos Argentina, Hungría, Polonia y Bolivia. Chile, el caso fundador, ha tenido que complementar su sistema de capitalización con un pilar solidario (2008) y aprobó en 2025 una reforma que introduce una cotización a cargo del empleador destinada en parte a un componente de seguro social con elementos de reparto. Pendiente de verificar el detalle de la reforma chilena de 2025.

b) El segundo pilar: la debilidad estructural española.

Descartada la sustitución, la alternativa que sí tiene amplio respaldo es el complemento: un segundo pilar de previsión social complementaria, preferentemente de empleo, que diversifique las fuentes de renta en la jubilación. El argumento teórico lo formularon Merton (1983) y, para las pensiones, Dutta, Kapur y Orszag (2000) o Settergren (2001): el reparto y la capitalización están expuestos a riesgos distintos, demográficos y salariales uno, financieros el otro, y no perfectamente correlacionados, lo que hace óptima una combinación (2.2).

España es una anomalía en este terreno. Según Inverco, el patrimonio total de los planes de pensiones rondaba los 138.000 millones de euros a finales de 2025, algo más del 8% del PIB: unos 96.000 millones en el sistema individual y unos 41.000 millones en el de empleo, con unos 3,1 millones de cuentas en este último. Son cifras muy lejanas de las de los Países Bajos, Dinamarca o Islandia, donde los activos de los fondos de pensiones superan el 100% o incluso el 150% del PIB, y de la media de la OCDE. La causa es, en parte, que el pilar público tiene tasas de sustitución muy elevadas (4.5), lo que reduce la necesidad de ahorro complementario, y en parte, la concentración del empleo en pymes sin capacidad para promover planes y la preferencia de los hogares por la vivienda como forma de ahorro para la vejez.

c) La reforma de 2021-2022: del incentivo fiscal individual al de empleo.

Entre 2021 y 2022, el Gobierno reorientó la política de previsión complementaria. Redujo el límite de reducción en el IRPF por aportaciones a planes individuales de 8.000 a 2.000 euros en 2021 y a 1.500 euros en 2022, y a la vez elevó el de los planes de empleo, con un límite conjunto de hasta 10.000 euros para aportaciones empresariales y de los trabajadores en ciertas condiciones. La Ley 12/2022, de 30 de junio, de regulación para el impulso de los planes de pensiones de empleo, creó los fondos de pensiones de empleo de promoción pública, de comisiones reducidas y gestionados por entidades privadas seleccionadas por concurso, y los planes simplificados para autónomos y sectores. Los quince fondos de promoción pública se inscribieron en diciembre de 2023, gestionados por cinco entidades (Caser, Gestión de Previsión y Pensiones, Ibercaja, Santander y VidaCaixa). A finales de 2025 había 45 planes simplificados registrados, con cerca de un millón de partícipes y unos 725 millones de euros de patrimonio, en buena parte por la creación de planes sectoriales como el de la construcción, pactado en convenio colectivo.

La justificación teórica del giro es doble. Por un lado, los incentivos fiscales a los planes individuales eran regresivos: la reducción en la base del IRPF beneficia más a quien tiene un tipo marginal más alto, y la mayor parte del beneficio fiscal se concentraba en las rentas altas. Es un ejemplo del "Estado de bienestar oculto" de Howard (1997) y de la "división social del bienestar" de Titmuss (1958), que vimos en 2.3. Por otro, la evidencia sugiere que los incentivos fiscales generan poco ahorro nuevo. Chetty, Friedman, Leth-Petersen, Nielsen y Olsen (2014, Quarterly Journal of Economics) mostraron con datos daneses que alrededor del 85% de los individuos son "pasivos" y no responden a los incentivos fiscales, y que por cada corona de subvención se generaba apenas un céntimo de ahorro nuevo. En cambio, las aportaciones automáticas del empleador sí aumentaban el ahorro total.

El resultado, hasta ahora, es ambivalente. Inverco y el sector han calificado la reforma de fracaso en su primera fase: las aportaciones a planes individuales se desplomaron tras la rebaja del límite y el crecimiento de los planes de empleo no las ha compensado. Las aportaciones brutas totales se situaron en torno a 3.250 millones de euros en 2025, y las netas de prestaciones apenas en 370 millones. El Gobierno sostiene que el despliegue de los planes sectoriales y de promoción pública está en sus primeras fases.

d) La inscripción automática: la lección de la economía conductual.

La evidencia internacional más sólida sobre cómo extender el segundo pilar procede de la economía conductual. Madrian y Shea (2001, Quarterly Journal of Economics) mostraron que cambiar la opción por defecto, de "no participar salvo que se elija" a "participar salvo que se renuncie", elevaba la participación en planes de empresa estadounidenses de menos del 50% a más del 85%. Thaler y Benartzi (2004, Journal of Political Economy), con el programa Save More Tomorrow, mostraron que comprometer de antemano aumentos futuros de la aportación, ligados a subidas salariales, multiplicaba la tasa de ahorro. El Reino Unido aplicó estos principios a escala nacional con la inscripción automática (auto-enrolment) desde 2012 y el fondo público NEST. La participación de los asalariados del sector privado elegibles en planes de pensiones pasó de en torno al 40% a más del 85%. Nueva Zelanda (KiwiSaver, 2007), Italia (para el TFR) y Polonia (PPK, 2019) han aplicado variantes. En España, la inscripción automática ha sido recomendada por la OCDE y debatida en el Pacto de Toledo, y el País Vasco la aplica en la práctica a través de las EPSV de empleo sectoriales como Geroa, de adhesión obligatoria por convenio en la mayoría de los sectores de Gipuzkoa. Es el principal ejemplo español de segundo pilar extendido.

e) La crítica macroeconómica: el segundo pilar no resuelve el envejecimiento.

Aquí hay que volver a la verdad incómoda del inicio. El segundo pilar diversifica riesgos y puede aumentar el ahorro, pero no elimina el coste del envejecimiento. Si una población más envejecida tiene más activos financieros, los venderá durante la jubilación a una población activa menor, y el precio de esos activos lo reflejará. Es la hipótesis de la caída de los precios de los activos (asset meltdown) de Mankiw y Weil (1989) y Poterba (2001, 2004), que vimos en 2.2. La evidencia empírica sugiere que ese efecto es moderado en economías abiertas, porque el capital puede invertirse en países más jóvenes. Esa es precisamente la ventaja de la diversificación internacional de los fondos, pero no es nula. Además, los planes de aportación definida trasladan al trabajador los riesgos financiero y de longevidad (el enigma de las rentas vitalicias de Yaari, 2.1), y sus comisiones pueden reducir sustancialmente el valor final del ahorro. El segundo pilar es un complemento útil, no una alternativa al primero.

\subsubsection{Base de financiación y la respuesta de fondo: empleo y productividad}
La última familia de alternativas actúa sobre los ingresos: ampliar la base de financiación más allá de los salarios. Ya se examinó en 2.2.3: financiación con impuestos generales (el modelo beveridgeano), IVA afectado o devaluación fiscal (Farhi, Gopinath e Itskhoki, 2014; De Mooij y Keen, 2013), contribuciones sociales generalizadas sobre todas las rentas (la CSG francesa) o imposición sobre la riqueza o el capital. En 4.2.4 añadimos el argumento del declive de la participación salarial y el debate sobre la imposición de la automatización. La reforma de 2021-2023 ya ha dado pasos en esa dirección, con las transferencias del Estado para gastos impropios y la cuota de solidaridad sobre salarios por encima de la base máxima, y las proyecciones de la AIReF implican transferencias adicionales de unos 2,4 puntos del PIB hasta 2050 si no se adoptan otras medidas.

La advertencia teórica es la misma que para las demás alternativas. Cambiar la base de financiación no crea recursos: los traslada. Financiar las pensiones con IVA en lugar de cotizaciones hace que los pensionistas contribuyan a través de su consumo. Financiarlas con impuestos sobre el capital traslada la carga a los propietarios, con posibles efectos sobre la inversión en una economía abierta. Financiarlas con deuda traslada la carga a las generaciones futuras. Cada opción tiene su distribución, sus costes de eficiencia y su viabilidad política, pero ninguna evita la restricción de fondo.

Esa restricción conduce a la respuesta de fondo, que no es un diseño alternativo sino una condición para que cualquier diseño funcione: más empleo y más productividad. Es la conclusión que se desprende de 4.2. Una tasa de empleo española convergente con la media europea reduciría significativamente la presión del gasto en la fase más crítica (4.2.2). Un crecimiento de la productividad cercano a las hipótesis del Ageing Report es la condición para que la caída proyectada de la tasa de sustitución no sea una caída del poder adquisitivo real, sino solo relativo (4.2.4). La política de pensiones, en este sentido, no puede separarse de la política de empleo, de la educativa, de la de innovación y de la de inmigración. Es la síntesis de Barr que ha guiado todo el bloque: la producción es lo central.

La valoración de conjunto que ofrece la literatura, desde la OCDE hasta el Banco de España, pasando por buena parte de la academia española, apunta a una combinación de medidas y no a una única alternativa:

- un vínculo parcial y socialmente modulado entre la edad de jubilación y la longevidad;
- una mayor contributividad de la fórmula, con una pensión mínima reforzada y financiada con impuestos;
- información individual transparente;
- un segundo pilar de empleo con inscripción automática;
- una base de financiación más amplia;
- y, sobre todo, un consenso político que dé durabilidad a las reglas.

La diferencia entre las posiciones no está tanto en la lista como en el peso de cada elemento, y en particular en cuánto del ajuste debe recaer en las prestaciones y cuánto en los ingresos. Esa es, al final, la pregunta política que ninguna técnica resuelve.

El apartado siguiente sitúa estas cuestiones en perspectiva comparada: cómo se sitúa España frente a los países de su entorno en gasto, generosidad, edad de jubilación y diseño, y qué enseñan las experiencias de otros sistemas.


No hay que presentar la tasa de reemplazo del 80% de la OCDE como prueba de que las pensiones españolas son "las más generosas del mundo". Es tentador, porque es el dato más citado, pero ese es el uso que más objeciones recibe en un tribunal técnico. Es una tasa teórica: la calcula la OCDE para un trabajador imaginario con carrera completa e ininterrumpida desde los 22 años, algo que muchos pensionistas españoles reales no tienen. Mide la generosidad de las reglas, no la de las pensiones efectivamente pagadas. El dato es valioso, pero solo si se explica qué mide. He construido 4.5 sobre esa distinción y he actualizado todo lo del documento: datos del Ageing Report 2021 y de Pensions at a Glance 2021, que el documento escribe con erratas, al Ageing Report 2024 y a Pensions at a Glance 2025, publicado en noviembre de 2025. También he detectado una errata relevante en la nota 28 del documento, que explico en su lugar.

\subsection{Perspectiva comparada: España frente a la UE y la OCDE}
Los apartados anteriores han analizado el sistema español en sus propios términos: su diagnóstico, sus reformas, su economía política. La comparación internacional añade una dimensión imprescindible, por tres razones. La primera es normativa: la afirmación de que un sistema es "generoso", "caro" o "insuficiente" solo tiene sentido en relación con algún referente, y los países de nuestro entorno son el más natural. La segunda es política: las instituciones europeas evalúan el sistema español con criterios comparados, y buena parte del debate público ("España gasta más que la media", "tiene la tasa de reemplazo más alta de la OCDE") se formula en esos términos. La tercera es de aprendizaje: otros países han ensayado mecanismos que España discute (vínculos a la esperanza de vida, cuentas nocionales, inscripción automática, fondos colectivos), y su experiencia es la mejor evidencia disponible sobre sus efectos.

[DOC] El documento advierte con razón de que los sistemas de pensiones tienen muchas características propias de cada país, por lo que la comparación homogénea es muy difícil y debe tomarse con cautela. Conviene concretar esa cautela en cinco problemas metodológicos, porque son el tipo de matiz que distingue una exposición de nivel.

El primero es el perímetro. El gasto "público" en pensiones no es comparable si en unos países las pensiones complementarias son obligatorias y privadas (Países Bajos, Dinamarca, Suiza, Australia) y en otros casi no existen (España, Italia, Grecia). Un país con un gasto público bajo puede tener un gasto total en pensiones, público más privado obligatorio, muy superior.

El segundo es bruto frente a neto. Las pensiones tributan de forma muy distinta: en algunos países pagan cotizaciones sociales y el IRPF completo; en otros, disfrutan de exenciones. Adema y Ladaique (lo vimos en 2.3.1) mostraron que las comparaciones de gasto bruto pueden invertir el orden de los países una vez se descuentan los impuestos que el Estado recupera.

El tercero es teórico frente a efectivo. Los indicadores de la OCDE sobre tasas de reemplazo son proyecciones teóricas para trabajadores tipo que empiezan hoy su carrera y aplican la legislación vigente hasta su jubilación. Los del Ageing Report se basan en datos reales de altas y en proyecciones con la población efectiva.

El cuarto es de momento. Los sistemas están en fases distintas de maduración y de reforma: las reformas aprobadas pero no aplicadas todavía aparecen en las proyecciones, pero no en el gasto actual.

El quinto, más conceptual, es el problema de Galton, clásico en la política comparada. Los países no son observaciones independientes: se imitan, se influyen, reciben las mismas recomendaciones de la OCDE y la Comisión. Un rasgo compartido puede deberse a la difusión y no a su eficacia. Por eso la literatura sobre la transferencia de políticas (Dolowitz y Marsh, 2000) advierte contra la importación acrítica de modelos: un mecanismo que funciona en Suecia, con su consenso político y su mercado de trabajo, puede fracasar en otro contexto. El factor de sostenibilidad español de 2013, inspirado en modelos nórdicos, es un buen ejemplo.

El apartado recorre cuatro dimensiones: el gasto y la presión demográfica (4.5.1), la generosidad y la suficiencia (4.5.2), la edad y la duración de la jubilación (4.5.3) y el diseño institucional y los mecanismos de ajuste (4.5.4). Concluye con las lecciones que la comparación ofrece para España (4.5.5).

\subsubsection{Gasto y presión demográfica}

a) El nivel de gasto.

[DOC] Según el Ageing Report de 2021, el gasto público en pensiones en España fue del 12,3% del PIB en 2019, por encima de la media simple de la UE, en torno al 10%. Los datos actualizados confirman el diferencial y lo amplían en las proyecciones. El Ageing Report 2024 sitúa el gasto bruto español en el 13,1% del PIB en 2022, frente a una media de la UE en torno al 11,4%. España está en el grupo de gasto alto junto con Italia, Grecia, Francia, Austria, Portugal, Finlandia y Bélgica, y muy por encima de los países con fuerte peso del segundo pilar (Países Bajos, Dinamarca, Suecia) o con sistemas asistenciales de base (Irlanda). Según Pensions at a Glance 2025, España sería el país de la OCDE que más PIB destinaría a pensiones en 2050, en torno al 17%, por delante de Italia (15,5%), Bélgica (14,8%), Portugal (14,6%), Grecia (14%) y Francia (13,7%).

La trayectoria es tan relevante como el nivel. A escala de la UE, el Ageing Report 2024 proyecta un aumento moderado del gasto, con un pico a mediados de la década de 2030 y una ligera reducción posterior. La descomposición media de la UE entre 2022 y 2070, según el análisis del BCE, atribuye unos +6,2 puntos a la tasa de dependencia, que se compensan con −2,9 por la ratio de prestación, −1,3 por la cobertura y −1,1 por el mercado de trabajo. España se aparta de ese patrón en dos sentidos. Su presión demográfica es mucho mayor, +10,3 puntos frente a +6,2 (4.2.5). Y sus factores compensadores, sobre todo la cobertura (−0,4 frente a −1,3), son más débiles, porque España no ha vinculado la edad a la longevidad y su ajuste por cobertura depende de los incentivos voluntarios. El BCE destacó que España era el país con la mayor revisión al alza del coste del envejecimiento respecto al Ageing Report 2021, en torno a 7 puntos del PIB en 2070, y lo atribuyó a la reversión parcial de las reformas anteriores.

b) La demografía comparada.

[DOC] Según el Ageing Report de 2021, la tasa de dependencia española, del 32,1% en 2019, estaba cerca de la media de la UE (32,7%), pero en 2050 estaría muy por encima (64,7% frente a 54,8%). Los datos del Ageing Report 2024 lo confirman: la tasa española casi se duplica, del 33,3% en 2022 a cerca del 64% en 2050, mientras la media de la UE sube bastante menos. Pensions at a Glance 2025 proyecta para 2054 unas 76 personas de 65 años o más por cada 100 en edad de trabajar en España, frente a 55 en la media de la OCDE; solo Corea, Japón e Italia estarían por encima. La razón es la que analizamos en 4.2.1: España combina una de las esperanzas de vida más altas del mundo con una de las fecundidades más bajas y un baby boom tardío. Es la demografía, más que la generosidad del sistema, la que explica por qué el gasto español crecerá más que el de sus socios.

c) El segundo pilar como variable omitida.

[DOC] El documento señala que en varios países el peso de las pensiones privadas (el segundo pilar) es muy importante, lo que dificulta la comparación. Suecia, Países Bajos o Reino Unido tienen sistemas privados cuyas prestaciones superan el 3% o 4% del PIB, mientras que en España apenas alcanzan el 0,4%. Los activos de los fondos de pensiones españoles equivalen a algo más del 10% del PIB, frente a más del 60% en la OCDE. Los datos de Inverco para finales de 2025 (unos 138.000 millones de euros, algo más del 8% del PIB solo en planes de pensiones, al que habría que sumar EPSV, mutualidades y seguros colectivos) confirman el diagnóstico (4.4.4). La consecuencia metodológica es importante. Cuando se dice que España "gasta más" en pensiones que los Países Bajos o Dinamarca, se compara el gasto público español con un gasto público holandés o danés que solo cubre una parte de la renta de los jubilados. El resto lo pagan fondos privados obligatorios o cuasi obligatorios. Si se considera el conjunto de las transferencias a los mayores, la distancia se reduce sustancialmente y, en algunos casos, se invierte.

\subsubsection{Generosidad y suficiencia}

a) Las tasas de reemplazo teóricas: el dato más citado y peor comprendido.

[Probable, a partir de la prensa que recoge Pensions at a Glance 2025] Según la edición de 2025 de Pensions at a Glance, la tasa de reemplazo bruta de un trabajador con salario medio y carrera completa en España es del 80,4%, la más alta de la OCDE, frente a una media del 52%. Supera a Grecia (79,6%), que ocupaba el primer puesto en la edición de 2023. La tasa de reemplazo neta, descontados impuestos y cotizaciones, es del 86,3% en España, frente al 63,2% de la media de la OCDE, y sitúa a España en torno al séptimo lugar, por detrás de Turquía, Países Bajos, Portugal, Grecia, Luxemburgo y Austria. La tasa bruta española ha subido en las sucesivas ediciones: en torno al 72% en la de 2019, el 74% en la de 2021 y el 80% en la de 2025. Refleja en buena medida la reforma de 2021-2023, que eliminó el factor de sostenibilidad y garantizó la revalorización con el IPC.

[DOC] El documento, con datos del Ageing Report de 2021, añade que la tasa de reemplazo, definida como la media de la primera pensión sobre el último salario, supera en España el 75%, cuando la media de la UE ronda el 45%, y que la tasa de acumulación de derechos (accrual rate) española también está por encima de la media de la Unión (2,4% frente a 1,4%). La tasa de acumulación es el porcentaje de la base reguladora que se gana por cada año cotizado. En España, con 15 años se alcanza el 50% y con 36,5 o 37 años el 100%, lo que da una media cercana al 2,3-2,4% anual, entre las más altas de Europa. Es la razón técnica de la elevada tasa de reemplazo: se acumulan derechos más deprisa y se llega antes al máximo.

Las tasas de reemplazo teóricas deben interpretarse con cuatro cautelas.

La primera es que suponen carreras completas e ininterrumpidas. El trabajador tipo de la OCDE empieza a cotizar a los 22 años y trabaja sin interrupción hasta la edad normal de jubilación. (OJO: la nota 28 del documento dice "para un hombre con una carrera completa desde los 12 años". Es una errata evidente: la hipótesis de la OCDE es la entrada en el mercado de trabajo a los 22 años. Hay que corregirla en el examen.) En España, con un mercado de trabajo marcado por el paro, la temporalidad y las carreras interrumpidas, sobre todo de las mujeres, la proporción de trabajadores que alcanza esa carrera completa es menor que en otros países. La tasa de reemplazo efectiva del pensionista medio es inferior a la teórica.

La segunda es que se calculan sobre la legislación vigente proyectada. Incorporan, por tanto, la revalorización con el IPC y la ausencia de factor de sostenibilidad, pero también la subida de la base máxima por encima de la pensión máxima (4.3.3). Para salarios altos, la tasa española cae de forma pronunciada.

La tercera es que la tasa es alta, en parte, porque el primer pilar es casi el único. En los Países Bajos o Dinamarca, la tasa del sistema público es baja, pero la de los sistemas obligatorios o cuasi obligatorios, sumados, es comparable o superior a la española. Una tasa de reemplazo pública alta no implica un nivel de vida en la jubilación superior: implica que el Estado asume una parte mayor de su financiación.

La cuarta es que la tasa de reemplazo en el momento de la jubilación no dice nada de la evolución posterior. Como vimos en 4.2.3, con indexación a los precios, la pensión pierde terreno relativo frente a los salarios año tras año. El Ageing Report 2024 proyecta que la tasa de reemplazo en la jubilación caerá en España del 77% al 64% en 2070, y la ratio de prestación, del 64% al 51%.

b) La suficiencia comparada: renta relativa y pobreza de los mayores.

La contrapartida de la generosidad es una de las mejores posiciones de Europa en suficiencia. Según Pensions at a Glance 2025 y Eurostat, la tasa de pobreza de ingresos de los mayores de 65 años en España es comparable o incluso inferior a la del resto de la población, a diferencia de lo que ocurre en Corea, los países bálticos o, en menor medida, la media de la OCDE. La renta disponible media de los mayores de 65 años en relación con la del conjunto de la población se sitúa en España entre las más altas de la OCDE, cerca del 90% o por encima [Suposición en la cifra exacta]. Es el reverso del debate intergeneracional de 4.4.2: la suficiencia de las pensiones españolas es real y ha convertido a los mayores en uno de los grupos con menor riesgo de pobreza, mientras la pobreza infantil está entre las más altas de la UE.

La brecha de género en pensiones (gender pension gap), definida como la diferencia porcentual entre la pensión media de hombres y mujeres, es elevada en España, en torno al 30% según Eurostat. Está cerca de la media de la UE, pero muy por encima de la de los países nórdicos o de Europa del Este. Refleja las carreras más cortas y los salarios más bajos de las mujeres de las cohortes que hoy están jubiladas. El complemento de brecha de género (3.2 y 4.3.3) es una respuesta que otros países han abordado con créditos de cuidados (care credits) más generosos, como Alemania con la Mütterrente o Francia con las bonificaciones por hijo.

\subsubsection{Edad y duración de la jubilación}

a) La edad legal.

Solemos pensar que la edad legal de jubilación en España, de unos 65 años, es algo superior a la media de la UE (64,2 años), en ambos casos para un hombre con carrera completa. Conviene precisarlo. En España conviven dos edades legales: 67 años desde 2027 con carácter general, y 65 para quienes acrediten al menos 38 años y 6 meses de cotización (4.3.1). La OCDE toma la edad aplicable a su trabajador tipo, que al empezar a los 22 años acumula carrera suficiente para jubilarse a los 65. En la edición de 2025, la edad normal de jubilación futura de la OCDE para quienes empezaron a trabajar en 2024 se sitúa de media en torno a 66,5 o 67 años, con máximos en los países que vinculan la edad a la longevidad: Dinamarca, en torno a 74 años hacia 2070 si se mantiene el vínculo; Estonia, Finlandia, Italia, los Países Bajos, Portugal y Suecia, por encima de 67. España figura entre los países con una edad futura más baja para carreras completas, en 65 años, porque no ha vinculado la edad a la longevidad y porque la excepción de las carreras largas se aplica al trabajador tipo.

b) La edad efectiva de salida.

[DOC] El documento indica que la edad efectiva de salida del mercado laboral en España es inferior a la media de la OCDE (60,7 años frente a 63,1), con datos de Pensions at a Glance 2021, en promedio simple de hombres y mujeres. Los datos más recientes muestran una mejora. La OCDE sitúa la edad media de salida en España en torno a 62-63 años, todavía algo por debajo de su media, que ronda los 64. El Ageing Report 2024 la estima en 64 años en 2022 y la proyecta en 66,4 en 2070 (4.2.2), aunque con una definición distinta: la edad media de salida de la población activa, no del empleo. Estas divergencias entre fuentes se deben a diferencias de definición: edad de salida del mercado de trabajo, edad de acceso a la pensión o edad media de jubilación en las altas. Es otro ejemplo de por qué la comparación debe hacerse con indicadores homogéneos.

c) La duración esperada de la jubilación.

[DOC] El documento señala que el número de años esperados de jubilación en España supera los 25, frente a una media de la OCDE de 21,7. Es la consecuencia combinada de una edad efectiva de salida relativamente temprana y de una de las esperanzas de vida a los 65 años más altas del mundo, solo superada en la OCDE por Japón y, en el caso de las mujeres, en torno a la de Francia. Es el dato que mejor resume el problema del sistema español en perspectiva comparada: no se trata tanto de una pensión anual más alta como de una pensión más alta durante más tiempo. El valor actual de las prestaciones que recibe un jubilado español medio a lo largo de su vida, en relación con lo cotizado, está entre los más altos de Europa, lo que conecta con los análisis de la TIR del sistema (4.1.3).

\subsubsection{Diseño institucional y mecanismos de ajuste}

a) La tipología de los sistemas.

La comparación de diseño parte de la tipología que presentamos en el bloque I: los sistemas bismarckianos o contributivos, con pensión proporcional al salario (Alemania, Francia, Italia, España, Austria, Bélgica), y los beveridgeanos, con pensión básica uniforme o asistencial complementada por sistemas privados (Reino Unido, Irlanda, Países Bajos, Dinamarca). Dentro de los contributivos, la literatura distingue a su vez entre los sistemas de prestación definida (España, Francia en el régimen general), los de puntos (Alemania, los complementarios franceses Agirc-Arrco, Eslovaquia) y los de cuentas nocionales (Suecia, Italia, Letonia, Polonia, Noruega). [DOC] En cuanto a la revalorización, el documento recoge que los países de la OCDE suelen indexar las pensiones a los precios, a los salarios o a una combinación de ambos, y que España lo hace según la inflación. La indexación a los precios es la más frecuente. Algunos países aplican fórmulas mixtas: Suiza, con un índice mixto del 50% de salarios y 50% de precios; Finlandia, con 80% de precios y 20% de salarios; y Reino Unido, con el triple lock.

b) Los mecanismos de ajuste automático: un mapa comparado.

[DOC] El documento recoge que varios países de la UE tienen mecanismos automáticos, y que en la mayor parte de los que aplican un factor de sostenibilidad la variable que genera el ajuste, en la pensión inicial, la edad de jubilación u otros parámetros, es el aumento de la esperanza de vida. En España, el sistema cuenta con un mecanismo semiautomático, el MEI, llamado así porque requiere aprobación parlamentaria. La OCDE (Pensions at a Glance 2021, cap. 2, y ediciones posteriores) contaba en torno a dos tercios de sus miembros con algún mecanismo automático. Pueden clasificarse en cuatro tipos, cada uno ilustrado con un caso.

El vínculo de la edad a la esperanza de vida (Dinamarca). Dinamarca aprobó en 2006, con un amplio consenso, vincular la edad legal a la esperanza de vida a los 60 años, con revisiones quinquenales y quince años de preaviso. La edad legal pasará a 68 años en 2030 y a 69 en 2035, y el Parlamento aprobó en 2025 la subida a 70 años en 2040. Es el caso más avanzado del mundo. Demuestra que el vínculo puede ser políticamente sostenible con consenso, aunque ha generado un debate creciente y fórmulas de jubilación anticipada para trabajadores con carreras largas (la Arne-pension de 2021). Pendiente de verificar la ley danesa de 2025.

El factor de sostenibilidad en la revalorización (Alemania). Alemania introdujo en 2004 un factor de sostenibilidad (Nachhaltigkeitsfaktor) que reduce la revalorización de las pensiones cuando empeora la relación entre pensionistas y cotizantes, junto con un límite máximo al tipo de cotización y un nivel mínimo de las pensiones (Haltelinie). La experiencia alemana muestra la misma fragilidad política que la española: el factor se ha suspendido o compensado en varias ocasiones con cláusulas de protección, y la Haltelinie del 48% del salario se ha prorrogado por ley. En 2025-2026 el Gobierno de coalición la ha extendido hasta 2031 y ha aprobado un "paquete de pensiones" que incluye la Aktivrente, una exención fiscal parcial para quienes sigan trabajando después de la edad legal. Pendiente de verificar el detalle de este paquete.

El mecanismo de equilibrio financiero en un sistema de cuentas nocionales (Suecia). El freno sueco (bromsen, 4.4.3) reduce la revalorización si el balance actuarial del sistema es negativo. Se activó en 2010, 2011 y 2014. El Gobierno compensó en parte los recortes con rebajas fiscales a los pensionistas, lo que muestra que ni siquiera el sistema considerado más "despolitizado" escapa a la economía política. La Agencia Sueca de Pensiones publica cada año un informe con el balance actuarial del sistema, que España no tiene (4.1.3).

El ajuste por la vía de los ingresos (España). El MEI y la cláusula de salvaguarda españoles son, en el panorama comparado, una rareza. Son de los pocos mecanismos automáticos que ajustan exclusivamente por la vía de las cotizaciones y cuya activación depende de una evaluación trienal que, como vimos en 4.3.4, tiene fuertes debilidades metodológicas. Los demás mecanismos reparten el ajuste entre cotizantes y pensionistas, o lo concentran en la edad o en la prestación.

c) Las experiencias recientes que enseñan algo.

Tres experiencias recientes son especialmente útiles como ganchos comparados.

Francia ilustra la dificultad de mover la edad de jubilación sin consenso. La reforma de 2023 elevó la edad legal de 62 a 64 años, se aprobó sin votación mediante el art. 49.3 de la Constitución y desencadenó meses de protestas. En el otoño de 2025, en plena crisis política y como condición para la supervivencia del Gobierno, el Ejecutivo anunció la suspensión del calendario de subida de la edad hasta después de las elecciones presidenciales de 2027. Pendiente de verificar la situación exacta de la suspensión y su tramitación. Es un eco del caso español de 2013-2018.

Los Países Bajos han emprendido la transición más ambiciosa de Europa en el segundo pilar. La Ley sobre el Futuro de las Pensiones (Wet toekomst pensioenen, en vigor desde julio de 2023) transforma sus fondos de empleo, de prestación definida, en un sistema de aportación definida colectiva, con reparto colectivo de riesgos, que debe completarse antes de 2028. Es la referencia actual del debate sobre cómo reforzar el segundo pilar sin trasladar todo el riesgo al individuo.

El Reino Unido ofrece dos lecciones opuestas. La inscripción automática (4.4.4) es el ejemplo más exitoso de extensión del segundo pilar mediante la economía conductual. El triple lock, que revaloriza la pensión estatal con el mayor de tres valores (inflación, crecimiento salarial o un 2,5%), es el ejemplo más claro de cómo la economía política de la edad puede producir una regla de indexación que, a largo plazo, aumenta sistemáticamente el gasto. La Oficina de Responsabilidad Presupuestaria británica ha señalado repetidamente su coste creciente, y ningún partido se atreve a eliminarla.

[DOC] Chile, que el documento trata en su nota 30 como el modelo más radicalmente opuesto al español, fue durante décadas la referencia de quienes defendían la capitalización frente al reparto. Implantado en 1980 bajo la dictadura de Pinochet por impulso del ministro José Piñera, sustituyó el reparto por cuentas individuales gestionadas por administradoras privadas (AFP). La insuficiencia de las pensiones finalmente percibidas por buena parte de los cotizantes provocó una fuerte contestación social y reformas sustanciales: el pilar solidario de 2008 y, [Probable], la reforma de 2025, que introdujo una cotización a cargo del empleador destinada en parte a un seguro social con elementos de reparto y solidaridad (4.4.4).

\subsubsection{Lecciones de la comparación para España}

La comparación internacional permite situar con precisión la singularidad española. España combina cuatro rasgos que, juntos, no se dan en ningún otro país de la OCDE. Tiene una de las presiones demográficas más intensas del mundo desarrollado. Tiene un sistema público contributivo que es casi la única fuente de renta en la jubilación, con tasas de reemplazo teóricas entre las más altas y un segundo pilar residual. Tiene una edad efectiva de salida algo inferior a la media y una esperanza de vida entre las más altas, lo que se traduce en una de las jubilaciones más largas. Y ha elegido, tras la reversión de 2013, un mecanismo de ajuste que actúa sobre los ingresos y no sobre la edad ni la prestación.

La lectura que se desprende no es unívoca, y conviene exponer las dos.

La lectura crítica subraya que España es un caso atípico de generosidad no ajustada: sin vínculo a la longevidad, con la tasa de reemplazo más alta de la OCDE y con el mayor aumento proyectado del gasto. Recuerda que los países de su entorno, con sistemas parecidos (Italia, Portugal, Grecia, Finlandia), introdujeron mecanismos automáticos sobre la edad o la prestación, y concluye que el sistema español tendrá que acabar haciéndolo.

La lectura favorable subraya que ese mismo sistema logra una de las menores tasas de pobreza de los mayores de Europa y que la comparación de gasto público está sesgada por la ausencia del segundo pilar. Añade que la experiencia comparada muestra la fragilidad política de los mecanismos que recaen sobre las prestaciones, como las suspensiones en Alemania, Suecia y Francia, y que la vía española de los ingresos es una elección legítima entre las posibles, siempre que se financie de forma creíble.

Ambas lecturas comparten una conclusión: la singularidad española no está tanto en el diseño como en la combinación de una demografía extrema con un único pilar. Los países que afrontan presiones demográficas comparables (Italia, Japón, Corea) lo hacen con ajustes intensos de la prestación o con tasas de pobreza de los mayores muy elevadas. Los que mantienen tasas de reemplazo altas (Países Bajos, Dinamarca) lo hacen con un segundo pilar maduro. España intenta mantener a la vez tasas de reemplazo altas, un único pilar público y una demografía muy desfavorable. La comparación sugiere que esa combinación solo es sostenible con aumentos continuados de los ingresos públicos, lo que en la práctica significa una presión fiscal creciente o un recurso creciente a la deuda, o con ajustes que, antes o después, afectarán a alguno de los otros dos elementos.

Con este apartado se cierra el bloque IV y el cuerpo del tema. Quedan la introducción general del tema y la conclusión (recapitulación, extensiones y relación con otras partes del temario, opinión e idea final), que según el plan acordado se redactan al final.

















































































\clearpage
\section{Conclusión} 




\subsection{Recapitulación}


\subsection{Extensiones y relación con otras partes del Temario}



\subsection{Opinión}

\subsubsection{Idea final (Salida o cierra)}

\clearpage
\pagenumbering{gobble}
\MyBibliography


- Akerlof, G. (1970), "The Market for 'Lemons': Quality Uncertainty and the Market Mechanism", Quarterly Journal of Economics, 84(3).
- Arrow, K. (1963), "Uncertainty and the Welfare Economics of Medical Care", American Economic Review, 53(5).
- Arts, W. y Gelissen, J. (2002), "Three Worlds of Welfare Capitalism or More? A State-of-the-Art Report", Journal of European Social Policy, 12(2).
- Barr, N. (2001), The Welfare State as Piggy Bank, Oxford University Press.
- Barr, N. y Diamond, P. (2008), Reforming Pensions: Principles and Policy Choices, Oxford University Press.
- Bonoli, G. (1997), "Classifying Welfare States: A Two-Dimension Approach", Journal of Social Policy, 26(3).
- Buchanan, J. (1975), "The Samaritan's Dilemma", en Phelps, E. (ed.), Altruism, Morality and Economic Theory, Russell Sage Foundation.
- Castles, F. y Mitchell, D. (1993), "Worlds of Welfare and Families of Nations", en Castles, F. (ed.), Families of Nations, Dartmouth.
- Chetty, R. y Finkelstein, A. (2013), "Social Insurance: Connecting Theory to Evidence", en Handbook of Public Economics, vol. 5, Elsevier.
- Chetty, R. et al. (2016), "The Association Between Income and Life Expectancy in the United States, 2001-2014", JAMA, 315(16).
- Coate, S. (1995), "Altruism, the Samaritan's Dilemma, and Government Transfer Policy", American Economic Review, 85(1).
- Diamond, P. (1977), "A Framework for Social Security Analysis", Journal of Public Economics, 8(3).
- Dworkin, R. (1981), "What is Equality? Part 2: Equality of Resources", Philosophy & Public Affairs, 10(4).
- Emmenegger, P., Häusermann, S., Palier, B. y Seeleib-Kaiser, M. (eds.) (2012), The Age of Dualization, Oxford University Press.
- Esping-Andersen, G. (1990), The Three Worlds of Welfare Capitalism, Polity Press.
- Esping-Andersen, G. (1999), Social Foundations of Postindustrial Economies, Oxford University Press.
- Estevez-Abe, M., Iversen, T. y Soskice, D. (2001), "Social Protection and the Formation of Skills", en Hall, P. y Soskice, D. (eds.), Varieties of Capitalism, Oxford University Press.
- Falkingham, J. y Hills, J. (eds.) (1995), The Dynamic of Welfare: The Welfare State and the Life Cycle, Prentice Hall/Harvester Wheatsheaf.
- Feldstein, M. (1985), "The Optimal Level of Social Security Benefits", Quarterly Journal of Economics, 100(2).
- Ferrera, M. (1996), "The 'Southern Model' of Welfare in Social Europe", Journal of European Social Policy, 6(1).
- Finkelstein, A. y Poterba, J. (2004), "Adverse Selection in Insurance Markets: Policyholder Evidence from the U.K. Annuity Market", Journal of Political Economy, 112(1).
- Friedman, M. (1962), Capitalism and Freedom, University of Chicago Press.
- Gal, J. (2010), "Is There an Extended Family of Mediterranean Welfare States?", Journal of European Social Policy, 20(4).
- Gordon, R. y Varian, H. (1988), "Intergenerational Risk Sharing", Journal of Public Economics, 37(2).
- Hall, P. y Soskice, D. (eds.) (2001), Varieties of Capitalism, Oxford University Press.
- Harsanyi, J. (1955), "Cardinal Welfare, Individualistic Ethics, and Interpersonal Comparisons of Utility", Journal of Political Economy, 63(4).
- Hayek, F. (1960), The Constitution of Liberty, University of Chicago Press.
- Kenworthy, L. (2011), Progress for the Poor, Oxford University Press.
- Korpi, W. y Palme, J. (1998), "The Paradox of Redistribution and Strategies of Equality", American Sociological Review, 63(5).
- Laibson, D. (1997), "Golden Eggs and Hyperbolic Discounting", Quarterly Journal of Economics, 112(2).
- Lewis, J. (1992), "Gender and the Development of Welfare Regimes", Journal of European Social Policy, 2(3).
- Lindert, P. (2004), Growing Public: Social Spending and Economic Growth since the Eighteenth Century, Cambridge University Press.
- Madrian, B. y Shea, D. (2001), "The Power of Suggestion: Inertia in 401(k) Participation and Savings Behavior", Quarterly Journal of Economics, 116(4).
- Marshall, T. H. (1950), Citizenship and Social Class, Cambridge University Press.
- Marx, I., Salanauskaite, L. y Verbist, G. (2013), "The Paradox of Redistribution Revisited: And That It May Rest in Peace?", IZA Discussion Paper 7414.
- Merton, R. (1983), "On the Role of Social Security as a Means for Efficient Risk Sharing in an Economy Where Human Capital Is Not Tradable", en Bodie, Z. y Shoven, J. (eds.), Financial Aspects of the United States Pension System, NBER/University of Chicago Press.
- Mitchell, O., Poterba, J., Warshawsky, M. y Brown, J. (1999), "New Evidence on the Money's Worth of Individual Annuities", American Economic Review, 89(5).
- Moreno, L. (2000), Ciudadanos precarios: la "última red" de protección social, Ariel.
- Moreno, L. (2001), "La 'vía media' española del modelo de bienestar mediterráneo", Papers. Revista de Sociologia, 63/64.
- Musgrave, R. (1959), The Theory of Public Finance, McGraw-Hill.
- Nozick, R. (1974), Anarchy, State, and Utopia, Basic Books.
- Okun, A. (1975), Equality and Efficiency: The Big Tradeoff, Brookings Institution.
- Orloff, A. (1993), "Gender and the Social Rights of Citizenship", American Sociological Review, 58(3).
- Palier, B. (ed.) (2010), A Long Goodbye to Bismarck? The Politics of Welfare Reform in Continental Europe, Amsterdam University Press.
- Pauly, M. (1968), "The Economics of Moral Hazard: Comment", American Economic Review, 58(3).
- Rawls, J. (1971), A Theory of Justice, Harvard University Press.
- Rothschild, M. y Stiglitz, J. (1976), "Equilibrium in Competitive Insurance Markets: An Essay on the Economics of Imperfect Information", Quarterly Journal of Economics, 90(4).
- Rueda, D. (2007), Social Democracy Inside Out, Oxford University Press.
- Sen, A. (1985), Commodities and Capabilities, North-Holland; (1999), Development as Freedom, Oxford University Press.
- Shiller, R. (1999), "Social Security and Institutions for Intergenerational, Intragenerational, and International Risk-Sharing", Carnegie-Rochester Conference Series on Public Policy, 50.
- Sinn, H.-W. (1995), "A Theory of the Welfare State", Scandinavian Journal of Economics, 97(4).
- Thaler, R. y Benartzi, S. (2004), "Save More Tomorrow: Using Behavioral Economics to Increase Employee Saving", Journal of Political Economy, 112(S1).
- Thaler, R. y Sunstein, C. (2008), Nudge, Yale University Press.
- Titmuss, R. (1974), Social Policy: An Introduction, Allen & Unwin.
- Van Parijs, P. (1995), Real Freedom for All, Oxford University Press.
- Yaari, M. (1965), "Uncertain Lifetime, Life Insurance, and the Theory of the Consumer", Review of Economic Studies, 32(2).
- [Ministerio de Hacienda (Chile): nueva cotización del empleador](https://www.hacienda.cl/noticias-y-eventos/noticias/implementacion-de-la-reforma-previsional-nueva-cotizacion-del-empleador-regira)
- [Superintendencia de Pensiones (Chile): Ley 21.735](https://www.spensiones.cl/portal/institucional/594/w3-article-16483.html)
- [Gobierno Federal de Alemania: paquete de pensiones de 2025](https://www.bundesregierung.de/breg-en/news/pension-package-2025-2397954)
- [Infobae: Francia suspende la reforma de pensiones hasta 2028](https://www.infobae.com/espana/2025/11/12/macron-logra-frenar-temporalmente-la-reforma-de-las-pensiones-la-asamblea-de-francia-respalda-la-suspension-hasta-2028/)
- [El Debate: el Congreso tumba el decreto ómnibus](https://www.eldebate.com/espana/20260127/pleno-congreso-preve-rechazo-revalorizacion-pensiones_378645.html)
- [Iberley: convalidación del RDL 3/2026](https://www.iberley.es/noticias/el-congreso-convalida-revalorizacion-las-pensiones-mientras-no-convalida-rd-ley-escudo-social-36131)
- [BOE: Real Decreto-ley 3/2026](https://www.boe.es/buscar/act.php?id=BOE-A-2026-2548)
- [Consejo General del Trabajo Social: aprobación definitiva de la reforma de dependencia](https://www.cgtrabajosocial.es/noticias/el-congreso-aprueba-definitivamente-la-reforma-de-las-leyes-de-dependencia-y-discapacidad/15086/view)
- Banco Mundial (1994), Averting the Old Age Crisis: Policies to Protect the Old and Promote Growth, Oxford University Press.
- Beveridge, W. (1942), Social Insurance and Allied Services, Cmd. 6404, HMSO.
- Beveridge, W. (1944), Full Employment in a Free Society, Allen & Unwin.
- Bonoli, G. (2005), "The Politics of the New Social Policies: Providing Coverage against New Social Risks in Mature Welfare States", Policy & Politics, 33(3).
- Bonoli, G., Cantillon, B. y Van Lancker, W. (2017), "Social Investment and the Matthew Effect", en Hemerijck, A. (ed.), The Uses of Social Investment, Oxford University Press.
- Cantillon, B. (2011), "The Paradox of the Social Investment State: Growth, Employment and Poverty in the Lisbon Era", Journal of European Social Policy, 21(5).
- Esping-Andersen, G., Gallie, D., Hemerijck, A. y Myles, J. (2002), Why We Need a New Welfare State, Oxford University Press.
- Feldstein, M. (1974), "Social Security, Induced Retirement, and Aggregate Capital Accumulation", Journal of Political Economy, 82(5).
- Hemerijck, A. (2013), Changing Welfare States, Oxford University Press.
- Hemerijck, A. (ed.) (2017), The Uses of Social Investment, Oxford University Press.
- Holzmann, R. y Hinz, R. (2005), Old-Age Income Support in the 21st Century: An International Perspective on Pension Systems and Reform, Banco Mundial.
- Holzmann, R. y Palmer, E. (eds.) (2006), Pension Reform: Issues and Prospects for Non-Financial Defined Contribution (NDC) Schemes, Banco Mundial.
- Morel, N., Palier, B. y Palme, J. (eds.) (2012), Towards a Social Investment Welfare State? Ideas, Policies and Challenges, Policy Press.
- OCDE (1981), The Welfare State in Crisis, OECD Publishing.
- OIT (1944), Declaración relativa a los fines y objetivos de la Organización Internacional del Trabajo (Declaración de Filadelfia).
- OIT (1952), Convenio n.º 102 sobre la seguridad social (norma mínima).
- Orszag, P. y Stiglitz, J. (1999), "Rethinking Pension Reform: Ten Myths About Social Security Systems", ponencia presentada en la conferencia New Ideas About Old Age Security, Banco Mundial.
- Ortiz, I., Durán-Valverde, F., Urban, S. y Wodsak, V. (eds.) (2018), Reversing Pension Privatizations: Rebuilding Public Pension Systems in Eastern Europe and Latin America, OIT.
- Pierson, P. (1994), Dismantling the Welfare State? Reagan, Thatcher and the Politics of Retrenchment, Cambridge University Press.
- Pierson, P. (1996), "The New Politics of the Welfare State", World Politics, 48(2).
- Polanyi, K. (1944), The Great Transformation, Farrar & Rinehart.
- Taylor-Gooby, P. (ed.) (2004), New Risks, New Welfare: The Transformation of the European Welfare State, Oxford University Press.
- Weaver, R. K. (1986), "The Politics of Blame Avoidance", Journal of Public Policy, 6(4).
- León XIII (1891), encíclica Rerum Novarum; Pío XI (1931), encíclica Quadragesimo Anno.
- Unión Europea: Reglamento (UE) 2020/672 del Consejo, de 19 de mayo de 2020, relativo a la creación de un instrumento europeo de apoyo temporal para atenuar los riesgos de desempleo en una emergencia (SURE).
- Chile: Decreto Ley 3.500 (1980); Ley 21.419 (2022), Pensión Garantizada Universal; Ley 21.735, de 26 de marzo de 2025.
- España (normas no citadas previamente):
  - Real Decreto de 5 de diciembre de 1883 (Comisión de Reformas Sociales);
  - Ley de 30 de enero de 1900, de Accidentes de Trabajo;
  - Ley de 27 de febrero de 1908 (Instituto Nacional de Previsión);
  - Real Decreto de 11 de marzo de 1919 (Retiro Obrero Obligatorio);
  - Ley de 14 de diciembre de 1942 (Seguro Obligatorio de Enfermedad);
  - Ley 62/1961 (Seguro Nacional de Desempleo);
  - Ley 193/1963, de 28 de diciembre, de Bases de la Seguridad Social;
  - Decreto 907/1966, de 21 de abril;
  - Ley 24/1972, de 21 de junio;
  - Decreto 2065/1974;
  - Real Decreto-ley 36/1978;
  - Ley 31/1984, de protección por desempleo;
  - Real Decreto Legislativo 1/1994, de 20 de junio;
  - Real Decreto-ley 8/2010;
  - Real Decreto-ley 28/2012, de 30 de noviembre;
  - STC 49/2015, de 5 de marzo;
  - Real Decreto-ley 5/2013;
  - Ley 6/2018, de Presupuestos Generales del Estado para 2018;
  - Real Decreto-ley 20/2020, de 29 de mayo, y Ley 19/2021, de 20 de diciembre (IMV);
  - Real Decreto-ley 3/2021;
  - Ley 12/2022, de 30 de junio;
  - Real Decreto-ley 13/2022, de 26 de julio;
  - Real Decreto-ley 2/2024, de 21 de mayo;
  - Real Decreto-ley 11/2024, de 23 de diciembre;
  - Real Decreto-ley 16/2025, de 23 de diciembre (no convalidado);
  - Real Decreto-ley 3/2026, de 3 de febrero.
- Sentencia del TJUE de 12 de diciembre de 2019, asunto C-450/18, WA c. INSS.
- [Moncloa.com: transferencias al País Vasco en materia de Seguridad Social (julio de 2026)](https://www.moncloa.com/2026/07/16/transferencias-pais-vasco-seguridad-social-puertos-3400750/)
- [Moncloa.com: el PSE-EE cuestiona la viabilidad de la Seguridad Social vasca](https://www.moncloa.com/2026/07/05/seguridad-social-vasca-deficit-pse-ee-3394997/)
- [BOE: ratificación de la Carta Social Europea revisada](https://www.boe.es/diario_boe/txt.php?id=BOE-A-2021-9719)
- [BOE: aplicación provisional del Protocolo de reclamaciones colectivas](https://www.boe.es/diario_boe/txt.php?id=BOE-A-2021-10683)
- [BOE: ratificación del Protocolo de reclamaciones colectivas](https://www.boe.es/diario_boe/txt.php?id=BOE-A-2022-17976)
- [BOE: Recomendación del Consejo de 21 de enero de 2025 sobre el plan fiscal-estructural de España](https://www.boe.es/buscar/doc.php?id=DOUE-Z-2025-70011)
- [Servimedia: 2026 comienza con los Presupuestos de 2023 prorrogados](https://www.servimedia.es/noticias/gobierno-inicia-2026-presupuestos-2023-prorrogados-promesa-presentarlos-primer-trimestre/1412284733)

---

### Referencias nuevas

- Auerbach, A., Gokhale, J. y Kotlikoff, L. (1991), "Generational Accounts: A Meaningful Alternative to Deficit Accounting", Tax Policy and the Economy, 5.
- Burke, E. (1790), Reflections on the Revolution in France, Dodsley.
- Ferrera, M. (2005), The Boundaries of Welfare: European Integration and the New Spatial Politics of Social Protection, Oxford University Press.
- Hohaus, R. (1938), "Equity, Adequacy, and Related Factors in Old Age Security", The Record of the American Institute of Actuaries, 27.
- Oates, W. (1972), Fiscal Federalism, Harcourt Brace Jovanovich.
- Scharpf, F. (1999), Governing in Europe: Effective and Democratic?, Oxford University Press.
- Summers, L. (1989), "Some Simple Economics of Mandated Benefits", American Economic Review, 79(2).
- Comité de Derechos Económicos, Sociales y Culturales de Naciones Unidas (2007), Observación General n.º 19: El derecho a la seguridad social (art. 9), E/C.12/GC/19.

Constitución y legislación española:
- Constitución Española, arts. 9.3, 10.2, 14, 31.3, 39, 41, 43, 49 (reformado el 15 de febrero de 2024), 50, 53.3, 86, 96, 129.1, 135 (reformado en 2011), 148.1.20ª, 149.1.1ª y 149.1.17ª.
- Real Decreto Legislativo 8/2015, de 30 de octubre, por el que se aprueba el texto refundido de la Ley General de la Seguridad Social, arts. 2, 3 y 42.
- Ley 47/2003, de 26 de noviembre, General Presupuestaria.
- Ley Orgánica 2/2012, de 27 de abril, de Estabilidad Presupuestaria y Sostenibilidad Financiera.
- Texto refundido de la Ley de Clases Pasivas del Estado (Real Decreto Legislativo 670/1987).
- Texto refundido de la Ley de Regulación de los Planes y Fondos de Pensiones (Real Decreto Legislativo 1/2002).
- Estatuto de Autonomía del País Vasco (Ley Orgánica 3/1979), art. 18.

Jurisprudencia constitucional:
- STC 103/1983 (viudedad de los varones).
- STC 65/1987 (la Seguridad Social como función del Estado).
- STC 134/1987 (topes de las pensiones).
- STC 124/1989 (régimen económico y caja única).
- STC 37/1994 (garantía institucional del art. 41).
- STC 185/1995 (prestaciones patrimoniales de carácter público).
- STC 239/2002 (complementos autonómicos a las pensiones no contributivas).
- STC 61/2013 (trabajadores a tiempo parcial).
- STC 91/2019 (coeficiente de parcialidad).

TJUE:
- Sentencia de 22 de noviembre de 2012, Elbal Moreno, asunto C-385/11.
- Sentencia de 8 de mayo de 2019, Villar Láiz, asunto C-161/18.

Derecho internacional:
- Pacto Internacional de Derechos Económicos, Sociales y Culturales (1966), art. 9.
- Código Europeo de Seguridad Social (Consejo de Europa, 1964).
- Carta Social Europea (1961) y Carta Social Europea revisada (1996): Instrumento de ratificación publicado en el BOE en 2021 (BOE-A-2021-9719).
- Protocolo Adicional de reclamaciones colectivas (1995): aplicación provisional en 2021 (BOE-A-2021-10683) e instrumento de ratificación de 2022 (BOE-A-2022-17976).
- Convenio Multilateral Iberoamericano de Seguridad Social (2007).

Derecho de la Unión Europea:
- Tratado de Funcionamiento de la Unión Europea, arts. 48 y 153; Carta de los Derechos Fundamentales de la Unión Europea, art. 34.
- Reglamento (CE) 883/2004 y Reglamento (CE) 987/2009, sobre coordinación de los sistemas de seguridad social.
- Directiva 79/7/CEE del Consejo, de 19 de diciembre de 1978, sobre igualdad de trato en materia de seguridad social.
- Pilar Europeo de Derechos Sociales (proclamación interinstitucional, Gotemburgo, 17 de noviembre de 2017).
- Recomendación del Consejo de 8 de noviembre de 2019, relativa al acceso a la protección social para los trabajadores por cuenta ajena y por cuenta propia.
- Recomendación del Consejo de 30 de enero de 2023, sobre una renta mínima adecuada que procure la inclusión activa.
- Directiva (UE) 2016/2341 (IORP II) y Reglamento (UE) 2019/1238 (PEPP).
- Reglamento (UE) 2021/241, por el que se establece el Mecanismo de Recuperación y Resiliencia.
- Reglamento (UE) 2024/1263, relativo a la coordinación eficaz de las políticas económicas y la supervisión presupuestaria multilateral.
- Recomendación del Consejo de 21 de enero de 2025, por la que se aprueba el plan fiscal-estructural nacional a medio plazo de España.
- Calmfors, L. y Wren-Lewis, S. (2011), "What Should Fiscal Councils Do?", Economic Policy, 26(68).
- Christensen, T. y Lægreid, P. (2007), "The Whole-of-Government Approach to Public Sector Reform", Public Administration Review, 67(6).
- Dunleavy, P. (1991), Democracy, Bureaucracy and Public Choice, Harvester Wheatsheaf.
- Feldstein, M. (1976), "Temporary Layoffs in the Theory of Unemployment", Journal of Political Economy, 84(5).
- Hood, C. (1991), "A Public Management for All Seasons?", Public Administration, 69(1).
- Kydland, F. y Prescott, E. (1977), "Rules Rather than Discretion: The Inconsistency of Optimal Plans", Journal of Political Economy, 85(3).
- Niskanen, W. (1971), Bureaucracy and Representative Government, Aldine-Atherton.
- Pollitt, C. y Bouckaert, G. (2011), Public Management Reform: A Comparative Analysis, 3.ª ed., Oxford University Press.
- Rogoff, K. (1985), "The Optimal Degree of Commitment to an Intermediate Monetary Target", Quarterly Journal of Economics, 100(4).
- Topel, R. (1983), "On Layoffs and Unemployment Insurance", American Economic Review, 73(4).
- Eurostat, ESSPROS Manual and User Guidelines (European System of Integrated Social PROtection Statistics).
- Naciones Unidas, Classification of the Functions of Government (COFOG).
- Reglamento (UE) 549/2013, relativo al Sistema Europeo de Cuentas Nacionales y Regionales de la Unión Europea (SEC 2010), sector S.13 y subsector S.1314.

Normativa española:
- Constitución Española, art. 136.
- Ley 27/2011, de 1 de agosto, disposición adicional séptima (agencia estatal de la Seguridad Social).
- Ley 35/2014, de 26 de diciembre, por la que se modifica el texto refundido de la Ley General de la Seguridad Social en relación con el régimen jurídico de las Mutuas de Accidentes de Trabajo y Enfermedades Profesionales de la Seguridad Social.
- Ley 40/2015, de 1 de octubre, de Régimen Jurídico del Sector Público.
- Ley Orgánica 6/2013, de 14 de noviembre, de creación de la Autoridad Independiente de Responsabilidad Fiscal.
- Real Decreto 2/2020, de 12 de enero, por el que se reestructuran los departamentos ministeriales.
- Instrucción de la Consellería de Sanidade de Galicia sobre gestión clínica de la incapacidad temporal, modificada en mayo de 2026 (Diario Oficial de Galicia de 20 de mayo de 2026).

Pendiente de verificar antes de fijar el texto:
- el número del art. del texto refundido que define las entidades gestoras;
- el número actual de mutuas colaboradoras;
- la norma de 2017 que regula el sistema de incentivos por baja siniestralidad;
- la fecha exacta de traspaso de la gestión de Clases Pasivas al Ministerio de Inclusión;
- la ubicación contable del mutualismo administrativo en la contabilidad nacional;
- las observaciones del Tribunal de Cuentas sobre el patrimonio neto negativo del sistema;
- el alcance exacto de las discrepancias de 2025 entre la AIReF y el Ministerio.

Queda cerrado el bloque I. ¿Pasamos al bloque II (financiación), empezando por 2.1 (teoría: reparto frente a capitalización, cotizaciones frente a impuestos)?

Fuentes de actualidad usadas:
- [Redacción Médica: MUFACE confirma a Adeslas y Asisa como únicas aseguradoras del concierto 2025-2027](https://www.redaccionmedica.com/secciones/sanidad-hoy/muface-confirma-a-adeslas-y-asisa-como-unicas-aseguradoras-de-su-concierto-2975)
- [65 y más: MUFACE adjudica el concierto 2025-2027](https://www.65ymas.com/economia/muface-aprueba-ofertas-asisa-adeslas-adjudica-concierto-2025-2027_69181_102.html)
- [Laboral Pensiones: ¿Por fin la Agencia Estatal de la Seguridad Social?](https://laboralpensiones.com/por-fin-la-agencia-estatal-de-la-seguridad-social/)
- [elEconomista: Escrivá acelera la creación de la Agencia Estatal de la Seguridad Social](https://www.eleconomista.es/economia/noticias/11929075/09/22/Escriva-acelera-la-creacion-de-la-Agencia-Estatal-de-la-Seguridad-Social.html)
- [Galiciapress: el médico de la mutua podrá dar altas, dicta la Xunta](https://www.galiciapress.es/articulo/economia/2026-05-28/5899086-medico-mutua-podra-dar-altas-dicta-xunta-ante-protestas-trabajadores)
- [elEconomista: CCOO plantea que las mutuas traten las lesiones físicas](https://www.eleconomista.es/economia/noticias/13843491/03/26/ccoo-plantea-que-las-mutuas-traten-las-lesiones-fisicas-y-que-seguridad-social-abone-el-coste.html)
- [El Blog Salmón: la Agencia Española de Empleo, bloqueada](https://www.elblogsalmon.com/economia/ministerio-trabajo-mantiene-bloqueada-agencia-espanola-empleo-dos-anos-despues-aprobar-reforma)

- Aaron, H. (1966), "The Social Insurance Paradox", Canadian Journal of Economics and Political Science, 32(3).
- Abel, A., Mankiw, N. G., Summers, L. y Zeckhauser, R. (1989), "Assessing Dynamic Efficiency: Theory and Evidence", Review of Economic Studies, 56(1).
- Barr, N. (2002), "Reforming Pensions: Myths, Truths, and Policy Choices", International Social Security Review, 55(2).
- Barro, R. (1974), "Are Government Bonds Net Wealth?", Journal of Political Economy, 82(6).
- Blanchard, O. (2019), "Public Debt and Low Interest Rates", American Economic Review, 109(4).
- Breyer, F. (1989), "On the Intergenerational Pareto Efficiency of Pay-as-you-go Financed Pension Systems", Journal of Institutional and Theoretical Economics, 145(4).
- Buchanan, J. (1963), "The Economics of Earmarked Taxes", Journal of Political Economy, 71(5).
- Chetty, R., Looney, A. y Kroft, K. (2009), "Salience and Taxation: Theory and Evidence", American Economic Review, 99(4).
- Conde-Ruiz, J. I. y González, C. I. (2016), "From Bismarck to Beveridge: The Other Pension Reform in Spain", SERIEs, 7(4).
- Daveri, F. y Tabellini, G. (2000), "Unemployment, Growth and Taxation in Industrial Countries", Economic Policy, 15(30).
- De Mooij, R. y Keen, M. (2013), "'Fiscal Devaluation' and Fiscal Consolidation: The VAT in Troubled Times", en Alesina, A. y Giavazzi, F. (eds.), Fiscal Policy after the Financial Crisis, University of Chicago Press/NBER.
- Diamond, P. (1965), "National Debt in a Neoclassical Growth Model", American Economic Review, 55(5).
- Diamond, P. (1998), "The Economics of Social Security Reform", NBER Working Paper 6719.
- Disney, R. (2004), "Are Contributions to Public Pension Programmes a Tax on Employment?", Economic Policy, 19(39).
- Farhi, E., Gopinath, G. e Itskhoki, O. (2014), "Fiscal Devaluations", Review of Economic Studies, 81(2).
- Geanakoplos, J., Mitchell, O. y Zeldes, S. (1998), "Would a Privatized Social Security System Really Pay a Higher Rate of Return?", en Arnold, D., Graetz, M. y Munnell, A. (eds.), Framing the Social Security Debate, Brookings Institution.
- Gruber, J. (1997), "The Incidence of Payroll Taxation: Evidence from Chile", Journal of Labor Economics, 15(3).
- Leimer, D. y Lesnoy, S. (1982), "Social Security and Private Saving: New Time-Series Evidence", Journal of Political Economy, 90(3).
- Mankiw, N. G. y Weil, D. (1989), "The Baby Boom, the Baby Bust, and the Housing Market", Regional Science and Urban Economics, 19(2).
- Merton, R., Dutta, J., Kapur, S. y Orszag, P. (2000), "A Portfolio Approach to the Optimal Funding of Pensions", Economics Letters, 69(2).
- Poterba, J. (2001), "Demographic Structure and Asset Returns", Review of Economics and Statistics, 83(4).
- Puviani, A. (1903), Teoria della illusione finanziaria, Sandron.
- Saez, E., Matsaganis, M. y Tsakloglou, P. (2012), "Earnings Determination and Taxes: Evidence from a Cohort-Based Payroll Tax Reform in Greece", Quarterly Journal of Economics, 127(1).
- Saez, E., Schoefer, B. y Seim, D. (2019), "Payroll Taxes, Firm Behavior, and Rent Sharing: Evidence from a Young Workers' Tax Cut in Sweden", American Economic Review, 109(5).
- Samuelson, P. (1958), "An Exact Consumption-Loan Model of Interest with or without the Social Contrivance of Money", Journal of Political Economy, 66(6).
- Settergren, O. y Mikula, B. (2005), "The Rate of Return of Pay-As-You-Go Pension Systems: A More Exact Consumption-Loan Model of Interest", Journal of Pension Economics and Finance, 4(2).
- Sinn, H.-W. (2000), "Why a Funded Pension System Is Useful and Why It Is Not Useful", International Tax and Public Finance, 7(4).
- Whitehouse, E. (2000), "Administrative Charges for Funded Pensions: An International Comparison and Assessment", Social Protection Discussion Paper 0016, Banco Mundial.
- AIReF (2019), Evaluación del gasto público 2018: Estudio sobre los incentivos a la contratación, Spending Review.
- OCDE, Taxing Wages (edición anual).
- Reglamento (UE) 549/2013 (SEC 2010), tabla complementaria sobre derechos de pensión devengados (tabla 29).

Pendiente de verificar antes de fijar el texto:
- la cita exacta de Samuelson en Newsweek (1967);
- el valor de la deuda implícita de pensiones de España en la última tabla 29 publicada;
- la cifra de la cuña fiscal española en la última edición de Taxing Wages;
- los valores de 2010 y 2022 de la base máxima y la pensión máxima, para fijar la corrección de la nota del documento sobre la reforma silenciosa;
- las conclusiones exactas de la revisión del gasto de la AIReF sobre incentivos a la contratación.
### Referencias nuevas

- García Díaz, M. Á. (2024), "La situación financiera del componente contributivo de las pensiones públicas en 2023", Apuntes/Estudios sobre la Economía Española, Fedea.
- Kleven, H., Knudsen, M., Kreiner, C., Pedersen, S. y Saez, E. (2011), "Unwilling or Unable to Cheat? Evidence from a Tax Audit Experiment in Denmark", Econometrica, 79(3).
- Tribunal de Cuentas (2026), Declaración sobre la Cuenta General del Estado del ejercicio 2024, con las observaciones sobre la financiación de la Seguridad Social mediante préstamos.
- Ministerio de Inclusión, Seguridad Social y Migraciones, Informe a las Cortes Generales sobre el Fondo de Reserva de la Seguridad Social (edición anual), y notas de prensa sobre su evolución (cierre de 2025 y abril de 2026).

Normativa:
- Real Decreto-ley 13/2022, de 26 de julio, por el que se establece un nuevo sistema de cotización para los trabajadores por cuenta propia o autónomos.
- Real Decreto-ley 1/2023, de 10 de enero, de medidas urgentes en materia de incentivos a la contratación laboral y mejora de la protección social de las personas artistas.
- Real Decreto 126/2026, de 18 de febrero, por el que se fija el SMI para 2026.
- Orden de cotización para 2026, que desarrolla las normas legales de cotización.
- Ley 28/2003, de 29 de septiembre, reguladora del Fondo de Reserva de la Seguridad Social, y Real Decreto 337/2004.
- Tarifa de primas por contingencias profesionales (disposición adicional cuarta de la Ley 42/2006, en su redacción vigente) y Clasificación Nacional de Actividades Económicas 2025.

Pendiente de verificar antes de fijar el texto:
- la base mínima general de 2026 en la orden de cotización;
- las cuantías de la tabla de autónomos de 2026;
- el importe de las transferencias por gastos impropios en 2024, 2025 y 2026;
- el porcentaje del PIB del máximo del Fondo de Reserva en 2011;
- las disposiciones del Fondo en 2016 y 2017;
- el límite anual de disposición del Fondo a partir de 2033;
- la estimación del déficit del componente contributivo en el trabajo de García Díaz.

¿Seguimos con 2.3 (presupuesto consolidado por subsectores y gasto comparado)?

Fuentes de actualidad usadas:
- [El Debate: el Tribunal de Cuentas critica que el Estado financie la Seguridad Social con préstamos (junio de 2026)](https://www.eldebate.com/economia/20260626/tribunal-cuentas-afea-estado-financie-seguridad-social-prestamos_433084.html)
- [Press Digital: el Tribunal de Cuentas pide sustituir los préstamos por transferencias](https://www.pressdigital.es/articulo/economia/2026-06-26/5933706-tribunal-cuentas-pide-sustituir-prestamos-seguridad-social-transferencias-estado)
- [El Independiente: el crecimiento de la deuda de la Seguridad Social duplica lo que entra en la hucha](https://www.elindependiente.com/economia/2026/04/14/crecimiento-deuda-seguridad-social-duplica-entra-hucha-pensiones/)
- [MISSM: el Fondo de Reserva cerrará 2025 por encima de los 14.000 millones](https://www.inclusion.gob.es/en/w/el-fondo-de-reserva-de-la-seguridad-social-cerrara-2025-por-encima-de-los-14.000-millones-de-euros)
- [MISSM: el Fondo de Reserva supera los 15.200 millones en abril de 2026](https://www.inclusion.gob.es/w/el-fondo-de-reserva-de-la-seguridad-social-supera-los-15.200-millones-de-euros-en-abril)
- [La Moncloa: ingresos por cotizaciones de enero a julio de 2026](https://www.lamoncloa.gob.es/serviciosdeprensa/notasprensa/inclusion/Paginas/2026/310726-ingresos-cotizaciones-ss.aspx?qfr=2)
- [El Español: previsión de ingresos por cotizaciones para 2026](https://www.elespanol.com/invertia/economia/empleo/20251208/gobierno-estima-seguridad-social-ingresara-cifra-record-subir-cotizaciones/1003744044881_0.html)
- [MISSM: afiliación a cierre de 2025](https://www.inclusion.gob.es/en/w/la-seguridad-social-suma-en-2025-medio-mill%C3%B3n-de-afiliados-por-cuarto-a%C3%B1o-consecutivo-y-se-acerca-a-los-21-9-millones-de-ocupados)
- [AFINOM: bases y tipos de cotización 2026](https://afinom.es/bases-y-tipos-de-cotizacion-2026/)
- [BOE: Real Decreto 126/2026 (SMI)](https://www.boe.es/buscar/doc.php?id=BOE-A-2026-3815)
- [65 y más: gastos impropios y déficit de la Seguridad Social](https://www.65ymas.com/futuro-pensiones/escriva-ahorra-38284-millones-gastos-impropios-deficit-no-baja-0-5-pib_46858_102.html)
- [Fedea (IDEAS/RePEc): García Díaz (2024)](https://ideas.repec.org/p/fda/fdafen/2024-14.html)


### Referencias nuevas

- Adema, W. y Ladaique, M. (2009), "How Expensive is the Welfare State? Gross and Net Indicators in the OECD Social Expenditure Database (SOCX)", OECD Social, Employment and Migration Working Papers, 92.
- Alesina, A. y Angeletos, G.-M. (2005), "Fairness and Redistribution", American Economic Review, 95(4).
- Alesina, A. y Glaeser, E. (2004), Fighting Poverty in the US and Europe: A World of Difference, Oxford University Press.
- Baumol, W. (1967), "Macroeconomics of Unbalanced Growth: The Anatomy of Urban Crisis", American Economic Review, 57(3).
- Howard, C. (1997), The Hidden Welfare State: Tax Expenditures and Social Policy in the United States, Princeton University Press.
- Iversen, T. y Soskice, D. (2006), "Electoral Institutions and the Politics of Coalitions: Why Some Democracies Redistribute More Than Others", American Political Science Review, 100(2).
- Korpi, W. (1983), The Democratic Class Struggle, Routledge & Kegan Paul.
- Lynch, J. (2006), Age in the Welfare State: The Origins of Social Spending on Pensioners, Workers, and Children, Cambridge University Press.
- Titmuss, R. (1958), "The Social Division of Welfare", en Essays on "the Welfare State", Allen & Unwin.
- Vanhuysse, P. (2013), Intergenerational Justice in Aging Societies: A Cross-National Comparison of 29 OECD Countries, Bertelsmann Stiftung.
- Eurostat, Government expenditure by function – COFOG y Government expenditure on social protection, Statistics Explained (datos de 2023 y 2024).
- Eurostat, ESSPROS Manual and User Guidelines.
- OCDE, Social Expenditure Database (SOCX) y Society at a Glance 2024.
- Fedea (2025), "La distribución del gasto público por políticas en el periodo 1995-2024", Fedea Policy Blog.
- Ministerio de Sanidad (2026), nota sobre el gasto sanitario público de 2024.
- Ministerio de Inclusión, Seguridad Social y Migraciones (2026), nota sobre el cierre presupuestario de la Seguridad Social en 2025.
- IGAE, Clasificación funcional del gasto de las Administraciones Públicas (COFOG), informe anual.

Pendiente de verificar antes de fijar el texto:
- las cifras COFOG de España para 2024 en protección social y sanidad, en porcentaje del PIB;
- el déficit de 2025 de las Administraciones de Seguridad Social en contabilidad nacional;
- el PIB nominal de 2025 utilizado para las ratios;
- la homogeneidad entre los datos del Ministerio de 2025 y los de liquidación de 2021 del documento;
- los datos de reducción de la pobreza por transferencias y las tasas de pobreza infantil y de los mayores de la última Encuesta de Condiciones de Vida;
- el gasto social público y neto de España en la última edición de SOCX.


Fuentes de actualidad usadas:
- [La Moncloa: el déficit de la Seguridad Social se reduce al 0,4% del PIB en 2025](https://www.lamoncloa.gob.es/serviciosdeprensa/notasprensa/inclusion/paginas/2026/310326-seguridad-social-deficit-pib.aspx)
- [Eurostat: gasto público por funciones (COFOG), 2024](https://ec.europa.eu/eurostat/statistics-explained/index.php?title=Government_expenditure_by_function_%E2%80%93_COFOG)
- [Eurostat: gasto público en protección social, 2023](https://ec.europa.eu/eurostat/statistics-explained/index.php?title=Government_expenditure_on_social_protection)
- [Fedea Policy Blog: distribución del gasto público por políticas, 1995-2024](https://policy.fedea.net/la-distribucion-del-gasto-publico-por-politicas-en-el-periodo-1995-2024-muestra-un-peso-cada-vez-mayor-de-las-prestaciones-sociales-en-especial-de-las-pensiones-y-la-sanidad-en-detrimento-del-rest/)
- [Ministerio de Sanidad: el gasto sanitario público de 2024](https://www.sanidad.gob.es/en/gabinete/notasPrensa.do?id=6908)
- [IGAE: clasificación funcional del gasto (COFOG)](https://www.igae.pap.hacienda.gob.es/sitios/igae/es-ES/Contabilidad/ContabilidadNacional/Publicaciones/paginas/iacogof.aspx)
- [OCDE: Social Expenditure Database (SOCX)](https://www.oecd.org/en/data/datasets/social-expenditure-database-socx.html)
### Referencias nuevas

- Boldrin, M., Jiménez-Martín, S. y Peracchi, F. (1999), "Social Security and Retirement in Spain", en Gruber, J. y Wise, D. (eds.), Social Security and Retirement around the World, University of Chicago Press/NBER.
- Currie, J. (2006), "The Take-Up of Social Benefits", en Auerbach, A., Card, D. y Quigley, J. (eds.), Public Policy and the Income Distribution, Russell Sage Foundation.
- Gruber, J. y Wise, D. (eds.) (1999), Social Security and Retirement around the World, University of Chicago Press/NBER.
- Gruber, J. y Wise, D. (eds.) (2004), Social Security Programs and Retirement around the World: Micro-Estimation, University of Chicago Press/NBER.
- Herd, P. y Moynihan, D. (2018), Administrative Burden: Policymaking by Other Means, Russell Sage Foundation.
- AIReF (2025), Cuarta Opinión sobre el IMV, julio de 2025.
- AIReF (2019), Los programas de rentas mínimas en España, Estudio.

Normativa y jurisprudencia:
- Real Decreto 39/2026, de 21 de enero, sobre limitación de la cuantía inicial de las pensiones públicas y revalorización de las pensiones para 2026.
- Real Decreto-ley 11/2024, de 23 de diciembre, para la mejora de la compatibilidad de la pensión de jubilación con el trabajo.
- Ley 2/2025, de 29 de abril, por la que se modifica el Estatuto de los Trabajadores en materia de extinción del contrato por incapacidad permanente.
- Texto refundido de la Ley General de la Seguridad Social, disposición transitoria trigésima cuarta (aplicación gradual de coeficientes reductores sobre pensiones que superen el límite).
- Sentencia del TJUE de 18 de enero de 2024, asunto C-631/22.
- Ley 19/2021, de 20 de diciembre, por la que se establece el IMV.

Pendiente de verificar antes de fijar el texto:
- la cifra actual de pensiones y gasto de Clases Pasivas;
- la esperanza de vida a los 65 años en la última tabla del INE;
- los parámetros exactos de la integración de lagunas para hombres y mujeres;
- la senda exacta de convergencia de las pensiones mínimas con el umbral de la pobreza;
- los porcentajes de la jubilación activa según los años de demora;
- la cuantía de la cantidad a tanto alzado de la jubilación demorada;
- la brecha de género en las pensiones de jubilación;
- la sentencia del TS de 2024 sobre la compatibilidad de la incapacidad permanente absoluta;
- las cuantías de 2026 del IMV y del complemento de ayuda para la infancia;
- el número de partícipes y el patrimonio de los fondos de pensiones de empleo de promoción pública.



Fuentes de actualidad usadas:
- [MISSM: la pensión media del sistema alcanza los 1.374,6 euros en septiembre de 2026](https://www.inclusion.gob.es/w/la-pension-media-del-sistema-de-la-seguridad-social-alcanza-los-1.374-6-euros-en-septiembre-un-4-6-mas-alta-que-hace-un-ano)
- [El Independiente: pensiones mínimas y máximas en 2026](https://www.elindependiente.com/economia/2026/09/23/pensiones-minimas-maximas-espana-2026/)
- [BOE: Real Decreto 39/2026 de revalorización](https://www.boe.es/diario_boe/txt.php?id=BOE-A-2026-1484)
- [WTW: cálculo de la base reguladora en 2026](https://www.wtwco.com/es-es/insights/2026/03/como-se-calcula-mi-base-reguladora-en-la-pension-de-jubilacion)
- [BBVA Mi Jubilación: reforma de la jubilación parcial, activa y demorada](https://www.bbvamijubilacion.es/blog/reforma-de-la-jubilacion-parcial-la-jubilacion-activa-y-la-demorada/)
- [BBVA Mi Jubilación: jubilación anticipada en 2026](https://www.bbvamijubilacion.es/blog/jubilacion-anticipada-en-2026-estos-son-los-requisitos-y-los-coeficientes-reductores/)
- [BOE: Real Decreto-ley 11/2024](https://www.boe.es/buscar/doc.php?id=BOE-A-2024-26917)
- [El Español: cuarta opinión de la AIReF sobre el IMV](https://www.elespanol.com/invertia/economia/20250709/airef-detecta-cobrar-ingreso-minimo-reduce-probabilidad-trabajar/1003743840443_0.html)
- [Laboral Pensiones: cuarta opinión de la AIReF sobre el IMV](https://laboralpensiones.com/64303-2/)
- [Teleprensa: respuesta del Ministerio a la AIReF](https://www.teleprensa.com/articulo/nacional-3/ministerio-inclusion-responde-airef-que-evaluacion-incentivo-empleo-imv-es-temprana/202507091522282166164.html)

- Askildsen, J. E., Bratberg, E. y Nilsen, Ø. A. (2005), "Unemployment, Labor Force Composition and Sickness Absence: A Panel Data Study", Health Economics, 14(11).
- Cools, S., Fiva, J. H. y Kirkebøen, L. J. (2015), "Causal Effects of Paternity Leave on Children and Parents", Scandinavian Journal of Economics, 117(3).
- Dahl, G., Løken, K. y Mogstad, M. (2014), "Peer Effects in Program Participation", American Economic Review, 104(7).
- De Quinto, A., Hospido, L. y Sanz, C. (2021), "The Child Penalty: Evidence from Spain", SERIEs, 12(4).
- Farré, L. y González, L. (2019), "Does Paternity Leave Reduce Fertility?", Journal of Public Economics, 172.
- Johansson, P. y Palme, M. (1996), "Do Economic Incentives Affect Work Absence? Empirical Evidence Using Swedish Micro Data", Journal of Public Economics, 59(2).
- Johansson, P. y Palme, M. (2005), "Moral Hazard and Sickness Insurance", Journal of Public Economics, 89(9-10).
- Kleven, H., Landais, C. y Søgaard, J. E. (2019), "Children and Gender Inequality: Evidence from Denmark", American Economic Journal: Applied Economics, 11(4).
- Kleven, H., Landais, C., Posch, J., Steinhauer, A. y Zweimüller, J. (2019), "Child Penalties across Countries: Evidence and Explanations", AEA Papers and Proceedings, 109.
- Lalive, R. y Zweimüller, J. (2009), "How Does Parental Leave Affect Fertility and Return to Work? Evidence from Two Natural Experiments", Quarterly Journal of Economics, 124(3).
- Markussen, S., Røed, K., Røgeberg, O. y Gaure, S. (2011), "The Anatomy of Absenteeism", Journal of Health Economics, 30(2).
- Markussen, S., Mykletun, A. y Røed, K. (2012), "The Case for Presenteeism: Evidence from Norway's Sickness Insurance Program", Journal of Public Economics, 96(11-12).
- Markussen, S., Røed, K. y Røgeberg, O. (2013), "The Changing of the Guards: Can Physicians Contain Social Insurance Costs?", Journal of Health Economics, 32(5).
- Patnaik, A. (2019), "Reserving Time for Daddy: The Consequences of Fathers' Quotas", Journal of Labor Economics, 37(4).
- Ruhm, C. (1998), "The Economic Consequences of Parental Leave Mandates: Lessons from Europe", Quarterly Journal of Economics, 113(1).
- Van Groezen, B., Leers, T. y Meijdam, L. (2003), "Social Security and Endogenous Fertility: Pensions and Child Allowances as Siamese Twins", Journal of Public Economics, 87(2).
- Ziebarth, N. y Karlsson, M. (2010), "A Natural Experiment on Sick Pay Cuts, Sickness Absence, and Labor Costs", Journal of Public Economics, 94(11-12).
- AIReF (2026), Evaluación del gasto público. Estudio: Incapacidad temporal, Spending Review 2022-2026, fase II.

Normativa:
- Ley Orgánica 1/2023, de 28 de febrero, por la que se modifica la Ley Orgánica 2/2010, de salud sexual y reproductiva y de la interrupción voluntaria del embarazo.
- Real Decreto 625/2014, de 18 de julio, sobre la gestión y control de la incapacidad temporal, y Real Decreto 1060/2022, de 27 de diciembre, que lo modifica.
- Real Decreto-ley 20/2012, de 13 de julio, de medidas para garantizar la estabilidad presupuestaria y de fomento de la competitividad.
- Real Decreto-ley 6/2019, de 1 de marzo, de medidas urgentes para garantía de la igualdad de trato y de oportunidades entre mujeres y hombres en el empleo y la ocupación.
- Real Decreto-ley 9/2025, de 29 de julio, por el que se amplía el permiso de nacimiento y cuidado.
- Directiva (UE) 2019/1158 del Parlamento Europeo y del Consejo, de 20 de junio de 2019, relativa a la conciliación de la vida familiar y la vida profesional de los progenitores y los cuidadores.

Pendiente de verificar antes de fijar el texto:
- la penalización por maternidad estimada para España por De Quinto, Hospido y Sanz;
- los requisitos de cotización de la prestación por nacimiento según la edad;
- el ámbito de edad vigente de la prestación por cuidado de menores con enfermedad grave;
- la tasa de fecundidad de 2025;
- el gasto español en la función de familia e infancia en porcentaje del PIB;
- la fecha exacta de la reforma neerlandesa que amplió a dos años el pago empresarial de las bajas.

¿Seguimos con 3.3 (el desempleo y el Sistema Nacional de Empleo)?

Fuentes de actualidad usadas:
- [AIReF: estudio de revisión del gasto sobre la incapacidad temporal (febrero de 2026)](https://www.airef.es/wp-content/uploads/2026/02/SpendingReview-2022-2026-fase-2/SR2226F2_IT.pdf)
- [ESdiario: la AIReF denuncia que el absentismo se ha disparado un 60\%](https://www.esdiario.com/economia/260204/178871/airef-denuncia-absentismo-laboral-disparado-60-falta-control.html)
- [La Moncloa: el Gobierno pone en marcha el Observatorio de la Incapacidad Temporal](https://www.lamoncloa.gob.es/serviciosdeprensa/notasprensa/inclusion/paginas/2026/090226-observatorio-incapacidad-temporal.aspx)
- [El Independiente: la propuesta de altas progresivas deja fríos a los agentes sociales](https://www.elindependiente.com/economia/2025/06/04/agentes-sociales-oponen-propuesta-altas-progresivas-seguridad-social/)
- [Infobae: el Gobierno propone una incorporación gradual de 30 días a media jornada](https://www.infobae.com/espana/agencias/2025/06/02/el-gobierno-propone-una-incorporacion-gradual-de-30-dias-a-media-jornada-tras-bajas-largas/)
- [Mundiario: la patronal presiona al Gobierno contra el incremento de las bajas](https://www.mundiario.com/articulo/economia/patronal-presiona-gobierno-que-actue-incremento-bajas-laborales/20260709123420393518.html)
- [BOE: Real Decreto-ley 9/2025](https://www.boe.es/buscar/doc.php?id=BOE-A-2025-15741)
- [Laboral Social: contenido del Real Decreto-ley 9/2025](https://www.laboral-social.com/real-decreto-ley-9-2025-permiso-nacimiento-cuidado)
- [La Moncloa: el Congreso aprueba la ampliación de los permisos por nacimiento](https://www.lamoncloa.gob.es/serviciosdeprensa/notasprensa/trabajo14/paginas/2025/090925-real-decreto-ley-permisos-nacimientos.aspx)
- Baily, M. N. (1978), "Some Aspects of Optimal Unemployment Insurance", Journal of Public Economics, 10(3).
- Bentolila, S., Cahuc, P., Dolado, J. J. y Le Barbanchon, T. (2012), "Two-Tier Labour Markets in the Great Recession: France versus Spain", Economic Journal, 122(562).
- Blanchard, O. y Tirole, J. (2008), "The Joint Design of Unemployment Insurance and Employment Protection: A First Pass", Journal of the European Economic Association, 6(1).
- Card, D., Chetty, R. y Weber, A. (2007), "Cash-on-Hand and Competing Models of Intertemporal Behavior: New Evidence from the Labor Market", Quarterly Journal of Economics, 122(4).
- Card, D., Kluve, J. y Weber, A. (2010), "Active Labour Market Policy Evaluations: A Meta-Analysis", Economic Journal, 120(548).
- Card, D., Kluve, J. y Weber, A. (2018), "What Works? A Meta Analysis of Recent Active Labor Market Program Evaluations", Journal of the European Economic Association, 16(3).
- Chetty, R. (2006), "A General Formula for the Optimal Level of Social Insurance", Journal of Public Economics, 90(10-11).
- Chetty, R. (2008), "Moral Hazard versus Liquidity and Optimal Unemployment Insurance", Journal of Political Economy, 116(2).
- Hopenhayn, H. y Nicolini, J. P. (1997), "Optimal Unemployment Insurance", Journal of Political Economy, 105(2).
- Katz, L. y Meyer, B. (1990), "The Impact of the Potential Duration of Unemployment Benefits on the Duration of Unemployment", Journal of Public Economics, 41(1).
- Landais, C., Michaillat, P. y Saez, E. (2018), "A Macroeconomic Approach to Optimal Unemployment Insurance: Theory", American Economic Journal: Economic Policy, 10(2).
- Madsen, P. K. (2004), "The Danish Model of 'Flexicurity': Experiences and Lessons", Transfer: European Review of Labour and Research, 10(2).
- Nekoei, A. y Weber, A. (2017), "Does Extending Unemployment Benefits Improve Job Quality?", American Economic Review, 107(2).
- Persson, T. y Tabellini, G. (1996), "Federal Fiscal Constitutions: Risk Sharing and Moral Hazard", Econometrica, 64(3).
- Rebollo-Sanz, Y. y Rodríguez-Planas, N. (2020), "When the Going Gets Tough... Financial Incentives, Duration of Unemployment, and Job-Match Quality", Journal of Human Resources, 55(1).
- Shavell, S. y Weiss, L. (1979), "The Optimal Payment of Unemployment Insurance Benefits over Time", Journal of Political Economy, 87(6).
- INE (2026), Encuesta de Población Activa. Segundo trimestre de 2026.

Normativa y jurisprudencia:
- Ley 3/2023, de 28 de febrero, de Empleo.
- Texto refundido de la Ley General de la Seguridad Social, título III (protección por desempleo).
- Real Decreto-ley 20/2012, de 13 de julio.
- Real Decreto-ley 32/2021, de 28 de diciembre, de medidas urgentes para la reforma laboral.
- Real Decreto-ley 16/2022, de 6 de septiembre, para la mejora de las condiciones de trabajo y de Seguridad Social de las personas trabajadoras al servicio del hogar.
- Real Decreto-ley 7/2023, de 19 de diciembre (no convalidado).
- Real Decreto-ley 2/2024, de 21 de mayo, para la simplificación y mejora del nivel asistencial de la protección por desempleo.
- Sentencia del Tribunal de Justicia de la Unión Europea de 24 de febrero de 2022, asunto C-389/20.

Pendiente de verificar antes de fijar el texto:
- la referencia y los resultados exactos de Rebollo-Sanz y Rodríguez-Planas;
- la duración del subsidio según la edad y las responsabilidades familiares;
- la base de cotización para la jubilación en el subsidio de mayores de 52 años;
- el gasto, el número de beneficiarios, la cuantía media y la tasa de cobertura de 2026;
- las conclusiones de la revisión del gasto de la AIReF sobre las políticas activas;
- la proporción de contrataciones intermediadas por los servicios públicos de empleo.

¿Seguimos con 3.4 (el Sistema Nacional de Salud)?

Fuentes de actualidad usadas:
- [INE: Encuesta de Población Activa, segundo trimestre de 2026](https://www.ine.es/dyngs/Prensa/EPA2T26.htm)
- [Iberley: las 15 claves del Real Decreto-ley 2/2024](https://www.iberley.es/revista/rdl-2-2024-las-15-claves-reforma-subsidio-desempleo-875)
- [La Moncloa: claves de la reforma del subsidio por desempleo](https://www.lamoncloa.gob.es/serviciosdeprensa/notasprensa/trabajo14/paginas/2024/210624-reforma-subsidio-desempleo.aspx)
- [Civio: segundo intento de reforma del subsidio](https://civio.es/el-boe-nuestro-de-cada-dia/2024/05/22/reforma-del-subsidio-por-desempleo-segundo-intento-con-algunos-cambios/)
- [Ayuda Trámites: cuantías del paro y de los subsidios en 2026](https://ayudatramites.es/noticias/cuantias-paro-2026/)
- [SEPE: cuantías anuales](https://www.sepe.es/HomeSepe/en/prestaciones-desempleo/Cuantias-anuales.html)
- Arrow, K. (1963), "Uncertainty and the Welfare Economics of Medical Care", American Economic Review, 53(5).
- Besley, T. y Case, A. (1995), "Incumbent Behavior: Vote-Seeking, Tax-Setting, and Yardstick Competition", American Economic Review, 85(1).
- Cutler, D. y McClellan, M. (2001), "Is Technological Change in Medicine Worth It?", Health Affairs, 20(5).
- Enthoven, A. (1978), "Consumer-Choice Health Plan", New England Journal of Medicine, 298(12-13); (1993), "The History and Principles of Managed Competition", Health Affairs, 12(suppl.).
- Fries, J. (1980), "Aging, Natural Death, and the Compression of Morbidity", New England Journal of Medicine, 303(3).
- Getzen, T. (2000), "Health Care Is an Individual Necessity and a National Luxury: Applying Multilevel Decision Models to the Analysis of Health Care Expenditures", Journal of Health Economics, 19(2).
- Kornai, J. (1986), "The Soft Budget Constraint", Kyklos, 39(1).
- Manning, W. et al. (1987), "Health Insurance and the Demand for Medical Care: Evidence from a Randomized Experiment", American Economic Review, 77(3).
- Newhouse, J. (1977), "Medical-Care Expenditure: A Cross-National Survey", Journal of Human Resources, 12(1).
- Newhouse, J. (1992), "Medical Care Costs: How Much Welfare Loss?", Journal of Economic Perspectives, 6(3).
- Pauly, M. (1968), "The Economics of Moral Hazard: Comment", American Economic Review, 58(3).
- Rodden, J. (2006), Hamilton's Paradox: The Promise and Peril of Fiscal Federalism, Cambridge University Press.
- Tobin, J. (1970), "On Limiting the Domain of Inequality", Journal of Law and Economics, 13(2).
- Tudor Hart, J. (1971), "The Inverse Care Law", The Lancet, 297(7696).
- Zweifel, P., Felder, S. y Meiers, M. (1999), "Ageing of Population and Health Care Expenditure: A Red Herring?", Health Economics, 8(6).
- Ministerio de Sanidad, Sistema de Información sobre Listas de Espera en el Sistema Nacional de Salud, datos a 31 de diciembre de 2025.
- Ministerio de Sanidad (2025), nota de prensa sobre el gasto en receta médica del SNS en 2024.

Normativa:
- Ley 14/1986, de 25 de abril, General de Sanidad.
- Ley 15/1997, de 25 de abril, sobre habilitación de nuevas formas de gestión del Sistema Nacional de Salud.
- Ley 16/2003, de 28 de mayo, de cohesión y calidad del Sistema Nacional de Salud.
- Real Decreto 1030/2006, de 15 de septiembre, por el que se establece la cartera de servicios comunes del Sistema Nacional de Salud.
- Ley 33/2011, de 4 de octubre, General de Salud Pública.
- Real Decreto-ley 16/2012, de 20 de abril.
- Real Decreto Legislativo 1/2015, de 24 de julio, texto refundido de la Ley de garantías y uso racional de los medicamentos y productos sanitarios.
- Real Decreto-ley 7/2018, de 27 de julio, sobre el acceso universal al Sistema Nacional de Salud.
- Real Decreto-ley de reforma del copago farmacéutico, convalidado el 28 de mayo de 2026.
- Proyecto de Ley de Gestión Pública e Integridad del Sistema Nacional de Salud (Boletín Oficial de las Cortes Generales, serie A, núm. 95-1, 2026).

Pendiente de verificar antes de fijar el texto:
- el número y la fecha del real decreto-ley de reforma del copago;
- la referencia exacta de los estudios de Puig-Junoy y coautores sobre la reforma de 2012;
- las proyecciones de gasto sanitario de España en el Ageing Report de 2024;
- el número de asegurados del seguro privado de salud;
- las cifras de gasto farmacéutico de 2025 con una metodología homogénea con la de 2024;
- la fecha exacta de la reversión del hospital de Alzira.

¿Seguimos con 3.5 (el Sistema para la Autonomía y Atención a la Dependencia)?

Fuentes de actualidad usadas:
- [Infobae: las listas de espera quirúrgica superan las 850.000 personas](https://www.infobae.com/espana/2026/04/16/las-listas-de-espera-quirurgica-superan-las-850000-personas-en-2025/)
- [Ministerio de Sanidad: listas de espera a 31 de diciembre de 2025](https://www.sanidad.gob.es/estadEstudios/estadisticas/inforRecopilaciones/docs/LISTAS_PUBLICACION_Dic_2025.pdf)
- [El Debate: el Congreso convalida la reforma del copago farmacéutico](https://www.eldebate.com/sociedad/20260528/congreso-da-luz-verde-reforma-copago-farmaceutico-nuevos-tramos-topes-mensuales_422629.html)
- [Ministerio de Sanidad: reforma del sistema de copago farmacéutico](https://www.sanidad.gob.es/gabinete/notasPrensa.do?id=6905)
- [Ministerio de Sanidad: aprobado el proyecto de Ley de Gestión Pública e Integridad del SNS](https://www.sanidad.gob.es/gabinete/notasPrensa.do?id=6906)
- [Congreso: Proyecto de Ley 121/000094](https://www.congreso.es/public_oficiales/L15/CONG/BOCG/A/BOCG-15-A-95-1.PDF)
- [Ministerio de Sanidad: el gasto sanitario público alcanzó los 101.739 millones en 2024](https://www.sanidad.gob.es/en/gabinete/notasPrensa.do?id=6908)
- [Ministerio de Sanidad: gasto en receta médica del SNS en 2024](https://www.sanidad.gob.es/en/gabinete/notasPrensa.do?id=6821)
- [Redacción Médica: el gasto farmacéutico del SNS en 2025](https://www.redaccionmedica.com/secciones/industria/el-sns-cerro-el-2025-con-una-subida-del-7-en-su-gasto-farmaceutico-7522)
- [Opinando sin anestesia: el nuevo modelo de financiación autonómica y el reparto sanitario](https://opinandosinanestesia.es/2026/09/luz-verde-al-nuevo-modelo-de-financiacion-autonomica-con-2-instrumentos-para-garantizar-el-reparto-sanitario/)
- [La Moncloa: propuesta del nuevo modelo de financiación autonómica](https://www.lamoncloa.gob.es/serviciosdeprensa/notasprensa/hacienda/paginas/2026/090126-modelo-financiacion-comunidades.aspx)
- Barr, N. (2010), "Long-Term Care: A Suitable Case for Social Insurance", Social Policy & Administration, 44(4).
- Bernheim, B. D., Shleifer, A. y Summers, L. (1985), "The Strategic Bequest Motive", Journal of Political Economy, 93(6).
- Bettio, F., Simonazzi, A. y Villa, P. (2006), "Change in Care Regimes and Female Migration: The 'Care Drain' in the Mediterranean", Journal of European Social Policy, 16(3).
- Brown, J. y Finkelstein, A. (2008), "The Interaction of Public and Private Insurance: Medicaid and the Long-Term Care Insurance Market", American Economic Review, 98(3).
- Cabrero, G. R. et al. (2022), estudio sobre el desarrollo del Sistema para la Autonomía y Atención a la Dependencia, citado en el documento de referencia (referencia completa pendiente de verificar).
- Da Roit, B. y Le Bihan, B. (2010), "Similar and Yet So Different: Cash-for-Care in Six European Countries' Long-Term Care Policies", Milbank Quarterly, 88(3).
- González Ortega, S. (2013), estudio sobre el impacto de la crisis en el sistema de dependencia, citado en el documento de referencia (referencia completa pendiente de verificar).
- Pauly, M. (1990), "The Rational Nonpurchase of Long-Term-Care Insurance", Journal of Political Economy, 98(1).
- Asociación Estatal de Directoras y Gerentes en Servicios Sociales (2026), XXVI Dictamen del Observatorio Estatal de la Dependencia, marzo de 2026.
- Ministerio de Derechos Sociales, Consumo y Agenda 2030 (2026), Panel del SAAD, segundo trimestre de 2026.
- Naciones Unidas (2006), Convención sobre los Derechos de las Personas con Discapacidad.

Normativa:
- Ley 39/2006, de 14 de diciembre, de Promoción de la Autonomía Personal y Atención a las personas en situación de dependencia.
- Real Decreto-ley 20/2012, de 13 de julio (medidas en materia de dependencia).
- Real Decreto-ley 6/2019, de 1 de marzo (recuperación de la cotización de los cuidadores no profesionales).
- Acuerdo del Consejo Territorial de Servicios Sociales y del SAAD sobre criterios comunes de acreditación y calidad de los centros y servicios (2022).
- Reforma de la Ley 39/2006 y del texto refundido de la Ley General de derechos de las personas con discapacidad, aprobada definitivamente en septiembre de 2026.

Pendiente de verificar antes de fijar el texto:
- el requisito exacto de residencia y el régimen de los extranjeros;
- el número de ley y la fecha de la reforma de 2026;
- el instrumento normativo del aumento de financiación de 6.200 millones;
- el gasto en cuidados de larga duración de la Unión Europea en porcentaje del PIB;
- la cifra de fallecidos en residencias durante la pandemia;
- las referencias completas de González Ortega (2013) y Cabrero y otros (2022).

Queda cerrado el bloque III. ¿Pasamos al bloque IV (la sostenibilidad del sistema público de pensiones), empezando por 4.1 (el concepto de sostenibilidad)?

Fuentes de actualidad usadas:
- [Ministerio de Derechos Sociales: la lista de espera de la dependencia baja un 21% en el último año](https://www.dsca.gob.es/es/comunicacion/notas-prensa/lista-espera-dependencia-acelera-su-descenso-baja-21-ultimo-ano)
- [Observatorio Estatal de la Dependencia: nota de prensa del XXVI Dictamen (marzo de 2026)](https://directoressociales.com/wp-content/uploads/2026/03/NdP-XXVI-Dictamen-13-03-26.pdf)
- [Público: personas fallecidas en las listas de espera de la dependencia en 2026](https://www.publico.es/sociedad/sanidad/13-503-personas-han-muerto-listas-espera-dependencia-primeros-cinco-meses-2026-insoportable.html)
- [Moncloa.com: el Congreso aprueba la reforma de dependencia y discapacidad](https://www.moncloa.com/2026/07/15/reforma-dependencia-discapacidad-congreso-aprobada-3399837/)
- [Consejo General del Trabajo Social: aprobación definitiva de la reforma](https://www.cgtrabajosocial.es/noticias/el-congreso-aprueba-definitivamente-la-reforma-de-las-leyes-de-dependencia-y-discapacidad/15086/view)
- [IMSERSO: estadísticas mensuales del SAAD](https://imserso.es/en/el-imserso/documentacion/estadisticas/sistema-autonomia-atencion-dependencia-saad/estadisticas-mensual)
### Referencias nuevas

- Blanchard, O. (1990), "Suggestions for a New Set of Fiscal Indicators", OECD Economics Department Working Papers, 79.
- Boado-Penas, M. C., Valdés-Prieto, S. y Vidal-Meliá, C. (2008), "The Actuarial Balance Sheet for Pay-As-You-Go Finance: Solvency Indicators for Spain and Sweden", Fiscal Studies, 29(1).
- Buiter, W. (1997), "Generational Accounts, Aggregate Saving and Intergenerational Distribution", Economica, 64(256).
- Devesa, J. E., Devesa, M., Meneu, R. y otros, trabajos sobre la tasa interna de rendimiento y el equilibrio actuarial del sistema español de jubilación (referencias exactas pendientes de verificar).
- Haveman, R. (1994), "Should Generational Accounts Replace Public Budgets and Deficits?", Journal of Economic Perspectives, 8(1).
- Kotlikoff, L. (1992), Generational Accounting: Knowing Who Pays, and When, for What We Spend, Free Press.
- Lee, R. y Mason, A. (eds.) (2011), Population Aging and the Generational Economy: A Global Perspective, Edward Elgar.
- Settergren, O. (2001), "The Automatic Balance Mechanism of the Swedish Pension System", Wirtschaftspolitische Blätter, 48(4).
- Comisión Europea (2024), 2024 Ageing Report: Economic and Budgetary Projections for the EU Member States (2022-2070), Institutional Paper 279.
- Comisión Europea, Debt Sustainability Monitor y Fiscal Sustainability Report (últimas ediciones), indicadores S1 y S2.
- Comisión Europea y Comité de Protección Social (2024), 2024 Pension Adequacy Report.
- AIReF (2026), Informe sobre la regla de gasto de pensiones.

Pendiente de verificar antes de fijar el texto:
- la clasificación de riesgo de sostenibilidad fiscal de España en el último Debt Sustainability Monitor;
- los resultados del balance actuarial de Boado-Penas, Valdés-Prieto y Vidal-Meliá y su actualización;
- las referencias y estimaciones exactas de las tasas internas de rendimiento de Devesa y coautores;
- los resultados de las cuentas nacionales de transferencias para España;
- la fecha de publicación del próximo Ageing Report.
- Aaron, H. (1966). "The Social Insurance Paradox". Canadian Journal of Economics and Political Science, 32(3).
- Abbott, R. y Bogenschneider, B. (2018). "Should Robots Pay Taxes? Tax Policy in the Age of Automation". Harvard Law & Policy Review, 12.
- Acemoglu, D. y Restrepo, P. (2019). "Automation and New Tasks: How Technology Displaces and Reinstates Labor". Journal of Economic Perspectives, 33(2).
- Autor, D., Dorn, D., Katz, L., Patterson, C. y Van Reenen, J. (2020). "The Fall of the Labor Share and the Rise of Superstar Firms". Quarterly Journal of Economics, 135(2).
- Ayuso, M., Bravo, J. M. y Holzmann, R. (2017). "Addressing Longevity Heterogeneity in Pension Scheme Design". Journal of Finance and Economics, 6(1).
- Banco de España (2019). Informe Anual 2018, cap. 4 (envejecimiento y migración).
- Bongaarts, J. y Feeney, G. (1998). "On the Quantum and Tempo of Fertility". Population and Development Review, 24(2).
- Comisión Europea (2024). The 2024 Ageing Report: Economic and Budgetary Projections for the EU Member States (2022-2070), Institutional Paper 279; y ficha de España.
- Diamond, P. (2004). "Social Security". American Economic Review, 94(1).
- Esping-Andersen, G. (2009). The Incomplete Revolution: Adapting to Women's New Roles. Polity.
- Gopinath, G., Kalemli-Özcan, Ş., Karabarbounis, L. y Villegas-Sánchez, C. (2017). "Capital Allocation and Productivity in South Europe". Quarterly Journal of Economics, 132(4).
- Guerreiro, J., Rebelo, S. y Teles, P. (2022). "Should Robots Be Taxed?". Review of Economic Studies, 89(1).
- INE (2026). Proyecciones de Población 2026-2076, nota de prensa del 17 de junio de 2026.
- Karabarbounis, L. y Neiman, B. (2014). "The Global Decline of the Labor Share". Quarterly Journal of Economics, 129(1).
- Kohler, H.-P., Billari, F. y Ortega, J. A. (2002). "The Emergence of Lowest-Low Fertility in Europe during the 1990s". Population and Development Review, 28(4).
- National Academies of Sciences (2015). The Growing Gap in Life Expectancy by Income. National Academies Press.
- Pérez Díaz, J. y MacInnes, J. (2009). "The Reproductive Revolution". The Sociological Review, 57(2).
- Piketty, T. (2014). Capital in the Twenty-First Century. Harvard University Press.
- Sanderson, W. y Scherbov, S. (2005). "Average Remaining Lifetimes Can Increase as Human Populations Age". Nature, 435; y (2010). "Remeasuring Aging". Science, 329.
- Spijker, J. y MacInnes, J. (2013). "Population Ageing: The Timebomb That Isn't?". BMJ, 347.
- Storesletten, K. (2000). "Sustaining Fiscal Policy through Immigration". Journal of Political Economy, 108(2).
- United Nations Population Division (2000). Replacement Migration: Is It a Solution to Declining and Ageing Populations?

Pendiente de verificar:
- el incremento adicional exacto de la pensión máxima en el RDL 2/2023;
- la definición de ingresos del AWG que subyace al déficit del 2,7% en 2070;
- el estado de tramitación de la regularización extraordinaria de inmigrantes;
- las cifras de la ficha de España del Ageing Report 2024, que he obtenido de un resumen y conviene contrastar en el PDF original antes de memorizarlas.

Fuentes de actualidad consultadas:
- [Ficha de España del Ageing Report 2024 (Comisión Europea)](https://economy-finance.ec.europa.eu/document/download/392994f2-c025-4066-bea8-f65cf4bcd411_en?filename=2024-ageing-report-country-fiche-Spain.pdf)
- [CaixaBank Research: el Ageing Report 2024 y la reforma de pensiones](https://www.caixabankresearch.com/sites/default/files/content/file/2024/05/16/34454/im05_24-07-ee-focus-8-en.pdf)
- [INE: Proyecciones de Población 2026-2076](https://www.ine.es/dyngs/Prensa/PROP20262076.htm)
- [INE: EPA del segundo trimestre de 2026](https://www.ine.es/dyngs/Prensa/EPA2T26.htm)
- [Ministerio de Inclusión: el Fondo de Reserva supera los 15.200 millones en abril](https://www.inclusion.gob.es/w/el-fondo-de-reserva-de-la-seguridad-social-supera-los-15.200-millones-de-euros-en-abril)
- [elEconomista: la ratio de trabajadores por pensionista llega a 2,35](https://www.eleconomista.es/economia/noticias/14012383/09/26/la-ratio-de-trabajadores-por-pensionista-marca-su-nivel-mas-alto-en-casi-20-anos-y-llega-a-235.html)
### Referencias nuevas (introducidas en 4.3)

- AIReF (2023). Opinión sobre la sostenibilidad del sistema de pensiones y evaluación de las reformas de 2021 y 2023.
- AIReF (2025). Informe sobre la regla de gasto de pensiones, 31 de marzo de 2025.
- AIReF (2026). Actualización del informe sobre la regla de gasto de pensiones, mayo de 2026.
- Beetsma, R. y Debrun, X. (eds.) (2018). Independent Fiscal Councils: Watchdogs or Lapdogs? CEPR Press.
- Comité de Expertos sobre el Factor de Sostenibilidad (2013). Informe del Comité de Expertos sobre el factor de sostenibilidad del sistema público de pensiones, 7 de junio de 2013.
- Conde-Ruiz, J. I. y González, C. I. (2013). "Reforma de pensiones 2011 en España". Hacienda Pública Española / Review of Public Economics, 204(1).
- De la Fuente, Á. (2025). "Sobre la aplicación de la cláusula de salvaguarda del MEI". Fedea, Apuntes 2025/18.
- De la Fuente, Á. y Doménech, R. (2013). "The financial impact of Spanish pension reform: a quick estimate". Journal of Pension Economics and Finance, 12(1).
- Díaz-Giménez, J. y Díaz-Saavedra, J. (2017). "The future of Spanish pensions". Journal of Pension Economics and Finance, 16(2).
- Lynch, J. (2006). Age in the Welfare State. Cambridge University Press.
- OCDE (2021). Pensions at a Glance 2021, cap. 2, "Automatic adjustment mechanisms in pension systems".
- Pierson, P. (1996). "The New Politics of the Welfare State". World Politics, 48(2).
- Sánchez, A. R. (2014). "The automatic adjustment of pension expenditures in Spain: an evaluation of the 2013 pension reform". Banco de España, Documento de Trabajo 1420.
- Sinn, H.-W. y Uebelmesser, S. (2002). "Pensions and the path to gerontocracy". European Journal of Political Economy, 19(1).
- Real Decreto 100/2025, de metodología de la cláusula de salvaguarda (pendiente de verificar el título exacto).

Pendiente de verificar:
- la naturaleza del informe de la AIReF de mayo de 2026;
- las cifras exactas de la revalorización de 2018-2019 fijadas en el acuerdo PP-PNV;
- la ratio de prestación proyectada en los Ageing Reports de 2015 y 2018 (tabla 5 del documento, ilegible en la versión de trabajo);
- el rendimiento anual del MEI en proporción del PIB.

Fuentes de actualidad consultadas:
- [AIReF: la regla de gasto de pensiones se cumple formalmente, pero no garantiza la sostenibilidad](https://www.airef.es/es/noticias/la-airef-ratifica-el-resultado-de-2025-la-regla-de-gasto-de-pensiones-se-cumple-formalmente-pero-no-garantiza-la-sostenibilidad/)
- [AIReF: nota de prensa del 31 de marzo de 2025](https://www.airef.es/wp-content/uploads/2025/03/Informe_regla_de_gasto_de_pensiones/Nota-de-prensa.pdf)
- [El Confidencial Digital: la AIReF constata el cumplimiento de la regla, pero avisa de que no mejora la sostenibilidad](https://www.elconfidencialdigital.com/articulo/actualidad/airef-constata-cumplimiento-regla-gasto-pensiones-avisa-mejora-sostenibilidad/20250331092215952260.html)
- [Canal Sur: la AIReF ratifica la regla de gasto (mayo de 2026)](https://www.canalsur.es/noticias/espana/airef-ratifica-reforma-pensiones-dice_1_1405297.html)
- [Fedea: De la Fuente (2025), sobre la aplicación de la cláusula de salvaguarda del MEI](https://documentos.fedea.net/pubs/ap/2025/ap2025-18.pdf)
- [Ficha de España del Ageing Report 2024](https://economy-finance.ec.europa.eu/document/download/392994f2-c025-4066-bea8-f65cf4bcd411_en?filename=2024-ageing-report-country-fiche-Spain.pdf)
- [CaixaBank Research: el Ageing Report 2024 y la reforma de pensiones](https://www.caixabankresearch.com/sites/default/files/content/file/2024/05/16/34454/im05_24-07-ee-focus-8-en.pdf)

### Referencias nuevas (introducidas en 4.4)

- Albertini, M. y Kohli, M. (2013). "The Generational Contract in the Family: An Analysis of Transfer Regimes in Europe". European Sociological Review, 29(4).
- Bonoli, G. (2000). The Politics of Pension Reform: Institutions and Policy Change in Western Europe. Cambridge University Press.
- Breyer, F. (1989). "On the Intergenerational Pareto Efficiency of Pay-as-you-go Financed Pension Systems". Journal of Institutional and Theoretical Economics, 145(4).
- Breyer, F. y Craig, B. (1997). "Voting on social security: Evidence from OECD countries". European Journal of Political Economy, 13(4).
- Bristow, J. (2015). Baby Boomers and Generational Conflict. Palgrave Macmillan.
- Browning, E. (1975). "Why the Social Insurance Budget Is Too Large in a Democracy". Economic Inquiry, 13(3). (OJO: la cita habitual es esta revista, no el Journal of Political Economy; pendiente de verificar.)
- Chetty, R., Friedman, J., Leth-Petersen, S., Nielsen, T. y Olsen, T. (2014). "Active vs. Passive Decisions and Crowd-Out in Retirement Savings Accounts: Evidence from Denmark". Quarterly Journal of Economics, 129(3).
- Conde-Ruiz, J. I. y Galasso, V. (2003). "Early Retirement". Review of Economic Dynamics, 6(1).
- Conde-Ruiz, J. I. y Galasso, V. (2005). "Positive Arithmetic of the Welfare State". Journal of Public Economics, 89(5-6).
- Cooley, T. y Soares, J. (1999). "A Positive Theory of Social Security Based on Reputation". Journal of Political Economy, 107(1).
- Disney, R. (2007). "Population Ageing and the Size of the Welfare State: Is There a Puzzle to Explain?". European Journal of Political Economy, 23(2). (Pendiente de verificar si se publicó en esta revista o en Economic Policy.)
- Galasso, V. y Profeta, P. (2002). "The Political Economy of Social Security: A Survey". European Journal of Political Economy, 18(1).
- Häusermann, S. (2010). The Politics of Welfare State Reform in Continental Europe: Modernization in Hard Times. Cambridge University Press.
- Kohli, M. (1999). "Private and Public Transfers between Generations: Linking the Family and the State". European Societies, 1(1).
- Mackenroth, G. (1952). "Die Reform der Sozialpolitik durch einen deutschen Sozialplan". Schriften des Vereins für Socialpolitik, 4.
- Madrian, B. y Shea, D. (2001). "The Power of Suggestion: Inertia in 401(k) Participation and Savings Behavior". Quarterly Journal of Economics, 116(4).
- Mastrobuoni, G. (2011). "The Role of Information for Retirement Behavior: Evidence Based on the Stepwise Introduction of the Social Security Statement". Journal of Public Economics, 95(7-8).
- Myles, J. y Pierson, P. (2001). "The Comparative Political Economy of Pension Reform", en Pierson (ed.), The New Politics of the Welfare State. Oxford University Press.
- Natali, D. (2008). Pensions in Europe, European Pensions. PIE Peter Lang.
- Sinn, H.-W. y Uebelmesser, S. (2002). "Pensions and the Path to Gerontocracy in Germany". European Journal of Political Economy, 19(1).
- Tabellini, G. (2000). "A Positive Theory of Social Security". Scandinavian Journal of Economics, 102(3).
- Thaler, R. y Benartzi, S. (2004). "Save More Tomorrow: Using Behavioral Economics to Increase Employee Saving". Journal of Political Economy, 112(S1).
- Vanhuysse, P. (2013). Intergenerational Justice in Aging Societies: A Cross-national Comparison of 29 OECD Countries. Bertelsmann Stiftung.
- Willetts, D. (2010). The Pinch: How the Baby Boomers Took Their Children's Future. Atlantic Books.
- Ley 12/2022, de 30 de junio, de regulación para el impulso de los planes de pensiones de empleo.

Pendiente de verificar:
- el estado de la revisión de las recomendaciones del Pacto de Toledo;
- el detalle de la reforma chilena de 2025;
- la disposición exacta de la Ley 27/2011 sobre información individual;
- las revistas de Browning (1975) y Disney (2007);
- el número de pensiones de viudedad y los datos de pobreza por edad de la última ECV.

Fuentes de actualidad consultadas:
- [Ley 12/2022 (BOE)](https://www.boe.es/buscar/act.php?id=BOE-A-2022-10852)
- [Ministerio de Inclusión: inscritos los 15 fondos de pensiones de empleo de promoción pública](https://www.inclusion.gob.es/en/w/seguridad-social-3)
- [Grupo Aseguranza: así cerraron 2025 los planes de pensiones (datos de Inverco)](https://www.grupoaseguranza.com/noticias-de-seguros/asi-cerraron-2025-planes-pensiones-patrimonio-rentabilidad)
- [Finect: Inverco sobre los efectos de la reforma de los planes de pensiones](https://www.finect.com/usuario/eduardogarcia/articulos/inverco-la-reforma-de-los-planes-de-pensiones-no-ha-tenido-los-efectos-deseados)
- [Congreso: Informe de evaluación y reforma del Pacto de Toledo (2020)](https://www.congreso.es/docu/comisiones/LegXIV/PactTole/Informe_Pacto_de_Toledo.pdf)
- [Demócrata: el Pacto de Toledo se lanza hacia la revisión de sus recomendaciones](https://www.democrata.es/actualidad/el-pacto-de-toledo-se-lanza-hacia-la-revision-de-sus-recomendaciones-del-sistema-de-pensiones/)

- Comisión Europea (2024). The 2024 Ageing Report: Economic and Budgetary Projections for the EU Member States (2022-2070).
- BCE (2024). "Ageing cost projections – new evidence from the 2024 Ageing Report". Economic Bulletin, 5/2024, recuadro.
- Dolowitz, D. y Marsh, D. (2000). "Learning from Abroad: The Role of Policy Transfer in Contemporary Policy-Making". Governance, 13(1).
- OCDE (2025). Pensions at a Glance 2025: OECD and G20 Indicators. OECD Publishing, noviembre de 2025.
- OCDE (2023). Pensions at a Glance 2023.
- Ross, M. H. y Homer, E. (1976). "Galton's Problem in Cross-National Research". World Politics, 29(1).
- Wet toekomst pensioenen (Países Bajos, 2023).

Pendiente de verificar:
- la ley danesa de 2025 que eleva la edad a 70 años en 2040;
- el paquete de pensiones alemán de 2025-2026 (Haltelinie y Aktivrente);
- la suspensión de la reforma francesa de 2023 y su tramitación;
- las cifras exactas de Pensions at a Glance 2025, que he obtenido de la prensa porque el informe de la OCDE no era accesible desde aquí: tasa de pobreza de los mayores, renta relativa, edades futuras y edad efectiva de salida;
- la media de gasto de la UE en el Ageing Report 2024.

Fuentes de actualidad consultadas:
- [OCDE: Pensions at a Glance 2025](https://www.oecd.org/en/publications/pensions-at-a-glance-2025_e40274c1-en.html)
- [OCDE: nota de prensa de Pensions at a Glance 2025](https://www.oecd.org/en/about/news/press-releases/2025/11/rapidly-ageing-populations-will-continue-to-put-pressure-on-pension-systems.html)
- [Vozpópuli: España, el país de la OCDE con las pensiones más generosas (80% del último salario)](https://www.vozpopuli.com/economia/espana-ya-es-el-pais-de-la-ocde-que-mas-paga-a-los-pensionistas-el-80-del-salario-al-jubilarse.html)
- [Mundiario: España será el país de la OCDE que más PIB destinará a pensiones en 2050](https://www.mundiario.com/articulo/economia/espana-sera-pais-ocde-que-mas-pib-destinara-pensiones-2050/20251127154345364780.html)
- [24h Economía: el coste de la generosidad](https://www.24heconomia.com/actualidad/economia/20251128/coste-generosidad-espana-enfrenta-mayor-reto-sostenibilidad-pensiones-historia-noticia-1221014359/)
- [BCE: proyecciones del coste del envejecimiento en el Ageing Report 2024](https://www.ecb.europa.eu/press/economic-bulletin/focus/2024/html/ecb.ebbox202405_08~5f9531042f.en.html)
- [Ficha de España del Ageing Report 2024](https://economy-finance.ec.europa.eu/document/download/392994f2-c025-4066-bea8-f65cf4bcd411_en?filename=2024-ageing-report-country-fiche-Spain.pdf)





% ANEXOS
\clearpage






\end{document}